% Options for packages loaded elsewhere
% Options for packages loaded elsewhere
\PassOptionsToPackage{unicode}{hyperref}
\PassOptionsToPackage{hyphens}{url}
\PassOptionsToPackage{dvipsnames,svgnames,x11names}{xcolor}
%
\documentclass[
  american,
  letterpaper,
  DIV=11,
  numbers=noendperiod]{scrreprt}
\usepackage{xcolor}
\usepackage{amsmath,amssymb}
\setcounter{secnumdepth}{5}
\usepackage{iftex}
\ifPDFTeX
  \usepackage[T1]{fontenc}
  \usepackage[utf8]{inputenc}
  \usepackage{textcomp} % provide euro and other symbols
\else % if luatex or xetex
  \usepackage{unicode-math} % this also loads fontspec
  \defaultfontfeatures{Scale=MatchLowercase}
  \defaultfontfeatures[\rmfamily]{Ligatures=TeX,Scale=1}
\fi
\usepackage{lmodern}
\ifPDFTeX\else
  % xetex/luatex font selection
\fi
% Use upquote if available, for straight quotes in verbatim environments
\IfFileExists{upquote.sty}{\usepackage{upquote}}{}
\IfFileExists{microtype.sty}{% use microtype if available
  \usepackage[]{microtype}
  \UseMicrotypeSet[protrusion]{basicmath} % disable protrusion for tt fonts
}{}
\makeatletter
\@ifundefined{KOMAClassName}{% if non-KOMA class
  \IfFileExists{parskip.sty}{%
    \usepackage{parskip}
  }{% else
    \setlength{\parindent}{0pt}
    \setlength{\parskip}{6pt plus 2pt minus 1pt}}
}{% if KOMA class
  \KOMAoptions{parskip=half}}
\makeatother
% Make \paragraph and \subparagraph free-standing
\makeatletter
\ifx\paragraph\undefined\else
  \let\oldparagraph\paragraph
  \renewcommand{\paragraph}{
    \@ifstar
      \xxxParagraphStar
      \xxxParagraphNoStar
  }
  \newcommand{\xxxParagraphStar}[1]{\oldparagraph*{#1}\mbox{}}
  \newcommand{\xxxParagraphNoStar}[1]{\oldparagraph{#1}\mbox{}}
\fi
\ifx\subparagraph\undefined\else
  \let\oldsubparagraph\subparagraph
  \renewcommand{\subparagraph}{
    \@ifstar
      \xxxSubParagraphStar
      \xxxSubParagraphNoStar
  }
  \newcommand{\xxxSubParagraphStar}[1]{\oldsubparagraph*{#1}\mbox{}}
  \newcommand{\xxxSubParagraphNoStar}[1]{\oldsubparagraph{#1}\mbox{}}
\fi
\makeatother

\usepackage{color}
\usepackage{fancyvrb}
\newcommand{\VerbBar}{|}
\newcommand{\VERB}{\Verb[commandchars=\\\{\}]}
\DefineVerbatimEnvironment{Highlighting}{Verbatim}{commandchars=\\\{\}}
% Add ',fontsize=\small' for more characters per line
\usepackage{framed}
\definecolor{shadecolor}{RGB}{254,254,254}
\newenvironment{Shaded}{\begin{snugshade}}{\end{snugshade}}
\newcommand{\AlertTok}[1]{\textcolor[rgb]{0.47,0.16,0.63}{#1}}
\newcommand{\AnnotationTok}[1]{\textcolor[rgb]{0.41,0.41,0.41}{#1}}
\newcommand{\AttributeTok}[1]{\textcolor[rgb]{0.65,0.35,0.00}{#1}}
\newcommand{\BaseNTok}[1]{\textcolor[rgb]{0.47,0.16,0.63}{#1}}
\newcommand{\BuiltInTok}[1]{\textcolor[rgb]{0.33,0.33,0.33}{#1}}
\newcommand{\CharTok}[1]{\textcolor[rgb]{0.00,0.50,0.00}{#1}}
\newcommand{\CommentTok}[1]{\textcolor[rgb]{0.41,0.41,0.41}{#1}}
\newcommand{\CommentVarTok}[1]{\textcolor[rgb]{0.41,0.41,0.41}{\textit{#1}}}
\newcommand{\ConstantTok}[1]{\textcolor[rgb]{0.85,0.12,0.09}{#1}}
\newcommand{\ControlFlowTok}[1]{\textcolor[rgb]{0.85,0.12,0.09}{#1}}
\newcommand{\DataTypeTok}[1]{\textcolor[rgb]{0.47,0.16,0.63}{#1}}
\newcommand{\DecValTok}[1]{\textcolor[rgb]{0.47,0.16,0.63}{#1}}
\newcommand{\DocumentationTok}[1]{\textcolor[rgb]{0.41,0.41,0.41}{\textit{#1}}}
\newcommand{\ErrorTok}[1]{\textcolor[rgb]{0.47,0.16,0.63}{#1}}
\newcommand{\ExtensionTok}[1]{\textcolor[rgb]{0.33,0.33,0.33}{#1}}
\newcommand{\FloatTok}[1]{\textcolor[rgb]{0.65,0.35,0.00}{#1}}
\newcommand{\FunctionTok}[1]{\textcolor[rgb]{0.02,0.16,0.49}{#1}}
\newcommand{\ImportTok}[1]{\textcolor[rgb]{0.33,0.33,0.33}{#1}}
\newcommand{\InformationTok}[1]{\textcolor[rgb]{0.41,0.41,0.41}{#1}}
\newcommand{\KeywordTok}[1]{\textcolor[rgb]{0.85,0.12,0.09}{#1}}
\newcommand{\NormalTok}[1]{\textcolor[rgb]{0.33,0.33,0.33}{#1}}
\newcommand{\OperatorTok}[1]{\textcolor[rgb]{0.00,0.46,0.62}{#1}}
\newcommand{\OtherTok}[1]{\textcolor[rgb]{0.85,0.12,0.09}{#1}}
\newcommand{\PreprocessorTok}[1]{\textcolor[rgb]{0.47,0.16,0.63}{#1}}
\newcommand{\RegionMarkerTok}[1]{\textcolor[rgb]{0.33,0.33,0.33}{#1}}
\newcommand{\SpecialCharTok}[1]{\textcolor[rgb]{0.00,0.46,0.62}{#1}}
\newcommand{\SpecialStringTok}[1]{\textcolor[rgb]{0.00,0.50,0.00}{#1}}
\newcommand{\StringTok}[1]{\textcolor[rgb]{0.00,0.50,0.00}{#1}}
\newcommand{\VariableTok}[1]{\textcolor[rgb]{0.65,0.35,0.00}{#1}}
\newcommand{\VerbatimStringTok}[1]{\textcolor[rgb]{0.00,0.50,0.00}{#1}}
\newcommand{\WarningTok}[1]{\textcolor[rgb]{0.41,0.41,0.41}{\textit{#1}}}

\usepackage{longtable,booktabs,array}
\usepackage{calc} % for calculating minipage widths
% Correct order of tables after \paragraph or \subparagraph
\usepackage{etoolbox}
\makeatletter
\patchcmd\longtable{\par}{\if@noskipsec\mbox{}\fi\par}{}{}
\makeatother
% Allow footnotes in longtable head/foot
\IfFileExists{footnotehyper.sty}{\usepackage{footnotehyper}}{\usepackage{footnote}}
\makesavenoteenv{longtable}
\usepackage{graphicx}
\makeatletter
\newsavebox\pandoc@box
\newcommand*\pandocbounded[1]{% scales image to fit in text height/width
  \sbox\pandoc@box{#1}%
  \Gscale@div\@tempa{\textheight}{\dimexpr\ht\pandoc@box+\dp\pandoc@box\relax}%
  \Gscale@div\@tempb{\linewidth}{\wd\pandoc@box}%
  \ifdim\@tempb\p@<\@tempa\p@\let\@tempa\@tempb\fi% select the smaller of both
  \ifdim\@tempa\p@<\p@\scalebox{\@tempa}{\usebox\pandoc@box}%
  \else\usebox{\pandoc@box}%
  \fi%
}
% Set default figure placement to htbp
\def\fps@figure{htbp}
\makeatother


% definitions for citeproc citations
\NewDocumentCommand\citeproctext{}{}
\NewDocumentCommand\citeproc{mm}{%
  \begingroup\def\citeproctext{#2}\cite{#1}\endgroup}
\makeatletter
 % allow citations to break across lines
 \let\@cite@ofmt\@firstofone
 % avoid brackets around text for \cite:
 \def\@biblabel#1{}
 \def\@cite#1#2{{#1\if@tempswa , #2\fi}}
\makeatother
\newlength{\cslhangindent}
\setlength{\cslhangindent}{1.5em}
\newlength{\csllabelwidth}
\setlength{\csllabelwidth}{3em}
\newenvironment{CSLReferences}[2] % #1 hanging-indent, #2 entry-spacing
 {\begin{list}{}{%
  \setlength{\itemindent}{0pt}
  \setlength{\leftmargin}{0pt}
  \setlength{\parsep}{0pt}
  % turn on hanging indent if param 1 is 1
  \ifodd #1
   \setlength{\leftmargin}{\cslhangindent}
   \setlength{\itemindent}{-1\cslhangindent}
  \fi
  % set entry spacing
  \setlength{\itemsep}{#2\baselineskip}}}
 {\end{list}}
\usepackage{calc}
\newcommand{\CSLBlock}[1]{\hfill\break\parbox[t]{\linewidth}{\strut\ignorespaces#1\strut}}
\newcommand{\CSLLeftMargin}[1]{\parbox[t]{\csllabelwidth}{\strut#1\strut}}
\newcommand{\CSLRightInline}[1]{\parbox[t]{\linewidth - \csllabelwidth}{\strut#1\strut}}
\newcommand{\CSLIndent}[1]{\hspace{\cslhangindent}#1}

\ifLuaTeX
\usepackage[bidi=basic]{babel}
\else
\usepackage[bidi=default]{babel}
\fi
% get rid of language-specific shorthands (see #6817):
\let\LanguageShortHands\languageshorthands
\def\languageshorthands#1{}
\ifLuaTeX
  \usepackage[english]{selnolig} % disable illegal ligatures
\fi


\setlength{\emergencystretch}{3em} % prevent overfull lines

\providecommand{\tightlist}{%
  \setlength{\itemsep}{0pt}\setlength{\parskip}{0pt}}



 


\KOMAoption{captions}{tableheading}
\makeatletter
\@ifpackageloaded{tcolorbox}{}{\usepackage[skins,breakable]{tcolorbox}}
\@ifpackageloaded{fontawesome5}{}{\usepackage{fontawesome5}}
\definecolor{quarto-callout-color}{HTML}{909090}
\definecolor{quarto-callout-note-color}{HTML}{0758E5}
\definecolor{quarto-callout-important-color}{HTML}{CC1914}
\definecolor{quarto-callout-warning-color}{HTML}{EB9113}
\definecolor{quarto-callout-tip-color}{HTML}{00A047}
\definecolor{quarto-callout-caution-color}{HTML}{FC5300}
\definecolor{quarto-callout-color-frame}{HTML}{acacac}
\definecolor{quarto-callout-note-color-frame}{HTML}{4582ec}
\definecolor{quarto-callout-important-color-frame}{HTML}{d9534f}
\definecolor{quarto-callout-warning-color-frame}{HTML}{f0ad4e}
\definecolor{quarto-callout-tip-color-frame}{HTML}{02b875}
\definecolor{quarto-callout-caution-color-frame}{HTML}{fd7e14}
\makeatother
\makeatletter
\@ifpackageloaded{caption}{}{\usepackage{caption}}
\AtBeginDocument{%
\ifdefined\contentsname
  \renewcommand*\contentsname{Table of contents}
\else
  \newcommand\contentsname{Table of contents}
\fi
\ifdefined\listfigurename
  \renewcommand*\listfigurename{List of Figures}
\else
  \newcommand\listfigurename{List of Figures}
\fi
\ifdefined\listtablename
  \renewcommand*\listtablename{List of Tables}
\else
  \newcommand\listtablename{List of Tables}
\fi
\ifdefined\figurename
  \renewcommand*\figurename{Figure}
\else
  \newcommand\figurename{Figure}
\fi
\ifdefined\tablename
  \renewcommand*\tablename{Table}
\else
  \newcommand\tablename{Table}
\fi
}
\@ifpackageloaded{float}{}{\usepackage{float}}
\floatstyle{ruled}
\@ifundefined{c@chapter}{\newfloat{codelisting}{h}{lop}}{\newfloat{codelisting}{h}{lop}[chapter]}
\floatname{codelisting}{Listing}
\newcommand*\listoflistings{\listof{codelisting}{List of Listings}}
\makeatother
\makeatletter
\makeatother
\makeatletter
\@ifpackageloaded{caption}{}{\usepackage{caption}}
\@ifpackageloaded{subcaption}{}{\usepackage{subcaption}}
\makeatother
\newcounter{quartocalloutnteno}
\newcommand{\quartocalloutnte}[1]{\refstepcounter{quartocalloutnteno}\label{#1}}
\usepackage{bookmark}
\IfFileExists{xurl.sty}{\usepackage{xurl}}{} % add URL line breaks if available
\urlstyle{same}
\hypersetup{
  pdftitle={Reading Genomes: A Companion to Bioinformatics and Genome Analysis},
  pdfauthor={Laurie Stevison},
  pdflang={en-US},
  colorlinks=true,
  linkcolor={blue},
  filecolor={Maroon},
  citecolor={Blue},
  urlcolor={Blue},
  pdfcreator={LaTeX via pandoc}}


\title{Reading Genomes: A Companion to Bioinformatics and Genome
Analysis}
\author{Laurie Stevison}
\date{}
\begin{document}
\maketitle

<a href="#quarto-content" class="skip-nav">Skip to main content</a>
<style>
.skip-nav {
  position: absolute;
  left: -10000px;
  top: auto;
  width: 1px;
  height: 1px;
  overflow: hidden;
}
.skip-nav:focus {
  position: fixed;
  top: 0;
  left: 0;
  width: auto;
  height: auto;
  padding: 8px;
  background: #000;
  color: #fff;
  text-decoration: none;
  z-index: 9999;
}
</style>

\renewcommand*\contentsname{Table of Contents}
{
\hypersetup{linkcolor=}
\setcounter{tocdepth}{2}
\tableofcontents
}

\chapter*{Preface}\label{preface}
\addcontentsline{toc}{chapter}{Preface}

\markboth{Preface}{Preface}

This is a Quarto book. This companion textbook is a draft teaching
resource created specifically for students enrolled in BIOL 5860/6860:
Bioinformatics and Genome Analysis at Auburn University during the
Spring 2026 term. \emph{It is not intended for general public use,
redistribution, or citation as a finalized scholarly work.}

The content integrates instructor-written material, published
references, and AI-assisted drafting that has been reviewed and edited
for use in this course, with sources and AI assistance acknowledged.
Access to this site may be temporarily open while technical settings are
configured, but it is intended to be password-protected and restricted
to enrolled students as soon as feasible.

To learn more about Quarto books visit
\url{https://quarto.org/docs/books}.

\section*{Why This Course? Why This
Book?}\label{why-this-course-why-this-book}
\addcontentsline{toc}{section}{Why This Course? Why This Book?}

\markright{Why This Course? Why This Book?}

This companion textbook supports \emph{BIOL 5860/6860: Bioinformatics
and Genome Analysis}, a hands-on survey of computational approaches to
``reading'' biological data. As biologists with genetics and
computational biology prerequisites, you already know how to frame
biological questions and wrangle basic code. This course builds directly
on that foundation to demystify the ``black box'' of bioinformatics
tools, showing you how to move from raw sequencing data to biological
insight. We will also read primary scientific literature, which you are
expected to be familiar with. If you need a refresher, please refer to
\href{appendixA.qmd}{Appendix A}.

The field moves fast---back when I was in graduate school, I first
learned 454 sequencing and applied it to my PhD work, which is now long
obsolete. Rather than teaching specific tools that may change by
graduation, this book emphasizes \textbf{adaptability}: the core
concepts, file formats, quality control practices, and workflow thinking
that transfer across platforms and projects. Labs progress from
exploratory (``find a genome paper on your favorite organism'') to
practical (GATK variant calling on HPC) to independent (semester-long
group research projects where you discover and master new tools
yourselves).

\section*{Who This Book Serves}\label{who-this-book-serves}
\addcontentsline{toc}{section}{Who This Book Serves}

\markright{Who This Book Serves}

This resource targets our diverse class:

\begin{itemize}
\item
  \textbf{Pre-med and future genetic counselors}: Clinical relevance
  through variant interpretation, population genetics, and ethical
  considerations in human genomics.
\item
  \textbf{Graduate students and researchers}: Practical pipelines,
  reproducibility best practices, and project management skills for real
  analyses.
\item
  \textbf{All students}: Demystification of the ``hidden
  curriculum''---writing reproducible workflows, applying for funding,
  and participating in grant review panels.
\end{itemize}

While a pre-requisite of Genetics is required, no prior genomics
experience is required. Similarly, the prerequisite of BIOL 5800/6800
ensures that everyone can keep pace with Linux, R, git, and basic
scripting.

\section*{How Does Computational Biology Differ from
Bioinformatics}\label{how-does-computational-biology-differ-from-bioinformatics}
\addcontentsline{toc}{section}{How Does Computational Biology Differ
from Bioinformatics}

\markright{How Does Computational Biology Differ from Bioinformatics}

\begin{tcolorbox}[enhanced jigsaw, rightrule=.15mm, leftrule=.75mm, colbacktitle=quarto-callout-note-color!10!white, colframe=quarto-callout-note-color-frame, toprule=.15mm, opacitybacktitle=0.6, title=\textcolor{quarto-callout-note-color}{\faInfo}\hspace{0.5em}{Figure 0.1 -- Fields of Research within Bioinformatics}, left=2mm, colback=white, bottomtitle=1mm, toptitle=1mm, bottomrule=.15mm, coltitle=black, titlerule=0mm, breakable, opacityback=0, arc=.35mm]

\begin{figure}[H]

\centering{

\includegraphics[width=0.5\linewidth,height=\textheight,keepaspectratio]{images/Fig_0_1_Fields.png}

}

\caption{\label{fig-fields}Venn Diagram with fields of study in the area
of bioinformatics}

\end{figure}%

Having taken a course in Computational Biology, it can often be
confusing of the specific distinction between the various fields and how
to define each.

This figure shows how bioinformatics lies at the interface of
statistics, computer science, and biology, integrating methods from all
three disciplines to analyze and interpret biological data. This Figure
highlights how related subfields overlap: biostatistics emerges from the
intersection of statistics and biology, computational biology from the
overlap of computer science and biology, and data science from the
overlap of statistics and computer science. At the center,
bioinformatics combines these quantitative, computational, and
biological approaches under one umbrella.

Importantly, how we define biology in this course is broadly
emcompassing all fields of empirical research. This would include
biology, biomedical research, such as pharamcology, kinesiology, and
veterinary sciences.

\end{tcolorbox}

\section*{Student Learning
Objectives}\label{student-learning-objectives}
\addcontentsline{toc}{section}{Student Learning Objectives}

\markright{Student Learning Objectives}

Upon completion of the course, you should be able to:

\begin{itemize}
\tightlist
\item
  Interpret data formats and quality of:

  \begin{itemize}
  \tightlist
  \item
    raw genome sequence data
  \item
    sequence alignments
  \item
    genetic variants
  \end{itemize}
\item
  Conduct genomic analysis using best practice pipelines
\item
  Analyze data using R and various bioinformatics tools
\item
  Critically evaluate genome analysis study designs
\item
  Be able to describe in depth the basic analysis for a variety of data
  types
\item
  Develop proficiency in scientific communication skills and
  reproducibility of research
\end{itemize}

\section*{Labs and Projects as a parallel
track}\label{labs-and-projects-as-a-parallel-track}
\addcontentsline{toc}{section}{Labs and Projects as a parallel track}

\markright{Labs and Projects as a parallel track}

Throughout the semester, you will encounter two recurring threads:

\begin{itemize}
\tightlist
\item
  A \textbf{comparative genomics project} (\href{appendixB.qmd}{Appendix
  B}) that applies the skills you are learning to a real research
  question.
\item
  A \textbf{yeast genomics dataset} (Labs 4, 7, 9;
  \href{appendixC.qmd}{Appendix C}) that takes you from FASTQ → BAM →
  VCF → IGV.
\end{itemize}

Each chapter includes ``Lab connections'' callouts that point to
relevant labs and appendices, so you can cross‑reference concepts,
commands, and datasets as you work.

\section*{Semester Schedule}\label{semester-schedule}
\addcontentsline{toc}{section}{Semester Schedule}

\markright{Semester Schedule}

\begin{longtable}[]{@{}
  >{\raggedright\arraybackslash}p{(\linewidth - 10\tabcolsep) * \real{0.1667}}
  >{\raggedright\arraybackslash}p{(\linewidth - 10\tabcolsep) * \real{0.1667}}
  >{\raggedright\arraybackslash}p{(\linewidth - 10\tabcolsep) * \real{0.1667}}
  >{\raggedright\arraybackslash}p{(\linewidth - 10\tabcolsep) * \real{0.1667}}
  >{\raggedright\arraybackslash}p{(\linewidth - 10\tabcolsep) * \real{0.1667}}
  >{\raggedright\arraybackslash}p{(\linewidth - 10\tabcolsep) * \real{0.1667}}@{}}
\toprule\noalign{}
\begin{minipage}[b]{\linewidth}\raggedright
Week
\end{minipage} & \begin{minipage}[b]{\linewidth}\raggedright
Date
\end{minipage} & \begin{minipage}[b]{\linewidth}\raggedright
Class topic
\end{minipage} & \begin{minipage}[b]{\linewidth}\raggedright
Primary readings
\end{minipage} & \begin{minipage}[b]{\linewidth}\raggedright
Description
\end{minipage} & \begin{minipage}[b]{\linewidth}\raggedright
Assessments
\end{minipage} \\
\midrule\noalign{}
\endhead
\bottomrule\noalign{}
\endlastfoot
1 & 8 Jan & Course overview and state of the field &
\href{ch01-reading-genomes.qmd}{Chapter 1};
Readings:\textsuperscript{1,2}; \href{appendixA.qmd}{Appendix A} &
Big‑picture framing of course. & Start \href{Lab01.qmd}{Lab 1} \\
2 & 13 Jan & What is a genome analysis? &
\href{ch02-sequencing-tech.qmd}{Chapter 2}; Reading:\textsuperscript{3}
& Sequencing platforms & Start \href{Lab02.qmd}{Lab 2} \\
& 15 Jan & Intro to NGS data and data QC &
\href{ch03-formats-qc-design.qmd}{Chapter 3} & Genomics File Formats
& \\
3 & 20 Jan & {Intro to semester‑long research project} &
\href{appendixB.qmd}{Appendix B} & Project framing and expectations.
& \\
& 22-Jan & {Lab day -- Genome browsers } &
\href{ch04-browsers-databases.qmd}{Chapter 4} & Explore different
browsers and customizations & \href{Lab03.qmd}{Lab 3} \\
4 & 27-Jan & Pairwise \& multiple sequence alignment &
\href{ch05-seq-alignment.qmd}{Chapter 5} & Intro to pairwise/MSA
concepts & \\
& 29-Jan & Genome sequence alignment &
\href{ch06-genome-align-assembly.qmd}{Chapter 6} & Scaling alignment
concepts and algorithms & \\
5 & 3-Feb & Project management \& writing methods in bioinformatics &
\href{ch07-project-mgmt.qmd}{Chapter 7} & Hidden curriculum emphasis &
{GP Step 1 Due } \\
& 5-Feb & {Lab day -- Indexing, alignment, and assessment} &
\href{appendixC.qmd}{Appendix C} & Introduce yeast dataset and practical
application & \href{Lab04.qmd}{Lab 4} \\
6 & 10-Feb & Genome assembly guest lecture &
\href{ch06-genome-align-assembly.qmd}{Chapter 6} & Algorithms and case
studies & \\
& 12-Feb & {Lab day -- BLAST on the command line } &
\href{ch05-seq-alignment.qmd}{Chapter 5} & BLAST concepts and practice
on the HPC & \href{Lab05.qmd}{Lab 5} \\
7 & 17-Feb & Various genome analysis workflows &
\href{ch08-alt-workflows.qmd}{Chapter 8} & Multi-omics workflows and
human case studies & {Meet with Dr.~S} \\
& 19-Feb & {Lab day -- HMMs for gene annotation} &
\href{ch09-hmms.qmd}{Chapter 9}; Reading:\textsuperscript{4} & HMM
intuition and gene‑finding & \href{Lab06.qmd}{Lab 6} \\
8 & 24-Feb & Intro to variant calling/filtering &
\href{ch10-variant-calling.qmd}{Chapter 10} & From alignment to variants
using GATK & \\
& 26-Feb & {\textbf{Midterm Exam}} & In class conceptual exam &
Assessment of Chapters 1-7; Labs 1-5; Appendices A-C & Exam \\
9 & 3-Mar & {Lab day -- Variant filtering and QC } &
\href{ch10-variant-calling.qmd}{Chapter 10};
\href{appendixC.qmd}{Appendix C} & Practical filtering and
interpretation & \href{Lab07.qmd}{Lab 7}; {GP Step 2 Due} \\
& 5-Mar & {Beyond GATK: LLMs and other Variant Calling Approaches} &
\href{ch10-variant-calling.qmd}{Chapter 10}; Assigned
Readings:\textsuperscript{5--9} & In depth comparison of variant calling
methods & \href{Lab08.qmd}{Lab 8}; Annotation Report Due \\
10 & 10-Mar & Spring break -- no class & --- & No assigned reading & \\
& 12-Mar & Spring break -- no class & --- & No assigned reading & \\
11 & 17-Mar & {Open lab -- group project work} &
\href{appendixB.qmd}{Appendix B} & Execute Bioinformatics Plan and Setup
GitHub & {GP Step 3}; Grad Assignment \\
& 19-Mar & {Lab Day -- Visualizing Genetic Variants} &
\href{ch10-variant-calling.qmd}{Chapter 10};
\href{appendixC.qmd}{Appendix C} & Visualizing variants in genome
viewers & \href{Lab09.qmd}{Lab 9} \\
12 & 24-Mar & How Science is Funded? -- Grant Funding \& Review Process
& \href{ch11-grants-funding.qmd}{Chapter 11} & Hidden curriculum
emphasis & \\
& 26-Mar & Population genomics & \href{ch12-popgen.qmd}{Chapter 12} &
Human evolutionary genomics & {GP Step 4 Due } \\
13 & 31-Mar & {Lab day -- Human Population Genetics } &
\href{ch12-popgen.qmd}{Chapter 12} & Hands‑on diversity and divergence
analysis & \href{Lab10.qmd}{Lab 10} \\
& 2-Apr & Genome scans \& sliding‑window analysis &
\href{ch12-popgen.qmd}{Chapter 12} & Outlier scans, patterns of
selection, and case studies & {Proposal Reviews Due} \\
14 & 7-Apr & {Grant panel day 1 } &
\href{ch11-grants-funding.qmd}{Chapter 11} & In‑class panel and critique
& {Mock Grant Panel} \\
& 9-Apr & {Grant panel day 2 } & \href{ch11-grants-funding.qmd}{Chapter
11} & In‑class panel and critique & {Mock Grant Panel} \\
15 & 14-Apr & Sequence motifs & \href{ch13-motifs.qmd}{Chapter 13} &
Finding sequence motifs in genomes & {Panel Summary Due } \\
& 16-Apr & {Open lab -- group project work} &
\href{appendixB.qmd}{Appendix B}; \href{ch14-where-next.qmd}{Chapter 14}
& Project‑oriented application of course content & \\
16 & 21-Apr & {\textbf{Final Exam} } & In class conceptual exam & Review
of Chapters 8-14; Labs 6-10 & Exam \\
& 23-Apr & {Open lab -- group project work} &
\href{appendixB.qmd}{Appendix B}; \href{ch14-where-next.qmd}{Chapter 14}
& Project‑oriented application of course content & {GP Step 5 Due } \\
17 & 1-May & {Final group presentations (10:30--12:30)} &
\href{appendixB.qmd}{Appendix B} & Capstone, reflections, and next steps
& {GP Step 6-7 Due } \\
\end{longtable}

\section*{Key Dates}\label{key-dates}
\addcontentsline{toc}{section}{Key Dates}

\markright{Key Dates}

\begin{longtable}[]{@{}
  >{\raggedright\arraybackslash}p{(\linewidth - 8\tabcolsep) * \real{0.2000}}
  >{\raggedright\arraybackslash}p{(\linewidth - 8\tabcolsep) * \real{0.2000}}
  >{\raggedright\arraybackslash}p{(\linewidth - 8\tabcolsep) * \real{0.2000}}
  >{\raggedright\arraybackslash}p{(\linewidth - 8\tabcolsep) * \real{0.2000}}
  >{\raggedright\arraybackslash}p{(\linewidth - 8\tabcolsep) * \real{0.2000}}@{}}
\toprule\noalign{}
\begin{minipage}[b]{\linewidth}\raggedright
Assignment
\end{minipage} & \begin{minipage}[b]{\linewidth}\raggedright
Description
\end{minipage} & \begin{minipage}[b]{\linewidth}\raggedright
Due Date
\end{minipage} & \begin{minipage}[b]{\linewidth}\raggedright
Points
\end{minipage} & \begin{minipage}[b]{\linewidth}\raggedright
Percent course completed
\end{minipage} \\
\midrule\noalign{}
\endhead
\bottomrule\noalign{}
\endlastfoot
Lab 1 & Find and describe a genome analysis paper & 15-Jan & 25 & 3\% \\
Lab 2 & Design a Genome Project & 23-Jan & 25 & 6\% \\
Lab 3 & Genome Browsers & 29-Jan & 25 & 9\% \\
GP Step 1 & Project Overview & 3-Feb & 20 & 12\% \\
Lab 4 & Genome Alignment/Indexing & 12-Feb & 25 & 15\% \\
Lab 5 & BLAST on the command line & 19-Feb & 25 & 18\% \\
GP Meeting & Meet with Dr.~Stevison to discuss research plan & 20-Feb &
NA & 18\% \\
Lab 6 & HMMs for gene annotation & 26-Feb & 25 & 21\% \\
Midterm & Exam on Chapters 1-7 & 26-Feb & 100 & 33\% \\
GP Step 2 & Detailed Bioinformatics Plan & 2-Mar & 30 & 36\% \\
Annotation Project & Manual Gene Annotation Report & 6-Mar & 100 &
48\% \\
Lab 7 & Variant Calling and Filtering & 10-Mar & 25 & 52\% \\
Lab 8 & Variant Calling Discussion & 12-Mar & 25 & 55\% \\
Grad Assignment & Guide for a bioinformatics tool! & 16-Mar & NA &
55\% \\
GP Step 3 & GitHub Repo with Prelim analysis & 23-Mar & 10 & 56\% \\
Lab 9 & IGV and Genome Viewers & 26-Mar & 25 & 59\% \\
GP Step 4 & Peer Review of Step 3 & 27-Mar & 20 & 61\% \\
Panel Review & Mock Panel Reviews (due BEFORE panel) & 3-Apr & 20 &
64\% \\
Lab 10 & FST, Tajima's D and Diversity Scans & 7-Apr & 25 & 67\% \\
Panel Discussion & Panel Discussion and Participation IN CLASS & 7-Apr &
30 & 70\% \\
Panel Summary & Detailed Summary of Panel Discussion & 14-Apr & 25 &
73\% \\
Final Exam & Exam on Chapters 8-14 & 21-Apr & 100 & 85\% \\
GP Step 5 & Final GitHub Repository & 24-Apr & 50 & 92\% \\
GP Step 6 & Final Presentation & 1-May & 50 & 98\% \\
GP Step 7 & Peer Review of Step 6 & 1-May & 20 & 100\% \\
\end{longtable}

*Note: The above schedule and these deadlines are subject to change.

\section*{Final Note to Students}\label{final-note-to-students}
\addcontentsline{toc}{section}{Final Note to Students}

\markright{Final Note to Students}

Bioinformatics feels like magic until you see the patterns: every
pipeline starts with QC, every analysis needs reproducibility, every
grant needs a clear workflow. By semester's end, you'll not only run
these analyses but build the confidence and comfort to explain them to
collaborators, reviewers, and future employers.

Let's read some genomes together.

\emph{Dr.~Laurie Stevison - Your Steward in Genomics} \emph{December
2025}

\section*{Acknowledgments}\label{acknowledgments}
\addcontentsline{toc}{section}{Acknowledgments}

\markright{Acknowledgments}

Portions of this book were drafted and revised with the assistance of an
AI‑based writing tool. The instructor reviewed, edited, and is
responsible for the final content.

\begin{center}\rule{0.5\linewidth}{0.5pt}\end{center}

\section*{References}\label{references}
\addcontentsline{toc}{section}{References}

\markright{References}

\part{Companion Chapters}

\chapter{Reading Genomes: What Bioinformatics Is
For}\label{reading-genomes-what-bioinformatics-is-for}

\section{Genomics timeline: from Sanger to the data
deluge}\label{sec-ch1}

Modern genome analysis traces back to the development of DNA sequencing
methods in the 1970s, which made it possible to determine the nucleotide
sequence of individual genes and small genomes for the first
time\textsuperscript{10};\textsuperscript{11}. Over subsequent decades,
improvements in chemistry and automation enabled large‑scale projects
such as the Human Genome Project, which produced the first high‑quality
human reference genome and marked the beginning of the ``genomic
era''\textsuperscript{10};\textsuperscript{11}.

In the roughly twenty‑five years since the first complete animal and
plant genomes were published, sequencing technologies have advanced from
low‑throughput Sanger methods to high‑throughput short‑read platforms
and, more recently, long‑read technologies that generate much longer
sequences at scale\textsuperscript{1};\textsuperscript{2}. These
advances have produced thousands of animal and hundreds of land‑plant
genome assemblies, but current references still cover only a small
fraction of described species and are unevenly distributed across the
tree of life\textsuperscript{1};\textsuperscript{2}.

\begin{tcolorbox}[enhanced jigsaw, rightrule=.15mm, leftrule=.75mm, colbacktitle=quarto-callout-note-color!10!white, colframe=quarto-callout-note-color-frame, toprule=.15mm, opacitybacktitle=0.6, title=\textcolor{quarto-callout-note-color}{\faInfo}\hspace{0.5em}{Figure 1.1 -- Milestones in genomics}, left=2mm, colback=white, bottomtitle=1mm, toptitle=1mm, bottomrule=.15mm, coltitle=black, titlerule=0mm, breakable, opacityback=0, arc=.35mm]

\begin{figure}[H]

\centering{

\includegraphics[width=0.8\linewidth,height=\textheight,keepaspectratio]{images/Fig_1_1_genomics_timeline.png}

}

\caption{\label{fig-genomics-timeline}Milestones in genomics}

\end{figure}%

This figure presents a timeline of key technological innovations and
landmark projects that shaped modern genomics, from the recognition of
DNA as the hereditary material and the description of its double‑helix
structure through the development of PCR, first‑generation Sanger
sequencing, and early automated sequencers. It then highlights the
launch and 2001 completion of the Human Genome Project and related
large‑scale efforts, such as the generation of reference genomes for
multiple organisms and subsequent population‑scale initiatives like the
1000 Genomes Project and the Human Microbiome Project, which
collectively established comprehensive genomic resources for research
and medicine. The upper portion of the figure tracks the transition from
low‑throughput first‑generation methods to second‑generation short‑read
platforms and, more recently, to third‑generation long‑read technologies
such as single‑molecule real‑time and nanopore sequencing, which now
enable routine production of near‑complete, reference‑quality
assemblies, including end‑to‑end human genomes.

This figure visually anchors the historical narrative for later chapters
where we will go into these technologies and research milestones in more
depth throughout this course. Created in https://BioRender.com.

\end{tcolorbox}

\section{Core concepts: genomes, omes, and basic
operations}\label{core-concepts-genomes-omes-and-basic-operations}

For this course, the \textbf{genome} is the complete DNA sequence of an
organism, typically organized into chromosomes, and including both
coding and noncoding regions\textsuperscript{10,12}. The
\textbf{transcriptome} is the collection of RNA molecules transcribed
from the genome under specific conditions, and the \textbf{proteome} is
the set of proteins produced from those transcripts; together they
connect static DNA sequence to dynamic cellular
phenotypes\textsuperscript{10,12}.

Beyond sequence alone, the \textbf{epigenome} captures chemical and
structural modifications (such as DNA methylation and chromatin state)
that affect how accessible different genomic regions are without
changing the underlying bases\textsuperscript{13}. Modern genome
analysis often integrates these layers to understand how genomic
sequence, gene expression, and epigenetic regulation combine to produce
organism‑level traits\textsuperscript{10,13}.

A \textbf{reference genome} is a community standard assembly for a
species (or strain) that serves as a coordinate system for mapping
reads, annotating genes, and comparing individuals\textsuperscript{1,2}.
A \textbf{genome assembly} is the specific reconstruction of a genome
from sequencing reads for a given project, usually summarized by its
size, contiguity (for example, N50), and completeness; an
\textbf{annotation} is the set of predicted and curated genomic features
such as genes, transcripts, and regulatory elements placed on that
assembly\textsuperscript{1,2}.

Two core computational operations recur throughout this course.
\textbf{Alignment} is a process that refers to the maping of reads or
sequences to a reference genome to determine their physical location in
the genome. This step is key in enabling variant calling, coverage
analysis, and many downstream inferences\textsuperscript{10}. \textbf{De
novo assembly} instead stitches reads together based on sequence overlap
or graph structure to reconstruct longer contiguous sequences without
relying on a close reference, which is essential for new or highly
divergent genomes and for resolving structural
complexity\textsuperscript{1}. Both processes are necessary to expand
our understanding of genomes and require various bioinformatics tools.

\begin{tcolorbox}[enhanced jigsaw, rightrule=.15mm, leftrule=.75mm, colbacktitle=quarto-callout-note-color!10!white, colframe=quarto-callout-note-color-frame, toprule=.15mm, opacitybacktitle=0.6, title=\textcolor{quarto-callout-note-color}{\faInfo}\hspace{0.5em}{Figure 1.2 -- Core genome concepts}, left=2mm, colback=white, bottomtitle=1mm, toptitle=1mm, bottomrule=.15mm, coltitle=black, titlerule=0mm, breakable, opacityback=0, arc=.35mm]

\begin{figure}[H]

\centering{

\includegraphics[width=0.75\linewidth,height=\textheight,keepaspectratio]{images/Figure 1.2 Core Concepts.png}

}

\caption{\label{fig-concepts}Core Concepts}

\end{figure}%

Typical eukaryotic gene model illustrating exons and introns along a
genomic region (A). Schematic of core ``omes'' for a single cell: DNA as
the genome, transcribed RNA molecules as the transcriptome, translated
polypeptides as the proteome, and chromatin with chemical modifications
representing the epigenome (B). Reference-based alignment, showing
sequencing reads stacked over a reference genome coordinate axis (C).
Genome assembly, in which clusters of overlapping reads form contigs
that are joined into an assembled scaffold (D). Created with
BioRender.com.

\end{tcolorbox}

\section{Why bioinformatics exists}\label{why-bioinformatics-exists}

As sequencing costs have fallen, generating genome‑scale data has become
easier and cheaper than analyzing, storing, and interpreting it,
creating a \textbf{bioinformatic bottleneck} in many
projects\textsuperscript{11,13} (see Figure~\ref{fig-bottleneck}). The
price of whole‑genome sequencing has dropped orders of magnitude since
the Human Genome Project, but high‑quality analysis still requires
computational expertise, careful data management, and appropriate
statistical methods\textsuperscript{13}. At the same time, public
repositories, such as GenBank, have grown to tens of trillions of bases
and billions of sequences, so simply storing and querying genomic data
has become a large‑scale data‑management challenge (genomic databases
will be revisited in Chapter 4).

This data deluge is often compared to Moore's law, which describes a
long-term trend in the computer hardware industry that involves the
doubling of `compute power' every two years. Technology improvements
that `keep up' with Moore's Law are widely regarded to be doing
exceedingly well, making it useful for comparison. Sequencing prices
have fallen by more than six orders of magnitude since the Human Genome
Project (see Figure~\ref{fig-genomics-timeline}), easily outpacing these
expected trends in improvements to general computing hardware
(\href{https://www.genome.gov/about-genomics/fact-sheets/DNA-Sequencing-Costs-Data}{NHGRI
DNA Sequencing Costs}).

At the same time, concerns about a \textbf{reproducibility crisis} in
biomedical and computational research have highlighted the need for
transparent, well‑documented workflows\textsuperscript{10}. In genomics,
this includes clearly recording data sources and metadata, tracking
software versions and parameters, using version control and scripted
pipelines where possible, and depositing both data and code in
accessible repositories so that others can repeat and extend
analyses\textsuperscript{1,2,10}.

The software and methods landscape in bioinformatics changes rapidly:
new alignment algorithms, variant callers, assemblers, and visualization
tools appear frequently, while older tools may become unsupported or
superseded\textsuperscript{10,12}. A central goal of this course is
therefore not mastery of any single program but
\textbf{adaptability}---the ability to read documentation, reason about
how a tool fits into a workflow, evaluate whether its assumptions match
a given dataset, and switch approaches as technologies and best
practices evolve\textsuperscript{12,13}. This rapid pace is also the
reason we do not have a standardized textbook for this course. Instead,
I have compiled various resources into this single course companion.
Much of this textbook relies heavily on an excellent textbook last
updated in 2015\textsuperscript{10}. You will see this text cited
throughout and some of the content comes directly from this text. But
much of that is outdated and has been significantly supplemented with
primary scientific literature as well as recent review articles. As much
as possible, I will try to provide the references so that you can dive
deeper on any of these topics.

\begin{tcolorbox}[enhanced jigsaw, rightrule=.15mm, leftrule=.75mm, colbacktitle=quarto-callout-note-color!10!white, colframe=quarto-callout-note-color-frame, toprule=.15mm, opacitybacktitle=0.6, title=\textcolor{quarto-callout-note-color}{\faInfo}\hspace{0.5em}{Figure 1.3 -- From data generation to interpretation}, left=2mm, colback=white, bottomtitle=1mm, toptitle=1mm, bottomrule=.15mm, coltitle=black, titlerule=0mm, breakable, opacityback=0, arc=.35mm]

\begin{figure}[H]

\centering{

\pandocbounded{\includegraphics[keepaspectratio]{images/Figure 1.3 Bioinformatic Bottleneck.png}}

}

\caption{\label{fig-bottleneck}Bioinformatics Bottleneck}

\end{figure}%

From data generation to interpretation and publication. Sequencing
instruments generate large volumes of raw DNA reads from prepared
libraries, a step that is now relatively easy to purchase and scale. The
resulting data must then pass through a bottleneck of data cleaning,
filtering, and analysis---where analysts perform quality control,
alignment or assembly, and downstream processing while maintaining
documentation, reproducible code, and version control. Finally, cleaned
and analyzed results are interpreted, visualized, and synthesized into
figures, genomic summaries, and manuscripts for publication and broader
communication. Created with BioRender.com.

\end{tcolorbox}

\section{What makes someone a
``bioinformatician''?}\label{what-makes-someone-a-bioinformatician}

Bioinformatics can be approached through many different workflows---web
interfaces, command‑line tools, scripted analyses in languages like R or
Python, and now AI‑assisted systems---and this course aims to give
students a broad survey of that landscape rather than endorsing a single
``best'' way to do bioinformatics.

\textbf{Multiple ways to ``do genomics''}: For some people, doing
bioinformatics means interactively exploring genome browsers or running
analyses through web portals like Galaxy; for others, it means managing
workflows on Linux clusters, writing code, or curating databases and
infrastructure, and all of these activities fall under the very broad
genomics umbrella.

Job descriptions in the field range from analysts and software
developers to computational biologists and AI bioinformatics engineers,
illustrating that bioinformatics is defined more by the biological
questions and data than by any fixed skillset or toolchain.

\textbf{Web tools, command line, and pipelines}: Web‑based resources
such as NCBI, Ensembl, UCSC Genome Browser, Galaxy, and similar portals
make it possible to query databases, visualize genomes, and run common
analyses without local installation, which is especially useful for new
users and occasional genomics work.

Command‑line tools, workflow managers (for example, Nextflow or
Snakemake), and scripted pipelines support large‑scale, automated, and
reproducible analyses when projects or datasets grow more complex, but
they demand more technical setup and comfort with Unix‑like
environments.

\textbf{R and recent advances}: Modern R ecosystems---exemplified by
Akalin's Computational Genomics with R\textsuperscript{12}, which
introduces Bioconductor workflows, genomic intervals, and
high‑throughput sequencing analyses---show how a general‑purpose
language can connect statistical modeling, visualization, and
domain‑specific genomics packages in a single environment.

In this course R will appear in selected examples and advanced modules
(rather than as a primary focus) to highlight these newer genomics
capabilities and to give interested students a starting point if they
wish to pursue more scripted, R‑based analysis later.

\textbf{AI as an emerging bridge}: New AI‑driven platforms and
multi‑agent systems can help design, execute, and document
bioinformatics workflows, lowering the learning curve for building
reproducible pipelines by suggesting tools, generating code, and
tracking parameters automatically.

These approaches do not replace traditional skills but offer another
path into genomics work, especially for students who are more
comfortable describing goals in natural language than writing full
pipelines from scratch.

\textbf{How this course uses these ecosystems}: This course was
originally designed before current AI and tooling advances, and newer
R‑based genomics workflows. Therefore, AI‑assisted workflows, and
updated R resources are being integrated incrementally while preserving
the core emphasis on concepts and data literacy.

Students will be exposed to web‑based tools, the Unix command line, and
examples of scripted analyses, including some in R, and are encouraged
to share outside resources that offer clearer or simpler solutions,
reinforcing that there are many valid ways to contribute to genomics and
that ``doing bioinformatics'' does not have to look the same for
everyone.

\section{Types of Genome Analysis Research
Projects}\label{types-of-genome-analysis-research-projects}

Genome analysis underpins a wide range of applications in
\textbf{medical genetics}, including identifying variants that
contribute to Mendelian disorders, studying complex disease risk, and
informing targeted therapies\textsuperscript{13}. In
\textbf{agriculture}, genomic selection and genome‑wide association
studies accelerate breeding by linking genetic markers to traits such as
yield, stress tolerance, and disease resistance in crops and
livestock\textsuperscript{2}. The role of computational biologists
varies significantly too, from support staff to project
leaders\textsuperscript{14}.

In \textbf{ecology and evolution}, reference genomes and
population‑scale sequencing allow researchers to study adaptation,
speciation, and biodiversity at unprecedented resolution, including for
non‑model organisms\textsuperscript{1,2}. In \textbf{public health},
genomic surveillance of pathogens supports outbreak tracking,
antimicrobial resistance monitoring, and the design of diagnostics and
vaccines\textsuperscript{13} (Figure~\ref{fig-equity}).

\subsection{Three perspectives on genomics and
bioinformatics}\label{three-perspectives-on-genomics-and-bioinformatics}

One way to organize this diversity is to think about three complementary
``views'' of bioinformatics and genomics, moving from small to large
scales\textsuperscript{10}.

\begin{itemize}
\item
  \textbf{Cell:} At the molecular level, genomics focuses on DNA (the
  genome), RNA (the transcriptome), and proteins (the proteome).
  Questions at this scale include: What genes and transcripts exist? How
  are they regulated? How do they interact in pathways and networks?
  Many functional genomics tools (e.g., RNA‑seq, ChIP‑seq) are designed
  to answer questions at this cellular scale\textsuperscript{10}.
\item
  \textbf{Organism:} At the level of an individual, genomics asks how
  the same genome produces different phenotypes across tissues,
  developmental stages, environments, and disease states. Typical
  projects compare gene expression or other genomic features across
  conditions (for example, healthy vs.~diseased tissue, or individuals
  exposed to different environments) to understand how genomes respond
  to internal and external signals\textsuperscript{10}.
\item
  \textbf{Tree of life:} At the largest scale, comparative genomics uses
  genome sequences from many species to study phylogeny, genome
  structure, and macroevolutionary processes such as gene family
  expansion, whole‑genome duplication, and chromosomal rearrangements.
  These studies emphasize both the shared molecular ``toolkit'' of life
  and the ways genomes are shaped by adaptation to different
  environments\textsuperscript{1,2,10}.
\end{itemize}

The required readings for this week highlight this tree‑of‑life
perspective. One paper surveys the current state of animal genome
sequencing across the animal kingdom, quantifying which lineages have
reference genomes and how assembly quality has improved over
time\textsuperscript{1}. The complementary plant genomics paper
summarizes 20 years of land‑plant genome sequencing, showing how
taxonomic coverage, assembly quality, and global participation have
changed, and pointing out important gaps and inequities in who is
leading and benefiting from plant genomics projects\textsuperscript{2}.

\subsection{Hypothesis‑driven and discovery‑driven
genomics}\label{hypothesisdriven-and-discoverydriven-genomics}

As sequencing has become cheaper and datasets have grown larger, most
impactful genomics studies now go beyond a single genome assembly or
alignment. Instead, they combine substantial data generation with either
explicit hypothesis testing, open‑ended discovery, or (ideally)
both\textsuperscript{10,15}.

\begin{itemize}
\item
  In \textbf{hypothesis‑driven} genomics, researchers start with
  specific predictions (for example, that a gene family will be more
  expanded in species from colder environments) and design sampling,
  sequencing, and statistical analyses to test those predictions. These
  projects emphasize careful experimental design, replication, and
  rigorous statistical models to distinguish signal from
  noise\textsuperscript{10}.
\item
  In \textbf{discovery‑driven} genomics, researchers focus on exploring
  new genomes or large datasets to see what patterns emerge, often
  without a single predefined ``main'' hypothesis. Examples include
  assembling genomes for under‑studied clades, cataloging structural
  variation, or visualizing complex expression or variant datasets to
  identify unexpected clusters or associations. This exploratory work is
  especially important in genomics, where high‑dimensional data can
  contain many unanticipated signals\textsuperscript{1,2,15}.
\end{itemize}

Both modes are valuable and interdependent. Exploratory analyses are
often where new ideas and hypotheses come from, while hypothesis‑driven
analyses are essential for testing those ideas and connecting genomic
patterns to mechanisms\textsuperscript{15}. Throughout the course, you
will see studies that move back and forth between these modes.

\begin{tcolorbox}[enhanced jigsaw, rightrule=.15mm, leftrule=.75mm, colbacktitle=quarto-callout-note-color!10!white, colframe=quarto-callout-note-color-frame, toprule=.15mm, opacitybacktitle=0.6, title=\textcolor{quarto-callout-note-color}{\faInfo}\hspace{0.5em}{Lab connections}, left=2mm, colback=white, bottomtitle=1mm, toptitle=1mm, bottomrule=.15mm, coltitle=black, titlerule=0mm, breakable, opacityback=0, arc=.35mm]

\begin{itemize}
\tightlist
\item
  \textbf{Lab01 (Section Section~\ref{sec-lab01})} asks you to find and
  describe a genome analysis paper, then classify it using the three
  perspectives introduced in this chapter and the ideas of
  hypothesis‑driven versus discovery‑driven
  genomics\textsuperscript{1,2,10}.
\item
  If you would like a refresher on how to read scientific papers,
  \href{appendixA.qmd}{Appendix A} walks through hypothesis
  identification, concept mapping, and designing follow‑up experiments
  using an example scientific article on albinism\textsuperscript{16}.
\end{itemize}

These lab activities reinforce the high‑level view of genomics developed
in this chapter and prepare you for later, more technical labs.

\end{tcolorbox}

\subsection{Looking ahead: case studies that span
scales}\label{looking-ahead-case-studies-that-span-scales}

Early next week, you will read a recent genomics paper on antifreeze
proteins in polar and deep‑sea fishes that illustrates how these ideas
come together\textsuperscript{3}. Without going into detail here, that
study:

\begin{itemize}
\tightlist
\item
  Uses high‑quality genome assemblies and comparative analyses across
  many species (tree‑of‑life perspective).
\item
  Examines how gene families and genome structure relate to
  environmental variables such as temperature and depth (organism and
  environment).
\item
  Combines exploratory genome‑scale analyses (discovery‑driven) with
  explicit statistical tests of how gene copy number and genomic
  location change across species and habitats (hypothesis‑driven).
\end{itemize}

As you read it, you will be able to connect the analyses in that paper
to the three perspectives (cell, organism, tree of life) and to the
balance between hypothesis‑driven and discovery‑driven genomics
discussed above.

These opportunities come with ethical, legal, and social questions
(Figure~\ref{fig-equity}). For human data, key concerns include informed
consent, privacy, data security, and the potential for misuse or
discrimination based on genetic information\textsuperscript{13}. For
non‑human genomics, recent assessments of animal and plant genome
projects have emphasized inequities in who leads sequencing efforts, how
benefits and credit are shared, and how historical patterns of
colonialism and ``parachute science'' continue to shape access to
biodiversity and genomic resources\textsuperscript{1,2}.

International agreements and community norms now guide data sharing and
\textbf{equity considerations}, encouraging collaborations that involve
local researchers, respect Indigenous knowledge, and support capacity
building in regions where study species are found\textsuperscript{1,2}.
Throughout the course, case studies such as long‑term animal and plant
genome surveys will be used to discuss how technical decisions, sampling
strategies, and authorship practices connect to these broader ethical
themes\textsuperscript{1,2,11}.

\begin{tcolorbox}[enhanced jigsaw, rightrule=.15mm, leftrule=.75mm, colbacktitle=quarto-callout-note-color!10!white, colframe=quarto-callout-note-color-frame, toprule=.15mm, opacitybacktitle=0.6, title=\textcolor{quarto-callout-note-color}{\faInfo}\hspace{0.5em}{Figure 1.4 -- Applications and equity}, left=2mm, colback=white, bottomtitle=1mm, toptitle=1mm, bottomrule=.15mm, coltitle=black, titlerule=0mm, breakable, opacityback=0, arc=.35mm]

\begin{figure}[H]

\centering{

\includegraphics[width=0.5\linewidth,height=\textheight,keepaspectratio]{images/Figure_1.4_Equity.jpeg}

}

\caption{\label{fig-equity}Equity}

\end{figure}%

A four‑quadrant schematic illustrates major application areas of
genomics: medical genetics (top left), agriculture and food security
(top right), ecology and conservation (bottom left), and public health
and pathogen genomics (bottom right). Each quadrant contains simple
domain‑specific icons (for example, DNA and a clinical symbol, crop
plants, ecosystems, and pathogens) to emphasize the diversity of
contexts in which genomic data are generated and used. Surrounding the
central square, icons for consent, privacy, benefit‑sharing, and global
participation highlight that technical choices in genomic research and
applications are embedded within ethical, legal, and social frameworks
that should be considered across all domains. Figure concept and artwork
generated with an AI‑assisted illustration tool.

\end{tcolorbox}

\section{Ethics spotlight -- timelines, cases, and
trust}\label{ethics-spotlight-timelines-cases-and-trust}

\emph{Note: The goal of this section is to situate genomics in its
historical and ethical context in a factual, objective way, without
assigning blame to students or asking them to endorse any particular
viewpoint.}

The timeline in Figure~\ref{fig-genomics-timeline} begins the story of
milestones in human genomics with the Belmont
Report\textsuperscript{17}, developed after a series of widely
criticized research practices including the Tuskegee syphilis research
study. In that study, conducted in Alabama, effective treatment was
withheld from Black men with syphilis; the case is now frequently cited
in bioethics as a reason for today's requirements for informed consent
and ongoing oversight of human‑subjects research\textsuperscript{18}.

Several later examples illustrate how questions about consent, privacy,
and benefit sharing continue to arise in genetics and genomics.
Henrietta Lacks's cervical tumor cells were collected in 1951 without
her knowledge or consent, leading to the HeLa cell line and, decades
later, to debates about privacy and family involvement when her cells
and genome sequence were shared widely\textsuperscript{19--21}. Members
of the Havasupai Tribe in Arizona provided blood samples for diabetes
research and later learned their samples had also been used in studies
of topics such as schizophrenia and population origins, prompting a
lawsuit, a 2010 settlement, and renewed discussion of how consent forms
describe future uses of samples\textsuperscript{22}. Human subject
research is now strictly regulated, even in the context of survey-based
work, through institutional review boards (IRB) that have guidelines to
protect consent and reduce harm\textsuperscript{23}. At Auburn, the AU
Human Research Protection Program is responsible for the ethical and
regulatory requirements related to the protection of human participants
in research. The Program includes the Institutional Official (IO; Senior
Vice President for Research \& Economic Development), the Institutional
Review Board (IRB) for the protection of human subjects in research, and
the Office of Research Integrity \& Compliance (ORIC).

Genomics today includes many different projects, from large biobanks to
smaller disease or population studies, and these efforts depend on
volunteered samples and trust that data will be used and shared in ways
consistent with what participants and communities were
told\textsuperscript{13}. International plant and animal genomics
projects also raise questions about benefit sharing, authorship, and
power imbalances between institutions and the regions or communities
that provide samples\textsuperscript{1,2}.

This course does not require students to agree with any particular
ethical position, but it does present these well‑documented cases as
background for understanding why modern genomics emphasizes informed
consent, privacy protections, and clear communication about how data may
be used. Later chapters briefly revisit these themes in connection with
topics like data sharing and genome‑wide studies; students are
welcome---but not required---to explore the broader ethical literature
or related historical discussions if they are interested.

\section{Course roadmap: ``reading genomes'' over the
semester}\label{course-roadmap-reading-genomes-over-the-semester}

This course is organized around the idea that \textbf{reading genomes}
is a multi‑step process: learning the language of sequences and formats,
understanding how different technologies generate data, and practicing
how to transform raw reads into biological
insight\textsuperscript{10};\textsuperscript{12}. Lectures and readings
provide conceptual foundations, while labs and projects give hands‑on
experience with real data and tools\textsuperscript{12}.

Early labs focus on basic data formats, quality control, and simple
workflows for short‑read data, building toward more complex analyses of
alignment, variant calling, and genome annotation later in the
semester\textsuperscript{10}. A semester‑long \textbf{genome analysis
project} asks each student or group to design and execute a focused
analysis (for example, on a chosen genome or dataset), document their
workflow in a reproducible way, and present their findings, mirroring
how modern genomics research is carried out.

In parallel, a \textbf{mock grant review panel} assignment exposes
students to how genomics proposals are evaluated: what makes a question
compelling, whether the proposed analyses are appropriate and feasible,
and how reviewers weigh innovation, rigor, and broader impacts.
Together, these components---lectures, integrated readings, labs, the
project, and the review panel---are designed to help students move from
passively reading about genomes to actively \textbf{reading, analyzing,
and critiquing} genomic data and studies.

\section{A Consistent Example}\label{a-consistent-example}

Throughout the semester, the class will return again and again to a
single yeast genome sequence dataset, following it from raw sequencing
reads all the way to biological interpretation. The dataset comes from
an experimental evolution study in \emph{Saccharomyces cerevisiae} in
which populations were propagated for many generations under defined
conditions, allowing standing genetic variation to be reshaped by
selection in a controlled laboratory environment\textsuperscript{24}. To
keep the full workflow manageable, the course uses the compact sacCer3
reference genome (about 12 Mb across 16 chromosomes), which makes it
feasible to run industry-standard bioinformatics workflows on the
high-performance computing cluster within realistic time and memory
limits.

In Lab 4, you will download the yeast reads from the NCBI SRA accession
SRR1693691, perform initial quality control with FastQC, trim away
low-quality bases and adapters, and align the cleaned reads to sacCer3,
ending with an indexed BAM file and basic coverage and mapping
statistics. In Lab 7, you will pick up the same alignment file to mark
PCR duplicates, run additional QC checks such as samtools flagstat and
MultiQC, and then use GATK's ``best practices'' variant calling workflow
to generate and filter a set of SNPs, comparing different hard-filtering
thresholds and visual summaries of variant quality metrics. In Lab 9,
you will load both the pre- and post-filtering BAM files and the
resulting VCFs into IGV, visually inspect specific regions of the
genome, and decide which candidate variants look like real evolutionary
changes versus sequencing or mapping artifacts, interpreting those
variants in terms of coding changes and potential functional
consequences for the yeast.

Most of the heavy lifting on the command line will be handled by
provided shell scripts, so students can focus on understanding what each
step does and what each output means rather than on debugging complex
pipelines from scratch. The overarching goals are to help them feel
comfortable working with real whole-genome data on a high-performance
cluster, to connect the dots between textbook methods and concrete files
like FASTQ, BAM, and VCF, and to practice reading genomes as biological
documents---writing clear methods, reporting variants in a biomedical
context, and reasoning about how specific genetic changes might shape
evolutionary trajectories in a model eukaryote.

More on the GATK workflow and these file formats can be found in
subsequent chapters, so do not worry so much about the details right
now. You can also read more about this semester dataset in
\href{appendixC.qmd}{Appendix C}.

\section{Looking ahead to Chapter 2}\label{looking-ahead-to-chapter-2}

Chapter 2 builds directly on this foundation by unpacking the sequencing
technologies behind the timeline in Figure~\ref{fig-genomics-timeline}.
It introduces first‑, second‑, and third‑generation sequencing
platforms, compares Illumina, PacBio, and Oxford Nanopore data, and
connects platform choice to alignment, de novo assembly, and
experimental design. The goal is to give you a single, integrated
reference for sequencing technologies and data so that later chapters on
algorithms and applications have a concrete foundation.

\begin{center}\rule{0.5\linewidth}{0.5pt}\end{center}

\section{Reflection questions}\label{reflection-questions}

You can use the following questions to reflect on the concepts and
papers introduced in this chapter. These are intended for after you have
finished the chapter and associated readings. Some of these will lead
directly into future lab assignments.

\begin{enumerate}
\def\labelenumi{\arabic{enumi}.}
\item
  \textbf{Three perspectives:} For a genomics study you have encountered
  (in this course or elsewhere), which perspective is most prominent:
  cell, organism, or tree of life? What specific datasets or analyses
  support your answer?
\item
  \textbf{Multiple scales:} How does that study connect at least two of
  the three perspectives (for example, using cellular‑scale data to
  answer organism‑level or tree‑of‑life questions)?
\item
  \textbf{Discovery vs.~hypothesis:} Identify one example of
  discovery‑driven analysis (exploration of data without a single,
  predefined main hypothesis) and one example of hypothesis‑driven
  analysis (explicit testing of a stated prediction) from the assigned
  readings (animal genomes\textsuperscript{1}, plant
  genomes\textsuperscript{2}, or the antifreeze‑protein paper you will
  read next week\textsuperscript{3}). How do these approaches complement
  each other?
\item
  \textbf{Data and methods:} In the animal and plant genomics overview
  papers, what kinds of data (for example, reference genomes, metadata)
  and methods (for example, phylogenetic analyses, summary statistics
  across taxa) are needed to draw conclusions about progress and gaps
  across the tree of life\textsuperscript{1,2}? How do these differ from
  the data and methods needed to investigate a specific gene family,
  such as antifreeze proteins\textsuperscript{3}?
\item
  \textbf{Equity and participation:} Both the animal and plant genomics
  overview papers highlight mismatches between where species are found,
  where sequencing work is done, and who leads that
  work\textsuperscript{1,2}. What are some concrete ways that future
  genomics projects could promote more equitable collaboration and
  benefit‑sharing?
\item
  \textbf{Future directions:} Imagine you are designing a new genomics
  project. How would you balance hypothesis‑driven goals (specific
  questions you want to test) with discovery‑driven goals (opening space
  for unexpected findings)? Which of the three perspectives (cell,
  organism, tree of life) would your project emphasize first, and how
  might you connect to the others as your work develops?
\end{enumerate}

\section{References}\label{references-1}

\chapter{Sequencing Technologies and Data
Formats}\label{sequencing-technologies-and-data-formats}

\section{Why this chapter?}\label{sec-ch2}

Chapter 1 framed genomics as the process of \emph{reading genomes} and
introduced core ideas such as reference assemblies, alignment, and de
novo assembly as recurring operations in this
course\textsuperscript{1,2,10}. In this chapter, the focus shifts to the
sequencing technologies themselves, the data they generate, and how
platform choice interacts with experimental design and downstream
bioinformatics workflows\textsuperscript{10,25,26}.

By the end of the chapter, you should be able to:

\begin{itemize}
\tightlist
\item
  Explain how first‑, second‑, and third‑generation sequencing methods
  work at a conceptual level.\\
\item
  Compare Illumina, PacBio, and Oxford Nanopore platforms in terms of
  read length, accuracy, throughput, and error
  profiles\textsuperscript{10,25,26}.\\
\item
  Decide when alignment to a reference is appropriate and when de novo
  assembly or hybrid approaches are needed\textsuperscript{10,25,26}.\\
\item
  Interpret key assembly quality metrics and understand how long‑range
  technologies (Hi‑C, optical maps, Strand‑seq) scaffold contigs to
  chromosomes\textsuperscript{25,26}.\\
\item
  Recognize how modern sequencing companies increasingly provide full
  ``sequencer‑to‑variants'' workflows and what that means for practicing
  bioinformatics\textsuperscript{10,26}.
\end{itemize}

Along the way, we connect the historical genomics timeline from Chapter
1 to current platforms and highlight ethical questions raised by the
rapid expansion of genomic data, especially in human
populations\textsuperscript{1,13}.

\section{From timelines to
laboratories}\label{from-timelines-to-laboratories}

In Chapter 1, Figure~\ref{fig-genomics-timeline} traced the history of
genomics from the discovery of DNA as the hereditary molecule through
Sanger sequencing, the Human Genome Project, high‑throughput short‑read
platforms, and the emergence of long‑read technologies such as PacBio
and Oxford Nanopore. Each milestone reflects two intertwined trends:
improvements in chemistry and instrumentation, and increasingly
sophisticated computational methods required to turn raw signals into
usable sequence data\textsuperscript{10,25}.

In practice, most modern genome projects follow a workflow similar to
the one you saw in Lecture 2: experimental design and library
preparation, sequencing on one or more platforms, quality control,
alignment or assembly, variant calling or annotation, visualization, and
long‑term data storage\textsuperscript{10}. This chapter focuses
primarily on the first half of that pipeline up to the point where reads
are ready to be aligned or assembled; later chapters return to
algorithms and downstream analyses in more detail.

\begin{tcolorbox}[enhanced jigsaw, rightrule=.15mm, leftrule=.75mm, colbacktitle=quarto-callout-note-color!10!white, colframe=quarto-callout-note-color-frame, toprule=.15mm, opacitybacktitle=0.6, title=\textcolor{quarto-callout-note-color}{\faInfo}\hspace{0.5em}{Connecting Chapters 1 and 2}, left=2mm, colback=white, bottomtitle=1mm, toptitle=1mm, bottomrule=.15mm, coltitle=black, titlerule=0mm, breakable, opacityback=0, arc=.35mm]

Chapter 1 emphasized \emph{what} we want from genomes---reference
assemblies, annotations, and interpretable variation---while treating
sequencing platforms as a black box. This chapter opens that box,
because understanding how the data are generated is essential for making
good decisions about coverage, library design, and downstream analysis
choices.

\end{tcolorbox}

\section{First‑generation sequencing: Sanger and
BAC‑by‑BAC}\label{firstgeneration-sequencing-sanger-and-bacbybac}

First‑generation sequencing refers to the Sanger dideoxy method, which
dominated from the late 1970s through the early 2000s and underpinned
the Human Genome Project\textsuperscript{10}. In Sanger sequencing, a
single‑stranded DNA template is copied by DNA polymerase in the presence
of a primer, the four normal deoxynucleotides (dNTPs), and small amounts
of chain‑terminating dideoxynucleotides (ddNTPs) labeled with distinct
dyes\textsuperscript{10}. A normal dNTP has a 3′‑OH group on the
deoxyribose sugar that allows the next nucleotide to be added, whereas a
ddNTP lacks this 3′‑OH, so once incorporated no further extension is
possible and the growing strand stops at that base.

Because polymerase chooses between dNTPs and ddNTPs at random, the
reaction generates a \textbf{nested set} of fragments that all start at
the same primer and end at every possible position where that base
occurs in the template\textsuperscript{10}. Earlier protocols ran four
separate reactions (one per ddNTP) and resolved the fragments on a
polyacrylamide gel, with bands visualized by autoradiography; later
instruments combined fluorescently labeled ddNTPs into a single reaction
and read fragment size and dye color as they migrated past a laser in a
capillary, automatically converting peaks into a base‑called
chromatogram\textsuperscript{10}. Modern capillary Sanger sequencing
routinely yields reads of 500--1000 bp with error rates below 1\%, which
is why it remains the ``gold standard'' for validating variants and
small loci today\textsuperscript{10}.

From a whole‑genome perspective, however, Sanger chemistry is inherently
\textbf{low throughput}. Each capillary run produces one read for one
template, so thousands to millions of reactions and capillaries are
required to cover a large genome\textsuperscript{10}. Even with
high‑density instruments like the ABI PRISM 3700, which multiplexed 96
capillaries and automated loading, producing enough sequence for a draft
human genome required years of continuous operation and substantial
personnel, reagent, and instrument costs\textsuperscript{10}.

\subsection{BAC‑by‑BAC genome
sequencing}\label{bacbybac-genome-sequencing}

To scale Sanger sequencing to large genomes, the Human Genome Project
and similar efforts used a \textbf{BAC‑by‑BAC} strategy. Genomic DNA was
first fragmented into \textasciitilde100--200 kb pieces and cloned into
bacterial artificial chromosomes (BACs), producing a library of
large‑insert clones that tiled across each
chromosome\textsuperscript{10}. Genetic linkage maps and physical maps
(for example, restriction fragment patterns) were then used to order
overlapping BACs along the chromosome, creating a scaffold or ``tiling
path'' that indicated which region of the genome each clone
represented\textsuperscript{10,25}.

Each BAC in this tiling path was then sequenced by \textbf{shotgun
Sanger sequencing}: the BAC insert was sheared into smaller fragments,
cloned again, and hundreds to thousands of individual Sanger reads were
generated and assembled into contigs using overlap‑layout‑consensus
algorithms\textsuperscript{10}. Finishing a BAC required closing gaps
and resolving repeats by designing new primers, resequencing troublesome
regions, and manually inspecting assemblies, often iterating many times
to obtain a contiguous, high‑quality sequence\textsuperscript{10}.
Finally, finished BAC sequences were stitched together along the tiling
path, using the BAC map to determine their order and orientation and
filling any remaining gaps as resources allowed\textsuperscript{10,25}.

This pipeline produced extremely accurate reference sequences, but it
illustrates why Sanger‑era genomics was \textbf{inefficient and costly}
for whole genomes. Every step---constructing BAC libraries, mapping
clones, performing thousands of Sanger reactions per BAC, and manually
finishing assemblies---was labor‑intensive and required large teams and
dedicated sequencing centers\textsuperscript{10,25}. The Human Genome
Project consumed over a decade and billions of dollars partly because
Sanger chemistry and BAC‑by‑BAC strategies, while reliable, simply could
not deliver base‑pair--level coverage at the scale and speed that later
high‑throughput technologies made routine\textsuperscript{10}.

From a modern perspective, Sanger's legacy is conceptual. It established
the idea that sequencing is a combination of \textbf{chemistry} (how
nucleotides are incorporated and terminated) and \textbf{signal
processing} (how fragment sizes and dye colors are turned into digital
base calls with quality scores)\textsuperscript{10}. Those same
ideas---polymerase‑based synthesis, chain termination or detection
events, and computational basecalling---reappear in second‑ and
third‑generation platforms even as the physics of detection and the
scale of data generation change\textsuperscript{25}.

\section{Second‑generation sequencing: short reads at
scale}\label{secondgeneration-sequencing-short-reads-at-scale}

Second‑generation or \textbf{``next‑generation'' sequencing (NGS)}
refers to platforms that produce millions to billions of short reads in
parallel by imaging clonal clusters or beads\textsuperscript{10}.
Several competing technologies appeared in the mid‑2000s---454
pyrosequencing (Roche), Solexa/Illumina sequencing‑by‑synthesis, Applied
Biosystems SOLiD sequencing‑by‑ligation, and later Ion Torrent's
semiconductor‑based chemistry---each with its own way of coupling
nucleotide incorporation to a detectable signal\textsuperscript{10}.

\subsection{Early NGS chemistries and why Illumina
won}\label{early-ngs-chemistries-and-why-illumina-won}

454 pyrosequencing amplified DNA on beads in emulsion PCR and flowed one
unlabeled dNTP across a picotiter plate at a time; incorporation
released pyrophosphate, which was converted to light via coupled
enzymatic reactions, so peak height corresponded to the length of a
homopolymer run\textsuperscript{10}. The platform's early strength was
\textbf{read length}: several hundred bases when early Illumina reads
were often \textless50~bp. However, homopolymer runs were error‑prone,
throughput and scalability lagged behind, and the cost per base remained
high, so 454 was eventually discontinued as Illumina read lengths and
yields improved\textsuperscript{10,25}.

SOLiD also used bead‑based clonal amplification but read sequence via
ligation of fluorescently labeled oligonucleotide probes, encoding each
base in \textbf{color space} rather than direct base
calls\textsuperscript{10}. This design provided theoretical
error‑checking redundancy, but it came with practical problems: file
sizes were large, interpretation required extra decoding steps, and the
overall error rate and effective read length lagged behind Illumina's
steadily improving platforms\textsuperscript{10}. Color‑space quality
scores and formats were cumbersome enough that many downstream tools
never fully supported them, which further reduced SOLiD's appeal.

Ion Torrent eliminated optics entirely and detected the pH change
(proton release) that accompanies nucleotide incorporation in
semiconductor wells, trading read length and homopolymer performance for
cheaper, smaller instruments\textsuperscript{10}. It remains useful in
some targeted amplicon and diagnostics settings but has not displaced
Illumina for large‑scale genomics.

In contrast, Illumina's (originally Solexa's) sequencing‑by‑synthesis
chemistry combined scalable throughput, a relatively simple short‑read
error profile dominated by substitutions, and continuous improvements in
read length, cost per base, and automation\textsuperscript{10,25}. Over
time this allowed Illumina to become the dominant short‑read platform
used in research, clinical, and agricultural genomics.

\subsection{Illumina
sequencing‑by‑synthesis}\label{illumina-sequencingbysynthesis}

In the canonical Illumina workflow, genomic DNA (or cDNA) is fragmented
and specialized adapters are ligated to both ends; these adapters
provide binding sites for primers, indices (barcodes), and short
sequences complementary to oligos covalently attached to the flow‑cell
surface\textsuperscript{10}. During \textbf{cluster generation},
fragments hybridize to one type of surface oligo and are copied by
polymerase; the original strand is washed away, and the copy bends over
to form a ``bridge'' to the second oligo type, where it is extended
again. Repeated cycles of this bridge amplification create dense
clusters of clonal copies of the original fragment, each tethered at
both ends to the glass slide\textsuperscript{10}.

After amplification, one strand in each cluster is selectively cleaved
and washed away and the remaining 3′ ends are blocked to prevent
unwanted priming, leaving a lawn of single‑stranded clusters ready for
sequencing. \textbf{Sequencing‑by‑synthesis} then proceeds in cycles: a
sequencing primer anneals to each cluster, and a mix of four
fluorescently labeled, 3′‑blocked reversible terminator nucleotides is
added\textsuperscript{10}. In each cycle, exactly one nucleotide is
incorporated into each cluster based on the template base; a
high‑resolution camera images the flow cell to record the color and
intensity at each cluster; then the dye and blocking group are
chemically removed so the next cycle can proceed. The number of cycles
determines the read length, typically 75--300 bases for most current
instruments\textsuperscript{10}.

Paired‑end libraries exploit the fact that both ends of each fragment
have adapters. After Read 1 is complete and index reads have been
acquired, the extension product is melted away, the intact template
folds over to the opposite surface oligo, and a new primer initiates
Read 2 from the other end of the fragment. Index 1 and index 2 reads are
interleaved into the run capture with the sample barcodes, allowing
hundreds of libraries to be pooled on a single lane of sequencing.

\begin{tcolorbox}[enhanced jigsaw, rightrule=.15mm, leftrule=.75mm, colbacktitle=quarto-callout-note-color!10!white, colframe=quarto-callout-note-color-frame, toprule=.15mm, opacitybacktitle=0.6, title=\textcolor{quarto-callout-note-color}{\faInfo}\hspace{0.5em}{Figure 2.1 -- Illumina sequencing workflow}, left=2mm, colback=white, bottomtitle=1mm, toptitle=1mm, bottomrule=.15mm, coltitle=black, titlerule=0mm, breakable, opacityback=0, arc=.35mm]

\begin{figure}[H]

\centering{

\includegraphics[width=0.75\linewidth,height=\textheight,keepaspectratio]{images/Fig 2-1 Illumina Sequencing.png}

}

\caption{\label{fig-illumina-workflow}Overview of Illumina library
preparation, bridge amplification, sequencing‑by‑synthesis, and
imaging.}

\end{figure}%

This figure summarizes the four main stages of an Illumina run: library
preparation with adapters, bridge amplification to form clusters, cyclic
sequencing‑by‑synthesis with fluorescent nucleotides, and imaging to
produce per‑cluster intensity traces that are converted into base calls.
Scientists receive FASTQ files from these reads which will be covered
more in the next chapter. Created with BioRender.com.

\end{tcolorbox}

\subsection{Throughput, evolution, and where sequencing happens
now}\label{throughput-evolution-and-where-sequencing-happens-now}

Modern Illumina instruments routinely generate paired‑end reads of
150--300~bp, with total outputs ranging from a few gigabases per run on
benchtop systems (e.g., MiSeq) to many terabases per run on high‑end
systems (e.g., NovaSeq) at very low per‑base error
rates\textsuperscript{10}. Each cycle's fluorescence image is converted
into intensity traces and then into base calls with associated
PHRED‑scaled quality scores, which downstream tools use to weight
alignments and variant calls; the FASTQ format that stores these calls
and quality scores is introduced formally in Chapter~3.

Illumina has continued to innovate even as long‑read platforms have
matured. Patterned flow cells that regularize cluster positions,
two‑channel or one‑channel dye schemes that reduce optical complexity,
steadily higher cluster densities, and improved chemistries have all
driven large gains in throughput and reductions in cost per base while
maintaining or improving accuracy\textsuperscript{25,26}. As a result,
short‑read NGS remains the workhorse for high‑sample‑number
studies---population resequencing, exomes, RNA‑seq, ChIP‑seq, and many
clinical assays---even when long‑read data are used for assembly or
structural variant discovery.

In the early NGS era, many universities and research institutes operated
their own sequencing cores with 454, SOLiD, and later Illumina
instruments, employing dedicated staff to handle library preparation,
runs, and first‑pass QC\textsuperscript{10}. As instruments became
larger, more expensive, and more specialized, and as commercial
providers began to offer highly competitive pricing and turnaround
times, many institutions shifted toward \textbf{outsourcing} most
sequencing. At Auburn, for example, several labs still maintain older
Illumina machines for small projects or pilot data, but most large‑scale
sequencing is sent to external facilities such as Novogene, HudsonAlpha
(Huntsville, Alabama), or similar providers that operate high‑throughput
Illumina platforms at scale.

\subsection{Strengths and limitations of short‑read
NGS}\label{strengths-and-limitations-of-shortread-ngs}

Short‑read NGS revolutionized genomics by making whole‑genome,
whole‑exome, and RNA‑seq experiments affordable and practical for many
labs\textsuperscript{10,13}. At the same time, the short read lengths
and PCR‑based library preparation introduce characteristic challenges
that you will see repeatedly in this course and that are discussed in
detail by van~Dijk~et~al.~2018\textsuperscript{25}:

\begin{enumerate}
\def\labelenumi{\arabic{enumi}.}
\item
  \textbf{Assembly}: De novo assembly is difficult when reads are much
  shorter than common repeats and structural variants. Short reads
  cannot uniquely span many repeat families, so assemblers break contigs
  at these points, resulting in fragmented genomes with hundreds to
  thousands of contigs and potential mis‑assemblies in repeat‑rich
  regions\textsuperscript{25,26}.
\item
  \textbf{Phasing}: In genetics, \emph{phasing} (or haplotype phasing)
  is the crucial process of determining which genetic variants are
  located together on the same haplotype during inheritance from each
  parent. Because short reads cover only a few hundred bases, they often
  cannot connect heterozygous variants that are far apart on the same
  chromosome copy. This makes it hard to determine which alleles
  co‑occur on the same haplotype (maternal vs paternal) and limits our
  ability to study inheritance patterns of compound heterozygosity and
  complex haplotypes in human and non‑human
  genomes\textsuperscript{25,26}.
\item
  \textbf{Transcript identification}: RNA‑seq with short reads can
  quantify gene‑level expression well, but reconstructing full‑length
  transcript isoforms is challenging. Short reads provide only local
  views of exon--exon junctions, so many possible isoform combinations
  are consistent with a given set of reads unless sequencing depth is
  extremely high; long‑read transcriptomics (e.g., PacBio Iso‑Seq or
  nanopore cDNA/direct RNA) addresses this limitation by reading
  complete transcripts\textsuperscript{25,26}.
\item
  \textbf{GC bias}: PCR amplification during library prep and cluster
  generation is less efficient at extreme GC contents. Regions that are
  very GC‑rich or GC‑poor therefore receive lower coverage, leading to
  false negatives in variant calling and uneven representation in
  assemblies and RNA‑seq experiments\textsuperscript{25}.
\end{enumerate}

These limitations motivated both the development of third‑generation
long‑read technologies and a wave of algorithmic advances---better
mappers, local re‑assembly around structural variants, graph‑based
references, and specialized isoform reconstruction tools---that you will
encounter in Chapters~3--5 and beyond. This is also why you might notice
that in Figure~\ref{fig-genomics-timeline}, the Human Genome is often
cited as having been completed in 2003, but the complete
\textbf{Telomere-to-Telomere (T2T)} genome assembly took nearly two more
decades to complete, largely assisted by third-generation sequencing
efforts.

\section{Third‑generation sequencing: long reads and single
molecules}\label{thirdgeneration-sequencing-long-reads-and-single-molecules}

Third‑generation or \textbf{long‑read sequencing (LRS)} describes
technologies that read individual DNA or RNA molecules in real time,
producing reads that can span tens to hundreds of
kilobases\textsuperscript{25,26}. Two major players are Pacific
Biosciences and \textbf{Oxford Nanopore Technologies (ONT)}, alongside
several synthetic and linked‑read approaches that attempted to bridge
the gap between short‑ and long‑read data\textsuperscript{25,27}.

\subsection{PacBio SMRT and HiFi}\label{pacbio-smrt-and-hifi}

\textbf{Pacific Biosciences (PacBio)} \textbf{single‑molecule real‑time
(SMRT)} sequencing uses tiny observation chambers called
\textbf{zero‑mode waveguides (ZMWs)}, each containing a single DNA
polymerase complexed with a circular ``SMRTbell''
template\textsuperscript{25}. Fluorescently labeled nucleotides diffuse
into the ZMW; as the polymerase incorporates bases, short pulses of
fluorescence are recorded in real time and converted into a continuous
long read\textsuperscript{25,26}.

Early SMRT runs produced continuous long reads (CLRs) tens of kilobases
long but with raw error rates around 10--15\%, dominated by random
insertions and deletions\textsuperscript{25}. This is why early genome
assembly papers using long read data were often completed with short
read sequences to improve accuracy. Modern workflows instead emphasize
\textbf{HiFi} (high‑fidelity) reads where the polymerase repeatedly
circles the same insert, generating multiple passes from which a
circular consensus sequence can be derived, trading some maximum length
for accuracies ≥99.9\% over 10--25~kb\textsuperscript{26,27}. Because
polymerase kinetics depend on base modifications, SMRT sequencing can
also directly detect certain DNA methylation events when analyzed with
specialized software\textsuperscript{25,26}.

Improvements in polymerase engineering, fluorescent chemistry, and
consensus algorithms have steadily moved PacBio from a noisy long‑read
platform toward one that can rival or exceed Illumina for small‑variant
calling in many contexts, while still providing long‑range information
for assembly, phasing, and structural‑variant
discovery\textsuperscript{25--27}.

\begin{tcolorbox}[enhanced jigsaw, rightrule=.15mm, leftrule=.75mm, colbacktitle=quarto-callout-note-color!10!white, colframe=quarto-callout-note-color-frame, toprule=.15mm, opacitybacktitle=0.6, title=\textcolor{quarto-callout-note-color}{\faInfo}\hspace{0.5em}{Figure~2.2 -- PacBio SMRT sequencing}, left=2mm, colback=white, bottomtitle=1mm, toptitle=1mm, bottomrule=.15mm, coltitle=black, titlerule=0mm, breakable, opacityback=0, arc=.35mm]

\begin{figure}[H]

\centering{

\includegraphics[width=0.75\linewidth,height=\textheight,keepaspectratio]{images/Fig_2.2_PacBio.png}

}

\caption{\label{fig-pacbio}Overview of PacBio Sequencing}

\end{figure}%

Overview of Pacific Biosciences' Single-Molecule Real-Time Sequencing
(SMRT) Technology. (A) Library preparation comprises the ligation of
hairpin adapters (yellow) to double-stranded DNA molecules (blue),
thereby creating circular molecules called `SMRTbells'. Next, a primer
(red) and a polymerase (green) are annealed to the adapter. (B)
Graphical representation of a zero mode waveguide (ZMW), a nanoscale
observation chamber. The polymerase--primer--SMRTbell complex binds to
the bottom of the ZMW through biotin--streptavidin chemistry. Note,
however, that not all ZMWs will contain a DNA molecule because the
library is loaded by diffusion. The polymerase incorporates
fluorescently labeled nucleotides emitting a fluorescent signal on
illumination from below. These signals are recorded by a camera in real
time in a process called a `movie'. \textbf{Middle box inset}:
Association formed between phospholinked nucleotide with the template in
the polymerase active site which leads to elevation of the fluorescence
signal on respective color channels. Phosphodiester bond formatted
further releases dye--linker--pyrophosphate product and diffuses it out
of ZMW, and ends the fluorescence pulse. The polymerase shifts to next
position, forms bond with nucleotide and generates next pulse. (C) In a
movie, not only the fluorescence color is registered, but also the time
between nucleotide incorporations, called the interpulse duration (IPD)
(black). The presence of an epigenetic modification, such as
6-methyladenosine (6 mA), results in a delayed IPD. (D) Multiple
`passes' of the circular library can be combined into a circular
consensus sequence (CCS) that increases in accuracy as the number of
passes increases. Accuracy is expressed as the quality value (QV). Note
that at 25 passes, the accuracy reaches 99.999\% (QV40), which is
similar to the accuracy of Illumina sequencing. At 50 passes, accuracy
can even reach 99.9999\% (QV50). The data indicated in the figure are
based on a bacterial genome run on the Sequel system with 2.1 chemistry
and 5.1 Sequel software. Image largely sourced from\textsuperscript{25}
with panels from\textsuperscript{28}.

\end{tcolorbox}

\subsection{Oxford Nanopore
Technology}\label{oxford-nanopore-technology}

Oxford nanopore sequencing takes an entirely different physical
approach. DNA or RNA molecules are threaded through protein nanopores
embedded in a membrane that separates two ionic buffer chambers; as a
stretch of nucleotides occupies the constriction of the pore, it
perturbs the ionic current in a sequence‑dependent
way\textsuperscript{25}. The resulting ``squiggle'' of current values is
decoded by basecalling software into nucleotide sequences, using
increasingly sophisticated deep‑learning models\textsuperscript{26,27}.

In principle, nanopore read length is limited only by input molecule
length and stability; with careful high‑molecular‑weight DNA extraction,
reads with N50 values around 100~kb and individual reads approaching or
exceeding one megabase have been reported\textsuperscript{25,26}. Early
nanopore chemistries had raw error rates of 20--30\%, with systematic
difficulties in homopolymers and certain motifs, but newer pore designs
(for example, R10), duplex sequencing of both strands, and improved
basecalling have raised simplex accuracies to around 99.6\% and duplex
accuracies toward 99.9\% in favorable conditions\textsuperscript{26,27}.
ONT can sequence native RNA directly as well as cDNA, and because the
ionic‑current signal is sensitive to base modifications, the same data
can be re‑analyzed to infer some DNA and RNA epigenetic
marks\textsuperscript{25--27}.

ONT's product line ranges from pocket‑sized MinION devices powered by a
laptop to large PromethION instruments with many flow cells running in
parallel, making it possible to deploy sequencing in field settings as
well as high‑throughput labs\textsuperscript{25,26}.

\begin{tcolorbox}[enhanced jigsaw, rightrule=.15mm, leftrule=.75mm, colbacktitle=quarto-callout-note-color!10!white, colframe=quarto-callout-note-color-frame, toprule=.15mm, opacitybacktitle=0.6, title=\textcolor{quarto-callout-note-color}{\faInfo}\hspace{0.5em}{Figure~2.3 -- Oxford nanopore sequencing}, left=2mm, colback=white, bottomtitle=1mm, toptitle=1mm, bottomrule=.15mm, coltitle=black, titlerule=0mm, breakable, opacityback=0, arc=.35mm]

\begin{figure}[H]

\centering{

\includegraphics[width=0.5\linewidth,height=\textheight,keepaspectratio]{images/Fig_2.3_ONT.png}

}

\caption{\label{fig-ont}Overview of ONT Sequencing}

\end{figure}%

Schematic Representation of Nanopore Technology. (A) In an Oxford
Nanopore Technologies (ONT) flow cell, two chambers (cis and trans)
filled with ionic solutions are separated by a membrane containing a
CsgG nanopore (blue; R9 chemistry). A nucleic acid (black) is
electrophoretically driven through the pore in a controlled manner owing
to the presence of a `motor' protein (green). Note that the nucleic acid
is unwound on translocation and only one strand passes through the pore.
As the DNA or RNA translocates through the pore, current shifts are
recorded in real time and are characteristic for particular k-mer
sequences. The current shifts are graphically represented in a `squiggle
plot'. (B) Typical nanopore library. Double-stranded DNA fragments
(blue) often undergo an optional DNA repair step, as single-stranded
nicks will lead to premature termination of nanopore sequencing. Then,
the extremities are processed to create suitable substrates for ligation
of adapters (brown). The adapters have 50 protruding ends to which a
`motor' protein is bound (red); this extremity will enter the pore first
and thus sequencing occurs in the 50-to-30 direction. To the other
strand of the adapter, an oligonucleotide with a cholesterol moiety
(black) is hybridized, which will tether the library molecules to the
membrane and increase the efficiency of nanopore sequencing. (C)
Different types of nanopores. The a-hemolysin pore, the MspA pore, and
the CsgG pore, which is currently being used by ONT (R9 chemistry). The
narrow `sensing regions' (SRs) of the different pores are indicated;
note that the MspA and CsgG pores have shorter SRs than a-hemolysin. As
a result, a smaller number of nucleotides contribute to the signal,
leading to more accurate base determination. Image sourced
from\textsuperscript{25}.

\end{tcolorbox}

\subsection{Synthetic and linked‑read
approaches}\label{synthetic-and-linkedread-approaches}

Between ``pure'' short‑read and true single‑molecule long‑read
platforms, several \textbf{synthetic long‑read} and \textbf{linked‑read}
methods attempted to recover long‑range information from short reads.
Examples include 10x~Genomics Chromium linked‑reads, which partitioned
high‑molecular‑weight DNA into microdroplets, barcoded fragments within
each droplet, and then sequenced the barcoded fragments on an Illumina
instrument so that short reads sharing a barcode could be associated
with the same original long molecule\textsuperscript{25}. Other
approaches (for example, Moleculo/Illumina TruSeq Synthetic Long Reads)
physically separated long molecules, amplified them, and computationally
reassembled short reads back into ``synthetic'' long
reads\textsuperscript{25}.

These methods were attractive because they worked with standard
short‑read sequencers and avoided some of the early high error rates of
PacBio and ONT. However, effective molecule lengths were constrained by
DNA quality, barcode collisions and uneven coverage complicated
analysis, and costs were not dramatically lower than true long‑read
sequencing as PacBio and ONT matured\textsuperscript{25--27}. As
consensus accuracies for HiFi and duplex nanopore reads climbed and
per‑base costs fell, many synthetic and linked‑read products were
discontinued or shifted toward niche structural‑variant and phasing
applications.

\subsection{Comparing major platforms}\label{comparing-major-platforms}

The table below summarizes key properties of the three platforms you
will encounter most often in this course. Three practical differences
matter: typical read length, error profile, and run
scale\textsuperscript{10,25,26}. Exact numbers vary with specific
instruments and chemistries, but the qualitative trade‑offs are stable.

\begin{longtable}[]{@{}
  >{\raggedright\arraybackslash}p{(\linewidth - 10\tabcolsep) * \real{0.0584}}
  >{\raggedright\arraybackslash}p{(\linewidth - 10\tabcolsep) * \real{0.1168}}
  >{\raggedright\arraybackslash}p{(\linewidth - 10\tabcolsep) * \real{0.1350}}
  >{\raggedright\arraybackslash}p{(\linewidth - 10\tabcolsep) * \real{0.1898}}
  >{\raggedright\arraybackslash}p{(\linewidth - 10\tabcolsep) * \real{0.2372}}
  >{\raggedright\arraybackslash}p{(\linewidth - 10\tabcolsep) * \real{0.2628}}@{}}
\toprule\noalign{}
\begin{minipage}[b]{\linewidth}\raggedright
Platform
\end{minipage} & \begin{minipage}[b]{\linewidth}\raggedright
Typical read length
\end{minipage} & \begin{minipage}[b]{\linewidth}\raggedright
Typical accuracy
\end{minipage} & \begin{minipage}[b]{\linewidth}\raggedright
Typical throughput per run
\end{minipage} & \begin{minipage}[b]{\linewidth}\raggedright
Key strengths
\end{minipage} & \begin{minipage}[b]{\linewidth}\raggedright
Key limitations
\end{minipage} \\
\midrule\noalign{}
\endhead
\bottomrule\noalign{}
\endlastfoot
Illumina NGS & 150--300~bp paired‑end & Very high (≈99.9\%) & 10s--1000s
of~Gb (instrument‑dependent) & Cost‑effective short reads; variant
calling; RNA‑seq quantification\textsuperscript{10} & Repeats and SVs;
limited phasing; GC bias\textsuperscript{25,26} \\
PacBio HiFi & 10--25~kb consensus reads & Very high (≥99.9\%) &
10s--100s of~Gb HiFi & Accurate long reads; assembly; SVs; phasing; some
methylation detection\textsuperscript{25,26} & Instrument cost;
relatively shorter than ultra‑long nanopore reads \\
ONT nanopore & 10--100+~kb; Mb possible & \textasciitilde99--99.9\%
(chemistry‑dependent) & 10s of~Gb (MinION) to Tb (PromethION) & Very
long reads; portable devices; direct DNA and RNA; base‑modification
detection\textsuperscript{25--27} & Higher error for some contexts;
protocol‑sensitive performance \\
\end{longtable}

Later chapters return to these platforms in the context of specific
analysis tasks such as structural‑variant discovery, genome assembly,
and full‑length transcriptomics, and to how improvements in chemistry
and analysis have kept long‑read technologies competitive with
short‑read NGS\textsuperscript{25--27}.

\section{Matching technology to biological
questions}\label{matching-technology-to-biological-questions}

Choosing a sequencing technology is fundamentally an
\textbf{experimental design} decision. The same genome can be sequenced
and analyzed in very different ways depending on whether the primary
goal is variant discovery, de novo assembly, transcriptome
characterization, or epigenetic profiling\textsuperscript{10,25,27}.

\begin{tcolorbox}[enhanced jigsaw, rightrule=.15mm, leftrule=.75mm, colbacktitle=quarto-callout-note-color!10!white, colframe=quarto-callout-note-color-frame, toprule=.15mm, opacitybacktitle=0.6, title=\textcolor{quarto-callout-note-color}{\faInfo}\hspace{0.5em}{Lab connections}, left=2mm, colback=white, bottomtitle=1mm, toptitle=1mm, bottomrule=.15mm, coltitle=black, titlerule=0mm, breakable, opacityback=0, arc=.35mm]

In \textbf{\href{Lab02.qmd}{Lab 2}} you will design a hypothetical
genome project, choosing sequencing platforms and workflows based on the
trade‑offs discussed in this chapter\textsuperscript{10,25,26}. You will
also locate your project within the five perspectives on genomics
described by Pevsner, which extend the three perspectives from Chapter
1.

Keep Lab02 in mind as you read: note which platforms and experimental
designs would best support the kinds of questions you are most
interested in. Lab02 will integrate the knowledge across Chapters 1-3
into a genome analysis narrative, so start thinking about this
assignment now and it will be due after we finish Chapter 5.

\end{tcolorbox}

In practice, many modern projects use \textbf{hybrid strategies}: long
reads provide the backbone for assembly or SV discovery, while
inexpensive short reads polish base‑level accuracy, improve
small‑variant calling, and support large sample
sizes\textsuperscript{10,25,26}.

\begin{tcolorbox}[enhanced jigsaw, rightrule=.15mm, leftrule=.75mm, colbacktitle=quarto-callout-note-color!10!white, colframe=quarto-callout-note-color-frame, toprule=.15mm, opacitybacktitle=0.6, title=\textcolor{quarto-callout-note-color}{\faInfo}\hspace{0.5em}{Lab connections}, left=2mm, colback=white, bottomtitle=1mm, toptitle=1mm, bottomrule=.15mm, coltitle=black, titlerule=0mm, breakable, opacityback=0, arc=.35mm]

\begin{itemize}
\tightlist
\item
  Labs 4, 7, and 9 (see \href{appendixC.qmd}{Appendix C}) use Illumina
  short‑read data and the sacCer3 reference genome to illustrate how
  platform choice, coverage, and error profiles affect QC, alignment,
  and variant calling.
\item
  Later projects (see \href{appendixB.qmd}{Appendix B}) highlight how
  long‑read and hybrid assemblies from vertebrates support comparative
  genomics questions about gene families such as PRDM9.
\end{itemize}

\end{tcolorbox}

\subsection{Alignment versus assembly in
practice}\label{alignment-versus-assembly-in-practice}

Chapter 1 introduced alignment and de novo assembly as core
computational operations\textsuperscript{1,10}. Here, we emphasize how
platform choice and experimental goals determine which operation is
primary.

\subsubsection{Alignment to a reference}\label{alignment-to-a-reference}

For organisms with high‑quality reference genomes, most short‑read
workflows begin by \textbf{aligning} reads to the reference using
algorithms such as BWA or Bowtie2, producing SAM/BAM files that record
alignment positions and quality metrics\textsuperscript{10}. Alignment
is the backbone of variant calling, read‑depth analysis, and many
RNA‑seq and ChIP‑seq methods; its performance depends on read length,
error rate, and the divergence between sample and
reference\textsuperscript{10}.

Long reads can also be aligned, and specialized mappers have been
developed to handle higher error rates and complex structural variation.
In lecture you saw how long‑read alignments make it easier to resolve
large insertions, tandem repeats, or transposable element insertions
that short reads may miss entirely\textsuperscript{25,26}.

\subsubsection{De novo assembly and graph‑based
references}\label{de-novo-assembly-and-graphbased-references}

When no good reference exists---or when the reference is known to be
incomplete or structurally divergent---\textbf{de novo assembly}
reconstructs the genome directly from reads. Classical short‑read
assemblers typically use de Bruijn graph formulations, breaking reads
into k‑mers and finding Eulerian paths that traverse the graph to
produce contigs\textsuperscript{10}. Assemblies from short reads alone
often fragment into many contigs because repeats longer than the read
length create ambiguous graph structures\textsuperscript{25}.

Long‑read assemblers instead work with overlap graphs or string graphs
that leverage long overlaps between reads, greatly reducing ambiguity
and producing much larger contigs\textsuperscript{25,26}. However, even
high‑quality long‑read assemblies typically need additional information
to be oriented and ordered into chromosome‑scale scaffolds. The field is
also moving toward \textbf{graph‑based pangenomes}, where multiple
haplotypes or assemblies are combined into a single graph structure that
better represents population‑level variation than any single linear
reference\textsuperscript{26}.

Later in the book, a guest lecture and accompanying chapter will dig
into assembly algorithms and pangenome graphs in more detail; at this
point the key idea is that sequencing technologies and algorithms
co‑evolve, and practical decisions about assembly strategy must take
both into account.

\subsection{Assessing assemblies: the three
C's}\label{assessing-assemblies-the-three-cs}

Genome assemblies are often evaluated using three high‑level criteria
sometimes called the ``three C's'': \textbf{contiguity},
\textbf{completeness}, and
\textbf{correctness}\textsuperscript{10,25,26}. This will be covered in
more depth later when we discuss genome assembly in depth, but for the
sake of helping you with the reading this week, I have provided an
overview of basic assessment criteria that are discussed in the paper.

\subsubsection{Contiguity}\label{contiguity}

Contiguity describes how fragmented an assembly is. The most widely
reported metric is \textbf{N50}, the contig (or scaffold) length such
that all contigs of that length or longer contain 50\% of the assembled
bases\textsuperscript{10,25}. Assemblies with high contig N50 and
scaffold N50 values are generally more useful for downstream analyses
and annotation.

However, N50 alone can be misleading. Aggressive assemblers may
artificially join contigs through repeats or low‑confidence overlaps,
inflating N50 at the cost of mis‑assemblies that misrepresent the true
genome structure\textsuperscript{25,26}. Later labs will show you how to
plot contig length distributions and use multiple metrics to judge
contiguity.

\subsubsection{Completeness}\label{completeness}

Completeness asks how much of the expected gene content is represented
in the assembly. A widely used tool is \textbf{BUSCO} (Benchmarking
Universal Single‑Copy Orthologs), which examines the presence and
integrity of conserved single‑copy genes expected for a given
clade\textsuperscript{29}. High BUSCO scores (e.g., ≥95\% complete)
suggest that most core genes are present; low scores can indicate
missing regions, assembly fragmentation, or annotation
issues\textsuperscript{10,26}.

\subsubsection{Correctness}\label{correctness}

Correctness encompasses both base‑level accuracy and larger‑scale
structural validity. Common approaches include:

\begin{itemize}
\tightlist
\item
  Mapping reads back to the assembly and examining mismatch and indel
  rates, discordant paired‑end mappings, and depth anomalies that
  suggest mis‑joins or collapses\textsuperscript{10}.\\
\item
  Comparing assemblies built from different technologies (e.g., long
  reads vs short reads) to identify consistent features and potential
  errors\textsuperscript{25,26}.\\
\item
  Using independent long‑range data such as Hi‑C, optical maps, or
  Strand‑seq to confirm or correct scaffold orientation and
  order\textsuperscript{26}.
\end{itemize}

\begin{tcolorbox}[enhanced jigsaw, rightrule=.15mm, leftrule=.75mm, colbacktitle=quarto-callout-note-color!10!white, colframe=quarto-callout-note-color-frame, toprule=.15mm, opacitybacktitle=0.6, title=\textcolor{quarto-callout-note-color}{\faInfo}\hspace{0.5em}{Connection to Assigned Reading}, left=2mm, colback=white, bottomtitle=1mm, toptitle=1mm, bottomrule=.15mm, coltitle=black, titlerule=0mm, breakable, opacityback=0, arc=.35mm]

In the paper you were assigned for this week\textsuperscript{3}, they
generated several novel genome assemblies. Look closely at the methods
they used. Which sequencing technologies did they use to generate the
genome assembly? Which metrics did they use for quality assessment of
the genome? Do you agree with their conclusions regarding quality of the
genomes they generated?

\end{tcolorbox}

\subsection{Long‑range scaffolding: 3D genome methods and optical
maps}\label{longrange-scaffolding-3d-genome-methods-and-optical-maps}

Even with long reads, de novo assemblies often consist of contigs rather
than full chromosomes. Long‑range scaffolding methods exploit physical
properties of chromosomes to order and orient
contigs\textsuperscript{25,26}.

\subsubsection{Hi‑C and other conformation capture
assays}\label{hic-and-other-conformation-capture-assays}

High‑throughput chromosome conformation capture (\textbf{Hi‑C}) measures
which regions of the genome are physically close to each other in the
nucleus by crosslinking chromatin, digesting and ligating DNA, and
sequencing junctions between interacting loci\textsuperscript{26}.
Because loci on the same chromosome interact more frequently than loci
on different chromosomes, and nearby loci interact more frequently than
distant ones, Hi‑C contact maps provide powerful information for
assigning contigs to chromosomes and arranging them along each
chromosome\textsuperscript{26}.

\subsubsection{Optical maps and
Strand‑seq}\label{optical-maps-and-strandseq}

\textbf{Optical mapping} labels specific sequence motifs along
ultra‑long DNA molecules, linearizes them in nanochannels, and images
the pattern of labels; consensus maps can then be used to validate and
scaffold assemblies\textsuperscript{26}. \textbf{Strand‑seq}, a
single‑cell sequencing method that preserves strand orientation, can
help distinguish homologous chromosomes, identify inversions, and phase
assemblies\textsuperscript{26}. Together with long reads, these
technologies enabled the recent telomere‑to‑telomere (T2T) human genome
assembly and are increasingly used in high‑quality assemblies for other
species\textsuperscript{26}.

\subsection{Experimental design and
coverage}\label{experimental-design-and-coverage}

All sequencing experiments must grapple with \textbf{coverage}, roughly
defined as:

Coverage = \(\frac{N \times L}{G}\) where \(G\) = genome size.

For Illumina WGS, 30× coverage is a common target for germline variant
discovery; higher coverage may be required for somatic mutation
detection or low‑frequency variants in mixed
samples\textsuperscript{10}. For long‑read assembly projects, 20--40×
coverage in HiFi or high‑quality nanopore reads is often recommended,
sometimes supplemented with additional short reads for
polishing\textsuperscript{25,26}. This metric will be covered in more
depth in the next chapter.

Because sequencing costs are typically proportional to total data
volume, experimental design must balance per‑sample coverage against the
number of samples and platforms used. Later chapters will revisit these
trade‑offs with concrete R code examples inspired by workflows in
\emph{Computational Genomics with R}\textsuperscript{12}, using simple
functions to explore how changes in coverage, sample size, and error
rate impact variant discovery power and assembly quality.

\begin{tcolorbox}[enhanced jigsaw, rightrule=.15mm, leftrule=.75mm, colbacktitle=quarto-callout-tip-color!10!white, colframe=quarto-callout-tip-color-frame, toprule=.15mm, opacitybacktitle=0.6, title=\textcolor{quarto-callout-tip-color}{\faLightbulb}\hspace{0.5em}{Box 2.1 -- From sequencers to full‑stack workflows}, left=2mm, colback=white, bottomtitle=1mm, toptitle=1mm, bottomrule=.15mm, coltitle=black, titlerule=0mm, breakable, opacityback=0, arc=.35mm]

In the early days of NGS, data analysis pipelines were largely assembled
by individual labs from a mosaic of third‑party tools: separate programs
for basecalling, demultiplexing, alignment, variant calling, and
visualization, each with its own configuration and output
formats\textsuperscript{10}. Reproducing an analysis often required
extensive documentation and command‑line literacy.

Over the past decade, major sequencing vendors have increasingly offered
\textbf{end‑to‑end workflow ecosystems}. Examples include Illumina's
cloud‑hosted pipelines, PacBio's SMRT Link suite for HiFi data, and
ONT's MinKNOW plus associated basecallers and analysis
apps\textsuperscript{26}. These platforms bundle basecalling, alignment,
variant calling, and sometimes annotation or structural variant
detection into graphical interfaces and preconfigured workflows,
lowering the barrier to entry for non‑specialists and making routine
clinical and agricultural genomics feasible in more
settings\textsuperscript{13,26}.

For bioinformatics, this trend is double‑edged. On one hand,
user‑friendly workflows reduce technical hurdles and standardize
analyses across labs; on the other, they can obscure algorithmic
assumptions, limit flexibility, and make it harder to debug or extend
pipelines\textsuperscript{10,12}. A core goal of this course is to help
you operate at both levels: you should be able to use vendor pipelines
productively when appropriate, but also understand enough about
alignment, assembly, and variant calling algorithms to interpret their
outputs critically and, when needed, construct your own reproducible
workflows.

\textbf{Connection to Assigned Reading}: In the paper you were assigned
for this week\textsuperscript{3}, they generated several novel genome
assemblies. Look closely again at their methods, but this time focus on
the bioinformatics software they used to generate their assemblies and
perform QC. What are the names of the tools? When were these first
described? Which are ``old'' established tools? Which are newer tools
that are recently developed? Are any unique to the seqeuncing platform
they used?

\end{tcolorbox}

\begin{center}\rule{0.5\linewidth}{0.5pt}\end{center}

\section{Practice Problem‑sets}\label{practice-problemsets}

\begin{enumerate}
\def\labelenumi{\arabic{enumi}.}
\item
  \textbf{Coverage calculations in R}: Write an R function that takes
  genome size (bp), planned read length, and number of reads as input
  and returns expected coverage. Explore how coverage changes as you
  vary read length from 150 bp to 20 kb while holding total sequenced
  bases constant, and interpret what this means for assembly contiguity
  versus variant calling in a 3 Gb genome\textsuperscript{10,25}.
\item
  \textbf{Platform choice for a real project}: Choose a published genome
  study of your organism of interest (or a related species). Based on
  the questions the authors asked, justify whether their choice of
  sequencing platform(s) and coverage seems appropriate. Propose at
  least one alternative design using a different mix of short‑ and
  long‑read technologies, and explain what new information it might
  reveal\textsuperscript{25--27}.
\item
  \textbf{Interpreting N50 and BUSCO}: You are given summary statistics
  for two assemblies of the same plant genome: Assembly A has higher
  contig N50 but lower BUSCO completeness; Assembly B has lower N50 but
  higher BUSCO and fewer structural inconsistencies when compared to an
  optical map. Which assembly you would choose as the main reference and
  what additional data or analyses you would request before making a
  final decision\textsuperscript{26,29}?
\item
  \textbf{Alignment versus assembly}: Consider a project to study
  structural variation in a non‑model fish species with only a draft
  short‑read assembly available. Outline two analysis strategies: (a)
  alignment‑centric using new short‑read WGS data, and (b)
  assembly‑centric using long reads and Hi‑C. For each, list at least
  three strengths and three limitations, and suggest which approach
  would better address large rearrangements and repeat
  expansions\textsuperscript{10,25,26}.
\item
  \textbf{Workflow transparency and vendor ecosystems}: Imagine your lab
  gains access to a turnkey sequencing‑plus‑analysis platform that
  produces variant calls through a graphical interface with minimal user
  control. Consider how you would develop a lab policy to ensure
  analysis transparency and reproducibility when using this system (for
  example, documenting software versions, exporting intermediate files,
  or reanalyzing subsets of data with open‑source
  tools)\textsuperscript{10,12,13}.
\end{enumerate}

\section{Reflection questions}\label{reflection-questions-1}

\begin{enumerate}
\def\labelenumi{\arabic{enumi}.}
\item
  \textbf{Data permanence and participant expectations}: Sequencing
  costs continue to fall, but data storage and long‑term stewardship
  remain challenging\textsuperscript{10,13}. When human genomes are
  sequenced today, they may be reanalyzed decades later with new
  algorithms and linked to other data sources. How should this
  possibility influence what participants are told at consent, and who
  should decide how data can be reused?
\item
  \textbf{Algorithmic opacity}: As vendor pipelines and
  machine‑learning--based basecallers become standard, fewer users will
  understand exactly how raw signals are transformed into sequences and
  variants\textsuperscript{12,26}. Is it acceptable to rely on
  ``black‑box'' algorithms in clinical or conservation genomics if they
  are empirically well validated? What safeguards or transparency
  requirements, if any, should be in place?
\item
  \textbf{Polymerase proofreading and design trade‑offs}: Many native
  DNA polymerases have a proofreading exonuclease domain that increases
  fidelity by removing incorrectly incorporated bases, but Illumina's
  sequencing‑by‑synthesis chemistry uses engineered polymerases without
  this proofreading activity so that all strands in a cluster advance in
  lock‑step each cycle. How might proofreading interfere with
  synchronized sequencing, and what does this tell you about the
  trade‑off between \emph{natural} accuracy mechanisms and
  \emph{engineered} control of the chemistry in high‑throughput
  platforms?
\item
  \textbf{What counts as a ``read''?}: PacBio HiFi consensus reads are
  built from multiple noisy passes around a circular template, whereas
  ONT reads are decoded directly from continuous current traces, and
  short‑read Illumina reads are aggregates of thousands of synchronized
  molecules in a cluster. In what sense are these three signals all
  called ``reads''? How does thinking carefully about what a read
  represents at the chemistry/signal level change how you interpret
  downstream metrics such as coverage, mapping quality, or read depth?
\item
  \textbf{Platform choice and hidden biases}: Illumina, PacBio HiFi, and
  ONT now all claim per‑base accuracies near or above 99\%, but their
  errors arise from different physical processes. For example, the
  chapter highlighted GC‑bias in short‑read PCR‑based methods and
  length‑ or context‑dependent biases in long‑read platforms. Imagine
  two groups sequencing the same genome: one with Illumina only, the
  other with a mix of PacBio HiFi and ONT. Without doing any
  computation, list several ways in which their biological conclusions
  could legitimately differ because of platform chemistry and error
  modes alone. How might you design an experiment or analysis to reveal,
  quantify, or mitigate those differences?
\end{enumerate}

\begin{center}\rule{0.5\linewidth}{0.5pt}\end{center}

\section{References}\label{references-2}

\chapter{File Formats and Quality Control in Genome
Workflows}\label{file-formats-and-quality-control-in-genome-workflows}

\section{Introduction to Genome File Formats}\label{sec-ch3}

Next-generation sequencing (NGS) generates massive datasets that flow
through standardized workflows from raw reads to biological
insights\textsuperscript{10}. Understanding file formats is essential
for experimental design, as each format encodes metadata critical for
downstream analysis decisions.

Figure Figure~\ref{fig-ngs-workflow} shows the standard NGS pipeline,
where FASTQ files capture raw sequencing data, SAM/BAM store alignments,
and VCF summarizes variants.

\begin{tcolorbox}[enhanced jigsaw, rightrule=.15mm, leftrule=.75mm, colbacktitle=quarto-callout-note-color!10!white, colframe=quarto-callout-note-color-frame, toprule=.15mm, opacitybacktitle=0.6, title=\textcolor{quarto-callout-note-color}{\faInfo}\hspace{0.5em}{Figure 3.1 -- NGS Workflow}, left=2mm, colback=white, bottomtitle=1mm, toptitle=1mm, bottomrule=.15mm, coltitle=black, titlerule=0mm, breakable, opacityback=0, arc=.35mm]

\begin{figure}[H]

\centering{

\includegraphics[width=0.25\linewidth,height=\textheight,keepaspectratio]{images/Fig_3_1_ngs_workflow.png}

}

\caption{\label{fig-ngs-workflow}Standard NGS workflow}

\end{figure}%

Workflow for next-generation sequence experiments: from experimental
design to data analysis with the standard file formats for each step.
Each file format is described in depth in this chapter and in the table
below. Created with BioRender.com.

\end{tcolorbox}

Additionally, visualizing next-generation sequence data as well as
genomic arithmetic can be done using genome viewer softwares, such as
Integrated Genome Viewer (IGV), which we will cover later in this
course. Another common file format in genomics is the BED format, which
can be used with genome views, to perform various genome analyses, and
to make custom tracks in genome browsers.

The goal of this chapter is to familiarize you with these four major
file formats, which are reviewed in this chapter:

\subsection{Common Genome Analysis File
Formats}\label{common-genome-analysis-file-formats}

\begin{longtable}[]{@{}
  >{\raggedright\arraybackslash}p{(\linewidth - 6\tabcolsep) * \real{0.2394}}
  >{\raggedright\arraybackslash}p{(\linewidth - 6\tabcolsep) * \real{0.2394}}
  >{\raggedright\arraybackslash}p{(\linewidth - 6\tabcolsep) * \real{0.2394}}
  >{\raggedright\arraybackslash}p{(\linewidth - 6\tabcolsep) * \real{0.2817}}@{}}
\toprule\noalign{}
\begin{minipage}[b]{\linewidth}\raggedright
Format
\end{minipage} & \begin{minipage}[b]{\linewidth}\raggedright
Acronym
\end{minipage} & \begin{minipage}[b]{\linewidth}\raggedright
Description
\end{minipage} & \begin{minipage}[b]{\linewidth}\raggedright
Primary Toolkit(s)
\end{minipage} \\
\midrule\noalign{}
\endhead
\bottomrule\noalign{}
\endlastfoot
FASTA & FASTA & Text-based format for representing nucleotide or protein
sequences with sequence identifiers & Biopython, BioPerl, seqtk,
samtools faidx, EMBOSS \\
FASTQ & FASTQ & Extension of FASTA that includes per-base quality scores
from sequencing instruments & FastQC, Trimmomatic, cutadapt, seqtk \\
SAM/BAM & Sequence Alignment/Map (Binary) & Stores aligned sequence
reads to a reference genome; BAM is compressed binary version &
SAMtools, Picard, GATK, sambamba \\
VCF/BCF & Variant Call Format (Binary) & Describes genetic variations
including SNPs, indels, and structural variants; BCF is binary version &
BCFtools, GATK, VCFtools, vcflib, R: VariantAnnotation, vcfR \\
BED & Browser Extensible Data & Tab-delimited format for genomic
intervals and annotations (chromosome, start, end) & BEDtools, UCSC
tools, pybedtools, R: GenomicRanges, rtracklayer \\
\end{longtable}

\subsubsection{Additional Common Formats covered
later}\label{additional-common-formats-covered-later}

\begin{longtable}[]{@{}
  >{\raggedright\arraybackslash}p{(\linewidth - 6\tabcolsep) * \real{0.2394}}
  >{\raggedright\arraybackslash}p{(\linewidth - 6\tabcolsep) * \real{0.2394}}
  >{\raggedright\arraybackslash}p{(\linewidth - 6\tabcolsep) * \real{0.2394}}
  >{\raggedright\arraybackslash}p{(\linewidth - 6\tabcolsep) * \real{0.2817}}@{}}
\toprule\noalign{}
\begin{minipage}[b]{\linewidth}\raggedright
Format
\end{minipage} & \begin{minipage}[b]{\linewidth}\raggedright
Acronym
\end{minipage} & \begin{minipage}[b]{\linewidth}\raggedright
Description
\end{minipage} & \begin{minipage}[b]{\linewidth}\raggedright
Primary Toolkit(s)
\end{minipage} \\
\midrule\noalign{}
\endhead
\bottomrule\noalign{}
\endlastfoot
GFF/GTF & General Feature Format / Gene Transfer Format & Annotation
formats describing genomic features like genes, transcripts, and exons &
gffread, bedtools, AGAT \\
CRAM & CRAM & Highly compressed alternative to BAM with reference-based
compression & SAMtools, htslib \\
bigWig/bigBed & - & Indexed binary formats for dense continuous data
(bigWig) and interval data (bigBed) & UCSC tools (bigWigToBedGraph),
deepTools \\
\end{longtable}

\begin{tcolorbox}[enhanced jigsaw, rightrule=.15mm, leftrule=.75mm, colbacktitle=quarto-callout-note-color!10!white, colframe=quarto-callout-note-color-frame, toprule=.15mm, opacitybacktitle=0.6, title=\textcolor{quarto-callout-note-color}{\faInfo}\hspace{0.5em}{Lab connections}, left=2mm, colback=white, bottomtitle=1mm, toptitle=1mm, bottomrule=.15mm, coltitle=black, titlerule=0mm, breakable, opacityback=0, arc=.35mm]

This chapter's file formats and QC concepts feed directly into the first
two labs:

\begin{itemize}
\tightlist
\item
  \textbf{\href{Lab01.qmd}{Lab 1}} asks you to identify how published
  genome analysis papers report assembly metrics and QC, and to relate
  those choices back to FASTQ, BAM, and VCF concepts.
\item
  \textbf{\href{Lab02.qmd}{Lab 2}} asks you to design a genome project
  that makes concrete decisions about coverage, platforms, and
  workflows---decisions that will be encoded in FASTQ, BAM, VCF, and BED
  files.
\end{itemize}

\end{tcolorbox}

\section{FASTA: The Foundation of Sequence
Data}\label{fasta-the-foundation-of-sequence-data}

Last chapter discussed Sanger sequencing and the generation of
\href{https://www.labxchange.org/library/pathway/lx-pathway:ff079975-ac71-418a-ac93-2e14b056f539/items/lb:LabXchange:22c08d85:html:1/58303?t=8c80f5c3-6a6b-49cf-9d3a-87f332b26c42}{\textbf{chromatograms}}---the
colorful trace files produced by these sequencing instruments. While
chromatograms preserve valuable quality information about individual
sequencing reactions, they are impractical for computational analysis
and database storage. This challenge led to the development of
text-based sequence formats that could efficiently represent biological
sequences while remaining human-readable and computationally tractable.
The most fundamental and enduring of these formats is FASTA.

The FASTA format was developed in \textbf{1985} by \textbf{William R.
Pearson} and \textbf{David J. Lipman} at the National Institutes of
Health\textsuperscript{30,31}. Originally created for their FASTP
protein sequence similarity search program, the format was designed to
handle the exponentially growing genetic databases of the 1980s when
computer memory and processing speed were severely limited. The name
``FASTA'' stands for ``FAST-All,'' reflecting the program's ability to
handle both protein (FAST-P) and nucleotide (FAST-N) sequences.

\subsection{FASTA File Structure}\label{fasta-file-structure}

A FASTA file consists of one or more sequence records, each with two
components:

\textbf{1. Definition line (header):}

\begin{itemize}
\tightlist
\item
  Begins with \texttt{\textgreater{}} character
\item
  First word (no spaces) is the \textbf{sequence identifier}
\item
  Remaining text is \textbf{optional description}
\item
  Must be unique within the file
\end{itemize}

\textbf{2. Sequence data:}

\begin{itemize}
\tightlist
\item
  Uses standard IUPAC single-letter codes
\item
  Typically wrapped to 60-80 characters per line (for readability)
\item
  Case-insensitive (though lowercase sometimes indicates masked/repeat
  regions)
\item
  Whitespace and numbers are ignored
\end{itemize}

FASTA files store only sequence data, not quality scores. Additionally,
compression with \texttt{gzip} is common for large genome FASTA files,
and many bioinformatics tools can read compressed FASTA files directly
(\texttt{.fasta.gz} or \texttt{.fa.gz}).

\subsection{From Individual Sequences to Genomic
Scale}\label{from-individual-sequences-to-genomic-scale}

While FASTA files are plain text and viewable in any text editor,
specialized tools enhance their utility. In the next chapter on Genome
Browsers, we will install \textbf{AliView}, a cross-platform tool that
can visualize FASTA files along with other sequence formats. AliView is
particularly useful for viewing multiple sequence alignments and
comparing sequences side-by-side.

The evolution from Sanger sequencing chromatograms to FASTA files to
modern high-throughput formats reflects the scaling challenges in
genomics. Where a single Sanger reaction produced one chromatogram file
(\textasciitilde1 KB) representing \textasciitilde800 bp of sequence,
modern Illumina sequencing produces billions of reads requiring formats
like FASTQ (for raw reads) and BAM (for alignments). Yet FASTA remains
the standard for reference genomes, gene sequences, and protein
databases---a testament to the format's elegance and utility established
nearly 40 years ago.

\section{FASTQ: Raw Sequencing Reads}\label{fastq-raw-sequencing-reads}

When genomic DNA is sequenced, raw image files are typically generated
that are used to interpret which nucleotide is called in a given
fragment or ``read''. For most users, the FASTQ files (rather than the
underlying image files) represent the raw reads from which other
analyses are performed. Like the FASTA format, the FASTQ format includes
a sequence string, consisting of the nucleotide sequence of each read.
FASTQ also includes an associated quality score for every base.

FASTQ extends FASTA by adding per-base quality scores, forming the
foundation of NGS analysis. Each record spans four
lines\textsuperscript{10,32}:

\begin{enumerate}
\def\labelenumi{\arabic{enumi}.}
\tightlist
\item
  \textbf{Header} (\texttt{@}): Encodes rich metadata about sequencing
  conditions
\item
  \textbf{Sequence}: Raw nucleotides (A/C/G/T/N)
\item
  \textbf{Separator} (\texttt{+}): Optional repeat of header
\item
  \textbf{Quality scores}: ASCII-encoded Phred values
\end{enumerate}

Each FASTQ file has records that are in blocks four lines long. The
first line, beginning with the @ symbol, identifies the record. It may
optionally include information about the sequence length or the machine
used for sequencing. The second line has the sequence (in upper case),
including the nucleotides G, A, T, C, and (as is the case here in the
second position) there may be an N for unknown nucleotide. The third
line begins with the + symbol and typically contains just that character
(as in this case), or it can have more information. The fourth line
includes the quality scores corresponding to every base. Each quality
score is assigned a single character, and the entire quality score
string must equal the length of the sequence string.

\begin{tcolorbox}[enhanced jigsaw, rightrule=.15mm, leftrule=.75mm, colbacktitle=quarto-callout-tip-color!10!white, colframe=quarto-callout-tip-color-frame, toprule=.15mm, opacitybacktitle=0.6, title=\textcolor{quarto-callout-tip-color}{\faLightbulb}\hspace{0.5em}{Sample FASTQ entry}, left=2mm, colback=white, bottomtitle=1mm, toptitle=1mm, bottomrule=.15mm, coltitle=black, titlerule=0mm, breakable, opacityback=0, arc=.35mm]

A four-line FASTQ entry showing a single read of a genome sequenced
sample:

\begin{verbatim}
@SRR1693691.1 HWI-ST591:122:D19EFACXX:1:1101:1948:2137 length=75
CTATAGATCGACGCATATGTATTATCTTAAATAGGTTGCAATATATTATTATTTTTTAACATTTATTGCGATGCT
+SRR1693691.1 HWI-ST591:122:D19EFACXX:1:1101:1948:2137 length=75
?@@DFDDFHA<CDHGIGJIHGHJIGIJJGIIIICHGHJJIGIHGHIJHIJGIJJJIIGGGIIGJGIDEDH<?BEE
\end{verbatim}

This represents a single Illumina sequencing read of 75 base pairs. The
first line gives the header, the next the sequence, the next repeats the
header, and the last line is hte ASCII-encode Phred quality scores for
each position.

\end{tcolorbox}

\subsection{FASTQ Header Format}\label{fastq-header-format}

FASTQ headers contain important metadata about each sequencing read and
typically uses \textbf{colons (\texttt{:})} as delimiters. However, the
exact format has evolved over time as Illumina updated their sequencing
platforms. Understanding these variations is important when working with
data from different sources so that reads can be paired correctly,
quality control tracks flowcell lanes properly, and to ensure tool
compatibility. Generally, Illumina headers contain important metadata
about the sequencing run, such as instrument, run, flowcell, lane,
etc.\textsuperscript{10}.

For example, imagine the first line of the FastQ file looks something
like this:

\textbf{``@HS2000-306\_201:6:1204:19922:79127/1''}

This header string is an example of the modern Illumina FASTQ header
fields (Casava 1.8+) and reveals the following metadata of the
sequencing reads in the corresponding file:

\begin{longtable}[]{@{}ll@{}}
\toprule\noalign{}
Information in header & Interpretation \\
\midrule\noalign{}
\endhead
\bottomrule\noalign{}
\endlastfoot
HS2000-306\_201 & the unique instrument name \\
6 & flowcell lane \\
1204 & tile number within the flowcell lane \\
19922 & x-coordinate of the cluster within the tile \\
79127 & y-coordinate of the cluster within the tile \\
1 & the member of a pair, 1 or 2 (paired-end) \\
\end{longtable}

When data is downloaded from NCBI's Sequence Read Archive (SRA), or from
older Illumina platforms, the header format differs, as you might notice
in the example FASTQ entry above.

\textbf{SRA Format in example above:}

\begin{itemize}
\tightlist
\item
  Begins with \textbf{SRA accession number} (e.g., SRR1693691) followed
  by the original Illumina header
\item
  The SRA accession is followed by a period and read number (\texttt{.1}
  or \texttt{.2} for paired reads) with a space before the instrument
  name instead of a colon
\item
  May include additional metadata like \texttt{length=75} at the end
\item
  Original Illumina information preserved after the SRA identifier
\item
  Older formats may also include a \textbf{flowcell ID} (e.g.,
  D19EFACXX) not present in modern format
\end{itemize}

The specific metadata encoded in the identifier varies by sequencing
platform and processing pipeline, but the fundamental purpose remains
the same: uniquely identifying each sequencing read and providing
traceability to its origin on the sequencer.

\subsection{Phred Quality Scores}\label{phred-quality-scores}

PHRED scores were introduced in 1988 by Phil Green and colleagues to
describe base quality scores from Sanger sequencing. This definition is
used in the Sanger FASTQ format. Characters are stored as ASCII
printable characters 33--126 (i.e., with an ASCII offset of 33) so the
range of possible quality scores is 0 to 93. These are listed in Figure
9.8. A value of 93 corresponds to a probability of 10\^{}--9.3 that a
base call occurred by chance, that is, the read is extremely likely to
be correct. At Q30 there is a 1:1000 (i.e., 10--3) error rate, a
threshold that is often set as a minimum for high-quality reads.

Quality scores use Phred scale: \(Q = -10\log_{10}(P_e)\) where \(P_e\)
is error probability\textsuperscript{10}. ASCII encoding adds +33 offset
for printable characters:

\begin{tcolorbox}[enhanced jigsaw, rightrule=.15mm, leftrule=.75mm, colbacktitle=quarto-callout-note-color!10!white, colframe=quarto-callout-note-color-frame, toprule=.15mm, opacitybacktitle=0.6, title=\textcolor{quarto-callout-note-color}{\faInfo}\hspace{0.5em}{Figure 3.2 -- Phred-to-ASCII conversion}, left=2mm, colback=white, bottomtitle=1mm, toptitle=1mm, bottomrule=.15mm, coltitle=black, titlerule=0mm, breakable, opacityback=0, arc=.35mm]

\begin{figure}[H]

\centering{

\includegraphics[width=0.75\linewidth,height=\textheight,keepaspectratio]{images/Figure_9_8_Pevsner.jpg}

}

\caption{\label{fig-phred-ascii}Phred-to-ASCII conversion (from Pevsner
9.8)}

\end{figure}%

FASTQ quality scores use ASCII coding symbols. The chart shows ASCII
symbols (Char columns indicate characters) corresponding to decimal
notation 1--127. The characters for 0--31 (which are not used for
printing) are not shown. Decimal 31 is a space, and subsequent
characters (decimal 33--126) are used to represent base quality scores
in the Sanger FASTQ format. For example, if base calls GATC have quality
scores of 28, 30, 25, and 31 then the symbols for their quality scores
are =?:@. Adapted from http://www.lookuptables.com.
Source\textsuperscript{10}.

\end{tcolorbox}

Because genome sequencing has so much data, the cutoffs for Phred scores
require added stringency. As is shown in the table below, a typical
quality score of 20, which is recommended in most pipelines, ensures a
99\% accuracy. Still, this equates to \textasciitilde100 million
sequencing errors in a standard dataset of 10 billion base pairs of DNA.
Later, in Chapter 10, we will discuss variant filtering in more detail
where additional parameters are used to filter data beyond simply Phred
Quality scores of individual base calls.

\begin{longtable}[]{@{}
  >{\raggedright\arraybackslash}p{(\linewidth - 4\tabcolsep) * \real{0.3288}}
  >{\raggedright\arraybackslash}p{(\linewidth - 4\tabcolsep) * \real{0.3425}}
  >{\raggedright\arraybackslash}p{(\linewidth - 4\tabcolsep) * \real{0.3288}}@{}}
\toprule\noalign{}
\begin{minipage}[b]{\linewidth}\raggedright
Phred Quality Score
\end{minipage} & \begin{minipage}[b]{\linewidth}\raggedright
Probability of Incorrect Base Call
\end{minipage} & \begin{minipage}[b]{\linewidth}\raggedright
Base Call Accuracy
\end{minipage} \\
\midrule\noalign{}
\endhead
\bottomrule\noalign{}
\endlastfoot
10 & 1 in 10 & 90\% \\
20 & 1 in 100 & 99\% \\
30 & 1 in 1,000 & 99.9\% \\
40 & 1 in 10,000 & 99.99\% \\
50 & 1 in 100,000 & 99.999\% \\
60 & 1 in 1,000,000 & 99.9999\% \\
\end{longtable}

\section{Quality Control: FastQC
Interpretation}\label{quality-control-fastqc-interpretation}

Several packages are useful to assess the quality of the sequence data.
There are many types of sequencing errors:

\begin{itemize}
\item
  Errors typically increase as a function of read length. For example,
  for Illumina technology, with each increasing cycle it becomes
  increasingly difficult to differentiate the signal to noise ratio for
  which nucleotide has been incorporated.
\item
  Errors occur as a function of GC content.
\item
  Errors may occur in homopolymer positions. Both pyrosequencing and Ion
  Torrent technologies are susceptible to this.
\end{itemize}

FastQC software (\textsuperscript{33}) provides quality control
statistics. Data are imported from FASTQ files, and are then analyzed
for qulality. The output includes the following:

\begin{itemize}
\item
  \emph{Basic statistics}. This includes information such as the range
  of sequence lengths and the percent GC content.
\item
  \emph{Per base sequence quality}. This shows quality scores (y axis)
  versus base position (x axis).
\item
  \emph{Per sequence quality scores} shows the number of sequences (y
  axis) versus mean sequence quality (Phred score; x axis).
\item
  \emph{Overrepresented sequences}. You can copy these sequences and
  save them to a text file. You can then try BLAT (particularly if you
  know the species of the overrepresented sequences), BLASTN (to search
  for matches across species), batch BLAST, or other tools.
\item
  \emph{k-mer content}: this shows data for a series of 5-mers (strings
  of 5 nucleotides) plotted by relative enrichment (y axis) versus
  position in read (in base pairs; x axis). The expected 5-mer frequency
  (determined from the base composition of the entire sequence) can be
  compared to the observed frequency. k-mer counts may be reduced (e.g.,
  when poor quality reads reduce the counts for duplicated sequences) or
  enriched (e.g., when 5-mers are overrepresented at particular
  locations along a read, such as in the vicinity of a tag added to the
  5′ end of sequencing reads).
\end{itemize}

\begin{tcolorbox}[enhanced jigsaw, rightrule=.15mm, leftrule=.75mm, colbacktitle=quarto-callout-note-color!10!white, colframe=quarto-callout-note-color-frame, toprule=.15mm, opacitybacktitle=0.6, title=\textcolor{quarto-callout-note-color}{\faInfo}\hspace{0.5em}{Figure 3.3 -- FastQC Sample Output}, left=2mm, colback=white, bottomtitle=1mm, toptitle=1mm, bottomrule=.15mm, coltitle=black, titlerule=0mm, breakable, opacityback=0, arc=.35mm]

\begin{figure}[H]

\centering{

\includegraphics[width=0.75\linewidth,height=\textheight,keepaspectratio]{images/Fig_3_3_fastqc.png}

}

\caption{\label{fig-fastqc-plots}FastQC Sample Output}

\end{figure}%

Sample FastQC output showing a sharp decline in Phred Quality scores
along sequence length. This is indicated with a red X in the main panel
and indicates a failed quality check. This would typically indicate that
trimming of read ends is required to improve overall quality of read
data prior to alignment.
\href{https://www.bioinformatics.babraham.ac.uk/projects/fastqc/}{Source}

\end{tcolorbox}

\begin{itemize}
\tightlist
\item
  \textbf{Per base sequence quality}: Boxplots by position; expect
  Q\textgreater30 early, decline to Q20+ at ends
\item
  \textbf{Per base sequence content}: Even A/C/G/T; 5' bias = adapters\\
\item
  \textbf{Adapter content}: Rising \% indicates contamination
\item
  \textbf{Sequence duplication}: \textgreater50\% duplicates flags PCR
  bias
\end{itemize}

\begin{tcolorbox}[enhanced jigsaw, rightrule=.15mm, leftrule=.75mm, colbacktitle=quarto-callout-tip-color!10!white, colframe=quarto-callout-tip-color-frame, toprule=.15mm, opacitybacktitle=0.6, title=\textcolor{quarto-callout-tip-color}{\faLightbulb}\hspace{0.5em}{Coverage Calculation using FastQC or Samtools}, left=2mm, colback=white, bottomtitle=1mm, toptitle=1mm, bottomrule=.15mm, coltitle=black, titlerule=0mm, breakable, opacityback=0, arc=.35mm]

Read depth (or depth of coverage) is a basic design consideration. If a
library is sequenced more often (e.g., analyzing it on multiple lanes of
a flow cell) this will offer greater depth and more statistical power to
detect variants. At the same time, obtaining deeper coverage is
relatively expensive. For a typical whole-genome sequence using Illumina
technology that generates 150 nucleotide paired end reads, coverage of
30× to 50× is obtained, meaning that on average any given base in the
genome is covered by 30--50 independent sequencing reads. For
whole-exome sequencing (which often spans ∼50 megabases or 1--2\% of the
genome), depth of coverage is often 100× or more. For targeted
sequencing, depth of sequencing ranges from 5× to 30× depending on
factors such as the number of samples that are run simultaneously.

The redundancy of coverage c is a function of the number of reads N, the
average length of each read L, and the length of the region (e.g.,
genome G) being sequenced\textsuperscript{10}:

Coverage = \(\frac{N \times L}{G}\) where \(G\) = genome size.

\textbf{Using Samtools depth to calculate coverage}: SAMtools is a
library and a software package (Li et al., 2009). We can use it to
analyze alignments in the SAM/BAM input, accomplishing various tasks.
For example, the tool \textbf{samtools depth} is used to give a coverage
amount for EVERY position in the genome. Lucky for you it is in a column
style format with the following columns:

\begin{itemize}
\tightlist
\item
  Column 1: ID\\
\item
  Column 2: Genomic position\\
\item
  Column 3: No of reads covering that position
\end{itemize}

Using a simple awk command, we can calculate the average coverage using
the third column:

\begin{Shaded}
\begin{Highlighting}[]
\NormalTok{samtools depth  *bamfile*  |  awk \textquotesingle{}\{sum+=$3\} END \{ print "Average = ",sum/GENOME\_SIZE\}\textquotesingle{}}
\end{Highlighting}
\end{Shaded}

\begin{quote}
I also recommend calculating genome size directly as this output does
not list sites with zero coverage. While this is less of an issue with
your small yeast genomes, for larger genomes it matters. Also note that
we have used the masked version of the genome. This simply converts all
sites in the genome file that are known to be repetitive (thus difficult
to map for short reads) to N's. These sites in particular are expected
to have zero coverage as well.
\end{quote}

You can do this using awk and the file ``sacCer3.fa.masked.fai'' which
has as its second column the total number of bases on each chromosome.

\textbf{Standard Deviation}: Though not required, you can also calculate
\textbf{standard deviation} to get an idea of how large of a range in
coverage exists across your sample. Try using awk and the formula for
standard deviation to accomplish this.

\href{http://standard-deviation.appspot.com/}{Standard Deviation
Formula}

\begin{Shaded}
\begin{Highlighting}[]
\NormalTok{\#hint you will need both the sum and the squared sum:}
\NormalTok{awk \textquotesingle{}\{sum+=$3; sumsq+=$3*$3\} END \{print "Stdev = ",sqrt(sumsq/GENOME\_SIZE {-} (sum/GENOME\_SIZE)**2)\}\textquotesingle{} depth.output}
\end{Highlighting}
\end{Shaded}

\end{tcolorbox}

\section{SAM/BAM: Aligned Reads}\label{sambam-aligned-reads}

The next major step in a basic genome analysis pipeline is alignment,
which we will go into a lot of depth in subsequent chapters. When FastQ
files are aligned to a reference genome, the output is a SAM (Sequence
Alignment/Map) file. This file format is a compressed file that stores
read alignments. It can be further compressed into its binoary form
BAM\textsuperscript{10}. Finally, an even more compressed version is a
CRAM (Compressed Reference-oriented Alignment Map) file which was
introduced in 2011. CRAM is a compressed columnar file format for
storing biological sequences aligned to a reference sequence. For the
purposes of this chpater, we will focus on the SAM format. BAM files
cannot be directly viewed due to their binary nature.

\begin{tcolorbox}[enhanced jigsaw, rightrule=.15mm, leftrule=.75mm, colbacktitle=quarto-callout-note-color!10!white, colframe=quarto-callout-note-color-frame, toprule=.15mm, opacitybacktitle=0.6, title=\textcolor{quarto-callout-note-color}{\faInfo}\hspace{0.5em}{Lab connections}, left=2mm, colback=white, bottomtitle=1mm, toptitle=1mm, bottomrule=.15mm, coltitle=black, titlerule=0mm, breakable, opacityback=0, arc=.35mm]

In \textbf{\href{Lab04.qmd}{Lab 4}} you will provide a practical
opportunity to have hands-on experience working with raw fastq files on
the Alabama Super Computer. You will also experience moving between
these formats as you generate a bam alignment. You will also index a
fasta formatted genome file.

This upcoming lab will also allow you to use popular bioinformatics
tools to both generate these types of files and perform standard quality
control analysis. You will also first encounter the yeast genomics
dataset we will be using throughout the semester and be able to practice
interpreting QC outputs for a real dataset\textsuperscript{24}. See
\href{appendixC.qmd}{Appendix C} for more detail.

\end{tcolorbox}

\subsection{The SAM/baM Format and
SAMtools}\label{the-sambam-format-and-samtools}

The SAM (sequence alignment/map) format is commonly used to store
next-generation sequence alignments. SAM files can be easily converted
to the BAM (binary alignment/map) format. BAM is a binary representation
of SAM, compressed by the BGZF library, and contains the same
information as the SAM file. Because they are easily interchangeable, we
may refer to the SAM/BAM format. This format is very popular, with many
datasets available in repositories (such as the Sequence Read Archive at
NCBI, the 1000 Genomes Project, and the Cancer Genome Atlas).

The SAM format includes a header section (having lines beginning with
the character @) and an alignment section. The file is tab-delimited,
and there are 11 mandatory fields (Table 9.4) which we examine in our
autism panel (Fig. 9.13). We use the samtools view command to display
the first of over a million rows (each row corresponding to a read that
has been aligned to a reference genome). Twelve fields are shown,
including the sequence (beginning AATCT\ldots) followed by the
corresponding quality scores. The CIGAR string refers to a notation
system for variants. Here string ``148M2S'' shows 148 matches and 2
soft-clipped (unaligned) bases. Standard CIGAR operations are M (match),
I (insertion), and D (deletion). Extended CIGAR options are N (skipped
bases on reference), S (soft clipping), H (hard clipping), and P
(padding). SAMtools is a library and a software package (Li et al.,
2009). We can use it to analyze alignments in the SAM/BAM input,
accomplishing the following tasks:

\begin{itemize}
\tightlist
\item
  convert from other alignment formats, or between SAM and BAM formats;
\item
  sort and merge alignments;
\item
  index alignments (once sorted, BAM file can be indexed generating a
  BAI file used in downstream analyses);
\item
  view alignments in the pileup format (as shown below with the samtools
  view command);
\item
  remove PCR duplicates (this procedure, called duplicate marking or
  ``dedupping,'' removes reads that are redundant); and
\item
  call two classes of variants: single-nucleotide polymorphisms (SNPs)
  and small indels.
\end{itemize}

\begin{tcolorbox}[enhanced jigsaw, rightrule=.15mm, leftrule=.75mm, colbacktitle=quarto-callout-tip-color!10!white, colframe=quarto-callout-tip-color-frame, toprule=.15mm, opacitybacktitle=0.6, title=\textcolor{quarto-callout-tip-color}{\faLightbulb}\hspace{0.5em}{Sample BAM File}, left=2mm, colback=white, bottomtitle=1mm, toptitle=1mm, bottomrule=.15mm, coltitle=black, titlerule=0mm, breakable, opacityback=0, arc=.35mm]

Using SAMtools to view sequence reads from a BAM file at genomic
coordinates of interest. By using the samtools tview command you enable
a genome-wide view of the reads. Using commands accessed via a help
menu, you can view the same reads colored by (a) base quality or (b)
mapping quality. The read depth is relatively low to the left side. The
quality scores are colored blue (for 0--9), green (10--19), yellow
(20--29), or white (≥30). Underlining represents secondary or orphan
reads. This viewer is useful to quickly assess the quality at genomic
loci of interest, such as positions having single-nucleotide variants.

\begin{figure}[H]

\centering{

\includegraphics[width=0.75\linewidth,height=\textheight,keepaspectratio]{images/Fig_3.4_SAMtools_tview.png}

}

\caption{\label{fig-sam}Samtools tview}

\end{figure}%

\end{tcolorbox}

While it is not typical to view BAM files natively, it is helpful to
review the anatomay of a sam file and the various ways to view this file
format using the bioinformatics software \texttt{samtools}.

\begin{tcolorbox}[enhanced jigsaw, rightrule=.15mm, leftrule=.75mm, colbacktitle=quarto-callout-note-color!10!white, colframe=quarto-callout-note-color-frame, toprule=.15mm, opacitybacktitle=0.6, title=\textcolor{quarto-callout-note-color}{\faInfo}\hspace{0.5em}{Figure 3.4 -- Anatomy of a SAM File}, left=2mm, colback=white, bottomtitle=1mm, toptitle=1mm, bottomrule=.15mm, coltitle=black, titlerule=0mm, breakable, opacityback=0, arc=.35mm]

\begin{figure}[H]

\centering{

\includegraphics[width=0.75\linewidth,height=\textheight,keepaspectratio]{images/Figure_9_13_Pevsner.jpg}

}

\caption{\label{fig-sam-anatomy}SAM alignment anatomy (Pevsner 9.13)}

\end{figure}%

Anatomy of a SAM file. The \texttt{samtools\ view} command was used in
Linux to view a BAM file, and the \texttt{\textbar{}\ less} command
sends the output to a single screen of data at a time. (A typical
SAM/BAM file has millions of rows.) A single record of the output is
shown, with 12 features as indicated.

The drop down menu above shows a sample output of
\texttt{samtools\ tview} on the yeast genome sample.

\end{tcolorbox}

The bioinformatics tool \href{https://www.htslib.org/}{SAMtools} is a
suite of programs for interacting with high-throughput sequencing
data\textsuperscript{34}. It consists of three separate repositories:

\begin{itemize}
\tightlist
\item
  Samtools: Reading/writing/editing/indexing/viewing SAM/BAM/CRAM format
\item
  BCFtools: Reading/writing BCF2/VCF/gVCF files and
  calling/filtering/summarising SNP and short indel sequence variants
\item
  HTSlib: A C library for reading/writing high-throughput sequencing
  data
\end{itemize}

Samtools and BCFtools both use HTSlib internally, but these source
packages contain their own copies of htslib so they can be built
independently.

These can be used on the HPC which we will do in labs this semester. But
you can also install these locally. On a Mac, you may use the package
manager \href{https://brew.sh/}{HomeBrew} to install \texttt{samtools}.
If you are on Windows, first install WSL (Windows Subsystem for Linux)
to get a full Linux environment. Then, the easiest approach would be to
use a pacakge manager like \texttt{conda} which will also handle
dependencies automatically.

\begin{tcolorbox}[enhanced jigsaw, rightrule=.15mm, leftrule=.75mm, colbacktitle=quarto-callout-tip-color!10!white, colframe=quarto-callout-tip-color-frame, toprule=.15mm, opacitybacktitle=0.6, title=\textcolor{quarto-callout-tip-color}{\faLightbulb}\hspace{0.5em}{Verification of Sam/Bam files using Samtools}, left=2mm, colback=white, bottomtitle=1mm, toptitle=1mm, bottomrule=.15mm, coltitle=black, titlerule=0mm, breakable, opacityback=0, arc=.35mm]

Another SAMtools tool, the \textbf{samtools flagstat} program is a very
useful way to analyze a bam file. From samtools manual it describes
flagstat's function as:

\begin{quote}
Does a full pass through the input file to calculate and print
statistics to stdout. Provides counts for each of 13 categories based
primarily on bit flags in the FLAG field. Each category in the output is
broken down into QC pass and QC fail, which is presented as ``\#PASS +
\#FAIL'' followed by a description of the category. The first row of
output gives the total number of reads that are QC pass and fail
(according to flag bit 0x200). For example: 122 + 28 in total (QC-passed
reads + QC-failed reads) Which would indicate that there are a total of
150 reads in the input file, 122 of which are marked as QC pass and 28
of which are marked as ``not passing quality controls''
\end{quote}

Sample output file from running flagstat on a bam file:

\begin{Shaded}
\begin{Highlighting}[]
\NormalTok{20968800 + 0 in total (QC{-}passed reads + QC{-}failed reads)}
\NormalTok{0 + 0 duplicates}
\NormalTok{20968800 + 0 mapped (100.00\%:nan\%)}
\NormalTok{20968800 + 0 paired in sequencing}
\NormalTok{11431237 + 0 read1}
\NormalTok{9537563 + 0 read2}
\NormalTok{71098 + 0 properly paired (0.34\%:nan\%)}
\NormalTok{6633306 + 0 with itself and mate mapped}
\NormalTok{14335494 + 0 singletons (68.37\%:nan\%)}
\NormalTok{0 + 0 with mate mapped to a different chr}
\NormalTok{0 + 0 with mate mapped to a different chr (mapQ\textgreater{}=5)}
\end{Highlighting}
\end{Shaded}

\end{tcolorbox}

\section{VCF: Variant Call Format}\label{vcf-variant-call-format}

Once FASTQ files from your sample have been aligned to a reference
genome, it may be assumed that variants (both single-nucleotide variants
and indels) can be called by inspecting the alignment and tabulating the
differences. The problem with this approach is that many sources of
errors occur in the process of sequencing and performing alignment: bias
may occur in how libraries are prepared and amplified; sequencing
technologies all have associated error rates; and mapping has error
rates.

VCF tabulates variants post-alignment
(GATK/SAMtools)\textsuperscript{10,35}.The VCF format was developed
under the 1000 Genomes Project and is currently maintained by the Global
Alliance for Genomics \& Health (GA4GH). For practical guidance on
working with VCF files using BCFtools, see\textsuperscript{36}.

\begin{tcolorbox}[enhanced jigsaw, rightrule=.15mm, leftrule=.75mm, colbacktitle=quarto-callout-note-color!10!white, colframe=quarto-callout-note-color-frame, toprule=.15mm, opacitybacktitle=0.6, title=\textcolor{quarto-callout-note-color}{\faInfo}\hspace{0.5em}{Lab connections}, left=2mm, colback=white, bottomtitle=1mm, toptitle=1mm, bottomrule=.15mm, coltitle=black, titlerule=0mm, breakable, opacityback=0, arc=.35mm]

In \textbf{\href{Lab07.qmd}{Lab 7}} you will have an opportunity to
continue your practical exposure to these file formats as you take the
next step in your genome analysis of the yeast genomics dataset. On the
Alabama Super Computer, you will get hands-on experience using GATK to
call variants from the BAM file you generate in
\textbf{\href{Lab04.qmd}{Lab 4}} creating a VCF file.

This later lab will also guide you in filtering and evaluating variants
for quality. You can refer back to this chapter whenever you encounter
new datasets or need to interpret QC plots and file structures in labs.

\end{tcolorbox}

\subsection{VCF File Format Details}\label{vcf-file-format-details}

A VCF file consists of:

\begin{itemize}
\tightlist
\item
  \textbf{Meta-information lines} (starting with \texttt{\#\#}) that
  define file format, reference genome, and field descriptions
\item
  \textbf{Header line} (starting with \texttt{\#CHROM}) that names the
  columns
\item
  \textbf{Data lines} containing variant information
\end{itemize}

\textbf{9 mandatory columns} + sample columns:

\begin{longtable}[]{@{}ll@{}}

\toprule\noalign{}
Column & Description \\
\midrule\noalign{}
\endhead
\bottomrule\noalign{}
\endlastfoot
CHROM & Chromosome \\
POS & 1-based variant start position \\
ID & dbSNP rsID or other identifier \\
REF & Reference allele \\
ALT & Alternate allele(s), comma-separated \\
QUAL & Phred-scaled quality score \\
FILTER & PASS if passes filters, or filter names \\
INFO & Semicolon-separated annotation fields \\
FORMAT & Colon-separated format of sample fields \\

\caption{\label{tbl-vcf-columns}VCF column definitions (Pevsner Table
9.6).}

\tabularnewline

\end{longtable}

\textbf{Genotype encoding}:

\begin{itemize}
\tightlist
\item
  \texttt{0/0} = homozygous reference
\item
  \texttt{0/1} = heterozygous
\item
  \texttt{1/1} = homozygous alternate
\item
  \texttt{1/2} = heterozygous with two different alternate alleles
\end{itemize}

\subsection{Example VCF Data}\label{example-vcf-data}

\begin{tcolorbox}[enhanced jigsaw, rightrule=.15mm, leftrule=.75mm, colbacktitle=quarto-callout-tip-color!10!white, colframe=quarto-callout-tip-color-frame, toprule=.15mm, opacitybacktitle=0.6, title=\textcolor{quarto-callout-tip-color}{\faLightbulb}\hspace{0.5em}{Sample VCF File}, left=2mm, colback=white, bottomtitle=1mm, toptitle=1mm, bottomrule=.15mm, coltitle=black, titlerule=0mm, breakable, opacityback=0, arc=.35mm]

Below is an example VCF file showing single nucleotide variants (SNVs)
from chromosome I of a yeast genome:

\begin{verbatim}
##fileformat=VCFv4.2
##ALT=<ID=NON_REF,Description="Represents any possible alternative allele at this location">
##FILTER=<ID=LowQual,Description="Low quality">
##FORMAT=<ID=AD,Number=R,Type=Integer,Description="Allelic depths for the ref and alt alleles">
##FORMAT=<ID=DP,Number=1,Type=Integer,Description="Approximate read depth">
##FORMAT=<ID=GQ,Number=1,Type=Integer,Description="Genotype Quality">
##FORMAT=<ID=GT,Number=1,Type=String,Description="Genotype">
##FORMAT=<ID=PL,Number=G,Type=Integer,Description="Normalized Phred-scaled likelihoods">
##INFO=<ID=AC,Number=A,Type=Integer,Description="Allele count in genotypes">
##INFO=<ID=AF,Number=A,Type=Float,Description="Allele Frequency">
##INFO=<ID=AN,Number=1,Type=Integer,Description="Total number of alleles">
##INFO=<ID=DP,Number=1,Type=Integer,Description="Approximate read depth">
##INFO=<ID=MQ,Number=1,Type=Float,Description="RMS Mapping Quality">
##INFO=<ID=QD,Number=1,Type=Float,Description="Variant Confidence/Quality by Depth">
##contig=<ID=chrI,length=230218>
##source=HaplotypeCaller
#CHROM  POS     ID      REF     ALT     QUAL    FILTER  INFO    FORMAT  SAMN03152093
chrI    136     .       G       A       812.01  .       AC=1;AF=1.00;AN=1;DP=25;MQ=46.32;QD=32.48     GT:AD:DP:GQ:PL  1:4,21:25:99:822,0
chrI    349     .       C       T       191.01  .       AC=1;AF=1.00;AN=1;DP=17;MQ=46.06;QD=11.24     GT:AD:DP:GQ:PL  1:6,11:17:99:201,0
chrI    457     .       CA      C       194.01  .       AC=1;AF=1.00;AN=1;DP=12;MQ=39.48;QD=32.33     GT:AD:DP:GQ:PL  1:0,6:10:99:204,0
chrI    610     .       G       A       80.01   .       AC=1;AF=1.00;AN=1;DP=2;MQ=33.14;QD=28.73      GT:AD:DP:GQ:PL  1:0,2:2:90:90,0
\end{verbatim}

\textbf{Key observations:}

\begin{itemize}
\tightlist
\item
  Position 136: G→A SNV with high quality (QUAL=812) and good mapping
  quality (MQ=46.32)
\item
  Position 457: 2bp deletion (CA→C) representing an indel
\item
  INFO field QD (Quality by Depth) helps assess variant confidence
  relative to coverage
\item
  FORMAT GT field shows genotype: \texttt{1} indicates homozygous
  alternate (haploid sample)
\end{itemize}

\end{tcolorbox}

\subsection{Quality Metrics}\label{quality-metrics}

Understanding VCF quality fields is crucial for variant filtering:

\begin{itemize}
\tightlist
\item
  \textbf{QUAL}: Phred-scaled probability that the variant call is
  incorrect (higher is better)
\item
  \textbf{QD (Quality by Depth)}: QUAL/DP ratio; helps identify false
  positives from low-coverage regions
\item
  \textbf{MQ (Mapping Quality)}: Root mean square of mapping quality
  scores; indicates alignment confidence
\item
  \textbf{DP (Depth)}: Number of reads covering the position
\item
  \textbf{GQ (Genotype Quality)}: Confidence in the assigned genotype
\end{itemize}

These metrics are typically used together to filter high-confidence
variants from sequencing artifacts, which we will get into in depth
later in the course as part of Chapter 10.

\section{Variant Types Beyond SNPs}\label{variant-types-beyond-snps}

The outline of a genome analysis workflow presented in Figure 9.6 is
relatively simplistic. We will use that workflow to gain experience with
basic data manipulation. A state-of-the-art workflow, used by many
experts, is the Genome Analysis Toolkit (GATK;\textsuperscript{37,38}).
We show the GATK workflow in Figure 9.7, and explain why its approaches
are crucial to sensitive, specific analyses of genome sequences.

The basic GATK workflow is used for variant discovery with the goal
output being a VCF. But in addition to SNPs, GATK can be used to call
other types of genomic variants. Specifically, GATK can be used to
identify indels and other structural variants\textsuperscript{10}:

\begin{tcolorbox}[enhanced jigsaw, rightrule=.15mm, leftrule=.75mm, colbacktitle=quarto-callout-note-color!10!white, colframe=quarto-callout-note-color-frame, toprule=.15mm, opacitybacktitle=0.6, title=\textcolor{quarto-callout-note-color}{\faInfo}\hspace{0.5em}{Figure 3.5 -- GATK Workflow}, left=2mm, colback=white, bottomtitle=1mm, toptitle=1mm, bottomrule=.15mm, coltitle=black, titlerule=0mm, breakable, opacityback=0, arc=.35mm]

\begin{figure}[H]

\centering{

\includegraphics[width=0.75\linewidth,height=\textheight,keepaspectratio]{images/Fig_3_4_gatk.jpg}

}

\caption{\label{fig-gatk-workflow}GATK pipeline: raw reads →
analysis-ready BAM → VCF.)}

\end{figure}%

\textbf{GATK Best Practices} (Figure @ref(fig:gatk-workflow)) is the
best genome workflow for model organisms, and for modern human genomics.
However, it may not be suitable for calling variants in non-model
systems, which will be covered later in the
course\textsuperscript{37,38}. Each step in the GATK pipeline reduces
the amount of data in a genome analysis, the goal is to improve the
overall quality of the results that are filtered along the way. For
example, we may expect to see up to 50\% of the data is removed during a
standard workflow from raw FastQ file to a VCF. This does not mean that
your data are of poor quality of that your experimental design was
flawed, but only that your resulting variants have been stringently
filtered to be of the highest confidence.

\end{tcolorbox}

\subsection{Variant Types in VCF
Files}\label{variant-types-in-vcf-files}

VCF files can represent several types of genetic variation:

\textbf{SNVs (Single Nucleotide Variants)}: Position 136 above shows a
simple G→A substitution

\textbf{Indels}:

\begin{itemize}
\tightlist
\item
  \textbf{Deletions}: Position 457 shows a deletion where CA in the
  reference becomes C (one base deleted)
\item
  \textbf{Insertions}: Can be represented as \texttt{A→ATAAATTTTTT}
  where bases are added
\end{itemize}

\begin{tcolorbox}[enhanced jigsaw, rightrule=.15mm, leftrule=.75mm, colbacktitle=quarto-callout-tip-color!10!white, colframe=quarto-callout-tip-color-frame, toprule=.15mm, opacitybacktitle=0.6, title=\textcolor{quarto-callout-tip-color}{\faLightbulb}\hspace{0.5em}{Additional Indel Examples}, left=2mm, colback=white, bottomtitle=1mm, toptitle=1mm, bottomrule=.15mm, coltitle=black, titlerule=0mm, breakable, opacityback=0, arc=.35mm]

\begin{verbatim}
chrI    17290   .       AACAC   A       80.01   .       AC=1;AF=1.00;AN=1;DP=4;MQ=36.85;QD=27.08;SOR=2.303  GT:AD:DP:GQ:PL  1:0,2:2:90:90,0
chrI    30097   .       A       ATAAATTTTTT     125.04  .       AC=1;AF=1.00;AN=1;DP=3;MQ=26.04;QD=21.53;SOR=1.179  GT:AD:DP:GQ:PL  1:0,3:3:99:135,0
chrI    31122   .       CATAT   C       69.01   .       AC=1;AF=1.00;AN=1;DP=6;MQ=60.00;QD=34.50;SOR=2.303  GT:AD:DP:GQ:PL  1:0,2:2:79:79,0
\end{verbatim}

\begin{itemize}
\tightlist
\item
  Position 17290: 4bp deletion (AACAC→A)
\item
  Position 30097: 10bp insertion (A→ATAAATTTTTT)
\item
  Position 31122: 4bp deletion (CATAT→C)
\end{itemize}

\end{tcolorbox}

\section{BED file format and
BEDtools}\label{bed-file-format-and-bedtools}

The Browser Extensible Data (BED) format is a tab-delimited text format
for storing genomic regions as coordinates and associated annotations.
Originally developed during the Human Genome Project for the UCSC Genome
Browser\textsuperscript{39}, BED has become a de facto standard in
bioinformatics for representing genomic intervals.

The key advantage of BED format is its use of genomic coordinates rather
than sequences, which optimizes computational efficiency when comparing
genomic regions. Its simplicity also makes it easy to manipulate using
standard command-line tools, scripting languages, or specialized
utilities like BEDTools\textsuperscript{40}.

\subsection{BED File Structure}\label{bed-file-structure}

A minimal BED file requires three tab-delimited fields:

\begin{longtable}[]{@{}lll@{}}

\toprule\noalign{}
Column & Description & Example \\
\midrule\noalign{}
\endhead
\bottomrule\noalign{}
\endlastfoot
chrom & Chromosome name & chr1 \\
chromStart & Start position (0-based) & 1000 \\
chromEnd & End position (1-based, exclusive) & 1500 \\

\caption{\label{tbl-bed-fields}Core BED format fields.}

\tabularnewline

\end{longtable}

\textbf{Important}: BED uses a 0-based coordinate system for the start
position and a 1-based coordinate system for the end position. This
means:

\begin{itemize}
\tightlist
\item
  The first nucleotide of a chromosome has
  \texttt{chromStart=0,\ chromEnd=1}
\item
  A feature's length is simply \texttt{chromEnd\ -\ chromStart}
\item
  The position chr1:1001-1500 in ``human'' coordinates becomes
  \texttt{chr1\ 1000\ 1500} in BED format
\end{itemize}

The BED format supports up to 12 standard fields, with the first three
being mandatory:

\begin{longtable}[]{@{}rll@{}}

\toprule\noalign{}
Field & Name & Description \\
\midrule\noalign{}
\endhead
\bottomrule\noalign{}
\endlastfoot
4 & name & Feature name or identifier \\
5 & score & Score (0-1000) \\
6 & strand & Strand: + or - \\
7 & thickStart & Start of thick drawing (e.g., coding region) \\
8 & thickEnd & End of thick drawing \\
9 & itemRgb & RGB color value (e.g., 255,0,0) \\
10 & blockCount & Number of blocks (exons) \\
11 & blockSizes & Comma-separated list of block sizes \\
12 & blockStarts & Comma-separated list of block starts \\

\caption{\label{tbl-bed-extended}Extended BED format fields (optional
fields 4-12).}

\tabularnewline

\end{longtable}

\subsection{Example BED Files}\label{example-bed-files}

\begin{tcolorbox}[enhanced jigsaw, rightrule=.15mm, leftrule=.75mm, colbacktitle=quarto-callout-tip-color!10!white, colframe=quarto-callout-tip-color-frame, toprule=.15mm, opacitybacktitle=0.6, title=\textcolor{quarto-callout-tip-color}{\faLightbulb}\hspace{0.5em}{Simple BED3 Format}, left=2mm, colback=white, bottomtitle=1mm, toptitle=1mm, bottomrule=.15mm, coltitle=black, titlerule=0mm, breakable, opacityback=0, arc=.35mm]

A basic three-column BED file showing genomic intervals:

\begin{verbatim}
chr1    1000    1500
chr1    2000    2300
chr2    5000    5800
chrX    1000000 1001000
\end{verbatim}

This represents four genomic regions. The first spans 500 bp from
position 1000 to 1499 on chromosome 1 (remember: end position is
exclusive).

\end{tcolorbox}

\subsection{Working with BED Files:
BEDTools}\label{working-with-bed-files-bedtools}

BEDTools\textsuperscript{40} is a powerful suite of utilities for
performing ``genome arithmetic'' on BED files. These tools enable
efficient comparison and manipulation of genomic features through set
operations like intersect, merge, and complement. BEDTools can be
combined with standard UNIX commands in pipelines, facilitating complex
genomic analyses without requiring web-based tools or large data
transfers.

\subsection{BED Format in Other Tools}\label{bed-format-in-other-tools}

Beyond BEDTools, BED format is supported by numerous bioinformatics
tools:

\begin{itemize}
\tightlist
\item
  \textbf{UCSC Genome Browser tools}: bedToBigBed, bigBedToBed for
  creating indexed binary formats
\item
  \textbf{R/Bioconductor}: \texttt{rtracklayer} for importing/exporting
  BED files, \texttt{GenomicRanges} for genomic interval operations
\item
  \textbf{Python}: \texttt{pybedtools} provides Python interface to
  BEDTools functionality
\item
  \textbf{IGV (Integrative Genomics Viewer)}: Native support for BED
  track display
\end{itemize}

\subsection{Connection to Genome
Browsers}\label{connection-to-genome-browsers}

BED files serve as a fundamental format for displaying custom annotation
tracks in genome browsers. The UCSC Genome Browser\textsuperscript{39}
uses BED format extensively, allowing researchers to:

\begin{enumerate}
\def\labelenumi{\arabic{enumi}.}
\tightlist
\item
  \textbf{Upload custom tracks} for visualization alongside reference
  annotations
\item
  \textbf{Define track display properties} through optional track lines
\item
  \textbf{Share results} by hosting BED files on web servers
\end{enumerate}

For large datasets, BED files can be converted to \textbf{bigBed}
format\textsuperscript{41}, a compressed, indexed binary format that
enables rapid visualization of millions of features without uploading
entire datasets to the browser.

As we transition to Chapter 4 on Genome Browsers, understanding BED
format provides the foundation for effectively visualizing and
interpreting genomic data. Genome browsers use BED-formatted tracks to
display genes, regulatory elements, variants, and experimental results
in their chromosomal context, making BED fluency essential for modern
genomics research.

\begin{center}\rule{0.5\linewidth}{0.5pt}\end{center}

\section{Practice Problem‑sets}\label{practice-problemsets-1}

\begin{enumerate}
\def\labelenumi{\arabic{enumi}.}
\item
  The chapter mentions that Q20 is a recommended cutoff for typical
  sequencing projects, but that this still equates to \textasciitilde100
  million errors in a standard sequencing project. This was based on a
  dataset with 100 million 100 base pair reads, or 10 billion reads.
  Imagine you increased the stringency from 99\% to 99.9\% by imposing a
  Q30 cutoff instead. \textbf{What would be the new number of expected
  sequencing errors resulting from this adjustment?}
\item
  Quality control is the critical first step in any sequencing analysis.
  Different research environments (local computers, web platforms,
  command-line servers, R environments) offer different tools for QC
  analysis. To better understand the strengths and limitations of each
  approach, compare FastQC implementations across different platforms.
  To do this, search the SRA accession \textbf{SRR518692}
  (\emph{Aeromonas hydrophila} whole genome sequencing) which was
  sequenced here at Auburn. In class, we will run FastQC using the
  Galaxy Web Platform. Here, test out a local FastQC Installation
  (Mac/Windows). First, download
  \href{https://www.bioinformatics.babraham.ac.uk/projects/download.html\#fastqc}{FastQC}
  for Mac/PC. Next, download the SRR518692 FASTQ files using Galaxy or
  SRA Toolkit. Third, run FastQC locally on the downloaded file.
  \textbf{Compare the output to your Galaxy results from class.}
\item
  As an extension of Problem \#2, look up the genome size of the species
  \emph{Aeromonas hydrophila} on NCBI. Then using the number of reads
  and the read length from the FastQC report, \textbf{calculate the
  expected coverage} of this sample.
\item
  As another example of tool implementation, try your hand at
  \textbf{running FastQC using the R/Bioconductor Approach}. There are
  several R pacakges that you can use. If you have installed FastQC,
  then you can use
  \href{https://cran.r-project.org/web/packages/fastqcr/readme/README.html}{fastqcr}.
  Alternatively, you can use an alternative pacakge called
  \href{https://kasperdanielhansen.github.io/genbioconductor/html/ShortRead.html}{ShortRead}.
  Once you have retrieved QC metrics using either R package,
  \emph{compare your results with that of Galaxy and FastQC locally.
  Also reflect on the difference in experience across these three
  platforms. Which was easiest to run? fastest? Which was easiest to
  interpret? Which one would you choose if you needed to run these QC
  metrics on a set of 100 fastqc files?}
\item
  Look back at the Lab from Computational Biology where you learned to
  use AWK to manipulate VCF files. Now having learned more about this
  file format, why do you think a VCF file worked so well for this lab
  assignment working with the tool \texttt{awk}?
\item
  \textbf{SAM/BAM CIGAR Strings}: Alignment information encoded in
  SAM/BAM files use CIGAR strings. CIGAR (Compact Idiosyncratic Gapped
  Alignment Report) strings provide a compressed representation of how
  reads align to a reference genome. Given these CIGAR strings, decode
  what they mean: a). \texttt{148M2S} (from chapter example); b).
  \texttt{50M1I49M}; c). \texttt{75M5D25M}; d). \texttt{10S90M}.
  \emph{Write out what they mean in plain English. How long is each read
  corresponding to each string? What types of potential variants do
  these reads indicate?}
\end{enumerate}

\section{Reflection Questions}\label{reflection-questions-2}

\begin{enumerate}
\def\labelenumi{\arabic{enumi}.}
\item
  \textbf{Format Evolution and Compression}: The chapter describes the
  evolution from FASTA → FASTQ → SAM → BAM → CRAM, with each format
  adding information or improving compression. Explain the trade-offs
  between human readability and computational efficiency in this
  progression. Why do you think FASTA has remained the standard for
  reference genomes for 40 years despite being uncompressed text, while
  CRAM has largely replaced SAM/BAM for aligned reads?
\item
  \textbf{Phred Quality Score as ASCII}: In the FastQ file format, Phred
  quality scores are encoded as ASCII (+33 offset) rather than raw
  numbers. Why do you think this format is preferred?
\item
  \textbf{Phred Quality Score Encoding}: FASTQ format encodes Phred
  quality scores as single ASCII characters with a +33 offset (e.g., Q30
  becomes `?') rather than storing them as numeric values (e.g.,
  ``30''). Why do you think was this ASCII encoding scheme chosen? What
  benefits do you think this conversion has in terms of file
  compression? For example, why might the string ``?????'' (five Q30
  scores) compress differently than ``30,30,30,30,30''? Would ASCII
  encoding help or hurt compression? What are the trade-offs between
  human-readability and computational efficiency?
\item
  \textbf{Different ASCII encoding}: The +33 offset specifically makes
  quality scores printable characters (! through \textasciitilde). The
  earlier Solexa/Illumina 1.0 format used +64 offset (@ through
  \textasciitilde). What practical problems might arise from having two
  different encoding schemes in the field? How would you detect which
  encoding a file uses?
\item
  \textbf{Coordinate Systems and Biological Meaning}: BED format uses
  0-based start positions but 1-based end positions (0-based, half-open
  intervals), while VCF uses 1-based positions for both start and end.
  Discuss why these different conventions exist and how they affect
  biological interpretation. What errors might occur if you incorrectly
  convert between these systems? Provide a specific example with
  coordinates.
\end{enumerate}

\section{References}\label{references-3}

\chapter{Genome Databases and
Browsers}\label{genome-databases-and-browsers}

\section{Genome Databases and
Browsers}\label{genome-databases-and-browsers-1}

Genome databases and browsers form the foundation of modern
bioinformatics, providing access to reference assemblies, annotations,
and experimental data across thousands of species. These resources
evolved from early sequence repositories like GenBank (established 1982)
to sophisticated platforms integrating comparative genomics, functional
annotations, and interactive visualization\textsuperscript{10}. This
chapter traces their historical development while emphasizing practical
skills for data retrieval, sharing policies, visualization, and
reproducible workflows.

Three main sites that have been responsible for storing nucleotide
sequence data from 1982 to the present. These are:

\begin{enumerate}
\def\labelenumi{\arabic{enumi}.}
\tightlist
\item
  GenBank at the National Center for Biotechnology Information (NCBI) of
  the National Institutes of Health (NIH);
\item
  The European Molecular Biology Laboratory (EMBL)-Bank Nucleotide
  Sequence Database (EMBL-Bank), part of the European Nucleotide Archive
  (ENA) at the European Bioinformatics Institute (EBI); and
\item
  The DNA Database of Japan (DDBJ) at the National Institute of
  Genetics.
\end{enumerate}

All three are coordinated by the International Nucleotide Sequence
Database Collaboration (INSDC), and they share their data daily.
GenBank, EMBL-Bank, and DDBJ are organized as databases within NCBI,
EBI, and DDBJ which offer many dozens of other resources for the study
of sequence data. GenBank exceeds 10\^{}12 bases (2025), with daily
releases.

In addition to these nucleotide sequence databases, there are three
major genome portals that we will cover in this course. First, NCBI
which is directly connected to GenBank. Second, is
\href{https://www.ensembl.org/index.html}{Ensembl}, started before the
human genome was completed with a goal of annotating genomes and
providing a comparative genomic database. Third, the University of
California, San Cruz \href{https://genome.ucsc.edu/}{(UCSC) Genome
Browser} which was started in 2000 concurrent with the human genome
project. The latter two are heavily focused on vertebrate genomes. In
addition to these three, there are several organism-specific databases.
Each database and browser has unique strengths in curation, annotation
depth, and visualization.

In addition to web genome browsers and databases, this chapter will
introduce tools for retrieving data from databases, tools for
visualizing data retrieved, and an online platform for genome analysis
called ``Galaxy''. Web interfaces like browsers enable exploration,
while platforms like Galaxy transition to scripted workflows for
scalability. Finally, this chapter concludes with a section on data
sharing policies.

\begin{tcolorbox}[enhanced jigsaw, rightrule=.15mm, leftrule=.75mm, colbacktitle=quarto-callout-note-color!10!white, colframe=quarto-callout-note-color-frame, toprule=.15mm, opacitybacktitle=0.6, title=\textcolor{quarto-callout-note-color}{\faInfo}\hspace{0.5em}{Lab connections}, left=2mm, colback=white, bottomtitle=1mm, toptitle=1mm, bottomrule=.15mm, coltitle=black, titlerule=0mm, breakable, opacityback=0, arc=.35mm]

\begin{itemize}
\tightlist
\item
  \href{Lab03.qmd}{Lab 3}) will involve an exploratory activity to
  engage with three different genome databases to give a practical
  hands-on opportunity to explore, customize and export genome data. You
  will get a chance to engage with UCSC and Ensembl, as well as an
  organism-specific database for fruit flies.
\end{itemize}

\end{tcolorbox}

\section{Historical Overview}\label{historical-overview}

The Human Genome Project (HGP, 1990--2003) catalyzed genome database
development, necessitating standardized repositories for massive
sequence data. GenBank initially stored \textasciitilde600 sequences; by
HGP completion, it held billions of bases. Parallel efforts produced
EMBL (Europe, 1980) and DDBJ (Japan, 1987) formed the International
Nucleotide Sequence Database Collaboration (INSDC).

Genome browsers emerged to visualize assemblies: UCSC Genome Browser
(2000) displayed the draft human genome with tracks for genes and
clones. Ensembl (2000) followed, focusing on automated annotation for
vertebrates. NCBI's Map Viewer (1999) integrated genetic/physical maps.
Organism-specific resources filled gaps: FlyBase (1992) for Drosophila,
WormBase (2000) for \emph{C. elegans}. These various browsers align data
to linear chromosomes using 1-based ``human-readable'' coordinates
(e.g.~chr1:1,000,000). The display of various tracks stack
hierarchically with the compression ranging from dense (blocks), to
squish (arrows), to a full view of annotation features.

\section{Major Databases}\label{major-databases}

\subsection{NCBI as a Genome Browser}\label{ncbi-as-a-genome-browser}

NCBI is best known as the home of GenBank, RefSeq, and the Sequence Read
Archive, but it also provides a full-featured genome browser called
\textbf{\href{https://www.ncbi.nlm.nih.gov/genome/gdv}{Genome Data
Viewer (GDV)}} that sits at the center of the NCBI ecosystem. Where
Ensembl is built around automated comparative annotation, NCBI's browser
is tightly integrated with RefSeq, SRA, GEO, ClinVar, and other NCBI
resources, so that you can move smoothly from a gene or variant record
to its precise genomic context and back again. In practice, GDV is
somewhat clunkier and less visually polished than UCSC or Ensembl, but
it often exposes the most up-to-date RefSeq annotations and assemblies,
including patch releases and alternate haplotypes, making it a valuable
default choice when you care about current reference models and NCBI's
own curation decisions.

NCBI's older ``Map Viewer'' historically provided karyotype-style
chromosome overviews and access to multiple map types, but this has been
superseded by Genome Data Viewer and the related Comparative Genome
Viewer. GDV retains the same basic idea: a high-level chromosome view
that you can use to orient yourself, combined with a detailed sequence
view with multiple tracks. The chromosome ideogram (or karyogram)
appears as a vertical strip, with bands arranged from top to bottom, in
contrast to UCSC's default left-to-right orientation that you will
encounter later in this chapter.

A typical workflow with GDV begins with choosing a species and assembly
using the tree or table on the left (for example, \emph{Homo sapiens}
GRCh38.p14) and then either type a gene symbol, coordinate range, or
variant ID into the search box to jump to your region of interest.
Alternatively, if you start from either the \textbf{Gene} database or
from a BLAST result page, there is usually a link labeled ``Open in
Genome Data Viewer'' (or a browser icon) that sends you directly into
GDV at the appropriate coordinates.

Once the browser is open to a particular locus, GDV's main panel is
divided into a \textbf{Sequence Viewer} region and several context
widgets. The Sequence Viewer shows tracks stacked horizontally, each
track corresponding to a different kind of annotation or data: RefSeq
gene models, mRNA and protein alignments, dbSNP and ClinVar variants,
structural variants from dbVar, expression or coverage tracks, and
optionally your own custom tracks or BLAST hits. Zoom controls allow you
to move from a gene-wide view down to individual exons, and the ``Exon
Navigator'' widget provides a schematic representation of exons for the
selected transcript; clicking on a particular exon in this widget
updates the Sequence Viewer to center on that exon and can even open a
pop-up with its genomic sequence in FASTA format. Right-click menus
within tracks provide additional options, such as viewing feature
details, linking back to the corresponding Gene or Nucleotide records,
or copying sequence with or without flanking regions.

\begin{tcolorbox}[enhanced jigsaw, rightrule=.15mm, leftrule=.75mm, colbacktitle=quarto-callout-note-color!10!white, colframe=quarto-callout-note-color-frame, toprule=.15mm, opacitybacktitle=0.6, title=\textcolor{quarto-callout-note-color}{\faInfo}\hspace{0.5em}{Connection to Assigned Reading}, left=2mm, colback=white, bottomtitle=1mm, toptitle=1mm, bottomrule=.15mm, coltitle=black, titlerule=0mm, breakable, opacityback=0, arc=.35mm]

In the paper you were assigned last week\textsuperscript{3}, they
generated several novel genome assemblies. In their methods, they list
the GCF accesssion id's for their six genomes. Look up ONE of these in
NCBI and open GDV. Once there, search for the antifreeze protein (AFP)
gene region to look at it more closely. \textbf{How does this
visualization help to better understand the major findings of this
study?}

\end{tcolorbox}

One of GDV's strengths is its tight integration with other NCBI analysis
tools. From within the Sequence Viewer toolbar, you can launch
\textbf{Primer-BLAST} to design primers flanking your region, access
BLAST to align a query sequence against the current assembly, or export
the underlying sequence and annotation for downstream use. The
``Tracks'' configuration dialog allows you to add or remove
NCBI-provided tracks such as dbSNP common variants, ClinVar pathogenic
variants, GWAS hits, and structural variants from dbVar, as well as
upload your own track files (e.g., BED or BigWig) or display data from
an SRA alignment. Because GDV draws on RefSeq as its primary annotation
source, the gene models you see there will usually represent NCBI's
current ``official'' interpretation of a gene's structure for that
assembly, and they are updated as new evidence accumulates.

For extracting data, GDV supports several levels of detail. At the most
basic, you can select a genomic interval and choose ``Download
sequence'' to obtain a FASTA file of the reference sequence for the
current view. For more structured downloads, GDV links out to
\textbf{NCBI Datasets}, a service that lets you download entire genomes,
chromosomes, annotation tables, or gene sets as zipped packages via a
web interface or a command-line client. You can also follow links back
to the Gene or Nucleotide databases, where the same region can be
exported in formats such as GenBank flat files or GFF3 for annotation
pipelines. For variation analysis, clicking on individual variant
markers (e.g., a SNP or structural variant in a GDV track) takes you to
dbSNP, ClinVar, or dbVar records with richer metadata and allele
frequency information. In this way, GDV serves both as a visualization
layer and as a routing hub into the broader NCBI resource network.

\begin{tcolorbox}[enhanced jigsaw, rightrule=.15mm, leftrule=.75mm, colbacktitle=quarto-callout-note-color!10!white, colframe=quarto-callout-note-color-frame, toprule=.15mm, opacitybacktitle=0.6, title=\textcolor{quarto-callout-note-color}{\faInfo}\hspace{0.5em}{Box 4.1: NCBI Genome Browser in the Gene Annotation Project}, left=2mm, colback=white, bottomtitle=1mm, toptitle=1mm, bottomrule=.15mm, coltitle=black, titlerule=0mm, breakable, opacityback=0, arc=.35mm]

In the manual gene annotation project later in the semester, you will
use NCBI's genome browser as the primary environment for building a
custom gene model for an unannotated PRDM9 locus in a non-model
reference genome. Starting from the appropriate assembly, you will use
the Genome Data Viewer's Sequence Viewer, exon navigation, and linked
tools to identify putative exons, splice sites, and coding regions,
applying basic principles of gene anatomy (start and stop codons, splice
donor and acceptor motifs, reading frame continuity, and domain
structure) to propose a gene model. You will then compare your proposed
annotation with existing RefSeq or GenBank entries and document the
evidence you used, practicing the same skills you will later apply at
scale using automated pipelines.

\end{tcolorbox}

\subsection{Ensembl and Comparative
Genomics}\label{ensembl-and-comparative-genomics}

\href{https://www.ensembl.org}{Ensembl} is a genome browser and database
that emphasizes automated annotation and comparative genomics across
many species, rather than only serving as a static repository of a
single reference genome\textsuperscript{42,43}. Where NCBI's RefSeq
project focuses on providing a curated, largely non-redundant set of
reference sequences, Ensembl's Compara pipeline systematically compares
gene and genome sequences across species to infer homology
relationships, build gene trees, and identify conserved syntenic blocks.
This makes Ensembl particularly useful when you are interested in
questions about gene family evolution, orthologous relationships between
species, or the conservation of genomic neighborhoods rather than just
the annotation of a single gene in a single organism.

A central concept in Ensembl is the distinction between orthologs and
paralogs. Orthologs are genes in different species that diverged through
a speciation event, and they often (though not always) retain similar
functions after divergence. Paralogs are genes that arose by duplication
within a species (or within a lineage) and may have diverged in
sequence, expression, or function as they evolve independently. Ensembl
infers these relationships using gene trees: for each gene, a
representative protein sequence is aligned against all other proteins, a
phylogenetic tree is built, and duplication versus speciation nodes are
labeled to classify relationships as orthologous or paralogous. These
inferred homologies are then exposed through the web browser (e.g.,
``Orthologues'' and ``Paralogues'' views), through BioMart for bulk
downloads, and via programmatic interfaces.

From a practical standpoint, using Ensembl for a gene-centric analysis
usually begins with the home page and the species selector. For example,
to investigate the human PRDM9 gene, you would choose \emph{Homo
sapiens} as the species, type ``PRDM9'' into the search box, and follow
the link to the \textbf{Gene} page. This gene page serves as a hub with
multiple tabs (or left-hand menu entries) such as ``Summary,''
``Transcript table,'' ``Sequence,'' ``Variation,'' and ``Comparative
genomics.'' The summary tab provides basic information about
coordinates, transcripts, protein domains, and links out to external
resources like UniProt, while the comparative genomics section is where
Ensembl's strengths in cross-species analysis become particularly
apparent.

Within the comparative genomics section of a gene page, the
``Orthologues'' view lists predicted orthologs of the selected gene
across other species, often grouped into taxonomic categories such as
primates, mammals, sauropsids, or teleost fish. For each ortholog,
Ensembl shows the target gene ID, percent identity, coverage statistics,
and an orthology type (e.g., one-to-one, one-to-many, many-to-many),
which can be used to infer if a gene has been duplicated in one lineage
but not another. The ``Paralogues'' view similarly lists duplicated
copies within the focal species, providing insight into gene family
expansions or subfunctionalization events. A separate ``Gene tree'' view
displays a phylogenetic tree of the gene family, marking speciation and
duplication nodes explicitly. Because these displays are driven by
Ensembl's Compara pipeline, they update as new assemblies and species
are added.

Ensembl also offers synteny and whole-genome alignment views that extend
comparative genomics beyond individual genes. Synteny, in this context,
refers to the conserved order of genomic blocks between species; Ensembl
calculates synteny from pairwise genome alignments and then displays it
in a dedicated ``Synteny'' view. In that view, the central chromosome
strip typically corresponds to your species of interest, and aligned
segments in a second species are shown as colored blocks on smaller
chromosome representations, indicating regions where gene order and
orientation are preserved. This is particularly useful when evaluating
whether a gene's genomic neighborhood is conserved across species
(syntenic orthologs) or has been rearranged, which can have implications
for regulatory conservation and for interpreting complex orthology
relationships that involve duplications and losses. For more
fine-grained comparisons, the ``Region in detail'' view can be
configured to display conservation tracks and pairwise or multiple
genome alignments, allowing you to visually inspect conservation at the
level of exons, introns, or regulatory regions.

Beyond these visualization tools, Ensembl provides several practical
ways to retrieve data. Directly from the gene page, you can download DNA
or protein sequences for a gene or a specific transcript in FASTA
format, including flanking regions or UTRs, using the ``Sequence'' or
``Export data'' options. You can also export alignments of a gene region
across multiple species as a multiple sequence alignment, which is
useful for downstream phylogenetic analysis or motif discovery. For
large-scale or systematic analyses, the BioMart interface (accessible
via the ``BioMart'' link in the top navigation) allows you to build
queries such as ``all human genes with mouse orthologs,'' ``all ortholog
pairs between human and zebrafish with one-to-one relationships,'' or
``all paralogs of a given gene family with their percent identity,'' and
then export results as tab-delimited tables suitable for R or Python.
Programmatic access via Ensembl's REST API and direct downloads of
Compara databases provide further options for scripting and reproducible
analyses, but those are beyond the scope of this chapter.

One limitation students will encounter, especially in research projects
that involve non-model species or unusual gene families, is that
Ensembl's coverage of orthology relationships is incomplete and varies
across taxa. Although Ensembl Genomes now includes plants, fungi, and
many other clades, not all assemblies are annotated at the same depth,
and not every gene has a confidently inferred ortholog in every species
of interest. The PRDM9 gene provides a useful example: in humans,
Ensembl offers a well-annotated PRDM9 gene with domain structure and
some orthology relationships, but if you search for PRDM9 orthologs in
many of the taxa used in the semester-long research
project---particularly among certain fish, amphibians, or non-model
vertebrates---you will find gaps, partial annotations, or missing
entries. This reflects real biological phenomena, such as repeated
losses and domain truncations of PRDM9 in vertebrate evolution, but it
also highlights technical limitations of current annotations and the
underlying assemblies. Students should therefore treat the absence of an
ortholog in Ensembl with caution: it may mean that the gene was lost,
that it is present but not annotated, or that it falls below Ensembl's
thresholds for confident homology prediction.

Finally, Ensembl's automated annotation and comparative resources are
most powerful when combined with a clear understanding of what kinds of
data you can and cannot extract. You can reliably obtain current gene
models for supported assemblies, predicted orthologs and paralogs within
the scope of Ensembl's Compara pipeline, whole-genome alignment-based
conservation scores, and synteny summaries for species that have been
included in pairwise or multiple alignments. You cannot expect Ensembl
to provide high-quality homology relationships for species with draft or
fragmented assemblies, to capture every lineage-specific gene, or to
resolve highly complex gene families perfectly across all taxa.

\subsection{UCSC Genome Browser}\label{ucsc-genome-browser}

The UCSC Genome Browser is one of the longest-standing and most widely
used genome browsers, distinguished by its rich catalogue of annotation
tracks and its highly configurable display. Where NCBI's Genome Data
Viewer emphasizes tight integration with RefSeq and the broader NCBI
ecosystem, UCSC focuses on flexible visualization of many independent
annotation layers on top of a reference assembly. In the browser, genes,
variants, regulatory elements, conservation scores, experimental
datasets, and user-defined features are all represented as tracks
aligned along the chromosome. You can turn tracks on and off, change
their display mode and order, and save track configurations as sessions,
which makes UCSC particularly powerful for exploratory work and for
troubleshooting complex analysis pipelines that generate interval-based
results such as BED or BigWig files.

UCSC's main interface is organized around a horizontal chromosome view,
in contrast to the vertical karyogram-style overview in NCBI's Genome
Data Viewer. After selecting a genome and assembly from the top menu
(for example, \emph{Homo sapiens} → GRCh38/hg38, or \emph{Drosophila
melanogaster} → dm6), you can search by gene name, genomic coordinates,
or keyword to move to a region of interest. The top of the page shows
the current coordinates and a navigational ``position bar'' that
includes zoom and scroll controls, while the main panel displays stacked
annotation tracks across this interval. Below the viewer is a long list
of track groups, organized into sections such as ``Mapping and
Sequencing,'' ``Genes and Gene Predictions,'' ``Variation,''
``Regulation,'' and ``Comparative Genomics,'' each with individual
tracks that can be set to hide, dense, squish, pack, or full display
modes. These modes control the level of detail and the vertical space
tracks occupy, allowing you to tailor the browser to your current
question---for example, switching gene tracks to ``pack'' to see exon
structures clearly, while showing dense conservation or variant tracks
as simple stacked bars.

One of the strengths you will encounter in Lab 3 is UCSC's suite of
conservation tracks, which are particularly useful. UCSC provides
multi-species alignments generated by tools like MultiZ, and from these
alignments computes conservation scores such as phastCons and phyloP for
each base across the genome. In the browser, the ``Comparative
Genomics'' section includes tracks like ``Conservation'' and ``MultiZ
alignments,'' which display conservation as a wiggle plot or as stacked
alignment blocks between the reference and other species. When you zoom
in on a gene such as \texttt{rdl} or \texttt{dsx} in \emph{D.
melanogaster}, turning on these conservation tracks reveals patterns
such as highly conserved exons contrasted with more rapidly evolving
introns and intergenic regions, or conserved noncoding elements that may
correspond to regulatory sequences.

Beyond the built-in tracks, UCSC's support for custom tracks and data
hubs allows you to visualize your own analysis results directly in the
browser, making it a useful companion to command-line tools like
bedtools. Any dataset that can be expressed as genomic intervals or
scores---classically in BED, bedGraph, BigBed, or BigWig formats---can
be uploaded and displayed as a custom track on top of the reference
genome. For small files, you can paste the BED content directly into the
``add custom tracks'' dialog or upload a file from your computer,
specifying track attributes such as name, description, color, and
visibility. For larger datasets, especially those hosted on a web or
cloud server, you can create ``track hubs'' that point UCSC to remotely
hosted BigBed or BigWig files; the browser streams only the portions
needed for the current view, keeping interactions responsive even with
gigabyte-scale datasets.

For this course, a critical connection is the interplay between
bedtools-based workflows and the UCSC browser, which can make it
immediately obvious whether the problem lies in coordinate conventions,
strand handling, file sorting, or an upstream step. Literally ``seeing''
where your intervals fall along the genome makes these debugging tasks
far more intuitive than staring at raw tables in a terminal.

As with Ensembl and NCBI, UCSC has limitations and design choices that
you should keep in mind. The browser is optimized for a curated set of
assemblies and track collections, historically biased toward human and
major model organisms; while many additional species and assemblies are
available, coverage and annotation depth are variable. The default gene
tracks prioritize UCSC and GENCODE models, which sometimes differ in
subtle ways from RefSeq or Ensembl annotations, so careful attention to
which gene set you are using is essential when interpreting exon
structures or designing experiments. Finally, while UCSC's flexibility
is a major strength, the sheer number of tracks and configuration
options can be overwhelming initially.

\subsection{Organism-Specific Model Organism
Databases}\label{organism-specific-model-organism-databases}

Before the rise of today's large, integrated portals, many community
databases emerged around individual model organisms, often inspired by
early efforts at The Institute for Genomic Research (TIGR). TIGR
developed a suite of gene indices and repeat databases in the late 1990s
and early 2000s, such as the TIGR Gene Indices and the TIGR Plant Repeat
Databases, to cluster ESTs, construct consensus ``virtual transcripts,''
and catalog repetitive elements in key plant genomes. These TIGR
resources provided some of the first species-focused, integrative views
of genes, repeats, and expression data, and they helped set the pattern
for organism-specific databases that tightly couple genome sequences
with curated biological knowledge for a particular research community.

Over the past three decades, a number of such organism-specific
databases have become indispensable anchors for model organism genetics
and genomics. A non-exhaustive list of major examples includes:

\begin{itemize}
\item
  \textbf{\href{https://flybase.org}{FlyBase}}, established early 1990s,
  funded from 1992 by the U.S. National Human Genome Research Institute
  of NIH\\
  FlyBase is the primary database for \emph{Drosophila} genetics and
  genomics, providing curated information on genes, alleles, phenotypes,
  stocks, transposons, and bibliography. It integrates genome sequence
  data with genetic maps, expression, and mutant information, and serves
  as the historical record for \emph{Drosophila} nomenclature and
  community resources.
\item
  \textbf{\href{https://www.informatics.jax.org}{Mouse Genome
  Informatics (MGI)}}, origins in late 1980s / early 1990s\\
  MGI is the central knowledgebase for the laboratory mouse, maintained
  at The Jackson Laboratory. It integrates genome features, gene
  function, expression data, phenotypes, and human disease models, and
  provides standardized gene, allele, and strain nomenclature. The
  resource evolved from early NIH-funded efforts to merge several
  Jackson Laboratory databases into a single visualization and analysis
  system for the mouse genome.
\item
  \textbf{\href{https://zfin.org}{Zebrafish Information Network
  (ZFIN)}}, established 1994 ZFIN is the community database for
  \emph{Danio rerio}, the zebrafish model system. It curates genetic,
  genomic, and developmental data, including mutants, transgenes,
  markers, expression patterns, and phenotypes. ZFIN also maintains
  zebrafish nomenclature standards and provides links to sequence and
  ortholog information in other genomic resources.
\item
  \textbf{\href{https://wormbase.org}{WormBase}}, origins in the 1990s,
  formal web-based resource by early 2000s\\
  WormBase began as a web interface to the ACeDB database for
  \emph{Caenorhabditis elegans} and now covers multiple nematode
  species. It integrates genetic and physical maps, gene models,
  expression, phenotypes, and RNAi and CRISPR resources, supporting both
  classical genetics and high-throughput functional genomics in worms.
\item
  \textbf{\href{https://www.yeastgenome.org}{Saccharomyces Genome
  Database (SGD)}} , established mid-1990s SGD is the primary database
  for the budding yeast \emph{Saccharomyces cerevisiae}, providing the
  reference genome sequence, gene catalog, functional annotations, and
  curated literature for each gene. It plays a major role in
  biocuration, systematically extracting experimental results from the
  yeast literature and capturing them in structured fields.
\item
  \textbf{\href{https://phytozome-next.jgi.doe.gov}{Phytozome}},
  launched late 2000s\\
  Phytozome, hosted by the U.S. Department of Energy Joint Genome
  Institute, is a comparative genomics portal for green plants. It
  aggregates genome assemblies and annotations for dozens of plant
  species, offering orthology assignments, gene family clustering, and
  synteny views that support both crop improvement and basic plant
  biology.
\end{itemize}

In addition to these, many other community resources exist, such as
PomBase for fission yeast, TAIR and Araport for \emph{Arabidopsis
thaliana}, Xenbase for \emph{Xenopus}, and Rat Genome Database (RGD) for
rat, each tailored to the needs and conventions of its respective
research community. Together, these organism-specific databases encode
not only sequences and annotations but also decades of accumulated
experimental knowledge, nomenclature standards, and community practices
that are not easily replicated in generic multi-species portals.

\subsubsection{Funding for Organism
Databases}\label{funding-for-organism-databases}

A common thread across these databases is their reliance on sustained
scientific funding, often through a mix of national research agencies
and institutional support. FlyBase, for example, has historically been
funded by NIH (via the National Human Genome Research Institute and
related institutes) to support the \emph{Drosophila} community; similar
NIH or international funding supports WormBase, ZFIN, MGI, and SGD.
These grants pay for database infrastructure, software development, and,
critically, expert biocurators who read the literature, integrate new
data types, and maintain nomenclature standards. Despite their central
importance, these community databases face recurring funding
uncertainties as grant cycles end, program priorities shift, and funders
weigh investments in organism-specific resources against broader,
centralized efforts. The long-term sustainability of some of these
databases remains unclear, especially for taxa with smaller research
communities or those outside the major human-health-focused models.

For researchers, however, the value of these organism-specific databases
is hard to overstate. They provide deep, historically informed context
that is often missing from more general resources: allele series and
classic mutant phenotypes in FlyBase, detailed recombinase resources and
phenotype ontologies in MGI, lineage and fate maps in ZFIN and WormBase,
and exhaustive gene function summaries in SGD. These resources
effectively encode the collective memory of their respective fields,
linking early classical genetics papers to modern genomic datasets and
making it possible to interpret new high-throughput experiments in light
of decades of accumulated knowledge.

\section{Sequence Read Archive (SRA)}\label{sequence-read-archive-sra}

The \href{https://www.ncbi.nlm.nih.gov/sra/docs/}{Sequence Read Archive}
(SRA) is NCBI's repository for \emph{raw} high-throughput sequencing
data and is part of the same ecosystem as GenBank and RefSeq. Whereas
GenBank focuses on assembled sequences and RefSeq on curated reference
sequences, SRA stores the unprocessed reads and alignments that underlie
those assemblies and publications. This design enables other researchers
to reanalyze published data, test new methods, and improve genome
annotations without having to regenerate the raw data themselves. Like
GenBank, SRA is synchronized with partner archives at EMBL-EBI and DDBJ
as part of the International Nucleotide Sequence Database Collaboration.

An SRA record is typically linked to a \textbf{BioProject} (study-level
description), one or more BioSamples (biological sources), and then to
one or more SRA runs that correspond to specific sequencing libraries or
lanes. In practice, you will usually interact with SRA at the ``run''
level via accessions like SRR518692 (Auburn \emph{Aeromonas hydrophila}
sequencing), which you can download as FASTQ and then process through
your own pipeline.

\subsection{Depositing Data into SRA}\label{depositing-data-into-sra}

For anyone with a free MyNCBI account, the web-based Submission Portal
Wizard is used to deposit sequencing data into SRA. First, prepare your
raw fastq files (*.fastq.gz or split based on read groups if paired),
and include a basic study description (project title, abstract-like
summary, organism). Next, you will need to include metadata for your
sample (strain, tissue, treatment, etc.) and the sequencing runs
(library information, sequencing platform, lane information, etc.). If
you have a lot of samples, consider downloading the BioSample worksheet
template that you can fill out locally.

Upon depsositing data, you will need to make a decision about data
access. I strongly recommend making an appointment with your
\href{https://libguides.auburn.edu/subjectlibrarians}{subject
librarian}. Subject librarians offer discipline-specific services
designed to help Auburn University students, faculty, and staff access
information resources and use them effectively. These include research
consultations and discipline-specific support for using the Innovation
and Research Commons (I\&RC). Alternatively, you can contact the
\href{https://www.lib.auburn.edu/people/alicia-krzton.php}{research data
management librarian} directly.

\subsection{Retrieving Data from SRA (SRR518692
Example)}\label{retrieving-data-from-sra-srr518692-example}

As an example of how to retrieve public data from SRA, let's walk
through downloading the \emph{Aeromonas hydrophila} dataset SRR518692
via the NCBI web interface (manual download). First, go to the
\href{https://www.ncbi.nlm.nih.gov/sra}{SRA home page} and paste
``SRR518692'' into the search box and hit search. The Run page for
SRR518692 shows links to the related BioProject/BioSample and basic
metadata (organism, instrument, library layout). In the ``Data access''
or ``Send to'' section, look for links to either ``FASTQ'' via the NCBI
Run Selector or ``SRA file download''. Keep in mind that for this small
file, you can download the compressed FASTQ directly through the
browser, but for reproducible workflows and large datasets you should
use the SRA Toolkit instead, which we will cover later in Lab 4.

\subsection{SRA Toolkit for Command Line
Retrieval}\label{sra-toolkit-for-command-line-retrieval}

The SRA Toolkit is a set of command-line programs that allow you to
download and convert SRA runs to FASTQ or FASTA. The two most relevant
tools for this course are:

\begin{itemize}
\tightlist
\item
  fastq-dump: older, widely used tool to convert SRA to FASTQ (now
  deprecated, but still available).
\item
  fasterq-dump: newer tool that is significantly faster and more
  memory-efficient, recommended by NCBI.
\end{itemize}

On HPC, always check local documentation for preferred options and
scratch storage guidelines.

\subsection{SRA Metadata and Its
Evolution}\label{sra-metadata-and-its-evolution}

SRA metadata has evolved alongside sequencing technologies and community
practices, gradually moving from sparse, free-text descriptions toward
more structured, machine-readable formats. Early submissions often
contained minimal or inconsistent sample information, which makes reuse
and integration across studies difficult, especially when key fields
such as strain, host, or geographic origin are missing or encoded only
in narrative text. In response, NCBI introduced structured
\textbf{BioSample packages} with controlled fields like host,
isolation\_source, and collection\_date, strengthened the links between
SRA runs and their associated BioProject and BioSample records, and
issued recommendations to rely more consistently on standard
vocabularies and ontologies where possible.

Despite these improvements, the overall metadata landscape in SRA
remains \textbf{heterogeneous}, particularly when older accessions are
considered alongside newer submissions. Historical datasets may lack
important fields or use idiosyncratic labels for the same concept (for
example, one group might use ``treatment'' while another uses
``condition'' for similar variables), and some submitters rely heavily
on file names or free-text notes to convey crucial information about
experimental design, which complicates downstream harmonization. As a
result, researchers reusing SRA data often need to spend substantial
effort cleaning, relabeling, and normalizing metadata before they can
perform robust comparative analyses across runs or projects.

For your own work, it is helpful to treat SRA metadata as part of the
primary scientific output rather than an afterthought. Choosing the
correct BioSample package for your experiment and filling out all
recommended fields produces records that are easier for both humans and
software to interpret. Consistent naming of samples across local files,
lab notebooks, and SRA accessions reduces confusion later. Thoroughly
documenting library preparation kits, read lengths, sequencing
platforms, and any unusual processing steps both in the SRA submission
and in your methods section further improves transparency and
reproducibility.

When mining existing SRA datasets for reanalysis, it is wise to assume
that metadata will require curation. A practical approach is to download
the run table or associated metadata, normalize field names to a
consistent internal schema, and explicitly check key variables such as
organism, library type, read layout, and sequencing platform before
combining runs into a single analysis. Older submissions, in particular,
should be treated cautiously; verifying sample identities and
experimental conditions by cross-referencing the associated publication,
or contacting the original authors when necessary, can prevent serious
misinterpretation of results. Ultimately, this variability in SRA
metadata is not just an annoyance but a concrete illustration of why
data literacy, careful documentation, and explicit provenance tracking
are essential skills for modern genome research.

\section{Visualizing FASTA Alignments and Tree
Files}\label{visualizing-fasta-alignments-and-tree-files}

Many of the databases introduced in this chapter provide downloadable
data that you can explore locally using desktop tools. At a minimum, you
will frequently encounter two key file types: FASTA files, which store
raw or aligned sequences, and \textbf{Newick files}, which encode
phylogenetic trees in a plain-text parenthetical format. FASTA files can
be downloaded from NCBI, Ensembl, and from organism-specific databases
such as FlyBase, SGD, or Phytozome. Newick tree files are produced by
many phylogenetic tools, but you can also download precomputed gene
trees from Ensembl's ``Gene tree'' view for individual gene families, or
species timetrees from resources like
\href{https://timetree.org/}{TimeTree}, which builds species-level
divergence trees from curated literature and provides them in Newick
format for downstream visualization.

A typical workflow is to obtain sequence FASTA from a genome database,
generate a multiple sequence alignment locally, infer a tree using a
phylogenetic software package, and then inspect the resulting Newick
tree in a dedicated tree viewer. For working with alignments, I
recommend AliView, a free, cross-platform alignment viewer and editor
designed to handle large phylogenetic datasets
efficiently\textsuperscript{44}. AliView reads many of the formats you
will encounter (FASTA, PHYLIP, NEXUS, etc.), renders the alignment with
color schemes that highlight nucleotide or amino acid patterns, and
supports basic editing operations such as trimming poorly aligned ends,
removing problematic sequences, and reordering taxa.

To use AliView in a typical gene-family analysis, you would first
download FASTA sequences for your gene of interest from NCBI, Ensembl,
or an organism-specific database, then open them in AliView and, if
necessary, pass them through an external aligner such as MAFFT (AliView
can be configured to call external programs). Once an alignment has been
generated and inspected, AliView can export the alignment in formats
suitable for phylogenetic inference, such as FASTA or PHYLIP, which can
then be used as input for tree-building software. Because AliView is
lightweight and easy to install on laptops, it serves as an accessible
bridge between database downloads and more advanced phylogenetic tools.

To move from an alignment to a tree, a commonly used tool is FastTree,
which infers approximately maximum-likelihood phylogenetic trees for
large nucleotide or protein alignments. FastTree is designed to be much
faster than traditional maximum-likelihood programs while maintaining
good accuracy, making it ideal for large gene families or species sets
where more computationally intensive methods would be impractical. A
typical command might look like running FastTree with a nucleotide
alignment in FASTA format, specifying an appropriate model such as GTR
for nucleotides, and writing the resulting tree to a file in Newick
format. AliView includes an ``external interface'' feature that can be
configured to send the current alignment to FastTree, wait for the
program to finish, and then automatically launch a tree viewer on the
resulting Newick file. In this way, you can go from a database-derived
FASTA file to a tree-ready Newick file in just a few integrated steps.

For visualizing and manipulating Newick trees, the standard tool
recommended here is FigTree, a graphical tree viewer and ``tree figure
drawing'' tool designed for interactive exploration and the production
of publication-ready figures. FigTree can open Newick files exported
from FastTree, BEAST, RAxML, or downloaded directly from Ensembl gene
trees or TimeTree species timetrees, and allows you to re-root trees,
collapse or expand clades, recolor branches, and label tips or internal
nodes based on metadata. A common workflow starts with a Newick tree
either generated locally from a multiple sequence alignment or
downloaded directly---for example, Ensembl's gene tree for PRDM9 or a
species tree from TimeTree built from a list of species names---and then
uses FigTree to adjust the layout (rectangular vs.~radial), annotate
clades, and export vector graphics for inclusion in reports or slides.
For gene trees, this allows you to interactively inspect relationships
among paralogs and orthologs, compare branch lengths, and connect the
patterns you see with the comparative genomics views in Ensembl and the
conservation tracks in genome browsers.

\begin{tcolorbox}[enhanced jigsaw, rightrule=.15mm, leftrule=.75mm, colbacktitle=quarto-callout-note-color!10!white, colframe=quarto-callout-note-color-frame, toprule=.15mm, opacitybacktitle=0.6, title=\textcolor{quarto-callout-note-color}{\faInfo}\hspace{0.5em}{Figure 4.1. Ali View External Tool Integration}, left=2mm, colback=white, bottomtitle=1mm, toptitle=1mm, bottomrule=.15mm, coltitle=black, titlerule=0mm, breakable, opacityback=0, arc=.35mm]

\begin{figure}[H]

\centering{

\pandocbounded{\includegraphics[keepaspectratio]{images/Fig_4.2_AliView.jpg}}

}

\caption{\label{fig-aliview}AliView Integration}

\end{figure}%

AliView, FastTree, and FigTree can all be combined into a simple,
\href{https://ormbunkar.se/aliview/\#FEATURE_LIST}{semi-automated
pipeline}. AliView's external command configuration supports defining
custom commands that take the current alignment, pipe it to FastTree to
infer a tree, and then open the resulting Newick file in FigTree. This
simple integration enables the ability to move fluidly from raw
sequences downloaded from NCBI, Ensembl, or model organism databases to
alignments and tree visualizations that you can manipulate and interpret
on your own laptop.

\end{tcolorbox}

\section{Galaxy as a web-based genome analysis
workbench}\label{galaxy-as-a-web-based-genome-analysis-workbench}

Though not a genome broswser, this chapter would not be complete without
touching on the web-based tool
\href{https://galaxy-main.usegalaxy.org/}{Galaxy}. Galaxy is a web-based
workbench that allows researchers to perform complex genome analyses
through a browser without writing code, addressing the ``informatics
crisis'' created by high-throughput sequencing\textsuperscript{45,46}.
By wrapping command-line tools in a consistent graphical interface and
automatically managing compute infrastructure and file formats, Galaxy
makes bioinformatics workflows accessible to experimentalists who might
never log into a cluster or install analysis
software\textsuperscript{47}.

\subsection{Historical development and design
goals}\label{historical-development-and-design-goals}

Galaxy originated in the mid-to-late 2000s in response to the rapid
expansion of microarray and next-generation sequencing data, when many
biologists lacked the training or infrastructure to analyze their own
data\textsuperscript{45,46}. Early papers described Galaxy as a
``web-based genome analysis tool for experimentalists'' and emphasized a
central design goal: making large-scale analyses biologist-friendly by
providing a unified, web-based interface to powerful tools while hiding
computational details\textsuperscript{47}.

\subsection{Doing bioinformatics in the
browser}\label{doing-bioinformatics-in-the-browser}

Galaxy's analysis workspace allows users to upload their own files or
import data from integrated sources, then chain tools together by
passing the output of one step directly into the next, constructing
multi-step workflows entirely via point-and-click
interfaces\textsuperscript{45}.

\begin{tcolorbox}[enhanced jigsaw, rightrule=.15mm, leftrule=.75mm, colbacktitle=quarto-callout-note-color!10!white, colframe=quarto-callout-note-color-frame, toprule=.15mm, opacitybacktitle=0.6, title=\textcolor{quarto-callout-note-color}{\faInfo}\hspace{0.5em}{Figure 4.2. Galaxy analysis workspace}, left=2mm, colback=white, bottomtitle=1mm, toptitle=1mm, bottomrule=.15mm, coltitle=black, titlerule=0mm, breakable, opacityback=0, arc=.35mm]

\begin{figure}[H]

\centering{

\pandocbounded{\includegraphics[keepaspectratio]{images/Fig_4.1_Galaxy.png}}

}

\caption{\label{fig-galaxy}Galaxy Workflow}

\end{figure}%

\textbf{Galaxy analysis workspace}. The Galaxy analysis workspace is
where users perform genomic analyses. The workspace has four areas: the
navigation bar, tool panel (left column), detail panel (middle column),
and history panel (right column). The navigation bar provides links to
Galaxy's major components, including the analysis workspace, workflows,
data libraries, and user repositories (histories, workflows, Pages). The
tool panel lists the analysis tools and data sources available to the
user. The detail panel displays interfaces for tools selected by the
user. The history panel shows data and the results of analyses performed
by the user, as well as automatically tracked metadata and
user-generated annotations. Every action by the user generates a new
history item, which can then be used in subsequent analyses, downloaded,
or visualized. Galaxy's history panel helps to facilitate
reproducibility by showing provenance of data and by enabling users to
extract a workflow from a history, rerun analysis steps, visualize
output datasets, tag datasets for searching and grouping, and annotate
steps with information about their purpose or importance. Here, a
download FASTQ file is being viewed in the detail panel by simply
clicking the eyeball button on the file in the history panel.

\end{tcolorbox}

The Galaxy toolkit illustrates how basic NGS data preparation can be
done entirely in the browser\textsuperscript{46}. Tools such as the
FASTQ Groomer, quality summary, trimmer, quality-based filters, and
paired-end joiner/splitter allow users to: convert between FASTA/QUAL
and all common FASTQ variants; compute per-base quality statistics and
visualize them; trim low-quality ends; filter reads by length or
quality; and perform complex pattern-based manipulations including
reverse complementing and adapter handling\textsuperscript{46}. These
tasks, traditionally the domain of custom scripts, become reproducible
``first steps'' in a teaching-friendly, GUI-based workflow that can be
reused across projects.

Beyond early QC, Galaxy supports complete analysis pipelines in many
domains, including RNA-seq, ChIP-seq, variant calling, and
metagenomics\textsuperscript{45}. Tool wrappers encapsulate standard
command-line software (e.g., mapping, peak calling, quantification) so
that a novice user can execute sophisticated end-to-end workflows by
selecting tools and parameters, rather than composing shell
commands\textsuperscript{45}.

\begin{tcolorbox}[enhanced jigsaw, rightrule=.15mm, leftrule=.75mm, colbacktitle=quarto-callout-note-color!10!white, colframe=quarto-callout-note-color-frame, toprule=.15mm, opacitybacktitle=0.6, title=\textcolor{quarto-callout-note-color}{\faInfo}\hspace{0.5em}{Box 4.2: Using SRA Toolkit inside Galaxy}, left=2mm, colback=white, bottomtitle=1mm, toptitle=1mm, bottomrule=.15mm, coltitle=black, titlerule=0mm, breakable, opacityback=0, arc=.35mm]

Many Galaxy instances provide command line tools like ``Faster Download
and Extract Reads in FASTQ'' (a wrapper for fasterq-dump) or ``NCBI SRA
Tools: fastq-dump''. For example, as mentioned above you retrieve the
\emph{Aeromonas hydrophila} dataset SRR518692 in Galaxy directly:

\begin{itemize}
\tightlist
\item
  Select the SRA tool in Galaxy.
\item
  Enter SRR518692 as the accession.
\item
  Choose ``Split files'' for paired-end.
\item
  Run the job to produce FASTQ datasets directly in your Galaxy history
  ( Figure~\ref{fig-galaxy}).
\end{itemize}

\end{tcolorbox}

\subsection{Reproducible workflows and
provenance}\label{reproducible-workflows-and-provenance}

A distinctive feature of Galaxy is its emphasis on reproducible
computational research\textsuperscript{45}. Every tool execution
generates provenance metadata that include the input datasets, tool
identity and version, parameter values, and resulting outputs, all
stored as steps in the user's history\textsuperscript{45}. Users can
inspect any intermediate dataset, rerun individual steps with modified
parameters, and branch analyses by copying histories, making it
straightforward to revise or expand an existing analysis.

\subsection{Transparency, sharing, and Galaxy
Pages}\label{transparency-sharing-and-galaxy-pages}

Galaxy also functions as a publishing medium for computational
analyses\textsuperscript{45}. Any dataset, history, or workflow can be
shared with collaborators, made accessible via a stable URL, or
published in a public repository where other users can search, view
metadata, and import items into their own
workspaces.\textsuperscript{45} This sharing model turns workflows into
community resources and promotes convergence on best-practice pipelines.

Galaxy Pages extend this concept by documenting an entire computational
experiment---from raw reads to final figures. For example, a metagenomic
study of ``windshield splatter'' used Galaxy to build a complete NGS
pipeline for phylogenetic profiling, then published a Page that exposes
the full workflow and histories as an executable online
supplement\textsuperscript{45}. In practice, this means that a methods
section and supplementary material can be instantiated as a live,
executable artifact rather than static text\textsuperscript{45}.

\section{Data Sharing Policies in
Genomics}\label{data-sharing-policies-in-genomics}

Genome data sharing policies have evolved from idealistic commitments to
rapid open release toward more nuanced systems that balance broad
access, participant protection, and the interests of data
producers\textsuperscript{48,49}. This section traces that trajectory,
uses the baboon genome as a case study in embargo-related bottlenecks,
and offers guidance on navigating current policies across human and
non‑human genomics.

\subsection{From Bermuda to managed
access}\label{from-bermuda-to-managed-access}

The Bermuda Principles (1996) emerged from the Human Genome Project
(HGP) and mandated that publicly funded human and model-organism
sequence be deposited in public databases within 24 hours, with daily
\textbf{pre‑publication} release into the public
domain\textsuperscript{48}. This approach framed genomic sequence as
part of a scientific commons and was explicitly justified as a way to
accelerate discovery and to limit restrictive gene patenting by using
rapid release as prior art.

Follow‑on statements such as the Fort Lauderdale agreement (2003)
reaffirmed early deposition but added clearer expectations that
downstream users respect the interests of data producers, especially
large sequencing centers and consortia. These early policies were
designed in an era when sequence data were mostly de‑identified, largely
non‑clinical, and generated by a relatively small number of centers,
which made it easier to sustain almost fully open release with only
informal norms to protect the ``first use'' of the data by generators.
This also meant that researchers largely understood the sequencing
bottleneck (see Figure~\ref{fig-bottleneck}) and were pro-active in
moving the science forward.

\subsection{Human genomics and the medical information
commons}\label{human-genomics-and-the-medical-information-commons}

As whole‑genome and exome sequencing became embedded in clinical care,
the notion of a medical information commons (MIC) reframed genome data
as one component of large, integrated resources containing genomic,
phenotypic, and environmental information. Cook‑Deegan and McGuire argue
that building a MIC requires going beyond Bermuda's simple ``sequence
into GenBank'' model toward federated systems that respect privacy,
regulatory heterogeneity, and data donors' expectations while still
enabling large‑scale reuse\textsuperscript{50,51}.

In the modern genome database landscape, ``open'' human genomic data are
increasingly rare; instead, most large‑scale human datasets are subject
to tiered or managed access, often governed by Data Access Committees
that vet applications, require data use agreements, and sometimes impose
limits on re‑identification, clinical interpretation, or re‑contact.
This shift reflects: (1) recognition that genomic data are inherently
identifying; (2) expansion to more diverse and often more vulnerable
populations; and (3) the embedding of genomic data within linked EHRs
and biobanks, making them qualitatively different from the relatively
``flat'' HGP sequences of the 1990s\textsuperscript{51}.

\subsection{Consortia, tiered access, and
embargoes}\label{consortia-tiered-access-and-embargoes}

Modern genomic consortia---spanning cancer, rare disease, population
genomics, agriculture, and microbial surveillance---use a wide variety
of access arrangements. Nicol, Nielsen, and Archer surveyed 98 genomic
consortia and identified twelve distinct access patterns, including
simple models such as open access, consortium‑only access, managed
access, and registered access, plus tiered combinations such as
``consortium + managed + open.'' Tiered arrangements often allow:
privileged access for consortium members, managed access for qualified
external users, and fully open release of a summary or low‑risk subset
of data\textsuperscript{49}.

Funding policies encode similar logic. The NIH Genomic Data Sharing
(GDS) Policy\textsuperscript{52} requires submission of large‑scale
human genomic data to approved repositories (e.g., dbGaP) and typically
defers broad release of controlled human data for about six months after
quality control, but explicitly rejects long publication embargoes in
favor of relatively rapid sharing with no formal guarantee of ``first
publication rights.'' For non‑human genomic data, NIH expects submission
no later than the time of publication, with fewer privacy constraints
but similar expectations that sequence and metadata be deposited in
recognized repositories.

Embargoes and ``privileged access'' periods remain widespread in
consortia as a way to recognize the investment of data producers. The
tiered access model can, in principle, honor rapid pre‑publication
release while giving generating groups a limited window to publish
primary analyses, but the effectiveness of this approach depends heavily
on clear documentation of embargo conditions and on the pace at which
data producers can move through their own analysis
bottlenecks\textsuperscript{49}.

\begin{tcolorbox}[enhanced jigsaw, rightrule=.15mm, leftrule=.75mm, colbacktitle=quarto-callout-important-color!10!white, colframe=quarto-callout-important-color-frame, toprule=.15mm, opacitybacktitle=0.6, title=\textcolor{quarto-callout-important-color}{\faExclamation}\hspace{0.5em}{Box 4.3: A cautionary tale - The Baboon Genome Project}, left=2mm, colback=white, bottomtitle=1mm, toptitle=1mm, bottomrule=.15mm, coltitle=black, titlerule=0mm, breakable, opacityback=0, arc=.35mm]

The baboon genome illustrates how informal embargo practices can stall
downstream science and even incentivize duplication of effort. Baylor
produced a reference assembly for the olive baboon (\emph{Papio anubis})
``Panu\_3.0'' and deposited under an informal embargo beginning around
2008, with access constrained under Fort Lauderdale--style expectations
that expected external researchers to defer primary genome‑scale
analyses to the generating group. While Rogers and colleagues at
Baylor's Human Genome Sequencing Center ultimately published this
reference genome, it was not until 2019\textsuperscript{53}, over a
decade after making it widely available.

In parallel, Wall et al.~published a 2016 study of admixture and
divergence in wild baboons using low‑coverage whole‑genome
sequencing\textsuperscript{54}. Because use of the Baylor reference
assembly for primary analyses was constrained by embargo conditions, the
group generated its own de novo reference, allowing it to publish
genome‑wide results without violating Fort Lauderdale norms or
undermining the Baylor team's planned flagship paper. A later paper from
the same team presented a new olive baboon assembly (Panubis1.0)
explicitly noting that Panu\_3.0 had ``negligible'' scientific influence
in part because it was effectively embargoed for more than a decade,
despite being the primary reference for the species\textsuperscript{55}.
The result was a situation in which NIH‑funded sequencing efforts
supported two separate reference‑quality assemblies and duplicated core
aspects of the work, with the later Panubis1.0 paper arguing that
several large syntenic discrepancies between Panu\_3.0 and Panubis1.0
reflect genuine assembly errors in Panu\_3.0.

This case study illustrates several key points about data sharing and
publication timing:

\begin{itemize}
\tightlist
\item
  Long, informal embargoes can effectively remove a reference genome
  from the scientific arena, even when the sequence is technically
  deposited at NCBI or similar repositories\textsuperscript{55}.
\item
  When primary producers cannot move from sequencing to analysis and
  publication quickly, other groups may either avoid the dataset
  entirely or reproduce the underlying sequencing effort in order to
  publish their own work\textsuperscript{54,55}.
\item
  Embargo‑protected references can interact poorly with the rapid pace
  of high‑throughput sequencing: the sequencing bottleneck has largely
  disappeared, but the analysis bottleneck remains, creating a period
  during which data exist but cannot be freely used for the most
  scientifically valuable questions\textsuperscript{50}.
\end{itemize}

\end{tcolorbox}

\subsection{Practical guidance for working with genome data
today}\label{practical-guidance-for-working-with-genome-data-today}

Across both human and non‑human genomics, the early‑deposit norm means
that reference genomes, resequencing data, and large cohort datasets
often appear in repositories months to years before the corresponding
flagship papers. This practice accelerates method development,
cross‑project integration, and secondary discovery, but it also raises
concrete risks for researchers who do not carefully assess data‑use
conditions (see Box 4.1 above for cautionary tale).

When working with human or non‑human genomic datasets, it is important
to:

\begin{itemize}
\tightlist
\item
  Check the data‑sharing policy and access tier. Repositories such as
  GenBank, SRA, ENA, and non‑human assemblies in NCBI usually host
  genuinely open sequence, but human genomic datasets in dbGaP or EGA
  are almost always subject to managed or registered access, with
  explicit Data Access Committee oversight and data use
  limitations\textsuperscript{49}.
\item
  Examine whether any publication embargo, moratorium, or ``first
  publication'' expectation applies, particularly for data labeled as
  coming from a consortium, a sequencing center, or a flagship
  project\textsuperscript{49}.
\item
  Look for documentation that specifies which analyses are reserved for
  the generating consortium (e.g., full genome‑wide association
  analyses, primary pan‑genome construction, or first comprehensive
  comparative genomics) and which are encouraged for external
  users\textsuperscript{49}.
\item
  Cross‑check whether closely related genome projects are underway that
  might produce overlapping analyses; in light of cases like the baboon
  reference, investigators considering large investments in new
  assemblies or resequencing may wish to contact data generators and
  funders early to clarify expectations and avoid inadvertent
  duplication\textsuperscript{55}.
\end{itemize}

For students and trainees, an important lesson is that ``open'' in
genomics is not a single category. Sequence reads in an NCBI project may
be technically accessible but still subject to norms, policies, or
formal embargoes that shape what is considered appropriate use. At the
same time, as Cook‑Deegan and McGuire emphasize, the long‑term goal
remains a more robust medical information commons and comparable shared
infrastructure for non‑human genomics, in which broad, responsible data
sharing supports both collective discovery and fair recognition of the
labor and resources that generate large genomic
datasets\textsuperscript{49--51}.

\section{Reflection Questions}\label{reflection-questions-3}

\begin{enumerate}
\def\labelenumi{\arabic{enumi}.}
\item
  \textbf{Genome database evolution and curation trade-offs}: Trace the
  evolution from GenBank to RefSeq; why does NCBI curate specific
  entries like HBB NM\_000518 rather than simply accepting all
  submissions? What are the trade-offs between GenBank's archival
  comprehensiveness and RefSeq's non-redundant, reviewed models, and
  when might you prefer one over the other in your own research?
\item
  \textbf{Orthology interpretation in comparative genomics}: Compare the
  PRDM9 orthologs shown in UCSC and Ensembl; why does Ensembl label some
  relationships as ``1:many'' rather than simple one-to-one, and how
  does this reflect the underlying gene tree inference process?
\item
  \textbf{Coordinate systems and variant reporting}: How do coordinate
  shifts between hg19 and hg38 (or other major builds) affect variant
  reporting, such as the position of HBB? Using a specific example from
  GDV or UCSC, explain how you would verify the exact position of a
  variant across builds and why this matters for reproducibility.
\item
  \textbf{Browser choice for specific questions}: NCBI GDV, Ensembl, and
  UCSC each excels in different areas; when would you choose GDV over
  Ensembl for annotating PRDM9 in a newly sequenced genome, and what
  limitations might make you pivot to UCSC for custom track
  visualization?
\item
  \textbf{Data sharing policies and rapid deposition}: Should we deposit
  ALL sequencing data immediately upon generation, even before
  publication? Why or why not? Is it surprising that early genomics
  followed the Bermuda Principles (24-hour data release), while modern
  consortia use tiered access? What factors changed (e.g., clinical
  privacy, commercial interests), and how does this impact
  reproducibility in model organism research?
\end{enumerate}

\section{Problem Sets}\label{problem-sets}

\begin{enumerate}
\def\labelenumi{\arabic{enumi}.}
\item
  \textbf{NCBI Gene and EDirect}: Search NCBI Gene for ``DSX
  Drosophila''; use EDirect to fetch the FASTA sequence for the primary
  transcript
  (\texttt{esearch\ -db\ gene\ -query\ "DSX{[}Gene{]}\ AND\ Drosophila"\ \textbar{}\ efetch\ -format\ fasta}).
  Load into AliView and identify the number of isoforms; document the
  stable RefSeq ID and how it differs from GenBank accessions.
\item
  \textbf{UCSC Conservation in Drosophila}: In UCSC (\emph{D.
  melanogaster} dm6), navigate to the \texttt{rdl} gene; enable MultiZ
  alignments for 4 species (\emph{D. simulans}, \emph{D. pseudoobscura},
  \emph{D. grimshawi}) and phyloP tracks. Zoom to exons and interpret
  how conservation profiles differ between coding vs.~non-coding
  regions; export a screenshot or BED of conserved elements and explain
  their potential function.
\item
  \textbf{Galaxy SRA Pipeline}: In Galaxy, retrieve SRR518692
  (\emph{Aeromonas hydrophila}); run FASTQC to perform quality control
  on the files downloaded.\\
\item
  \textbf{FlyBase Synteny}: In FlyBase JBrowse (\emph{D. melanogaster}),
  explore Muller elements synteny table; compare chromosome
  rearrangements across \emph{D. grimshawi} and \emph{D. pseudoobscura}.
  Export orthologs for ADH as FASTA and align in AliView.
\item
  \textbf{Phytozome Pectinase Phylogeny}: In Phytozome, search
  ``Pectinesterase'' in Brassicaceae; download FASTA for 10-20 hits.
  Download the protein sequences and import into AliView. As shown in
  Figure~\ref{fig-aliview} align the downloaded sequences and run
  FastTree → FigTree. Compare the resulting gene tree to species tree;
  hypothesize duplication timing and test with Ensembl orthologs.
\end{enumerate}

\section{References}\label{references-4}

\chapter{Sequence Alignment}\label{sequence-alignment}

Pairwise and Multiple Sequence Alignment Algorithms

\hfill\break

\section{Learning goals}\label{learning-goals}

By the end of this chapter you should be able to:

\begin{itemize}
\tightlist
\item
  Explain why biologists align sequences and distinguish
  \textbf{similarity} from homology, orthologs, and paralogs.
\item
  Describe how pairwise alignment scores are computed from substitution
  matrices and gap penalties.
\item
  Outline and interpret the Needleman--Wunsch global alignment algorithm
  and contrast it with Smith--Waterman local alignment.
\item
  Explain the difference between a dynamic versus heuristic programming.
\item
  Explain what a multiple sequence alignment (MSA) is, why it is more
  informative than pairwise alignment, and summarize major MSA algorithm
  families.
\item
  Read alignment output critically, recognizing poor alignments and the
  influence of parameter choices.
\end{itemize}

\section{Why align sequences?}\label{why-align-sequences}

Sequence alignment is a core operation in genomics, used whenever you
want to compare DNA or protein sequences to ask biological
questions\textsuperscript{10}. Typical motivations include inferring
gene or protein function for a new sequence by comparing it to known
sequences, identifying orthologous genes across species for phylogeny
reconstruction, and detecting conserved motifs and domains that may
underlie shared structure or regulation\textsuperscript{10}. This
chapter covers pairwise sequence alignment (both local and global)
before extending these concepts to multiple sequence alignment, where
information from many sequences is integrated to reveal patterns
invisible in pairwise comparisons\textsuperscript{10,56}.

An algorithm is a step-by-step set of instructions used to solve a
problem---a computational protocol. Understanding the algorithms
commonly used in sequence analysis helps you recognize their key
assumptions and limitations, unveiling the black box mentioned in
earlier chapters\textsuperscript{10}.

\subsection{Similarity, homology, orthologs,
paralogs}\label{similarity-homology-orthologs-paralogs}

Alignment provides a \emph{similarity} score, but what biologists
usually care about is \textbf{homology}, meaning shared ancestry rather
than just high percent identity\textsuperscript{10}. Two sequences are
either homologous or not; there is no partial homology, only stronger or
weaker evidence for it\textsuperscript{10}. \textbf{Orthologs} are
homologous genes in different species that diverged via a speciation
event and usually retain similar function, whereas \textbf{paralogs} are
homologous genes within a genome that arose by duplication and often
diverge in function over time\textsuperscript{10}. When you use
alignment to infer orthology or paralogy, remember that homology is
about evolutionary history, not just alignment score, and typically
requires integrating tree information as well.

\subsection{Amino acids versus
nucleotides}\label{amino-acids-versus-nucleotides}

With nucleotides offering only four different bases, the information
content for homology is low, and codon sequences can be degenerate in
how they code for amino acids\textsuperscript{10}. In contrast, protein
sequences offer up to 20 different amino acid types, vastly increasing
the information content\textsuperscript{10}. For identifying homologs,
protein sequence comparison is therefore considered superior, providing
a deeper look back in time and enabling identification of homologs
across the tree of life\textsuperscript{10}. Exceptions include
noncoding DNA regions such as introns and intergenic sequences, as well
as studies of DNA polymorphisms that distinguish different
populations\textsuperscript{10}.

\subsection{Types of sequence changes}\label{types-of-sequence-changes}

Single nucleotide substitutions are often referred to as either SNPs, as
you saw earlier in this course. Another term you might encounter is
\textbf{single nucleotide variants} (SNVs), a term used more recently to
indicate a low frequency variant, often used more in clinical settings.
Additionally, you might encounter the term multi-nucleotide
polymorphisms (MNPs). In addition to substitutions, there are short
insertion/deletion variants, or \textbf{indels}. Large scale
rearrangements are much harder to study given that they violate a lot of
the assumptions of most sequencing alignment algorithms. Therefore, we
will focus on small changes. It is also worth noting that more
algorithms only allow 2 possible variants at any one position. It is
important to understand the limitations of the algorithms you use when
working with datasets that are complex.

\begin{tcolorbox}[enhanced jigsaw, rightrule=.15mm, leftrule=.75mm, colbacktitle=quarto-callout-note-color!10!white, colframe=quarto-callout-note-color-frame, toprule=.15mm, opacitybacktitle=0.6, title=\textcolor{quarto-callout-note-color}{\faInfo}\hspace{0.5em}{Lab Connections}, left=2mm, colback=white, bottomtitle=1mm, toptitle=1mm, bottomrule=.15mm, coltitle=black, titlerule=0mm, breakable, opacityback=0, arc=.35mm]

Later in \href{Lab09.qmd}{Lab 9}, we will visualize sequence variants
relative to a reference genome in IGV and describe these sequence
changes using Human Genome Variation Society (HGVS) standardized
nomenclature\textsuperscript{10}. Specifically, a sequence change in the
DNA may or may not impact the protein sequence. Substitutions are
designated by a ``\textgreater{}'' character with the original base
before and the substitution after, prefixed by ``g.'' for genomic
sequence (e.g., \texttt{g.69A\textgreater{}T} for A to T substitution at
position 69) and ``p.'' for protein sequence (e.g., \texttt{p.K23Y} for
lysine to tyrosine substitution at position 23).

\end{tcolorbox}

\section{Pairwise Sequence Alignment}\label{pairwise-sequence-alignment}

Pairwise sequence alignment arranges two DNA or protein sequences to
assess similarity and infer common evolutionary
origin\textsuperscript{10,57,58}. Given a pair of sequences, an
alignment algorithm inserts gaps as needed to maximize a numerical score
defined by a scoring scheme, allowing different alignments to be
compared quantitatively\textsuperscript{10}. \textbf{Global alignment}
attempts to align sequences from end to end, most appropriate when
sequences are expected to be collinear and of similar
length\textsuperscript{10,57}. \textbf{Local alignment} searches for the
highest-scoring subsequences without requiring full alignment, ideal
when proteins share only conserved domains or when short motifs are
embedded in otherwise unrelated context\textsuperscript{10,58}.

\begin{tcolorbox}[enhanced jigsaw, rightrule=.15mm, leftrule=.75mm, colbacktitle=quarto-callout-note-color!10!white, colframe=quarto-callout-note-color-frame, toprule=.15mm, opacitybacktitle=0.6, title=\textcolor{quarto-callout-note-color}{\faInfo}\hspace{0.5em}{Figure 5.1 -- Interpreting Sequence Alignments}, left=2mm, colback=white, bottomtitle=1mm, toptitle=1mm, bottomrule=.15mm, coltitle=black, titlerule=0mm, breakable, opacityback=0, arc=.35mm]

\begin{figure}[H]

\centering{

\includegraphics[width=0.75\linewidth,height=\textheight,keepaspectratio]{images/Fig_5.1_Pairwise_PRDM9.png}

}

\caption{\label{fig-exSeq}Sequence Alignment of PRDM9}

\end{figure}%

In the example pairwise sequence above, we see the results of a global
sequence alignment between the PRDM9 protein from humans and that of
\emph{Lacerata agilis} which was manually annotated by Chipper Johnson,
course LA. Here, vertical lines \texttt{\textbar{}} indicate identical
resides. Substitutions are indicated by one or two dots, indicating the
degree of dissimilarity. Gaps can relfect either an insertion in one of
the sequences or a deletion in the other one. These are indicated by the
\texttt{-} symbol. This output is from the
\href{https://www.ebi.ac.uk/jdispatcher/psa/emboss_needle}{EMBOSS Needle
tool} that you will be using for your manual annotation projects this
semester.

\end{tcolorbox}

In genome alignment (more on this in the next chapter), it is common to
describe a ``sample sequence'' being compared to a ``reference genome
sequence,'' but in more general alignment problems the two inputs are
called the \textbf{query} sequence and the \textbf{target} (or subject)
sequence\textsuperscript{10}.

\subsection{Scoring Alignments}\label{scoring-alignments}

Alignment scores combine a substitution matrix assigning scores to
aligned residue pairs and a gap penalty scheme for insertions and
deletions\textsuperscript{10,59,60}. For proteins, substitution matrices
encode evolutionary and biochemical information: conservative
substitutions between similar residues receive positive scores, while
rare or structurally disruptive substitutions receive negative
scores\textsuperscript{10,59}. Not all identities are scored equally;
matching two tryptophans is more informative than matching two alanines
because tryptophan is rare and strongly
conserved\textsuperscript{10,60}.

PAM (Point Accepted Mutation) matrices, introduced by Margaret Dayhoff,
are derived from curated alignments of closely related protein families
and model substitution probabilities as proteins diverge over
evolutionary time\textsuperscript{10,59}. PAM1 reflects frequencies
after one accepted mutation per 100 residues, while higher PAM matrices
(e.g., PAM250) approximate more divergent
sequences\textsuperscript{10,59}. BLOSUM (Blocks Substitution Matrix)
series from Steven and Jorja Henikoff focuses on conserved blocks of
distantly related families, with BLOSUM62 widely used as a default in
database searches\textsuperscript{10,60}. Low PAM numbers and high
BLOSUM numbers emphasize close relationships, while high PAM numbers and
low BLOSUM numbers better detect distant
homology\textsuperscript{10,59,60}.

\begin{tcolorbox}[enhanced jigsaw, rightrule=.15mm, leftrule=.75mm, colbacktitle=quarto-callout-note-color!10!white, colframe=quarto-callout-note-color-frame, toprule=.15mm, opacitybacktitle=0.6, title=\textcolor{quarto-callout-note-color}{\faInfo}\hspace{0.5em}{Margaret Dayhoff and the Foundations of Scoring Matrices}, left=2mm, colback=white, bottomtitle=1mm, toptitle=1mm, bottomrule=.15mm, coltitle=black, titlerule=0mm, breakable, opacityback=0, arc=.35mm]

Margaret Oakley Dayhoff (1925--1983) was a physical chemist and
computational biologist whose work laid the groundwork for modern
sequence alignment scoring\textsuperscript{10,59}. From carefully
curated multiple sequence alignments of closely related proteins, she
defined accepted point mutations and quantified amino-acid substitution
frequencies, leading to the PAM model and the first empirically grounded
mutation probability matrices\textsuperscript{10,59}. By converting
these probabilities into log-odds scores, Dayhoff introduced a rigorous
framework that rewards substitutions observed more often than expected
by chance and penalizes unlikely changes\textsuperscript{10,59}. This
logic underlies essentially all modern substitution matrices and
directly informs algorithms such as Needleman--Wunsch, Smith--Waterman,
and BLAST\textsuperscript{10,57,58}.

\begin{figure}[H]

\centering{

\includegraphics[width=0.25\linewidth,height=\textheight,keepaspectratio]{images/dayhoff.jpg}

}

\caption{\label{fig-dayhoff}Margaret Dayhoff)}

\end{figure}%

\end{tcolorbox}

Nucleotide substitution matrices are usually simpler, because there are
only four bases and fewer distinct biochemical classes than for amino
acids\textsuperscript{10}. In many applications, DNA alignments are
scored with uniform match/mismatch schemes (for example, +1 for a match
and -1 for a mismatch) or with slightly more refined scoring where
transitions (A↔G, C↔T) receive different penalties from transversions
(purine↔pyrimidine), reflecting the higher frequency of transitions in
real genomic data\textsuperscript{10}. More elaborate nucleotide scoring
systems, such as matrices derived from observed substitution patterns in
genomic regions, exist but are less commonly used than protein
substitution matrices\textsuperscript{10}. Because protein sequences
carry higher information content and can leverage sophisticated
amino-acid substitution matrices, it is often preferable to translate
coding DNA sequences and perform alignments at the protein level when
the goal is to detect homology across moderate or large evolutionary
distances\textsuperscript{10}.

\subsection{Gap penalties}\label{gap-penalties}

Insertions and deletions (indels) in evolutionary history appear as gaps
in sequence alignments and must be penalized so that algorithms do not
introduce them without strong justification\textsuperscript{10}. Without
gap penalties, an alignment algorithm could trivially achieve perfect
matches by inserting arbitrarily many gaps, so a realistic scoring
system must balance the benefits of aligning additional residues against
the cost of adding gaps\textsuperscript{10}. Biologically, single long
indels are often more plausible than many short indels, and gap penalty
schemes are designed to reflect that intuition\textsuperscript{10}.

In the simplest scheme, a constant (or linear) gap penalty subtracts the
same score for each gap position, so a gap of length \((k)\) incurs a
cost proportional to \((k)\)\textsuperscript{10}. In this case, the cost
of opening a gap and the cost of extending it are identical, so a gap of
length 5 is exactly five times more costly than a gap of length
1\textsuperscript{10}. More realistic alignments, especially for
proteins, generally use an \textbf{affine gap penalty}, which separates
the penalty for initiating a gap from the penalty for extending
it\textsuperscript{10}. Under an affine model, the cost of a gap of
length \((k)\) is given by \[gap cost = a + b k,\]

where \((a)\) is a gap opening penalty and \((b)\) is a gap extension
penalty per residue, with \((|a|)\) typically larger than
\((|b|)\)\textsuperscript{10}. This favors fewer, longer gaps over many
short ones, mirroring the expectation that a single deletion of several
residues is more likely than multiple independent single-residue
events\textsuperscript{10}.

When you run alignment tools on the web or command line, default gap
opening and extension penalties are often optimized for particular
scoring matrices and use cases\textsuperscript{10}. Understanding how
these parameters influence alignments helps explain phenomena such as
long conserved blocks separated by large single gaps or the tendency of
some parameter combinations to fragment alignments into many short local
matches\textsuperscript{10}. In many practical workflows, end gaps
(leading and trailing gaps in global alignment) are either not penalized
or penalized differently from internal gaps, especially when one
sequence is known to be truncated relative to the
other\textsuperscript{10}.

\subsection{Needleman--Wunsch: Global
Alignment}\label{needlemanwunsch-global-alignment}

The Needleman--Wunsch algorithm, introduced in 1970, is a
\textbf{dynamic programming} method that guarantees an optimal global
alignment for two sequences given a specified substitution matrix and
gap penalty scheme\textsuperscript{10,57}. The core idea is to
systematically evaluate all possible ways of aligning prefixes of the
two sequences by filling a matrix of scores and then to recover the best
full-length alignment by tracing back through that
matrix\textsuperscript{10,57}. This approach avoids exhaustive
enumeration of all possible alignments, which would grow combinatorially
with sequence length, while still ensuring that the final alignment is
optimal under the chosen scoring system\textsuperscript{10,57}. Below is
a description of this algorithm, but if you are more of a visual
learner, there is a great 15 minute
\href{https://www.youtube.com/watch?v=LhpGz5--isw}{YouTube video} on the
topic\textsuperscript{61}.

Given sequences of lengths (m) and (n), Needleman--Wunsch constructs a
matrix of dimensions ((m+1) \times (n+1)), where the cell (F\_\{i,j\})
represents the best possible global alignment score for the first (i)
residues of sequence 1 and the first (j) residues of sequence
2\textsuperscript{10,57}. The first row and first column are initialized
with cumulative gap penalties, allowing either sequence to begin with a
terminal gap if that yields a higher overall score\textsuperscript{10}.
The matrix is then filled row by row or column by column using a
recurrence relation that considers the three possible ways to reach each
cell: aligning the current pair of residues (diagonal move), inserting a
gap in one sequence (vertical move), or inserting a gap in the other
(horizontal move)\textsuperscript{10,57}.

In its simplest form with linear gap penalties, the recurrence for
\((F_{i,j})\) can be written as

\[F_{i,j} = \max \begin{cases}
F_{i-1,j-1} + s(x_i, y_j) \\
F_{i-1,j} + g \\
F_{i,j-1} + g
\end{cases}
\]

where \((s(x_i, y_j))\) is the substitution score for aligning residues
\((x_i)\) and \((y_j)\), and \((g)\) is the (negative) gap
penalty\textsuperscript{10,57}. More elaborate formulations using affine
gap penalties introduce additional matrices or states to distinguish gap
opening from gap extension, but the underlying dynamic programming
principle is the same\textsuperscript{10}. Once the matrix is filled,
the optimal global alignment score is found in the bottom-right cell
\((F_{m,n})\), and a traceback procedure follows the path of choices
that led to this score, reconstructing the explicit alignment with
matches, mismatches, and gaps\textsuperscript{10,57}.

\begin{tcolorbox}[enhanced jigsaw, rightrule=.15mm, leftrule=.75mm, colbacktitle=quarto-callout-note-color!10!white, colframe=quarto-callout-note-color-frame, toprule=.15mm, opacitybacktitle=0.6, title=\textcolor{quarto-callout-note-color}{\faInfo}\hspace{0.5em}{Figure 5.2 -- Pairwise Sequence Alignment}, left=2mm, colback=white, bottomtitle=1mm, toptitle=1mm, bottomrule=.15mm, coltitle=black, titlerule=0mm, breakable, opacityback=0, arc=.35mm]

\begin{figure}[H]

\centering{

\includegraphics[width=0.75\linewidth,height=\textheight,keepaspectratio]{images/Figure_3.20_Pevsner.jpg}

}

\caption{\label{fig-pairwise}Pairwise Sequence Alignment (from Pevsner
3.20)}

\end{figure}%

Pairwise alignment of two amino acid sequences using a dynamic
programming algorithm of Needleman and Wunsch (1970) for global
alignment. (a) Two sequences can be assigned a diagonal path through the
matrix and, when necessary, the path can deviate horizontally or
vertically, reflecting gaps that are introduced into the alignment. (b)
Two identical sequences form a path on the matrix that fits a diagonal
line. (c) If there is a mismatch (or multiple mismatches), the path
still follows a diagonal, although a scoring system may penalize the
presence of mismatches. If the alignment includes a gap in (d) the first
sequence or (e) the second sequence, the path includes a vertical or
horizontal line.

\end{tcolorbox}

\subsection{Smith--Waterman: Local
Alignment}\label{smithwaterman-local-alignment}

The Smith--Waterman algorithm, introduced in 1981, adapts the
\textbf{dynamic programming} framework to local rather than global
alignment, focusing on the highest-scoring subsequences of the two
inputs\textsuperscript{10,58}. Instead of forcing an alignment from the
beginning to the end of both sequences, Smith--Waterman identifies the
best local region of similarity, which is especially useful for
detecting conserved domains, motifs, or exons embedded in longer, less
related sequence contexts\textsuperscript{10,58}. This local strategy is
also central to database search tools such as BLAST and FASTA, which
conceptually perform many local alignments between a query and entries
in a database\textsuperscript{10}. As with global alignment, below is a
description of this algorithm, but if you are more of a visual learner,
there is a short 2 minute
\href{https://www.youtube.com/watch?v=F1YId_z4pc8}{YouTube video} on the
topic\textsuperscript{62}.

Smith--Waterman uses a similar ((m+1) \times (n+1)) dynamic programming
matrix as Needleman--Wunsch, but modifies the recurrence relation in two
key ways\textsuperscript{10,58}. First, the first row and first column
are initialized to zero rather than cumulative gap penalties, reflecting
the fact that local alignments may begin anywhere in the sequences
without penalty\textsuperscript{10,58}. Second, at each cell the
algorithm introduces a fourth option in addition to the three used for
global alignment: the score zero. The recurrence becomes

\[F_{i,j} = \max \begin{cases}
0 \\
F_{i-1,j-1} + s(x_i, y_j) \\
F_{i-1,j} + g \\
F_{i,j-1} + g
\end{cases}\]

ensuring that scores never become negative\textsuperscript{10,58}. Any
path that would yield a negative score is instead truncated at zero,
effectively ``resetting'' the alignment and allowing new local
alignments to begin internally\textsuperscript{10,58}.

Once the matrix is filled, the starting point for traceback in
Smith--Waterman is not the bottom-right cell but the cell with the
highest score anywhere in the matrix, which represents the end of the
best local alignment\textsuperscript{10,58}. Traceback then proceeds
diagonally, vertically, or horizontally following the choices that
produced that score, terminating when a cell with score zero is
reached\textsuperscript{10,58}. The resulting alignment spans only the
high-scoring segment pair (or pairs) and ignores lower-scoring flanking
regions, often producing a shorter alignment with higher percent
identity and similarity than the corresponding global alignment of the
same sequences\textsuperscript{10,58}.

In protein database searching, local alignments are essential because
many protein families share only one or a few conserved domains while
differing elsewhere, and because databases contain partial sequences,
fragments, and alternative isoforms that would not align well under a
strict global scheme\textsuperscript{10}. Heuristic methods such as
FASTA and BLAST approximate Smith--Waterman's local alignment behavior
while greatly reducing runtime by focusing computation on promising
regions of the search space, enabling rapid comparison of a query
sequence to millions of database entries\textsuperscript{10,58}.

\subsection{Basic Local Alignment Search Tool
(BLAST)}\label{basic-local-alignment-search-tool-blast}

Introduced on the first day of class, BLAST is the most used
bioinformatics tool. No doubt ALL of you have used this tool at some
point prior to taking this class at least once if not multiple times.
The basic local alignment search tool was first published in 1990 with
nearly 80,000 citations to date (last updated 1/19/2026).

\begin{tcolorbox}[enhanced jigsaw, rightrule=.15mm, leftrule=.75mm, colbacktitle=quarto-callout-tip-color!10!white, colframe=quarto-callout-tip-color-frame, toprule=.15mm, opacitybacktitle=0.6, title=\textcolor{quarto-callout-tip-color}{\faLightbulb}\hspace{0.5em}{Getting citation metrics from Web of Science}, left=2mm, colback=white, bottomtitle=1mm, toptitle=1mm, bottomrule=.15mm, coltitle=black, titlerule=0mm, breakable, opacityback=0, arc=.35mm]

One of the best ways to assess the importance of a bioinformatics tool
is to check the citation metrics. The best to to use here is not google
but instead \href{https://www.webofscience.com/wos/}{Web of Science}.

Here, you can search a paper (e.g.~``Altschul 1990 basic local alignment
search tool''). Once you have the paper, you can check the box next to
the result that you want to analyze. Then, in the top right, click
``Citation Report''. In addition to the total number of citations
(e.g.~79,208), you can also analyze the citations in depth by clicking
the word ``Analyze'' next to the number of citing articles. This new
page offers a drop down menu with several ways to analyze the citations,
either by time (e.g.~Publication Years), journal (e.g.~Publication
Titles), by Web of Science Categories, Funding Agencies, etc.

I find the analysis over time to be particularly useful. For example,
you might find a particular tool has a LOT of citations, but that the
number have dwindled in recent years. A deeper dive might find that a
newer tool has largely replaced this tool. This information together
would suggest that you should try to use the newer tool.

In the bioinformatics approach for your group project, I expect you to
use these types of metrics to justify your selected bioinformatics
tools, including ones you might have considered and decided against.

\end{tcolorbox}

BLAST, first published in 1990, takes a query sequence and searches a
database for entries sharing significant local similarity, outputting a
ranked list of hits annotated with scores and statistics such as bit
score and expect value (E value)\textsuperscript{10,63}.

\begin{tcolorbox}[enhanced jigsaw, rightrule=.15mm, leftrule=.75mm, colbacktitle=quarto-callout-note-color!10!white, colframe=quarto-callout-note-color-frame, toprule=.15mm, opacitybacktitle=0.6, title=\textcolor{quarto-callout-note-color}{\faInfo}\hspace{0.5em}{Connection to Lab 5}, left=2mm, colback=white, bottomtitle=1mm, toptitle=1mm, bottomrule=.15mm, coltitle=black, titlerule=0mm, breakable, opacityback=0, arc=.35mm]

In \href{Lab05.qmd}{Lab 5}, we will use BLAST on the command line using
the Alabama Super Computer\textsuperscript{64}. The command-line version
(BLAST+) allows custom database searches, bulk queries, and complex
search strategies essential for large-scale genomic
data\textsuperscript{64}.

\end{tcolorbox}

BLAST builds on Smith--Waterman's local alignment framework but uses
\textbf{heuristic programming} rather than exhaustive dynamic
programming, trading small losses in theoretical optimality for massive
gains in speed\textsuperscript{10,63}. The algorithm proceeds in three
main phases\textsuperscript{10,63,64}:

\textbf{Phase 1 (List)}: BLAST analyzes the query and compiles short
words of fixed length \((w)\) (typically 3 for
proteins)\textsuperscript{10,63}. Each possible word pair is scored
using the substitution matrix, and only words scoring at or above
threshold \((T)\) are retained\textsuperscript{10,63}. This filters for
exact or conservative substitutions characteristic of conserved
motifs\textsuperscript{10}.

\textbf{Phase 2 (Scan and Extend)}: BLAST scans the database for
retained words, using a ``two-hit'' strategy requiring two
non-overlapping hits within window \((A)\) to initiate
extension\textsuperscript{10}. Extensions create high-scoring segment
pairs (HSPs), first ungapped then optionally gapped\textsuperscript{10}.
HSPs exceeding reporting cutoff \((S)\) are
retained\textsuperscript{10}.

\textbf{Phase 3 (Traceback)}: BLAST reconstructs explicit local
alignments for retained HSPs, assigning matches, mismatches, and gaps
according to scoring parameters\textsuperscript{10}.

Word size \((w)\) and threshold \((T)\) control sensitivity--speed
trade-offs\textsuperscript{10,63}. For proteins, default \((w=3)\)
balances sensitivity and computational cost; lowering \((T)\) increases
hits and extensions but requires more time\textsuperscript{10}. For
nucleotides, BLASTN requires exact word matches with larger default word
size (e.g., 11), with smaller words increasing sensitivity and larger
words speeding searches\textsuperscript{10}.

BLAST employs rigorous statistical evaluation based on extreme value
distribution (EVD) theory\textsuperscript{10,63}. The E value represents
the expected number of distinct alignments with scores at least as high
as \((S)\) occurring by chance in a database search\textsuperscript{10}:

\[E = K m n e^{-\lambda S}\]

where \((m)\) is query length, \((n)\) is effective database length, and
\((K)\) and \(\lambda\) are Karlin--Altschul
statistics\textsuperscript{10}. Smaller E values indicate more
significant matches, and E values depend on search space
size\textsuperscript{10}. The related p value is \((p = 1 - e^{-E})\);
for very small E values, \((E)\) and \((p)\) are nearly
identical\textsuperscript{10}.

In practice, BLAST users adopt thresholds such as \((E \le 0.05)\) as
minimal significance standards, with more stringent criteria (e.g.,
\((10^{-3})\), \((10^{-5})\)) used in large-scale genome
analyses\textsuperscript{10}. Biological interpretation always requires
combining E values with domain knowledge: similar sequence lengths,
conserved motifs, compatible domain structures, and congruent genomic
context strengthen homology inferences\textsuperscript{10}.

\begin{longtable}[]{@{}ll@{}}
\caption{Relationship of E to p values in BLAST. Small E values (0.05 or
less) correspond closely to p values.}\label{tbl-evalues}\tabularnewline
\toprule\noalign{}
E value & p value \\
\midrule\noalign{}
\endfirsthead
\toprule\noalign{}
E value & p value \\
\midrule\noalign{}
\endhead
\bottomrule\noalign{}
\endlastfoot
10 & 0.99995460 \\
5 & 0.99326205 \\
2 & 0.86466472 \\
1 & 0.63212056 \\
0.1 & 0.09516258 \\
0.05 & 0.04877058 \\
0.001 & 0.00099950 \\
0.0001 & 0.00010000 \\
\end{longtable}

\subsubsection{Interpreting BLAST in Biological
Context}\label{interpreting-blast-in-biological-context}

Although BLAST's statistical framework is powerful, biological
interpretation of BLAST results always requires combining E values and
scores with domain knowledge about the sequences
involved\textsuperscript{10}. A ``true positive'' hit---one that is
homologous to the query and descended from a common ancestor---should
have a convincingly low E value, but additional criteria often matter:
similar sequence lengths, the presence of conserved motifs or
signatures, compatible domain structures, and congruent genomic context
all strengthen the inference of homology\textsuperscript{10}.
Conversely, BLAST results can include false positives with modest
sequence identity and acceptable E values that actually belong to
unrelated protein families, particularly when alignments are short or
concentrated in low-complexity segments\textsuperscript{10}.

For these reasons, it is common to follow up marginal BLAST matches with
reciprocal searches (using the putative hit as a new query), multiple
sequence alignment and phylogenetic analysis, structural comparisons, or
domain-based methods such as profile HMMs and position-specific scoring
matrices\textsuperscript{10}. In practice, BLAST serves as the first,
essential step: it rapidly maps a query into the landscape of known
sequences, providing a statistically grounded set of candidate
relationships that subsequent analyses can
refine\textsuperscript{10,63}.

\section{From pairwise to multiple sequence
alignment}\label{from-pairwise-to-multiple-sequence-alignment}

Multiple sequence alignment (MSA) extends pairwise alignment to three or
more sequences, allowing homologous residues to be aligned in columns
across a protein family or related DNA sequences\textsuperscript{10}. A
carefully constructed MSA encodes shared evolutionary history, conserved
structure, and functional motifs, enabling prediction of uncharacterized
protein functions, identification of catalytic residues, and
phylogenetic tree reconstruction\textsuperscript{10,56,65}. Compared to
pairwise alignment, MSA is more sensitive for detecting distant homologs
because information from many sequences is integrated rather than
relying on a single query--target comparison\textsuperscript{10,56}.

Conceptually, an MSA arranges homologous sequences so residues in the
same column are presumed homologous in both evolutionary and structural
senses\textsuperscript{10,65}. Residues aligned in a given column
descended from a shared ancestral residue and, in proteins, occupy
comparable positions in the three-dimensional
fold\textsuperscript{10,65}. In practice, achieving a uniquely
``correct'' alignment for divergent families is rarely possible because
sequence divergence often outpaces structural or functional data
availability; algorithms aim to approximate optimal alignments under
explicit criteria such as maximizing total alignment score or
sum-of-pairs objective functions\textsuperscript{10}.

\subsection{Criteria and Uses of MSAs}\label{criteria-and-uses-of-msas}

MSAs are built using different implicit criteria reflecting how residues
ought to align\textsuperscript{10}. \textbf{Structural similarity}:
residues occupying comparable three-dimensional positions (e.g., within
the same helix or beta strand) should align\textsuperscript{10,65}.
\textbf{Evolutionary similarity}: aligned residues should be
phylogenetically homologous, tracing back to a single ancestral
residue\textsuperscript{10}. \textbf{Functional similarity}: residues
with the same biochemical role (e.g., catalytic residues,
ligand-coordinating residues) should align even when surrounding regions
have diverged\textsuperscript{10}. \textbf{Sequence similarity}:
maximizing identical or conservatively substituted residues in each
column---the criterion directly optimized by most algorithms because it
is easily formalized and computed\textsuperscript{10}.

MSA algorithms largely implement the sequence-similarity criterion, but
success is judged by how well alignments respect structural and
functional constraints when that information is
available\textsuperscript{10}. For example, in the globin family, MSAs
reveal highly conserved histidine residues coordinating the heme group,
even though overall pairwise identities between some globins fall near
20--30\% identity\textsuperscript{10,65}.

MSAs inform protein structure and function
prediction\textsuperscript{10}. Profiles and hidden Markov models (e.g.,
DELTA-BLAST, HMMER) are derived from MSAs and exploit conservation
patterns to detect remote homologs\textsuperscript{10,56}. Conservation
patterns inform variant interpretation: amino-acid changes at highly
conserved positions across species are more likely deleterious, and many
variant effect predictors incorporate cross-species MSAs to quantify
evolutionary constraint\textsuperscript{10}. MSAs form the raw data for
phylogenetic inference methods; tree-building algorithms assume each
column corresponds to characters evolved from a common ancestral state,
so alignment inaccuracies propagate directly into tree
errors\textsuperscript{10}.

\begin{tcolorbox}[enhanced jigsaw, rightrule=.15mm, leftrule=.75mm, colbacktitle=quarto-callout-note-color!10!white, colframe=quarto-callout-note-color-frame, toprule=.15mm, opacitybacktitle=0.6, title=\textcolor{quarto-callout-note-color}{\faInfo}\hspace{0.5em}{Figure 5.3 -- Multiple Sequence Alignment of Insulin}, left=2mm, colback=white, bottomtitle=1mm, toptitle=1mm, bottomrule=.15mm, coltitle=black, titlerule=0mm, breakable, opacityback=0, arc=.35mm]

\begin{figure}[H]

\centering{

\pandocbounded{\includegraphics[keepaspectratio]{images/Figure_7.3_Pevsner.jpg}}

}

\caption{\label{fig-msa}Multiple Sequence Alignment (from Pevsner 7.3)}

\end{figure}%

Since the 1950s, studies of insulin have facilitated our understanding
of molecular evolution. (a) The human insulin molecule consists of a
signal peptide (required for intracellular transport; amino acid
residues 1--24), the B chain, the C peptide, and the A chain. Dibasic
residues (amino acids RR, KR) flank the C peptide and are the sites at
which proteases cleave the protein. The A chain and B chain are then
covalently linked through disulfide bridges, forming mature insulin. (b)
Multiple sequence alignment of insulin from 12 species. Amino acid
substitutions occur in nonrandom patterns. Note that within the A chain
of insulin the amino acid residues are almost perfectly conserved
between different species, except for three divergent columns of amino
acids (A chain, residues in green font). However, the rate of nucleotide
substitution is about six-fold higher in the region encoding the
intervening C peptide than in the region encoding the B and A chains
(Kimura, 1983), and gaps in the multiple sequence alignment are evident
here. Disulfide bridges between cysteine residues are indicated by
dashed lines. The accession numbers are NP\_000198.1 (human),
NP\_001008996.1 (chimpanzee), NP\_062003.1 (rat), NP\_001123565.1 (dog),
NP\_001172013.1 (mouse), NP\_001075804.1 (rabbit), NP\_001103242.1
(pig), NP\_990553.1 (chicken), NP\_001172055.1 (cow), P01318.2 (sheep),
XP\_003422420.1 (elephant), and P67974.1 (sperm whale). (c) Guinea pig
(Cavia porcellus, accession NP\_001166362.1) and coypu (Myocastor
coypus, P01330.1) insulins evolve about seven-fold faster than insulin
from other species. Human, mouse, guinea pig, and coypu insulins are
aligned. Arrows indicate 18 amino acid positions at which the guinea pig
sequence (purple) varies from that of human and/or mouse. Figure
source:\textsuperscript{10}.

\end{tcolorbox}

\subsection{Conserved Residues, Motifs, and Databases of
MSAs}\label{conserved-residues-motifs-and-databases-of-msas}

Highly conserved residues often mark structurally or functionally
critical positions such as catalytic residues, metal-binding residues,
or disulfide-forming cysteines\textsuperscript{10}. In the insulin
family, an MSA across mammals shows strongly conserved A and B chains,
particularly cysteines forming disulfide bonds, whereas the intervening
C peptide is far more variable with numerous indels\textsuperscript{10}.
This pattern reflects different functional constraints: the mature
hormone requires precise disulfide-bond geometry, whereas the C peptide
primarily functions as a processing intermediate and tolerates more
substitutions\textsuperscript{10}.

Databases such as Pfam and SMART curate large collections of MSAs
representing protein families and domains, each associated with profile
hidden Markov models for searching domain instances in new
sequences\textsuperscript{10,66,67}. In Pfam, each family has a manually
curated seed alignment capturing representative sequences and a larger
full alignment extending coverage to additional
homologs\textsuperscript{10,66}. Integrated resources such as InterPro
and NCBI Conserved Domain Database (CDD) combine multiple
alignment-based resources, allowing single query proteins to be
annotated against Pfam, PROSITE, SMART, TIGRFAMs, and other domain
collections\textsuperscript{10,68,69}.

\begin{tcolorbox}[enhanced jigsaw, rightrule=.15mm, leftrule=.75mm, colbacktitle=quarto-callout-note-color!10!white, colframe=quarto-callout-note-color-frame, toprule=.15mm, opacitybacktitle=0.6, title=\textcolor{quarto-callout-note-color}{\faInfo}\hspace{0.5em}{Connection to Chapter 4 and Lab 3}, left=2mm, colback=white, bottomtitle=1mm, toptitle=1mm, bottomrule=.15mm, coltitle=black, titlerule=0mm, breakable, opacityback=0, arc=.35mm]

In Chapter 4, you explored genome browsers including UCSC and Ensembl,
which display whole-genome MSAs through conservation tracks such as
PhastCons, phyloP, and GERP\textsuperscript{70--72}. These tracks derive
from multi-species alignments (e.g., Multiz) spanning entire chromosomes
and highlight conserved exons and regulatory
elements\textsuperscript{73,74}. In \href{Lab03.qmd}{Lab 3}, you will
explore these conservation tracks directly to identify functional
elements within specific genomic loci, connecting protein-level MSAs
covered in this chapter to genome-scale comparative genomics.

\end{tcolorbox}

\subsection{Algorithms for Multiple Sequence
Alignment}\label{algorithms-for-multiple-sequence-alignment}

Constructing accurate MSAs is algorithmically challenging because the
search space grows combinatorially with sequence number, and exact
optimization by dynamic programming becomes infeasible beyond a small
number of sequences\textsuperscript{10}. Practical MSA methods use
heuristic strategies trading exact optimality for computational
tractability while maintaining high accuracy\textsuperscript{10}. Five
main approaches exist: exact methods, progressive alignment, iterative
refinement, consistency-based approaches, and structure-based methods
(Table~\ref{tbl-msa-algorithms})\textsuperscript{10}.

\textbf{Exact methods} generalize Needleman--Wunsch to multiple
dimensions, maximizing sum-of-pairs or related objective
functions\textsuperscript{10}. Time and space requirements scale roughly
as (O(2\^{}N L\^{}N)) for (N) sequences of average length (L), making
them impractical for more than four or five
sequences\textsuperscript{10}. Exact approaches play mainly conceptual
or benchmarking roles\textsuperscript{10}.

\textbf{Progressive alignment} builds MSAs stepwise by computing all
pairwise alignments, constructing a guide tree based on pairwise
distances, and aligning sequences in tree-specified
order\textsuperscript{10,75,76}. CLUSTALW exemplifies this approach: it
uses global pairwise alignments (via Needleman--Wunsch) to compute
percent identities, converts similarities to distances, uses clustering
(UPGMA or neighbor-joining) to build a guide tree, and aligns the most
closely related pair first, progressively adding more sequences until
all are included\textsuperscript{10,77,78}. A key heuristic is ``once a
gap, always a gap,'' giving highest weight to initial closest-sequence
alignments when placing gaps\textsuperscript{10,77}. CLUSTALW
down-weights very similar sequence clusters so they do not dominate
scoring and uses position-specific gap penalties to encourage block-like
gap patterns\textsuperscript{10,77}. While historically the most popular
web-based progressive alignment program, experts now recommend newer
programs such as MAFFT, ProbCons, and MUSCLE for improved
performance\textsuperscript{10}. CLUSTALW remains valuable pedagogically
for illustrating progressive alignment and is implemented directly in
AliView, introduced in Chapter 4\textsuperscript{10,75--78}.

\textbf{Iterative methods} refine an initial MSA by repeatedly modifying
and re-aligning sequence subsets or profiles to improve an objective
score\textsuperscript{10,79,80}. MAFFT and MUSCLE illustrate this
strategy\textsuperscript{10,79,80}. MAFFT begins with fast progressive
alignment using k-mer counting (e.g., 6-tuples) to estimate pairwise
distances without full alignments, then refines by recalculating
distances from the current MSA and performing additional progressive
passes and iterative profile--profile realignments to maximize
sum-of-pairs scores\textsuperscript{10,80,81}. MAFFT's command-line
interface allows integration into reproducible pipelines and
high-performance computing environments, complementing web-based tools
for exploratory work\textsuperscript{10,80,81}. MUSCLE follows a
three-stage protocol: building a draft progressive alignment,
constructing an improved guide tree using corrected distance estimates
(e.g., Kimura distances), and performing iterative refinement by
partitioning the tree into subprofiles, realigning them, and accepting
changes increasing sum-of-pairs scores\textsuperscript{10,79,82}. In
both cases, iterative refinement can correct early errors from the first
progressive pass, overcoming the main limitation of strictly progressive
methods\textsuperscript{10,79,80}.

\textbf{Consistency-based approaches} introduce additional information
by using the emerging MSA to guide pairwise alignments respecting
transitive relationships among residues\textsuperscript{10,83,84}. If
residue \((x_i)\) in sequence \((x)\) aligns with residue \((z_k)\) in
sequence \((z)\), and \((z_k)\) aligns with \((y_j)\) in sequence
\((y)\), this supports aligning \((x_i)\) with
\((y_j)\)\textsuperscript{10,83,84}. ProbCons implements this via a
pair-HMM framework: computing posterior probability matrices for each
sequence pair, transforming these using consistency across all
sequences, building a guide tree based on expected alignment accuracy,
and performing progressive alignment with optional iterative
refinement\textsuperscript{10,84}. T-COFFEE builds an extended library
of alignments combining global and local pairwise alignments, reweights
residue pairs using consistency, and applies progressive alignment with
custom substitution matrices derived from the
library\textsuperscript{10,83}. Benchmarks show consistency-based
methods often yield more accurate alignments than purely progressive
methods, especially for moderately or highly divergent protein
families\textsuperscript{10,84}.

\textbf{Structure-based methods} incorporate three-dimensional
structural information, recognizing that tertiary structures are better
conserved than primary sequences\textsuperscript{10,85}. The Expresso
mode in T-COFFEE begins by identifying structural templates for input
sequences via BLAST searches against the Protein Data Bank, then uses
pairwise structural alignments to guide multiple alignment, improving
alignment of secondary structure elements and functionally important
residues\textsuperscript{10,85}. PRALINE can integrate profiles from
PSI-BLAST searches and predicted secondary structure, using this
additional information to refine MSAs beyond what sequence alone
dictates\textsuperscript{10}. Structural information is also crucial in
benchmarking: curated reference alignments based on superposed protein
structures serve as gold standards for evaluating algorithmic accuracy
using metrics such as sum-of-pairs and column
scores\textsuperscript{10,86,87}.

\begin{longtable}[]{@{}
  >{\raggedright\arraybackslash}p{(\linewidth - 4\tabcolsep) * \real{0.3824}}
  >{\raggedright\arraybackslash}p{(\linewidth - 4\tabcolsep) * \real{0.3235}}
  >{\raggedright\arraybackslash}p{(\linewidth - 4\tabcolsep) * \real{0.2941}}@{}}
\caption{Algorithm families with representative examples you will
encounter in lab exercises.}\label{tbl-msa-algorithms}\tabularnewline
\toprule\noalign{}
\begin{minipage}[b]{\linewidth}\raggedright
Method type
\end{minipage} & \begin{minipage}[b]{\linewidth}\raggedright
Core idea
\end{minipage} & \begin{minipage}[b]{\linewidth}\raggedright
Examples
\end{minipage} \\
\midrule\noalign{}
\endfirsthead
\toprule\noalign{}
\begin{minipage}[b]{\linewidth}\raggedright
Method type
\end{minipage} & \begin{minipage}[b]{\linewidth}\raggedright
Core idea
\end{minipage} & \begin{minipage}[b]{\linewidth}\raggedright
Examples
\end{minipage} \\
\midrule\noalign{}
\endhead
\bottomrule\noalign{}
\endlastfoot
Exact & Multi-dimensional dynamic programming; guaranteed optimal but
feasible only for very few short sequences. & Used mainly in
benchmarking studies. \\
Progressive & Compute all pairwise global alignments (often with
Needleman--Wunsch), build a guide tree from pairwise distances, then
align the most similar sequences first and add others stepwise; ``once a
gap, always a gap.'' & CLUSTALW (in AliView). \\
Iterative & Start with an initial MSA, then repeatedly modify and
re-score it, keeping changes that improve an objective until
convergence. & MAFFT, MUSCLE. \\
Consistency-based & Combine information from many pairwise local
alignments or profiles so that if A aligns to B and B to C, the
algorithm encourages A to align consistently with C. & T-COFFEE,
ProbCons. \\
Structure-based & Incorporate 3D structural information (when available)
alongside sequence when aligning, giving special weight to structurally
conserved cores. & Expresso (T-COFFEE module). \\
\end{longtable}

\begin{tcolorbox}[enhanced jigsaw, rightrule=.15mm, leftrule=.75mm, colbacktitle=quarto-callout-note-color!10!white, colframe=quarto-callout-note-color-frame, toprule=.15mm, opacitybacktitle=0.6, title=\textcolor{quarto-callout-note-color}{\faInfo}\hspace{0.5em}{NCBI COBALT Tool}, left=2mm, colback=white, bottomtitle=1mm, toptitle=1mm, bottomrule=.15mm, coltitle=black, titlerule=0mm, breakable, opacityback=0, arc=.35mm]

NCBI's COBALT (Constraint-Based Alignment Tool) provides a web-based MSA
tool integrating pairwise sequence similarity with additional
constraints such as conserved motifs and domain
information\textsuperscript{10}. COBALT uses RPS-BLAST searches against
the Conserved Domain Database to identify shared domains and motifs,
then incorporates this domain-level information into the alignment
process so residues within the same domain are preferentially
aligned\textsuperscript{10,69}. For classroom use, COBALT offers an
accessible way to perform MSA in a browser, visualize conserved regions,
and connect sequence-level conservation to domain annotations without
requiring local software installation\textsuperscript{10}.

\end{tcolorbox}

\subsection{Benchmarking and Databases of
MSAs}\label{benchmarking-and-databases-of-msas}

Because different MSA algorithms produce noticeably different
alignments, especially for divergent protein families, benchmarking is
essential for assessing accuracy and guiding tool
choice\textsuperscript{10,86}. Benchmarking datasets such as BAliBASE,
PREFAB, and SABmark consist of curated reference alignments, many
derived from three-dimensional structural superpositions rather than
sequence alone\textsuperscript{10,86}. Alignment accuracy is typically
quantified using sum-of-pairs scores (measuring how many residue pairs
aligned in the reference are also aligned in the test) and column scores
(measuring the fraction of correctly recovered alignment
columns)\textsuperscript{10,86,87}. These metrics have
limitations---column scores are very stringent, and sum-of-pairs scores
do not constitute true distance metrics---but provide standardized ways
to compare methods across test cases\textsuperscript{10,87}.

Benchmarking studies consistently show no single algorithm is best in
all situations: progressive methods such as CLUSTALW perform well on
closely related sequences and are easy to interpret, whereas iterative
and consistency-based methods such as MAFFT, MUSCLE, ProbCons, and
T-COFFEE generally achieve higher accuracy on more divergent families at
greater computational cost\textsuperscript{10,79,80,83,84}.
Structure-based methods like Expresso provide further gains when
reliable templates are available, particularly for alignments used in
homology modeling and structure prediction\textsuperscript{10,85}. For
practical work, a good strategy is often to generate alignments with at
least two different high-performing tools and compare them in regions of
biological interest, possibly making manual adjustments where structural
or functional information suggests better gap
placement\textsuperscript{10}.

Specialized MSA databases such as Pfam, SMART, and InterPro, together
with tools like CDD and DELTA-BLAST, provide a curated layer on top of
raw alignment algorithms by maintaining expert-reviewed MSAs and
associated models for thousands of protein
families\textsuperscript{10,66--69}. These resources are evaluated and
refined over time using benchmarking and structural data, so the MSAs
they expose can be trusted as high-quality starting points for
functional annotation and comparative analysis\textsuperscript{10,66}.

\subsection{Whole-Genome MSAs and Genome
Browsers}\label{whole-genome-msas-and-genome-browsers}

MSA concepts extend beyond individual proteins to entire genomic
regions, where the goal is to align orthologous loci across multiple
species to identify conserved exons, regulatory elements, and regions
under selection\textsuperscript{10,70,73,74}. Tools such as Multiz,
PhastCons, phyloP, and GERP are central to this domain, and their
results are visualized in genome browsers such as UCSC and Ensembl,
which you explored in Chapter 4\textsuperscript{10,70--73}.

In the UCSC Genome Browser, the Vertebrate Multiz Alignment and
Conservation track for the human beta-globin (HBB) locus shows
alignments of dozens of vertebrate genomes across exons and introns,
with PhastCons and phyloP scores highlighting regions of high
conservation or accelerated evolution\textsuperscript{10,70,71}. Exons
appear as highly conserved blocks with strong conservation scores,
whereas surrounding noncoding regions show a mixture of poorly conserved
segments and discrete peaks corresponding to putative regulatory
elements\textsuperscript{10}. GERP scores provide complementary
information by estimating rejected substitutions at each alignment
column, identifying constrained elements evolving more slowly than
expected under neutrality\textsuperscript{10,72}. Through the UCSC Table
Browser and tools such as Galaxy, these whole-genome MSAs can be
exported in MAF format, converted to FASTA, and further analyzed or
re-aligned with conventional MSA tools\textsuperscript{10,88}.

The Ensembl browser implements a related but distinct pipeline based on
Enredo, Pecan, and Ortheus (EPO) to generate phylogeny-aware MSAs and
reconstruct ancestral sequences at internal nodes of the species
tree\textsuperscript{10,74,89}. Enredo partitions genomes into colinear
segments accounting for rearrangements and duplications, Pecan uses a
consistency-based algorithm to align these segments across species, and
Ortheus infers ancestral sequences and insertion--deletion histories
using a phylogenetic model\textsuperscript{10,74}. For the HBB locus,
EPO alignments in Ensembl provide not only multi-species alignment of
extant sequences but also reconstructed ancestral sequences clarifying
whether particular gaps represent insertions or deletions and what
ancestral nucleotides likely occupied a given
position\textsuperscript{10,74}. These whole-genome MSAs complement
protein-level MSAs by revealing conservation patterns in exons, splice
sites, untranslated regions, and distant regulatory elements, connecting
naturally back to earlier chapters on genome
browsers\textsuperscript{10,74}.

\section{Problem Sets}\label{problem-sets-1}

\begin{enumerate}
\def\labelenumi{\arabic{enumi}.}
\item
  \textbf{Global Protein Alignment}: Implement the Needleman-Wunsch
  algorithm on this toy example. Below is the PAM matrix already set up
  based on a gap penalty of -6. First, select the optimal alignment
  between the pair of sequences. On a blank sheet of paper, work through
  the matrix to align the two sequences, remembering that vertical and
  horizontal moves insert gaps and reduce your overall score as in
  Figure~\ref{fig-pairwise}. Then use the scoring matrix to compute the
  overall score of your resulting optimal alignment. When done, write
  out the resulting pairwise sequence alignment of the two sequences as
  in Figure~\ref{fig-exSeq}.

  \begin{longtable}[]{@{}lllllll@{}}
  \caption{Scoring Matrix}\tabularnewline
  \toprule\noalign{}
  - & S & P & E & A & R & E \\
  \midrule\noalign{}
  \endfirsthead
  \toprule\noalign{}
  - & S & P & E & A & R & E \\
  \midrule\noalign{}
  \endhead
  \bottomrule\noalign{}
  \endlastfoot
  \textbf{S} & 2 & 1 & 0 & 1 & 0 & 0 \\
  \textbf{H} & -1 & 0 & 1 & -1 & 2 & 1 \\
  \textbf{A} & 1 & 1 & 0 & 2 & -2 & 0 \\
  \textbf{K} & 0 & -1 & 0 & -1 & 3 & 0 \\
  \textbf{E} & 0 & -1 & 4 & 0 & -1 & 4 \\
  \end{longtable}
\item
  \textbf{Multiple Sequence Alignment with COBALT}: Download the FASTA
  file \href{datasets/N.fa}{N.fa}. Go to the
  \href{https://www.ncbi.nlm.nih.gov/tools/cobalt/re_cobalt.cgi}{COBALT
  tool} at NCBI, upload the FASTA file, and click ``Align''.

  \begin{itemize}
  \tightlist
  \item
    How \textbf{similar} are these sequences? Would you say they are
    highly similar (likely from the same species) or dissimilar (likely
    from different species as in Figure~\ref{fig-msa})? These sequences
    are actually from the nucleocapsid phosphoprotein of the SARS-CoV-2
    genome.
  \item
    \textbf{Configure the view} by changing the color scheme to
    ``BLOSUM62''. How does this help? Configure the view again by
    changing to ``Frequency-based difference''. How does this change
    your interpretation? What do some of the other settings do?
  \item
    How \textbf{long} is this protein?
  \item
    Click the button at the top to generate a \textbf{phylogenetic
    tree}. Examine the tree topology and branch lengths.
  \end{itemize}
\item
  \textbf{Comparative MSA analysis}: Repeat the steps above with
  \textbf{this FASTA file: \href{datasets/S.fa}{S.fa}}. This second set
  of sequences is from the SARS-CoV-2 surface glycoprotein, also called
  the ``spike'' protein. Notice how the similarity of these sequences
  differs from the first set.

  \begin{itemize}
  \tightlist
  \item
    Again configure the view settings and explore the length.
  \item
    How does the length of this protein compare to the first?
  \item
    How do the two trees compare to each other? Look specifically at the
    scale bar at the bottom. Which sequence represents MORE genetic
    change?
  \item
    Look at the times these sequences were collected in the sequence
    names. Do the sequences cluster by year collected?
  \item
    Find the node associated with the cluster of sequences from
    2019/2020 and reroot the tree. Do these clusters in time make sense?
  \end{itemize}
\item
  \textbf{Algorithm comparison}: Download the FASTA file
  \href{datasets/hbb.fa}{hbb.fa} and go to the
  \href{https://www.ebi.ac.uk/jdispatcher/msa}{multiple alignment} page
  on the EBI platform. Select an MSA tool (e.g., CLUSTALW, MAFFT,
  MUSCLE) and upload the sequences to submit the job.

  \begin{itemize}
  \tightlist
  \item
    Like with COBALT, explore different color schemes for viewing the
    MSA.
  \item
    Look at the resulting phylogenetic tree. Work with others in the
    class to compare your results across different alignment tools.
  \item
    Compare the results to conservation patterns you would expect for
    hemoglobin based on its function. Do highly conserved regions
    correspond to functionally important domains?
  \item
    Try the same alignment with a second tool. Where do the two
    alignments differ? Which differences might matter for downstream
    phylogenetic or structural analyses?
  \end{itemize}
\item
  \textbf{Margaret Dayhoff tribute}: Margaret Dayhoff is credited with
  developing the single-letter abbreviation system for amino acids,
  similar to ASCII codes used for quality scores of genome sequences in
  FASTQ files covered in Chapter 3.

  \begin{itemize}
  \tightlist
  \item
    For fun, convert your name to a ``protein.'' If any letters in your
    name do not match amino acids, either remove them or choose your
    middle or last name (or a nickname---this is really just for fun!).
    \emph{For example, I translated my name to ``LAVRIE,'' replacing the
    U for the letter V.}
  \item
    Take your ``name protein'' and submit it to BLAST. Select the tool
    Protein BLAST or ``blastp''. For the Search set, leave it as the
    ``ClusteredNR'' which is the non-redundant set of protein sequences
    in NCBI.
  \item
    Examine your top hits. Do they make biological sense, or are they
    spurious matches driven by low complexity? What does this exercise
    reveal about the importance of query length and composition in
    database searches?
  \end{itemize}
\end{enumerate}

\section{Reflection Questions}\label{reflection-questions-4}

\begin{enumerate}
\def\labelenumi{\arabic{enumi}.}
\item
  \textbf{Substitution matrices and biological interpretation}: How does
  understanding where substitution matrices and gap penalties come from
  change the way you interpret BLAST or MSA output? Think about a gene
  or protein you care about---would you learn more from global or local
  alignments when comparing it across species, and why?
\item
  \textbf{Conservation patterns and functional inference}: In the MSA
  problem sets above, how does your knowledge of the N and S proteins
  from SARS-CoV-2 match the conservation patterns you observed in their
  alignments? Most PCR-based tests for SARS-CoV-2 detection target the N
  protein, with few targeting the S protein. Why do you think the N
  protein makes a better universal target for detection? In cases where
  the spike protein is used for detection, what information do you think
  it provides beyond positive or negative for infection?
\item
  \textbf{Algorithm persistence and modern genomics}: Many
  genome-alignment tools were developed decades after Needleman--Wunsch
  and Smith--Waterman but still rely on their core ideas. What do you
  think enabled these early algorithms to remain so influential in
  modern genomics, even as sequencing technologies and computational
  power have advanced dramatically?
\item
  \textbf{MSA tool choice and uncertainty}: When you see conflicting
  alignments from different tools (e.g., CLUSTALW vs.~MAFFT vs.~COBALT),
  how will you decide which to trust, and how much does that uncertainty
  matter for the biological inferences you intend to draw? What
  additional information (structural data, domain annotations,
  phylogenetic context) might help resolve ambiguities?
\item
  \textbf{Cross-chapter connections}: How do the conservation tracks you
  explored in UCSC and Ensembl (Chapter 4) relate to the MSA algorithms
  and benchmarking datasets discussed in this chapter? What are the
  advantages and limitations of using whole-genome MSAs displayed in
  browsers compared to protein-level MSAs you generate yourself with
  tools like MAFFT or COBALT?
\end{enumerate}

\section{References}\label{references-5}

\chapter{Genome Alignment and Genome
Assembly}\label{genome-alignment-and-genome-assembly}

\section{Genome alignment and genome
assembly}\label{genome-alignment-and-genome-assembly-1}

Genome alignment and genome assembly are now central to almost every
area of genome research, from variant discovery to comparative genomics
and conservation biology\textsuperscript{90--92}. This chapter connects
short‑read alignment algorithms such as BWA, modern long‑read mappers,
and telomere‑to‑telomere (T2T) genome assembly, with an emphasis on data
literacy and critical interpretation of algorithmic
choices\textsuperscript{90,92,93}. Throughout, the Bogan et~al.~(2025)
article cited in your Week~1 readings serves as a recurring example of
how experimental design, alignment strategy, and assembly quality
interact in real genome projects\textsuperscript{3}.

This chapter builds directly on earlier chapters, where sequencing
technologies, genome assembly assessment, and the ``three C's'' of
assemblies---continuity, completeness, and correctness---were
introduced. We also revisit Figure~\ref{fig-genomics-timeline}, where
the Human Genome Project, the 2004 euchromatic ``completion'' and the
2022 T2T human assembly were introduced\textsuperscript{92,94}. Here,
those concepts are revisited in the concrete context of read mapping and
assembly pipelines, giving you tools to critically evaluate alignments
and assemblies in papers such as Bogan et~al.~and in your own lab
work\textsuperscript{93,95}.

\section{Pre‑alignment quality control and read
trimming}\label{prealignment-quality-control-and-read-trimming}

Before any alignment, raw sequencing reads must be assessed for quality,
adapter contamination, and systematic errors that could bias downstream
analyses\textsuperscript{32,33}. Common pre‑processing steps include
per‑base quality profiling, removal of adapters and low‑quality trailing
bases, and filtering out reads that are too short or dominated by
ambiguous bases\textsuperscript{32,33}. In practice these steps are
assessed with FastQC and related tools introduced in Chapter 3. These
early decisions about filtering and trimming directly affect coverage,
error profiles, and the reliability of biological conclusions,
highlighting data literacy as an essential skill for modern genome
scientists\textsuperscript{93}. Though there are many tools for
trimming, including via Galaxy, the one we will use in this course is
\textbf{Trimmomatic}, which we will use on the ASC in
\href{Lab04.qmd}{Lab 4}.

The impact of trimming on downstream analyses has been systematically
evaluated for Illumina data, where over‑aggressive trimming can reduce
mapping rates and distort coverage, whereas moderate quality‑aware
trimming typically improves variant calling
accuracy\textsuperscript{96}. In practice, students should learn to
interpret quality score distributions, relate observed base‑quality
decay and adapter content to library preparation protocols, and justify
trimming parameters in the context of their experimental goals and
consider carefully the trade‑offs between depth, complexity, and read
quality\textsuperscript{93}.

\begin{tcolorbox}[enhanced jigsaw, rightrule=.15mm, leftrule=.75mm, colbacktitle=quarto-callout-tip-color!10!white, colframe=quarto-callout-tip-color-frame, toprule=.15mm, opacitybacktitle=0.6, title=\textcolor{quarto-callout-tip-color}{\faLightbulb}\hspace{0.5em}{Figure 6.1: Read QC and trimming}, left=2mm, colback=white, bottomtitle=1mm, toptitle=1mm, bottomrule=.15mm, coltitle=black, titlerule=0mm, breakable, opacityback=0, arc=.35mm]

\begin{figure}[H]

\centering{

\pandocbounded{\includegraphics[keepaspectratio]{images/Figure_6_1_trimming.png}}

}

\caption{\label{fig-trimming}Trimming}

\end{figure}%

Side‑by‑side FastQC plots showing per‑base sequence quality, adapter
content, and per base sequence content before and after trimming for a
sample dataset from my own research\textsuperscript{33}. Here you can
see the improvements after trimming. It is typical that after trimming,
FastQC report will result in an error for read length distribution
simply due to trimming, which can be ignored.

\end{tcolorbox}

\section{Foundations of short‑read genome
alignment}\label{foundations-of-shortread-genome-alignment}

Short‑read alignment is a specialized case of approximate string
matching, in which each read must be placed on a reference genome while
allowing for sequencing errors and true biological
differences\textsuperscript{90}. Because dynamic‑programming algorithms
such as Smith--Waterman scale with the product of read length \((m)\)
and genome size \((n)\), direct alignment of billions of reads to
gigabase genomes is computationally infeasible without algorithmic
shortcuts\textsuperscript{10,90}. Modern aligners therefore rely on a
combination of indexing and filtering strategies to rapidly identify
candidate loci, followed by a more precise local alignment step to
verify and score each candidate\textsuperscript{90}.

In this chapter, several computer‑science ideas appear repeatedly:
\textbf{hash tables} (fast lookup tables built on hashing functions),
\textbf{\((k)\)‑mers} (all possible substrings of length \((k)\) from a
sequence), and permutations of strings, including \textbf{reversible
permutations} such as the Burrows--Wheeler transform that let us
rearrange a sequence but still reconstruct the original
later\textsuperscript{10,90}.

Two key error models dominate short‑read mapping: the \textbf{Hamming
distance}, a metric for string comparison that counts sequence
mismatches only, and the \textbf{edit distance}, which accounts for
mismatches plus insertions and deletions\textsuperscript{10,90}. Many
practical aligners use weighted versions of edit distance, incorporating
position‑specific penalties and base‑quality scores so that low‑quality
mismatches are penalized less severely than confident
ones\textsuperscript{90}. Earlier, you saw how BLAST uses scoring
matrices and E‑values to distinguish biologically meaningful alignments
from hits expected by chance in a large database. Genome mappers, such
as
\textbf{\href{https://bio-bwa.sourceforge.net/bwa.shtml}{Burrows-Wheeler
Alignment (BWA) tool}}, use the same basic idea at a different scale:
each read--reference alignment is summarized by an alignment score and a
CIGAR string describing the pattern of matches, mismatches, insertions,
and deletions, and this score is then converted into a mapping quality
that reflects how likely it is that the read truly originated from that
genomic locus rather than from a similar sequence elsewhere. Thinking of
BWA as a ``genome‑scale BLAST with CIGAR strings'' provides an entry
point for understanding why the scoring scheme, gap penalties, and
thresholds strongly influence both the reported alignment path (the
CIGAR operations) and your confidence in each mapped
read\textsuperscript{10}.

\begin{tcolorbox}[enhanced jigsaw, rightrule=.15mm, leftrule=.75mm, colbacktitle=quarto-callout-note-color!10!white, colframe=quarto-callout-note-color-frame, toprule=.15mm, opacitybacktitle=0.6, title=\textcolor{quarto-callout-note-color}{\faInfo}\hspace{0.5em}{Figure 6.2 - Algorithm Strategies of Short Read Alignment Tools}, left=2mm, colback=white, bottomtitle=1mm, toptitle=1mm, bottomrule=.15mm, coltitle=black, titlerule=0mm, breakable, opacityback=0, arc=.35mm]

\begin{figure}[H]

\centering{

\includegraphics[width=0.75\linewidth,height=\textheight,keepaspectratio]{images/Fig_5_18_Pevsner.jpg}

}

\caption{\label{fig-alignment}Alignment Algorithms}

\end{figure}%

Two strategies by which short read aligners take a large number of short
reads (e.g., 500 million reads, each of length 150 base pairs) and align
them to a reference genome (e.g., the haploid human genome reference
sequence of ∼3 billion base pairs). (a) Spaced seed indexing algorithms
use hash tables. The reference genome and the short reads are cut into
equal-sized segments called seeds. Seeds from the short reads are paired
and stored in a lookup table (hash table) and used to scan the
reference. Matches to the seed index have an assigned genomic location.
(b) The Burrows--Wheeler transform is used to efficiently represent the
reference genome in software such as BWA2 and Bowtie. The reference
genome is concatenated into a string, transformed using the
Burrows--Wheeler transform (see Figure~\ref{fig-bwt}), and indexed.
Reads are aligned beginning with the base at the 3′ end and continuing
towards the 5′ end. This approach is very fast relative to spaced seed
approaches. This figure also serves as your first comparison of two
major families of short‑read aligners: hash‑table--based ``seed and
scan'' tools versus BWT/FM‑index--based tools like BWA and
Bowtie\textsuperscript{97}.

\end{tcolorbox}

\subsection{Burrows--Wheeler transform and
FM‑index}\label{burrowswheeler-transform-and-fmindex}

BWA is a canonical example of a Burrows--Wheeler transform (BWT)‑based
aligner that indexes the \textbf{reference genome} using an FM‑index,
and has become one of the most influential tools in bioinformatics, with
37,028 citations in Web of Science as of January
2026\textsuperscript{91}. At a high level, BWA replaces the naive ``scan
every base of the genome'' approach with an indexed representation that
makes it possible to search for a pattern of length \((k)\) in time
roughly proportional to \((k)\), largely independent of the genome
length\textsuperscript{90,97,98}.

Before we get to the BWT itself, it is helpful to expand two of the
computer‑science ideas introduced earlier that underlie many read
mappers: \textbf{hashing} and \textbf{reversible permutations}. A hash
table lets us map keys (for example all possible \((k)\)‑mers in the
genome) to values (their positions) in a way that supports very fast
lookup, and early mappers such as MAQ and SOAP built indices by hashing
short words from the reference or the reads\textsuperscript{90,97}. A
reversible permutation is a way of rearranging the characters of a
string (or rows of a matrix) so that we can later recover the original
order; the BWT is one such permutation of the genome sequence that
groups similar local contexts together and, crucially, can be inverted
to recover the original text\textsuperscript{90,98}.

The Burrows--Wheeler transform works by conceptually forming all cyclic
rotations of the genome (terminated with a special end‑of‑string
symbol), sorting those rotations lexicographically, and then taking the
last column of this sorted matrix as the transformed
string\textsuperscript{90}. Because rows with similar prefixes end up
adjacent after sorting, characters that share similar left‑hand contexts
in the genome tend to cluster in the BWT, which is why it can be
compressed effectively and searched efficiently\textsuperscript{10,90}.
The FM‑index augments this transformed string with two additional data
structures: a table of cumulative character counts and a way to compute,
for any position, how many times a particular character has appeared up
to that position (a ``rank'' operation). Together, these allow us to
perform \textbf{backward search}---matching a pattern from right to
left---by repeatedly narrowing a range of rows in the sorted matrix that
are consistent with the part of the pattern seen so
far\textsuperscript{90,97,98}.

\begin{tcolorbox}[enhanced jigsaw, rightrule=.15mm, leftrule=.75mm, colbacktitle=quarto-callout-note-color!10!white, colframe=quarto-callout-note-color-frame, toprule=.15mm, opacitybacktitle=0.6, title=\textcolor{quarto-callout-note-color}{\faInfo}\hspace{0.5em}{Lab connection: BWA and the yeast genome (Lab~4)}, left=2mm, colback=white, bottomtitle=1mm, toptitle=1mm, bottomrule=.15mm, coltitle=black, titlerule=0mm, breakable, opacityback=0, arc=.35mm]

In \textbf{\href{Lab04.qmd}{Lab 4}}, you will use BWA to build an
FM‑index of the \emph{Saccharomyces cerevisiae} reference genome and
then map Illumina reads to it. This lab gives you practical experience
with the concepts introduced here: you will see how indexing is
performed in practice, inspect the index files that BWA creates, run
alignments, and examine how CIGAR strings and mapping qualities reflect
the underlying backward‑search and LF‑mapping operations discussed in
this section.

\end{tcolorbox}

Concretely, to search for a pattern like ``GAT'' in a BWT‑indexed
genome, we start from the last character ``T'' and use the cumulative
counts and ranks to find all rows whose first character is T, giving us
an interval in the matrix that corresponds to all occurrences of ``T''
in the genome\textsuperscript{90}. We then prepend ``A'', updating the
interval so that it now corresponds to ``AT'', and then prepend ``G'',
updating again so that the final interval corresponds exactly to all
rows that begin with ``GAT''; the size of this interval tells us how
many matches exist and its boundaries can be mapped back to coordinates
in the original genome using LF‑mapping\textsuperscript{90,91}. From an
algorithm perspective, this is powerful because every extension step is
\(O(1)\), so the total time is proportional to the pattern length and
not the genome size, and because the index is compressed, we can keep
the entire structure in memory even for large
genomes\textsuperscript{90,97,98}.

\begin{tcolorbox}[enhanced jigsaw, rightrule=.15mm, leftrule=.75mm, colbacktitle=quarto-callout-note-color!10!white, colframe=quarto-callout-note-color-frame, toprule=.15mm, opacitybacktitle=0.6, title=\textcolor{quarto-callout-note-color}{\faInfo}\hspace{0.5em}{Figure 6.3 - BWT and backward search}, left=2mm, colback=white, bottomtitle=1mm, toptitle=1mm, bottomrule=.15mm, coltitle=black, titlerule=0mm, breakable, opacityback=0, arc=.35mm]

\begin{figure}[H]

\centering{

\includegraphics[width=0.75\linewidth,height=\textheight,keepaspectratio]{images/Fig_5_19_Pevsner.jpg}

}

\caption{\label{fig-bwt}Alignment Algorithms}

\end{figure}%

The Burrows--Wheeler Transform (BWT) shown for a short input string
\(S\) of length \(N = 6\). Panel~A lists all cyclic rotations of \(S\)
and their sorted order; Panel~B highlights the first (F) and last (L)
columns that define the BWT string; Panel~C indicates how this
representation can be compressed and later decompressed to recover the
original \(S\)\textsuperscript{10}. Use this figure when you work
through the BANANA\$ exercise in class to map each step in your
pencil‑and‑paper calculation onto rows, columns, and LF‑mapping arrows
in the diagram.

\end{tcolorbox}

LF‑mapping (last‑to‑first mapping) is the concrete mechanism that makes
the BWT a reversible permutation of the genome: it uses the fact that
the first column of the sorted matrix is just the characters of the
genome in sorted order, and that each occurrence of a character in the
last column corresponds to the same relative occurrence in the first
column\textsuperscript{90}. In practice, this means that once we know
which rows correspond to our pattern, we can follow LF‑mapping steps to
walk backwards along each row and reconstruct the original genomic
coordinates, which is how BWA turns a range in the FM‑index into a
concrete alignment position\textsuperscript{91}. When mismatches or
small gaps are allowed, BWA explores nearby paths in this index by
controlled backtracking---trying alternative characters at a limited
number of positions---rather than by rescanning the entire reference,
which explains its efficiency compared with naive dynamic programming
over the whole genome\textsuperscript{90,91}.

\begin{longtable}[]{@{}
  >{\raggedright\arraybackslash}p{(\linewidth - 12\tabcolsep) * \real{0.1667}}
  >{\raggedright\arraybackslash}p{(\linewidth - 12\tabcolsep) * \real{0.2667}}
  >{\centering\arraybackslash}p{(\linewidth - 12\tabcolsep) * \real{0.1167}}
  >{\centering\arraybackslash}p{(\linewidth - 12\tabcolsep) * \real{0.1333}}
  >{\centering\arraybackslash}p{(\linewidth - 12\tabcolsep) * \real{0.1333}}
  >{\centering\arraybackslash}p{(\linewidth - 12\tabcolsep) * \real{0.1000}}
  >{\centering\arraybackslash}p{(\linewidth - 12\tabcolsep) * \real{0.0833}}@{}}
\caption{Classes of short‑read mappers Adapted from Li \& Homer (2010),
Briefings in Bioinformatics 11: 473--483\textsuperscript{97}. Include
columns for representative tools, indexing strategy (hash‑based
vs.~BWT/FM‑index), seeding and error model, typical read lengths, and
primary use cases. BWA was highlighted in lecture along with with the
BWT/FM‑index (see Figure~\ref{fig-alignment}). \textsuperscript{a}Work
well for Sanger and 454 reads, allowing gaps and clipping;
\textsuperscript{b}Paired end mapping.\\
\textsuperscript{c}Make use of base quality in alignment;
\textsuperscript{d}BWA trims the primer base and the first color for a
color read; \textsuperscript{e}Long-read alignment implemented in the
BWA-SW module; \textsuperscript{f}MAQ only does gapped alignment for
Illumina paired-end reads; \textsuperscript{g}Free executable for
non-profit projects only.
}\label{tbl-short-read-aligners}\tabularnewline
\toprule\noalign{}
\begin{minipage}[b]{\linewidth}\raggedright
Program
\end{minipage} & \begin{minipage}[b]{\linewidth}\raggedright
Algorithm
\end{minipage} & \begin{minipage}[b]{\linewidth}\centering
SOLiD
\end{minipage} & \begin{minipage}[b]{\linewidth}\centering
Long\textsuperscript{a}
\end{minipage} & \begin{minipage}[b]{\linewidth}\centering
Gapped
\end{minipage} & \begin{minipage}[b]{\linewidth}\centering
PE\textsuperscript{b}
\end{minipage} & \begin{minipage}[b]{\linewidth}\centering
Q\textsuperscript{c}
\end{minipage} \\
\midrule\noalign{}
\endfirsthead
\toprule\noalign{}
\begin{minipage}[b]{\linewidth}\raggedright
Program
\end{minipage} & \begin{minipage}[b]{\linewidth}\raggedright
Algorithm
\end{minipage} & \begin{minipage}[b]{\linewidth}\centering
SOLiD
\end{minipage} & \begin{minipage}[b]{\linewidth}\centering
Long\textsuperscript{a}
\end{minipage} & \begin{minipage}[b]{\linewidth}\centering
Gapped
\end{minipage} & \begin{minipage}[b]{\linewidth}\centering
PE\textsuperscript{b}
\end{minipage} & \begin{minipage}[b]{\linewidth}\centering
Q\textsuperscript{c}
\end{minipage} \\
\midrule\noalign{}
\endhead
\bottomrule\noalign{}
\endlastfoot
Bfast & hashing ref. & Yes & No & Yes & Yes & No \\
Bowtie & FM-index & Yes & No & No & Yes & No \\
\textbf{BWA} & \textbf{FM-index} & \textbf{Yes}\textsuperscript{d} &
\textbf{Yes}\textsuperscript{e} & \textbf{Yes} & \textbf{Yes} & No \\
MAQ & hashing reads & Yes & No & Yes\textsuperscript{f} & Yes & Yes \\
Mosaik & hashing ref. & Yes & Yes & Yes & Yes & No \\
Novoalign\textsuperscript{g} & hashing ref. & No & No & Yes & Yes &
Yes \\
\end{longtable}

\subsection{Short‑read aligners in practice: BWA‑aln, BWA‑mem, and
ancient
DNA}\label{shortread-aligners-in-practice-bwaaln-bwamem-and-ancient-dna}

Even though BWA was introduced in 2009, its algorithms remain central to
current human genomics and ancient DNA research. BWA implements three
major algorithms: \textbf{bwa‑aln}, an older backtracking aligner
optimized for short reads; \textbf{bwa‑sw}, an early
Smith--Waterman‑style local aligner; and \textbf{bwa‑mem}, a more recent
algorithm designed for longer reads that uses maximal exact matches
(MEMs) as seeds\textsuperscript{91,97}. Bwa‑aln does not rely on long
exact seeds; instead, it systematically explores the space of
approximate matches with bounded numbers of mismatches, guaranteeing
that all alignments within a specified edit distance are found, at the
cost of increased computational effort for long
reads\textsuperscript{90,91}. In contrast, bwa‑mem heuristically
searches for MEMs of length at least 19~bp by default, which works well
for modern 150~bp Illumina reads but can fail for highly degraded
ancient DNA (aDNA) reads that harbor many mismatches and
deamination‑induced substitutions\textsuperscript{99}. This newest
algorithm is the standard recommended aligner for short reads through
the genome analyis toolkit (GATK) introduced in Chapter 4 of this
textbook. You will get a chance to use \textbf{bwa‑mem} in
\href{Lab04.qmd}{Lab 4} to map short reads from a yeast experimental
evolution experiment to the yeast genome.

For example, in the case of a 32~bp read with two widely separated
mismatches, the longest exact match may be only 11~bp, shorter than
bwa‑mem's default minimum seed length, causing the true alignment to be
missed entirely even when it would be acceptable under the overall
mismatch threshold\textsuperscript{99}. Reducing the seed length would
in principle restore sensitivity, but an 11‑mer occurs many times in the
human genome, so seeding on all such matches would render bwa‑mem
impractically slow\textsuperscript{99}. \textbf{Bwa‑aln} therefore
remains widely used in modern workflows using ancient DNA (aDNA) because
its non‑exact seeding strategy maintains sensitivity for short,
error‑prone reads without requiring extremely short exact seeds,
illustrating how algorithmic assumptions and parameter defaults embed
implicit priors about read length and error
profiles\textsuperscript{91,99}.

Mapping quality (MAPQ) estimation further complicates comparisons
between aligners. Bwa‑aln assigns a MAPQ~≥~25 when no additional
alignment with one extra mismatch is observed, effectively encoding a
heuristic notion of uniqueness that many downstream pipelines adopt as a
hard filter\textsuperscript{99,100}. This leads to a close coupling
between bwa‑aln's internal heuristics and community practices such as
``MAPQ~25'' cut‑offs, which may inadvertently penalize alternative
aligners whose MAPQ models respond differently to the number of
mismatches\textsuperscript{90,99}.

\section{Long‑read sequencing and the evolution of alignment
algorithms}\label{longread-sequencing-and-the-evolution-of-alignment-algorithms}

In Chapter~2 you saw how third‑generation sequencing technologies such
as PacBio and Oxford Nanopore Technologies (ONT) extend read lengths
from a few hundred bases to many kilobases, at the cost of substantially
higher error rates in their early
implementations\textsuperscript{25--27}. Those same properties that make
long reads powerful for resolving repeats and structural variants also
force us to rethink alignment algorithms: a mapper designed around
short, mostly accurate Illumina reads will struggle when almost every
base has a non‑trivial chance of being wrong\textsuperscript{90}. As
long‑read chemistries have improved---from noisy continuous long reads
to highly accurate PacBio HiFi and ONT duplex reads---alignment
algorithms have evolved in parallel to exploit the new balance between
length and accuracy\textsuperscript{28,92}.

Conceptually, long‑read mapping started from the same
\textbf{seed‑and‑extend} idea you encountered for short reads: identify
promising seed matches between read and reference, then extend them with
dynamic programming\textsuperscript{90}. However, for long reads almost
every read contains many thousands of potential seeds, so the bottleneck
shifts from ``can we find any seed?'' to ``can we select and connect a
small number of informative seeds into a coherent alignment
path?''\textsuperscript{101}. Modern long‑read aligners therefore adopt
a \textbf{seed‑and‑chain} paradigm: they choose a sparse subset of seeds
(for example minimizers or \((k)\)‑min‑mers) and then use fast chaining
algorithms in this reduced ``minimizer space'' to build an approximate
alignment path before running any expensive base‑level
alignment\textsuperscript{101}.

\textbf{Minimap2} is a prototypical example and the recommended aligner
of choice for GATK with long read datasets. It computes minimizers from
both reads and reference, discards most of them to keep only a sparse
sketch, and then uses dynamic programming in two dimensions (position in
read vs.~position in reference) to chain collinear seeds into long
candidate alignments, which are then refined with banded alignment
around the chain\textsuperscript{101}. Other tools, such as Winnowmap2,
modify the seeding step to down‑weight repetitive \((k)\)‑mers so that
the index is dominated by more informative seeds, which reduces spurious
matches in repeat‑rich regions\textsuperscript{101}. Other aligners,
such as NGMLR and LRA, adjust their scoring schemes and chaining
heuristics to be particularly sensitive to long insertions, deletions,
and complex rearrangements, which is important when long reads are used
primarily for structural variant discovery rather than single‑nucleotide
variant calling\textsuperscript{101,102}.

\begin{longtable}[]{@{}
  >{\raggedright\arraybackslash}p{(\linewidth - 6\tabcolsep) * \real{0.0549}}
  >{\raggedright\arraybackslash}p{(\linewidth - 6\tabcolsep) * \real{0.1832}}
  >{\raggedright\arraybackslash}p{(\linewidth - 6\tabcolsep) * \real{0.3516}}
  >{\raggedright\arraybackslash}p{(\linewidth - 6\tabcolsep) * \real{0.4103}}@{}}
\caption{Long‑read alignment software Representative long‑read aligners
and their typical use cases, based on recent surveys and benchmarking
studies\textsuperscript{101,102}.}\label{tbl-long-read-aligners}\tabularnewline
\toprule\noalign{}
\begin{minipage}[b]{\linewidth}\raggedright
Software tool
\end{minipage} & \begin{minipage}[b]{\linewidth}\raggedright
Description
\end{minipage} & \begin{minipage}[b]{\linewidth}\raggedright
Key features
\end{minipage} & \begin{minipage}[b]{\linewidth}\raggedright
Benchmarking insights
\end{minipage} \\
\midrule\noalign{}
\endfirsthead
\toprule\noalign{}
\begin{minipage}[b]{\linewidth}\raggedright
Software tool
\end{minipage} & \begin{minipage}[b]{\linewidth}\raggedright
Description
\end{minipage} & \begin{minipage}[b]{\linewidth}\raggedright
Key features
\end{minipage} & \begin{minipage}[b]{\linewidth}\raggedright
Benchmarking insights
\end{minipage} \\
\midrule\noalign{}
\endhead
\bottomrule\noalign{}
\endlastfoot
Minimap2 & Versatile aligner for long nucleotide sequences. & Efficient
alignment of PacBio and ONT reads; supports spliced alignment. &
Computationally lightweight; recommended for large‑scale use and as a
default mapper in many pipelines. \\
GraphMap2 & Aligner optimized for long, error‑prone reads. & High
sensitivity; suitable for complex genomes. & High computational burden;
not ideal for whole human‑genome exploration. \\
LRA & Designed for accurate alignment of long reads. & Balances speed
and accuracy; handles high‑error‑rate reads and complex structural
variants. & Effective with PacBio data; offers strong structural‑variant
performance in human genomics. \\
NGMLR & Tailored for structural variant detection. & Efficiently aligns
reads with large insertions or deletions and complex breakpoints. &
Resource‑intensive and time‑consuming; produces reliable SV alignments
and is useful as a supplementary tool. \\
Winnowmap2 & Aligner for long reads with repetitive sequences. &
Utilizes minimizer‑based seeding with repeat‑aware weighting; reduces
spurious alignments. & Computationally efficient; complements Minimap2
by offering different insights in repeat‑rich regions. \\
\end{longtable}

From an algorithmic standpoint, you can view these tools as different
points along a design space defined by three main choices: (1) how seeds
are selected (uniform minimizers versus repeat‑aware seeds), (2) how
chains are scored and constructed (emphasis on collinearity
vs.~tolerance for large jumps), and (3) how expensive base‑level
alignment is used to refine candidate chains\textsuperscript{101}. For
example, a tool tuned for transcriptomics might favor seeds and chains
that support spliced alignments across introns, while a tool optimized
for structural variants may accept non‑collinear seed chains that jump
across complex breakpoints\textsuperscript{101,102}. These design
decisions echo the trade‑offs between sensitivity and speed that you saw
earlier for BLAST and BWA: more seeds and more permissive chains usually
increase sensitivity, but they also increase runtime and memory
usage\textsuperscript{10,90}.

\begin{tcolorbox}[enhanced jigsaw, rightrule=.15mm, leftrule=.75mm, colbacktitle=quarto-callout-tip-color!10!white, colframe=quarto-callout-tip-color-frame, toprule=.15mm, opacitybacktitle=0.6, title=\textcolor{quarto-callout-tip-color}{\faLightbulb}\hspace{0.5em}{Benchmarking long‑read mappers}, left=2mm, colback=white, bottomtitle=1mm, toptitle=1mm, bottomrule=.15mm, coltitle=black, titlerule=0mm, breakable, opacityback=0, arc=.35mm]

Benchmarking is the process of systematically evaluating software tools
on curated datasets using agreed‑upon metrics, so that performance
claims are supported by evidence rather than anecdotes. In the context
of long‑read alignment, typical benchmarking metrics include alignment
recall and precision (how many true read placements are recovered and
how many mappings are correct), structural‑variant detection accuracy,
runtime, and memory usage\textsuperscript{102}. Datasets may be fully
synthetic (where the true answer is known), based on simulated reads
from real genomes, or high‑coverage ``gold standard'' genomes such as
the Genome in a Bottle resources.

In their 2023 study, LoTempio and colleagues benchmarked several
long‑read genome sequence alignment tools for human genomics
applications, including Minimap2, GraphMap2, LRA, NGMLR, and
Winnowmap2\textsuperscript{102}. They found substantial trade‑offs
between computational efficiency, alignment success rates, and the
ability to detect structural variants: no single aligner performed best
on all metrics, and different tools excelled on different variant
classes and coverage regimes\textsuperscript{102}. Their results support
a \textbf{combination strategy} rather than a single‑tool mindset: in
many human‑genome projects it is reasonable to use Minimap2 as a fast,
baseline mapper, add Winnowmap2 to improve performance in repeat‑rich
regions, and optionally incorporate NGMLR or LRA when highly accurate
structural‑variant discovery is a central goal\textsuperscript{102}.

In \textbf{\href{Lab08.qmd}{Lab~8}}, you will take on a similar role for
variant calling tools. Different student groups will read and discuss
papers that evaluate a variety of variant‑calling tools on reference
datasets such as Genome in a Bottle, focusing on how study design,
metrics, and datasets shape the conclusions. In class, we will
synthesize these results to formulate tool recommendations for different
project goals, mirroring the way LoTempio et~al.~used systematic
benchmarking to motivate their alignment recommendations.

\end{tcolorbox}

For you as a practitioner, the key message is that long‑read mappers are
not interchangeable black boxes. When you design a pipeline---whether
for a T2T‑style assembly evaluation, a structural‑variant survey, or an
ecological genomics project like Bogan et~al.---you should explicitly
match your choice of aligner (and its preset or parameter set) to the
sequencing technology, read accuracy, and biological
question\textsuperscript{3,92}. Just as you learned to justify trimming
thresholds and short‑read aligner settings earlier in the chapter, you
should be prepared to justify why a particular long‑read mapper, or
combination of mappers, is appropriate for your data.

\section{From alignment to assembly: concepts and
algorithms}\label{from-alignment-to-assembly-concepts-and-algorithms}

Alignment assumes that a reasonably complete reference genome already
exists; assembly tackles the inverse problem of reconstructing that
genome from a noisy collection of reads\textsuperscript{90,92}. In
\textbf{de novo assembly}, reads are connected to each other using
overlaps or shared \((k)\)‑mers, forming graphs whose paths represent
candidate genome reconstructions\textsuperscript{10,90,92}. This section
develops the main graph formulations, links them to core
computer‑science ideas, and prepares you for the long‑read and
pan‑genome methods that follow.

\subsection{Key computer‑science ideas for
assembly}\label{key-computerscience-ideas-for-assembly}

Assembly algorithms borrow heavily from graph theory. A \textbf{graph}
is a collection of vertices (nodes) connected by edges (links); in
assembly, vertices represent reads or \((k)\)‑mers and edges encode
evidence that two sequences are adjacent in the
genome\textsuperscript{90,92}. A \textbf{directed graph} is one in which
each edge has a direction, such as A → B; assembly graphs are directed
because DNA has a defined 5′→3′ orientation, and edges respect that
orientation\textsuperscript{90,92}. A \textbf{path} is an ordered
sequence of edges that visits a series of vertices, and an assembled
contig corresponds to a path through the assembly graph that spells out
a candidate genomic sequence\textsuperscript{90}.

The term \textbf{contig} refers to a maximal contiguous stretch of
sequence derived from the graph that contains no unresolved branches; in
many assemblers this corresponds to a maximal non‑branching path and is
sometimes called a ``unitig''\textsuperscript{92}. For the purposes of
this book, it is enough to view contigs as the basic building blocks
that assemblers output before higher‑level scaffolding connects them
into chromosome‑scale sequences. A related idea is that of a
\textbf{transitive edge}: an edge A → C is transitive if the path A → B
→ C already exists. Removing such edges (a process called
\textbf{transitive reduction}) simplifies overlap graphs without losing
information and makes later path‑finding more
efficient\textsuperscript{90,92}.

You also saw \((k)\)‑mers in earlier chapters: all substrings of length
\((k)\) in a sequence. In assembly, \((k)\)‑mers play a dual role as
basic units of overlap between reads and as vertices in de~Bruijn
graphs\textsuperscript{10,92}. Choosing \((k)\) therefore affects both
sensitivity to low‑coverage regions and the ability to distinguish
repeats, in much the same way that word size influences BLAST
sensitivity.

\subsection{Classical graph
formulations}\label{classical-graph-formulations}

Two classical graph paradigms dominate assembly algorithms.

In an \textbf{overlap graph} (and its compressed form, the string
graph), each read is a vertex and a directed edge connects reads whose
suffix and prefix overlap by at least some minimum
length\textsuperscript{90,92}. After transitive edges are removed and
reads that are fully contained within longer reads are discarded, long
non‑branching paths are merged into contigs that represent contiguous
genomic segments\textsuperscript{92}. This overlap--layout--consensus
approach matches the original Sanger‑era view of assembly: identify
overlaps, lay out reads in order, then compute a consensus sequence.

In a \textbf{de~Bruijn graph (DBG)}, the basic unit is not a whole read
but a \((k)\)‑mer. Vertices represent all \((k)\)‑mers observed in the
reads, and a directed edge connects two vertices when their \((k)\)‑mers
occur consecutively in a read and therefore share \((k-1)\)
bases\textsuperscript{10,92}. A path through the DBG corresponds to a
sequence of \((k)\)‑mers whose overlaps reconstruct an assembled
sequence, and contigs arise from maximal non‑branching paths in this
graph. Repeats appear as cycles or branching structures, and the choice
of \((k)\) determines how easily the graph distinguishes true genomic
branches from noise: small \((k)\) tends to tangle the graph in repeats
but maintains connectivity in low‑coverage regions, whereas large
\((k)\) resolves many repeats at the cost of fragmenting low‑coverage
areas\textsuperscript{10,92}.

\begin{tcolorbox}[enhanced jigsaw, rightrule=.15mm, leftrule=.75mm, colbacktitle=quarto-callout-note-color!10!white, colframe=quarto-callout-note-color-frame, toprule=.15mm, opacitybacktitle=0.6, title=\textcolor{quarto-callout-note-color}{\faInfo}\hspace{0.5em}{Figure 6.4 -- Genome assembly and repeats}, left=2mm, colback=white, bottomtitle=1mm, toptitle=1mm, bottomrule=.15mm, coltitle=black, titlerule=0mm, breakable, opacityback=0, arc=.35mm]

\begin{figure}[H]

\centering{

\includegraphics[width=0.5\linewidth,height=\textheight,keepaspectratio]{images/Figure_9.10_Pevsner.jpg}

}

\caption{\label{fig-assembly-repeats}Genome Assembly}

\end{figure}%

Efficient assembly of repetitive DNA using a de~Bruijn graph. Adapted
from Pevzner et~al.~as reproduced in Pevsner (Figure~9.10). Panels
illustrate (A) a genomic region with multiple copies of a repeat, (B) an
overlap/layout graph with many possible paths through the repeat, (C)
the conceptual ``gluing'' of repeated segments, and (D) the resulting
de~Bruijn graph, in which each repeat is represented as an edge rather
than a large tangle of nodes. This figure reinforces how graph design
influences our ability to resolve repeats in
assembly\textsuperscript{10}.

\end{tcolorbox}

Early Sanger‑based whole‑genome assemblies and many modern long‑read
assemblers operate directly on overlap or string graphs, because read
lengths are long enough that explicit read--read overlaps are highly
informative\textsuperscript{90,92}. Short‑read assemblers typically
favor de~Bruijn graphs, which avoid an all‑against‑all overlap search by
working in \((k)\)‑mer space and are therefore more computationally
efficient for billions of short reads\textsuperscript{10}. Conceptually,
both paradigms encode the same idea---local sequence adjacency---but
they distribute the computational work differently and respond
differently to repeats and sequencing errors\textsuperscript{90,92}.

Repeats longer than the read length remain a central difficulty in both
frameworks\textsuperscript{90,92}. If no read (or pair of reads) bridges
a repeat, the graph contains ambiguous branches that cannot be resolved
unambiguously, leading to breaks in contigs or potential mis‑joins. This
limitation motivated both the move to longer reads and the development
of auxiliary long‑range data types that you will see in the
telomere‑to‑telomere section below.

\subsection{Assembly metrics and the three
Cs}\label{assembly-metrics-and-the-three-cs}

Assembly quality is summarized using the ``three Cs'' introduced in
Chapter~2: continuity, completeness, and correctness.
\textbf{Continuity} is often described by contig and scaffold N50
values, defined as the length \((L)\) such that contigs or scaffolds of
length at least \((L)\) together cover 50~\% of the
assembly\textsuperscript{10,92}. \textbf{Completeness} is assessed by
total assembly size, recovery of known genes, and gene‑set metrics such
as BUSCO or compleasm, which count how many conserved single‑copy
orthologs are present, fragmented, or duplicated\textsuperscript{29,92}.
\textbf{Correctness} is evaluated using read‑back alignment, k‑mer
spectrum analysis (for example, KAT, Merqury, yak), and checks for
structural mis‑assemblies or false duplications relative to high‑quality
references or genetic maps\textsuperscript{92}.

Historically, assembler benchmarks have shown that different tools
applied to the same short‑read dataset can yield assemblies with very
different N50, error rates, and gene recovery, even when all nominally
target the same genome\textsuperscript{10}. Some assemblies achieve
impressive continuity at the cost of increased mis‑joins; others are
conservative but fragmented. As with alignment, these differences
reflect underlying algorithmic choices and parameter settings,
reinforcing the need to interpret assembly metrics in the context of the
sequencing strategy and biological question rather than as absolute
scores\textsuperscript{10,92}.

\begin{longtable}[]{@{}
  >{\raggedright\arraybackslash}p{(\linewidth - 10\tabcolsep) * \real{0.1725}}
  >{\raggedright\arraybackslash}p{(\linewidth - 10\tabcolsep) * \real{0.1338}}
  >{\raggedright\arraybackslash}p{(\linewidth - 10\tabcolsep) * \real{0.1056}}
  >{\raggedright\arraybackslash}p{(\linewidth - 10\tabcolsep) * \real{0.1127}}
  >{\raggedright\arraybackslash}p{(\linewidth - 10\tabcolsep) * \real{0.1514}}
  >{\raggedright\arraybackslash}p{(\linewidth - 10\tabcolsep) * \real{0.3239}}@{}}
\caption{Genome assembly eras and
strategies}\label{tbl-assembly-strategies}\tabularnewline
\toprule\noalign{}
\begin{minipage}[b]{\linewidth}\raggedright
Era / exemplar assembly
\end{minipage} & \begin{minipage}[b]{\linewidth}\raggedright
Dominant sequencing technology
\end{minipage} & \begin{minipage}[b]{\linewidth}\raggedright
Typical read length
\end{minipage} & \begin{minipage}[b]{\linewidth}\raggedright
Main assembly paradigm
\end{minipage} & \begin{minipage}[b]{\linewidth}\raggedright
Typical continuity (contig N50)
\end{minipage} & \begin{minipage}[b]{\linewidth}\raggedright
Major limitations and notes
\end{minipage} \\
\midrule\noalign{}
\endfirsthead
\toprule\noalign{}
\begin{minipage}[b]{\linewidth}\raggedright
Era / exemplar assembly
\end{minipage} & \begin{minipage}[b]{\linewidth}\raggedright
Dominant sequencing technology
\end{minipage} & \begin{minipage}[b]{\linewidth}\raggedright
Typical read length
\end{minipage} & \begin{minipage}[b]{\linewidth}\raggedright
Main assembly paradigm
\end{minipage} & \begin{minipage}[b]{\linewidth}\raggedright
Typical continuity (contig N50)
\end{minipage} & \begin{minipage}[b]{\linewidth}\raggedright
Major limitations and notes
\end{minipage} \\
\midrule\noalign{}
\endhead
\bottomrule\noalign{}
\endlastfoot
Early clone‑based assemblies (e.g.~\emph{C. elegans}, early human) &
Sanger, hierarchical BAC tiling\textsuperscript{10} & 500--1000~bp &
Overlap--layout--consensus (OLC) & Hundreds of kb to a few
Mb\textsuperscript{10} & Labor‑intensive clone maps; gaps in repeats and
heterochromatin; extensive manual curation. \\
Short‑read WGS human and model genomes (e.g.~GRCh37/38) & Illumina short
reads\textsuperscript{92} & 75--150~bp & Primarily de~Bruijn graphs &
Tens to low hundreds of kb\textsuperscript{92} & Fragmented in long
repeats; many unresolved segmental duplications and centromeric gaps. \\
Early long‑read hybrid assemblies & Mix of Illumina and noisy PacBio CLR
/ ONT\textsuperscript{92} & Illumina 150~bp; long reads 10--30~kb
(∼10~\% error) & OLC or DBG plus hybrid polishing & Hundreds of kb to a
few Mb\textsuperscript{92} & Better contiguity but frequent collapsed
repeats and consensus errors in low‑accuracy long reads. \\
HiFi‑based diploid assemblies (no ultra‑long) & PacBio HiFi; optional
ONT duplex\textsuperscript{92} & 10--20~kb, ≤0.5~\% error &
HiFi‑optimized overlap or multiplex DBG & Multi‑Mb contigs, often
chromosome arms\textsuperscript{92} & Many repeats and long haplotype
blocks remain partially unresolved; scaffolding needed for full
chromosomes. \\
Near T2T assemblies (e.g.~T2T‑CHM13 human genome) & PacBio HiFi + ONT
ultra‑long + Hi‑C / Strand‑seq / trio\textsuperscript{92,103} & HiFi
10--20~kb; ONT ≥100~kb & Graph‑based assembly with ultra‑long
integration and long‑range phasing & Chromosome‑scale, often
telomere‑to‑telomere\textsuperscript{92,103} & Remaining challenges in
very long tandem arrays and highly repetitive regions; phasing complex
polyploids. \\
\end{longtable}

\section{Assembly in the telomere‑to‑telomere
era}\label{assembly-in-the-telomeretotelomere-era}

The advent of single‑molecule long‑read sequencing around 2010, and
especially the arrival of accurate PacBio HiFi and ONT duplex reads
around 2019, transformed assembly from fragmented drafts into near
telomere‑to‑telomere reconstructions for many
species\textsuperscript{92}. A standard ``T2T recipe'' has emerged with
four main steps: error correction of accurate long reads, construction
of an initial assembly graph (overlap or de~Bruijn), simplification of
that graph using ultra‑long ONT reads, and final scaffolding and phasing
with long‑range data such as Hi‑C, Strand‑seq, or trio
information\textsuperscript{92}. Accurate long reads provide the
backbone, ultra‑long reads thread through remaining tangles, and
long‑range contacts connect phased blocks into chromosome‑scale
scaffolds\textsuperscript{92}.

\begin{tcolorbox}[enhanced jigsaw, rightrule=.15mm, leftrule=.75mm, colbacktitle=quarto-callout-note-color!10!white, colframe=quarto-callout-note-color-frame, toprule=.15mm, opacitybacktitle=0.6, title=\textcolor{quarto-callout-note-color}{\faInfo}\hspace{0.5em}{Figure 6.5 -- Telomere‑to‑telomere assembly strategy}, left=2mm, colback=white, bottomtitle=1mm, toptitle=1mm, bottomrule=.15mm, coltitle=black, titlerule=0mm, breakable, opacityback=0, arc=.35mm]

\begin{figure}[H]

\centering{

\includegraphics[width=0.8\linewidth,height=\textheight,keepaspectratio]{images/Fig_6.5_LiDurbin_Assembly.jpg}

}

\caption{\label{fig-t2t-assembly}T2T Genome Assembly}

\end{figure}%

Strategy for near telomere‑to‑telomere assembly of large eukaryotic
genomes. Adapted from Li and Durbin (2024). Accurate long reads are
first error‑corrected and assembled into an initial graph; ultra‑long
ONT reads are then threaded through the graph to resolve tangles and
patch small gaps; finally, long‑range data such as Hi‑C, Strand‑seq, or
trio information scaffold and phase the assembly into chromosome‑scale
haplotypes\textsuperscript{92}.

\end{tcolorbox}

Modern assemblers such as HiCanu, hifiasm, Verkko, and LJA embody this
strategy using different graph and error‑correction
choices\textsuperscript{92}. HiCanu and hifiasm operate in overlap or
string‑graph space and exploit exact overlaps between HiFi reads to
distinguish repeat copies and phase haplotypes\textsuperscript{92}.
Verkko and LJA instead build multiplex de~Bruijn graphs that combine
information from multiple \((k)\) values and use minimizer‑based
sparsification to keep the graph tractable while still separating long,
similar repeats\textsuperscript{92}. In all cases, familiar ideas from
the alignment section---\((k)\)‑mers, minimizers, and efficient
indices---reappear, but now applied to read--read relationships and
graph traversal rather than read--reference
matching\textsuperscript{92}.

\begin{tcolorbox}[enhanced jigsaw, rightrule=.15mm, leftrule=.75mm, colbacktitle=quarto-callout-note-color!10!white, colframe=quarto-callout-note-color-frame, toprule=.15mm, opacitybacktitle=0.6, title=\textcolor{quarto-callout-note-color}{\faInfo}\hspace{0.5em}{Case Study: The 2022 telomere‑to‑telomere human genome}, left=2mm, colback=white, bottomtitle=1mm, toptitle=1mm, bottomrule=.15mm, coltitle=black, titlerule=0mm, breakable, opacityback=0, arc=.35mm]

The T2T‑CHM13 assembly published in 2022 provided the first essentially
complete sequence of a human genome, filling in centromeres, acrocentric
short arms, and other highly repetitive regions that were missing or
only crudely modeled in earlier references such as
GRCh38\textsuperscript{92,103}. The project combined PacBio HiFi reads,
ONT ultra‑long reads, and complementary long‑range data to build,
simplify, and phase an assembly graph to true telomere‑to‑telomere
continuity on most chromosomes\textsuperscript{92,103}.

From a three‑Cs perspective, T2T‑CHM13 dramatically improved continuity
(full chromosomal paths), completeness (recovery of large satellite
arrays and segmental duplications), and correctness (better resolution
of structural variants and fewer false duplications), providing a much
richer substrate for downstream alignment, variant calling, and
pan‑genome construction\textsuperscript{92,103}. Later work extended
this to the human Y chromosome and to multiple diploid human assemblies,
illustrating how the same T2T strategies scale from a haploid‑like cell
line to more complex diploid genomes\textsuperscript{92}. In the
practice problems at the end of this chapter, you will revisit T2T‑CHM13
using the three‑Cs framework and compare it to GRCh38.

\end{tcolorbox}

\subsection{Long‑read advances and the assembly--alignment feedback
loop}\label{longread-advances-and-the-assemblyalignment-feedback-loop}

The same long‑read technologies that motivated new alignment algorithms
in the previous section also drive the transition to T2T assembly.
Accurate long reads transform repeats and haplotypes from intractable
obstacles into resolvable graph features: small sequence differences
between repeat copies or homologous chromosomes become informative edges
and bubbles rather than sources of unavoidable
collapse\textsuperscript{92}. Assemblers exploit exact long overlaps or
long \((k)\)‑mers to separate these structures, while aligners such as
Minimap2 and Winnowmap2 then map long reads or assembled contigs back to
references to validate assemblies, call variants, and support polishing
and quality assessment\textsuperscript{92}.

Haplotypes in diploid genomes can be viewed as long, highly similar
repeats, so algorithms that are good at resolving repeats are naturally
good at phasing and vice versa\textsuperscript{92}. Tools such as
hifiasm and Verkko build phased assembly graphs in which bubbles
represent heterozygous loci, and long‑range data plus ultra‑long reads
help choose consistent paths through these bubbles, producing
chromosome‑scale haplotype‑resolved assemblies\textsuperscript{92}.
Downstream, graph‑aware aligners and k‑mer‑based evaluators close the
loop by aligning reads and assemblies back to these graphs to quantify
correctness and identify residual gaps or mis‑joins\textsuperscript{92}.

\section{From assemblies to pan‑genomes and graph
mapping}\label{from-assemblies-to-pangenomes-and-graph-mapping}

In Chapter~2 you encountered the idea of a \textbf{pan‑genome}: a
reference that represents the diversity of many genomes in a species or
clade rather than a single linear haplotype. T2T assemblies and
population‑scale projects now provide raw material for such references,
while graph‑based mapping algorithms extend alignment ideas to this new
setting\textsuperscript{92}. Instead of a single linear coordinate
system, a pan‑genome graph encodes shared sequence as common paths and
diverging alleles or structural variants as alternative branches,
allowing reads or assembled haplotypes from new individuals to be mapped
into a richer, more representative structure\textsuperscript{92}.

Tools such as minigraph build reference graphs from multiple
high‑quality assemblies using minimizer‑based sketches, and
sequence‑to‑graph aligners then map new reads or assemblies to these
graphs\textsuperscript{92}. Conceptually, this extends the
seed‑and‑chain paradigm from linear references into graphs: seeds are
found as minimizers or \((k)\)‑mers shared between reads and graph
nodes, chains are constructed as paths through the graph, and dynamic
programming refines alignments locally along candidate
paths\textsuperscript{92}. It may be helpful to recall that these graph
mapping approaches are natural descendants of the overlap and de~Bruijn
graphs introduced here, repurposed from assembling a single genome to
representing and aligning many genomes at once.

\section{Reflection questions}\label{reflection-questions-5}

\begin{enumerate}
\def\labelenumi{\arabic{enumi}.}
\item
  \textbf{Revisiting Bogan et~al.~Table~1}: Return to Table~1 in Bogan
  et~al., which summarizes their genome assemblies and key metrics such
  as contig N50, total size, and BUSCO completeness\textsuperscript{3}.
  Using the three Cs framework from Chapter~2 and this chapter's
  discussion of assembly evaluation, argue whether you consider each
  assembly ``fit for purpose'' for their antifreeze‑protein questions.
  Where would you be most cautious in interpreting their downstream
  analyses?
\item
  \textbf{Graph choices and repeats}: Consider a medium‑sized vertebrate
  genome with many recent transposon insertions. Describe how an
  overlap/string‑graph assembler and a de~Bruijn‑graph assembler would
  represent a family of 6~kb LINE elements, and explain how read length
  and \emph{k} influence whether those repeats become resolvable paths
  or persistent tangles\textsuperscript{90,92}.
\item
  \textbf{From short‑read drafts to T2T references}: Compare a
  short‑read draft assembly of a genome to a hypothetical T2T assembly
  of the same species. Using continuity, completeness, and correctness,
  list at least three types of analyses (for example, structural‑variant
  discovery, repeat evolution, gene copy‑number estimation) that become
  more reliable in the T2T setting, and explain why\textsuperscript{92}.
\item
  \textbf{Haplotypes as repeats}: Explain why a heterozygous diploid
  genome can be viewed as a special case of a repeat‑rich genome, and
  how this perspective motivates phasing algorithms that operate on
  chains of bubbles in an assembly graph\textsuperscript{92}. How does
  this viewpoint change the way you interpret ``duplicated'' BUSCO genes
  or excess k‑mer multiplicity in an assembly report?
\item
  \textbf{Pan‑genomes and reference bias}: Using the concept of a
  pan‑genome introduced in Chapter~2, describe how a genome graph that
  incorporates multiple high‑quality assemblies can reduce reference
  bias compared with a single linear reference\textsuperscript{92}. What
  kinds of variants or loci do you expect to benefit most from a
  graph‑based representation?
\end{enumerate}

\section{Practice Problem Set
Questions}\label{practice-problem-set-questions}

\begin{enumerate}
\def\labelenumi{\arabic{enumi}.}
\tightlist
\item
  \textbf{Interpreting a real assembly: Bogan et~al}:

  \begin{enumerate}
  \def\labelenumii{\alph{enumii}.}
  \tightlist
  \item
    Using Table~1 in Bogan et~al., classify each reported assembly in
    terms of continuity, completeness, and
    correctness\textsuperscript{3}.\\
  \item
    Identify one assembly that you would trust for gene‑level analyses
    but not for structural‑variant analyses, and justify your answer
    using specific metrics.\\
  \item
    Propose one additional evaluation (for example, k‑mer spectrum,
    alignment of raw reads, or comparison to a related reference) that
    would strengthen their assessment of assembly quality.
  \end{enumerate}
\item
  \textbf{Designing an assembly strategy}: You are given a budget to
  assemble the genome of a non‑model amphibian with an estimated genome
  size of 5~Gb and moderate heterozygosity.

  \begin{enumerate}
  \def\labelenumii{\alph{enumii}.}
  \tightlist
  \item
    Propose a data‑generation plan (sequencing technologies, approximate
    coverage, long‑range data) aimed at a high‑quality but not
    necessarily T2T assembly\textsuperscript{92}.\\
  \item
    Choose an assembly paradigm (overlap/string graph vs.~de~Bruijn
    vs.~multiplex DBG) and two candidate assemblers, and explain how
    their graph strategies relate to the genome's repeat structure.\\
  \item
    Outline how you would evaluate the resulting assembly using both
    gene‑based and k‑mer‑based metrics.
  \end{enumerate}
\item
  \textbf{Algorithm thought experiment: \emph{k}‑mer size}: Imagine
  assembling the same bacterial genome with three de~Bruijn graphs at (k
  = 31), (k = 101), and (k = 501), using a long‑read dataset with mean
  read length 10~kb\textsuperscript{10,92}.

  \begin{enumerate}
  \def\labelenumii{\alph{enumii}.}
  \tightlist
  \item
    Predict how the number of contigs and presence of unresolved repeats
    will change as (k) increases.\\
  \item
    Explain why extremely large (k) values may still lead to fragmented
    graphs in low‑coverage regions, even when repeats are fully
    resolvable in principle.
  \end{enumerate}
\item
  \textbf{T2T assembly planning for a Bogan‑like system}: Suppose you
  want to move from the assemblies in Bogan et~al.~to a T2T‑style
  assembly for a closely related polar fish
  species\textsuperscript{3,92}.

  \begin{enumerate}
  \def\labelenumii{\alph{enumii}.}
  \tightlist
  \item
    Specify additional data types and coverage you would request (for
    example, HiFi, ONT ultra‑long, Hi‑C, trio sequencing) and how each
    would help resolve particular challenges in their genome.\\
  \item
    Describe how you would integrate these data in the four‑step T2T
    recipe (error correction, graph construction, ultra‑long
    integration, scaffolding/phasing).\\
  \item
    List two concrete scientific questions about antifreeze gene
    evolution that might become answerable only once a T2T or near‑T2T
    reference is available.
  \end{enumerate}
\item
  \textbf{Pan‑genome mapping exercise (conceptual)}: Consider a human
  pan‑genome graph built from GRCh38, T2T‑CHM13, and several
  high‑quality diploid assemblies\textsuperscript{92,103}.

  \begin{enumerate}
  \def\labelenumii{\alph{enumii}.}
  \tightlist
  \item
    Sketch how a structural variant present in only one haplotype (for
    example, a 20~kb insertion) would appear in the graph relative to a
    linear reference.\\
  \item
    Explain how a minimizer‑based seed‑and‑chain algorithm could map a
    new HiFi read spanning this insertion onto the graph.\\
  \item
    Discuss how mapping reads to this graph might change the variant
    calls you would obtain compared with mapping to GRCh38 alone.
  \end{enumerate}
\end{enumerate}

\section{References}\label{references-6}

\chapter{Project Management and Writing
Methods}\label{project-management-and-writing-methods}

A Guide to Reproducible Computational Research

\hfill\break

This chapter delves into some topics that are not covered in typical
bioinformatics course, but are essential to \emph{doing} bioinformatics.
Specifically, we dive deep into two major components of reproducible
bioinformatics research - project management and methods sections. As
such, this chapter is split into these two major sections to cover these
critical best practices and reveal these aspects of hidden curriculum in
this field.

\section{Section 1: Project Management for Bioinformatics
Research}\label{section-1-project-management-for-bioinformatics-research}

Bioinformatics projects are inherently complex: they involve large
datasets, many computational steps, diverse software tools, and often
distributed teams. Digital work is fragile---files are easy to
overwrite, analysis steps are easy to forget, and key decisions are
rarely obvious months later. Systematic project organization and
documentation are therefore not optional extras but the \emph{foundation
of reproducible research} and your own future
productivity\textsuperscript{104}. A small investment in organization at
the start of a project pays off every time you need to rerun, extend, or
defend an analysis.

\begin{tcolorbox}[enhanced jigsaw, rightrule=.15mm, leftrule=.75mm, colbacktitle=quarto-callout-note-color!10!white, colframe=quarto-callout-note-color-frame, toprule=.15mm, opacitybacktitle=0.6, title=\textcolor{quarto-callout-note-color}{\faInfo}\hspace{0.5em}{Figure 7.1 -- Importance of Digital Organization}, left=2mm, colback=white, bottomtitle=1mm, toptitle=1mm, bottomrule=.15mm, coltitle=black, titlerule=0mm, breakable, opacityback=0, arc=.35mm]

\begin{figure}[H]

\centering{

\includegraphics[width=0.5\linewidth,height=\textheight,keepaspectratio]{images/Figure_7-1_Organization.png}

}

\caption{\label{fig-organization}Contrasting cluttered and organized lab
benches; how does organization make you feel?}

\end{figure}%

Just as a cluttered bench slows down wet‑lab work, a chaotic directory
structure makes simple computational tasks stressful and error‑prone.
Your goal is to make your digital lab bench feel like the image on the
right and now those on the left.

\end{tcolorbox}

\subsection{Foundational Principles}\label{foundational-principles}

At the heart of effective project management in bioinformatics lie
several core principles that guide all subsequent practices. First among
these is the principle that everything you do, \textbf{you will probably
have to do over again}\textsuperscript{104}. This may sound pessimistic,
but it reflects the iterative nature of computational research. Reviewer
requests for new analyses or access to additional samples are much
easier to address when your earlier work is well documented.

The second principle recognizes that project organization exists not
just for human readers, but to facilitate automation. Bioinformatics
routinely involves processing hundreds or thousands of files, and the
ability to \textbf{programmatically} access and manipulate these files
is essential. This automation depends on predictable directory
structures and consistent file naming schemes. When your sequencing
files follow a systematic naming pattern like
\texttt{sample-A\_replicate-1\_read-R1.fastq.gz}, writing a script to
process all samples becomes straightforward. When files are
inconsistently named like \texttt{sampleA-1.fastq} and
\texttt{sample\_A\_rep2\_R1.fq.gz}, the same task becomes frustrating
and time intensive\textsuperscript{105}.

The third principle emphasizes that documentation should be
comprehensive yet lightweight. Every analytical step that produces
results used in your work needs to be documented somewhere, including
the full command lines, all parameter values (even defaults), software
versions, and the provenance of all input data. However, this
documentation need not be burdensome. Plain text formats like
\textbf{Markdown} allow you to quickly record information in a format
that is readable, searchable, and future-proof.

\begin{tcolorbox}[enhanced jigsaw, rightrule=.15mm, leftrule=.75mm, colbacktitle=quarto-callout-tip-color!10!white, colframe=quarto-callout-tip-color-frame, toprule=.15mm, opacitybacktitle=0.6, title=\textcolor{quarto-callout-tip-color}{\faLightbulb}\hspace{0.5em}{Workflow Thinking in Bioinformatics}, left=2mm, colback=white, bottomtitle=1mm, toptitle=1mm, bottomrule=.15mm, coltitle=black, titlerule=0mm, breakable, opacityback=0, arc=.35mm]

Modern genomic analyses involve complex chains of computational
steps---downloading data, quality control, alignment, variant calling,
and downstream analysis. Rather than writing monolithic shell scripts
that hardcode file paths and parameters, workflow thinking means viewing
your analysis as a \textbf{directed acyclic graph (DAG)} where each
computational step is a node and dependencies are edges. This approach
naturally leads to modular, reproducible, and parallelizable analyses.
Workflow managers like Snakemake, Nextflow, and Common Workflow Language
(CWL) implement these principles by letting you separate configuration
from code, explicitly declare inputs and outputs for each step, and
automatically parallelize independent tasks. For example, the
\textbf{snpArcher pipeline} demonstrates how this thinking enables
consistent reanalysis of diverse datasets\textsuperscript{106}. By
defining samples in structured tables and parameters in configuration
files, the authors processed datasets ranging from 6 to 306 individuals
using identical variant calling logic---what would take weeks of manual
script-wrangling became a single command per dataset.

The key insight is to design for reproducibility from the start: each
workflow step should be \textbf{idempotent}, or an operation or function
that produces the same result regardless of how many times it is
executed. Configuration should live in YAML or CSV files rather than
scattered through code, and the DAG structure should make
parallelization opportunities obvious. When snpArcher splits genome
intervals at assembly gaps rather than processing chromosomes
sequentially, the workflow manager automatically distributes hundreds of
small tasks across available compute resources, reducing runtime from
days to hours. This same workflow runs unchanged on laptops, HPC
clusters, or cloud platforms by simply switching resource profiles.
Students new to workflow thinking should start by converting a single
analysis pipeline to a workflow manager, focusing first on getting the
basic DAG working before optimizing for parallelization. The investment
pays immediate dividends: analyses become faster, more reliable, and
truly reproducible---collaborators can recreate exact results years
later because both code and configuration are versioned together.

\end{tcolorbox}

\subsection{Project Directory
Structure}\label{project-directory-structure}

The foundation of any well-organized bioinformatics project is a clear
\textbf{directory structure} that logically separates different types of
content while maintaining everything related to a single project under a
common root directory. This seemingly simple practice has profound
implications for project management. When all files live within a single
directory hierarchy, the project becomes portable between systems,
collaborators can easily obtain all necessary materials, and version
control systems can track the entire project as a cohesive
unit\textsuperscript{105}.

A practical directory structure balances logical organization at the top
level with chronological organization for experimental work. Consider a
project investigating SNP calling in maize. At the root level, you might
have directories for \texttt{data/} to store raw and intermediate
datasets, \texttt{scripts/} for analysis code, \texttt{results/} for
computational experiments, and \texttt{docs/} for manuscripts and
documentation. This top-level logical organization makes it immediately
clear where different types of content belong. Within the \texttt{data/}
and results/ directories, however, chronological organization often
works better than purely logical schemes. Creating subdirectories named
by date (using the \texttt{YYYY-MM-DD} format for proper sorting)
provides an automatic record of when work was performed and allows you
to easily track the evolution of your analyses\textsuperscript{104}.

This chronological approach prevents a common pitfall where researchers
create elaborate logical directory structures that quickly become
outdated as projects evolve. What seems like a well-planned set of
analysis categories at the project's start may bear little resemblance
to the final structure of completed work. Chronologically organized
directories adapt naturally to this evolution, and when combined with
descriptive README files in each directory, provide both structure and
flexibility. For example, a directory named
\texttt{2024-01-15-quality-control/} clearly indicates when this work
occurred, while its README file explains what quality control steps were
performed and why.

Within individual analysis directories, you can reintroduce logical
organization as needed. If a single experiment generates many output
files, creating \textbf{subdirectories} for different file types
maintains clarity. The key is to let the organization emerge from the
actual structure of your work rather than trying to impose a rigid
scheme in advance. This organic approach, guided by consistent
documentation, creates project structures that remain understandable
months or years later\textsuperscript{107}.

One critical detail in project organization concerns the use of
\textbf{relative versus absolute paths}. When scripts or configuration
files reference other files in your project, always use relative paths
like \texttt{../data/reference.fasta} rather than absolute paths like
\texttt{/home/username/projects/maize-snps/data/reference.fasta}.
Relative paths ensure that your project remains functional regardless of
where it is located on the filesystem, making it truly portable between
collaborators and systems. Absolute paths, by contrast, embed
assumptions about directory structures that break as soon as the project
is moved\textsuperscript{105}.

\subsection{File Naming Conventions}\label{file-naming-conventions}

File naming in bioinformatics is more than a matter of personal
preference; it directly impacts your ability to work efficiently with
your data. The Unix command line and most programming languages provide
powerful tools for programmatically accessing files, but these tools
become useless when filenames are inconsistent or contain problematic
characters. The transition from graphical user interfaces, where
\textbf{spaces in filenames} pose no problems, to command-line
environments, where spaces must be carefully escaped, catches many
researchers off guard. A filename like \texttt{raw\ sequences.fastq}
appears innocuous but creates constant headaches in Unix environments
where spaces separate command arguments\textsuperscript{105}.

Best practice dictates using only letters, numbers, underscores, and
dashes in filenames, while avoiding spaces, special characters, and
symbols. This restriction may seem limiting, but it ensures your
filenames work seamlessly across different operating systems and
command-line environments. Additionally, including \textbf{file
extensions} that clearly indicate content type greatly aids
interpretation. A file named \texttt{zmays-genes.fasta} immediately
signals that it contains sequence data in FASTA format, while a file
named simply \texttt{zmays-genes} provides no such clarity.

Consistency in naming schemes becomes particularly important when
working with multiple related files. Consider a project with paired-end
sequencing data for three maize samples. Filenames like
\texttt{zmaysA\_R1.fastq}, \texttt{zmaysA\_R2.fastq},
\texttt{zmaysB\_R1.fastq}, \texttt{zmaysB\_R2.fastq},
\texttt{zmaysC\_R1.fastq}, and \texttt{zmaysC\_R2.fastq} follow a
predictable pattern where sample names and read pairs occupy fixed
positions in the filename. This consistency enables powerful pattern
matching using shell wildcards. The expression \texttt{zmaysB*} matches
all files for sample B regardless of read pair, while
\texttt{zmays*\_R1.fastq} matches all forward reads across all samples.
These seemingly simple pattern matching capabilities become invaluable
when processing hundreds of samples. For numbered files, use leading
zeros (gene-001.txt) so alphabetical sorting matches numeric
order\textsuperscript{105}.

The consequences of poor file naming extend beyond inconvenience. In
collaborative projects, confusion over which file represents the current
version of an analysis can cascade into wasted effort or, worse, serious
errors that make it into publications. The small investment in
establishing and maintaining consistent naming conventions provides
insurance against these costly mistakes\textsuperscript{108}.

\subsection{Documentation and
Metadata}\label{documentation-and-metadata}

Without adequate documentation, a project directory becomes an
archaeological site where future researchers (often your future self)
must attempt to reconstruct what happened and why, with only the
artifacts of computation as evidence. Therefore, you should
\textbf{document everything} about how you obtained and processed your
data, because the details that seem obvious today will be mysterious
later\textsuperscript{109}. At the most fundamental level, you must
document the origin of all data in your project. This includes not just
experimental data you generated, but also reference genomes, annotation
files, sequence databases, and any other external resources. For each
dataset, record where you obtained it, when you downloaded it, how you
downloaded it, and any version information available. A simple
\textbf{README} file in your data directory might note that the mouse
reference genome version GRCm38 was downloaded from Ensembl on a
specific date using a specific \textbf{wget} command. This level of
detail may seem excessive, but external databases change over time, file
versions are updated, and resources sometimes disappear entirely, making
this documentation invaluable years later\textsuperscript{105}.

Version information deserves special attention because even minor
software updates can produce different results. Good documentation
records not just the name of each program used, but its version number,
the parameters specified (including default values), and ideally the
exact command line used. Many modern bioinformatics tools provide a flag
to display version information (often \texttt{-\/-version} or
\texttt{-v}), and this version string should be captured in your
documentation. For software managed through version control systems like
Git, the specific commit hash provides an unambiguous identifier of the
exact code version used\textsuperscript{109}.

The chronological aspects of data acquisition and analysis also matter
for reproducibility. When you download a dataset from a database that is
regularly updated, the current version of that database may differ from
what was available when you performed your original analysis. Recording
the \textbf{download date} provides context for understanding potential
discrepancies. Similarly, documenting when each analysis was performed
helps track the evolution of your understanding and can be critical for
debugging when results change unexpectedly\textsuperscript{104}.

Beyond raw provenance information, documentation should capture the
\textbf{reasoning behind your choices}. Why did you choose one assembly
tool over another? What led you to select particular quality filtering
thresholds? Which preliminary analyses informed your final approach?
This narrative documentation need not be lengthy, but it should provide
enough context that someone reading your notes can understand your
thought process. Such documentation often proves most valuable when it
records what did not work, as this prevents future repetition of failed
approaches and provides crucial negative results that rarely appear in
publications but are essential for scientific
progress\textsuperscript{107}. These habits will save you time when you
later write up your methods and data availability statements (see
Section 2 below).

\subsection{Version Control
Integration}\label{version-control-integration}

Version control systems, particularly \textbf{Git}, have become
indispensable tools for managing bioinformatics projects. At its core,
version control maintains a complete history of your project, allowing
you to track changes over time, recover from mistakes, and maintain
multiple versions of your work simultaneously\textsuperscript{107}.

The discipline of \textbf{regular commits} to version control serves
multiple purposes in project management. Each commit represents a
snapshot of your project at a specific point in time, creating a
detailed historical record of how your analysis evolved. This becomes
invaluable when you need to reproduce exactly what was done three months
ago, or when you need to determine when a particular change was made
that affected your results. The commit messages themselves form a
chronological log of your project's development, supplementing your
README documentation with a more granular record of
changes\textsuperscript{109}.

However, not all files in your project directory should be placed under
version control; only the scripts and documentation needed to regenerate
your workflow need to be tracked\textsuperscript{104}. Large files and
computational results that can be regenerated from documented workflows
do not need version tracking, and can be ignored by adding these files
and directories to your \texttt{.gitignore} file.

The \textbf{collaborative features} of Git and GitHub extend the value
of version control beyond individual project
management\textsuperscript{110}. When working with collaborators, a
shared Git repository becomes the single source of truth for project
code and documentation. Team members can work on different features
simultaneously through branching, merge their changes systematically,
and track who made which modifications and why. This collaborative
workflow dramatically reduces the email attachments, confused file
versions, and communication overhead that plague unstructured
collaborative projects\textsuperscript{108}. Clear commit histories and
tagged releases also make it possible to cite the exact version of your
code in the methods section (see Section 2).

\subsection{Agile Project Management and Kanban
Boards}\label{agile-project-management-and-kanban-boards}

While the bioinformatics literature has primarily focused on technical
aspects of project organization, project management methodologies from
software engineering offer valuable frameworks for organizing
computational research. Agile methodologies, particularly when
implemented through visual tools like \textbf{Kanban boards}, provide
structure for managing the iterative, exploratory nature of
bioinformatics projects while maintaining the flexibility that research
demands.

The core insight of Agile approaches is that complex projects are better
managed through iterative development cycles rather than rigid upfront
planning. In a bioinformatics context, this aligns naturally with how
research actually proceeds. You might begin with a broad research
question, but the specific analyses required only become clear as you
explore the data and obtain preliminary results. Rather than attempting
to plan every analysis step months in advance, Agile methodologies
embrace this uncertainty by breaking work into \textbf{small, manageable
tasks} that can be completed in short iterations\textsuperscript{107}.

Kanban boards provide a visual system for tracking these tasks through
their lifecycle. At its simplest, a Kanban board has three columns: ``To
Do,'' ``In Progress,'' and ``Done.'' Each task is represented by a card
that moves through these columns as work progresses. This visualization
serves several purposes. It provides an immediate overview of project
status, makes work-in-progress explicit (preventing the common problem
of starting too many tasks and finishing none), and creates a satisfying
sense of progress as cards move to the ``Done'' column. The physical act
of moving a card provides positive reinforcement that sustains
motivation through long projects.

\begin{tcolorbox}[enhanced jigsaw, rightrule=.15mm, leftrule=.75mm, colbacktitle=quarto-callout-note-color!10!white, colframe=quarto-callout-note-color-frame, toprule=.15mm, opacitybacktitle=0.6, title=\textcolor{quarto-callout-note-color}{\faInfo}\hspace{0.5em}{Figure 7.2 -- Example KanBan Board}, left=2mm, colback=white, bottomtitle=1mm, toptitle=1mm, bottomrule=.15mm, coltitle=black, titlerule=0mm, breakable, opacityback=0, arc=.35mm]

\begin{figure}[H]

\centering{

\includegraphics[width=0.5\linewidth,height=\textheight,keepaspectratio]{images/Figure_7-2_Kanban.png}

}

\caption{\label{fig-kanban}Example KanBan board for the life cycle of a
research project}

\end{figure}%

Here is an example of how you can extend the simple three‑column model
into stages that better match a research project lifecycle, illustrating
how flexible these tools can be. This research project Kanban board has
customized swim lanes for stages of a research project, from when they
start out as simple ideas to when they are accepted for publication.

\end{tcolorbox}

GitHub Projects integrates Kanban functionality directly into the
platform where many bioinformatics projects already live. This
integration is particularly powerful because it connects tasks directly
to code, documentation, and issues. A card on your Kanban board might
represent implementing a new quality control step, and that card can
link to the specific GitHub issue discussing the approach, the pull
request containing the implementation, and the commit that ultimately
merged the code. This creates a rich web of connections between project
planning and project execution that greatly aids both doing the work and
documenting what was done\textsuperscript{107}.

The level of granularity for tasks on a Kanban board requires some
judgment. Tasks should be small enough to complete in a reasonable
timeframe (hours to a few days rather than weeks), but not so small that
the overhead of managing tasks exceeds their value. In practice, tasks
like ``download reference genome and verify checksums,'' ``run quality
control on batch 1 samples,'' or ``implement GC content calculation
function'' represent appropriate granularity. Larger objectives like
``complete genome assembly'' should be broken down into their
constituent steps, each becoming a separate card.

One valuable practice is to include cards for documentation and quality
assurance alongside cards for analysis tasks. A card might represent
``document quality control parameters in README'' or ``review and test
assembly script.'' This ensures that documentation and quality control
remain integral to your workflow rather than afterthoughts. The
visibility of these tasks on your board reinforces their importance and
helps prevent the common pattern where documentation is perpetually
postponed until some mythical future time when the ``real work'' is
finished\textsuperscript{108}.

\subsection{Practical Workflow
Integration}\label{practical-workflow-integration}

Understanding principles and tools is valuable, but seeing how they
integrate into an actual workflow brings these concepts to life.
Consider a realistic bioinformatics project investigating genetic
variants associated with drought resistance in a crop plant. The project
begins with planning and setup, proceeds through data generation and
analysis, and concludes with interpretation and communication of
results. Throughout this lifecycle, the project management practices we
have discussed work together to maintain organization and
reproducibility.

Project initialization establishes the foundation. You create a project
directory with a descriptive name like \texttt{prdm9-protein-evolution},
initialize it as a Git repository, and establish the basic directory
structure:

\begin{verbatim}
project-root/
├── data/           # Project datasets
│   ├── README.md
│   ├── fastq/
│   │   ├── sample1_R1.fastq.gz
│   │   ├── sample1_R2.fastq.gz
│   │   └── ...
│   └── fasta/
│       └── GRCh38.fa
├── analysis/        # Computational experiments
├── docs/           # Documentation and manuscripts
├── scripts/        # Scripts and code
│   ├── 01_fetch_sra.sh
│   ├── 02_index_genome.sh
│   └── 03_trimming.sh
└── README.md       # Project overview
\end{verbatim}

A root-level \textbf{README} file documents the project's objectives,
your collaborators, and high-level organization. You create a GitHub
repository and push your initial structure, then set up a GitHub Project
with a Kanban board. Initial cards capture the major milestones:
obtaining sequences, quality control, variant calling, association
analysis, and manuscript writing. Each of these will later be broken
down into smaller tasks\textsuperscript{105}.

\begin{tcolorbox}[enhanced jigsaw, rightrule=.15mm, leftrule=.75mm, colbacktitle=quarto-callout-note-color!10!white, colframe=quarto-callout-note-color-frame, toprule=.15mm, opacitybacktitle=0.6, title=\textcolor{quarto-callout-note-color}{\faInfo}\hspace{0.5em}{Lab Connections}, left=2mm, colback=white, bottomtitle=1mm, toptitle=1mm, bottomrule=.15mm, coltitle=black, titlerule=0mm, breakable, opacityback=0, arc=.35mm]

In \href{Lab04.qmd}{Lab 4} you will get access to a Sample Genome
Analysis working directory that is organized using the same directory
struture described above. As you work through the yeast genomics dataset
this semester, you will continue to build this sample working directory
which can serve as an example moving forward from this course.

\end{tcolorbox}

Data acquisition requires careful documentation. When you download the
reference genome from a public database, you record the exact URL,
download date, genome version, and calculate \textbf{checksums} to
verify integrity. This information goes into \texttt{data/README.md}
along with the specific wget command used. The reference genome file
itself remains outside version control due to its size, but your
documentation ensures anyone can obtain the identical file. The same
applies to your experimental sequencing data; detailed \textbf{metadata}
about samples, sequencing runs, and quality metrics are documented even
though the large FASTQ files themselves are not
versioned\textsuperscript{109}.

As analysis proceeds, your workflow becomes a rhythm of task management,
execution, and documentation. Each morning, you review your Kanban board
to select the next task. Perhaps today's card is ``implement read
quality filtering.'' You create a new script in the \texttt{scripts/}
directory with a descriptive name like \texttt{filter\_reads.py},
following your established naming conventions. As you write the script,
you document your approach with inline comments and ensure all
parameters are configurable rather than hard-coded. When the script
works correctly, you commit it to Git with a clear commit message
explaining what the script does and why certain filtering thresholds
were chosen. You move the card to ``Done'' on your Kanban board and add
a new card for ``apply quality filtering to all
samples''\textsuperscript{107}.

Running analyses at scale demonstrates the value of consistent file
organization. Because your FASTQ files follow the naming pattern
\texttt{sample-ID\_replicate\_read.fastq.gz}, you can write a simple
loop that processes all samples without manually specifying each
filename. The output goes into a dated subdirectory within
\texttt{results/}, like \texttt{results/2026-01-28-fastqc-report/}. Each
results directory contains its own README documenting exactly what was
run, which version of each script was used, what parameters were
specified, and which input files were processed. The full command line
is recorded, often by redirecting the script's verbose output to a log
file that becomes part of the permanent record\textsuperscript{105}.

When you encounter problems, version control and documentation prove
their worth. A script that worked last week now produces errors. You use
Git to compare the current version against the last working version,
quickly identifying the problematic change. Or perhaps a collaborator
questions an analysis result. You can point them to the specific commit
containing the analysis script, the README documenting parameters, and
the log files showing exactly what was executed. This
\textbf{transparency builds trust} and facilitates the kind of critical
examination that strengthens science\textsuperscript{109}. These logs
often become the raw material you will transform into formal methods
text later in the manuscript.

Collaboration scenarios showcase GitHub's strengths. A colleague wants
to contribute a new analysis approach. They fork your repository, create
a branch for their work, and submit a pull request when ready. You
review their code, discuss approaches through GitHub's comment system,
and merge their contribution once satisfied. The entire discussion and
decision process becomes part of the project's permanent record.
Meanwhile, your Kanban board reflects everyone's contributions, with
cards assigned to different team members and status visible to
all\textsuperscript{108}.

\begin{tcolorbox}[enhanced jigsaw, rightrule=.15mm, leftrule=.75mm, colbacktitle=quarto-callout-note-color!10!white, colframe=quarto-callout-note-color-frame, toprule=.15mm, opacitybacktitle=0.6, title=\textcolor{quarto-callout-note-color}{\faInfo}\hspace{0.5em}{License agreements}, left=2mm, colback=white, bottomtitle=1mm, toptitle=1mm, bottomrule=.15mm, coltitle=black, titlerule=0mm, breakable, opacityback=0, arc=.35mm]

When you first set up a new GitHub repository, you can choose to add a
license. Many people skip this in the haste of getting started,
especially if it is a private repository. But if the goal is to
eventually make the code public, it is a good idea to add a license
agreement. A standard licence is the MIT license, which GitHub offers as
one of their templates. It is as easy as checking a box.

If you already have a repository and you want to add a license, simply
add a file and name it License.md and GitHub will automatically ask you
if you want to use one of their templates. You can learn more on the
\href{https://libguides.auburn.edu/scholcomm}{open access page} of the
AU Library.

\end{tcolorbox}

\subsection{Common Pitfalls and
Solutions}\label{common-pitfalls-and-solutions}

Even with good intentions and solid practices, certain challenges recur
in bioinformatics project management. Recognizing these common pitfalls
and knowing how to address them can save considerable frustration. One
frequent problem is the proliferation of slightly different versions of
analysis scripts or data files, often with names like
\texttt{analysis\_v2.py}, \texttt{analysis\_v2\_final.py},
\texttt{analysis\_v2\_final\_revised.py}, and so on. This typically
occurs when researchers are hesitant to overwrite files for fear of
losing previous work. The solution is version control; when all changes
are tracked through Git, you can confidently update files knowing that
previous versions remain accessible through the repository
history\textsuperscript{107}.

Another common issue is the accumulation of poorly documented
\textbf{intermediate files} whose purpose becomes mysterious over time.
A \texttt{results/} directory fills with files like
\texttt{output\_temp.txt}, \texttt{test\_results\_2.csv}, and
\texttt{filtered\_data\_new.tsv}, and six months later no one remembers
what they contain or whether they are still needed. This problem stems
from inadequate documentation during exploratory analysis. The solution
is discipline around naming (even temporary files should have
descriptive names) and documentation (README files should explain even
exploratory analyses). Some researchers also adopt a practice of regular
cleanup sessions where they review intermediate files, document any that
are valuable, and delete those that are not\textsuperscript{104}.

The challenge of \textbf{parameter proliferation} also trips up many
projects. A typical bioinformatics pipeline might involve dozens of
parameters across multiple tools, and tracking which parameter values
were used for which analysis becomes a nightmare. Some researchers try
to solve this by embedding parameters in filenames
(\texttt{results\_minscore30\_minlen50.txt}), but this quickly becomes
unwieldy. A better approach is to write scripts that log all parameters
to output files, or to use configuration files that are versioned
alongside scripts. Many researchers now adopt workflow management
systems specifically to handle parameter tracking
systematically\textsuperscript{109}.

Data organization can break down when projects grow beyond their initial
scope. The temptation to reorganize should be resisted, as
\textbf{reorganization breaks the connections} between old documentation
and current file locations. If reorganization is essential, it should be
done systematically with careful documentation of the changes and
updating of all relevant README files and analysis
scripts\textsuperscript{105}.

The integration of automated testing, a standard practice in software
development, remains underutilized in bioinformatics despite its value.
When you write a script to process sequencing data, how do you know it
works correctly? Manual testing with a few examples is better than
nothing, but automated tests that run every time you modify the script
provide much stronger assurance. Even simple checks, like verifying that
output files contain the expected number of records or that numerical
results fall within sensible ranges, can catch many errors before they
propagate through an analysis pipeline\textsuperscript{107}.

\section{Conclusion Part 1}\label{conclusion-part-1}

Effective project management in bioinformatics is ultimately about
making life easier for \textbf{your future self} and your collaborators
while advancing the broader goal of reproducible science. The practices
described in this chapter---systematic directory organization,
consistent file naming, comprehensive documentation, version control,
and structured task management---may seem like overhead when you are
eager to dive into analysis. However, this investment pays dividends
throughout a project's lifetime and beyond. The time spent documenting a
data download is trivial compared to the time lost trying to remember
months later where data came from. The discipline of following naming
conventions seems tedious until you need to programmatically process
hundreds of files. The ritual of regular Git commits feels unnecessary
until you need to recover from a mistake or track down when a change was
introduced\textsuperscript{104}.

Moreover, these practices scale naturally from individual work to
collaborative projects. When you maintain clear documentation and
organization for your own benefit, bringing a collaborator up to speed
becomes straightforward. When your scripts are version controlled and
well documented, sharing them with the broader community requires
minimal additional effort. When your data provenance is meticulously
recorded, responding to questions from reviewers or readers becomes a
matter of consulting your notes rather than reconstructing history from
fragments\textsuperscript{109}.

The specific tools and approaches will continue to evolve. New version
control systems will emerge, project management platforms will add
features, and workflow management systems will grow more sophisticated.
However, the \textbf{fundamental principles remain constant}: organize
systematically, document comprehensively, track changes explicitly, and
automate where possible. These principles transcend any particular tool
or technology and will serve you well regardless of how the technical
landscape changes.

Perhaps most importantly, good project management \textbf{makes science
better}. When analyses are reproducible, reviewers can verify results
and readers can build on your work. When methods are thoroughly
documented, other researchers can apply your approaches to new datasets
or extend them in novel directions. When negative results are documented
alongside positive findings, the scientific community avoids redundant
dead ends. In an era where reproducibility concerns have shaken
confidence in published findings, the practices of rigorous project
management represent not just personal virtue but \textbf{professional
responsibility}\textsuperscript{108}.

The path to mastery of these practices is \textbf{incremental}. You need
not implement everything at once, nor achieve perfection on your first
project. Start with the basics: create a clear directory structure,
maintain README files, use version control for your scripts, and name
files consistently. As these habits become second nature, layer on
additional practices like Kanban boards for task management or more
sophisticated documentation approaches. Over time, good project
management stops feeling like an external discipline imposed on your
work and becomes simply how you work. This transformation, from viewing
organization as a chore to experiencing it as an enabling foundation,
marks the transition from novice to experienced bioinformatician.

\section{Section 2: Writing Methods Sections for Bioinformatics
Research}\label{section-2-writing-methods-sections-for-bioinformatics-research}

Though it may feel daunting right now to think of writing a thesis,
dissertation, or research paper, at some point you will need to
communicate your findings in written form---and the project management
practices from Section 1 are what make that possible. The methods
section is the cornerstone of scientific
reproducibility\textsuperscript{111}. In bioinformatics and
computational biology, where complex analytical workflows involving
multiple software packages, databases, and statistical procedures are
routine, clear methods documentation becomes even more critical.

This section addresses the unique challenges of documenting
computational research, including how to cite software and databases,
manage supplementary materials, and ensure that your analytical
workflows can be understood and reproduced by other researchers. Whether
you're analyzing genomic data, building phylogenetic trees, or
developing machine learning models for biological applications, the
principles outlined here will help you communicate your methods
effectively.

The core principle underlying effective methods writing is simple:
\emph{provide sufficient detail for others to reproduce your work} while
organizing content logically and clearly. This requires balancing
completeness with readability, justifying your methodological choices
with citations, documenting all software versions and parameter
settings, and making data and code accessible through appropriate
repositories.

\begin{tcolorbox}[enhanced jigsaw, rightrule=.15mm, leftrule=.75mm, colbacktitle=quarto-callout-note-color!10!white, colframe=quarto-callout-note-color-frame, toprule=.15mm, opacitybacktitle=0.6, title=\textcolor{quarto-callout-note-color}{\faInfo}\hspace{0.5em}{Lab Connections}, left=2mm, colback=white, bottomtitle=1mm, toptitle=1mm, bottomrule=.15mm, coltitle=black, titlerule=0mm, breakable, opacityback=0, arc=.35mm]

The deliverable from Lab 4 will be for you to use what you learn in this
section to write a methods paragraph documenting the bioinformatics
steps you take in the upcoming lab assignment. Specifically, you will
need to parse the scripts provided to figure out what they do, and then
describe them as if you were writing the methods section of a paper.

\end{tcolorbox}

\subsection{From Workflow Diagram to Written
Methods}\label{from-workflow-diagram-to-written-methods}

Bioinformatics research typically involves transforming raw biological
data through complex computational pipelines to generate biological
insights. The challenge in writing methods sections is converting these
complex workflows---often represented as flowcharts with dozens of
steps---into clear, \textbf{readable prose} that allows others to
replicate your work. Consider the transformation from a typical
bioinformatics workflow diagram (raw reads → quality control → alignment
→ variant calling → annotation → statistical analysis) to a published
methods section. The key is to present this information chronologically
while providing sufficient technical detail at each step. If you have
maintained organized directories, clear file names, and README files as
described in Section 1, reconstructing this chronological narrative
becomes much simpler.

At each step of your workflow, you must include the \textbf{software
name and version number}, a citation to the software's primary
publication, key parameters (especially if different from defaults),
rationale for software choice when multiple options exist, and relevant
input and output file formats. For instance, rather than simply stating
\texttt{"reads\ were\ aligned\ to\ the\ reference\ genome"}, you should
specify which aligner was used, what version, which genome build, and
any non-default parameters employed.

\subsection{Citing Published Workflows and
Software}\label{citing-published-workflows-and-software}

When you follow an established protocol or workflow, \textbf{proper
citation} is essential. However, vague references are insufficient and
hinder reproducibility. A poor practice example would be:
\texttt{"Samples\ were\ processed\ similar\ to\ methods\ described\ previously\ in\ Smith\ et\ al.\ (2002)."}
This statement is problematic because it provides no specific details
about which parts were followed, and may require accessing paywalled
content to understand the methods. A better approach provides specific
detail:
\texttt{"RNA-seq\ libraries\ were\ prepared\ following\ the\ protocol\ of\ Smith\ et\ al.\ (2002),\ with\ modifications\ to\ the\ adapter\ ligation\ step.\ Specifically,\ we\ used\ Y-shaped\ adapters\ (Illumina\ TruSeq)\ instead\ of\ the\ original\ T-tailed\ adapters\ to\ improve\ library\ yield\ for\ degraded\ RNA\ samples."}
This clearly states what was followed, what was changed, and why the
modification was made.

Proper software citation requires three elements: the software name, the
specific version used, and a citation to the primary publication.
Version numbers are critical because software behavior can change
substantially between versions. For example, in the antifreeze protein
study\textsuperscript{3}, the authors clearly state:
\texttt{"Haplotype-phased\ genomes\ of\ all\ species\ except\ for\ O.\ amberensis\ were\ assembled\ from\ long-read\ libraries\ using\ hifiasm\ v0.19.5-r592\ (Cheng\ et\ al.\ 2021)\ set\ to\ default\ parameters."}
This complete citation includes the tool name (hifiasm), the exact
version number (v0.19.5-r592), the primary publication citation (Cheng
et al.~2021), and the parameters used (defaults). This level of detail
enables others to replicate the assembly process exactly.

\begin{tcolorbox}[enhanced jigsaw, rightrule=.15mm, leftrule=.75mm, colbacktitle=quarto-callout-tip-color!10!white, colframe=quarto-callout-tip-color-frame, toprule=.15mm, opacitybacktitle=0.6, title=\textcolor{quarto-callout-tip-color}{\faLightbulb}\hspace{0.5em}{Reference management}, left=2mm, colback=white, bottomtitle=1mm, toptitle=1mm, bottomrule=.15mm, coltitle=black, titlerule=0mm, breakable, opacityback=0, arc=.35mm]

Reference management software like \textbf{Zotero} and \textbf{EndNote}
allow you to collect, managing and cite references and software. These
and other tools have browser plug-ins that can allow you to build your
library of citations throughout your dissertation. Then, once you are
ready to write your thesis, dissertation or scientific papers, these
tools integrate seamlessly with Microsoft Word and Google Docs via free
plugins or add-ons, enabling ``Cite While You Write''. This allows you
to insert citations directly into your documents, change the format to
match journal styles, and create formatted bibliography sections.

Once inserted, citations and bibliographies update dynamically,
supporting collaboration with co-authors who are able to cite from
shared/group libraries or changing styles if you need to reformat for a
new journal. You can learn more about different citation managers
\href{https://libguides.auburn.edu/citationmanagers}{here}.
\href{https://libguides.auburn.edu/zotero}{Auburn Libraries} also offers
regular workshops to learn how to use these softwares.

\end{tcolorbox}

\subsection{Documenting Parameter Changes and
Deviations}\label{documenting-parameter-changes-and-deviations}

When you deviate from default parameters, you must document what changed
and why. This justification demonstrates methodological rigor and helps
readers understand your analytical decisions. Consider this example of
parameter documentation:
\texttt{"Variant\ calling\ was\ performed\ using\ GATK\ HaplotypeCaller\ v4.2.0\ (McKenna\ et\ al.\ 2010).\ We\ reduced\ the\ minimum\ base\ quality\ score\ from\ 10\ (default)\ to\ 5\ to\ increase\ sensitivity\ for\ low-coverage\ regions,\ accepting\ a\ modest\ increase\ in\ false\ positives\ that\ would\ be\ filtered\ in\ downstream\ analyses."}

Deviations from established protocols are common and often necessary.
The key is to document these changes transparently and \textbf{explain
the rationale}, demonstrating careful experimental design rather than ad
hoc modification. Parameter logs, configuration files, and
version‑controlled scripts created during project management become your
primary sources for this level of detail. You should also justify
deviations when you change published protocols or standard workflows,
use non-default software parameters, apply methods to non-standard
organisms or data types, omit standard quality control or filtering
steps, or use alternative software from what is commonly employed.

\subsection{Citing Public Data and
Databases}\label{citing-public-data-and-databases}

Bioinformatics research frequently incorporates publicly available data,
including reference genomes, variant databases, and previously published
datasets. Proper citation of these resources is essential for
reproducibility and gives credit to data generators. For any public
dataset or database, you should include the resource name and version or
build number, accession numbers or identifiers, URL or DOI for access,
date accessed for databases that update frequently, and citation to the
primary publication if available. For example, there is a great
\href{https://support.nlm.nih.gov/knowledgebase/article/KA-03391/en-us}{knowledge
base article} on citing NCBI Sequence Read Archive data. Specifically,
be sure to provide the BioProject, BioSample, and SRA accession numbers.
For databases and software without formal publications, cite the URL and
include the access date. When DOIs are available, use them as they
provide permanent, citable links. For example:
\texttt{"Protein\ domain\ architecture\ was\ analyzed\ using\ the\ SMART\ database\ (http://smart.embl.de/,\ accessed\ March\ 2024)."}

Modern scientific practice requires that data and code be made available
to enable reproducibility. The \textbf{Data availability} section
provides multiple layers of access: the raw data in appropriate domain
repositories (NCBI), as well as the code and intermediate files on
GitHub for active development. It is worth noting the limitations of
GitHub (which can change or disappear, and the importance of archiving a
snapshot or release on Zenodo, which will generate a permanent DOI for
long-term citation.

Different types of data should be deposited in appropriate
domain-specific repositories. NCBI SRA and GEO are appropriate for
high-throughput sequencing data and microarray data. GenBank handles
individual sequences and annotations. Dryad and Zenodo serve as general
research data repositories across disciplines. Figshare hosts datasets,
figures, and supplementary materials. GitHub and GitLab are suitable for
code, scripts, and small datasets under 100 MB. Institutional
repositories like Auburn's \href{http://aurora.auburn.edu/repo/}{AUrora}
can handle large datasets that exceed GitHub's limitations.

For GitHub repositories, best practices include providing a
comprehensive README with installation instructions, documenting
software dependencies and version requirements, including example data
and expected outputs, and using meaningful commit messages and version
tags.

\subsection{Balancing Main Text and Supplementary
Methods}\label{balancing-main-text-and-supplementary-methods}

Deciding what goes in the main methods section versus supplementary
materials requires balancing readability with completeness. The goal is
to provide sufficient detail in the main text for readers to understand
your approach while reserving extensive \textbf{technical details for
supplements}. The main text methods should include an overview of the
experimental and analytical workflow, key software and versions for
major analytical steps, sample sizes and experimental design,
statistical approaches and significance thresholds, novel or modified
protocols, and justifications for methodological choices.

\begin{tcolorbox}[enhanced jigsaw, rightrule=.15mm, leftrule=.75mm, colbacktitle=quarto-callout-note-color!10!white, colframe=quarto-callout-note-color-frame, toprule=.15mm, opacitybacktitle=0.6, title=\textcolor{quarto-callout-note-color}{\faInfo}\hspace{0.5em}{Figure 7.3 - Main text → Supplement → GitHub}, left=2mm, colback=white, bottomtitle=1mm, toptitle=1mm, bottomrule=.15mm, coltitle=black, titlerule=0mm, breakable, opacityback=0, arc=.35mm]

\begin{figure}[H]

\centering{

\includegraphics[width=0.75\linewidth,height=\textheight,keepaspectratio]{images/Figure_7-3_Detail_Strategies.jpg}

}

\caption{\label{fig-detail}Continuum of methods detail from brief
main-text description, to more detailed supplement, to most detailed
external documentation (e.g., GitHub).}

\end{figure}%

This chapter assumes a tiered approach to documentation where the main
text provides a readable overview, the supplement adds enough detail for
reproduction, and external resources such as GitHub hold the most
granular, evolving record.

\end{tcolorbox}

\subsection{Comparing Your Methods to Published
Studies}\label{comparing-your-methods-to-published-studies}

When reanalyzing published data or using established methods, comparing
your results to previous studies provides validation and helps identify
potential issues. Choose comparison metrics that make sense for your
analysis. The Bogan et al.~(2025) study provides an excellent example
when they validated their AFP identification approach:
\texttt{"To\ determine\ if\ multiple\ AFP\ translocations\ within\ a\ species\ arose\ from\ misassembly,\ translocated\ AFP\ III\ copies\ outside\ of\ the\ ancestral\ syntenic\ AFP\ III\ region\ were\ manually\ reviewed\ for\ evidence\ of\ misassembly.\ This\ was\ achieved\ by\ aligning\ translocated\ AFP\ III\ regions\ to\ (i)\ a\ known,\ characterized\ AFP\ translocation\ in\ A.\ lupus\ (Desjardins\ et\ al.\ 2012)\ and\ (ii)\ additional\ regions\ of\ translocation\ within\ single\ species\ if\ they\ possessed\ multiple\ translocation\ events\ to\ identify\ artificial\ duplications."}

When differences exist between your results and previous studies,
provide potential explanations. For example, you might note:
\texttt{"Re-analysis\ of\ the\ Smith\ et\ al.\ (2020)\ dataset\ using\ our\ updated\ pipeline\ identified\ 1,247,892\ SNPs,\ compared\ to\ 1,238,456\ reported\ in\ the\ original\ study\ (0.8\%\ difference).\ This\ minor\ discrepancy\ likely\ reflects\ differences\ in\ variant\ caller\ versions\ (GATK\ 4.2\ vs.\ 3.8)\ and\ updated\ reference\ genome\ annotation\ (GRCh38.p13\ vs.\ GRCh38.p7)."}

\subsection{Detailed Components of Bioinformatics
Methods}\label{detailed-components-of-bioinformatics-methods}

\subsubsection{Ethics Statements and Institutional
Approvals}\label{ethics-statements-and-institutional-approvals}

Research involving human subjects or vertebrate animals requires
institutional review and approval before any data collection begins.
This is not merely a formality---it is a legal and ethical requirement
for publication in reputable journals. Approval must be obtained
\textbf{before starting} any research activities. You must include the
specific protocol number from your institution, describe measures taken
to ensure welfare or confidentiality, state which institutional review
board granted approval, and confirm compliance with relevant regulations
such as HIPAA or IACUC guidelines.

For human subjects research, an appropriate statement might read:
\texttt{"This\ study\ was\ approved\ by\ the\ Auburn\ University\ Institutional\ Review\ Board\ (Protocol\ \#21-123\ EP\ 2104).\ All\ participants\ provided\ written\ informed\ consent.\ Patient\ data\ were\ de-identified\ prior\ to\ analysis\ to\ ensure\ confidentiality\ in\ accordance\ with\ HIPAA\ regulations."}
For animal research:
\texttt{"All\ animal\ procedures\ were\ approved\ by\ the\ Auburn\ University\ Institutional\ Animal\ Care\ and\ Use\ Committee\ (IACUC\ Protocol\ \#2021-3456)\ and\ conducted\ in\ accordance\ with\ the\ Guide\ for\ the\ Care\ and\ Use\ of\ Laboratory\ Animals."}

\subsubsection{Data Sources and
Collection}\label{data-sources-and-collection}

Clearly documenting where your data came from and how it was generated
is fundamental to reproducibility. This applies whether you generated
the data yourself or obtained it from public repositories. For newly
generated data, you should document the source of biological material
including species, strain, and tissue type; explain why this source was
selected with supporting literature citations; describe collection
methods and protocols; detail husbandry or growth conditions; specify
sample size and biological replicates; and note age, developmental
stage, or other relevant characteristics.

When using previously published or public data, provide the repository
name and URL, accession numbers for each dataset, version number or
access date, citation to the original publication, and a brief
description of how data were originally generated. The authors describe
their use of existing genomic resources:
\texttt{"Publicly-available\ long-read\ genomes\ used\ in\ this\ study\ included\ Cebidichthys\ violaceus\ (GCA\_023349555.1),\ Gasterosteus\ aculeatus\ (GCA\_016920845.1),\ Leptoclinus\ maculatus\ (GCA\_032191485.1)..."}
providing both the species names and their NCBI accession numbers.

For environmental or oceanographic data, similar principles apply. For
example, in Bogan et al 2025, they describe their environmental data
sources:
\texttt{"Depth-informed\ temperatures\ were\ mined\ from\ the\ Copernicus\ Global\ Ocean\ Reanalysis\ (GLORYS)\ database\ (Jean-Michel\ et\ al.\ 2021)\ within\ ±1°\ latitude/longitude\ of\ species\textquotesingle{}\ coordinate\ of\ maximum\ likelihood\ of\ occurrence\ and\ the\ expected\ depth\ range\ of\ each\ species\ spanning\ a\ 10-year\ period\ between\ June\ 1,\ 2011\ and\ June\ 1,\ 2021."}\textsuperscript{3}

For wet-lab procedures, provide enough detail for replication while
avoiding excessive minutiae. Reference established protocols when
possible, but document any modifications. For nucleic acid extraction,
specify the extraction method or kit with manufacturer and catalog
number, any modifications to the manufacturer's protocol, quality
control measures such as NanoDrop or Bioanalyzer assessments, and
quality thresholds for sample inclusion. When describing library
preparation or other molecular procedures, include manufacturer
information, catalog numbers, and specific protocol modifications. Cite
supporting literature that justifies your choice of tissue, protocol, or
approach for your specific research goals.

Document the sequencing technology, instrument, and parameters used, as
this information affects data quality and interpretation. You should
specify the platform or technology (Illumina, PacBio, Oxford Nanopore),
the specific instrument model, the sequencing center or facility, read
type (single-end versus paired-end), read length, and target coverage
depth or number of reads per sample.

\subsubsection{Bioinformatics Analysis
Workflow}\label{bioinformatics-analysis-workflow}

This section is often the most detailed for bioinformatics papers.
Organize chronologically and use subsections to improve readability.
Statistical analyses must be described with sufficient detail for
replication, including specific tests used, multiple testing
corrections, and significance thresholds.

When methods are complex, use subsections organized by
topic\textsuperscript{111}. Material within each subsection should be
ordered from most to least important while maintaining chronological
flow. For example, the Bogan et al.~(2025) study uses clear subsections
including Genome Assembly, AFP III and sasa/b Annotation, Environmental
Data, Species Tree Construction and Phylogenetic Modeling, Analyses of
AFP III Gene Family Evolution and Synteny, and Models of Selection on
AFP III Copy Number and Sequence.

Within each subsection, they maintain a logical flow. For example, in
the genome assembly section, they proceed from specimen collection to
DNA extraction to library preparation to sequencing to assembly to
quality assessment. This chronological organization helps readers follow
the workflow naturally.

\subsection{Detail Level and Supplementary
Materials}\label{detail-level-and-supplementary-materials}

Determining how much detail to provide and where to place it represents
one of the most challenging aspects of methods writing. The goal is to
ensure \textbf{reproducibility while maintaining readability}. Journal
page limits often necessitate supplementary materials, but their use
must be strategic. Excessive supplementary material can undermine the
scientific process by making comprehensive review difficult or
impossible\textsuperscript{112}.

The main text should include the experimental design and rationale,
sample sizes with justification, key software and databases with
versions for all major steps, novel or modified protocols in sufficient
detail, statistical methods and significance thresholds, justifications
for methodological choices, and clear references to supplementary
methods when appropriate.

Supplementary methods should extend and detail the main text, not
replace it. They are appropriate for technical details that would
interrupt the flow of the main narrative, such as complete command-line
parameters for all software, detailed protocols for standard procedures,
extended descriptions of quality control procedures, alternative
analyses or sensitivity tests, custom scripts or detailed algorithm
descriptions, comprehensive lists like all sample accession numbers, and
detailed mathematical derivations or model specifications.

There are several critical problems with how supplementary materials are
currently used in scientific publishing\textsuperscript{112}. One of the
most serious issues is that references cited only in supplementary
materials are not tracked by major citation indices. This means software
developers and methodologists whose work is cited only in supplements do
not receive proper credit for their contributions.

To address this citation problem, include all \textbf{critical software
citations in the main text} reference list. If supplements cite many
tools, summarize key citations in the main text. Some journals now offer
supplementary methods sections that appear in the main PDF. Follow
journal policies on supplementary citations when they exist, though many
journals lack clear guidance. As researchers, we should also advocate
for better citation tracking of supplementary materials.

Extensive supplementary materials are often reviewed superficially or
not at all with some supplements missing critical element, suggesting
inadequate review\textsuperscript{112}. This gap threatens scientific
validity when critical methodological details are relegated to
unreviewed supplements. To ensure supplements are reviewable, keep them
focused and well-organized, use clear headings and cross-references from
main text, limit supplement length to what can reasonably be reviewed,
highlight critical methodological details in cover letters to editors,
and consider whether extremely long supplements indicate problems with
study design.

Supplementary materials are often poorly integrated with main text,
making it difficult for readers to find specific information. Pop and
Salzberg (2015) describe ``grab-bag'' supplements where vague references
like \texttt{"see\ Supplementary\ Material"} point to hundreds of pages
of unorganized content\textsuperscript{112}. Improve integration by
using specific cross-references such as
\texttt{"Supplementary\ Methods\ Section\ 2.3,"} creating a
supplementary table of contents, using hyperlinks in digital versions,
organizing supplements with clear section headings, and ensuring
supplement numbering aligns with main text references.

Many journals provide supplementary tables only as PDFs, making data
extraction difficult. Large image files may be compressed, reducing
quality. These format restrictions undermine the utility of
supplementary materials. Follow format best practices by providing
tabular data in machine-readable formats like CSV, TSV, or Excel;
including PDF versions for readability alongside data files; using
appropriate image resolution (300+ DPI for publication quality);
compressing images intelligently (PNG for line art, JPEG for photos);
and including README files explaining file formats and contents.

\subsubsection{Solutions and Best
Practices}\label{solutions-and-best-practices}

While systemic problems with supplementary materials require journal
policy changes, individual researchers can adopt practices that improve
reproducibility and accessibility. Some journals have addressed
supplementary material issues better than others. When choosing
journals, consider those with no arbitrary page limits for online
versions, online methods sections, integrated supplementary materials in
main PDF, support for interactive figures and data, clear policies on
supplementary citations, and required data and code deposition.
Electronic publishing enables features that improve methods
accessibility. Some journals like PNAS provide interactive online
articles with hover-over citations, hyperlinked figures, and integrated
supplements.

GitHub provides an excellent platform for detailed methods
documentation. GitHub's version control and collaborative features make
it ``like a laboratory notebook for scientific
computing''\textsuperscript{113}. A GitHub repository can contain
complete analysis workflows with documentation that exceeds what
journals allow. For comprehensive methods documentation on GitHub, host
complete executable analysis pipelines, provide detailed README files
with step-by-step instructions, include example data and expected
outputs, document software dependencies and versions, enable community
feedback and error reporting via Issues, and version-control methods
updates and corrections.

\section{Conclusion Part 2}\label{conclusion-part-2}

Writing effective methods sections for bioinformatics research requires
attention to detail, awareness of reproducibility standards, and
strategic use of available publishing tools. As computational approaches
become increasingly central to biological research, clear methods
documentation becomes ever more critical. Remember that the methods
section serves multiple audiences: reviewers judging your work's
validity, researchers attempting to reproduce your results, and future
scientists building upon your findings. Each of these audiences requires
complete, accessible information.

By following the principles outlined in this guide---providing
sufficient detail, documenting software versions, justifying
methodological choices, and making data and code accessible---you
contribute to a more transparent and reproducible scientific enterprise.
The extra effort invested in comprehensive methods documentation pays
dividends in the impact and longevity of your research contributions. As
demonstrated by the antifreeze protein study discussed throughout this
guide, high-quality methods documentation not only enables others to
validate and build upon your work but also showcases the rigor and
thoughtfulness of your research approach.

The evolution of scientific publishing continues to create new
opportunities for methods transparency. Long-read sequencing
technologies, version-controlled code repositories, and interactive
online publications are transforming how we document and share
computational methods. By embracing these tools and maintaining rigorous
documentation standards, bioinformatics researchers can ensure their
work contributes to a cumulative, reproducible body of scientific
knowledge.

\section{Practice Problem Sets}\label{practice-problem-sets}

\begin{enumerate}
\def\labelenumi{\arabic{enumi}.}
\tightlist
\item
  \textbf{Repository Analysis}: Examine the following two GitHub repos
  from my lab for data organization principles (e.g., directory
  structure, READMEs, file naming). Note 2-3 strengths (e.g.,
  chronological folders) and 2 improvements (e.g., consistent naming for
  automation). How would you adapt their structures for your semester
  project's SNP analysis, including relative paths and a root README?
\end{enumerate}

\begin{itemize}
\tightlist
\item
  \href{https://github.com/StevisonLab/Arctoides-Hybridization}{Primate
  Genital Evolution Project}
\item
  \href{https://github.com/StevisonLab/Diet-and-recombination-in-Dmel}{Recombination
  Manipulation in Flies}
\end{itemize}

\begin{enumerate}
\def\labelenumi{\arabic{enumi}.}
\setcounter{enumi}{1}
\item
  \textbf{Data Availability Review}: Review the Bogan et al.~(2025) Data
  Availability statement and their GitHub repo Long\_AFP
  (https://github.com/snbogan/Long\_AFP). Identify how it exemplifies
  methods reproducibility (e.g., code versioning, intermediate files).
  Critique one strength and one gap per the standards introduced here.
\item
  \textbf{Project Setup Exercise}: In advance of Step 3 of your group
  project, practice: Create a personal GitHub repo (e.g.,
  username/snp-mock-analysis). Add directories (data/, scripts/,
  results/, docs/), populate data/README.md with mock FASTQ provenance
  (URLs, checksums, dates), scripts/01-fetch-data.sh with relative
  paths, and root README.md outlining objectives. Commit descriptively,
  push, and share the link. Explain one automation win (e.g., for f in
  data/*.fastq.gz; do \ldots).
\item
  \textbf{Methods Writing Practice}: Read Bogan et al.~(2025) methods
  section (focus on haplotype phasing, assembly). Reflect on a prior
  script (class/research): (a) How often is it called by other code? (b)
  If wrong, how does it impact results? (c) How noticeable are errors?
  Rewrite its usage as a 150-word methods paragraph, citing
  tool/version/parameters (e.g., ``BWA-MEM v0.7.17 Li 2009, -T 30''),
  justifying choices, and noting deviations from defaults. Add to
  docs/methods.md in your mock repo from \#3.
\item
  \textbf{Kanban and Version Control Workflow}: In your mock repo (\#3),
  create a GitHub Project Kanban board with columns (To Do, In Progress,
  Done). Add cards: ``Download ref genome + checksums'', ``Write QC
  script'', ``Document params''. Branch for ``feature-qc'', implement a
  simple FastQC wrapper script (log version/params), PR it with review
  comments, merge. Move cards accordingly. Document in
  results/YYYY-MM-DD-kanban/README.md how this prevents ``v2final.py''
  chaos and aids collaboration.
\end{enumerate}

\section{Reflection Questions}\label{reflection-questions-6}

\begin{enumerate}
\def\labelenumi{\arabic{enumi}.}
\item
  \textbf{AntiFreeze Paper Methods}: Reread the methods section and any
  methods‑related supplementary material for the Bogan et
  al.~study\textsuperscript{3}. As you read, sketch the implicit
  workflow using the categories introduced in this chapter: data
  acquisition, reference genome preparation, read alignment, variant
  calling, variant filtering, functional annotation, and downstream
  statistical analyses. For each category, note any missing information.
  Does the paper specify the reference genome accession and version? Are
  the alignment and variant calling tools named, and are versions
  indicated? Are there places where it is unclear whether default
  parameters were used, or where non‑default settings are mentioned
  without justification?
\item
  \textbf{Software Justifcation}: The tool \texttt{Exonerate} used in
  the\textsuperscript{3} paper is similar to GeneWise, which you are
  using for your manual annotation of PRDM9. The primary difference is
  that Exonerate can be run on the command line and GeneWise is
  restricted to the web.
\end{enumerate}

\begin{itemize}
\tightlist
\item
  Why do you think the authors of the antifreeze paper chose to use
  Exonerate over GeneWise?
\item
  Why do you think we are using GeneWise in your manual annotation
  project instead of Exonerate?
\end{itemize}

\begin{enumerate}
\def\labelenumi{\arabic{enumi}.}
\setcounter{enumi}{2}
\item
  How does adopting chronological directory structures (e.g., YYYY-MM-DD
  subfolders) change your approach to debugging failed analyses months
  later, compared to purely logical organization?
\item
  Why is documenting software versions and exact command lines in README
  files more critical for bioinformatics than traditional wet-lab
  methods, and how does this tie into reproducibility crises in
  computational biology?
\item
  When balancing main text vs.~supplementary methods, how might vague
  citations like ``processed as in Smith et al.'' undermine peer review,
  and what specific elements from this chapter would you prioritize in
  your own papers?
\end{enumerate}

\section{References}\label{references-7}

\chapter{Alternative Post‑Alignment
Workflows}\label{alternative-postalignment-workflows}

\section{Introduction}\label{sec-ch08-intro}

The previous chapters have focused on variant calling workflows using
tools like GATK, providing a foundation for understanding how sequencing
reads are processed, aligned, and analyzed to identify genomic variants.
However, the power of next-generation sequencing extends far beyond
simple variant discovery. The same fundamental technologies---Illumina
short-read sequencing, PacBio long-read sequencing, and Oxford Nanopore
sequencing---can be deployed to answer diverse biological questions
across multiple scales of organization\textsuperscript{10}. In this
chapter, we explore alternative post-alignment workflows that leverage
these sequencing technologies for applications ranging from population
genetics to transcriptomics, epigenomics, metagenomics, and integrative
multi-omic analyses.

As introduced in Chapter 1 (Figure~\ref{fig-genomics-timeline}), the
history of genomics has been marked by technological innovation that has
progressively expanded our capacity to interrogate biological systems.
The completion of the Human Genome Project in 2003 was followed by the
development of next-generation sequencing platforms that dramatically
reduced both the cost and time required for genome
sequencing\textsuperscript{10}. This democratization of sequencing
technology has enabled researchers to apply genome sequencing to address
five key perspectives that shape modern genomics research.

\section{Five Perspectives on Genome
Sequencing}\label{sec-five-perspectives}

Pevsner outlines five complementary perspectives that frame how we
approach genome sequencing projects\textsuperscript{10}. These
perspectives provide a conceptual scaffold for understanding the diverse
applications of sequencing technologies and serve as organizing
principles for the remainder of this chapter.

The first perspective involves cataloging genomic information at the
most fundamental level. This includes basic features such as genome
size, chromosome number, GC content, gene count, and the identification
of repetitive elements. Genome browsers and annotation pipelines
represent the primary tools for organizing and presenting these
catalogs. The reference genomes maintained by consortia like the Genome
Reference Consortium provide the foundation for comparative analyses
across individuals and populations.

The second perspective emphasizes comparative genomics, recognizing that
our understanding of any single genome is dramatically enhanced through
comparison with related genomes. Questions about orthology, synteny, and
the timing of speciation events require whole-genome alignments and
phylogenetic reconstruction. Tools like the UCSC Genome Browser and
Ensembl Genomes facilitate these comparisons by providing pre-computed
alignments and evolutionary annotations.

The third perspective focuses on biological principles, asking how
genome content and organization serve the functions of organisms and how
evolutionary forces shape genome architecture. This includes
consideration of gene birth and death, polyploidization, positive and
negative selection, and the role of epigenetic regulation. Population
genetic analyses using tools like BLAST and molecular phylogenetic
methods help address these fundamental questions.

The fourth perspective centers on human disease relevance, examining how
genomic variation contributes to disease susceptibility and how
pathogens exploit host biology. Single nucleotide polymorphisms, copy
number variants, and structural variants are analyzed through linkage
studies, genome-wide association studies, and functional genomics
approaches to identify disease-causing mutations and therapeutic
targets.

The fifth perspective addresses bioinformatics infrastructure, including
the databases, software tools, and computational workflows required to
analyze and visualize genomic data. The explosion of sequencing data has
necessitated development of specialized file formats, compression
algorithms, and cloud computing platforms to manage and distribute
genomic information at scale.

These five perspectives are not mutually exclusive but rather represent
complementary lenses through which we view genomic data. The alternative
workflows discussed in this chapter draw upon multiple perspectives
simultaneously, reflecting the integrative nature of modern genomics
research.

\begin{tcolorbox}[enhanced jigsaw, rightrule=.15mm, leftrule=.75mm, colbacktitle=quarto-callout-note-color!10!white, colframe=quarto-callout-note-color-frame, toprule=.15mm, opacitybacktitle=0.6, title=\textcolor{quarto-callout-note-color}{\faInfo}\hspace{0.5em}{Lab Connections}, left=2mm, colback=white, bottomtitle=1mm, toptitle=1mm, bottomrule=.15mm, coltitle=black, titlerule=0mm, breakable, opacityback=0, arc=.35mm]

Look back over the assignment you submitted for \href{Lab02.qmd}{Lab 2}
where you framed a genome sequencing project based on these
perspectives. While you may have previously only considered a
traditional NGS workflow, now consider the alternative post-alignment
workflows and other sequencing platforms introduced below. What might
you consider changing if you were to re-submit this narrative? How might
some you improve your ability to test your hypotheses by adding another
sequencing strategy?

\end{tcolorbox}

\section{DNA Sequencing Applications}\label{sec-dna-sequencing}

DNA sequencing workflows that go beyond standard variant calling
encompass several major application areas, including population
genetics, ancient DNA analysis, and whole-genome resequencing projects.
Each application presents unique computational challenges and biological
insights.

\subsection{Population Genetics and Resequencing
Projects}\label{sec-population-genetics}

Population genetics leverages DNA sequencing to understand allele
frequency distributions, genetic structure, linkage disequilibrium
patterns, and demographic history across populations. Whole-genome
resequencing of multiple individuals from the same or different
populations enables inference of selection pressures, migration
patterns, and population size changes over evolutionary time. The
computational workflow typically involves variant calling across all
individuals simultaneously, followed by population genetic analyses that
estimate parameters such as nucleotide diversity and estimating natural
selection\textsuperscript{114}. This topic will be covered in more depth
in Chapter 12.

\begin{tcolorbox}[enhanced jigsaw, rightrule=.15mm, leftrule=.75mm, colbacktitle=quarto-callout-note-color!10!white, colframe=quarto-callout-note-color-frame, toprule=.15mm, opacitybacktitle=0.6, title=\textcolor{quarto-callout-note-color}{\faInfo}\hspace{0.5em}{Case Study: The 1000 Genomes Project}, left=2mm, colback=white, bottomtitle=1mm, toptitle=1mm, bottomrule=.15mm, coltitle=black, titlerule=0mm, breakable, opacityback=0, arc=.35mm]

The 1000 Genomes Project, initiated in 2008 and completed in 2015,
represents a landmark achievement in population
genomics\textsuperscript{115}. This international consortium sequenced
the genomes of over 2,500 individuals from 26 populations across Africa,
Europe, Asia, and the Americas. The project's primary goal was to create
a comprehensive catalog of human genetic variation, discovering more
than 95\% of variants with minor allele frequencies as low as 1\% across
the genome.

The project employed a multi-phase strategy beginning with three pilot
studies to optimize sequencing depth and sample selection. The first
pilot sequenced 180 individuals at low coverage (2×), the second
sequenced two parent-offspring trios at high coverage (20×), and the
third focused on deep sequencing of exons in 1,000 genes. This design
allowed the consortium to balance the trade-offs between sequencing
depth and the number of individuals sampled.

The final dataset comprises over 88 million variants, including single
nucleotide polymorphisms, short insertions and deletions, and structural
variants. The project revealed that each individual carries
approximately 4-5 million variants compared to the reference genome,
with the vast majority being rare variants found at low frequency in the
population. Many rare variations were restricted to closely related
geographic groups, highlighting the importance of sampling diverse
populations to capture the full spectrum of human genetic variation.

The 1000 Genomes data has become an indispensable resource for
genome-wide association studies, enabling researchers to impute untyped
variants and increase statistical power for disease mapping. The
International Genome Sample Resource continues to host and expand upon
the dataset, making it freely accessible to researchers worldwide.

\end{tcolorbox}

The statistical methods developed to analyze 1000 Genomes data---joint
genotyping, imputation panels, and phasing algorithms---have become
standard tools in population genetics. The project's design trade-offs
between breadth (number of individuals) and depth (coverage per
individual) continue to inform contemporary study designs, with the
recognition that low-coverage sequencing of many individuals can
efficiently capture common and moderately rare variation when combined
with sophisticated imputation methods\textsuperscript{116}.

Population resequencing projects require careful consideration of
variant calling strategies to minimize false positives while maintaining
sensitivity for rare variants. Joint calling approaches, where all
samples are genotyped simultaneously, help distinguish true variants
from sequencing artifacts by leveraging information across the entire
cohort. Phasing algorithms such as BEAGLE and SHAPEIT reconstruct
haplotypes from genotype data, enabling analysis of linkage
disequilibrium and inference of recombination history.

\subsection{Ancient DNA Analysis}\label{sec-ancient-dna}

Ancient DNA (aDNA) analysis represents a specialized application of DNA
sequencing that requires modified computational and laboratory protocols
to address the unique challenges of degraded genetic material. DNA
molecules recovered from archaeological and paleontological specimens
are typically fragmented, chemically modified through deamination of
cytosine bases, and contaminated with DNA from environmental
microorganisms and modern humans. These characteristics necessitate
specialized bioinformatic pipelines that account for increased error
rates, short read lengths, and the presence of post-mortem damage
patterns\textsuperscript{117}.

The aDNA workflow begins with careful authentication procedures to
verify that recovered DNA sequences are genuinely ancient rather than
modern contaminants. Damage patterns, particularly C-to-T transitions at
the 5' ends of reads resulting from cytosine deamination, serve as
molecular signatures of authenticity. Tools like mapDamage quantify
these damage patterns and can be used to trim or correct affected bases
during variant calling. Low sequencing coverage typical of ancient
samples requires statistical models that account for uncertainty in
genotype calls, often using genotype likelihoods rather than called
genotypes for downstream population genetic analyses.

\begin{tcolorbox}[enhanced jigsaw, rightrule=.15mm, leftrule=.75mm, colbacktitle=quarto-callout-note-color!10!white, colframe=quarto-callout-note-color-frame, toprule=.15mm, opacitybacktitle=0.6, title=\textcolor{quarto-callout-note-color}{\faInfo}\hspace{0.5em}{Case Study: Ancient Hominins and Ötzi the Iceman}, left=2mm, colback=white, bottomtitle=1mm, toptitle=1mm, bottomrule=.15mm, coltitle=black, titlerule=0mm, breakable, opacityback=0, arc=.35mm]

The sequencing of archaic hominin genomes has revolutionized our
understanding of human evolution. High-coverage genomes from
Neanderthals and Denisovans revealed that these extinct human lineages
interbred with anatomically modern humans, with non-African populations
carrying 1-4\% Neanderthal ancestry and some Oceanian populations
carrying Denisovan ancestry. These findings fundamentally changed our
view of human origins from a simple ``out of Africa'' replacement model
to a more complex picture involving admixture between divergent
lineages\textsuperscript{10}.

Ötzi the Iceman, a naturally mummified human discovered in the Italian
Alps in 1991, provides a different window into ancient
genomics\textsuperscript{118}. The 5,300-year-old individual represents
one of the oldest complete human genomes sequenced to date. Initial
genome sequencing in 2012 used DNA extracted from the inner part of
Ötzi's femur to minimize modern contamination, obtaining approximately
20 nanograms of genomic DNA---hundreds of times less than typically used
for modern whole-genome sequencing. Despite these challenges,
researchers identified over 2 million sequence variants and traced
Ötzi's ancestry to populations from the Mediterranean region.

More recent reanalysis using improved sequencing methods and comparison
with a growing database of over 10,000 prehistoric European genomes has
refined our understanding of Ötzi's genetic makeup. The updated analysis
reveals that he showed unusually high early Neolithic farmer ancestry
and provides new insights into his physical appearance, including darker
skin pigmentation than previously estimated. The ability to extract and
sequence DNA from Ötzi's gut contents has also enabled reconstruction of
his diet and the microbial communities in his digestive system,
demonstrating the potential for ancient DNA analysis to provide
multifaceted insights into past human populations.

These ancient DNA studies illustrate how degraded genetic material can
be successfully recovered and analyzed when appropriate computational
methods account for the unique characteristics of ancient samples,
including fragmentation, chemical damage, and contamination.

\end{tcolorbox}

The insights gained from ancient DNA extend beyond human evolution to
include extinct megafauna, ancient pathogens, and the domestication
history of crops and livestock. Each application requires tailored
approaches to address the specific preservation conditions and
biological questions of interest.

\subsection{Copy Number Variation and Structural Variant
Analysis}\label{sec-cnv-analysis}

Copy number variants (CNVs) and larger structural variants represent an
important class of genomic variation that is incompletely captured by
standard SNP-based variant calling\textsuperscript{119,120}. These
variants, which include deletions, duplications, inversions, and
translocations, can have profound effects on gene dosage and regulation.
Detection of CNVs from short-read sequencing data relies on multiple
orthogonal signals: read depth, which reflects the number of copies
present; paired-end mapping, which identifies discordant read pairs
spanning breakpoints; split reads, which span breakpoints at base-pair
resolution; and de novo assembly, which can reconstruct variant
haplotypes directly. The diversity of structural variant classes
(Figure~\ref{fig-structural-variation}) reflects the multiple molecular
mechanisms that reshape genome architecture. While deletions and
duplications alter gene dosage through copy number changes, inversions
and translocations rearrange gene order without necessarily changing
copy number, and mobile element insertions introduce new sequence
content. Each variant class leaves distinct signatures in sequencing
data that computational tools exploit for detection.

\begin{tcolorbox}[enhanced jigsaw, rightrule=.15mm, leftrule=.75mm, colbacktitle=quarto-callout-note-color!10!white, colframe=quarto-callout-note-color-frame, toprule=.15mm, opacitybacktitle=0.6, title=\textcolor{quarto-callout-note-color}{\faInfo}\hspace{0.5em}{Figure 8.1 -- Classes of genome structural variation}, left=2mm, colback=white, bottomtitle=1mm, toptitle=1mm, bottomrule=.15mm, coltitle=black, titlerule=0mm, breakable, opacityback=0, arc=.35mm]

\begin{figure}[H]

\centering{

\includegraphics[width=0.9\linewidth,height=\textheight,keepaspectratio]{images/Figure_8.1_SVs.png}

}

\caption{\label{fig-structural-variation}Genome structural variation
schematics}

\end{figure}%

This figure summarizes the major classes of structural variants that can
be detected in whole‑genome sequencing data. The top row illustrates
events that change copy number within a chromosome: deletions remove a
contiguous block of sequence, tandem duplications create an immediately
adjacent copy of a segment, and interspersed duplications insert the
copied segment at a distant location. The middle row highlights events
that introduce new sequence or move segments between chromosomes,
including novel sequence insertions, balanced translocations that
exchange segments between loci, and inter‑chromosome insertions in which
a segment from one chromosome is inserted into another. The bottom row
shows additional mechanisms that remodel genome structure:
mobile‑element insertions add transposable elements at new sites,
inversions flip the orientation of an internal segment, and
intra‑chromosome insertions move a block of sequence to a different
position on the same chromosome. Together, these categories provide a
conceptual framework for thinking about how genome architecture can
change and how different algorithms leverage read depth, split reads,
and discordant pair mappings to discover each class of event in
sequencing‑based analyses. Created in \url{https://BioRender.com}.

\end{tcolorbox}

\begin{tcolorbox}[enhanced jigsaw, rightrule=.15mm, leftrule=.75mm, colbacktitle=quarto-callout-note-color!10!white, colframe=quarto-callout-note-color-frame, toprule=.15mm, opacitybacktitle=0.6, title=\textcolor{quarto-callout-note-color}{\faInfo}\hspace{0.5em}{Case Study: The Human Amylase Locus}, left=2mm, colback=white, bottomtitle=1mm, toptitle=1mm, bottomrule=.15mm, coltitle=black, titlerule=0mm, breakable, opacityback=0, arc=.35mm]

Recent work on the human amylase locus exemplifies the power of
long-read sequencing for resolving complex structural variation. The
amylase gene family includes salivary amylase (AMY1) and pancreatic
amylase (AMY2A and AMY2B) genes that encode enzymes responsible for
breaking down dietary starch. Previous studies using short-read
sequencing identified substantial copy number variation in AMY1, with
individuals carrying between 2 and 17 copies per diploid genome, and
suggested that this variation might represent an adaptive response to
agricultural diets high in starch.

However, the amylase locus comprises highly similar segmental
duplications with greater than 99\% sequence identity, making accurate
assembly from short reads nearly impossible. The amylase locus exhibits
multiple structural variant classes simultaneously
(Figure~\ref{fig-structural-variation}): tandem duplications create
arrays of AMY1 genes, interspersed duplications move segments to distant
locations, and inversions flip the orientation of entire haplotype
blocks. A recent study combined PacBio long-read sequencing and optical
genome mapping to reconstruct the amylase locus at nucleotide resolution
in 98 individuals, identifying 30 distinct structural
haplotypes\textsuperscript{121}. This high-resolution map revealed that
the coding sequences of AMY1 copies are evolving under negative
selection to maintain essential enzymatic function. The study identified
two distinct mutational mechanisms---nonallelic homologous recombination
(NAHR) and microhomology-mediated break-induced replication
(MMBIR)---that generate the observed copy number variation.

Phylogenetic analysis of the structural haplotypes revealed that a
common three-copy haplotype dates back approximately 800,000 years,
predating the split between modern humans and Neanderthals. Analysis of
ancient human genomes showed that hunter-gatherers already exhibited
highly variable AMY1 copy numbers as early as 45,000 years ago. However,
haplotypes with more than three AMY1 copies increased significantly in
frequency among European farmers over the past 4,000 years, suggesting
that selection acted on existing copy number variation in response to
agricultural diets.

This case study demonstrates how emerging sequencing technologies and
computational methods are revealing the complex evolutionary history of
structurally variable regions that were previously intractable to
analysis\textsuperscript{122}. It also illustrates the importance of
combining population-scale data with ancient DNA to distinguish between
the initial generation of variation and subsequent selection on that
variation.

\end{tcolorbox}

Third-generation long-read sequencing from PacBio and Oxford Nanopore
has dramatically improved our ability to detect and characterize
structural variants. These platforms can directly sequence through
repetitive regions that confound short-read approaches and can phase
variants over long distances to reconstruct complete haplotypes. Optical
genome mapping provides an orthogonal technology that labels specific
sequence motifs, allowing visualization of large-scale genome structure
and validation of sequence-based assemblies.

\section{RNA Sequencing Applications}\label{sec-rna-sequencing}

RNA sequencing has revolutionized transcriptomics by enabling
comprehensive measurement of gene expression with single-nucleotide
resolution. Unlike DNA sequencing, which interrogates the static genomic
template, RNA-seq captures the dynamic transcriptome, reflecting which
genes are active in a given cell type, developmental stage, or
environmental condition. The computational workflow for RNA-seq differs
substantially from DNA sequencing, requiring specialized tools to handle
spliced alignments, quantify transcript abundance, and identify
differentially expressed genes\textsuperscript{123}.

\subsection{Transcriptome Analysis and Differential
Expression}\label{sec-transcriptomics}

The standard RNA-seq workflow begins with quality control of sequencing
reads, followed by alignment to a reference genome or transcriptome
using splice-aware aligners such as HISAT2, STAR, or Salmon. Unlike
genomic DNA alignments, RNA-seq alignments must account for splicing,
allowing reads to span exon-exon junctions. The choice between genome
and transcriptome alignment affects downstream quantification, with
transcriptome alignment enabling faster quantification but potentially
missing novel transcripts, while genome alignment allows discovery of
unannotated genes at the cost of increased computational
requirements\textsuperscript{124,125}.

Quantification of gene expression involves counting the number of reads
that map to each gene or transcript. Raw read counts must be normalized
to account for differences in sequencing depth between samples and for
systematic biases such as gene length. Methods like RPKM (reads per
kilobase per million) and TPM (transcripts per million) provide
normalized expression values suitable for within-sample comparisons,
while methods based on scaling factors, such as those implemented in
DESeq2 and edgeR, are preferred for differential expression analysis
between samples.

Differential expression analysis identifies genes whose expression
levels change significantly between experimental conditions. This
typically involves fitting statistical models---often negative binomial
distributions that account for overdispersion common in count data---to
estimate fold changes and test for significant differences. Multiple
testing correction using false discovery rate methods controls for the
large number of hypotheses tested. Downstream analysis includes gene
ontology enrichment to identify biological processes associated with
differentially expressed genes and gene set enrichment analysis to
detect coordinated changes in functionally related gene groups.

RNA-seq data can also be used to identify alternative splicing events,
quantify allele-specific expression, discover novel transcripts, and
characterize non-coding RNAs. Each application requires specialized
computational tools and statistical approaches tailored to the specific
biological question and experimental design. The compgenomr textbook by
Uyar and colleagues provides comprehensive R-based workflows for each
stage of this analysis, from quality assessment through functional
enrichment, with reproducible code examples that parallel the conceptual
framework outlined here.

\subsection{Exome Sequencing for Variant
Discovery}\label{sec-exome-sequencing}

Exome sequencing represents a targeted application of RNA-seq technology
to identify coding variants in clinical and research
contexts\textsuperscript{126}. By enriching for exonic sequences prior
to sequencing, whole-exome sequencing (WES) captures approximately 1-2\%
of the genome where the majority of disease-causing mutations reside.
This targeted approach reduces sequencing costs while maintaining high
coverage of protein-coding regions, making it particularly attractive
for clinical diagnostics and studies of Mendelian
disorders\textsuperscript{127}.

The computational workflow for exome sequencing closely parallels that
for whole-genome sequencing, including alignment, duplicate marking,
base quality score recalibration, and variant calling. However, exome
sequencing presents unique challenges including uneven coverage due to
capture efficiency variation, off-target capture of intronic and
intergenic sequences, and limited ability to detect structural variants
that span capture probe boundaries\textsuperscript{116}. Variant
filtering must account for these technical artifacts while prioritizing
rare, protein-altering variants likely to be pathogenic.

Clinical exome sequencing has proven highly effective for diagnosing
rare genetic disorders, with diagnostic yields typically ranging from
25-50\% depending on the phenotype and patient population.
Interpretation of candidate variants requires integration of multiple
lines of evidence including allele frequency in population databases,
computational predictions of functional impact, evolutionary
conservation, segregation with disease in families, and functional
studies in model systems. Standards and guidelines from organizations
like the American College of Medical Genetics and Genomics provide
frameworks for classifying variants as pathogenic, likely pathogenic,
uncertain significance, likely benign, or benign. Guidelines for
clinical whole genome sequencing interpretation emphasize that variant
classification is an iterative process requiring multidisciplinary
expertise\textsuperscript{126}. As population databases grow and
functional assays improve, variants initially classified as uncertain
significance may be reclassified, highlighting the importance of
periodic reanalysis of unsolved cases

The advent of long-read RNA sequencing using PacBio Iso-Seq and Oxford
Nanopore direct RNA sequencing enables sequencing of full-length
transcripts without assembly, providing unprecedented resolution of
transcript isoforms and enabling detection of RNA modifications directly
from sequencing data. These emerging technologies promise to further
expand the applications of RNA sequencing in the coming years.

\section{Epigenetic Sequencing Strategies}\label{sec-epigenetics}

Epigenetic modifications---heritable changes in gene expression that do
not involve alterations to the DNA sequence itself---play fundamental
roles in development, cellular differentiation, and disease.
Sequencing-based methods to interrogate epigenetic states have become
essential tools for understanding gene regulation and chromatin
organization.

At the molecular level, chromatin structure exists along a continuum
from open, accessible configurations that permit transcription to
condensed, inaccessible states that maintain gene silencing
(Figure~\ref{fig-epigenetics-overview}). Histone acetyltransferases and
methyltransferases catalyze opposing modifications---acetylation of
lysine residues loosens DNA-histone contacts while methylation promotes
chromatin compaction. Sequencing-based approaches leverage these
structural differences to map regulatory landscapes across the genome.

\begin{tcolorbox}[enhanced jigsaw, rightrule=.15mm, leftrule=.75mm, colbacktitle=quarto-callout-note-color!10!white, colframe=quarto-callout-note-color-frame, toprule=.15mm, opacitybacktitle=0.6, title=\textcolor{quarto-callout-note-color}{\faInfo}\hspace{0.5em}{Figure 8.2 -- Epigenetic regulation of chromatin state}, left=2mm, colback=white, bottomtitle=1mm, toptitle=1mm, bottomrule=.15mm, coltitle=black, titlerule=0mm, breakable, opacityback=0, arc=.35mm]

\begin{figure}[H]

\centering{

\includegraphics[width=0.8\linewidth,height=\textheight,keepaspectratio]{images/Figure_8.2_Epigenetics.png}

}

\caption{\label{fig-epigenetics-overview}Epigenetic modifications and
chromatin state}

\end{figure}%

This figure provides an overview of how covalent histone modifications
reshape chromatin structure and influence transcriptional activity. On
the left, an open chromatin fiber is associated with histone
acetylation: histone acetyltransferases (HATs) transfer acetyl groups
from acetyl‑CoA to lysine residues on histone tails, neutralizing their
positive charge and loosening interactions with DNA. The resulting
relaxed nucleosome array permits recruitment of RNA polymerase II and
transcriptional machinery, favoring active gene expression. On the
right, a condensed chromatin fiber is linked to histone methylation: the
PRC2 complex uses S‑adenosylmethionine (SAM) as a methyl donor to
deposit trimethyl marks on specific lysines, promoting the formation of
compact nucleosome structures that exclude RNA polymerase II and
maintain transcriptional repression. By juxtaposing these two states,
the figure emphasizes that epigenetic modifications do not change the
underlying DNA sequence but provide a dynamic, reversible layer of
regulation that helps cells interpret genomic information in a
context‑dependent manner. Created in \url{https://BioRender.com}.

\end{tcolorbox}

\subsection{ChIP-seq and Protein-DNA Interactions}\label{sec-chipseq}

Chromatin immunoprecipitation followed by sequencing (ChIP-seq) enables
genome-wide mapping of protein-DNA interactions, including transcription
factor binding sites and histone modifications. The technique involves
crosslinking proteins to DNA, fragmenting chromatin, immunoprecipitating
specific proteins using antibodies, and sequencing the associated DNA
fragments\textsuperscript{128}. Regions of the genome bound by the
protein of interest show enrichment of sequencing reads compared to
background.

\begin{tcolorbox}[enhanced jigsaw, rightrule=.15mm, leftrule=.75mm, colbacktitle=quarto-callout-note-color!10!white, colframe=quarto-callout-note-color-frame, toprule=.15mm, opacitybacktitle=0.6, title=\textcolor{quarto-callout-note-color}{\faInfo}\hspace{0.5em}{Figure 8.3 -- ChIP‑seq workflow for mapping protein--DNA interactions}, left=2mm, colback=white, bottomtitle=1mm, toptitle=1mm, bottomrule=.15mm, coltitle=black, titlerule=0mm, breakable, opacityback=0, arc=.35mm]

\begin{figure}[H]

\centering{

\includegraphics[width=0.9\linewidth,height=\textheight,keepaspectratio]{images/Fig_8.3_ChipSeq.png}

}

\caption{\label{fig-chipseq-workflow}ChIP-seq experimental workflow}

\end{figure}%

This figure outlines the major laboratory and computational steps of a
ChIP‑seq experiment used to identify genome‑wide binding sites of
transcription factors or histone modifications. In the first steps,
DNA‑binding proteins interact with chromatin in vivo, and cells are
treated with crosslinking agents so that proteins remain covalently
attached to their bound DNA fragments. The chromatin is fragmented and
subjected to immunoprecipitation using an antibody specific to the
protein or histone mark of interest, enriching DNA fragments from
occupied sites while unbound regions are largely washed away. After
reversing crosslinks, the recovered DNA is purified, ligated to
sequencing adapters, and amplified to generate a sequencing library.
Short reads produced by next‑generation sequencing are then aligned to a
reference genome, and regions with significant enrichment of aligned
reads relative to background are called as peaks. The final panel
illustrates these peaks as local maxima over the genome, which can be
linked to nearby genes, used for motif discovery, and integrated with
RNA‑seq or ATAC‑seq data to infer regulatory networks. Created in
\url{https://BioRender.com}.

\end{tcolorbox}

Following library preparation ans sequencing
Figure~\ref{fig-chipseq-workflow}, the computational analysis of
ChIP-seq data begins with alignment of reads to the reference genome,
followed by peak calling to identify regions of significant
enrichment\textsuperscript{129}. Peak-calling algorithms such as MACS2
employ statistical models that account for local background rates and
estimate false discovery rates to distinguish true binding sites from
random fluctuations. The shape of enrichment profiles differs depending
on the protein assayed: sequence-specific transcription factors
typically show sharp, punctate peaks centered on their binding motifs,
while histone modifications often show broader domains spanning multiple
kilobases.

Downstream analysis includes motif discovery to identify the DNA
sequences recognized by transcription factors, peak annotation to
determine the genomic context of binding sites relative to genes and
regulatory elements, and integration with other datasets such as RNA-seq
to correlate binding with gene expression changes. ChIP-seq has revealed
fundamental principles of gene regulation including the combinatorial
logic of transcription factor binding, the chromatin states associated
with active and repressed genes, and the three-dimensional organization
of the genome through methods like Hi-C that map chromatin contacts
genome-wide.

\subsection{DNA Methylation and Bisulfite
Sequencing}\label{sec-bisulfite-seq}

DNA methylation, the addition of methyl groups to cytosine bases,
represents one of the most extensively studied epigenetic modifications
in mammalian genomes. Methylation occurs predominantly at CpG
dinucleotides and plays critical roles in gene silencing, X-chromosome
inactivation, genomic imprinting, and suppression of transposable
elements. Bisulfite sequencing leverages the differential chemical
reactivity of methylated and unmethylated cytosines to enable
single-base resolution mapping of methylation patterns across the
genome\textsuperscript{12,130}.

The bisulfite conversion process treats unmethylated cytosines with
sodium bisulfite, which deaminates them to uracil, while methylated
cytosines remain protected. Following PCR amplification, uracils are
replaced with thymines, creating a permanent record of the methylation
state in the DNA sequence. Alignment of bisulfite-treated reads requires
specialized aligners that account for the C-to-T conversions, and
quantification involves calculating the proportion of reads showing
cytosine versus thymine at each CpG site.

Computational analysis of bisulfite sequencing data identifies
differentially methylated regions between samples, which often correlate
with changes in gene expression\textsuperscript{131}. Methylation
patterns show characteristic signatures at different genomic elements:
CpG islands at gene promoters are typically unmethylated when genes are
active and become methylated during silencing, while gene bodies of
actively transcribed genes often show elevated methylation. Integration
of methylation data with chromatin state maps and gene expression
profiles provides comprehensive views of the regulatory landscape.

\subsection{ATAC-seq and Chromatin Accessibility}\label{sec-atacseq}

While ChIP-seq maps where specific proteins bind and bisulfite
sequencing reveals methylation patterns, neither directly measures
chromatin accessibility---the fundamental property that determines
whether regulatory DNA is available for transcription factor binding.
The assay for transposase-accessible chromatin using sequencing
(ATAC-seq) fills this gap by profiling open chromatin genome-wide.

The assay for transposase-accessible chromatin using sequencing
(ATAC-seq) profiles open chromatin by exploiting a hyperactive Tn5
transposase that simultaneously cuts DNA and inserts sequencing adapters
into accessible regions of the genome
(Figure~\ref{fig-atacseq-workflow}). Because nucleosomes and tightly
bound protein complexes occlude DNA from transposase integration, the
density of insertion sites provides a high-resolution map of
nucleosome-free regions, transcription factor binding sites, and
nucleosome positions.

\begin{tcolorbox}[enhanced jigsaw, rightrule=.15mm, leftrule=.75mm, colbacktitle=quarto-callout-note-color!10!white, colframe=quarto-callout-note-color-frame, toprule=.15mm, opacitybacktitle=0.6, title=\textcolor{quarto-callout-note-color}{\faInfo}\hspace{0.5em}{Figure 8.4 -- ATAC‑seq workflow for profiling chromatin accessibility}, left=2mm, colback=white, bottomtitle=1mm, toptitle=1mm, bottomrule=.15mm, coltitle=black, titlerule=0mm, breakable, opacityback=0, arc=.35mm]

\begin{figure}[H]

\centering{

\includegraphics[width=0.9\linewidth,height=\textheight,keepaspectratio]{images/Fig_8.4_ATAC-Seq.png}

}

\caption{\label{fig-atacseq-workflow}ATAC-seq experimental workflow}

\end{figure}%

This figure depicts the assay for transposase‑accessible chromatin using
sequencing (ATAC‑seq), which uses a hyperactive Tn5 transposase to
interrogate open chromatin regions genome‑wide. The workflow begins with
chromatin that contains a mixture of nucleosome‑occluded and
nucleosome‑free DNA; Tn5 loaded with sequencing adapters preferentially
binds and inserts into accessible regions, simultaneously fragmenting
and tagging these sites in a single tagmentation step. The resulting
short, adapter‑flanked DNA fragments are purified and PCR‑amplified to
create a sequencing library enriched for accessible chromatin. After
next‑generation sequencing, the aligned reads form sharp peaks at
nucleosome‑free regions such as active promoters, enhancers, and
regulatory elements, while nucleosome‑protected DNA yields larger
fragments that can be used to infer nucleosome positioning. The final
schematic shows these peaks superimposed on a cartoon chromatin fiber,
emphasizing how ATAC‑seq converts differences in chromatin compaction
into quantitative sequencing signals that can be integrated with
ChIP‑seq and RNA‑seq data to reconstruct cell‑type‑specific regulatory
landscapes. Created in \url{https://BioRender.com}.

\end{tcolorbox}

ATAC-seq requires relatively few cells and has a simple, rapid protocol
compared to earlier chromatin accessibility assays such as DNase-seq,
MNase-seq, and FAIRE-seq, making it widely adopted for mapping
regulatory landscapes across cell types and conditions. Computationally,
ATAC-seq data are processed much like ChIP-seq, with alignment followed
by peak calling to identify accessible regions, but downstream analyses
often focus on footprinting to infer transcription factor occupancy,
nucleosome phasing patterns, and integration with RNA-seq to connect
accessibility changes to gene expression programs\textsuperscript{132}.
Single-cell ATAC-seq extends this framework to individual cells,
enabling reconstruction of cell-type-specific chromatin accessibility
profiles and regulatory trajectories that are especially powerful when
jointly analyzed with single-cell RNA-seq in multi-omic
workflows\textsuperscript{133,134}.

\begin{tcolorbox}[enhanced jigsaw, rightrule=.15mm, leftrule=.75mm, colbacktitle=quarto-callout-note-color!10!white, colframe=quarto-callout-note-color-frame, toprule=.15mm, opacitybacktitle=0.6, title=\textcolor{quarto-callout-note-color}{\faInfo}\hspace{0.5em}{Case Study: Chromatin Landscapes and Transposable Elements}, left=2mm, colback=white, bottomtitle=1mm, toptitle=1mm, bottomrule=.15mm, coltitle=black, titlerule=0mm, breakable, opacityback=0, arc=.35mm]

A recent study by Huang and colleagues used comparative epigenomics to
show how species-specific chromatin landscapes shape the evolutionary
impact of transposable elements (TEs)\textsuperscript{135}. TEs are
selfish genetic elements that copy and insert throughout host genomes,
where they can disrupt genes, introduce novel regulatory sequences, and
promote chromosomal rearrangements\textsuperscript{135}. Many eukaryotes
counteract this activity through small RNA--guided pathways that deposit
repressive histone marks such as H3K9me2/3 at TE insertions, thereby
silencing their transcription. However, this host-directed silencing can
``spill over'' into flanking chromatin, inadvertently repressing nearby
genes and imposing a fitness cost on the host.

Huang et al.~combined population genomics with ChIP-seq profiling of
H3K9me2 across multiple \emph{Drosophila} species to quantify these
TE-mediated epigenetic effects\textsuperscript{135}. They showed that
individual TE insertions differ in the magnitude and spatial extent of
H3K9me2 spreading into neighboring chromatin, and that insertions with
stronger or broader repressive effects tend to be rarer in natural
populations, consistent with stronger purifying selection. By comparing
TE families across species, they further demonstrated that variation in
host chromatin machinery and in the propensity of TEs to recruit
repressive marks contributes to species-specific differences in overall
TE burden. These results link chromatin-level measurements from ChIP-seq
to classical population genetic signatures, illustrating how epigenetic
mapping can reveal the hidden fitness consequences of regulatory
perturbations and help explain why genomes with similar TE activity can
evolve very different TE landscapes\textsuperscript{135}.

\end{tcolorbox}

\subsection{Direct Detection of Base Modifications with Long-Read
Sequencing}\label{sec-base-modifications}

Recall from the introduction that third-generation sequencing platforms
fundamentally differ from short-read technologies in their ability to
sequence native DNA without amplification. This property extends beyond
merely generating longer reads---it enables direct detection of chemical
modifications to DNA bases during sequencing. Specifically, PacBio and
Oxford Nanopore sequencing technologies detect base modifications
directly from native DNA without bisulfite conversion or other chemical
treatments\textsuperscript{136}. These platforms measure kinetic signals
during sequencing that differ between modified and unmodified bases,
enabling simultaneous readout of genetic sequence and epigenetic
modifications from the same DNA molecule.

PacBio sequencing detects modifications through changes in polymerase
kinetics as the enzyme incorporates fluorescently labeled nucleotides.
The interpulse duration---the time between successive nucleotide
incorporations---differs for modified versus unmodified bases, providing
a signature of modification status. This approach enables detection of
5-methylcytosine, 5-hydroxymethylcytosine, and N6-methyladenine in
bacterial genomes.

Oxford Nanopore sequencing measures ionic current as DNA strands pass
through protein nanopores. Modified bases alter the current signal in
characteristic ways, enabling their identification through machine
learning models trained on known modified sequences. Recent advances
have improved the accuracy of modification calling and expanded the
repertoire of detectable modifications.

The ability to detect modifications directly on long reads enables
phasing of epigenetic marks with genetic variants over large genomic
distances, revealing how genetic variation influences epigenetic states.
This integrated view of sequence and modification has proven
particularly valuable for studying imprinted genes, where maternal and
paternal alleles show differential methylation, and for characterizing
bacterial restriction-modification systems where methylation patterns
provide strain-specific signatures.

\section{Metagenomic Sequencing}\label{sec-metagenomics}

Metagenomics extends sequencing approaches to complex mixtures of
organisms, enabling culture-independent characterization of microbial
communities from environmental samples. Rather than sequencing a single
isolated genome, metagenomic sequencing captures DNA from all organisms
present in a sample, providing a comprehensive view of community
composition and functional potential\textsuperscript{137}.

\subsection{16S rRNA Amplicon Sequencing}\label{sec-16s-sequencing}

The most widely used approach for characterizing bacterial communities
involves targeted amplification and sequencing of the 16S ribosomal RNA
gene\textsuperscript{138}. This gene is present in all bacteria and
archaea, contains both highly conserved regions suitable for universal
PCR primers and variable regions that enable taxonomic discrimination,
and has been extensively characterized through decades of Sanger
sequencing. Amplicon sequencing generates millions of 16S sequences from
a sample, which are then clustered into operational taxonomic units
(OTUs) or resolved into exact sequence variants (ASVs) and taxonomically
classified by comparison to reference databases.

Computational analysis of 16S data involves quality filtering, denoising
to correct sequencing errors, chimera removal to eliminate PCR
artifacts, and taxonomic assignment. Alpha diversity metrics quantify
the number and relative abundance of taxa within samples, while beta
diversity metrics compare community composition between samples.
Statistical methods test for differential abundance of taxa between
experimental groups and identify associations between microbial taxa and
environmental or clinical variables.

While 16S amplicon sequencing provides a cost-effective method for
taxonomic profiling, it has inherent limitations including PCR bias,
limited taxonomic resolution beyond the genus level, and inability to
provide functional information about community capabilities. These
limitations have driven adoption of shotgun metagenomic approaches that
sequence all DNA in a sample.

\subsection{Shotgun Metagenomics and Functional
Profiling}\label{sec-shotgun-metagenomics}

Shotgun metagenomic sequencing generates short reads from all DNA in a
sample without targeted amplification, enabling both taxonomic and
functional characterization of microbial
communities\textsuperscript{139}. Computational analysis involves either
assembling reads into longer contigs that may represent complete or
partial genomes, or directly mapping reads to reference databases of
genes and genomes.

Assembly-based approaches use specialized metagenome assemblers that
account for the presence of multiple related organisms with varying
abundances. Contigs can be binned into metagenome-assembled genomes
(MAGs) representing individual taxa, which are then annotated to predict
their metabolic capabilities. Reference-based approaches map reads to
databases of known genes or genomes, quantifying the abundance of
specific taxa and functional categories. Hybrid approaches combine
assembly and mapping to maximize recovery of both known and novel
sequences. Methodological reviews detail that the choice between
assembly-based and reference-based approaches depends on the study
goals: assembly is preferred when discovering novel organisms or
characterizing strain-level variation, while reference mapping provides
faster quantification when comprehensive databases exist for the
community of interest\textsuperscript{139}.

\begin{tcolorbox}[enhanced jigsaw, rightrule=.15mm, leftrule=.75mm, colbacktitle=quarto-callout-note-color!10!white, colframe=quarto-callout-note-color-frame, toprule=.15mm, opacitybacktitle=0.6, title=\textcolor{quarto-callout-note-color}{\faInfo}\hspace{0.5em}{Case Study: The Human Microbiome Project}, left=2mm, colback=white, bottomtitle=1mm, toptitle=1mm, bottomrule=.15mm, coltitle=black, titlerule=0mm, breakable, opacityback=0, arc=.35mm]

The Human Microbiome Project (HMP), launched by the NIH Common Fund in
2008, aimed to characterize the microbial communities inhabiting the
human body and understand their roles in health and
disease\textsuperscript{140,141}. The first phase focused on developing
reference resources and profiling healthy adults, while the second phase
(integrated HMP or iHMP) examined dynamic changes in the microbiome
associated with specific conditions including pregnancy, inflammatory
bowel disease, and onset of type 2 diabetes.

The HMP generated an unprecedented dataset including 16S rRNA gene
sequences from over 300 healthy adults sampled at up to 18 body sites
over multiple time points, along with whole metagenome shotgun sequences
from a subset of samples. This effort produced the world's largest
metagenomic sequence dataset from a single human cohort and resulted in
sequencing of approximately 3,000 reference bacterial genomes isolated
from human-associated sites.

Key findings from the HMP revealed substantial interpersonal variation
in microbiome composition, with different body sites harboring distinct
microbial communities characterized by different dominant phyla
(Figure~\ref{fig-microbiome-distribution}). Despite this compositional
variation, functional profiles---the genes and metabolic pathways
present---showed greater consistency across individuals, suggesting
functional redundancy where different taxa perform similar metabolic
roles.

Figure 8.5 -- The human microbiome across body sites

\begin{figure}[H]

\centering{

\includegraphics[width=0.5\linewidth,height=\textheight,keepaspectratio]{images/Figure_8.5_Microbiome.png}

}

\caption{\label{fig-microbiome-distribution}Human microbiome
distribution}

\end{figure}%

This figure illustrates the distribution of microbial communities across
major human body sites, highlighting that the human body harbors
approximately 10 to 100 trillion symbiotic bacterial cells representing
more than 10,000 distinct species. The microbiome is not uniformly
distributed but instead shows strong habitat specialization, with the
gastrointestinal tract hosting approximately 30\% of the total microbial
biomass, followed by the oral cavity (26\%), skin (21\%), airways
(14\%), and urogenital apparatus (9\%). Each body site presents a
distinct ecological niche characterized by unique pH, oxygen
availability, nutrient composition, and immune surveillance, leading to
site‑specific microbial community structures. The figure also emphasizes
an important clinical consideration: although these bacteria are
typically symbiotic and contribute to host health through functions such
as nutrient metabolism, immune system training, and pathogen exclusion,
they can become pathogenic if translocated to inappropriate body sites,
underscoring the importance of understanding both taxonomic composition
and spatial organization of the human microbiome. This conceptual
framework motivates metagenomic studies like the Human Microbiome
Project that profile microbial communities across body sites to
establish baseline reference data for health and to identify dysbiosis
associated with disease states. Created in \url{https://BioRender.com}.

The iHMP incorporated multi-omics profiling including
metatranscriptomics, metabolomics, and host immune profiling to capture
the dynamic interplay between microbiome and host. Longitudinal sampling
revealed that individual microbiomes are generally stable over time but
can undergo rapid shifts in response to perturbations such as infection
or dietary changes. Integration of microbiome data with host molecular
profiles has revealed associations between microbial metabolism and host
physiology, opening new avenues for microbiome-targeted therapeutics.

The HMP has established foundational resources, analytical tools, and
ethical frameworks that continue to support microbiome research. Public
availability of the data has enabled hundreds of secondary analyses and
methodological developments, exemplifying the value of large-scale,
open-access genomic resources.

\end{tcolorbox}

Metagenomic analysis faces several computational challenges including
the vast diversity of microorganisms, the presence of previously
uncharacterized organisms not represented in reference databases,
strain-level variation within species, and horizontal gene transfer that
creates mosaic genomes. Ongoing efforts to expand reference databases
through large-scale cultivation efforts and single-cell genomics aim to
reduce the fraction of sequences that cannot be classified.

\section{Multi-Omic Integration}\label{sec-multi-omics}

The ultimate goal of many genomics studies is to develop integrated
models that connect genome sequence to molecular phenotypes and
organismal traits\textsuperscript{142}. Multi-omic approaches combine
data from different molecular layers---genome, transcriptome, epigenome,
proteome, metabolome---to build comprehensive views of biological
systems Figure~\ref{fig-multiomics-workflow}. Integration of these
diverse data types presents both technical and conceptual challenges but
promises to reveal emergent properties not apparent from any single data
type\textsuperscript{143}.

\begin{tcolorbox}[enhanced jigsaw, rightrule=.15mm, leftrule=.75mm, colbacktitle=quarto-callout-note-color!10!white, colframe=quarto-callout-note-color-frame, toprule=.15mm, opacitybacktitle=0.6, title=\textcolor{quarto-callout-note-color}{\faInfo}\hspace{0.5em}{Figure 8.6 -- Multi‑omic integration from genotype to phenotype}, left=2mm, colback=white, bottomtitle=1mm, toptitle=1mm, bottomrule=.15mm, coltitle=black, titlerule=0mm, breakable, opacityback=0, arc=.35mm]

\begin{figure}[H]

\centering{

\includegraphics[width=0.35\linewidth,height=\textheight,keepaspectratio]{images/Figure_8.6_MultiOmics.png}

}

\caption{\label{fig-multiomics-workflow}Multi-omics workflow}

\end{figure}%

This figure depicts the hierarchical flow of biological information from
DNA sequence through multiple molecular layers to observable phenotypic
outcomes, illustrating the conceptual framework underlying multi‑omic
integration studies. At the top, the genomic layer captures the static
DNA sequence that encodes genes and regulatory elements. The
transcriptomic layer represents the dynamic subset of the genome that is
transcribed into mRNA in a given cell type or condition, reflecting both
the presence of regulatory signals and the activity of transcription
factors and chromatin modifiers. The proteomic layer shows the
translation of mRNA into polypeptide chains, which undergo
post‑translational modifications and assembly into functional protein
complexes that carry out cellular processes. The metabolomic layer
comprises the small molecule products and substrates of enzymatic
reactions, providing a functional readout of cellular biochemistry and
physiological state. Finally, the integration of these molecular layers
produces phenotypic changes observable at the cellular, tissue, or
organismal level. The central red arrow emphasizes that information
flows downward from genome to phenome, but also that feedback and
regulation operate at every transition: transcript levels do not
perfectly predict protein abundance, and protein levels do not perfectly
predict metabolic flux. Multi‑omic studies aim to measure all of these
layers simultaneously in the same samples, enabling researchers to
dissect regulatory networks, identify causal relationships between
molecular features and phenotypes, and build predictive models that
account for the complexity of biological systems. Created in
\url{https://BioRender.com}.

\end{tcolorbox}

\subsection{Integration Strategies and Data
Harmonization}\label{sec-integration-strategies}

Successful multi-omic integration requires careful attention to
experimental design, data harmonization, and appropriate statistical
methods. Ideally, all molecular measurements are made on the same
samples under identical conditions, enabling direct correlation of
signals across data types. Sample preparation protocols must balance the
requirements of different assays, and batch effects must be carefully
controlled or computationally corrected. Data harmonization involves
mapping features from different platforms to common coordinate systems,
normalizing measurements to comparable scales, and accounting for the
different statistical properties of different data
types\textsuperscript{144}.

Integration methods range from simple correlation approaches that
identify associations between layers, to more sophisticated methods that
model causal relationships or identify modules of co-regulated features
across data types. Network-based approaches construct multi-layer
networks where nodes represent genes, proteins, metabolites, or other
molecular entities and edges represent relationships within and between
layers. Machine learning methods can integrate multiple data types as
features for prediction tasks, while statistical methods adapted from
other fields provide frameworks for joint modeling. The compgenomr
textbook chapter on multi-omics analysis by Ronen provides R
implementations of several integration approaches, including matrix
factorization methods that identify latent factors driving covariation
across data types and network-based methods that construct multi-layer
graphs connecting molecular features\textsuperscript{144}.

\subsection{Applications of Multi-Omic
Analysis}\label{sec-multi-omic-applications}

Multi-omic integration has proven particularly powerful for
understanding complex diseases where dysfunction spans multiple
molecular layers. Cancer genomics studies routinely integrate
whole-genome sequencing, RNA-seq, methylation arrays, and proteomics to
characterize tumors, identify driver mutations, classify subtypes, and
predict therapeutic response. The Cancer Genome Atlas provides a model
for systematic multi-omic profiling of large cohorts, demonstrating how
integrated analysis reveals connections between genetic alterations and
downstream molecular changes.

In systems biology, multi-omic profiling enables reconstruction of
regulatory networks and metabolic pathways, revealing how genetic
variation is translated into phenotypic diversity. Studies of gene
expression regulation integrate ChIP-seq data on transcription factor
binding, histone modifications, and chromatin accessibility with RNA-seq
data on transcript levels and ribosome profiling data on translation to
dissect the multiple control points that determine protein abundance.
Similarly, integration of genomic, transcriptomic, and metabolomic data
allows mapping of the relationships between genetic variation, enzyme
expression, and metabolic flux through specific pathways.

The decreasing cost of sequencing technologies and the development of
single-cell multi-omic methods that measure multiple molecular layers in
the same cell are expanding the scope and resolution of integrative
studies. These approaches promise to reveal cell-type-specific
regulatory programs, trace cellular differentiation trajectories, and
ultimately enable predictive models that connect genotype to phenotype
across scales of biological organization.

\section{Synthesis: Connecting Sequencing Applications to Genomic
Perspectives}\label{sec-synthesis}

The diverse sequencing applications discussed in this chapter exemplify
how modern genomics transcends traditional boundaries between
disciplines and scales of inquiry. DNA sequencing projects like the 1000
Genomes Project and ancient DNA studies illuminate human evolutionary
history while also providing practical resources for medical genetics.
RNA sequencing bridges molecular and cellular perspectives by revealing
how genetic programs are executed differently across cell types and
conditions. Epigenetic sequencing uncovers regulatory mechanisms that
integrate environmental signals with genome function. Metagenomic
approaches extend our view beyond individual organisms to communities
and ecosystems.

The figures throughout this chapter illustrate this integration:
structural variants (Figure~\ref{fig-structural-variation}) connect
DNA-level variation to gene dosage changes; epigenetic modifications
(Figure~\ref{fig-epigenetics-overview}) mediate between genome sequence
and transcriptional output; experimental workflows
(Figure~\ref{fig-chipseq-workflow}; Figure~\ref{fig-atacseq-workflow})
show how molecular assays are translated into quantitative data; the
human microbiome (Figure~\ref{fig-microbiome-distribution}) represents
the challenge of analyzing communities rather than individuals; and
multi-omic integration (Figure~\ref{fig-multiomics-workflow}) captures
the hierarchical flow from genotype to phenotype through intermediate
molecular layers.

Each application draws upon multiple perspectives from Pevsner's
framework. Population resequencing projects simultaneously catalog
genetic variation (perspective 1), enable comparative analysis of
populations (perspective 2), test evolutionary hypotheses about
selection and demography (perspective 3), identify disease-associated
variants (perspective 4), and drive development of new computational
tools (perspective 5). Similarly, multi-omic integration necessarily
spans all five perspectives, from cataloging molecular features through
understanding biological principles to developing computational
infrastructure for data integration.

The computational workflows for these alternative applications share
common foundations with the variant calling pipelines introduced earlier
in this textbook: quality control, alignment to reference genomes, and
statistical inference from noisy data. However, each application also
requires specialized tools and analytical approaches tailored to its
specific biological questions and data characteristics. Mastery of
modern genomics requires not only understanding these specialized
methods but also recognizing when and how to apply them appropriately.

Looking forward, emerging sequencing technologies and analytical methods
continue to expand the frontiers of what is possible. Long-read
sequencing is enabling complete, gapless genome assemblies and revealing
structural variation previously hidden by repetitive sequences.
Single-cell methods are dissecting cellular heterogeneity at
unprecedented resolution. Spatial transcriptomics is adding location
information to molecular measurements. Each advance opens new questions
and applications, ensuring that genome sequencing will remain a dynamic
and evolving field for decades to come.

\section{Practice Problem Sets}\label{sec-problems}

\begin{enumerate}
\def\labelenumi{\arabic{enumi}.}
\item
  \textbf{Structural variant calling}: The human amylase locus exhibits
  complex copy number variation. Compare and contrast three
  complementary signals used to detect structural variants from
  short-read sequencing data: read depth, discordant paired-end mapping,
  and split reads. What are the strengths and limitations of each
  approach? Why are long reads advantageous for resolving complex
  structural variation?
\item
  \textbf{Phasing and haplotype reconstruction}: Explain why haplotype
  information (knowing which variants occur together on the same
  chromosome) is important for population genetic inference. Describe
  how statistical phasing algorithms like BEAGLE reconstruct haplotypes
  from unphased genotype data. What limitations do these methods have,
  and how do long reads and family-based designs address these
  limitations?
\item
  \textbf{Splice-aware alignment}: RNA-seq reads can span exon-exon
  junctions that do not exist in the genomic DNA. Explain how
  splice-aware aligners like HISAT2 or STAR handle these spliced
  alignments. Why is it important to provide splice junction annotations
  to these aligners, even though they can discover novel junctions?
\item
  \textbf{Exome sequencing variant interpretation}: You identify a novel
  missense variant in a disease gene through exome sequencing of a
  patient with an undiagnosed disorder. What lines of evidence would you
  use to evaluate whether this variant is likely pathogenic? How do the
  ACMG guidelines for variant classification integrate these different
  types of evidence?
\item
  \textbf{Bisulfite sequencing analysis}: Bisulfite treatment converts
  unmethylated cytosines to uracil but leaves methylated cytosines
  unchanged. Explain how this chemical property enables quantification
  of DNA methylation. What complications arise from incomplete bisulfite
  conversion? How would you estimate the conversion rate from the
  sequencing data itself?
\item
  \textbf{Integration of ChIP-seq and RNA-seq}: You have ChIP-seq data
  for a transcription factor and RNA-seq data measuring gene expression
  changes upon depletion of that factor. Design an analysis to evaluate
  whether genes with promoters bound by the transcription factor show
  changes in expression. What additional controls or analyses would
  strengthen causal inference about the regulatory role of this factor?
\item
  \textbf{Long-read epigenetics}: PacBio and Oxford Nanopore sequencing
  enable direct detection of base modifications without bisulfite
  conversion. What advantages does this provide compared to bisulfite
  sequencing? Describe a biological question that would specifically
  benefit from the ability to phase modifications with genetic variants
  on the same long reads.
\item
  \textbf{Human Microbiome Project analysis}: The HMP found that while
  taxonomic composition varies substantially between individuals,
  functional profiles are more consistent. Propose a hypothesis to
  explain this observation. How would you test your hypothesis using
  existing HMP data or by designing a new experiment?
\item
  \textbf{Contamination in metagenomic studies}: Metagenomic sequencing
  of clinical samples can be contaminated with human DNA, environmental
  microbes, and reagent contaminants. Design a computational pipeline to
  identify and filter potential contaminants. What negative controls
  should be included in metagenomic studies?
\item
  \textbf{Causal inference in multi-omics}: Correlation between
  molecular layers does not imply causation. What experimental designs
  or analytical approaches can help distinguish causal relationships
  from correlations in multi-omic data? Consider the role of time series
  data, genetic variants as instruments, and perturbation experiments.
\item
  \textbf{Single-cell multi-omics}: New technologies enable simultaneous
  measurement of chromatin accessibility and gene expression in the same
  single cells (SHARE-seq, 10X Multiome). What biological questions
  become tractable with paired measurements that could not be addressed
  by single-cell RNA-seq alone? What computational challenges arise from
  the increased dimensionality and sparsity of single-cell multi-omic
  data?
\end{enumerate}

\section{Reflection Questions}\label{sec-reflection}

\begin{enumerate}
\def\labelenumi{\arabic{enumi}.}
\item
  \textbf{Technology and biological questions}: The sequencing
  applications in this chapter all use variations of the same underlying
  sequencing technologies but address vastly different biological
  questions. Reflect on how the experimental design, sample preparation,
  and computational analysis must be tailored to each application. Are
  there fundamental limits to what can be learned from sequencing data
  alone, or do clever experimental designs always enable answering new
  questions?
\item
  \textbf{Reference genomes and bias}: Most analyses in this chapter
  rely on alignment to reference genomes. Reflect on how the choice of
  reference genome might introduce bias in population genetics studies,
  especially for underrepresented populations. How do pangenome
  references, which include multiple individuals' genomes, address this
  concern? What new computational challenges do pangenomes introduce?
\item
  \textbf{Ancient DNA and modern relevance}: Ancient DNA studies reveal
  patterns of human migration, admixture, and adaptation. How do these
  insights inform our understanding of present-day human genetic
  diversity? What are the ethical considerations in studying ancient
  human remains? How should researchers balance scientific inquiry with
  respect for descendant communities?
\item
  \textbf{Epigenetics and heritability}: Epigenetic modifications like
  DNA methylation can change in response to environmental exposures and
  potentially be transmitted across generations. How does this
  complicate the concepts of genetic and environmental contributions to
  phenotypic variation? Can sequencing-based approaches distinguish
  epigenetic changes that are cause versus consequence of phenotypic
  changes?
\item
  \textbf{Microbiome and the definition of ``organism''}: Metagenomic
  studies reveal that humans are colonized by trillions of
  microorganisms whose genes outnumber our own. How does this change our
  conception of what constitutes a human organism or human genome?
  Should microbiome composition be considered when studying human
  genetic variation and disease?
\item
  \textbf{Multi-omic integration and complexity}: As the number of
  measured molecular layers increases, integration becomes increasingly
  complex. Reflect on the trade-offs between comprehensive molecular
  profiling (which captures more biology but is more expensive and
  complex to analyze) versus focused studies of specific pathways or
  processes. When is multi-omic integration worth the additional effort
  and cost?
\item
  \textbf{From correlation to causation}: Much of genomics research
  identifies correlations between molecular features and phenotypes.
  Reflect on what additional information or experiments are needed to
  establish causation. How do emerging technologies like CRISPR genome
  editing change our ability to test causal hypotheses suggested by
  genomic data?
\end{enumerate}

\section{References}\label{sec-references}

\chapter{Hidden Markov Models in
Biology}\label{hidden-markov-models-in-biology}

\section{Introduction}\label{introduction}

Biological sequences contain rich information about the structure and
function of genes, proteins, and genomes. However, extracting this
information computationally requires sophisticated statistical models
that can capture the complex patterns hidden within strings of
nucleotides or amino acids. Often, the features we care most
about---such as whether a nucleotide belongs to an exon or an
intron---are not directly observed and must be inferred from the
sequence itself\textsuperscript{4,145}.

The fundamental challenge in biological sequence analysis often reduces
to a labeling problem. When examining a genomic sequence, we want to
identify which nucleotides belong to exons, which to introns, and which
to intergenic regions. When aligning protein sequences, we seek to
associate residues in a query sequence with homologous positions in
related proteins. In each case, we observe the sequence---the
``visible'' data---but the meaningful biological annotations remain
hidden and must be inferred\textsuperscript{4}.

\textbf{Hidden Markov models} (HMMs) provide a powerful probabilistic
framework for these problems. HMMs explicitly model both the observed
sequences and the hidden states that generate them, allowing us to
integrate diverse sources of information and compute not only the most
likely annotation, but also our confidence in that
annotation\textsuperscript{4,145}. In this chapter, we will first
introduce the core concepts behind HMMs, then focus on their
applications in biology: gene prediction, profile HMMs, pair HMMs, and
more advanced uses such as local ancestry inference. Along the way, we
will emphasize connections to earlier chapters and highlight how HMMs
underpin tools you have already encountered, such as ProbCons and
GeneWise\textsuperscript{146,147}.

\section{Understanding Hidden Markov
Models}\label{understanding-hidden-markov-models}

At their core, hidden Markov models represent biological sequences as
being generated by an underlying process that transitions between
discrete states, each of which emits observable symbols with
characteristic probabilities\textsuperscript{4,145}. The model consists
of \emph{two interrelated stochastic processes}: \textbf{a Markov chain
of hidden states} that we cannot directly observe, and \textbf{a visible
sequence of emitted symbols}---nucleotides, amino acids, or other
characters---that we can measure.

This dual-layer architecture is well-suited to biological problems
because the features we care most about, such as whether a nucleotide
participates in protein coding or lies within a regulatory element, are
precisely the hidden information we wish to recover from observable
sequence data\textsuperscript{4}. HMMs give us a principled way to
represent our assumptions about how these hidden states are organized
along the sequence and how they influence the observed
symbols\textsuperscript{145}.

To see how HMMs work in practice, consider a \textbf{simplified example}
of identifying splice sites in eukaryotic genes. Suppose we have a DNA
sequence that begins in an exon, contains one 5′ splice site, and ends
in an intron. The biological question is straightforward: where does the
transition from exon to intron occur? Different regions of genes have
distinct statistical properties that we can exploit. Exons might display
roughly uniform base composition, introns tend to be AT-rich, and the
consensus nucleotide at the 5′ splice site is almost always a
\texttt{G}\textsuperscript{4}. An HMM for this problem would contain
three states corresponding to these three functional categories: exon,
splice site, and intron. Each state would be characterized by emission
probabilities that reflect the expected base composition for that
functional class, and transition probabilities that describe the linear
order in which we expect to encounter these states as we move along the
sequence\textsuperscript{4,148}. Because this is an ordered sequence we
can add a start and stop state to either end as shown in
Figure~\ref{fig-toy-hmm}.

\begin{tcolorbox}[enhanced jigsaw, rightrule=.15mm, leftrule=.75mm, colbacktitle=quarto-callout-note-color!10!white, colframe=quarto-callout-note-color-frame, toprule=.15mm, opacitybacktitle=0.6, title=\textcolor{quarto-callout-note-color}{\faInfo}\hspace{0.5em}{Figure 9.1 - A simplified HMM for gene annotation}, left=2mm, colback=white, bottomtitle=1mm, toptitle=1mm, bottomrule=.15mm, coltitle=black, titlerule=0mm, breakable, opacityback=0, arc=.35mm]

\begin{figure}[H]

\centering{

\includegraphics[width=0.9\linewidth,height=\textheight,keepaspectratio]{images/Figure_9.1_ToyHMM_2.jpg}

}

\caption{\label{fig-toy-hmm}A toy HMM for a single exon--intron gene}

\end{figure}%

A toy hidden Markov model (HMM) for a simple eukaryotic gene with one
exon and one intron. The model has four emitting states---Exon, 5′
splice site (5′ SS), Intron, and Stop---plus non-emitting Start and Stop
states. Each emitting state is associated with a nucleotide emission
distribution (shown above the state), and arrows between states indicate
transition probabilities. The Exon and Intron states have high
self-transition probabilities (0.9), reflecting variable-length coding
and intronic regions, whereas transitions from Exon to 5′ SS and from
Intron to Stop occur with lower probability (0.1), modeling rare
boundary events. This compact model illustrates how HMMs can encode both
sequence composition and gene structure in a single probabilistic
framework, following the teaching example in Weisstein et
al.~(2021)\textsuperscript{148}.

\end{tcolorbox}

When an HMM generates a sequence, it produces two parallel streams of
information. One is the state path---the sequence of hidden states
visited during generation---which corresponds to the biological labels
we ultimately want to assign. The other is the observed sequence itself,
with each nucleotide emitted according to the emission probabilities of
whatever state the model occupied at that position. The state path forms
a \textbf{Markov chain} because the choice of the next state depends
only on the current state, not the entire history. Since we typically
observe only the nucleotide sequence and not the underlying state path,
this state sequence is \textbf{hidden}, giving the model its
name\textsuperscript{4}.

\begin{tcolorbox}[enhanced jigsaw, rightrule=.15mm, leftrule=.75mm, colbacktitle=quarto-callout-note-color!10!white, colframe=quarto-callout-note-color-frame, toprule=.15mm, opacitybacktitle=0.6, title=\textcolor{quarto-callout-note-color}{\faInfo}\hspace{0.5em}{Note \ref*{nte-hmm}: Key Terminology for Hidden Markov Models}, left=2mm, colback=white, bottomtitle=1mm, toptitle=1mm, bottomrule=.15mm, coltitle=black, titlerule=0mm, breakable, opacityback=0, arc=.35mm]

\quartocalloutnte{nte-hmm} 

Several fundamental concepts appear repeatedly when working with HMMs.
The \textbf{states} of an HMM represent the hidden functional categories
or labels we wish to assign to sequence positions. Each state is
characterized by two types of probabilities. \textbf{Emission
probabilities} \(e(x|i)\) specify the probability that state \(i\) will
produce symbol \(x\), and these probabilities must sum to one across all
possible symbols. \textbf{Transition probabilities} \(t(i,j)\) describe
the probability of moving from state \(i\) to state \(j\), and these
also must sum to one across all possible destination
states\textsuperscript{4,145}. A basic explanation using the example of
weather is provided as a
\href{https://auburn.hosted.panopto.com/Panopto/Pages/Viewer.aspx?id=b8c5095c-c69a-41a0-8b37-b3ed017c3e49}{video
introduction to these concepts}.

The collection of all states, emission probabilities, and transition
probabilities fully specifies an HMM. The joint probability that an HMM
with parameters \(\theta\) generates both a particular state path
\(\pi\) and an observed sequence \(S\) is the product of all the
relevant emission and transition probabilities. Because HMMs are fully
probabilistic models, we can use standard tools from probability and
statistics to manipulate these quantities, optimize parameters, and
interpret scores as probabilities rather than arbitrary
units\textsuperscript{4}.

In practice, we can now use the toy gene model in
Figure~\ref{fig-toy-hmm} to score concrete hypotheses about how an
observed sequence was generated. For example, given the sequence
\texttt{CTTCATGTGAAAGCAGACGTAAGTCA}, one candidate explanation is that
all bases before the splice site come from the Exon state, the splice
dinucleotide comes from the 5′~SS state, and all remaining bases come
from the Intron state, corresponding to the state path
\texttt{EEEE…E5II…I} shown in Figure 1 of Eddy 2004\textsuperscript{4}.
The probability of this particular path is obtained by multiplying the
appropriate transition probabilities (for example, the repeated
Exon→Exon self-transitions, followed by Exon→5′~SS and 5′~SS→Intron) and
emission probabilities for each base in the sequence given its state,
and then taking the log to obtain the log~P reported in the right
margin. We conceptually repeat this calculation for every possible state
path that could generate the sequence and select the most likely one
(the path with the highest probability, or equivalently the least
negative log~P), which in this toy example identifies the most probable
5′~splice site position\textsuperscript{4}.

\end{tcolorbox}

Constructing an HMM begins by specifying the symbol alphabet (for
example, \texttt{A}, \texttt{C}, \texttt{G}, and \texttt{T} for DNA),
the number and types of states in the model, the emission probabilities
associated with each state, and the transition probabilities between
states\textsuperscript{145}. Any model satisfying these basic properties
qualifies as an HMM, which means that designing a new model often
reduces to drawing an intuitive picture that represents the biological
system of interest\textsuperscript{4}. This graphical simplicity makes
HMMs a natural fit for building complex models of gene structure,
protein families, or evolutionary processes\textsuperscript{147}.

\section{Hidden Markov Models for Gene
Prediction}\label{hidden-markov-models-for-gene-prediction}

One of the most successful applications of HMMs in computational biology
is the prediction of gene structures in genomic
sequences\textsuperscript{147}. The challenge of computational gene
finding is far from trivial, particularly in eukaryotes where genes are
typically interrupted by introns that can span thousands of base pairs.
Early gene prediction programs relied on simple pattern matching or
sliding window approaches, but these methods struggled to integrate
multiple sources of information and often produced inconsistent
predictions with high false positive rates\textsuperscript{147}.

The introduction of HMM-based gene finders transformed the field by
providing a principled statistical framework for combining diverse
signals---splice site consensus sequences, codon usage bias, exon and
intron length distributions, and open reading frame constraints---into a
single coherent model\textsuperscript{145,147}. Rather than building an
ad hoc scoring system, we specify how likely it is that coding regions
emit given DNA patterns, how often splice sites occur, and how exons and
introns are arranged. The HMM then computes the most probable gene
structure consistent with the observed sequence and model parameters.

\subsection{Generalized HMMs and Gene
Structure}\label{generalized-hmms-and-gene-structure}

Most modern gene prediction programs employ a variant of HMMs called
generalized hidden Markov models (GHMMs). While ordinary HMMs emit a
single symbol per state, GHMMs allow each state to emit a sequence of
variable length. This modification is essential for gene finding because
functional elements like exons and introns span multiple nucleotides,
and their lengths vary considerably from one instance to
another\textsuperscript{147}.

In a GHMM for gene prediction, states correspond to biological features
such as exons, introns, splice donor regions, splice acceptor regions,
and intergenic sequences. Each state is associated with a probability
model that describes both the composition and the length distribution of
sequences that the state can emit\textsuperscript{147}. For example,
exons are modeled using codon statistics that reflect nonuniform codon
usage, while introns may be modeled with simpler background nucleotide
distributions and length distributions that match the organism's
intronic architecture\textsuperscript{145}.

Splice sites provide a concrete illustration of this modeling strategy.
The first few nucleotides at intron boundaries show strong
position-specific preferences. In many eukaryotes, the canonical donor
site begins with the dinucleotide \texttt{GT}, and the third and fourth
positions show characteristic biases toward purines and adenine,
respectively\textsuperscript{147}. By compiling these position-specific
frequencies from a training set of confirmed splice sites, we can build
a weight matrix model that assigns a probability to any candidate splice
site sequence. In the GHMM, this weight matrix becomes the emission
model for the splice donor state and helps distinguish true splice sites
from spurious \texttt{GT} dinucleotides elsewhere in the
genome\textsuperscript{147}.

The middle regions of exons and introns require other probability
models. For exon states, gene finders typically model the sequence using
a probabilistic description that captures codon usage bias and period-3
base composition patterns\textsuperscript{145}. One common approach uses
separate states for the first, second, and third positions of codons,
each with its own emission probabilities. For intron states, the model
might use a simpler background nucleotide distribution, with emission
probabilities estimated from intronic training sequences. Length
distributions for exons, introns, and intergenic regions are also
modeled, often using geometric or more flexible distributions tailored
to the organism\textsuperscript{147}.

Dual-genome and multi-genome \emph{de novo} gene predictors extend this
idea by using alignments between a target genome and one or more related
``informant'' genomes\textsuperscript{147}. Instead of treating the
observed data as a single DNA sequence, these methods model aligned
pairs (or tuples) of nucleotides and indels, allowing them to leverage
evolutionary conservation to improve gene predictions. Conserved coding
regions show characteristic patterns in alignments, such as mismatches
that preserve amino acid sequences or gaps that maintain reading frame,
and these patterns can be encoded in the emission models of a
GHMM\textsuperscript{145,147}.

\subsection{AUGUSTUS and Constraint-Based
Annotation}\label{augustus-and-constraint-based-annotation}

AUGUSTUS is a widely used gene prediction system that exemplifies the
power and flexibility of GHMM-based gene finding\textsuperscript{146}.
At its core, AUGUSTUS models the sequence around splice sites, branch
point regions, translation start sites, coding regions, non-coding
regions, and the length distributions of exons, introns, and intergenic
regions. The model is trained on known genes from a given species,
producing species-specific parameter sets that capture typical gene
architecture\textsuperscript{146}.

Beyond its probabilistic modeling, AUGUSTUS introduces an important
practical feature: the ability to incorporate user-defined constraints.
Users can specify positions of known exons, introns, splice sites,
translation start sites, or stop codons, based on experimental evidence
or other computational predictions\textsuperscript{146}. The system then
searches for the most probable gene structure that is consistent with
these constraints. This capability is particularly valuable when partial
experimental data (such as an RT-PCR-validated intron or an expressed
sequence tag alignment) is available. Rather than ignoring such data or
relying on manual interpretation alone, AUGUSTUS integrates it directly
into the probabilistic framework. This combination of a rich
probabilistic model and constraint-based refinement makes AUGUSTUS an
ideal tool for bridging the gap between fully automated and manual
annotation\textsuperscript{146,147}.

\begin{tcolorbox}[enhanced jigsaw, rightrule=.15mm, leftrule=.75mm, colbacktitle=quarto-callout-note-color!10!white, colframe=quarto-callout-note-color-frame, toprule=.15mm, opacitybacktitle=0.6, title=\textcolor{quarto-callout-note-color}{\faInfo}\hspace{0.5em}{Lab Connections}, left=2mm, colback=white, bottomtitle=1mm, toptitle=1mm, bottomrule=.15mm, coltitle=black, titlerule=0mm, breakable, opacityback=0, arc=.35mm]

In \href{Lab06.qmd}{Lab 6}, you will be working with a toy model of
\emph{ab initio} gene prediction. Similarly, in your Manual Annotation
Project, you will be working with the homology-based annotation tool
using GeneWise. AUGUSTUS can be viewed as a natural extension of those
ideas to full-genome scale, providing a realistic context for
understanding how HMMs underpin modern annotation
pipelines\textsuperscript{146}.

\end{tcolorbox}

\subsection{Manual vs.~Automated
Annotation}\label{manual-vs.-automated-annotation}

The distinction between manual and automated gene annotation reflects a
fundamental tension in genomics between throughput and
accuracy\textsuperscript{147}. Automated annotation pipelines use
computational gene finders like AUGUSTUS to generate predictions across
entire genomes, enabling the rapid processing of newly sequenced
species. These automated predictions serve as crucial starting points
for functional genomics studies and can be updated as new evidence
becomes available\textsuperscript{146}.

Manual annotation, by contrast, involves human curators examining
computational predictions in the context of multiple lines of evidence,
including expressed sequence tag alignments, protein homology, RNA-seq
data, and comparative genomics information\textsuperscript{147}.
Curators use this evidence to refine gene models, identify alternative
splice forms, and resolve ambiguous cases where computational methods
disagree. This process is time-consuming and expensive, but it generally
produces more accurate gene structures than purely computational
approaches\textsuperscript{147}.

The most reliable annotation strategy combines both approaches.
Automated gene finders provide initial predictions, which are then
manually curated for genes of particular interest or in cases where
multiple predictions conflict. In your laboratory work, when you
manually annotate genes using GeneWise, you are effectively acting as a
human integrator of evidence, complementing the purely statistical
reasoning embodied by HMM-based gene finders.

\section{Profile Hidden Markov
Models}\label{profile-hidden-markov-models}

While the basic HMM architecture works well for many sequence analysis
problems, biological sequences often display position-specific patterns
that require more specialized models. Profile hidden Markov models, or
profile HMMs, address this need by representing the consensus properties
of a multiple sequence alignment\textsuperscript{145}. Rather than
modeling individual sequences, profile HMMs capture the statistical
patterns that characterize an entire family of related sequences, making
them particularly powerful for tasks like protein family classification,
remote homology detection, and motif discovery.

The architecture of a profile HMM differs from general-purpose HMMs in
its strictly linear, left-to-right structure without cycles. The model
repetitively uses three types of states to describe each position in the
alignment. Match states represent conserved positions in the alignment,
with emission probabilities reflecting the observed base or amino acid
frequencies at that position. Insert states handle additional residues
that appear in some sequences but not others, accommodating length
variation among family members. Delete states are silent, non-emitting
placeholder states that allow sequences to skip conserved positions,
thereby modeling deletions relative to the
consensus\textsuperscript{145}.

Constructing a profile HMM from a multiple sequence alignment is
conceptually straightforward. Consider an alignment with a certain
number of columns representing a conserved domain or motif. For each
column, we must decide whether it should be modeled by a match state or
an insert state. A simple heuristic compares the number of residues to
the number of gaps in each column: if residues outnumber gaps, we treat
the column as a conserved position represented by a match state, and any
gaps are modeled by the corresponding delete state. If gaps outnumber
residues, we view the residues as insertions and represent the column
with an insert state. Once we have assigned state types to all columns,
we estimate emission probabilities by counting the frequency of each
symbol at each position, and we estimate transition probabilities by
counting state-to-state transitions across the sequences in the
alignment, often with pseudocounts to avoid zero
probabilities\textsuperscript{145}.

Profile HMMs have become workhorses of modern protein sequence
analysis\textsuperscript{145}. Software packages such as HMMER and SAM
provide complete toolkits for building profile HMMs from alignments,
searching sequence databases with profile HMMs, and comparing new
sequences against libraries of pre-built models. Large-scale resources
like the Pfam database and PROSITE have compiled profile HMMs for
thousands of known protein domains and families, enabling rapid
functional annotation of newly sequenced proteins. When a researcher
obtains a new protein sequence of unknown function, searching it against
Pfam can reveal which known domains it contains, providing immediate
insights into its likely biochemical activities and cellular
roles\textsuperscript{145}.

The applications of profile HMMs extend beyond simple sequence searches.
Comparing two profile HMMs---representing two different multiple
alignments or sequence families---can detect remote homology
relationships that would be missed by pairwise sequence
comparison\textsuperscript{145}. Tools such as HHsearch generalize
traditional pairwise alignment algorithms to find optimal alignments
between pairs of profile HMMs, greatly extending the depth of homology
detection. Researchers have also adapted profile HMMs to model
non-sequence features, such as protein secondary structure patterns,
where the model emits secondary structure symbols (helix, strand, coil)
rather than amino acids. These secondary-structure profile HMMs can
recognize three-dimensional protein folds based solely on predicted
secondary structure, demonstrating the versatility of the HMM
framework\textsuperscript{145}.

\section{Pair Hidden Markov Models}\label{pair-hidden-markov-models}

While profile HMMs model a family of related sequences, pair hidden
Markov models, or pair HMMs, address the fundamental problem of
comparing two individual sequences\textsuperscript{145}. Sequence
alignment is central to nearly all biological sequence analysis, as the
degree of similarity between sequences often indicates functional or
evolutionary relationships. Pair HMMs provide a formal probabilistic
framework for finding sequence alignments, computing alignment
probabilities, and evaluating the significance of aligned
regions\textsuperscript{145}.

The key innovation of pair HMMs is that they simultaneously generate two
aligned sequences rather than a single sequence. Consider a simple pair
HMM with three states. An insert state for sequence 1 emits a single
symbol in the first sequence but nothing in the second sequence,
representing an insertion in the first sequence or, equivalently, a
deletion in the second. Similarly, an insert state for sequence 2 emits
a symbol only in the second sequence. A match/mismatch state emits a
pair of symbols, one in each sequence, representing positions that are
aligned to each other. As the model transitions among these three
states, it builds up a pairwise alignment between the two output
sequences\textsuperscript{145}.

There is a one-to-one correspondence between a pair HMM state path and a
sequence alignment. Each state in the path specifies whether the current
position in each sequence is emitted as an aligned pair or as an
insertion in one sequence or the other. This means that finding the
optimal alignment between two sequences is mathematically equivalent to
finding the optimal state path through a pair HMM. A modification of the
Viterbi dynamic programming algorithm can find this optimal path
efficiently, with computational complexity proportional to the product
of the two sequence lengths\textsuperscript{145}. This approach recovers
alignments comparable to those produced by classical dynamic programming
algorithms, but the pair HMM framework provides additional capabilities
that standard alignment methods lack.

One major advantage of pair HMMs is their ability to compute the overall
probability that two sequences are related, rather than simply
identifying a single best alignment. When sequences are weakly similar,
there may be many plausible alignments with similar scores, and focusing
on only the optimal alignment can be misleading. The pair HMM framework
allows us to sum the probabilities over all possible alignments to
obtain the total probability of observing the two sequences given the
model. This calculation can be performed efficiently using a modified
version of the forward algorithm. Additionally, pair HMMs can compute
position-specific alignment probabilities, quantifying the probability
that any particular pair of residues are aligned to each other across
all possible alignments. These posterior probabilities provide a natural
measure of alignment confidence and can identify regions where the
alignment is uncertain\textsuperscript{145}.

Pair HMMs have found extensive applications beyond basic pairwise
alignment. Many state-of-the-art multiple sequence alignment programs,
including ProbCons, use pair HMMs as a core
component\textsuperscript{145}. ProbCons employs a pair HMM to compute
posterior probabilities for all pairwise alignments among the input
sequences, then uses these probabilities to guide the progressive
assembly of the full multiple alignment. Instead of directly using the
optimal Viterbi alignment for each sequence pair, ProbCons identifies
the alignment that maximizes the expected number of correctly aligned
positions based on the posterior probabilities. The algorithm further
improves accuracy by incorporating a probabilistic consistency
transformation, which considers information from alignments to
intermediate sequences when scoring each pairwise alignment. This
consistency-based approach has demonstrated substantial improvements
over traditional progressive alignment methods and has influenced the
design of subsequent multiple alignment tools\textsuperscript{145}.

\begin{tcolorbox}[enhanced jigsaw, rightrule=.15mm, leftrule=.75mm, colbacktitle=quarto-callout-note-color!10!white, colframe=quarto-callout-note-color-frame, toprule=.15mm, opacitybacktitle=0.6, title=\textcolor{quarto-callout-note-color}{\faInfo}\hspace{0.5em}{ProbCons and Multiple Sequence Alignment}, left=2mm, colback=white, bottomtitle=1mm, toptitle=1mm, bottomrule=.15mm, coltitle=black, titlerule=0mm, breakable, opacityback=0, arc=.35mm]

In Chapter 5, we explored various algorithms for multiple sequence
alignment, a fundamental challenge in comparative genomics and
evolutionary analysis. ProbCons represents a significant advance in this
area by leveraging pair HMMs to bring probabilistic rigor to the
alignment process\textsuperscript{145}. Traditional progressive
alignment methods construct multiple alignments by iteratively combining
pairwise alignments according to a guide tree, but they lack a
principled framework for quantifying alignment uncertainty or
propagating information across multiple sequences.

ProbCons addresses these limitations through three key innovations.
First, it uses a pair HMM to compute posterior alignment probabilities
for all sequence pairs, providing a confidence estimate for each
potential aligned position. Second, it applies probabilistic consistency
transformation, which refines pairwise alignment scores by considering
paths through intermediate sequences. If sequence A aligns well to
sequence C, and sequence C aligns well to sequence B, this provides
additional evidence that A and B should be aligned in a consistent
manner. Third, ProbCons uses these refined scores to identify alignments
that maximize expected accuracy rather than simply optimizing a scoring
function.

The performance improvements from ProbCons have been substantial,
particularly for divergent sequences where alignment uncertainty is
high. This probabilistic framework has become a template for modern
multiple alignment algorithms and exemplifies how HMM-based approaches
can elevate the quality and reliability of fundamental bioinformatics
operations\textsuperscript{145}.

\end{tcolorbox}

Pair HMMs also play a central role in comparative gene prediction.
Algorithms such as SLAM and TWAIN use generalized pair HMMs to
simultaneously compare two genomic sequences and annotate their gene
structures. These methods combine the alignment capabilities of pair
HMMs with the gene structure modeling of gene-finding HMMs, creating
integrated models that can leverage conservation patterns between
species to improve gene prediction accuracy\textsuperscript{145}. The
resulting predictions are typically more accurate than single-genome
gene finders because evolutionary conservation provides strong evidence
for functional elements.

\section{Other Extensions and
Applications}\label{other-extensions-and-applications}

Beyond gene prediction and sequence alignment, HMMs have proven valuable
for analyzing biological sequences with more complex structural and
evolutionary features\textsuperscript{145}. RNA sequences present
particular challenges because their function depends not only on primary
sequence but also on secondary structure, which arises from base pairing
interactions between distant positions. Standard HMMs assume that each
emission depends only on the current state, making it difficult to model
the long-range correlations induced by base pairing. Several extensions
to the basic HMM framework have been developed to address this
limitation\textsuperscript{145}.

Context-sensitive HMMs (csHMMs) extend the HMM architecture by allowing
certain states to use contextual information from previous emissions
when determining emission probabilities. The model includes three types
of states: standard single-emission states, pairwise-emission states
that store emitted symbols in memory, and context-sensitive states that
retrieve stored symbols and adjust their emission probabilities
accordingly. For modeling RNA secondary structures, pairwise-emission
states can be placed at positions where one half of a base pair occurs,
storing the emitted base for later use. When the model reaches the
complementary base-pairing position, a corresponding context-sensitive
state retrieves the stored base and emits a complementary nucleotide
with high probability. This mechanism allows csHMMs to represent a wide
range of base-pairing patterns, including certain complex structures
such as pseudoknots\textsuperscript{145}.

Profile context-sensitive HMMs (profile-csHMMs) combine the RNA modeling
capabilities of csHMMs with the profile representation of multiple
alignments. These models can represent the conserved sequence features
and base-pairing patterns that characterize families of functional RNAs.
Constructing a profile-csHMM from an RNA alignment with structural
annotation is analogous to building a standard profile HMM, but match
states can be designated as single-emission, pairwise-emission, or
context-sensitive depending on whether they represent unpaired positions
or the two sides of conserved base pairs. Profile-csHMMs have been
applied to RNA homology searches and structural alignment, providing
powerful tools for analyzing non-coding RNAs\textsuperscript{145}.

The computational cost of RNA analysis with sophisticated HMM variants
can be substantial, as modeling base-pairing interactions typically
requires algorithms with cubic time complexity or worse. Various
acceleration strategies have been developed to make these methods
practical for genome-scale searches. One effective approach uses
standard profile HMMs, which consider only sequence information and run
quickly, to prescreen a database and filter out sequences that are
unlikely to match the full structural model. The more expensive
structure-aware model is then applied only to sequences that pass the
initial filter, dramatically reducing the average computational cost
while maintaining high sensitivity\textsuperscript{145}.

\subsection{HMMs for Local Ancestry and Association
Studies}\label{hmms-for-local-ancestry-and-association-studies}

Another important application domain for HMMs is local ancestry
inference in admixed populations. When individuals have genetic ancestry
from multiple continental populations, their genomes become mosaics of
chromosomal segments inherited from different ancestral groups.
HMM-based methods can model this mosaic as a sequence of hidden ancestry
states along the genome, with the observed genotypes (or haplotypes)
providing evidence about which ancestry is most likely at each
position\textsuperscript{145,149}.

In a typical local ancestry HMM, the hidden states correspond to
ancestry labels (for example, African, European, Indigenous American),
and the emission probabilities describe the likelihood of observing the
genotypes in a given segment assuming each possible ancestry, given
reference panels from those populations\textsuperscript{149}. Transition
probabilities between ancestry states reflect recombination, with higher
recombination rates leading to more frequent ancestry switches. Once
local ancestry has been inferred, these segment-level assignments can be
integrated into genetic association analyses, admixture mapping, and
polygenic risk score construction to improve power and interpretability
in diverse populations\textsuperscript{149}.

Recent reviews highlight both the opportunities and challenges of local
ancestry-based association studies, including issues of reference panel
selection, model misspecification, and computational scaling as sample
sizes grow\textsuperscript{149}. Nevertheless, the HMM framework remains
central because it provides a flexible way to combine recombination
models, ancestry information, and genotype likelihoods within a coherent
probabilistic system.

\begin{tcolorbox}[enhanced jigsaw, rightrule=.15mm, leftrule=.75mm, colbacktitle=quarto-callout-note-color!10!white, colframe=quarto-callout-note-color-frame, toprule=.15mm, opacitybacktitle=0.6, title=\textcolor{quarto-callout-note-color}{\faInfo}\hspace{0.5em}{Case Study: Local Ancestry Inference and Neanderthal Introgression}, left=2mm, colback=white, bottomtitle=1mm, toptitle=1mm, bottomrule=.15mm, coltitle=black, titlerule=0mm, breakable, opacityback=0, arc=.35mm]

Hidden Markov models have recently illuminated one of the most
fascinating questions in human evolutionary history: the extent and
timing of interbreeding between modern humans and Neanderthals. As
modern humans migrated out of Africa and encountered Neanderthal
populations in Eurasia, genetic material was exchanged through
interbreeding, leaving traces of Neanderthal ancestry in the genomes of
present-day non-African populations\textsuperscript{150,151}.
Identifying these introgressed segments and studying their properties
over time requires sophisticated computational methods to distinguish
Neanderthal ancestry from the predominant modern human background.

Two recent studies approached this problem using HMM-based local
ancestry methods on ancient and present-day genomes. Iasi et al.~used a
specialized HMM called admixfrog to catalog Neanderthal ancestry
segments in more than 300 ancient and present-day human genomes spanning
the past 50,000 years\textsuperscript{150}. The method employs an HMM
where each genomic window can be assigned to one of three ancestries:
Neanderthal, Denisovan (another archaic human group), or modern human.
The model uses high-coverage Neanderthal and Denisovan genomes and a
panel of present-day African individuals as references representing
these ancestry sources. For each target individual, the HMM evaluates
small genomic windows and estimates the most likely ancestry based on
the pattern of alleles present and their frequencies in the reference
populations\textsuperscript{150}.

In parallel, Sümer et al.~analyzed some of the earliest modern human
genomes outside Africa to constrain the timing of Neanderthal admixture
events\textsuperscript{151}. By combining HMM-based identification of
archaic segments with demographic modeling, they showed that early
modern humans carried long, distinctive Neanderthal segments that are
not fully shared with present-day non-Africans, pointing to multiple
introgression events and helping to narrow the window during which
Neanderthal gene flow occurred\textsuperscript{151}. Together, these
studies illustrate how applying HMMs to both ancient and present-day
genomes allows researchers to disentangle complex admixture histories
and relate them to archaeological evidence.

The work by Iasi et al.~revealed that the vast majority of Neanderthal
ancestry in modern humans derives from a single extended period of gene
flow that occurred between roughly 50,500 and 43,500 years ago in the
common ancestors of all non-African populations\textsuperscript{150}.
This conclusion emerged from several lines of evidence enabled by the
HMM analysis. First, Neanderthal segments are extensively shared among
present-day populations, with the sharing patterns reflecting the
overall population structure of modern humans rather than indicating
independent admixture events. Second, the Neanderthal segments in
post-Ice Age individuals show similar levels of divergence from
sequenced Neanderthal genomes, consistent with introgression from a
single or closely related group. Third, by measuring how Neanderthal
segments decay in length over time due to recombination, the researchers
could estimate both the timing and duration of the admixture
period\textsuperscript{150}.

The HMM framework also enabled the researchers to study natural
selection on Neanderthal ancestry by tracking how the frequency of
introgressed segments changed over time. They identified genomic regions
where Neanderthal variants rose to high frequency rapidly, suggesting
positive selection for beneficial alleles affecting skin pigmentation,
metabolism, and immune function\textsuperscript{150}. Conversely, they
found that genomic regions depleted of Neanderthal ancestry---so-called
``archaic deserts''---were established almost immediately after the gene
flow event, indicating strong negative selection against Neanderthal
variants in functionally constrained regions. The X chromosome showed an
especially striking pattern, with large deserts of Neanderthal ancestry
overlapping regions that experienced selective sweeps in non-African
populations\textsuperscript{150,151}.

This case study illustrates how HMMs can address complex questions in
population genetics by formally modeling the mosaic ancestry structure
of admixed genomes. The ability to assign ancestry to small genomic
segments, quantify uncertainty in these assignments, and track changes
in ancestry patterns over time has transformed our understanding of
human evolutionary history and adaptation\textsuperscript{149--151}.

\end{tcolorbox}

\subsection{The Connection Between HMMs and
MCMC}\label{the-connection-between-hmms-and-mcmc}

Hidden Markov models and Markov chain Monte Carlo (MCMC) methods share
conceptual foundations in Markov chains, but they serve fundamentally
different roles in computational biology\textsuperscript{145}. HMMs are
generative models that describe how observed sequences are produced by
underlying state sequences, and we use them to infer the most likely
hidden states or to compute probabilities of observed
data\textsuperscript{4}. MCMC, by contrast, is a computational technique
for drawing samples from complex probability distributions that cannot
be analyzed analytically, and it plays a central role in Bayesian
inference and statistical modeling\textsuperscript{145}.

The connection becomes clearest when considering parameter estimation
for HMMs. When we have a collection of related sequences and want to
build an HMM that represents them, we need to estimate the transition
probabilities, emission probabilities, and sometimes the structure of
the model itself. The Baum--Welch algorithm, which is an
expectation--maximization (EM) approach, provides one standard method
for parameter estimation\textsuperscript{4,145}. However, EM algorithms
can be sensitive to local optima and may not fully characterize
uncertainty in parameter estimates.

Bayesian approaches to HMM parameter estimation use MCMC methods to
sample from the posterior distribution of model parameters given the
observed sequences. Rather than finding a single best set of parameters,
MCMC methods generate many parameter sets drawn from their posterior
probability distribution, naturally quantifying uncertainty and avoiding
over-commitment to local optima. The Markov chain in MCMC refers to a
sequence of parameter configurations where each new configuration is
proposed based only on the current configuration, satisfying the Markov
property. By carefully designing the proposal mechanism and acceptance
criteria, MCMC algorithms ensure that the chain eventually explores
parameter space in proportion to the posterior probability, providing
valid statistical inference\textsuperscript{145}.

Monte Carlo EM algorithms represent a hybrid approach that uses sampling
to approximate difficult-to-compute expectations within the EM
framework. When working with complex HMMs where computing exact
expectations is intractable, MCEM uses Monte Carlo sampling to
approximate these quantities, enabling parameter estimation for broader
classes of models\textsuperscript{145}. This interplay between HMMs and
MCMC methods highlights how different computational techniques can be
combined to address challenging problems in biological sequence
analysis.

\section{Practice Problem Sets}\label{practice-problem-sets-1}

\begin{enumerate}
\def\labelenumi{\arabic{enumi}.}
\item
  \textbf{Basic HMM Concepts}: Consider a simple HMM with two states: C
  (CpG island) and N (non-CpG island). The emission probabilities for
  state C give equal probability (0.25) to all four nucleotides, while
  state N is AT-rich with \(e_N(A) = e_N(T) = 0.35\) and
  \(e_N(C) = e_N(G) = 0.15\). The transition probability from C to N is
  0.1, and from N to C is 0.05. Calculate the probability that this HMM
  generates the sequence ACGT starting in state C and ending in state N
  along the state path C--C--C--N.
\item
  \textbf{Gene Finding with GHMMs}: Explain why generalized hidden
  Markov models are more appropriate than standard HMMs for eukaryotic
  gene prediction. What specific biological features of genes
  necessitate the variable-length emissions allowed by GHMMs? How would
  you modify a GHMM to model alternative splicing?
\item
  \textbf{Profile HMM Construction}: Given a multiple sequence alignment
  with seven columns, three sequences, and the following pattern:\\
  \texttt{Seq1:\ ACGT-GA}\strut \\
  \texttt{Seq2:\ ACG-\/-GA}\strut \\
  \texttt{Seq3:\ ACGTTGA}\strut \\
  Determine which columns should be represented by match states versus
  insert states using a simple gap/residue heuristic. Sketch the
  resulting profile HMM structure and estimate the emission
  probabilities for the match states.
\item
  \textbf{Pair HMM Alignment}: A pair HMM for protein sequence alignment
  has three states: M (match/mismatch), X (insert in sequence 1), and Y
  (insert in sequence 2). If you observe an alignment of length 10 that
  contains 7 matched positions and 3 positions where residues in
  sequence 1 are gapped, write out a plausible state path through the
  pair HMM. How does this state path correspond to the standard
  representation of a gapped alignment?
\item
  \textbf{Model Selection}: You have built two different HMMs to model a
  protein family---one simple model with 15 states and one complex model
  with 45 states. The complex model achieves higher likelihood on your
  training sequences. What concerns might you have about using the
  complex model for searching databases? How could you evaluate which
  model generalizes better to new sequences?
\item
  \textbf{RNA Secondary Structure}: Context-sensitive HMMs can model
  base-pairing interactions by using memory to store emitted bases and
  later retrieve them at complementary positions. Sketch a csHMM that
  could model a simple hairpin structure with a 6-nucleotide stem (three
  base pairs) and a 4-nucleotide loop. Label the states as
  pairwise-emission, context-sensitive, or single-emission, and indicate
  which states use memory.
\end{enumerate}

\section{Reflection Questions}\label{reflection-questions-7}

\begin{enumerate}
\def\labelenumi{\arabic{enumi}.}
\item
  \textbf{Integration of Evidence}: In your Manual Annotation Project,
  you performed manual gene annotation using GeneWise to make
  homology-based predictions. Reflect on how an HMM-based gene finder
  like AUGUSTUS integrates different types of evidence (splice signals,
  codon bias, length distributions). What are the advantages and
  limitations of having the computer automatically weight these
  different information sources compared to a human curator making
  judgment calls? When might you trust the HMM more than your own
  assessment, and when might your expert judgment carry more weight?
\item
  \textbf{The Hidden Information Problem}: Throughout this chapter, we
  emphasized that HMMs model two parallel processes---a visible sequence
  and a hidden state path. Think about other biological problems beyond
  those discussed where you observe some data but care primarily about
  hidden underlying processes. For example, consider ChIP-seq data
  showing protein--DNA binding across a genome, or Hi-C data showing
  chromosome conformation. \emph{Could HMMs be useful for these
  problems? What would the states represent? What would the observations
  be?}
\item
  \textbf{Probabilistic vs.~Deterministic Approaches}: Compare the HMM
  framework to simpler, deterministic approaches for sequence analysis,
  such as scanning for exact motif matches or using fixed score
  thresholds. What are the conceptual and practical advantages of a
  fully probabilistic approach? In what situations might a simpler
  deterministic method actually be preferable despite being less
  statistically sophisticated?
\item
  \textbf{Model Design and Biological Understanding}: One appealing
  feature of HMMs is that designing a model often involves drawing an
  intuitive picture. Reflect on whether this simplicity could be a
  double-edged sword. Does the ease of sketching an HMM architecture
  risk oversimplifying biology? How would you guard against building
  models that are mathematically elegant but biologically unrealistic?
\item
  \textbf{Ancient DNA and Temporal Genomics}: The Neanderthal ancestry
  case study illustrates how ancient DNA provides a time series of
  genomic data. How does having observations from multiple time points
  fundamentally change what we can infer about evolutionary processes?
  Consider other questions in evolutionary biology or medical genetics
  where access to historical samples might provide decisive evidence.
  What are the practical and ethical challenges of obtaining such
  samples?
\item
  \textbf{The Limits of HMMs}: We mentioned that HMMs struggle with
  long-range correlations because each emission depends only on the
  current state. Yet we have seen various
  workarounds---context-sensitive states for RNA, triplet-emitting
  states for codons. Reflect on whether these extensions maintain the
  conceptual elegance of HMMs or whether they indicate that
  fundamentally different approaches (like deep learning models) might
  be needed for some problems. What would we lose by abandoning the HMM
  framework, and what might we gain?
\end{enumerate}

\section{References}\label{references-8}

\chapter{Variant Calling and Filtering with
GATK}\label{variant-calling-and-filtering-with-gatk}

\section{From Aligned Reads to Biological Variants}\label{sec-ch10}

The journey from raw sequencing reads to biological insights requires
multiple processing steps, each building on the previous work. In
Chapter 6, you learned how short-read aligners like BWA use
sophisticated indexing strategies to map billions of reads to reference
genomes, producing BAM files that represent the best estimate of where
each sequenced fragment originated in the genome. These aligned reads,
however, are not the final product of a genome analysis pipeline.
Rather, they serve as the raw material for variant calling---the process
of identifying positions where sequenced individuals differ from the
reference genome or from each other. Understanding this distinction is
essential: alignment tells us where reads came from, while variant
calling tells us how those reads differ from expectation, and those
differences are often the biological signal we
seek\textsuperscript{37,38}.

\begin{tcolorbox}[enhanced jigsaw, rightrule=.15mm, leftrule=.75mm, colbacktitle=quarto-callout-note-color!10!white, colframe=quarto-callout-note-color-frame, toprule=.15mm, opacitybacktitle=0.6, title=\textcolor{quarto-callout-note-color}{\faInfo}\hspace{0.5em}{Figure 10.1 -- GATK Best Practices Workflow}, left=2mm, colback=white, bottomtitle=1mm, toptitle=1mm, bottomrule=.15mm, coltitle=black, titlerule=0mm, breakable, opacityback=0, arc=.35mm]

\begin{figure}[H]

\centering{

\includegraphics[width=0.9\linewidth,height=\textheight,keepaspectratio]{images/Figure_10.1_Workflow.png}

}

\caption{\label{fig-gatk-workflow2}GATK Best Practices Workflow}

\end{figure}%

This figure summarizes the end-to-end workflow you will implement in
this course, starting from raw Illumina reads and ending with a filtered
VCF that can be explored in IGV. The upper portion shows the familiar
pre-processing steps from earlier chapters: raw data are evaluated with
FastQC, trimmed to remove adapters and low-quality bases, aligned to a
reference genome with BWA, and converted into coordinate-sorted BAM
files using SAMtools. The lower portion highlights the GATK best
practices pipeline that this chapter focuses on: duplicate reads are
marked, quality control is collected on the resulting BAM files,
variants are called with HaplotypeCaller in reference-confidence mode to
produce gVCFs, joint genotyping (where applicable) produces a cohort
VCF, and hard-filtering steps refine this callset before variants are
examined in IGV. Created with BioRender.com.

\end{tcolorbox}

The Genome Analysis Toolkit, commonly known as GATK, has become the
\emph{de facto} standard for germline variant discovery in human
genomics and is widely applied to model organisms as well (this workflow
is shown above in Figure~\ref{fig-gatk-workflow2}, but see also
Figure~\ref{fig-gatk-workflow}). Originally developed at the Broad
Institute as a MapReduce framework for parallelizing variant calling
across large cohorts, GATK implements a series of ``best practices''
workflows that have been refined through years of application to
projects such as the 1000 Genomes Project and The Cancer Genome
Atlas\textsuperscript{37,38}. These workflows balance sensitivity---the
ability to discover real variants---with specificity---the ability to
avoid calling sequencing errors or alignment artifacts as true
biological variation. In this chapter, you will learn the logic behind
GATK's approach to variant calling, beginning with post-alignment
quality improvements, moving through the mechanics of the
HaplotypeCaller algorithm, and concluding with the critical step of
variant filtering that separates high-confidence calls from
noise\textsuperscript{114}.

The yeast experimental evolution dataset introduced in Appendix C
provides a concrete example throughout this chapter. Recall that this
dataset represents pooled genomic DNA from a \emph{Saccharomyces
cerevisiae} population that experienced hundreds of generations of
selection, so the variants you will call reflect both standing genetic
variation present in the ancestral population and allele frequency
shifts driven by adaptation\textsuperscript{24}. Because this is a
pooled sample rather than a single diploid individual, some of the
assumptions built into GATK's germline variant calling model do not
apply perfectly, but the core workflow---post-processing, variant
calling, filtering, and visual inspection---remains the same and
provides essential training for human and biomedical genomics where
these methods were originally developed.

\begin{tcolorbox}[enhanced jigsaw, rightrule=.15mm, leftrule=.75mm, colbacktitle=quarto-callout-note-color!10!white, colframe=quarto-callout-note-color-frame, toprule=.15mm, opacitybacktitle=0.6, title=\textcolor{quarto-callout-note-color}{\faInfo}\hspace{0.5em}{Lab connectionsL Hands-on Experience with GATK Best Practices}, left=2mm, colback=white, bottomtitle=1mm, toptitle=1mm, bottomrule=.15mm, coltitle=black, titlerule=0mm, breakable, opacityback=0, arc=.35mm]

\textbf{\href{Lab07.qmd}{Lab 7}}: You will use the Alabama Supercomputer
to run the full GATK variant calling workflow on the yeast dataset
SRR1693691, starting from the aligned BAM file you generated in Lab 4.
You will mark duplicates, call variants with HaplotypeCaller, apply hard
filters, and generate quality metrics using VCFtools and MultiQC. By
comparing unfiltered and filtered VCFs, you will see firsthand how
filtering parameters affect the trade-off between sensitivity and
specificity.

\end{tcolorbox}

\section{Post-Alignment Processing: Preparing BAM Files for Variant
Calling}\label{post-alignment-processing-preparing-bam-files-for-variant-calling}

The BAM file produced by BWA represents the aligner's best guess about
where each read originated, but it is not yet ready for variant calling.
Post-alignment processing steps refine this initial alignment by
identifying and marking artifacts that can mislead downstream analyses,
and by adjusting quality scores to better reflect the true probability
of error at each base\textsuperscript{38}. These steps do not change
read positions or CIGAR strings---the reads remain where BWA placed
them---but they add metadata and recalibrate scores so that the variant
caller can make better-informed decisions about which differences
between reads and reference are likely to be real biological variants
versus technical noise as shown visually in
Figure~\ref{fig-yeast-igv-overview}. Recall in Chapter 3, that even
using stringent quality score cutoffs of 40, there are still tens of
thousands of potential sequencing errors present in a typical genome
sequencing dataset. These next steps are critical in providing
additional evaluation to detect biological signal through the many
possible sequencing errors present.

\begin{tcolorbox}[enhanced jigsaw, rightrule=.15mm, leftrule=.75mm, colbacktitle=quarto-callout-note-color!10!white, colframe=quarto-callout-note-color-frame, toprule=.15mm, opacitybacktitle=0.6, title=\textcolor{quarto-callout-note-color}{\faInfo}\hspace{0.5em}{Figure 10.2 -- IGV reveals need for Post-Alignment Processing}, left=2mm, colback=white, bottomtitle=1mm, toptitle=1mm, bottomrule=.15mm, coltitle=black, titlerule=0mm, breakable, opacityback=0, arc=.35mm]

\begin{figure}[H]

\centering{

\includegraphics[width=1\linewidth,height=\textheight,keepaspectratio]{images/Figure_10.2_Pre-processing.jpg}

}

\caption{\label{fig-yeast-igv-overview}Post-alignment processing}

\end{figure}%

This IGV view shows a broad 23 kb region of the sacCer3 yeast genome
from the pooled experimental evolution sample introduced in Appendix C.
The top gray histogram track reports depth of coverage across the
window, illustrating how read depth fluctuates but remains generally
high enough for confident variant calling. The lower track displays
individual reads as horizontal bars colored by base identity, revealing
dense clusters of mismatches that correspond to SNPs and small indels
scattered across the region. This global perspective offers insight into
just how much noise there is in the initial BAM output and how many
putative variant sites need to be evaluated by the GATK pipeline.

\end{tcolorbox}

\subsection{Marking Duplicate Reads}\label{marking-duplicate-reads}

Nearly all second-generation sequencing platforms include a PCR
amplification step during library preparation, and this amplification
can create multiple copies of the same original DNA
molecule\textsuperscript{37}. When these PCR duplicates are sequenced,
they appear as independent reads with identical or nearly identical
start and end coordinates on the reference genome, but they do not
represent independent sampling of the genome---they are replicates of a
single molecule. If these duplicates are treated as independent
observations during variant calling, they can artificially inflate our
confidence in alleles that may actually stem from a single sequencing
error that was copied during PCR. This is particularly problematic at
low-coverage sites, where a single PCR-duplicated molecule carrying an
error could dominate the read pileup and cause a false-positive variant
call\textsuperscript{114}.

\begin{tcolorbox}[enhanced jigsaw, rightrule=.15mm, leftrule=.75mm, colbacktitle=quarto-callout-note-color!10!white, colframe=quarto-callout-note-color-frame, toprule=.15mm, opacitybacktitle=0.6, title=\textcolor{quarto-callout-note-color}{\faInfo}\hspace{0.5em}{Figure 10.3 -- Why Mark Duplicate Reads?}, left=2mm, colback=white, bottomtitle=1mm, toptitle=1mm, bottomrule=.15mm, coltitle=black, titlerule=0mm, breakable, opacityback=0, arc=.35mm]

\begin{figure}[H]

\centering{

\includegraphics[width=0.8\linewidth,height=\textheight,keepaspectratio]{images/Figure_10.3_MarkDups.png}

}

\caption{\label{fig-mark-duplicates}Why Mark Duplicate Reads?}

\end{figure}%

This figure illustrates why duplicate reads must be identified and
marked before variant calling. On the left, a single library molecule
has been over-amplified during PCR, producing many reads that share
identical start and end positions in the alignment. A sequencing error
(red X) present in the original molecule is propagated to all
duplicates, creating the illusion of strong support for a false variant.
MarkDuplicates flags these non-independent observations so that only one
representative read contributes to genotype likelihood calculations,
mitigating the impact of library artifacts on downstream analyses. In
practical terms, this step prevents overconfident false positives at
low-coverage sites and provides a useful summary metric---the percentage
of duplicate reads---for evaluating library complexity. Adapted from
GATK training materials and modified with the use of BioRender AI Tools.

\end{tcolorbox}

Duplicate marking algorithms identify clusters of reads that likely
arose from the same original fragment and flag all but one
representative read as duplicates. The most common approach, implemented
in both Picard MarkDuplicates and GATK MarkDuplicates, examines read
pairs with identical 5′ start coordinates on both mates. For paired-end
data, this means looking for pairs of reads where both the forward and
reverse read map to the exact same genomic positions, which is unlikely
to occur by chance if the reads truly represent independent DNA
molecules\textsuperscript{37}. Optical duplicates---clusters of reads
that arise from a single cluster on the flow cell being incorrectly
called as multiple clusters---are also flagged by examining pixel
distances in the original sequencing images, though this requires access
to tile and coordinate metadata in the FASTQ headers.

Crucially, GATK best practices recommend \emph{marking} duplicates
rather than \emph{removing} them from the BAM file\textsuperscript{38}.
Marking sets a flag in the BAM record (the 0x400 bit in the FLAG field)
that tells downstream tools to ignore these reads during variant
calling, but the reads remain in the file so that you can visually
inspect them in genome browsers like IGV and verify that the
duplicate-calling algorithm behaved reasonably. Very high duplication
rates (for example, more than 20 percent of reads flagged as duplicates)
can indicate problems with library complexity---perhaps the input DNA
was limited, or the library was over-amplified---but moderate
duplication is normal and expected\textsuperscript{37}. In the yeast
dataset, you can expect duplication rates around 5 to 10 percent, which
reflects standard Illumina library preparation.

After marking duplicates, quality control metrics from
\texttt{samtools\ flagstat} allow you to quantify the fraction of
duplicate reads across your samples. Plotting these percentages as a
histogram across multiple samples can reveal outliers that may need
additional scrutiny. In Lab 7, you will use GATK's MarkDuplicates tool
(previously done using Picard Tools, which are now part of GATK) and
compare flagstat reports before and after duplicate marking to see how
this step changes your \textbf{effective coverage}\textsuperscript{37}.

\begin{tcolorbox}[enhanced jigsaw, rightrule=.15mm, leftrule=.75mm, colbacktitle=quarto-callout-note-color!10!white, colframe=quarto-callout-note-color-frame, toprule=.15mm, opacitybacktitle=0.6, title=\textcolor{quarto-callout-note-color}{\faInfo}\hspace{0.5em}{Base Quality Score Recalibration (BQSR)}, left=2mm, colback=white, bottomtitle=1mm, toptitle=1mm, bottomrule=.15mm, coltitle=black, titlerule=0mm, breakable, opacityback=0, arc=.35mm]

Base quality scores, introduced in Chapter 3, are Phred-scaled
probabilities that a given base call is incorrect. These scores are
assigned by the sequencing instrument's base-calling software and are
supposed to reflect the true error rate, but systematic biases---related
to sequence context, machine cycle, or flow-cell chemistry---often cause
the reported quality scores to be systematically
miscalibrated\textsuperscript{38}. For example, base calls near
homopolymer runs or in high-GC contexts may be consistently assigned
quality scores that are too high (overly optimistic) or too low (overly
pessimistic), and quality scores often decline systematically with read
position as reagents are depleted over sequencing cycles. Additionally,
reads collected from separate runs on different lanes or at separate
times from pilots to larger scale sequencing projects should be treated
as separate experimental batches. BQSR helps to normalize quality scores
across batches for downstream variant calling.

GATK's Base Quality Score Recalibration (BQSR) procedure corrects these
systematic biases by building an empirical model of how quality scores
correlate with true error rates, conditional on various covariates such
as read position, sequence context, and the dinucleotide preceding each
base\textsuperscript{37,38}. The BQSR workflow has two passes: first,
the \textbf{BaseRecalibrator} tool scans the aligned BAM file and builds
a recalibration table by comparing observed mismatches to the reference
with reported quality scores, while masking out known variant sites (so
that real biological variation is not counted as sequencing error). In
the second pass, \textbf{ApplyBQSR} uses this recalibration table to
adjust the quality scores in place, writing out a new BAM file with
empirically corrected qualities\textsuperscript{152}.

In this course, we will \emph{not} require you to run BQSR on the yeast
dataset in Lab 7, primarily because BQSR requires a database of known
variants to mask true polymorphisms during the recalibration process,
and such databases are most complete for humans\textsuperscript{38}. For
model organisms like yeast, high-quality variant databases are less
comprehensive, and building one from your own data introduces
circularity (you need good variants to recalibrate, but you need good
quality scores to call good variants). Nonetheless, understanding BQSR
is essential for interpreting human genomics studies and for
appreciating why quality score accuracy matters in variant calling.

\end{tcolorbox}

\section{Variant Calling with GATK
HaplotypeCaller}\label{variant-calling-with-gatk-haplotypecaller}

With duplicates marked and optionally quality scores recalibrated, the
BAM file is now ready for variant discovery. GATK's
\textbf{HaplotypeCaller} is a sophisticated variant caller that differs
fundamentally from earlier, simpler approaches that examined each
genomic position independently\textsuperscript{37,38}. Instead of asking
``does this single position differ from the reference?'' HaplotypeCaller
takes a more holistic approach: it identifies regions of the genome that
show evidence of variation, performs local \emph{de novo} assembly of
the reads spanning those regions into candidate haplotypes, realigns
each read against these candidate haplotypes, and then uses a Bayesian
model to determine the most likely genotype at each variant site
conditional on the full haplotype context\textsuperscript{38}. This
haplotype-aware strategy dramatically improves accuracy, particularly
for insertions and deletions (indels), which are notoriously difficult
to call correctly because they disrupt local read alignment and create
characteristic error patterns in naive pileup-based
callers\textsuperscript{114}.

\begin{tcolorbox}[enhanced jigsaw, rightrule=.15mm, leftrule=.75mm, colbacktitle=quarto-callout-note-color!10!white, colframe=quarto-callout-note-color-frame, toprule=.15mm, opacitybacktitle=0.6, title=\textcolor{quarto-callout-note-color}{\faInfo}\hspace{0.5em}{Figure 10.4 -- HaplotypeCaller Local Assembly and Genotyping}, left=2mm, colback=white, bottomtitle=1mm, toptitle=1mm, bottomrule=.15mm, coltitle=black, titlerule=0mm, breakable, opacityback=0, arc=.35mm]

\begin{figure}[H]

\centering{

\includegraphics[width=0.9\linewidth,height=\textheight,keepaspectratio]{images/Figure_10.4_HaplotypeCaller.png}

}

\caption{\label{fig-haplotypecaller-workflow}HaplotypeCaller Local
Assembly and Genotyping}

\end{figure}%

This schematic decomposes the HaplotypeCaller algorithm into its major
conceptual stages. First, the tool scans the genome to identify
ActiveRegions where the pileup of reads shows evidence of divergence
from the reference sequence. Within each ActiveRegion, reads are
assembled into a small De Bruijn graph to generate a finite set of
candidate haplotypes that might explain the observed mismatches and
indels. HaplotypeCaller then uses a pair-HMM to compute, for each read,
the likelihood that it was generated from each candidate haplotype given
the observed bases and their quality scores. Finally, these likelihoods
are combined into genotype likelihoods for each possible genotype, and
the most probable genotype is chosen while variant-level annotations are
recorded for filtering. This haplotype-aware strategy improves accuracy
for indels and complex variants compared to simple pileup-based callers.
Adapted from Broad Institute GATK tutorial figures; reproduced for
educational use.

\end{tcolorbox}

\subsection{ActiveRegions and Local
Assembly}\label{activeregions-and-local-assembly}

HaplotypeCaller does not attempt to reassemble every base of the genome.
Instead, it scans the BAM file to identify
\textbf{ActiveRegions}---intervals where the reads show sufficient
evidence of variation compared to the reference that local reassembly is
warranted\textsuperscript{38}. Evidence for variation can come from
mismatches, soft-clipped bases, or indels in the initial BWA alignment.
Once an ActiveRegion is identified, HaplotypeCaller extracts all reads
overlapping that region and performs a localized \emph{de novo} assembly
using a De Bruijn graph, much like the genome assemblers discussed in
Chapter 6\textsuperscript{38}. This produces a set of candidate
haplotypes---plausible sequences that could explain the observed reads.
Typically, one of these haplotypes will be the reference sequence
itself, and others will carry one or more variants (SNPs, insertions, or
deletions).

HaplotypeCaller then realigns every read in the ActiveRegion against
each candidate haplotype using a pair-HMM (pair Hidden Markov Model)
likelihood calculation\textsuperscript{38}. This step asks: ``what is
the probability of observing this read sequence and its quality scores,
given that the true underlying sequence is haplotype H?'' The pair-HMM
naturally accounts for sequencing errors and small indels, and by
evaluating likelihoods across all candidate haplotypes, HaplotypeCaller
can distinguish reads that support different alleles even when they span
complex or compound variants\textsuperscript{37,38}.

After computing read likelihoods for each candidate haplotype,
HaplotypeCaller applies a Bayesian genotyping model to call the most
likely genotype (for a diploid sample, this means calling homozygous
reference, heterozygous, or homozygous alternate) at each variant
position\textsuperscript{38}. The output for each site in the VCF file
includes not only the called genotype, but also a wealth of
annotations---read depth, allele depth, mapping quality, strand bias,
and more---that will be essential for filtering in the next step. For
pooled population samples like the yeast dataset, the diploid assumption
does not strictly apply (you have an unknown mixture of haplotypes from
many cells), but HaplotypeCaller can still detect variants and will
report them as heterozygous calls when multiple alleles are
present\textsuperscript{24}.

\subsection{Reference Confidence Mode and gVCF
Files}\label{reference-confidence-mode-and-gvcf-files}

One challenge in variant calling is distinguishing true variant sites
from regions where no variant was called simply because coverage was too
low to make a confident call. The classic expression ``absence of
evidence is not evidence of absence'' captures this point perfectly. A
standard VCF file lists only sites where a variant was detected, which
means that absence of a variant call is ambiguous: it could mean the
site matches the reference, or it could mean we lacked data to evaluate
the site\textsuperscript{153}. This ambiguity becomes a major problem in
joint calling scenarios, where you want to combine variant calls from
multiple samples to improve power and accuracy by sharing information
across individuals\textsuperscript{153,154}.

GATK addresses this with \textbf{reference confidence mode}, enabled by
the \texttt{-ERC\ GVCF} flag (or \texttt{-\/-emit-ref-confidence\ GVCF}
in GATK4) when running HaplotypeCaller\textsuperscript{153}. In this
mode, HaplotypeCaller outputs a \textbf{gVCF} (genomic VCF) file that
lists not only variant sites but also blocks of reference-matching sites
along with a confidence estimate for each
block\textsuperscript{154,155}. These reference blocks indicate ``we
examined this region and found no evidence of variation, with this level
of confidence,'' which allows downstream joint genotyping tools to
distinguish confidently homozygous-reference sites from sites that were
simply not evaluated\textsuperscript{155}.

\begin{tcolorbox}[enhanced jigsaw, rightrule=.15mm, leftrule=.75mm, colbacktitle=quarto-callout-note-color!10!white, colframe=quarto-callout-note-color-frame, toprule=.15mm, opacitybacktitle=0.6, title=\textcolor{quarto-callout-note-color}{\faInfo}\hspace{0.5em}{Figure 10.5 -- Reference Confidence reveals clinical relavance}, left=2mm, colback=white, bottomtitle=1mm, toptitle=1mm, bottomrule=.15mm, coltitle=black, titlerule=0mm, breakable, opacityback=0, arc=.35mm]

\begin{figure}[H]

\centering{

\includegraphics[width=0.5\linewidth,height=\textheight,keepaspectratio]{images/Figure_10.5_RefConfidence.png}

}

\caption{\label{fig-reference-confidence}Why Reference Confidence Mode
Matters?}

\end{figure}%

This figure shows two patient VCFs produced by HaplotypeCaller in
reference-confidence mode at two variants in the APOE gene region, which
is associated with Alzheimer's Disease. A standard VCF lists only
positions where a variant was called, leaving ambiguity about whether
other sites truly match the reference or were simply not evaluated. In
contrast, reference confidence mode provides additional context for
homozygous-reference genotypes in the genome. The gVCF format enables
joint genotyping workflows by preserving information about both variant
and confidently non-variant sites for each sample, allowing downstream
tools like GenotypeGVCFs to distinguish true reference calls from
missing data. In a clinical context, this distinction is key in
reporting negative results to a potential patient about their risk for
Alzheimer's. In this toy example, we can report to the first patient
they have little risk based on their genotype at these two sites, but
for the second patient, we really cannot report anything due to a lack
of evidence in this region. Adapted from GATK reference confidence
tutorials; reproduced for educational purposes.

\end{tcolorbox}

Producing a gVCF is now the recommended first step in GATK's germline
variant calling workflow, even for single-sample projects, because it
preserves maximum information for future reanalysis or for adding new
samples later\textsuperscript{153,155}. The gVCF files can be quite
large because they cover the entire genome, but GATK provides tools to
consolidate or database multiple gVCFs efficiently in preparation for
joint genotyping\textsuperscript{153}.

\begin{tcolorbox}[enhanced jigsaw, rightrule=.15mm, leftrule=.75mm, colbacktitle=quarto-callout-note-color!10!white, colframe=quarto-callout-note-color-frame, toprule=.15mm, opacitybacktitle=0.6, title=\textcolor{quarto-callout-note-color}{\faInfo}\hspace{0.5em}{Figure 10.6 -- gVCF Blocks and Variant Records}, left=2mm, colback=white, bottomtitle=1mm, toptitle=1mm, bottomrule=.15mm, coltitle=black, titlerule=0mm, breakable, opacityback=0, arc=.35mm]

\begin{figure}[H]

\centering{

\includegraphics[width=0.85\linewidth,height=\textheight,keepaspectratio]{images/Figure_10.6_gVCF.png}

}

\caption{\label{fig-gvcf-blocks}Difference between VCF and gVCF file
formats}

\end{figure}%

This figure zooms in on how gVCF records represent continuous genomic
sequence for a single sample. Rather than writing an explicit record for
every base, HaplotypeCaller groups consecutive homozygous-reference
positions with similar confidence into non-variant blocks, each
annotated with depth and genotype quality for that interval. When a
variant is encountered---such as a SNP or small indel---the block is
interrupted by a standard variant record that includes alternate alleles
and detailed annotations. During joint genotyping, these blocks allow
GenotypeGVCFs to reconstruct site-by-site genotype likelihoods across
many samples without inflating file size, efficiently distinguishing
true reference positions from low-confidence or uncovered sites. Adapted
from Broad Institute GATK documentation on gVCF format; used here for
instruction.

\end{tcolorbox}

\subsection{Joint Genotyping with
GenotypeGVCFs}\label{joint-genotyping-with-genotypegvcfs}

Once you have gVCF files for one or more samples, the next step is
\textbf{joint genotyping}. This is performed by the
\textbf{GenotypeGVCFs} tool, which takes one or more gVCF inputs and
produces a standard multisample VCF that includes genotype calls at
every site where any sample showed evidence of
variation\textsuperscript{153,155}. Joint genotyping improves calling
accuracy in two key ways: first, it allows the statistical model to
borrow strength across samples, so that marginal evidence of a variant
in one individual can be confirmed or refuted by looking at the same
site in related individuals. Second, it provides a unified
representation of variation across the cohort, which is essential for
population genetic analyses, case-control studies, or trio-based
analyses\textsuperscript{153,154}.

For large cohorts (hundreds to thousands of samples), GATK4 introduced
\textbf{GenomicsDBImport}, which consolidates gVCFs into a scalable
tile-based database structure that enables efficient querying across
samples\textsuperscript{155}. GenotypeGVCFs can then operate directly on
this database, avoiding the need to load and parse thousands of
individual gVCF files for each genomic interval\textsuperscript{155}.
For smaller projects like your yeast dataset with a single sample, you
can skip the GenomicsDBImport step and run GenotypeGVCFs directly on the
gVCF, which will convert the per-base reference confidence blocks into a
standard VCF listing only variant sites and homozygous-reference sites
that passed thresholds\textsuperscript{155}.

The output of GenotypeGVCFs is a raw, unfiltered VCF that contains all
putative variants GATK identified, including many false positives. The
next critical step is filtering this raw callset to retain only
high-confidence variants.

\section{Variant Filtering and Quality
Control}\label{variant-filtering-and-quality-control}

Raw variant calls from HaplotypeCaller, while sophisticated, still
contain a mixture of true biological variants and technical artifacts
arising from sequencing errors, mapping errors, and alignment
ambiguities\textsuperscript{38,114}. Effective filtering is the key to
transforming a raw VCF into a high-quality callset suitable for
downstream biological interpretation. GATK offers two main filtering
strategies: \textbf{Variant Quality Score Recalibration (VQSR)}, a
machine-learning approach that models the relationship between variant
annotations and true-positive status; and \textbf{hard filtering}, a
simpler rule-based approach that applies explicit thresholds to
individual annotation metrics\textsuperscript{156,157}. For this course,
you will apply hard filtering in Lab 7, but understanding VQSR provides
important context for why certain annotations matter and how they are
used in large-scale human genomics projects.

\subsection{Variant Quality Score Recalibration
(VQSR)}\label{variant-quality-score-recalibration-vqsr}

VQSR is designed for projects with large cohorts (typically hundreds to
thousands of samples) and access to high-quality training resources such
as HapMap, 1000 Genomes, or dbSNP\textsuperscript{38,156}. The approach
uses machine learning (specifically, a Gaussian mixture model) to learn
the annotation profiles of known true-positive variants from the
training set, and then applies this model to score all variants in your
callset with a new quality metric called the \textbf{VQSLOD} (Variant
Quality Score Log-Odds)\textsuperscript{156}. Sites with high VQSLOD
scores resemble the training variants and are likely true positives;
sites with low or negative VQSLOD scores have annotation profiles
consistent with artifacts\textsuperscript{156}.

The power of VQSR lies in its ability to integrate multiple annotation
dimensions---such as read depth, mapping quality, strand bias, allele
balance, and proximity to other variants---into a single calibrated
score without requiring you to manually choose thresholds for each
metric\textsuperscript{38,156}. By adjusting a \textbf{sensitivity
cutoff} (for example, retaining the top 99 percent of SNPs that overlap
with HapMap), you can tune the trade-off between sensitivity (keeping
real variants) and specificity (excluding false positives) in a
principled way\textsuperscript{156}.

Because VQSR requires large training sets of known variants and performs
best with hundreds or thousands of samples, it is not well suited to
small projects or non-model organisms\textsuperscript{156,157}. The
yeast dataset, with a single pooled sample and limited high-confidence
variant databases, falls into this category. For such cases, hard
filtering provides a practical alternative.

\begin{tcolorbox}[enhanced jigsaw, rightrule=.15mm, leftrule=.75mm, colbacktitle=quarto-callout-note-color!10!white, colframe=quarto-callout-note-color-frame, toprule=.15mm, opacitybacktitle=0.6, title=\textcolor{quarto-callout-note-color}{\faInfo}\hspace{0.5em}{Variant Quality Score Recalibration (VQSR)}, left=2mm, colback=white, bottomtitle=1mm, toptitle=1mm, bottomrule=.15mm, coltitle=black, titlerule=0mm, breakable, opacityback=0, arc=.35mm]

VQSR is ideal for large cohorts and model organisms with comprehensive
variant databases. The approach builds a statistical model from known
high-quality variants (for example, HapMap or 1000 Genomes for humans)
and uses that model to assign a calibrated quality score (VQSLOD) to
each variant in your dataset based on its annotation
profile\textsuperscript{156}. You can then filter at different
sensitivity levels---retaining, for example, 99.5 percent or 99.0
percent of the known true positives---depending on whether your project
prioritizes sensitivity or specificity\textsuperscript{156}.

VQSR processes SNPs and indels separately because their annotation
distributions differ significantly (for instance, indels tend to have
different strand-bias and read-position-bias profiles than
SNPs)\textsuperscript{156}. Each variant type requires its own training
set and Gaussian mixture model, and the final filtered VCF merges the
results\textsuperscript{156,157}.

While VQSR represents the state-of-the-art for human genomics, it is not
always appropriate: small sample sizes give the model insufficient data
to learn accurate distributions, and projects on non-model organisms
often lack high-confidence training resources\textsuperscript{156,157}.
In these settings, hard filtering offers a transparent, interpretable
alternative that you will apply in Lab 7.

\end{tcolorbox}

\subsection{Hard Filtering: Thresholds and
Rationale}\label{hard-filtering-thresholds-and-rationale}

Hard filtering applies explicit cutoff values to individual variant
annotations, marking variants that fail one or more filters as
low-quality and retaining only those that pass all
filters\textsuperscript{114,157}. GATK provides the
\textbf{VariantFiltration} tool, which adds filter tags to the FILTER
column of the VCF without removing variants from the file, so you can
inspect borderline calls and adjust thresholds
iteratively\textsuperscript{158}. The key challenge in hard filtering is
choosing appropriate thresholds for your data, which requires
understanding what each annotation measures and how it relates to
variant quality.

\textbf{Quality by Depth (QD)} is one of the most important filtering
metrics. QD normalizes the variant quality score (QUAL) by the read
depth at the site, so that high-quality calls are not penalized simply
because they occur in low-coverage regions\textsuperscript{114,157}. A
QD threshold of 2.0 is often recommended for both SNPs and indels;
variants with QD below 2 typically reflect noisy regions where many weak
reads contribute ambiguous signal\textsuperscript{158}.

\textbf{Mapping Quality (MQ)} measures the average root-mean-square
mapping quality of reads supporting the variant\textsuperscript{157}.
Low MQ values (for example, MQ below 40) indicate that reads at this
site aligned poorly or aligned to multiple locations, suggesting the
variant may arise from mismapping rather than true biological
variation\textsuperscript{114}. For SNPs, GATK recommends MQ greater
than or equal to 40; for indels, the threshold may be slightly relaxed
because indels can locally disrupt mapping even when the broader region
is unique\textsuperscript{158}.

\textbf{Strand Bias} metrics, including \textbf{FS (FisherStrand)} and
\textbf{SOR (StrandOddsRatio)}, detect situations where the variant
allele is supported predominantly by reads from one
strand\textsuperscript{114,157}. True variants should appear on both
forward and reverse strands in roughly equal proportions (modulo
stochastic sampling); strong strand bias often indicates systematic
sequencing errors or alignment artifacts that affect one strand more
than the other\textsuperscript{114}. GATK recommends FS less than 60.0
for SNPs (less than 200.0 for indels) and SOR less than 3.0 for SNPs
(less than 10.0 for indels)\textsuperscript{158}.

\textbf{Read Position Rank Sum (ReadPosRankSum)} evaluates whether
variant alleles tend to occur at the ends of reads versus in the middle,
which can indicate sequencing artifacts or incorrect
alignments\textsuperscript{157}. Negative values indicate the variant
allele is enriched at read ends. GATK suggests ReadPosRankSum greater
than -8.0 for both SNPs and indels\textsuperscript{158}.

\textbf{MQRankSum (Mapping Quality Rank Sum)} compares the mapping
qualities of reads supporting the reference allele versus the alternate
allele\textsuperscript{157}. Negative values mean reads supporting the
alternate allele have lower mapping quality, which may indicate the
variant is an artifact. A threshold of MQRankSum greater than -12.5 is
commonly used\textsuperscript{158}.

In practice, you should begin by examining the distributions of these
metrics in your raw VCF using VCFtools or custom scripts, plotting
density histograms to identify natural break points and
outliers\textsuperscript{114}. This exploratory step ensures that
thresholds are \textbf{tailored to your dataset} and sequencing
platform, rather than applying default values that may not suit your
data. In Lab 7, you will generate such plots for the yeast dataset and
experiment with different threshold combinations, comparing the number
of retained variants and their overall quality metrics to see how
filtering decisions impact the final callset.

Because SNPs and indels have different error profiles, GATK recommends
filtering them separately\textsuperscript{157}. You first use
\textbf{SelectVariants} to extract SNPs and indels into separate VCF
files, apply VariantFiltration with appropriate thresholds to each file,
and then optionally merge them back together, or continue with separate
filtered VCFs depending on your downstream
analyses\textsuperscript{158}.

\subsection{VCFtools for Quality Assessment and
Filtering}\label{vcftools-for-quality-assessment-and-filtering}

While GATK provides the tools to annotate and filter variants,
\textbf{VCFtools} is a complementary toolkit designed specifically for
manipulating, summarizing, and quality-checking VCF
files\textsuperscript{35}. VCFtools offers a wide range of utilities:
you can calculate per-site or per-sample depth, extract specific genomic
regions using BED files, compute allele frequencies, measure linkage
disequilibrium, and convert VCF files to other formats for use in
population genetics software\textsuperscript{35}. In the context of
variant filtering, VCFtools is particularly useful for generating
summary statistics and coverage reports that help you evaluate the
impact of filtering thresholds.

For example, the \texttt{-\/-depth} option in VCFtools computes mean
depth per individual across all sites in the VCF, while
\texttt{-\/-site-depth} reports total depth at each variant site across
all samples\textsuperscript{35}. Comparing these depth statistics
between raw and filtered VCFs shows whether filtering disproportionately
removes low-coverage sites (which is generally desirable) or also
discards well-covered sites (which may indicate overly stringent
thresholds). The \texttt{-\/-site-quality} option extracts the QUAL
scores for each site, and plotting these distributions can reveal how
filtering changes the quality profile of the
callset\textsuperscript{35}.

VCFtools can also directly filter VCF files by applying thresholds (for
instance, \texttt{-\/-minQ}, \texttt{-\/-minDP},
\texttt{-\/-max-missing}) and subsetting to specific regions or
individuals\textsuperscript{35}. While GATK's VariantFiltration is often
preferred because it preserves all variants and simply annotates them
with filter tags, VCFtools offers a lightweight alternative for quick
filtering and reformatting tasks, and its output files are easily loaded
into R for custom plotting and analysis\textsuperscript{35}.

In Lab 7, you will use VCFtools to compute coverage summaries, extract
filtered variants passing all quality thresholds, and generate
tab-delimited reports that you can visualize as density plots or
histograms. You will also use VCFtools to create BED files of
high-confidence variant regions for downstream IGV inspection.

\section{Genome Browsers: Visual Inspection of Variants in
IGV}\label{genome-browsers-visual-inspection-of-variants-in-igv}

Automated filtering can remove many false positives, but no filtering
strategy is perfect. Visual inspection of alignments at candidate
variant sites remains an essential validation step, particularly for
variants of special interest---those with potential functional impact,
those driving phenotypic differences, or those that fall in critical
genes\textsuperscript{159}. The \textbf{Integrative Genomics Viewer
(IGV)} developed at the Broad Institute is the most widely used tool for
this purpose\textsuperscript{159}. IGV allows you to load reference
genomes, BAM alignment files, VCF variant calls, and genome annotation
tracks (such as gene models in GFF format) into a unified graphical
interface where you can zoom to any genomic locus and examine the read
evidence supporting or refuting each variant\textsuperscript{159}.

I have recorded a brief
\href{https://auburn.hosted.panopto.com/Panopto/Pages/Viewer.aspx?id=51da88a9-e569-4453-b414-b29d0125a0ac}{video
tutorial of IGV} and describe the use of this tool below in detail. You
will also get hands-on experience using IGV in Lab 9.

\subsection{Loading Data into IGV}\label{loading-data-into-igv}

To inspect variants in IGV, you first load the reference genome (in
FASTA format) along with its index (.fai file), which tells IGV how to
navigate the chromosomes\textsuperscript{159}. Next, you load one or
more BAM files---typically both the pre-processing BAM and the
duplicate-marked BAM---so you can compare alignment quality before and
after post-processing steps. IGV requires BAM files to be
coordinate-sorted and indexed (.bai files must be present), conditions
you satisfied during the alignment workflow in Lab
4\textsuperscript{159}. Finally, you load the filtered VCF file, which
IGV will display as a track showing the position and genotype of each
called variant\textsuperscript{159}.

IGV's interface presents a genome-wide overview at the top, showing the
positions of all loaded features across chromosomes, and a zoomed-in
view below that displays individual reads and annotations at base-pair
resolution once you zoom in sufficiently\textsuperscript{159}. Track
display options allow you to color reads by various attributes: pairing
status, mapping quality, read strand, or whether reads are marked as
duplicates\textsuperscript{159}. For variant inspection, coloring by
read strand is particularly useful because it immediately reveals strand
bias---if all reads supporting the alternate allele come from one
strand, the variant is likely an artifact\textsuperscript{114,159}.

\begin{tcolorbox}[enhanced jigsaw, rightrule=.15mm, leftrule=.75mm, colbacktitle=quarto-callout-note-color!10!white, colframe=quarto-callout-note-color-frame, toprule=.15mm, opacitybacktitle=0.6, title=\textcolor{quarto-callout-note-color}{\faInfo}\hspace{0.5em}{Figure 10.7 -- Inspecting a Yeast Coding SNP in IGV}, left=2mm, colback=white, bottomtitle=1mm, toptitle=1mm, bottomrule=.15mm, coltitle=black, titlerule=0mm, breakable, opacityback=0, arc=.35mm]

\begin{figure}[H]

\centering{

\includegraphics[width=1\linewidth,height=\textheight,keepaspectratio]{images/Figure_10.7_IGV.jpg}

}

\caption{\label{fig-yeast-igv-coding-snp}Using IGV to Inspect Putative
Functional Impact of Filtered Variants}

\end{figure}%

This IGV screenshot focuses on a single high-confidence SNP in a coding
region of the yeast genome. Two BAM tracks are shown: the upper track
corresponds to the original alignments, while the lower track displays
the duplicate-marked alignment used for variant calling, with reads
colored by strand to reveal any strand bias. At the highlighted
position, a vertical column of C bases indicates consistent support for
the alternate allele across many independent reads on both strands, with
uniform mapping and base quality. The gene annotation and translation
tracks at the bottom of the window show that this nucleotide change
alters a codon in the HXT14 open reading frame, leading to a specific
amino-acid substitution. This example illustrates how IGV can be used to
validate the technical quality of variant calls and immediately assess
their potential functional impact in the context of coding sequences.

\end{tcolorbox}

\subsection{Interpreting Read Evidence at Variant
Sites}\label{interpreting-read-evidence-at-variant-sites}

When you navigate to a variant site in IGV, you see a vertical alignment
or ``pileup'' of all reads covering that position\textsuperscript{159}.
Each read is displayed as a horizontal bar, with matches to the
reference shown in gray and mismatches highlighted in color (A=green,
C=blue, G=yellow/orange, T=red by default). Insertions appear as purple
I markers, and deletions as black dashes\textsuperscript{159}. The
height and density of colored bases at a given position immediately
convey the read support for each allele\textsuperscript{159}.

For high-quality SNPs, you expect to see the alternate allele
represented on both forward and reverse strands, distributed across
multiple reads that have high mapping quality (thick bars rather than
translucent, low-MQ bars), and supported by bases with high per-base
quality (IGV can optionally shade bases by quality
score)\textsuperscript{114,159}. For heterozygous calls in a diploid
sample, you expect roughly equal numbers of reads supporting reference
and alternate alleles, though stochastic sampling means small deviations
are normal\textsuperscript{38}.

Artifacts often exhibit characteristic patterns: PCR duplicates appear
as stacks of identical reads with identical start and end coordinates,
all supporting the same allele; strand-biased artifacts show the variant
allele only on forward or only on reverse reads; mapping artifacts
cluster in regions with low mapping quality or high mismatch rates
across all reads; and sequencing errors tend to occur at low-quality
bases near the ends of reads\textsuperscript{114,159}. By learning to
recognize these patterns, you develop an intuition for which variants
are trustworthy and which require additional scrutiny or experimental
validation.

\begin{tcolorbox}[enhanced jigsaw, rightrule=.15mm, leftrule=.75mm, colbacktitle=quarto-callout-note-color!10!white, colframe=quarto-callout-note-color-frame, toprule=.15mm, opacitybacktitle=0.6, title=\textcolor{quarto-callout-note-color}{\faInfo}\hspace{0.5em}{Lab connections: Completing the Yeast Workflow (Lab 9)}, left=2mm, colback=white, bottomtitle=1mm, toptitle=1mm, bottomrule=.15mm, coltitle=black, titlerule=0mm, breakable, opacityback=0, arc=.35mm]

In \textbf{\href{Lab09.qmd}{Lab 9} (Genome Visualization with IGV)}, you
will complete the end-to-end workflow introduced in
Figure~\ref{fig-gatk-workflow2} by moving from filtered variant calls to
biological interpretation. Starting with the filtered VCFs you generated
in \textbf{Lab 7}, you will load the sacCer3 reference genome, the
pre-filtering BAM (\texttt{SRR1693691.sorted.bam}), the duplicate-marked
BAM (\texttt{SRR1693691.markdup.bam}), and one or more filtered VCF
files into IGV. With these tracks side-by-side, you will navigate to
specific SNPs identified as high-confidence by your filtering criteria,
visually evaluate the read evidence supporting each variant, and decide
which calls you trust, which you question, and why.

A central goal of Lab 9 is to connect technical quality metrics to
biological meaning in the \emph{Saccharomyces cerevisiae} experimental
evolution dataset described in \href{appendixC.qmd}{Appendix C}. By the
end of the lab, you will have followed the yeast dataset all the way
from raw FASTQ files in Lab 4 through alignment, GATK variant calling,
filtering in Lab 7, and finally visual, function-focused interpretation
in IGV---mirroring the complete variant discovery workflow used in
modern genomics.

\end{tcolorbox}

\subsection{Using IGV to Evaluate Functional
Impact}\label{using-igv-to-evaluate-functional-impact}

Beyond verifying the technical quality of variant calls, IGV helps
assess biological significance by displaying gene annotation tracks
alongside alignments\textsuperscript{159}. If you load a GFF or GTF file
describing yeast gene models (such as those available from the
Saccharomyces Genome Database), IGV will show where variants fall
relative to exons, introns, and untranslated
regions\textsuperscript{159}. Variants that fall within coding exons are
candidates for nonsynonymous changes that alter protein sequence, while
variants in promoter regions or untranslated regions may affect gene
regulation\textsuperscript{159}.

IGV can translate coding regions and display amino acid sequences in a
separate track, allowing you to immediately see whether a coding variant
is synonymous (silent) or nonsynonymous (changes the amino
acid)\textsuperscript{159}. For pooled population samples like the yeast
dataset, functional variants are of particular interest because they may
underlie adaptive changes that rose to high frequency during the
evolution experiment\textsuperscript{24}. In \href{Lab09.qmd}{Lab 9},
you will use IGV to navigate to specific variants identified in your
filtered VCF, evaluate their technical quality, determine their
functional context, and consider hypotheses about their potential role
in adaptation.

\section{Beyond GATK: Alternative Variant Calling
Workflows}\label{beyond-gatk-alternative-variant-calling-workflows}

While GATK has become the standard for germline variant calling in human
genomics, it is not the only option, and for some applications or
organisms, alternative tools may be more
appropriate\textsuperscript{114}. Understanding the landscape of variant
callers and the trade-offs they embody is essential for making informed
choices in your own projects and for critically evaluating the methods
sections of published papers. Broadly, variant callers differ in their
algorithmic approach (pileup-based versus haplotype-aware versus
assembly-based), their optimization for particular data types (short
reads versus long reads, germline versus somatic), and their
computational requirements\textsuperscript{114}. The past decade has
seen extensive benchmarking efforts comparing these tools on real and
simulated datasets, revealing that no single caller dominates across all
metrics and that caller performance depends strongly on sequencing
depth, variant type, genomic context, and the population ancestry
represented in training data\textsuperscript{160--164}.

\begin{tcolorbox}[enhanced jigsaw, rightrule=.15mm, leftrule=.75mm, colbacktitle=quarto-callout-note-color!10!white, colframe=quarto-callout-note-color-frame, toprule=.15mm, opacitybacktitle=0.6, title=\textcolor{quarto-callout-note-color}{\faInfo}\hspace{0.5em}{Lab 8: Critical Evaluation of Variant Calling Benchmarks}, left=2mm, colback=white, bottomtitle=1mm, toptitle=1mm, bottomrule=.15mm, coltitle=black, titlerule=0mm, breakable, opacityback=0, arc=.35mm]

In \href{Lab08.qmd}{Lab 8}, you will work in small groups to critically
read and present recent benchmarking studies that compare variant
calling pipelines on diverse datasets and
organisms\textsuperscript{5--9}. Your task will be to extract key
findings, evaluate the study design and metrics used, identify which
tools performed best under which conditions, and synthesize practical
recommendations for future researchers. This exercise mirrors real-world
bioinformatics practice, where no single ``best'' tool exists and
informed decisions require understanding the literature, the data, and
the biological question.

\end{tcolorbox}

\subsection{Pileup-Based Callers: Samtools and
bcftools}\label{pileup-based-callers-samtools-and-bcftools}

The earliest generation of variant callers, including \textbf{samtools
mpileup} and \textbf{bcftools call}, operate on a simple per-position
pileup model: they examine each genomic position independently, count
the number of reads supporting each allele, and apply a statistical
model (often a Bayesian model assuming binomial or multinomial sampling)
to call genotypes\textsuperscript{114}. These tools are fast and require
minimal memory, making them suitable for large genomes and
high-throughput projects. However, they struggle with indels and with
complex variants (such as nearby SNPs that form haplotypes), because
each site is evaluated in isolation without considering the broader
sequence context\textsuperscript{114}.

Recent benchmarking in non-human species suggests that bcftools can
still perform competitively with GATK HaplotypeCaller for SNP discovery
under realistic coverage and error profiles, especially when calling
variants on high-quality, well-assembled
references\textsuperscript{165}. For projects where computational
resources are limited or where indel calling is not a priority,
samtools/bcftools therefore remains a viable option, particularly for
bacterial genomes or other small genomes where variants are sparse and
mostly SNPs\textsuperscript{114,165}. The bcftools suite also offers
powerful filtering and manipulation utilities that complement GATK's
VCFtools, and many researchers use bcftools for post-processing even if
they call variants with GATK\textsuperscript{114,165}.

\subsection{Haplotype-Aware Callers: FreeBayes, Platypus, and
DeepVariant}\label{haplotype-aware-callers-freebayes-platypus-and-deepvariant}

FreeBayes and Platypus are haplotype-aware variant callers that, like
GATK HaplotypeCaller, consider local haplotype structure when calling
variants\textsuperscript{114}. \textbf{FreeBayes} operates by
enumerating all possible haplotypes (combinations of alleles) within a
local window and computing the posterior probability of each haplotype
given the observed reads; it then reports the most probable genotypes at
each variant site\textsuperscript{114}. FreeBayes is known for its
ability to call complex variants and multi-allelic sites (where more
than two alleles are present at a locus), and it does not require a
separate local assembly step like HaplotypeCaller, which can make it
faster in some settings\textsuperscript{114}. Early benchmarking studies
found that FreeBayes and GATK showed complementary strengths, with
FreeBayes sometimes detecting variants missed by GATK in low-complexity
regions, though both tools exhibited platform-specific biases when
applied to data from different sequencing
instruments\textsuperscript{160,166}.

\textbf{Platypus} similarly uses local assembly and realignment but is
optimized for speed and low memory usage, making it attractive for
whole-genome calling on large cohorts or non-model
organisms\textsuperscript{114}. Platypus has been shown to produce
high-quality SNP and indel calls with lower computational cost than
GATK, though its performance relative to other callers depends on
coverage depth and the stringency of downstream
filtering\textsuperscript{161,162}. Both FreeBayes and Platypus produce
VCF output compatible with standard filtering and annotation pipelines,
and continue to be used in projects where GATK's computational demands
or specific assumptions (such as diploid genotypes) are
limiting\textsuperscript{114}.

A more recent addition to the haplotype-aware caller landscape is
\textbf{DeepVariant}, developed by Google and released in 2017, which
represents a fundamentally different approach: it uses deep
convolutional neural networks trained on labeled examples to classify
candidate variants from ``pileup images'' that encode read alignments,
base qualities, and mapping qualities as visual
tensors\textsuperscript{167}. DeepVariant has been trained on diverse
datasets including \textbf{Genome in a Bottle} reference samples and can
be retrained or fine-tuned for specific sequencing platforms, making it
particularly effective for PacBio HiFi and Oxford Nanopore data where
traditional callers struggle\textsuperscript{167}. Because DeepVariant
learns patterns from data rather than relying on hand-coded heuristics,
it often achieves state-of-the-art accuracy for both SNPs and indels,
though at the cost of requiring GPU resources for efficient execution
and lacking the interpretability of probabilistic
models\textsuperscript{167}.

\subsection{Ensemble and Voting-Based
Approaches}\label{ensemble-and-voting-based-approaches}

An emerging best practice, particularly in precision medicine and
clinical genomics, is to call variants with multiple independent tools
and take the intersection (or weighted consensus) of the resulting
callsets\textsuperscript{114}. The rationale is that different callers
have different strengths and failure modes: one caller may excel at SNPs
but miss indels, another may be sensitive to strand-biased artifacts,
and a third may handle low-complexity regions better. By requiring that
a variant be called by two or more independent algorithms, you can
dramatically reduce \textbf{false positives} while retaining most
\textbf{true positives}, a strategy that has been validated across
whole-genome, whole-exome, and targeted sequencing
projects\textsuperscript{160,161}.

Ensemble approaches require more computational effort (you run multiple
callers on the same data) and more sophisticated post-processing (you
must reconcile VCF files that may represent the same variant with
slightly different coordinates or alleles), but the payoff in callset
quality can be substantial\textsuperscript{114}. Simple voting
schemes---where a variant is retained if it appears in at least two out
of three callers---are common, but more sophisticated methods can weight
callers by their known precision and sensitivity for different variant
classes or use machine-learning meta-classifiers to integrate
annotations from all input callsets\textsuperscript{160}. Early studies
demonstrated that ensemble calls combining GATK, FreeBayes, and samtools
substantially reduced false-discovery rates compared to any single
caller, particularly for indels and variants in difficult genomic
regions\textsuperscript{160,162}. More recent work has extended these
principles to somatic variant calling in cancer genomics, where ensemble
methods combining tools like \textbf{Mutect2}, \textbf{VarScan}, and
\textbf{Strelka} are now considered best practice for high-confidence
mutation detection in tumor-normal pairs\textsuperscript{163}.

\subsection{The Importance of the Full Pipeline: Alignment,
Preprocessing, and Population
Context}\label{the-importance-of-the-full-pipeline-alignment-preprocessing-and-population-context}

Benchmarking studies over the past decade have consistently emphasized
that variant calling accuracy depends not only on the choice of caller
but on the entire upstream pipeline, including read alignment, duplicate
marking, base quality recalibration, and the choice of reference
genome\textsuperscript{160,163,164,166}. Different aligners (BWA-MEM,
Bowtie2, STAR for RNA-seq) can produce subtly different placements of
reads in repetitive or low-complexity regions, and these differences
propagate into variant calls\textsuperscript{166}. Similarly, the
decision to perform or skip BQSR, or the choice of known-variant
databases used during recalibration, can shift the
sensitivity-specificity trade-off in ways that interact with the
caller\textsuperscript{160,163}.

Population ancestry and reference genome choice add another layer of
complexity. Most variant callers and filtering models have been
optimized using data from individuals of European ancestry, and several
studies have shown that performance degrades when applied to genomes
from underrepresented populations, particularly in regions where the
reference genome is a poor match or where population-specific variants
are incorrectly flagged as low-quality\textsuperscript{163,164}.
Recognizing these biases and adapting pipelines---by using
population-matched training sets for machine-learning callers, adjusting
filter thresholds, or building population-specific references---is an
active area of research and an important consideration for equitable
genomic medicine\textsuperscript{163}.

\subsection{Long-Read and Structural Variant
Callers}\label{long-read-and-structural-variant-callers}

GATK HaplotypeCaller and most other tools discussed so far are optimized
for short-read data (100--300 bp Illumina reads)\textsuperscript{114}.
For long-read sequencing platforms like PacBio HiFi or Oxford Nanopore,
specialized callers are required that can handle the different error
profiles and leverage the increased read length\textsuperscript{167}.
\textbf{DeepVariant} has been successfully adapted for long reads and is
widely used for PacBio HiFi data, while \textbf{Clair3} and
\textbf{PEPPER-Margin-DeepVariant} offer state-of-the-art performance
for Oxford Nanopore sequencing\textsuperscript{167}. Long reads provide
richer information about haplotype structure and can span repetitive
regions that are inaccessible to short reads, enabling more accurate
calling of structural variants, copy-number variants, and phased
haplotypes\textsuperscript{167}.

Structural variant (SV) calling---detecting large deletions, insertions,
inversions, and translocations---requires entirely different algorithmic
approaches than SNP/indel calling, often relying on paired-end mapping
signatures (discordant pairs, split reads) or depth-of-coverage
signals\textsuperscript{114}. Tools like \textbf{Manta}, \textbf{Delly},
and \textbf{Lumpy} specialize in SV detection from short reads, while
long-read SV callers like \textbf{Sniffles} and \textbf{cuteSV} leverage
the ability of long reads to span entire structural
variants\textsuperscript{114}. Understanding when and how to apply these
specialized tools is part of the broader skill set required for modern
genome analysis.

\section{Reflection Questions}\label{reflection-questions-8}

\begin{enumerate}
\def\labelenumi{\arabic{enumi}.}
\item
  \textbf{Interpreting duplication rates}: Suppose you mark duplicates
  on three samples from the same sequencing run and observe duplication
  rates of 5 percent, 8 percent, and 35 percent. What biological or
  technical explanations might account for the outlier sample, and what
  quality control steps would you take before proceeding with variant
  calling?
\item
  \textbf{The value of reference confidence}: Explain why a gVCF file
  that includes blocks of homozygous-reference calls with confidence
  estimates is more informative than a standard VCF that lists only
  variant sites. Give an example of a downstream analysis (such as
  estimating missing heritability or validating negative results) where
  this additional information is essential.
\item
  \textbf{Filtering trade-offs}: You apply hard filters to a raw VCF and
  retain 85 percent of SNPs and 60 percent of indels. Discuss whether
  this suggests your filters are too strict, appropriately balanced, or
  too lenient, and describe what additional information you would need
  to make a definitive judgment.
\item
  \textbf{Strand bias as a diagnostic}: Describe a scenario in which a
  true biological variant might exhibit strand bias, and explain how you
  would distinguish such a variant from a technical artifact that also
  shows strand bias. What additional lines of evidence would you seek in
  IGV or through external validation?
\item
  \textbf{When to use VQSR versus hard filtering}: Compare the strengths
  and limitations of VQSR and hard filtering. Under what circumstances
  would you choose VQSR despite its added complexity, and when would
  hard filtering be not only simpler but actually preferable?
\item
  \textbf{Evaluating ensemble calling}: Imagine you call variants on the
  same sample using GATK HaplotypeCaller, FreeBayes, and Platypus, and
  observe that 70 percent of SNPs are called by all three tools, 20
  percent by two tools, and 10 percent by only one tool. How would you
  decide which variants to retain for downstream analysis, and what
  factors (such as annotation scores or functional impact) would
  influence your decision?
\end{enumerate}

\section{Problem Set}\label{problem-set}

\begin{enumerate}
\def\labelenumi{\arabic{enumi}.}
\item
  Duplicate Marking and Coverage Impact: You have aligned 50 million
  paired-end 150 bp reads from a diploid human exome to the reference
  genome. Before marking duplicates, \texttt{samtools\ flagstat} reports
  that 48 million reads are properly paired. After running GATK
  MarkDuplicates, you find that 15 percent of reads are flagged as
  duplicates.

  \begin{enumerate}
  \def\labelenumii{\alph{enumii}.}
  \item
    Calculate the effective coverage over the exome (assume the exome is
    50 Mb) before and after accounting for duplicates.
  \item
    Explain why GATK recommends marking but not removing duplicates, and
    describe one situation in which you might choose to remove
    duplicates instead.
  \item
    If you observe 30 percent duplication in one sample but 10 percent
    in all others from the same sequencing run, propose two possible
    explanations and outline how you would investigate which explanation
    is correct.
  \end{enumerate}
\item
  Interpreting HaplotypeCaller Annotations: Below is an excerpt from a
  VCF file produced by HaplotypeCaller for a diploid sample:

\begin{verbatim}
#CHROM  POS     REF ALT QUAL    FILTER  INFO     FORMAT  SAMPLE1
chr1    12345   A   G   456.77  .        DP=45;QD=10.15;MQ=60.0;FS=0.0;SOR=0.5;ReadPosRankSu m=0.2;MQRankSum=0.0   GT:AD:GQ    0/1:22,23:99
chr1    67890   C   T   89.5    .       DP=8;QD=11.19;MQ=35.2;FS=12.3;SOR=1.8;ReadPosRankSum=-2.1;MQRankSum=-1.5 GT:AD:GQ    0/1:4,4:89
chr1    11111   G   A   1205.3  .       DP=60;QD=20.09;MQ=58.7;FS=85.4;SOR=5.2;ReadPosRankSum=-7.8;MQRankSum=-0.3 GT:AD:GQ    0/1:15,45:99
\end{verbatim}

  \begin{enumerate}
  \def\labelenumii{\alph{enumii}.}
  \item
    For each variant, determine whether it would pass standard GATK hard
    filters (QD ≥ 2.0, MQ ≥ 40.0, FS ≤ 60.0, SOR ≤ 3.0, ReadPosRankSum ≥
    -8.0, MQRankSum ≥ -12.5). Identify which filter(s), if any, each
    variant fails.
  \item
    For the variant at position 67890, the mapping quality (MQ) is 35.2,
    below the recommended threshold of 40. Explain what this low MQ
    suggests about the reads at this site and whether you would trust
    this variant call.
  \item
    For the variant at position 11111, the strand bias metric FS is 85.4
    (above the recommended 60.0 threshold). Describe what you would look
    for in IGV to decide whether this variant is a true polymorphism or
    a technical artifact.
  \end{enumerate}
\item
  Designing a Filtering Strategy: You have called variants on a
  non-model plant species with no high-quality variant database, so VQSR
  is not an option. Your raw VCF contains 1.2 million putative SNPs and
  150,000 putative indels.

  \begin{enumerate}
  \def\labelenumii{\alph{enumii}.}
  \item
    Outline a hard-filtering workflow: which tools (GATK SelectVariants,
    GATK VariantFiltration, VCFtools) would you use, in what order, and
    with what initial threshold values (you may use GATK-recommended
    defaults as a starting point).
  \item
    Describe how you would use VCFtools to generate summary statistics
    (depth, quality scores, and allele frequency distributions) before
    and after filtering, and what plots you would create to evaluate
    whether your thresholds are appropriate.
  \item
    Suppose that after applying your initial filters, you retain only 50
    percent of SNPs but 80 percent of indels. Interpret this result:
    does it suggest that your SNP filters are too strict, your indel
    filters too lenient, or that SNPs and indels in your dataset have
    genuinely different quality profiles? What additional information
    would help you decide?
  \end{enumerate}
\item
  Visual Inspection in IGV: You load your filtered VCF and
  duplicate-marked BAM into IGV and navigate to a heterozygous SNP that
  passed all filters. At this site, you observe the following:

  \begin{itemize}
  \tightlist
  \item
    Total depth: 40 reads
  \item
    Reference allele (A): 10 reads, all from the forward strand
  \item
    Alternate allele (G): 30 reads, 28 from the reverse strand and 2
    from the forward strand
  \item
    Mapping quality: most reads have MQ = 60, but the 2 forward-strand
    alternate reads have MQ = 10
  \end{itemize}

  \begin{enumerate}
  \def\labelenumii{\alph{enumii}.}
  \item
    Based on this visual evidence, do you believe this variant is a true
    polymorphism or an artifact? Justify your answer by discussing
    strand bias, allele balance, and mapping quality.
  \item
    Describe one additional piece of evidence you could gather (from IGV
    or another tool) to strengthen your conclusion.
  \item
    If you decide this variant is likely an artifact despite passing
    automated filters, explain what this reveals about the limitations
    of hard filtering and how such cases might be handled in a
    high-throughput pipeline.
  \end{enumerate}
\item
  Exploring Alternative Callers: Your lab is planning a variant-calling
  project on a cohort of 200 samples from a non-model fish species.
  Computational resources are limited, and you want to maximize
  sensitivity for both SNPs and indels.

  \begin{enumerate}
  \def\labelenumii{\alph{enumii}.}
  \item
    Compare GATK HaplotypeCaller, FreeBayes, and samtools/bcftools in
    terms of computational cost, sensitivity for indels, and ease of
    use. Based on this comparison, recommend one caller for your project
    and justify your choice.
  \item
    Propose an ensemble calling strategy that combines two of the three
    callers. Explain how you would reconcile their outputs (for
    instance, requiring a variant to be called by both callers, or using
    a weighted voting scheme) and what benefits you expect from this
    approach.
  \item
    Describe how you would validate your variant calls experimentally.
    If you could Sanger-sequence 50 sites, how would you choose which
    sites to validate (for example, concordant calls, discordant calls,
    high-impact variants)?
  \end{enumerate}
\item
  Joint Genotyping and Population Genetics: You have generated gVCF
  files for 50 individuals from a single population using
  HaplotypeCaller in reference-confidence mode. You now plan to run
  GenotypeGVCFs to produce a multisample VCF.

  \begin{enumerate}
  \def\labelenumii{\alph{enumii}.}
  \item
    Explain the advantages of joint genotyping compared to merging
    single-sample VCFs produced without reference-confidence mode. Give
    one example of a population genetic analysis (such as allele
    frequency estimation or linkage disequilibrium calculation) that
    benefits from joint genotyping.
  \item
    After joint genotyping, you observe that some sites that were called
    as heterozygous in individual-sample gVCFs are now called as
    homozygous reference in the multisample VCF. Explain how joint
    genotyping can revise individual genotype calls and why this is
    usually an improvement.
  \item
    Propose a filtering strategy for the multisample VCF that accounts
    for both site-level quality (QD, MQ, FS) and sample-level
    missingness (the fraction of individuals with no-call genotypes at
    each site). What threshold for missingness would you apply, and how
    would you justify it?
  \end{enumerate}
\end{enumerate}

\section{References}\label{references-9}

\chapter{Grant Funding and the Scientific
Enterprise}\label{grant-funding-and-the-scientific-enterprise}

\section{Introduction}\label{introduction-1}

This chapter reveals another layer of the hidden curriculum in
scientific research---one that is rarely discussed in coursework but
that shapes nearly every research career in biology: the pursuit of
grant funding. Securing funding is not merely a bureaucratic necessity;
it is the engine that powers scientific discovery, sustains research
programs, and opens doors for trainees at every career stage. Just as
Chapter 7 equipped you with the skills to manage your research projects
and communicate your methods, this chapter equips you with the
conceptual and practical tools to understand where scientific money
comes from, how to ask for it compellingly, how it is evaluated by your
peers, and where to begin your own funding journey as a trainee.

\section{The History of Federal Funding for Scientific Innovation in the
United
States}\label{the-history-of-federal-funding-for-scientific-innovation-in-the-united-states}

\subsection{The Origins of Organized Science
Support}\label{the-origins-of-organized-science-support}

The relationship between the United States federal government and the
scientific community is one of the most consequential partnerships in
modern history. For most of the nineteenth century, the federal
government played only a modest role in supporting basic research. The
Marine Hospital Service maintained a small laboratory as early as 1887
to investigate infectious diseases---a one-room operation that would
eventually grow into the National Institutes of Health
(NIH)\textsuperscript{168}. Yet it was the crucible of the Second World
War, more than any other event, that fundamentally transformed how the
nation understood and funded science\textsuperscript{169}. The
mobilization of universities, private laboratories, and government
agencies to solve wartime problems---from radar to penicillin to the
atomic bomb---demonstrated beyond any doubt that basic scientific
research could be harnessed in service of national goals.

As the war drew to a close, political leaders and scientists alike
recognized that the government-science partnership could not simply
dissolve with the armistice. The most influential voice in this
conversation was Vannevar Bush, who had directed the wartime Office of
Scientific Research and Development. At the request of President
Roosevelt, Bush produced in 1945 a landmark report titled
\emph{Science---The Endless Frontier}, which articulated the case that
sustained federal investment in basic research at universities would pay
dividends not only in national security but in economic prosperity and
public health\textsuperscript{169}. Bush argued that government should
fund the best science at the best institutions, free from political
interference, and that discoveries made at the frontier of knowledge
would inevitably find practical application---even when no one could
predict exactly how or when. This philosophy, sometimes called the
``linear model'' of innovation, became the intellectual foundation upon
which both the NIH and the National Science Foundation (NSF) were built.

\subsection{The Rise of NIH}\label{the-rise-of-nih}

The NIH's origins trace back to 1887 and the modest Hygienic Laboratory
of the Marine Hospital Service, which was created primarily to
investigate cholera and other infectious diseases arriving at American
ports\textsuperscript{168}. Over the following decades, Congress
gradually expanded the laboratory's mandate and resources, formally
renaming it the National Institute of Health in 1930 and later
broadening it into the National Institutes of Health as multiple
disease-specific institutes were added. The landmark National Cancer Act
of 1937 established a National Cancer Institute and introduced the model
of funding extramural research through grants to university
scientists---a mechanism that would come to define the NIH's identity.
Today, NIH comprises 27 institutes and centers, each with its own
scientific mission and grant portfolio, and its annual budget exceeds
\$47 billion, making it the world's largest public funder of biomedical
research\textsuperscript{168}.

The culture of NIH funding was shaped by a basic bargain: the government
would provide funds, and scientists at universities and research
institutions would conduct the work, retaining intellectual freedom to
follow the science wherever it led. This model distinguished NIH from
purely mission-driven agencies, though the tension between pure
curiosity and translational goals has never fully resolved. Throughout
the Cold War, congressional enthusiasm for biomedical research
intensified as disease conquest became a symbol of American
technological prowess. Successive administrations doubled the NIH budget
multiple times over, creating what many researchers nostalgically call
the ``golden era'' of American science funding that stretched through
the late 1990s and culminated in the NIH budget doubling from 1998 to
2003\textsuperscript{168}.

\subsection{The Founding and Evolution of the
NSF}\label{the-founding-and-evolution-of-the-nsf}

The National Science Foundation took a longer and more contentious path
to existence. The debates that delayed its founding from 1945 to 1950
were, at their core, debates about democratic accountability versus
scientific autonomy. Bush and his allies wanted an agency governed by
scientists, insulated from politics, and dedicated to funding only the
most excellent basic research without regard for geographic
distribution\textsuperscript{169}. Senator Harley Kilgore of West
Virginia pushed back, arguing that federal money should be distributed
more broadly---geographically and disciplinarily---and that the agency
should be accountable to elected officials, not to an unelected
scientific elite. When President Truman vetoed the first version of the
legislation in 1947 on the grounds that it was insufficiently
democratic, it took three more years of negotiation to find workable
compromises\textsuperscript{169}. NSF was ultimately signed into law by
Truman on May 10, 1950, with a National Science Board of 24 part-time
members providing oversight and a presidentially appointed director as
chief executive officer\textsuperscript{170}. Critically, graduate
research fellowships were written into NSF's mission from the very
start: the Graduate Research Fellowship Program launched in 1952 and
remains the agency's longest continuously operating program, reflecting
how central trainee support has been to the agency's identity since its
inception\textsuperscript{170,171}.

In its early years, NSF was a modestly funded agency with an initial
appropriation of just \$225,000, dwarfed by the research programs of NIH
and the Department of Defense. The shock of Sputnik in October 1957
changed everything: when the Soviet Union placed the first artificial
satellite in orbit, Congress more than doubled the NSF budget almost
overnight---from \$40 million to \$134 million in a single year---and
passed the National Defense Education Act of 1958, which poured federal
money into science and mathematics education at every
level\textsuperscript{170}. The decades that followed saw NSF's
portfolio expand well beyond the physical sciences. A 1968 amendment to
the authorization act made explicit the agency's mandate to support the
social sciences and applied research and to foster development of
computer technologies\textsuperscript{170}. By 1985, NSF had established
NSFNET---the high-speed network linking its national supercomputing
centers---which became the direct technical forerunner of the internet,
one of the most consequential infrastructure investments in the agency's
history\textsuperscript{170}. New directorates were established as
scientific fields matured: Engineering in 1981, Computer and Information
Science and Engineering in 1985--86, and Social, Behavioral and Economic
Sciences in 1991. In 2001, NSF coined the now-ubiquitous acronym STEM,
and in 2022 it established its first new directorate in thirty
years---the Directorate for Technology, Innovation and Partnerships
(TIP)---specifically designed to bridge curiosity-driven basic research
and use-inspired applied innovation, grow industry, and cultivate a more
equitable STEM workforce\textsuperscript{170}.

NSF-funded science has produced a remarkable record of foundational
discoveries, many of which are woven into the fabric of everyday life:
the bacterium isolated in Yellowstone hot springs that yielded the DNA
polymerase behind PCR; the NSFNET backbone that gave rise to the
internet; CRISPR technology characterized in 2012; the LIGO
gravitational-wave observatory that confirmed Einstein's general theory
of relativity in 2016; and the Event Horizon Telescope image of the
first black hole, unveiled in 2019\textsuperscript{170}. These outcomes
were not predictable at the time of the original investments---which is
precisely Bush's point in \emph{Science---The Endless Frontier} and the
empirical foundation for the return-on-investment argument for basic
science funding\textsuperscript{171,172}. Today NSF remains distinct
from NIH in that it funds a broader range of disciplines---including the
mathematical and physical sciences, engineering, computer science, the
social sciences, and biology---and places particular emphasis on basic
research with no predetermined practical goal, while increasingly also
investing in translation of discoveries toward societal benefit through
programs like TIP and I-Corps\textsuperscript{170}.

\subsection{Beyond NIH and NSF: The Broader Funding
Ecosystem}\label{beyond-nih-and-nsf-the-broader-funding-ecosystem}

While NIH and NSF dominate the federal funding landscape for biology and
biomedical science, they are far from the only game in town. A
researcher working in computational biology or bioinformatics may find
equally relevant funding from the Department of Energy (DOE), which has
long supported genomics research through its Joint Genome Institute and
the work surrounding the Human Genome Project. The Department of Defense
funds biological research through its Defense Advanced Research Projects
Agency (DARPA) and service-specific research programs, particularly in
areas touching on infectious disease, biotechnology, and biosecurity.
The United States Department of Agriculture (USDA), through its National
Institute of Food and Agriculture (NIFA), is the primary federal funder
of agricultural and food science research, and for faculty and students
at land-grant universities like Auburn it represents one of the most
relevant and accessible funding streams available.

Beyond federal agencies, private foundations have played an enormous
role in shaping American science. The Howard Hughes Medical Institute
(HHMI) supports some of the most eminent biomedical scientists in the
country through long-term investigator awards explicitly designed to
give researchers the freedom to take risks without the pressure of
annual grant renewals. The Simons Foundation invests heavily in
mathematics, computational science, and autism research. The Bill and
Melinda Gates Foundation has redirected global health research
priorities through strategic grantmaking at a scale that rivals
government funding for certain disease areas. Professional scientific
societies---from the American Society for Microbiology to the Genetics
Society of America---also provide smaller grants and travel awards that,
while modest in dollar terms, can be transformative for early-career
researchers and trainees who are just establishing their independent
scientific identities\textsuperscript{173}.

\begin{tcolorbox}[enhanced jigsaw, rightrule=.15mm, leftrule=.75mm, colbacktitle=quarto-callout-note-color!10!white, colframe=quarto-callout-note-color-frame, toprule=.15mm, opacitybacktitle=0.6, title=\textcolor{quarto-callout-note-color}{\faInfo}\hspace{0.5em}{Land-Grant Institutions and the USDA-NIFA Funding Legacy}, left=2mm, colback=white, bottomtitle=1mm, toptitle=1mm, bottomrule=.15mm, coltitle=black, titlerule=0mm, breakable, opacityback=0, arc=.35mm]

The story of federal science funding in the United States cannot be told
without acknowledging the remarkable experiment in public higher
education that preceded the NIH and NSF by nearly a century. The Morrill
Act of 1862, signed into law by President Abraham Lincoln on July 2,
1862---in the middle of the Civil War---granted each state 30,000 acres
of federal land per member of Congress to establish colleges focused on
agriculture, mechanic arts, and the practical sciences. The goal was
explicitly democratic: to open higher education to farmers and working
people who had historically been excluded from the narrow classical
curriculum of existing colleges. Auburn University was designated
Alabama's land-grant institution in 1872 and has operated under that
tripartite mission of teaching, research, and public service ever since.

The original 1862 Act, however, did not extend its democratic promise
equally. With slavery still in force when the law passed and segregation
entrenched in its aftermath, Black students were effectively excluded
from most of the institutions it created. The Second Morrill Act of
1890, signed by President Benjamin Harrison on August 30, 1890, was a
direct response to this injustice. It provided direct federal financial
support to land-grant colleges and explicitly prohibited racial
discrimination in the distribution of those funds---requiring states
that maintained segregated systems to either admit Black students to
existing institutions or establish separate land-grant colleges for
them. This provision gave rise to the 19 institutions now known as the
1890 Land-Grant University System, all of which are Historically Black
Colleges and Universities (HBCUs). Alabama is home to two of these 1890
institutions, giving it a distinctive place in land-grant history.
Alabama A\&M University (AAMU) was founded in 1875 by Dr.~William Hooper
Councill in Huntsville as the Huntsville Normal School, and by 1891 was
receiving federal Morrill Act funds to support instruction in
agriculture, engineering, and architecture. Tuskegee University, founded
by Booker T. Washington in 1881, holds a unique status as the only
private institution designated as a land-grant university under the 1890
framework. Together, Auburn, AAMU, and Tuskegee constitute Alabama's
three land-grant institutions---a trio that traces the full arc of the
land-grant story, from its founding democratic ideals to the legislative
corrections that Congress later found necessary to extend those ideals
more broadly.

The agricultural research mission of land-grant universities was
formalized and funded by the Hatch Act of 1887, which established a
national system of Agricultural Experiment Stations at each land-grant
institution and provided direct federal appropriations to support their
work. These Hatch Act funds---administered today through USDA-NIFA---are
formula-based grants distributed annually to experiment stations,
supporting research programs in crop science, animal science, natural
resources, and increasingly, genomics and computational biology as these
tools have become essential to modern agricultural research. The
Smith-Lever Act of 1914 extended the land-grant mission further by
creating the Cooperative Extension Service, which translates university
research into practical guidance for farmers, communities, and families
across each state. NIFA also administers dedicated programs for 1890
institutions to support capacity building in research and extension,
though chronic underfunding of these programs relative to their 1862
counterparts has remained a documented and ongoing equity concern.

Today, NIFA administers a portfolio of competitive and capacity-based
grants that extend well beyond agriculture in the traditional sense.
Competitive programs such as the Agriculture and Food Research
Initiative (AFRI) fund fundamental and applied research in plant and
animal genomics, food safety, climate adaptation, and bioinformatics.
For students and faculty at Auburn working on topics at the intersection
of biology, genetics, and agricultural systems, NIFA represents a
historically grounded and institutionally supported pathway to federal
research funding that is worth understanding alongside NIH and NSF.

\end{tcolorbox}

The diversity of funding sources reflects the diversity of scientific
questions and the different priorities that animate different sponsors.
Understanding this landscape---knowing which agencies fund which kinds
of science, and what each funder's current programmatic priorities
are---is itself a professional skill that experienced scientists
cultivate over the course of their careers. No two grant applications
are identical because no two funders are identical, and the successful
proposal writer is always writing for a specific audience with a
specific mission in mind\textsuperscript{174}.

Underlying all of these mechanisms is a simple but powerful empirical
case: federal investment in scientific research generates substantial
returns for the national economy and public welfare. Economists estimate
that every dollar invested in basic science by the federal government
produces roughly \$2.56 in economic growth over a ten-year horizon, a
return that may be severalfold larger over multi-decadal time
scales---and one that individual investors or companies would rarely
make on their own because the benefits accrue broadly and unpredictably
across society rather than returning to a single
funder\textsuperscript{171,172,175}. This is precisely the economic
logic that animated Vannevar Bush's argument in \emph{Science---The
Endless Frontier} and that continues to be supported by empirical
research: curiosity-driven public investment creates knowledge that no
market would reliably produce, because the value of a discovery often
cannot be anticipated or captured by whoever paid for
it\textsuperscript{176,177}. Understanding this return-on-investment
argument is not merely academic background; it is the core rationale
that scientists are increasingly called upon to articulate when making
the case for their work to funding agencies, university administrators,
legislators, and the public\textsuperscript{178}.

\section{The Anatomy of a Grant
Proposal}\label{the-anatomy-of-a-grant-proposal}

\begin{tcolorbox}[enhanced jigsaw, rightrule=.15mm, leftrule=.75mm, colbacktitle=quarto-callout-note-color!10!white, colframe=quarto-callout-note-color-frame, toprule=.15mm, opacitybacktitle=0.6, title=\textcolor{quarto-callout-note-color}{\faInfo}\hspace{0.5em}{Figure 11.1 -- Anatomy of a Grant Proposal}, left=2mm, colback=white, bottomtitle=1mm, toptitle=1mm, bottomrule=.15mm, coltitle=black, titlerule=0mm, breakable, opacityback=0, arc=.35mm]

\begin{figure}[H]

\centering{

\includegraphics[width=0.7\linewidth,height=\textheight,keepaspectratio]{images/Fig_11.1_Grant_Anatomy.png}

}

\caption{\label{fig-grantanatomy}Anatomy of a Grant Proposal}

\end{figure}%

A successful grant proposal moves through five tightly linked
components, each building on the last. It begins not with methods but
with an \textbf{idea}---a central scientific question or hypothesis that
frames everything that follows. The \textbf{background} section is not a
literature review but a curated argument that establishes significance
and identifies the gap your work will fill. \textbf{Specific aims}
translate that argument into discrete, achievable goals stated with
strong action verbs. The \textbf{research approach} provides the
step-by-step experimental plan for each aim, complete with a
mini-background, anticipated challenges, and alternative strategies.
Finally, the \textbf{conclusion} closes the argument by articulating the
impact of the potential results---explaining what the field, and the
world, will gain if the work succeeds.

\end{tcolorbox}

\subsection{Beginning with an Idea}\label{beginning-with-an-idea}

Every successful grant proposal begins long before a single word is
written. It begins with a scientific idea---a question or gap in
knowledge that you find genuinely compelling and that you believe is
both important and tractable\textsuperscript{179}. The process of
converting that idea into a fundable proposal is a kind of translation:
you must move from the personal, often inarticulate conviction that
something is worth investigating to a clear, persuasive public argument
that a community of scientific peers should agree it is worth funding.
This is not a dishonest process, but it is a rhetorical one. Grant
writing is a genre of scientific communication with its own conventions,
its own vocabulary, and its own logic, and mastery of that genre
requires practice and attention just as learning to write clear methods
sections (as discussed in Chapter 7) requires practice and
attention\textsuperscript{173}.

The most common mistake beginning grant writers make is starting with
the methods rather than the idea. A proposal built around
techniques---``I will use RNA-seq and CRISPR to study gene X''---reads
as a solution looking for a problem. Reviewers respond far more
enthusiastically to proposals built around a scientific question---``We
do not understand how organism Y responds to environmental stress Z, and
answering this question has significant implications for W''---in which
the methods are clearly in service of the question rather than the other
way around. Before you write anything, you should be able to articulate
your central idea in plain language to a scientifically literate
non-specialist, because that is roughly the perspective of the reviewers
who will read your proposal\textsuperscript{174}.

\subsection{Building the Necessary
Background}\label{building-the-necessary-background}

Once your central idea is clear, the next task is constructing the
background section of your proposal. This section exists for a specific
purpose: to give reviewers the intellectual context they need in order
to understand why your question is significant and why it has not
already been answered. A common misconception is that the background
section is a comprehensive literature review. It is not. It is a curated
argument. You are selecting the most important prior findings,
synthesizing them in a way that points toward the gap you intend to
fill, and establishing that you have command of the relevant
literature\textsuperscript{180}. Every piece of background information
included should earn its place by doing one of three things:
establishing the importance of the problem, demonstrating what is
already known, or revealing what remains unknown.

The final paragraph of a background section should pivot decisively from
the known to the unknown. The reader should feel the intellectual
tension---here is what we know, here is what we do not know, and here is
why that gap matters. This pivot sets up the specific aims that follow
and gives reviewers a reason to keep reading with enthusiasm rather than
obligation. A well-crafted background does not just describe the field;
it tells a story with an unresolved ending that your proposed research
will complete\textsuperscript{181}.

\subsection{Specific Aims and the Power of Action
Verbs}\label{specific-aims-and-the-power-of-action-verbs}

The specific aims section---sometimes called objectives, or research
goals---is arguably the most important page in any grant proposal. It
distills your entire scientific vision into a small number of discrete,
achievable goals that collectively address your central
question\textsuperscript{182}. Many experienced grant writers treat the
specific aims as the backbone of the entire proposal and draft it first,
before writing any other section, because it forces clarity of thought
about what the project will actually accomplish.

Each aim must be stated using a strong action verb that describes what
you will do and what outcome is expected: \emph{determine}, \emph{test},
\emph{establish}, \emph{quantify}, \emph{evaluate}. Passive or vague
formulations---``aim one will examine aspects of the relationship
between X and Y''---communicate uncertainty and imprecision, while
active formulations---``Aim 1 will determine how transcription factor X
modulates the expression of stress-response genes in response to heat
shock in \emph{Drosophila}''---communicate confidence and
specificity\textsuperscript{173}. A proposal that promises to
``examine'' a phenomenon can never truly fail, because any result counts
as examination. A proposal that promises to ``test the hypothesis that X
causes Y'' succeeds or fails based on whether that hypothesis is
rigorously tested---which is exactly the kind of intellectual
accountability that reviewers want to see. This is not to say that
exploration driven work is not prioritized. As stated early in Chapter
1, discovery-based research is still a valid research goal. But again,
the verbs you use to describe these goals matters, so an aim promising
to \emph{characterize} something can be perceived as descriptive, which
can come across as weak. Instead, for discovery based work, use verbs
such as \emph{identify}, \emph{develop}, etc to strengthen your specific
aims descriptions. The choice of verb matters because it implicitly
defines what success looks like. In a lot of ways, these verbs can be
similar to verbs used in writing learning objectives for students based
on what is referred to as Bloom's Taxonomy. Simply put,
\href{https://www.utica.edu/academic/Assessment/new/Blooms\%20Taxonomy\%20-\%20Best.pdf}{the
verbs} match the expected level of understanding and these same
principles can be useful in constructing strong specific aims
statements.

It is also important that the aims be logically related without being
fatally dependent on one another. If Aim 2 can only be pursued if Aim 1
succeeds, then a reviewer will worry that the entire project collapses
if the first aim encounters unexpected obstacles. Experienced proposal
writers often structure their aims so that each one can yield meaningful
scientific results independently, even while they collectively address a
larger question\textsuperscript{174}. This architecture demonstrates
scientific maturity and reassures reviewers that the proposed work is
both ambitious and realistic.

\subsection{Structuring the Research
Approach}\label{structuring-the-research-approach}

For each aim, the research approach section must do several things in
sequence. It typically begins with a brief aim-specific background that
reminds the reader why this particular aim is important and what prior
evidence motivates the experimental strategy. This mini-background is
not redundant with the overall background section; rather, it provides
the immediate justification for the specific experiments that follow by
citing the most directly relevant prior work and, often, the applicant's
own preliminary data\textsuperscript{182}.

After establishing context, the approach presents the detailed
experimental steps in logical order. Each step should specify not just
what you will do but how you will do it, what you will measure, and what
criteria you will use to interpret the results. Reviewers reading an
approach section should be able to mentally simulate the experiments:
they should understand the design, anticipate the likely range of
outcomes, and evaluate whether the proposed methods are appropriate for
the stated goals\textsuperscript{174}. A common weakness in proposals
from early-career investigators is that the methods are described in
conceptual terms without enough procedural detail to give reviewers
confidence that the investigator knows how to actually execute the work.

Every well-constructed aim should include an honest discussion of
anticipated challenges and potential solutions\textsuperscript{179}.
This is not an invitation to undermine your own proposal by cataloguing
everything that could go wrong; rather, it is a demonstration that you
have thought deeply about the problem and have contingency plans in
place. A proposal that anticipates no obstacles reads as naive. A
proposal that anticipates specific, scientifically grounded
challenges---and responds with specific alternative approaches---reads
as the work of an experienced scientist who has genuinely engaged with
the complexity of the question. Reviewers are themselves working
scientists who know that research rarely proceeds exactly as planned,
and they are reassured when they see evidence that the applicant knows
this too.

\subsection{Concluding with Impact}\label{concluding-with-impact}

A strong grant proposal does not end with the last experimental
protocol. It ends with a statement of what the results will mean---for
the field, for human health, for society more
broadly\textsuperscript{180}. This conclusion is not a formality; it is
the payoff of the entire argument you have been building. You began by
establishing that a question is important and unanswered. You proposed a
rigorous strategy for answering it. Now you must explain what the world
will look like when you have succeeded. Will your findings establish a
new paradigm? Will they open new therapeutic avenues? Will they provide
foundational tools or datasets that will enable the broader community to
do science that is currently impossible?

The best closing statements are specific rather than generic. ``These
studies will advance our understanding of biology'' is not a compelling
statement of impact. ``Identifying the molecular mechanism by which
organism X repairs DNA under conditions of oxidative stress will provide
the first functional model for this process in a non-model eukaryote and
will directly inform the development of cancer therapies targeting DNA
repair pathways in human tumor cells'' tells a reader exactly what will
be gained and why it matters\textsuperscript{174}. Thinking carefully
about your impact statement also has a useful feedback effect: if you
find it difficult to articulate the significance of your results, that
may be a sign that your specific aims need to be sharpened or that the
framing of your project needs revision.

\section{Modern Peer Review of
Grants}\label{modern-peer-review-of-grants}

\subsection{How Grants Are Submitted and
Assigned}\label{how-grants-are-submitted-and-assigned}

After a proposal is written, polished, and submitted through the
relevant grants management system---Grants.gov for NIH and Research.gov
for NSF, for example---it enters a formal review process that is far
more structured than many applicants realize. Both NSF and NIH rely on
panels of expert volunteer reviewers drawn from the scientific
community, but the two agencies organize that review somewhat
differently in ways that reflect their distinct cultures and missions.
Understanding both systems is valuable, because a researcher in biology
or bioinformatics will likely encounter both over the course of a
career.

At NIH, most investigator-initiated research grant applications
(including the workhorse R01 mechanism) are submitted to the Center for
Scientific Review (CSR), which serves as the primary intake point for
the vast majority of NIH applications\textsuperscript{168}. The CSR
assigns each application to an appropriate standing study section---a
standing panel of expert reviewers organized around a scientific
topic---and also assigns it to one or more NIH institutes or centers
that might fund it based on programmatic fit. At NSF, proposals are
submitted directly to specific programs within a directorate, and the
cognizant program officer plays a much more active role from the outset
in deciding how and by whom a proposal will be reviewed. NSF programs
use a combination of ad hoc mail review and face-to-face panel review
depending on the program, though panel-based review has become
increasingly common\textsuperscript{169}. For both agencies, assigned
reviewers read their applications in advance and prepare written
evaluations, and the panel meeting is where those individual assessments
are synthesized into a collective recommendation.

\subsection{NIH Review Criteria: Scoring for
Impact}\label{nih-review-criteria-scoring-for-impact}

NIH review is organized around five core criteria: Significance,
Investigator, Innovation, Approach, and
Environment\textsuperscript{183}. Each criterion receives its own score
on a one-to-nine scale, and reviewers also provide an Overall Impact
score---a holistic judgment of how likely the proposed work is to exert
a sustained, powerful influence on the research field. The overall
impact score is not simply an average of the five criterion scores; it
is a deliberate synthesis that asks reviewers to weigh the relative
importance of each criterion in the context of the specific application.
An early-career investigator with a slightly less impressive publication
record but a highly innovative idea in an understudied area might score
excellently on Significance and Innovation in a way that compensates for
a somewhat weaker Investigator score, for example.

At NIH, the Significance criterion asks whether the problem addressed is
important, whether current knowledge will be advanced by the proposed
work, and whether the scientific premise---the existing evidence base
that motivates the hypothesis---is rigorously
established\textsuperscript{183}. The Innovation criterion asks whether
the application challenges existing paradigms or develops new
methodologies, and it rewards proposals that are conceptually bold
rather than incremental. The Approach criterion---often the most heavily
weighted in reviewer discussions---evaluates whether the experimental
design is rigorously conceived, whether the applicant has adequately
accounted for potential problems, and whether the statistical framework
and use of appropriate controls are sound. Rigor and reproducibility
have become increasingly explicit components of NIH review, including
requirements that applicants address sex as a biological variable in
animal and human studies and that they demonstrate the rigor of prior
research underlying their hypotheses.

\subsection{NSF Review Criteria: Intellectual Merit and Broader
Impacts}\label{nsf-review-criteria-intellectual-merit-and-broader-impacts}

NSF review is organized around two overarching criteria that have shaped
how American scientists think and write about their work for decades:
Intellectual Merit and Broader Impacts\textsuperscript{169}. These two
criteria are deceptively simple in their formulation but require genuine
depth and specificity in practice. Intellectual Merit asks whether the
proposed activities have the potential to advance knowledge---within the
field and potentially across fields---and whether the proposed work is
creative, original, or potentially transformative. Broader Impacts asks
whether the proposed activities have the potential to benefit society
and contribute to desired societal outcomes, a deliberately broad
criterion that encompasses education, diversity, public engagement, and
the training of the next generation of scientists.

The NSF panel process differs from NIH in some important structural
ways. In a typical NSF review panel, a proposal is assigned a lead
reviewer and a scribe---roles roughly equivalent to NIH's primary and
secondary assigned reviewers---who present their assessments to the
panel and frame the discussion. The lead reviewer opens with a summary
of strengths and weaknesses and states their rating, the scribe adds
complementary or qualifying observations, and then the panel discusses
the proposal openly. NSF uses a rating scale of Excellent, Very Good,
Good, Fair, and Poor rather than NIH's numerical one-to-nine scale, and
the outcome of the panel review is typically a written panel summary
that is provided to applicants as feedback. Crucially, as with NIH, NSF
panelists do not make funding decisions. The program officer considers
the panel's ratings and summaries alongside programmatic priorities,
balance across subfields, and available budget before making a funding
recommendation to the division director.

\subsection{What Happens at a Study Section or Review
Panel}\label{what-happens-at-a-study-section-or-review-panel}

Whether the panel follows the NIH or NSF model, the meeting itself
unfolds through a characteristic sequence that is worth understanding in
detail. The designated official responsible for the meeting---a
Scientific Review Officer (SRO) at NIH, or a program officer at
NSF---opens the session, reviews procedures, manages conflicts of
interest, and ensures that all discussion is conducted in accordance
with agency guidelines\textsuperscript{183}. Conflicts of interest are
treated with great seriousness at both agencies: reviewers who have a
financial, personal, or professional relationship with an applicant that
could bias their evaluation are recused from that application's
discussion. In virtual meetings, conflicted reviewers are typically
moved to a waiting room for the duration of the relevant discussion and
returned only after scoring is complete.

\begin{tcolorbox}[enhanced jigsaw, rightrule=.15mm, leftrule=.75mm, colbacktitle=quarto-callout-note-color!10!white, colframe=quarto-callout-note-color-frame, toprule=.15mm, opacitybacktitle=0.6, title=\textcolor{quarto-callout-note-color}{\faInfo}\hspace{0.5em}{Figure 11.2 -- Structure of a Scientific Grant Review Panel}, left=2mm, colback=white, bottomtitle=1mm, toptitle=1mm, bottomrule=.15mm, coltitle=black, titlerule=0mm, breakable, opacityback=0, arc=.35mm]

\begin{figure}[H]

\centering{

\includegraphics[width=0.8\linewidth,height=\textheight,keepaspectratio]{images/Fig_11.2_Review_Panel.png}

}

\caption{\label{fig-grantpanel}Structure of a Scientific Grant Review
Panel}

\end{figure}%

Whether modeled on NIH or NSF, all peer review panels share the same
fundamental architecture. A Scientific Review Officer (SRO) or Program
Officer (PO) chairs the meeting and ensures that agency
regulations---including confidentiality and conflict-of-interest (COI)
rules---are followed throughout. Three or more reviewers are assigned to
each proposal in advance: the lead (or primary) reviewer opens
discussion with a summary of strengths, weaknesses, and a preliminary
score or rating; the scribe (or secondary) reviewer adds complementary
observations; and together they frame the conversation for the full
panel. Any panelist who is not in conflict with the application may then
participate in open discussion, and scores or ratings frequently shift
based on what emerges. Assigned reviewers state their final scores to
define the range within which the rest of the NIH panel votes. At NIH,
the SRO records the discussion whereas at NSF, the reviewers then draft
a Panel Summary for the written summary that applicants ultimately
receive. At NIH, overall impact is reported on a 1--9 numerical scale;
at NSF, ratings follow an Excellent, Very Good, Good, Fair, or Poor
scale---but the deliberative process is remarkably similar across both
agencies.

\end{tcolorbox}

Not every application receives a full discussion. At NIH, applications
that receive preliminary scores in the less competitive range may be
triaged before the meeting and receive only written critiques without
oral discussion\textsuperscript{183}. This triaging process exists
because study sections often receive more applications than can be
meaningfully discussed in available time; focusing the discussion on the
most competitive applications uses reviewers' time most efficiently and
provides the most useful collective feedback for those applications. At
NSF, all assigned proposals typically receive written reviews and a
panel summary, though the depth of discussion varies with the
competitiveness of the application. Receiving a ``not discussed''
outcome on an NIH submission can be discouraging, but the written
critiques from assigned reviewers provide actionable feedback for
revision.

For applications that do receive full discussion, the pattern at both
agencies is similar: assigned reviewers present their assessments,
discussion opens to the full panel, and scores or ratings may shift
substantially based on what emerges from the
conversation\textsuperscript{183}. Observers of these discussions are
consistently surprised by two things. First, scores can move
dramatically---applications that opened with uniformly enthusiastic
preliminary scores can end with mediocre final scores once unassigned
panelists raise weaknesses that the assigned reviewers had not fully
weighted, and vice versa. Second, the discussion is genuinely
deliberative: reviewers argue, push back, and recalibrate their
assessments based on the expertise and perspectives of their colleagues.
This deliberative quality is not a flaw in the system; it is the system
working as designed. The goal is not consensus but clarity about the
strengths and weaknesses of each proposal, so that reviewers can submit
final scores or ratings that accurately reflect the panel's collective
scientific judgment. At the end of a well-run panel, applicants receive
both a quantitative assessment and a written narrative that together
constitute some of the most targeted scientific feedback they will ever
receive about their ideas.

\subsection{The Culture of Peer
Review}\label{the-culture-of-peer-review}

Sitting on a study section or NSF review panel---whether as an assigned
reviewer or a general panelist---is one of the most educational
experiences available to a scientist. Experienced investigators
consistently report that reviewing grants improved their own grant
writing because the experience makes the reviewer's perspective
viscerally real in a way that no amount of reading about review criteria
can replicate. Watching how assigned reviewers present their critiques,
hearing how unassigned panelists respond to a scientific argument, and
observing how panel discussion shapes and sometimes overturns
preliminary scores gives insights into the social dynamics of peer
review that are genuinely difficult to acquire any other
way\textsuperscript{183}.

One important thing to understand is that peer review is not designed to
produce consensus. When reviewers disagree fundamentally about whether a
proposal's central hypothesis is well-supported by prior evidence, the
chair's job is not to mediate until everyone agrees but rather to ensure
that all major perspectives are clearly articulated so that the range of
scores and the written critiques accurately reflect the genuine
scientific disagreement. Funding rates at both NSF and NIH are
sobering---competitive NSF programs may fund only five to ten percent of
submitted proposals, and NIH funding rates for R01 grants at many
institutes hover below twenty percent---which means that the bar for a
``fundable'' proposal is genuinely high, and the difference between
funded and unfunded applications is often a matter of precision and
depth rather than the overall quality of the science. Bias in peer
review is a recognized problem that both agencies actively work to
address through reviewer training, diverse panel composition, and
policies designed to keep reviewers focused on the scientific content of
applications.

\begin{tcolorbox}[enhanced jigsaw, rightrule=.15mm, leftrule=.75mm, colbacktitle=quarto-callout-note-color!10!white, colframe=quarto-callout-note-color-frame, toprule=.15mm, opacitybacktitle=0.6, title=\textcolor{quarto-callout-note-color}{\faInfo}\hspace{0.5em}{Course Connections: The Mock Grant Review Panel}, left=2mm, colback=white, bottomtitle=1mm, toptitle=1mm, bottomrule=.15mm, coltitle=black, titlerule=0mm, breakable, opacityback=0, arc=.35mm]

In the associated assignment for this chapter, you will participate in a
mock grant review panel that simulates the NSF panel process described
above. The proposals you will be reviewing are two-page NSF Graduate
Research Fellowship Program (GRFP) applications written by real graduate
students---the same fellowship discussed in Section 4 of this chapter.
The panel will follow the NSF framework, with a primary reviewer and a
second `scribe' for each proposal, open panel discussion, and a final
class-wide funding vote. While this chapter presents both the NIH and
NSF review frameworks because both are important for your long-term
development as a scientist, the class panel favors the NSF format, using
Intellectual Merit as the primary criterion and an Excellent-to-Poor
rating scale rather than NIH's numerical system. Reading the guidance in
this chapter on what reviewers look for---clear questions, compelling
background, rigorous approach, and honest acknowledgment of
challenges---is the best preparation you can do for that assignment.

\end{tcolorbox}

\section{Funding Opportunities for Trainees in Biology, Bioinformatics,
and the Biomedical
Sciences}\label{funding-opportunities-for-trainees-in-biology-bioinformatics-and-the-biomedical-sciences}

\subsection{Why Trainees Should Think about Funding
Early}\label{why-trainees-should-think-about-funding-early}

Many students encounter the idea of applying for their own grant funding
only late in their graduate careers, if at all. This is a missed
opportunity. Competitive fellowship awards are among the most powerful
credentials a trainee can accumulate, and not only because they come
with stipends and research allowances\textsuperscript{173}. A fellowship
communicates to future employers, collaborators, and mentors that a
trainee has already been evaluated by a community of scientific peers
and found meritorious---that their ideas are not just personally
interesting but publicly compelling. The process of applying, even
unsuccessfully, also provides invaluable training in articulating a
scientific vision, writing for a specific audience, and synthesizing a
body of literature into a coherent argument. These are the same skills
that will be required throughout a research career, and there is no
better time to begin developing them than during training when the
stakes are relatively low and mentorship is readily available.

There is also growing evidence that early career funding has lasting
consequences for scientific trajectories. Research on what is sometimes
called the \emph{Matthew effect} in science funding suggests that
early-career setbacks in funding can have disproportionate and
compounding negative effects on long-term research productivity, whereas
early success creates momentum that tends to
persist\textsuperscript{184,185}. When a program officer, foundation, or
fellowship panel invests in a junior scientist, they have a demonstrated
interest in seeing that investment through---meaning that initial
funding often opens doors to mentorship networks, collaborative
opportunities, and subsequent funding that would not otherwise have been
accessible\textsuperscript{186}. NSF survey data confirm that graduate
students and postdoctoral researchers who hold fellowships report higher
rates of research productivity, broader professional networks, and
stronger career outcomes than their unfunded peers\textsuperscript{187}.
The practical implication is straightforward: the time to begin building
your funding portfolio is not after you have something impressive to
show, but during the process of developing it.

\subsection{Opportunities for Undergraduate
Researchers}\label{opportunities-for-undergraduate-researchers}

Undergraduate students at research-active universities often have access
to funding mechanisms specifically designed for them, though these
opportunities require some navigation. At Auburn University, the
Department of Biology administers an Excellence Fund that provides
financial support to undergraduate students engaged in research with a
faculty mentor; students must be nominated and sponsored by their
research supervisor, making the mentoring relationship itself the entry
point into this funding stream. More broadly available is the
\href{https://our.auburn.edu/}{Office of Undergraduate Research's}
Undergraduate Research Fellowship
(\href{https://our.auburn.edu/undergraduate-research-fellowship/}{URF})
Program, which supports students who wish to pursue a self-designed
research project under faculty mentorship. Applications typically
require a brief project description, a faculty endorsement, and a
budget, giving undergraduates their first experience with the basic
architecture of a grant proposal in a low-stakes, supportive context.

Beyond the institutional level, the NSF Graduate Research Fellowship
Program (GRFP)---which will be discussed in more detail below---is open
to undergraduate seniors who are about to begin graduate school, making
it one of the few prestigious national competitions that can be entered
before a student has formally begun their graduate
training\textsuperscript{169}. Undergraduates who are thinking about
research careers should investigate the GRFP early enough to allow
adequate preparation time, because the strongest applications typically
require many drafts, significant input from mentors, and the kind of
intellectual clarity that develops over semesters of hands-on research
experience.

\subsection{NSF Graduate Research Fellowship
Program}\label{nsf-graduate-research-fellowship-program}

The
\href{https://www.nsf.gov/funding/opportunities/grfp-nsf-graduate-research-fellowship-program}{NSF
Graduate Research Fellowship Program (GRFP)} is the most prestigious
national fellowship competition for graduate students in science,
technology, engineering, and mathematics, and it is the one that
biology, bioinformatics, and computational biology students are most
likely to encounter early in their graduate
careers\textsuperscript{169}. The fellowship provides three years of
support, including an annual stipend and an education allowance intended
to offset tuition and fees. It is open to U.S. citizens, nationals, and
permanent residents who are in their first or second year of graduate
school (or who are applying as undergraduate seniors or
post-baccalaureate students), and importantly, each applicant may only
submit one application in their lifetime---making timing and preparation
especially critical.

GRFP applications are evaluated using the same two criteria that govern
all NSF review: Intellectual Merit (the potential of the proposed
activities to advance knowledge) and Broader Impacts (the potential to
benefit society and contribute to desired societal
outcomes)\textsuperscript{169}. A strong GRFP application requires three
components to work together: a personal statement that narrates the
applicant's scientific development and identity, a previous research
experience essay that demonstrates concrete accomplishments and
scientific thinking, and a research proposal that shows the applicant
can formulate an independent scientific question and design a credible
approach to answering it\textsuperscript{173}. The research proposal is
often the most challenging component for students to write well, because
it requires exercising exactly the skills discussed in Section 2 of this
chapter---identifying a genuine gap, building a compelling background,
and articulating specific, achievable aims with the kinds of strong
action verbs and realistic contingency planning that distinguish mature
scientific thinking from enthusiasm alone.

\subsection{NIH Individual Fellowship
Awards}\label{nih-individual-fellowship-awards}

For graduate students and postdoctoral researchers in the biomedical
sciences, the NIH offers a family of individual fellowship mechanisms
known collectively as the F series. The F31 award, the Ruth L.
Kirschstein National Research Service Award for Individual Predoctoral
Fellows, supports doctoral candidates pursuing dissertations in
health-related research\textsuperscript{173}. There is also an F31
Diversity supplement specifically designed to enhance representation of
individuals from groups underrepresented in biomedical research,
including racial and ethnic minorities and individuals with
disabilities. The F32 award supports postdoctoral researchers, providing
a structured opportunity to develop independent research skills under
the mentorship of a sponsor during the critical transition period
between doctoral training and independent faculty positions.

Unlike the GRFP, which is evaluated primarily on the strength of the
applicant's ideas and potential, NIH fellowship applications involve a
more elaborate review process that takes into account the proposed
research project, the applicant's background and qualifications, the
mentor's qualifications and track record, the training plan, and the
environment in which the training will occur\textsuperscript{183}. This
means that the quality of the mentorship relationship---and the mentor's
willingness to invest time in co-developing the fellowship
application---matters enormously. Students planning to apply for an F31
should initiate conversations with their mentor about the application
well in advance of any deadline, ideally at least a year before the
intended submission, because developing a strong individual development
plan and a coherent research proposal requires sustained collaboration
between the trainee and their sponsor.

\subsection{NNF Graduate and Postgraduate
Fellowship}\label{nnf-graduate-and-postgraduate-fellowship}

For students whose research sits at the intersection of biology and
agricultural or food systems, the USDA-NIFA National Needs Graduate and
Postgraduate Fellowship (NNF) program represents a directly relevant and
often overlooked federal opportunity. Unlike the GRFP or NIH F-series
awards, NNF fellowships are administered through individual colleges and
universities rather than applied for directly by students---meaning that
the entry point is a conversation with your department or graduate
program about whether your institution holds an active NNF grant in your
area. Awards provide stipends and travel support for master's and
doctoral students in national-need subject areas identified by NIFA,
which have historically included genomics, bioinformatics applied to
agriculture, food safety, environmental science, and natural resources.
For students at a land-grant institution like Auburn, where the USDA
relationship is institutionally deep and research programs in
agricultural genomics and computational biology are well established, it
is worth asking whether an NNF training grant is available in your
college before looking exclusively to NIH or NSF for fellowship support.

\subsection{Professional Society Grants and Other
Opportunities}\label{professional-society-grants-and-other-opportunities}

Beyond the major federal fellowships, a rich ecosystem of professional
society grants and foundation-sponsored awards provides additional
funding opportunities for graduate students and early-career
researchers. Sigma Xi, the scientific research honor society, offers
Grants in Aid of Research (GIAR) that provide modest awards---typically
up to one thousand dollars---for travel to field sites or the purchase
of equipment necessary to complete a specific research
project\textsuperscript{173}. These awards are small enough that they
fall within reach of students early in their training, but competitive
enough that earning one represents a genuine credential. The Society for
Study of Evolution, the Genetics Society of America, the Society for
Integrative and Comparative Biology, and many other discipline-specific
societies offer analogous small grants, travel awards, and presentation
fellowships.

The Phycological Society of America, for example, offers grants in aid
of research for graduate students and postdoctoral researchers
conducting work in phycology, while the American Association of
University Women provides dissertation fellowships and international
fellowships supporting women in science at the doctoral level. National
databases maintained by organizations such as Pathways to Science and
UCLA's graduate funding portal maintain searchable listings of
opportunities for MS and PhD students across disciplines, including
travel grants, writing grants, and research awards that many students
never discover simply because they do not know to look for them. The
lesson here is that the funding landscape is both broader and more
diverse than most students realize, and even awards that seem peripheral
to your primary research area can contribute to your professional
development and your CV in meaningful ways.

Making a habit of monitoring funding opportunities throughout your
graduate training---building a personal calendar of relevant deadlines,
discussing opportunities with your mentor and with advanced graduate
students in your program, and setting aside time each semester to
evaluate which competitions are most aligned with your current stage of
training and scientific focus---is itself a form of professional
practice that will pay dividends throughout your career. The scientists
who build the most robust and sustained research programs are almost
universally those who learned early to think strategically about
funding, to write clearly and compellingly about their ideas, and to
treat the peer review process not as an adversarial obstacle but as a
community conversation about what science is most worth doing.

\section{Conclusion}\label{conclusion}

Grant funding is, at its core, a conversation between scientists about
what questions matter and who is best positioned to answer them.
Understanding the history of that conversation---from the Morrill Act's
vision of practical education for working people to Vannevar Bush's
vision of science as an endless frontier to the global network of
federal agencies, private foundations, and professional societies that
funds research today---helps you see yourself not as a supplicant asking
for money but as a participant in a long and consequential tradition of
publicly supported inquiry. The skills developed by writing grant
proposals, by serving as a reviewer, and by navigating the peer review
system are among the most durable and transferable skills in science,
and there is no better time to begin developing them than during your
training. The hidden curriculum revealed in this chapter is simply this:
the scientists who thrive are those who learn to communicate the
significance of their work not just to themselves but to the entire
community that depends on good science being done well.

\section{Practice Problem Sets}\label{practice-problem-sets-2}

\begin{enumerate}
\def\labelenumi{\arabic{enumi}.}
\item
  \textbf{Proposal Identification Exercise}: Return to
  the\textsuperscript{3} paper we have read throughout the semester and,
  using the \href{https://www.nsf.gov/awardsearch/simple-search/}{NSF
  Award database}, identify one of the funded grants that supported the
  work described. Examine the project abstract available through NSF.
  What elements of the proposal structure described in
  Figure~\ref{fig-grantanatomy} can you identify in the abstract alone?
  What is missing from the abstract that would be present in the full
  proposal? What other publications have been produced as a result of
  this research?
\item
  \textbf{Specific Aims Practice}: Return to your narrative submitted
  for \href{Lab02.qmd}{Lab 2}. Visit the link of action verbs and try to
  write specific aims for your hypothetical research project. Expand
  your narrative to a specific aims page, which is a 1-page mini
  proposal. Your aims page should include a brief statement of the
  central problem and gap in knowledge, two to three specific aims each
  written with a strong action verb, and a closing statement of expected
  impact. Exchange your draft with a classmate and provide written
  feedback specifically addressing whether the aims are logically
  independent, whether the action verbs are sufficiently specific, and
  whether the impact statement is convincing.
\item
  \textbf{Reviewer Perspective Exercise}: Obtain two or three funded
  grant abstracts from \href{https://report.nih.gov/}{NIH Reporter} in
  your area of interest and evaluate each one against NIH's five core
  review criteria (Significance, Investigator, Innovation, Approach,
  Environment). Write a brief mock critique for each, identifying one
  major strength and one potential weakness for each criterion. Then
  compare your critiques with a classmate: where do you agree, and where
  do you disagree? What does the disagreement tell you about the
  subjective elements of peer review?
\item
  \textbf{Study Section Role Play}: In preparation for the class mock
  study section, read the assigned proposal carefully and prepare both a
  written critique and a two-minute oral summary of the major
  score-driving issues. Your oral summary should state your preliminary
  score or rating, describe the most important strengths, identify the
  most important weaknesses, and explain how you weighted those factors
  in arriving at your assessment. After the panel discussion, reflect in
  writing on whether your rating changed, why or why not, and what
  aspect of the discussion most influenced your thinking.
\item
  \textbf{Funding Landscape Map}: Using the resources described in
  Section 11.5 plus the \href{pathwaystoscience.org}{Pathways to Science
  database}, compile a personal funding calendar listing at least five
  fellowship or grant opportunities relevant to your stage of training
  and research area. For each opportunity, record the award name,
  sponsoring organization, eligibility requirements, award amount, and
  typical deadline. Present your calendar to a small group and discuss
  which opportunity you would prioritize and why.
\end{enumerate}

\section{Reflection Questions}\label{reflection-questions-9}

\begin{enumerate}
\def\labelenumi{\arabic{enumi}.}
\item
  Vannevar Bush argued in 1945 that government should fund the ``best
  science'' without regard for geographic distribution, while Senator
  Kilgore argued for broad-based, democratically accountable support.
  Both positions have continued to shape funding policy for more than 75
  years. Where do you see the tension between these two visions playing
  out in the current funding landscape, and which perspective do you
  find more compelling? How might your answer change depending on
  whether you are thinking as a scientist, a taxpayer, or a citizen of a
  small state?
\item
  The Morrill Act of 1862 established land-grant universities with an
  explicitly democratic mission: to open higher education to farmers and
  working people by focusing on practical agriculture, science, and
  engineering. How does this original mission connect to the research
  you are doing or hope to do? In what ways does the modern research
  university, funded partly through NIFA and Hatch Act capacity grants,
  honor or depart from the vision that Justin Morrill and Abraham
  Lincoln made law?
\item
  Section 11.3 describes the specific aims page as arguably the most
  important page in a grant proposal. Think about a piece of research
  you are currently involved in or are familiar with from the
  literature. How would you translate that work into three specific
  aims, each framed around a strong action verb? What does the exercise
  reveal about the difference between describing what scientists did and
  proposing what scientists will do?
\item
  Imagine a study section or review panel where an application opened
  with uniformly enthusiastic scores from assigned reviewers but ended
  with mediocre scores after unassigned reviewers pointed out multiple
  independent weaknesses. What would this sequence of events reveal
  about the potential gap between an author's perception of their own
  work and a panel's collective evaluation of it? How might you use this
  insight to improve your own grant proposals or scientific writing more
  broadly?
\item
  This chapter presents both the NIH and NSF review frameworks because
  scientists in biology and bioinformatics encounter both during their
  careers. After participating in the class mock panel using NSF
  criteria, reflect on the differences you notice between evaluating
  proposals through the lens of Intellectual Merit alone versus the
  five-criteria NIH framework. Which system felt more natural to you as
  a reviewer, and which would you find more useful as feedback on your
  own proposals? What does your preference reveal about your scientific
  values or communication strengths?
\end{enumerate}

\section{References}\label{references-10}

\chapter{Population Genomics}\label{population-genomics}

\section{Introduction}\label{introduction-2}

As first introduced in Chapter 8, \textbf{population genomics} typically
describes a DNA resequencing project of multiple individuals from the
same population of a species. As a field, population genomics brings
together evolutionary theory and genome‑scale data to ask how mutation,
recombination, genetic drift, migration, and natural selection shape
patterns of genetic variation across the
genome\textsuperscript{188,189}. Rather than focusing on single loci,
population genomics characterizes the ``genomic landscape'' of diversity
and differentiation, providing a richer picture of how genomes evolve in
natural populations\textsuperscript{188}. This chapter introduces the
historical roots of population genomics, the key summary statistics and
evolutionary processes that govern genomic variation, and modern
genome‑wide analysis strategies such as sliding‑window and spline‑based
approaches.

The goal is to provide a broad conceptual overview rather than a
mathematically exhaustive treatment. Readers interested in formal
derivations and deeper theory are encouraged to consult dedicated
population genetics textbooks, such as Graham Coop's \emph{Population
and Quantitative Genetics} notes, which develop the underlying models
and statistical machinery in detail\textsuperscript{189}. Here, I
emphasize intuitive explanations, connections to earlier chapters (e.g.,
sequence analysis, variant calling, and alignment), and links to labs in
which you will interpret population genomic statistics such as
nucleotide diversity \(\pi\), Tajima's D, and \(F_{ST}\).

\section{From the Modern Synthesis to Population
Genomics}\label{from-the-modern-synthesis-to-population-genomics}

\subsection{Biological evolution and the forces of
change}\label{biological-evolution-and-the-forces-of-change}

Evolution is the change in allele frequency over time, typically
measured in the context of a genetic population. A population is made up
of a set of interbreeding individuals whose collective genomes
constitute its genetic composition\textsuperscript{189}. This
composition shifts because individuals are born and die, because some
individuals migrate in or out, and because individuals vary in how many
offspring they leave. Every new individual born into the population
subtly changes things: their genome is a unique combination of their
parents' genomes, shuffled by segregation and recombination during
meiosis, and possibly altered by new mutations. These individual events
seem small, but the accumulation of changes across individuals and
generations --- compounded over tens, hundreds, and millions of
generations --- produces the remarkable diversity of life on
Earth\textsuperscript{189}.

\textbf{Population genetics} is the study of the genetic composition of
natural populations and its evolutionary causes and
consequences\textsuperscript{189}. Population genomics extends this to
genome-wide data, asking how mutation, recombination, genetic drift,
migration, and natural selection together shape patterns of variation
across entire genomes\textsuperscript{188}. To understand how these
forces interact, population geneticists develop mathematical models. As
the statistician George Box famously put it, ``All models are wrong but
some are useful'' --- and the models of population genetics, despite
their simplifying assumptions, yield surprisingly accurate predictions
about patterns of genetic diversity in nature\textsuperscript{189}.

The five major forces that shape allele frequencies
are\textsuperscript{189,190}:

\begin{itemize}
\tightlist
\item
  \textbf{Mutation}: Introduces new alleles at a low rate
  (\textasciitilde{}\(10^{-8}\) per base pair per generation in humans);
  the ultimate source of all genetic variation.
\item
  \textbf{Genetic drift}: Random fluctuations in allele frequencies
  caused by the finite number of individuals that survive and reproduce
  each generation; strongest in small populations.
\item
  \textbf{Migration (gene flow)}: The movement of individuals and their
  alleles between populations; tends to homogenize allele frequencies
  across geographic space.
\item
  \textbf{Recombination}: Shuffles alleles within chromosomes during
  meiosis, breaking down associations between nearby variants.
\item
  \textbf{Natural selection}: Non-random, directional change in allele
  frequencies based on differences in survival and reproduction among
  genotypes.
\end{itemize}

Understanding population genomics means understanding how these forces
operate individually, how they interact, and how to distinguish their
signatures in genomic data.

\subsection{The modern synthesis: reconciling Mendel and
Darwin}\label{the-modern-synthesis-reconciling-mendel-and-darwin}

The conceptual foundations of population genomics lie in the early
twentieth-century \textbf{modern synthesis}, which reconciled Mendelian
genetics with Darwin's theory of natural selection\textsuperscript{189}.
Darwin's theory required heritable variation and differential
reproductive success, but he had no mechanistic theory of inheritance.
Mendel provided particulate, discrete heritable factors (what we now
call alleles), but his `factors' initially seemed incompatible with the
continuous phenotypic variation --- height, weight, skin tone ---
observed in nature. The breakthrough came from Fisher, Wright, and
Haldane, who showed mathematically that Mendelian segregation at many
loci, each with small effects, produces smooth continuous phenotypic
distributions, and that changes in allele frequencies under selection
follow well-defined rules\textsuperscript{189}.

It is worth appreciating how conceptually remarkable this is. Many of
the core models of population genetics were developed in the
1920s--1940s, before anyone knew that DNA was the molecule of
heredity\textsuperscript{189}. Yet these models have only grown in power
and relevance as molecular data became available. This is because they
are built on the logic of Mendelian transmission rules combined with
probability theory --- a foundation that holds regardless of the
physical nature of heredity.

Fisher formalized the relationship between genetic variation and
evolution in his \textbf{Fundamental Theorem of Natural Selection},
which links the rate of increase in mean fitness to the amount of
additive genetic variance for fitness\textsuperscript{190}. This
synthesis laid the groundwork for connecting patterns of genetic
variation in populations to the action of selection, drift, mutation,
and migration. Building on these ideas came the modern evolutionary
synthesis more broadly, integrating insights from paleontology,
systematics, and ecology into a unified picture of how both short-term
evolutionary change and the long-term diversification of life can be
explained by the gradual accumulation of genetic changes within and
among populations\textsuperscript{189}.

\begin{tcolorbox}[enhanced jigsaw, rightrule=.15mm, leftrule=.75mm, colbacktitle=quarto-callout-note-color!10!white, colframe=quarto-callout-note-color-frame, toprule=.15mm, opacitybacktitle=0.6, title=\textcolor{quarto-callout-note-color}{\faInfo}\hspace{0.5em}{Eugenics, Population Genetics, and Scientific Racism}, left=2mm, colback=white, bottomtitle=1mm, toptitle=1mm, bottomrule=.15mm, coltitle=black, titlerule=0mm, breakable, opacityback=0, arc=.35mm]

\emph{Note: The goal of this section is to situate the field of
population genetics in its historical context in a factual, objective
way, without assigning blame or asking students to endorse any
particular viewpoint.}

The history of genetics and evolutionary biology is intertwined with the
history of eugenics and scientific racism. Francis Galton coined the
term \emph{eugenics} in 1883 to describe ``human improvement'' through
controlled breeding\textsuperscript{189,191}. Eugenic thinking rested on
the idea that certain groups were genetically ``inferior'' and that
social inequalities could be justified as biological
facts\textsuperscript{191,192}. In the late nineteenth and early
twentieth centuries, eugenicists used emerging tools in statistics and
heredity to support policies restricting immigration and discouraging or
preventing reproduction by people labeled
``unfit''\textsuperscript{189,191}. In the United States, compulsory
sterilization laws disproportionately affected Black, Latino,
Indigenous, poor, and disabled individuals\textsuperscript{189,193};
coerced sterilization of Native American women continued into the
1970s\textsuperscript{194}, California prison inmates were sterilized
without lawful consent as recently as 2010\textsuperscript{195}, and a
Tennessee judge was formally reprimanded in 2017 for offering reduced
sentences in exchange for sterilization\textsuperscript{196}.

Many early geneticists engaged actively with eugenics, and major
scientific institutions have roots in the movement. The Cold Spring
Harbor Laboratory hosted a Eugenics Record Office from 1910 to
1940\textsuperscript{197}, and the journal now known as \emph{Annals of
Human Genetics} was called \emph{Annals of Eugenics} until
1954\textsuperscript{189}. R.A. Fisher devoted substantial sections of
his landmark 1930 book to eugenic arguments\textsuperscript{189,198},
and H.J. Muller framed genetic diversity in some populations as a burden
to be managed through his work on mutational load\textsuperscript{199}.
The Mendelian framework was itself selectively distorted to support
eugenic claims before geneticists pushed back on its
misuse\textsuperscript{200}.

At the same time, geneticists including Thomas Hunt Morgan and Lancelot
Hogben criticized eugenic claims, emphasizing the importance of
environment and the complexity of inheritance\textsuperscript{189}.
After World War II, explicit state-sponsored eugenics programs lost
formal legitimacy, though coercive practices persisted for decades.
Lewontin's landmark 1972 analysis showed that roughly 85\% of human
genetic variation occurs \emph{within} populations rather than between
them\textsuperscript{189,201,202}, reaffirmed many times since with
larger genomic datasets. This work helped establish that human genetic
differences, while real and scientifically interesting, do not map
neatly onto rigid racial categories or justify treating any group as
inherently superior or inferior.

Communicating this accurately remains an ongoing responsibility for
geneticists, educators, and clinicians --- including genetic counselors
--- who encounter genomic information in professional practice. Some
groups continue to selectively misuse population-genetic results, making
clear and contextualized science communication
essential\textsuperscript{189,203}. The field of genetic counseling
itself emerged, in part, as a deliberate effort to move clinical
genetics away from eugenic paternalism and toward patient-centered,
non-directive care --- a foundational value documented in the
profession's own history\textsuperscript{192,204}.

\end{tcolorbox}

\subsection{The language of population genetics: loci, alleles, and
frequencies}\label{the-language-of-population-genetics-loci-alleles-and-frequencies}

Before going further, it helps to establish the basic vocabulary of
population genetics precisely, because these terms are used more
carefully here than in a typical introductory genetics course.

A \textbf{locus} (plural: \emph{loci}) is a specific position in the
genome\textsuperscript{189}. It could be a single nucleotide base pair
or an entire gene region. The different genetic variants that can occupy
a locus are called \textbf{alleles}. If every individual in the
population carries the same allele at a locus, that locus is
\textbf{monomorphic} --- there is no genetic variation at that position.
If two or more alleles exist in the population, the locus is
\textbf{polymorphic} (also called a \textbf{segregating
site})\textsuperscript{189}.

We call the less common allele at a polymorphic site the \textbf{minor
allele} and the more common one the \textbf{major allele}. We also
distinguish the \textbf{derived allele} (the evolutionarily newer one,
which arose by mutation) from the \textbf{ancestral allele} (the older
one, typically inferred by comparison to an outgroup
species)\textsuperscript{189}. This distinction matters when computing
the site frequency spectrum and interpreting patterns of variation.

Because we typically work with diploid organisms, every individual
carries two alleles at each autosomal locus. The \textbf{allele
frequency} is the proportion of all copies of the locus in the
population that carry a given allele. If we genotype \(N\) diploid
individuals at a biallelic locus. Based on some locus \(A\), we can
denote the number of homozygous individuals for the \(A_1\) allele as
\(N_{11}\) and the number of heterozygous individuals as \(N_{12}\),
meaning that they have both the \(A_1\) allele and the \(A_2\) allele.
The frequency of \(A_1\), is then based on the number of \(N_{11}\)
homozygotes for allele \(A_1\) and \(N_{12}\) heterozygotes, which is
calculated using the following equation:

\[
p = \frac{2N_{11} + N_{12}}{2N},
\]

and the frequency of the alternate allele \(A_2\) is simply
\(q = 1 - p\). This formula follows directly from counting allele copies
and holds regardless of Hardy--Weinberg equilibrium.

\subsection{Allozymes and the discovery of abundant molecular
variation}\label{allozymes-and-the-discovery-of-abundant-molecular-variation}

The advent of protein gel electrophoresis in the 1960s brought the first
systematic surveys of molecular genetic variation. Hubby and Lewontin,
working in \emph{Drosophila pseudoobscura}, and Harris, in humans, used
\textbf{allozyme} variation --- electrophoretically distinguishable
forms of enzymes --- to estimate heterozygosity and the proportion of
polymorphic loci without pre-selecting genes known to be
variable\textsuperscript{190}. Their surveys were a revelation: many
loci were polymorphic and individuals were heterozygous at a substantial
fraction of loci, in stark contrast to the earlier view that wild
populations were nearly monomorphic\textsuperscript{190}.

These studies introduced key summary statistics that remain central
today\textsuperscript{190}:

\begin{itemize}
\tightlist
\item
  \textbf{Proportion polymorphic \(P\)}: The fraction of loci with more
  than one allele in the sample.
\item
  \textbf{Average heterozygosity \(\bar{H}\)}: The expected probability
  that a randomly chosen individual is heterozygous at a randomly chosen
  locus, averaged across loci.
\end{itemize}

A striking and puzzling pattern emerged: levels of allozyme variation
were broadly similar across species with vastly different census
population sizes\textsuperscript{190,205}. A sea squirt with billions of
individuals and a lynx with far fewer showed far less difference in
heterozygosity than their census sizes would predict. This helped
motivate the \textbf{neutral theory of molecular evolution} (Kimura,
1968), which proposes that most molecular variants are selectively
neutral or nearly neutral and that genetic drift, rather than natural
selection, governs their frequencies\textsuperscript{189,205}.

Allozymes were also instrumental in formalizing population structure.
Wright's \(F_{ST}\) quantifies genetic differentiation among
populations. Lewontin's early applications to human allozymes showed
that the vast majority (\textasciitilde85\%) of human genetic variation
occurs within populations rather than between them --- a finding with
profound implications for how we understand the relationship between
human genetic diversity and socially constructed racial
categories\textsuperscript{190}.

\subsection{From allozymes to DNA sequences and the ADH
locus}\label{from-allozymes-to-dna-sequences-and-the-adh-locus}

While allozymes provided a first glimpse of molecular variation, they
detected only a subset of amino-acid changes and gave no information
about synonymous or noncoding variation. DNA sequencing changed this.
Kreitman's 1983 study of nucleotide polymorphism at the \emph{Drosophila
melanogaster} alcohol dehydrogenase (\emph{Adh}) locus was among the
first direct views of DNA-level polymorphism in a natural population,
revealing extensive synonymous and nonsynonymous
variation\textsuperscript{190}. \emph{Drosophila} was well-positioned
for these early molecular studies in part because its giant polytene
chromosomes --- visible in larval salivary glands --- had long allowed
cytogeneticists to map genes and inversions directly onto physical
chromosome structure\textsuperscript{206}. This cytogenetic foundation
identified conserved chromosome arms across \emph{Drosophila} species
(Muller's elements), providing a comparative scaffold that would later
anchor genome-wide sequence comparisons --- one of the earliest examples
of what we now call comparative genomics\textsuperscript{206}.

The \emph{Adh} locus illustrates the vocabulary of molecular variation
clearly. At some sites, individuals within \emph{D. melanogaster} differ
from each other --- these are \textbf{polymorphisms} (segregating
sites). At other sites, \emph{D. melanogaster} carries one allele while
related species (\emph{D. simulans}, \emph{D. yakuba}) carry another,
but variation is absent within each species --- these are \textbf{fixed
differences}\textsuperscript{189}. Comparing the ratio of nonsynonymous
to synonymous variants at polymorphic sites versus fixed differences
between species forms the basis of the McDonald--Kreitman test, one of
the most widely used tests for detecting adaptive protein
evolution\textsuperscript{190}.

More broadly, nucleotide sequence data from many loci established that:

\begin{itemize}
\tightlist
\item
  Most nonsynonymous variants are rare within populations, consistent
  with purifying selection removing them.
\item
  A small fraction of loci show elevated diversity at intermediate
  frequencies, consistent with balancing selection maintaining multiple
  alleles.
\item
  A substantial component of molecular evolution in some taxa (e.g.,
  \emph{Drosophila}) appears driven by positive selection, while in
  others (e.g., humans) adaptive substitutions appear less frequent.
\end{itemize}

These single-locus insights set the stage for asking the same questions
across entire genomes\textsuperscript{190,205}.

\subsection{The omics era and the rise of population
genomics}\label{the-omics-era-and-the-rise-of-population-genomics}

High-throughput sequencing technologies transformed population genetics
into population genomics. This transition into the ``omics era'' is
characterized by large-scale, quantitative, and integrative
analyses\textsuperscript{188}. Genome sequencing of non-model organisms
is now routine, and whole-genome resequencing of many individuals per
species enables true population genomics: simultaneous analysis of
polymorphism and divergence across every part of the
genome\textsuperscript{188}.

This genomic perspective enables researchers to ask questions that were
impossible with single loci:

\begin{itemize}
\tightlist
\item
  What fraction of the genome evolves under purifying selection versus
  neutrally?
\item
  How prevalent are adaptive substitutions genome-wide?
\item
  What demographic events --- bottlenecks, expansions, admixture ---
  have shaped a population's history, and where are their genomic
  signatures?
\item
  How does recombination rate variation shape the landscape of diversity
  and differentiation?
\end{itemize}

One striking insight from these surveys is that neutral diversity
\(\pi\) varies over a far narrower range across species than census
population sizes would predict\textsuperscript{205}. In a survey of 167
species across 14 phyla by Leffler et al.~(2012), \(\pi\) ranged only
about 800-fold --- from the Eurasian lynx (\(\pi \approx 0.01\%\), or
roughly 1 difference per 10,000 bases) to the sea squirt \emph{Ciona
savignyi} (\textasciitilde1 difference per 12 bases) --- even though
census population sizes span many more orders of
magnitude.\textsuperscript{189,205} This ``old riddle'' suggests that
forces beyond census size --- including natural selection reducing
effective population sizes --- constrain genetic diversity, and it
remains an active research frontier\textsuperscript{205}.

\section{Evolutionary Processes and Genomic Patterns of
Variation}\label{evolutionary-processes-and-genomic-patterns-of-variation}

Classical population genetics provides simple null models for how allele
and genotype frequencies behave under specified assumptions, and
population genomics extends these ideas across whole genomes. The goal
in this section is to connect key statistics --- \(\pi\), Tajima's D,
and \(F_{ST}\) --- to the underlying evolutionary processes and to show
how these forces leave recognizable signatures in genomic
data\textsuperscript{188,189}. It is worth noting that many of these
calculations can be on genomics datasets using toolkits like
\textbf{VCFtools} introduced in Chapter 3, working directly with VCF
files.

\subsection{Hardy--Weinberg equilibrium: a null model for genotype
frequencies}\label{hardyweinberg-equilibrium-a-null-model-for-genotype-frequencies}

The simplest population model assumes random mating, infinite population
size, and no mutation, migration, or selection. Under these conditions,
genotype frequencies are entirely predictable from allele frequencies
--- a result known as \textbf{Hardy--Weinberg equilibrium
(HWE)}\textsuperscript{189}.

Consider a biallelic locus where allele \(A_1\) is at frequency \(p\)
and allele \(A_2\) is at frequency \(q = 1 - p\). In a randomly mating
population, gametes fuse independently, so the probability of drawing
two \(A_1\) copies is \(p \times
p = p^2\). The three genotype frequencies are therefore:

\[
P(A_1A_1) = p^2, \quad P(A_1A_2) = 2pq, \quad P(A_2A_2) = q^2.
\]

A key subtlety is that we only need to assume \textbf{random mating at
the focal locus} for these frequencies to hold among newly formed
zygotes\textsuperscript{189}. Natural selection, genetic drift, and
other evolutionary processes can be acting elsewhere in the genome
without disrupting Hardy--Weinberg at any given locus, as long as
individuals mate without regard to their genotype at that location. This
is why HWE often holds remarkably well in natural populations:
individuals do not mate globally at random, but they typically do mate
at random with respect to the genotype at most individual
loci\textsuperscript{189}.

\textbf{Departures from HWE are informative:} An excess of homozygotes
relative to HWE expectations (i.e., a deficit of heterozygotes) can
arise from inbreeding, positive assortative mating, or population
subdivision (the Wahlund effect). An excess of heterozygotes can arise
from disassortative mating or balancing selection\textsuperscript{189}.
In bioinformatics pipelines, HWE tests are also used as quality control:
large-scale deviations at many loci can indicate genotyping errors,
sample contamination, or data artifacts rather than real biology.

\textbf{Generalized HWE and inbreeding:} When related individuals mate,
their offspring can inherit two copies of the same allele descended from
a common ancestor. We call such alleles \textbf{identical by descent
(IBD)}, and an individual carrying them is said to be \textbf{homozygous
by descent} or \textbf{autozygous}\textsuperscript{189}. We quantify
inbreeding using the \textbf{inbreeding coefficient} \(F\), the
probability that both alleles in an individual are IBD. When \(F > 0\),
genotype frequencies deviate from standard HWE in a predictable way:

\[
P(A_1A_1) = (1-F)p^2 + Fp, \quad
P(A_1A_2) = (1-F)\,2pq, \quad
P(A_2A_2) = (1-F)q^2 + Fq.
\]

Inbreeding inflates homozygote frequencies and reduces heterozygosity.
In conservation genetics this matters enormously: the Mexican wolf
(\emph{Canis lupus baileyi}), reduced to just five wild individuals
before a captive breeding program was established, now has an estimated
average inbreeding coefficient of \(\hat{F} \approx 0.33\), meaning
roughly a third of each wolf's genome is homozygous by
descent\textsuperscript{189}.

\textbf{Runs of homozygosity:} Because related parents transmit alleles
in large intact blocks --- recombination has had few generations to
break up IBD segments --- inbred individuals show long genomic
\textbf{runs of homozygosity (ROH)}: stretches of the genome where an
individual is homozygous at essentially every SNP\textsuperscript{189}.
The proportion of the genome in ROH segments (\(F_{ROH}\)) estimates the
inbreeding coefficient, and the length distribution of ROH encodes
timing information: very long ROH reflect recent inbreeding, while many
shorter ROH reflect accumulation over many
generations.\textsuperscript{189} In a study of 2,500 dogs across many
breeds, Doberman Pinschers showed more genome in very long ROH
(reflecting recent intensive breeding), while English bulldogs had more
shorter ROH (reflecting a longer cumulative
history)\textsuperscript{189}. ROH analysis is now standard in
conservation and livestock genomics.

\subsection{Measures of genetic variability at the sequence
level}\label{measures-of-genetic-variability-at-the-sequence-level}

When we sequence DNA from multiple individuals, several complementary
statistics summarize the genetic variation we observe. I have recorded a
\href{https://auburn.hosted.panopto.com/Panopto/Pages/Viewer.aspx?id=bf9c8b7c-92e2-436d-8ddf-b40a0139e33b}{video
explanation} of how to calculate several of the measures introduced
below.

\textbf{Nucleotide diversity (\(\pi\)):} The most widely used measure is
\textbf{nucleotide diversity \(\pi\)}, defined as the average number of
pairwise nucleotide differences per site between two randomly chosen
sequences from the population\textsuperscript{189}. If at site \(i\) the
derived allele has estimated frequency \(\hat{p}_i\) across \(L\) sites,
a convenient estimator is:

\[
\hat{\pi} = \frac{1}{L} \sum_{i=1}^{L} 2\hat{p}_i(1 - \hat{p}_i).
\]

Under a standard neutral model with constant effective population size
\(N_e\) and per-site mutation rate \(\mu\) per generation, the expected
value is:

\[
E[\pi] = 4N_e\mu.
\]

This relationship is fundamental: it means \(\pi\) provides a window
into the product \(\theta = 4N_e\mu\), sometimes called the
\textbf{population-scaled mutation rate}. Species with larger long-term
\(N_e\) are expected to harbor more diversity, because more individuals
means more opportunities for new mutations to arise and persist before
being lost to drift\textsuperscript{189}.

\textbf{Watterson's estimator:} A complementary measure of variation is
the total number of \textbf{segregating sites} \(S\) in the sample. The
number of segregating sites grows with sample size, so we standardize
using \textbf{Watterson's estimator}:

\[
\hat{\theta}_W = \frac{S}{a_n}, \quad \text{where } a_n = \sum_{k=1}^{n-1}
\frac{1}{k}.
\]

Both \(\hat{\pi}\) and \(\hat{\theta}_W\) estimate the same underlying
quantity \(\theta = 4N_e\mu\) under the standard neutral model, but they
weight variants differently. \(\hat{\pi}\) gives more weight to variants
at intermediate frequency (because those contribute more to pairwise
differences), while \(\hat{\theta}_W\) gives equal weight to every
segregating site, including rare singletons. This difference is the
foundation of Tajima's D\textsuperscript{190}.

\textbf{The site frequency spectrum:} Beyond counting polymorphic sites,
we can ask about the \textbf{distribution of allele frequencies} across
sites --- a summary called the \textbf{site frequency spectrum
(SFS)}\textsuperscript{189}. We count how many sites in our sample have
a derived allele frequency of \(1/n\), \(2/n\), \(3/n\), and so on up to
\((n-1)/n\). Alleles found in only one copy are called
\textbf{singletons}; those in two copies, \textbf{doubletons}, and so
on. Under neutrality with constant population size, the SFS has a
characteristic shape: many rare variants and fewer variants at high
frequency. Deviations from this shape encode information about
demographic history and selection, and the SFS underlies many of the
most powerful modern methods for demographic inference.

\subsection{Genetic drift and effective population
size}\label{genetic-drift-and-effective-population-size}

\textbf{Genetic drift} is the random change in allele frequencies caused
by the finite number of individuals that contribute genes to the next
generation\textsuperscript{189,205}. In any generation, only a sample of
all possible gametes actually form the next generation, and by chance
some alleles leave more descendants than others --- regardless of their
fitness effects. The strength of drift is inversely proportional to
population size: in a large population, random fluctuations average out
and allele frequencies change slowly; in a small population, random
changes can be large.

The clearest consequence of drift is the \textbf{loss of heterozygosity}
over time. In the absence of new mutations, drift will eventually fix
one allele and eliminate all others. The rate of this loss is:

\[
H_t = \left(1 - \frac{1}{2N_e}\right)^t H_0 \approx H_0\, e^{-t/2N_e},
\]

where \(H_0\) is the initial heterozygosity and \(H_t\) is
heterozygosity after \(t\) generations\textsuperscript{189}. For a
population of \(N_e = 1{,}000\), heterozygosity falls to about 37\% of
its original value after roughly 2,000 generations. For \(N_e = 100\),
the same decline takes only \textasciitilde200 generations.

The conservation implications are stark. The black-footed ferret
(\emph{Mustela nigripes}) crashed to just 18 individuals in 1985 (of
which only 7 reproduced), causing a sharp drop in heterozygosity that is
still visible today even after the population has recovered to over 300
individuals.\textsuperscript{189} Because all living black-footed
ferrets are descended from those 7 founders, genetic diversity remains
low and this reduced variation limits the population's long-term
adaptive potential.

\textbf{Effective population size (\(N_e\)):} Real populations rarely
behave like an idealized constant-size population of \(N\) individuals
all contributing equally to reproduction. Variance in reproductive
success, fluctuations in census size, unequal sex ratios, and geographic
structure all cause genetic drift to proceed faster than the census size
\(N\) would suggest. The \textbf{effective population size} \(N_e\) is
defined as the size of an idealized Wright--Fisher population that would
experience the same rate of random allele-frequency change as the real
population\textsuperscript{189}. Almost always, \(N_e < N\).

\(N_e\) is strongly affected by population bottlenecks. Because the
harmonic mean of a set of numbers is dominated by the smallest values, a
population that numbers one million individuals 99\% of the time but
crashes to 1,000 every hundred generations will have an \(N_e\) far
closer to 1,000 than to one million.\textsuperscript{189} For humans,
the global census size exceeds 8 billion, but \(N_e\) estimated from
genomic diversity is only about 10,000--30,000, reflecting ancient
bottlenecks (including the out-of-Africa dispersal) and variance in
reproductive success\textsuperscript{189,205}.

Importantly, \(N_e\) also differs among parts of the genome. Autosomes,
the X chromosome, the Y chromosome, and mitochondrial DNA (mtDNA) have
different modes of inheritance and different sensitivities to variance
in male versus female reproductive success\textsuperscript{189}. In many
animal species, fewer males than females reproduce successfully (e.g.,
in harem-forming species), which reduces the effective number of
contributing Y chromosomes relative to mtDNA. A striking demonstration
comes from ancient horse DNA: comparing Scythian (2,300 year-old) to
modern horses, levels of Y chromosome diversity collapsed far more
dramatically than mtDNA diversity, reflecting intense selection of a
small number of stallions during domestication\textsuperscript{189}.

\textbf{Mutation--drift balance:} When drift removes variation while
mutation continually introduces new alleles, these forces can reach a
dynamic equilibrium. At equilibrium, the expected heterozygosity is:

\[
H \approx 4N_e\mu = \theta,
\]

the same population-scaled mutation rate estimated by \(\hat{\pi}\) and
\(\hat{\theta}_W\). Notice that diversity depends on the \emph{product}
of \(N_e\) and \(\mu\), not on either alone. This is partly why species
with similar census sizes can have very different \(\pi\), and why the
narrow range of \(\pi\) across species with very different census sizes
is so puzzling --- additional factors such as linked selection (see
below) must be suppressing \(N_e\) in large-census-size species.

\subsection{Tajima's D and departures from the standard neutral
model}\label{tajimas-d-and-departures-from-the-standard-neutral-model}

\textbf{Tajima's D} exploits the fact that \(\hat{\pi}\) and
\(\hat{\theta}_W\) estimate the same \(\theta\) under the standard
neutral model but are sensitive to different parts of the allele
frequency distribution\textsuperscript{189}. A
\href{https://www.youtube.com/watch?v=wiyay4YMq2A}{great video
explanation} of Tajima's D and its relation to the site frequency
spectrum was recorded a while back by my own PhD Mentor Dr.~Mohamed
Noor. The statistic is calculated according to the following formula:

\[
D = \frac{\hat{\pi} - \hat{\theta}_W}{\sqrt{\operatorname{Var}(\hat{\pi} -
\hat{\theta}_W)}}.
\]

When the site frequency spectrum matches neutral expectations,
\(D \approx 0\). Departures in either direction are informative:

\begin{itemize}
\tightlist
\item
  \textbf{Negative \(D\)} (excess of rare variants): occurs after recent
  population expansion (new mutations enter as singletons), after a
  recent positive selective sweep (nearby variation is eliminated and
  then regenerates as rare variants), or under ongoing purifying
  selection\textsuperscript{188,189}.
\item
  \textbf{Positive \(D\)} (deficit of rare variants, excess at
  intermediate frequencies): occurs after a population bottleneck (rare
  alleles are lost preferentially), or under balancing selection (which
  actively maintains alleles at intermediate frequencies for long
  periods)\textsuperscript{188,189}.
\end{itemize}

Empirical examples anchor these interpretations. Genome-wide Tajima's D
in many human populations tends to be slightly negative, consistent with
recent population expansion that produced many new rare
variants\textsuperscript{205}. Maize (\emph{Zea mays}), domesticated
from wild teosinte roughly 10,000 years ago, shows a shift toward
positive Tajima's D relative to teosinte, consistent with the
domestication bottleneck eliminating rare alleles and enriching
intermediate-frequency variants\textsuperscript{189}. At loci under
balancing selection --- such as some MHC immune-gene loci in vertebrates
--- local Tajima's D can be strongly positive with elevated \(\pi\)
compared to neighboring neutral regions\textsuperscript{189,190}.

In genomic analyses, Tajima's D is typically computed in \textbf{sliding
windows} across chromosomes, producing a landscape of departures from
neutrality that can be compared to recombination rates, gene density,
and functional annotations. Windows with strongly negative \(D\) and
reduced \(\pi\) are candidate selective sweep regions; windows with
strongly positive \(D\) are candidate balancing selection or bottleneck
regions.

\subsection{\texorpdfstring{Population structure and
\(F_{ST}\)}{Population structure and F\_\{ST\}}}\label{population-structure-and-f_st}

Population subdivision is a pervasive feature of natural populations.
Individuals tend to mate with nearby neighbors rather than with randomly
chosen members of the entire species, and physical or ecological
barriers can further divide populations. The degree to which allele
frequencies have diverged among subpopulations is measured by
\(F_{ST}\)\textsuperscript{190}.

One standard formulation compares heterozygosity within versus across
populations:

\[
F_{ST} = \frac{H_T - H_S}{H_T},
\]

where \(H_T\) is the expected heterozygosity in the total pooled
population and \(H_S\) is the average expected heterozygosity within
subpopulations\textsuperscript{190}. Values range from 0 (no
differentiation; all populations have identical allele frequencies) to 1
(complete fixation of different alleles in different populations, so all
diversity is among populations rather than within them).

\textbf{Intuition from identity by descent:} \(F_{ST}\) can also be
understood through inbreeding. When a population is subdivided, an
individual within one subpopulation is more likely to mate with another
member of the same subpopulation than with a random individual from the
whole species. From the perspective of the total population, individuals
within subpopulations are therefore somewhat ``inbred'' relative to the
global average. \(F_{ST}\) measures the deficit of heterozygosity within
subpopulations compared to the global expectation --- in other words,
the fraction of total genetic variance explained by differences among
populations\textsuperscript{189}.

Under a simple model where one ancestral population of size \(N_e\)
split \(T\) generations ago into two isolated daughter populations with
no subsequent gene flow, \(F_{ST}\) grows over time as drift
differentiates allele frequencies:

\[
F_{ST} \approx \frac{T}{T + 4N_e}.
\]

This shows that differentiation depends on the ratio of divergence time
to \(N_e\): populations diverged long ago, or with small \(N_e\),
accumulate higher \(F_{ST}\).

\textbf{Wright's F-statistics hierarchy:} Wright extended this to a
three-level hierarchy. \(F_{IS}\) measures inbreeding of individuals
relative to their subpopulation; \(F_{ST}\) measures differentiation of
subpopulations relative to the total; \(F_{IT}\) measures overall
inbreeding of individuals relative to the total population. These
satisfy:

\[
(1 - F_{IT}) = (1 - F_{IS})(1 - F_{ST}).
\]

\textbf{Empirical values:} Among major human continental groups,
genome-wide \(F_{ST}\) is typically 0.05--0.15, indicating that most
genetic variation is within, not between,
populations\textsuperscript{205}. Among strongly isolated island
populations or distinct ecotypes, \(F_{ST}\) can exceed
0.3\textsuperscript{188}. Within a species, \(F_{ST}\) can vary sharply
along the genome: regions containing inversions, structural variants, or
loci under divergent selection often show elevated \(F_{ST}\) compared
to the genomic background, making windowed \(F_{ST}\) scans a standard
first step in searching for regions involved in local
adaptation\textsuperscript{188}.

\subsection{Recombination, linkage disequilibrium, and linked
selection}\label{recombination-linkage-disequilibrium-and-linked-selection}

Recombination each generation breaks down associations between alleles
at different loci. The degree of non-random association between alleles
at different loci is called \textbf{linkage disequilibrium (LD)}. For a
pair of biallelic loci with alleles \(A\)/\(a\) and \(B\)/\(b\), the LD
coefficient is:

\[
D_{AB} = p_{AB} - p_A p_B,
\]

where \(p_{AB}\) is the haplotype frequency of
\(AB\)\textsuperscript{189}. If the loci are statistically independent,
\(D = 0\). LD decays over time as:

\[
D_t = (1 - c)^t D_0,
\]

where \(c\) is the recombination fraction between the two
loci\textsuperscript{189}. For tightly linked loci, LD can persist for
thousands of generations; for loci on different chromosomes
(\(c = 0.5\)), LD is halved every generation. A striking demonstration
of LD decay comes from Neanderthal admixture: modern Eurasians carry
\textasciitilde1--4\% Neanderthal-derived sequence, and LD between
Neanderthal-ancestry alleles in modern Europeans decays exponentially
with genetic distance in a way consistent with admixture
\textasciitilde35,000--47,000 years ago\textsuperscript{189}.

In practice, LD is often summarized using \(r^2\), the squared
correlation of allele states between loci. Genome sequences and linkage
maps together reveal highly heterogeneous recombination landscapes with
hotspots and cold spots\textsuperscript{188}.

\textbf{Linked selection:} Recombination also modulates how selection at
one site affects neutral variation at nearby sites --- a phenomenon
called \textbf{linked selection}. Two major forms
are\textsuperscript{188,189}:

\begin{itemize}
\tightlist
\item
  \textbf{Background selection}: Purifying selection against deleterious
  alleles reduces diversity at linked neutral sites. When a deleterious
  allele is eliminated, it drags nearby neutral variation to extinction
  along with it.
\item
  \textbf{Selective sweeps}: Positive selection on a beneficial allele
  causes nearby neutral variation to ``hitchhike'' to high frequency or
  fixation, reducing diversity in a window around the selected site.
\end{itemize}

In regions of low recombination, both forms can substantially depress
neutral diversity, creating ``cold spots'' in the landscape of \(\pi\)
even when mutation rates are not unusually
low\textsuperscript{188,189,205}. In \emph{Drosophila}, neutral
diversity is often lowest near centromeres and telomeres, consistent
with strong background selection and sweeps in low-recombination
regions\textsuperscript{188}. This has a practical consequence for
interpreting population genomic scans: a dip in \(\pi\) near a gene does
not necessarily mean \emph{that} gene was the direct target of selection
--- it may simply be in LD with a nearby selected site.

\subsection{Natural selection: how fitness differences change allele
frequencies}\label{natural-selection-how-fitness-differences-change-allele-frequencies}

\textbf{Natural selection} changes allele frequencies non-randomly,
based on fitness differences among genotypes. We quantify these
differences using the \textbf{selection coefficient} \(s\), which
measures the difference in relative fitness between
alleles\textsuperscript{189}.

In the simplest case --- a haploid population where allele \(A_1\) has
fitness 1 and \(A_2\) has fitness \(1 - s\) --- the frequency of \(A_1\)
after \(\tau\) generations starting from initial frequency \(p_0\) is
approximately:

\[
p_\tau \approx \frac{p_0}{p_0 + q_0\, e^{-s\tau}}.
\]

This logistic equation produces the classic S-shaped trajectory of a
beneficial allele rising toward fixation. The allele increases slowly
when rare (it is mostly in heterozygotes, where its advantage may be
diluted), accelerates as it becomes common, then slows again as the
alternate allele nears loss.

\textbf{How long does a sweep take?} The time for a beneficial allele to
rise from a single copy to near-fixation in a population of size \(N\)
is approximately:

\[
\tau \approx \frac{2}{s}\log(N).
\]

For a selection coefficient of \(s = 0.01\) (a 1\% fitness advantage) in
a population of \(N = 10{,}000\), this is about 1,840 generations. The
real-world example of \textbf{lactase persistence} in humans illustrates
this well: the allele allowing lactase expression in adults was absent
in Central Europe more than 5,000 years ago and is now at frequencies
above 70\% in many European populations, consistent with a strong
selection coefficient (\(s \approx 0.01\)--\(0.05\)) that arose with
dairy farming\textsuperscript{189}.

\textbf{Diploid selection and dominance.} In diploid organisms, the
fitness of a heterozygote need not be intermediate between the two
homozygotes. The \textbf{dominance coefficient} \(h\) captures this:

\begin{longtable}[]{@{}llll@{}}
\toprule\noalign{}
Genotype & \(A_1A_1\) & \(A_1A_2\) & \(A_2A_2\) \\
\midrule\noalign{}
\endhead
\bottomrule\noalign{}
\endlastfoot
Fitness & \(1\) & \(1 - hs\) & \(1 - s\) \\
\end{longtable}

When \(h = 0\), allele \(A_1\) is fully dominant (the heterozygote has
the same fitness as the better homozygote). When \(h = 1\), \(A_1\) is
fully recessive\textsuperscript{189}. Dominance profoundly affects how
selection acts on rare alleles: a recessive beneficial allele is nearly
invisible to selection when rare because it is almost always found in
heterozygotes.

\textbf{Balanced polymorphisms:} Not all selection drives alleles to
fixation or loss. When heterozygotes have higher fitness than either
homozygote --- \textbf{heterozygote advantage} or overdominance ---
selection maintains both alleles at a stable equilibrium. The classic
example is sickle-cell anemia: heterozygotes are protected from malaria
relative to both the normal homozygote (malaria susceptible) and the
sickle-cell homozygote (severe anemia). A second compelling example
comes from Soay sheep: two alleles at the \emph{RXFP2} locus affect horn
size, with each homozygote suffering fitness costs at different
life-history stages (viability vs.~reproductive success), while
heterozygotes achieve the highest combined fitness\textsuperscript{189}.
Both alleles are maintained at intermediate frequency, and this type of
locus shows elevated \(\pi\) and positive Tajima's D --- the genomic
signature we expect under balancing selection.

\subsection{Drift, migration, and demographic
history}\label{drift-migration-and-demographic-history}

Genetic drift and migration jointly shape the distribution of variation
within and among populations. Drift reduces diversity over time and can
fix neutral or mildly deleterious alleles; migration (gene flow) tends
to homogenize allele frequencies and counteract
differentiation\textsuperscript{189}. Their interplay generates a
spectrum of patterns from smooth \textbf{isolation-by-distance} clines
--- where genetic differentiation increases gradually with geographic
separation --- to sharp \textbf{contact zones} between diverging
lineages\textsuperscript{188}.

Population genomics leverages whole-genome data to reconstruct
\textbf{demographic histories} by examining the site frequency spectrum,
distributions of pairwise coalescence times, and patterns of haplotype
sharing\textsuperscript{188}. Bottlenecks, expansions, and admixture all
leave characteristic genomic signatures. Because these signatures can
sometimes resemble those produced by selection at particular loci, it is
essential to account for demography when interpreting outlier regions in
genome-wide scans.

\textbf{Coalescent theory} provides the formal framework for relating
patterns of genomic variation to demographic history by reasoning
backward in time\textsuperscript{189}. Rather than tracking allele
frequencies forward through generations, coalescent theory asks: looking
backward, when did the ancestors of our sampled sequences last share a
common ancestor? This event is called a \textbf{coalescence}. In a
population of effective size \(N_e\), the expected coalescence time for
a pair of lineages is \(2N_e\) generations. For a sample of \(k\)
lineages, the expected time until the first coalescence event is:

\[
E[T_k] = \frac{2N_e}{\binom{k}{2}}.
\]

The total expected time to the most recent common ancestor (MRCA) of all
\(n\) sampled sequences is approximately \(4N_e(1 - 1/n)\) generations.
A key insight is that neutral mutations do not affect the coalescent
process --- they are simply scattered onto the coalescent tree at
random. This means we can think about the tree structure first (shaped
by drift and demography) and then superimpose mutations (shaped by
mutation rate)\textsuperscript{189}.

Long-term changes in population size alter coalescence rates in
predictable ways. A recent population bottleneck accelerates coalescence
(many lineages find common ancestors rapidly in the recent past), while
population expansion slows it (lineages are slow to coalesce in the
larger ancestral population). Methods such as PSMC, SMC++, and
fastsimcoal2 exploit these relationships to infer detailed demographic
histories from genomic data.

Under the island-mainland model --- a small island population of size
\(N_I\) receiving migrants from a large mainland at rate \(m\) per
generation --- differentiation between island and mainland reaches an
equilibrium where:

\[
F_{IM} = \frac{1}{1 + 4N_I m}.
\]

This elegant result shows that even a small number of migrants per
generation (\(N_I m \gtrsim 1\)) is enough to substantially suppress
differentiation. The key parameter is not the migration rate alone, but
the product of effective population size and migration rate.

\begin{tcolorbox}[enhanced jigsaw, rightrule=.15mm, leftrule=.75mm, colbacktitle=quarto-callout-note-color!10!white, colframe=quarto-callout-note-color-frame, toprule=.15mm, opacitybacktitle=0.6, title=\textcolor{quarto-callout-note-color}{\faInfo}\hspace{0.5em}{Lab Connection: 1000 Genomes Sliding Windows}, left=2mm, colback=white, bottomtitle=1mm, toptitle=1mm, bottomrule=.15mm, coltitle=black, titlerule=0mm, breakable, opacityback=0, arc=.35mm]

In Lab 10, you will use an R Shiny app that I built based on 1000
Genomes data to visualize sliding‑window estimates of \(F_{ST}\),
\(\pi\), and Tajima's D across several human genomic regions that were
generated using VCFtools. For each assigned region, you will:

\begin{itemize}
\tightlist
\item
  Examine windowed \(F_{ST}\) between population pairs to identify peaks
  of differentiation.\\
\item
  Compare windowed \(\pi\) and Tajima's D within populations to assess
  diversity and allele frequency shifts.\\
\item
  Use the UCSC Genome Browser to connect these statistical patterns to
  specific genes and known aspects of human evolution.
\end{itemize}

This lab provides a concrete example of how the summary statistics
introduced in this section are used in practice to link evolutionary
processes, genomic variation, and gene function.

\end{tcolorbox}

\section{Sliding Window Analyses, Genome-wide Scans, and Data
Challenges}\label{sliding-window-analyses-genome-wide-scans-and-data-challenges}

Population genomics often proceeds by scanning the genome for regions
where summary statistics deviate from a background pattern, then asking
what biological processes might explain those deviations. Sliding‑window
approaches are central to this strategy: they smooth noisy site‑level
statistics, highlight clusters of extreme values, and provide a natural
bridge between population‑genetic theory and genome
annotation\textsuperscript{188,207}.

In a sliding‑window analysis, the genome is divided into consecutive
windows, defined either by physical distance (for example, 50 kb) or by
a fixed number of SNPs (for example, 100 variants). A statistic of
interest---such as mean \(F_{ST}\), average \(\pi\), or Tajima's D---is
calculated in each window, and the window is then moved along the
chromosome, often with overlap, to produce a continuous
profile\textsuperscript{207}. Because evolutionary processes such as
selection and linked selection typically affect stretches of the genome
rather than isolated sites, window‑based statistics are well suited to
detecting their footprints.

Different evolutionary scenarios lead to characteristic window‑level
signatures. Regions affected by recent positive selection in one
population may show a local reduction in \(\pi\), negative Tajima's D,
and elevated \(F_{ST}\) compared to the genome‑wide background. Regions
under balancing selection can show elevated \(\pi\) and positive
Tajima's D, and sometimes lower \(F_{ST}\) if similar polymorphisms are
maintained across populations. Regions strongly influenced by background
selection may display reduced \(\pi\) without dramatic changes in
\(F_{ST}\) or Tajima's D\textsuperscript{188,207}. Genome‑wide scans
exploit these patterns by computing many such statistics in windows and
then flagging windows that fall in the tails of their empirical
distributions.

Haasl and Payseur review a wide variety of statistics used in
genome‑wide scans for selection and emphasize that no single metric is
definitive\textsuperscript{207}. Site‑frequency‑based tests (such as
Tajima's D), haplotype‑based tests, and differentiation measures (like
\(F_{ST}\) or related quantities) each respond differently to selection
and demography. Combining them can increase power and specificity. For
example, a region with high \(F_{ST}\) but otherwise unremarkable
\(\pi\) and Tajima's D might reflect drift or demographic history,
whereas a region with high \(F_{ST}\), strongly reduced \(\pi\) in one
population, and negative Tajima's D is more consistent with a localized
selective sweep.

Practical choices about window design matter. Larger windows average
over more sites, reducing sampling noise but potentially diluting narrow
signals of selection; smaller windows increase resolution but yield
noisier estimates and more extreme values by
chance\textsuperscript{207}. Windows defined by physical distance align
naturally with gene models and recombination maps, while SNP‑based
windows equalize the number of variants per window but can span very
different physical lengths in regions of high or low SNP density.
Overlapping windows produce smoother profiles but create strong
correlations among neighboring windows; non‑overlapping windows reduce
this correlation but may slice through true
signals\textsuperscript{207}.

Once candidate windows have been identified, the next step is to connect
them to biological features. A typical workflow is to merge adjacent
outlier windows into broader candidate regions, intersect these regions
with gene annotations and regulatory elements, and then prioritize genes
based on known functions, expression patterns, or previous association
and QTL studies\textsuperscript{188,207}. In many systems, such
genome‑wide scans have successfully highlighted loci involved in traits
such as high‑altitude adaptation, host--pathogen interactions,
pigmentation, and domestication.

\subsection{Missing data, ascertainment, and robustness of window-based
scans}\label{missing-data-ascertainment-and-robustness-of-window-based-scans}

The reliability of sliding‑window analyses depends strongly on data
quality and study design. Missing genotypes, uneven coverage, and
ascertainment bias can all distort estimates of \(F_{ST}\), \(\pi\),
Tajima's D, and related statistics\textsuperscript{207,208}. These
issues are particularly important in reduced‑representation datasets
(such as RADseq, GBS, and SNP chips), low‑coverage sequencing, and
ancient DNA.

Missing data arise for many reasons. Random dropout, low read depth, and
sequencing errors can cause some genotypes to be uncalled or miscalled.
Reduced‑representation methods sample only a subset of the genome by
design, often with allele dropout at polymorphic restriction sites. Some
variants may also be effectively ``missing'' because they occur in
unsampled or extinct ``ghost'' populations that have contributed
ancestry to the sampled groups\textsuperscript{208}. Each of these
mechanisms can affect window‑based statistics in different ways.

Systematic differences in coverage or missingness among populations can
bias between‑population statistics. If one population has lower average
depth or a higher fraction of missing genotypes, variant callers may
preferentially retain high‑frequency alleles in that population, making
it appear less diverse and exaggerating \(F_{ST}\) in some
windows\textsuperscript{207,208}. Ascertainment bias---such as SNPs
discovered in a subset of populations and then genotyped elsewhere---can
skew the site frequency spectrum, reducing the apparent number of rare
variants and altering \(\hat{\theta}_W\), \(\hat{\pi}\), and Tajima's
D\textsuperscript{207}.

Common responses to missing data include discarding sites or individuals
with too much missingness, treating missing genotypes as random and
ignoring them in likelihood calculations, or imputing missing genotypes
from reference panels or local haplotype structure\textsuperscript{208}.
These strategies can work reasonably well when missingness is modest and
approximately random, but can introduce biases when missingness is
structured across populations or along the genome. Recent work has
developed corrected estimators of summary statistics that explicitly
account for genotyping errors and missing data, and methods that
integrate missingness into demographic or selection
models\textsuperscript{208}.

For sliding‑window analyses, these considerations suggest several best
practices. Filters on depth, genotype quality, and call rate should be
applied consistently across populations, and sensitivity analyses should
be used to check whether apparent peaks in \(F_{ST}\) or dips in \(\pi\)
are robust to reasonable changes in filtering
thresholds.\textsuperscript{207} Where possible, neutral simulations or
resampling approaches that reproduce the observed missing‑data pattern
can be used to assess how often similar extrema would arise by
chance\textsuperscript{207,208}. Interpreting outlier windows is most
convincing when statistical signals align with regions of good data
quality and with independent evidence from gene function, expression, or
other genomic features.

Overall, sliding‑window analyses, genome‑wide scans for selection, and
careful attention to missing data together form a powerful framework for
reading the evolutionary history of populations from their genomes. The
same summary statistics introduced earlier in the chapter---\(\pi\),
Tajima's D, \(F_{ST}\), and related measures---play central roles in
these analyses, but their interpretation always depends on the
biological context, the demographic background, and the limitations of
the data.

\begin{tcolorbox}[enhanced jigsaw, rightrule=.15mm, leftrule=.75mm, colbacktitle=quarto-callout-note-color!10!white, colframe=quarto-callout-note-color-frame, toprule=.15mm, opacitybacktitle=0.6, title=\textcolor{quarto-callout-note-color}{\faInfo}\hspace{0.5em}{Placeholder: Case Study -- Human Population Genomics}, left=2mm, colback=white, bottomtitle=1mm, toptitle=1mm, bottomrule=.15mm, coltitle=black, titlerule=0mm, breakable, opacityback=0, arc=.35mm]

\textbf{To be added by the author.} This box will present a detailed
case study of human population genomics that parallels Lab 9. It will
walk through:

\begin{itemize}
\tightlist
\item
  Computing windowed \(\pi\), \(F_{ST}\), and Tajima's D in human
  populations.
\item
  Identifying genomic regions with signatures of selection and/or
  complex demography.
\item
  Integrating these results with knowledge of human evolutionary history
  and trait biology.
\end{itemize}

The case study will serve as a capstone example, demonstrating how the
concepts introduced in this chapter come together in a real‑world
dataset.

\end{tcolorbox}

\section{Problems}\label{problems}

\begin{enumerate}
\def\labelenumi{\arabic{enumi}.}
\tightlist
\item
  \textbf{Nucleotide diversity and heterozygosity:} Consider a sample of
  \(n\) haploid genomes from a population. For each of \(L\) SNPs, you
  estimate the derived allele frequency \(\hat{p}_i\).

  \begin{enumerate}
  \def\labelenumii{\alph{enumii}.}
  \tightlist
  \item
    Write down the estimator of nucleotide diversity \(\hat{\pi}\) in
    terms of the \(\hat{p}_i\). Explain in words what \(\hat{\pi}\)
    measures and how it relates to the concept of expected
    heterozygosity at a locus.\\
  \item
    Suppose that across a set of allozyme loci in a diploid species, the
    average heterozygosity \(\bar{H}\) is 0.12. What does this imply
    about the expected fraction of loci at which a randomly chosen
    individual is heterozygous?\\
  \item
    In a resequenced 1 kb gene region, you estimate
    \(\hat{\pi} = 0.004\) per site. How would you compare this gene's
    diversity to the genome‑wide allozyme heterozygosity estimate, and
    what additional context would you need to make a fair comparison?
  \end{enumerate}
\item
  \textbf{Interpreting Tajima's D:} You compute Tajima's D in 10 kb
  windows across a genome and observe many windows with strongly
  negative D values.

  \begin{enumerate}
  \def\labelenumii{\alph{enumii}.}
  \tightlist
  \item
    Describe at least two evolutionary scenarios (one demographic, one
    selective) that could produce genome‑wide negative Tajima's D.\\
  \item
    How might you distinguish between these scenarios using additional
    statistics or data (e.g., recombination maps, functional
    annotations, or outgroup comparisons)?\\
  \item
    In a specific genomic region, you observe negative Tajima's D
    together with markedly reduced \(\pi\) relative to the genomic
    background. Under what conditions would you interpret this as
    evidence for a recent selective sweep?
  \end{enumerate}
\item
  \textbf{\(F_{ST}\) and genomic landscapes of differentiation:} You
  have whole‑genome SNP data from two partially isolated populations.
  You compute windowed \(F_{ST}\) (100 kb windows) and find that most
  windows have low \(F_{ST}\) (\textasciitilde0.05), but a small
  fraction have much higher values (\textgreater0.3).

  \begin{enumerate}
  \def\labelenumii{\alph{enumii}.}
  \tightlist
  \item
    Describe at least two possible explanations for these
    high‑\(F_{ST}\) windows (one involving selection, one involving
    neutral processes).\\
  \item
    How could you use additional statistics such as \(\pi\), \(d_{XY}\),
    and Tajima's D to help distinguish between these explanations?\\
  \item
    What role might recombination rate variation play in generating
    heterogeneous \(F_{ST}\) across the genome?
  \end{enumerate}
\item
  \textbf{Sliding windows vs.~spline‑based windows:} You are designing a
  genome‑wide scan for selection using pooled sequencing data from two
  populations and SNP‑wise \(F_{ST}\) estimates.

  \begin{enumerate}
  \def\labelenumii{\alph{enumii}.}
  \tightlist
  \item
    Compare fixed‑size sliding windows (e.g., 50 SNPs per window) with
    the spline‑based windowing implemented in GenWin. What are the main
    advantages and disadvantages of each approach?\\
  \item
    In a simulation study, you find that spline‑based windows with the
    \(W\) statistic yield a higher true‑positive to false‑positive ratio
    than fixed windows, even though both approaches detect a similar
    number of true selected regions. Explain why this might be the
    case.\\
  \item
    Describe how you would use neutral simulations or permutation tests
    to determine an empirical significance threshold for \(W\) in your
    dataset.
  \end{enumerate}
\item
  \textbf{Missing data and summary statistics:} Consider a
  reduced‑representation dataset (e.g., RADseq) with substantial missing
  data and variable coverage across individuals and populations. You
  wish to estimate windowed \(\pi\) and \(F_{ST}\).

  \begin{enumerate}
  \def\labelenumii{\alph{enumii}.}
  \tightlist
  \item
    Discuss at least two ways in which missing data can bias estimates
    of \(\pi\) and \(F_{ST}\), particularly when missingness differs
    systematically among populations.\\
  \item
    How might you modify your analysis to mitigate these biases (e.g.,
    filtering strategies, using genotype likelihoods, or correcting
    estimators)?\\
  \item
    Reflect on how the missing data issues discussed here relate to
    earlier chapters on variant calling and GATK's reference confidence
    mode. How do choices made at the variant‑calling stage propagate
    into population genomic analyses?
  \end{enumerate}
\item
  \textbf{Ghost populations and demographic inference:} You are
  analyzing a species with a complex history, including suspected gene
  flow from an unsampled (``ghost'') population.

  \begin{enumerate}
  \def\labelenumii{\alph{enumii}.}
  \tightlist
  \item
    Conceptually, explain how unmodeled gene flow from a ghost
    population can bias estimates of \(\pi\), \(F_{ST}\), and Tajima's D
    in the sampled populations.\\
  \item
    Propose one or more strategies (conceptual or methodological) for
    detecting the influence of ghost populations using genomic data.\\
  \item
    How might sliding‑window or spline‑based analyses help you localize
    regions of the genome most affected by ghost admixture?
  \end{enumerate}
\end{enumerate}

\section{Reflection Questions}\label{reflection-questions-10}

\begin{enumerate}
\def\labelenumi{\arabic{enumi}.}
\item
  \textbf{From loci to genomes:} Earlier genetic studies focused on a
  handful of loci, whereas population genomics embraces whole genomes.
  Reflect on how this shift changes the kinds of questions we can ask
  about evolution and how confident we can be in our answers. What new
  challenges arise when moving from a few markers to millions?
\item
  \textbf{Linked selection and interpretation:} Linked selection
  (background selection and selective sweeps) complicates the
  relationship between diversity, recombination, and selection. How does
  this complexity affect our ability to identify specific targets of
  selection? Under what circumstances might you be cautious about
  interpreting low diversity and high \(F_{ST}\) as evidence of local
  adaptation?
\item
  \textbf{Model choice and null expectations:} Many population genomic
  statistics (Tajima's D, \(F_{ST}\), windowed \(\pi\)) are interpreted
  relative to simple null models (e.g., neutrality with constant
  population size). How should we update our ``null expectations'' in
  light of the rich demographic histories and pervasive linked selection
  revealed by genomic data? What does it mean to choose an appropriate
  null model in this context?
\item
  \textbf{Data quality and inference:} Think back to earlier chapters on
  alignment, variant calling, and quality control. In what ways do
  choices made during these early steps constrain or enable the kinds of
  population genomic analyses we can perform? How might you design a
  sequencing and analysis pipeline differently if your primary goal from
  the outset is population genomics rather than variant discovery in a
  single genome?
\item
  \textbf{Ethics and communication:} Population genomics of human
  populations can touch on sensitive issues of ancestry, identity, and
  history. How can scientists responsibly communicate findings about
  human genetic diversity and population structure to avoid reinforcing
  misconceptions or harmful ideologies? What roles do transparency,
  community engagement, and historical context play in this process?
\item
  \textbf{Where to go deeper:} This chapter has provided a conceptual
  overview of population genomics. If you were to take a dedicated
  course in this field, what topics would you be most interested in
  exploring in more depth (e.g., coalescent theory, inference of
  selection, demographic modeling, ancient DNA, graph‑based models of
  population structure)? How might these topics intersect with your own
  research interests or future work?
\end{enumerate}

\section{References}\label{references-11}

\chapter{Sequence Motif Analysis}\label{sequence-motif-analysis}

\section{Overview}\label{overview}

Short paragraph linking to previous chapter and framing why this topic
matters in the course.

\section{Learning objectives}\label{learning-objectives}

\begin{itemize}
\tightlist
\item
  Bullet list of 3--5 concrete things students should be able to do.
\end{itemize}

\section{Narrative}\label{narrative}

\subsection{Key concepts}\label{key-concepts}

\begin{itemize}
\tightlist
\item
  2--4 subsections that walk through your bullet points (e.g., ``FASTQ
  and Phred scores'', ``QC metrics'', ``Design decisions based on QC'').
\end{itemize}

\subsection{Figures and examples}\label{figures-and-examples}

\begin{itemize}
\tightlist
\item
  Simple diagrams you sketch yourself (e.g., read‑quality plot,
  schematic of FASTQ fields, coverage cartoon).
\item
  One or two short, self‑contained R or shell code blocks where
  appropriate.
\end{itemize}

\section{Activities and labs}\label{activities-and-labs}

\begin{itemize}
\tightlist
\item
  Describe how the corresponding lab (e.g., Galaxy vs command line QC)
  connects to the chapter, and give just enough detail that students can
  re‑orient later.
\end{itemize}

\section{Further reading}\label{further-reading}

\begin{itemize}
\tightlist
\item
  List your resources with 1--2‑sentence annotations each, rather than
  reproducing any text.
\end{itemize}

\chapter{Where to Go Next}\label{where-to-go-next}

\section{Overview}\label{overview-1}

Short paragraph linking to previous chapter and framing why this topic
matters in the course.

\section{Learning objectives}\label{learning-objectives-1}

\begin{itemize}
\tightlist
\item
  Bullet list of 3--5 concrete things students should be able to do.
\end{itemize}

\section{Narrative}\label{narrative-1}

\subsection{Key concepts}\label{key-concepts-1}

\begin{itemize}
\tightlist
\item
  2--4 subsections that walk through your bullet points (e.g., ``FASTQ
  and Phred scores'', ``QC metrics'', ``Design decisions based on QC'').
\end{itemize}

\subsection{Figures and examples}\label{figures-and-examples-1}

\begin{itemize}
\tightlist
\item
  Simple diagrams you sketch yourself (e.g., read‑quality plot,
  schematic of FASTQ fields, coverage cartoon).
\item
  One or two short, self‑contained R or shell code blocks where
  appropriate.
\end{itemize}

\section{Activities and labs}\label{activities-and-labs-1}

\begin{itemize}
\tightlist
\item
  Describe how the corresponding lab (e.g., Galaxy vs command line QC)
  connects to the chapter, and give just enough detail that students can
  re‑orient later.
\end{itemize}

\section{Further reading}\label{further-reading-1}

\begin{itemize}
\tightlist
\item
  List your resources with 1--2‑sentence annotations each, rather than
  reproducing any text.
\end{itemize}

\part{Lab Manual}

\chapter{Lab01: Find and describe a Genome Analysis
Paper}\label{lab01-find-and-describe-a-genome-analysis-paper}

\section{Lab 1}\label{sec-lab01}

In class we will read and discuss a recent genome analysis paper that
illustrates the three perspectives on genomics (cell, organism, tree of
life) and the balance between hypothesis‑driven and discovery‑driven
research introduced in Chapter 1\textsuperscript{1,2,10}.

After class, you will search PubMed for a scientific study that includes
some form of genome analysis. Avoid ``double‑dipping'' by reusing papers
assigned for other courses or lab meetings. Ideally, choose:

\begin{itemize}
\tightlist
\item
  a genome paper on your study organism, or\\
\item
  a paper that follows a workflow similar to one you plan to use in your
  own research, or
\item
  a recent paper that excited you and presents a genome analysis.
\end{itemize}

Before you start, skim the ``Types of Genome Analysis Research
Projects'' section in Chapter 1 and, if needed, review Appendix
Section~\ref{sec-papers} for a refresher on reading primary
literature\textsuperscript{1,2,10,24}.

By the end of this lab, you should be able to:

\begin{itemize}
\tightlist
\item
  Identify the main hypothesis (or questions) of a genomics paper.
\item
  Classify the study using the three perspectives from Chapter 1 (cell,
  organism, tree of life).
\item
  Describe whether the study leans more hypothesis‑driven,
  discovery‑driven, or combines both.
\item
  Summarize key genome features (assembly size, chromosome number, basic
  QC metrics) reported in the paper.
\end{itemize}

\section{Deliverables}\label{deliverables}

\begin{enumerate}
\def\labelenumi{\arabic{enumi}.}
\tightlist
\item
  \textbf{Paper and methods diagram}

  \begin{itemize}
  \tightlist
  \item
    Upload the PDF of your chosen paper to Canvas.\\
  \item
    Upload a one‑page diagram of the methods used to test the main
    hypothesis. Give your diagram a descriptive title and include:

    \begin{itemize}
    \tightlist
    \item
      data types (e.g., short‑read WGS, long‑read assembly, RNA‑seq),\\
    \item
      main analysis steps (alignment, assembly, variant calling,
      comparative analyses),\\
    \item
      which of the three perspectives your paper emphasizes
    \item
      whether it is hypothesis-drive, discovery-based, or both.
    \end{itemize}
  \end{itemize}
\item
  \textbf{2--3 minute video summary}

  \begin{itemize}
  \tightlist
  \item
    Record an informal 2--3 minute Panopto video summarizing the paper's
    goals, methods, and main findings, and how it fits into the
    perspectives from Chapter 1.
  \item
    Upload the video to the ``Lab Activity \#1'' Panopto assignments
    folder as described on Canvas.
  \end{itemize}
\end{enumerate}

Please refer to the syllabus for late policies. Detailed step‑by‑step
instructions for recording and uploading Panopto videos are provided on
Canvas rather than here.

\section{Resources and Refreshers:}\label{resources-and-refreshers}

If you do not have much experience finding and reading primary
literature, complete Appendix Section~\ref{sec-papers} (``How to Read
Scientific Literature'') before or alongside this lab. That appendix
walks through hypothesis identification, concept mapping, and designing
follow‑up experiments using a concrete example.

\chapter{Lab02: Design a Genome
Project}\label{lab02-design-a-genome-project}

\section{Lab 2}\label{sec-lab02}

In Chapter 1, we introduced three broad perspectives on bioinformatics
and genomics (cell, organism, tree of life)\textsuperscript{10}. Later
in the textbook, Pevsner expands these 3 perspectives into five basic
perspectives on genome sequencing.

For this lab, you will use these 5 perspectives, along with your new
understanding of sequencing technologies from Chapter 2, to sketch a
1--2 page narrative for a hypothetical genome
project\textsuperscript{10,25,26}.

The five perspectives on genomics discussed in Pevsner Chapter 15 can be
summarized as:

\begin{itemize}
\tightlist
\item
  \textbf{Catalog genomic information}: basic genome features (size,
  chromosomes, GC content, repeats, gene counts), based on sequencing,
  assembly, and annotation.
\item
  \textbf{Catalog comparative genomic information}: whole‑genome
  comparisons to related species, orthologs, divergence times, lateral
  gene transfer, using whole‑genome alignments and genome browsers.
\item
  \textbf{Biological principles}: how genome structure and variation
  underlie development, metabolism, behavior, and evolutionary processes
  (e.g., genome size evolution, polyploidization, gene birth/death).
\item
  \textbf{Human disease relevance}: how genomes relate to disease in
  humans or plants, including SNPs, linkage and association studies, and
  host--pathogen interactions.
\item
  \textbf{Bioinformatics aspects}: databases, software, and
  visualization tools that make genome analysis possible.
\end{itemize}

These map naturally onto the three perspectives in Chapter 1: cell‑scale
questions, organism‑level questions, and tree‑of‑life/comparative
questions\textsuperscript{10}.

\subsection{The Assignment}\label{the-assignment}

Write a 1--2 page narrative of the genome analysis project you have
designed. To guide you, begin by reflecting on the following questions:

\begin{itemize}
\tightlist
\item
  \emph{If you could sequence the genomes of 100 individuals from any
  species, which species would you choose?}
\item
  \emph{What hypotheses would you test, how would you perform data
  analyses, and what resources would you require in terms of hardware,
  software, and collaborators?}
\item
  \emph{What ethical issues might arise in sequencing these genomes?}
\end{itemize}

Specifically, think about:

\begin{itemize}
\tightlist
\item
  Which sequencing platforms and coverage would you choose (Chapter 2)?
\item
  Which file formats and QC steps would your workflow rely on (Chapter
  3)?\textsuperscript{10,25,26}
\end{itemize}

For your design narrative:

\begin{itemize}
\tightlist
\item
  Choose a species and justify why 100 genomes from this species would
  be informative. - State one or two main hypotheses, or guiding
  questions if discovery-based.
\item
  Identify which of the five perspectives (and which of the three
  Chapter 1 perspectives) your project emphasizes.
\item
  Outline your sequencing strategy (platforms, approximate coverage,
  sample types) using concepts from Chapter 2.
\item
  Sketch the main analysis steps (from raw FASTQ through
  alignment/assembly and downstream analyses) using ideas from Chapters
  1 and 3.
\item
  Briefly note any major ethical issues if humans or other sensitive
  species are involved.
\end{itemize}

The focus of your design narrative should be on \textbf{motivation and
approach}. You may refer to assigned readings
(e.g.,\textsuperscript{1};\textsuperscript{2};\textsuperscript{3}) in
your narrative, but you do not need a formal reference list.

\chapter{Lab03: Genome Browsers and
Databases}\label{lab03-genome-browsers-and-databases}

\section{Using the Jump-2-gene (J2G) feature on
FlyBase}\label{using-the-jump-2-gene-j2g-feature-on-flybase}

\begin{center}
\includegraphics[width=0.5\linewidth,height=\textheight,keepaspectratio]{images/Lab3_FlyBase.png}
\end{center}

\begin{enumerate}
\def\labelenumi{\arabic{enumi}.}
\tightlist
\item
  Go to the Drosophila Organismal Genomic Database
  \href{https://flybase.org/}{``FlyBase''} and search for the
  \textbf{DSX} gene using the search box in the top-right corner.
\end{enumerate}

\begin{center}

\includegraphics[width=0.75\linewidth,height=\textheight,keepaspectratio]{images/Lab3_J2G.png}
\end{center}

\begin{enumerate}
\def\labelenumi{\arabic{enumi}.}
\setcounter{enumi}{1}
\item
  Which chromosome is this gene on? \emph{Hint: In the right hand panel,
  you can jump between sections. Here you want to find the ``Genomic
  Location''.}
\item
  What do you think the ``R'' in the chromosome name means?
\item
  In the section \textbf{Genomic Location} use the
  
\includegraphics[width=0.1\linewidth,height=\textheight,keepaspectratio]{images/Lab3_JBrowse.png}
  button to view ``DSX'' in the genome browser.
\item
  On the left, turn on the \textbf{RNA} and \textbf{CDS} tracks:
\end{enumerate}

\begin{center}
\includegraphics[width=1\linewidth,height=\textheight,keepaspectratio]{images/Lab3_browser.jpeg}
\end{center}

\begin{enumerate}
\def\labelenumi{\arabic{enumi}.}
\setcounter{enumi}{5}
\item
  Once you configure the browser, answer the following questions:

  \begin{itemize}
  \tightlist
  \item
    How many different transcripts does this gene have?
  \item
    How many exons are in each protein?
  \item
    What strand of DNA is it on? (Hint: Right click on features to get
    details.)
  \end{itemize}
\item
  Click on one of the protein isoforms and identify which protein
  domains are present. \emph{Note: There are several options, so just
  pick one. The goal of this assignment is to \textbf{explore}.}
\item
  Scroll down to \textbf{Transcript Level Features} and turn on these
  tracks, clicking the \texttt{?} for each to read about the evidence:

  \begin{itemize}
  \tightlist
  \item
    TSS (modENCODE, embryo)
  \item
    RNA-Seq exon junctions
  \item
    What evidence has been used to confirm the different genetic
    features of this gene?
  \end{itemize}
\item
  View these evidence tracks in the genome browser and hover over
  features to see read counts, then answer:

  \begin{itemize}
  \tightlist
  \item
    How many genomic ``reads'' support each prediction?
  \item
    TSS (modENCODE, embryo)
  \item
    RNA-Seq exon junctions
  \end{itemize}
\item
  In the \textbf{Expression} section under \textbf{RNA-Seq → Tissues},
  turn on expression for \textbf{Gonads and male accessory glands}.
\end{enumerate}

\begin{center}
\includegraphics[width=0.75\linewidth,height=\textheight,keepaspectratio]{images/Lab3_ExprTrack.png}
\end{center}

\begin{enumerate}
\def\labelenumi{\arabic{enumi}.}
\setcounter{enumi}{10}
\tightlist
\item
  Configure and maximize vertical spacing as in the image above, then
  examine the expression data. \textbf{What predictions do you make
  about this gene?}
\end{enumerate}

\section{2. Compare orthologous sequences using two-way BLAST at
NCBI}\label{compare-orthologous-sequences-using-two-way-blast-at-ncbi}

\begin{enumerate}
\def\labelenumi{\arabic{enumi}.}
\tightlist
\item
  In the \href{https://www.ncbi.nlm.nih.gov/}{NCBI} \textbf{Gene}
  database, find the \textbf{ADH} gene in \emph{Drosophila grimshawi}.
  \emph{Hint: On the main NCBI page, use the dropdown menu to select the
  ``Gene'' database. Then search for the specific gene and species.}
\end{enumerate}

\begin{center}
\includegraphics[width=0.4\linewidth,height=\textheight,keepaspectratio]{images/Lab3_Dgrim.jpg}
\end{center}

\begin{enumerate}
\def\labelenumi{\arabic{enumi}.}
\setcounter{enumi}{1}
\item
  From the gene page, what is the Gene Symbol? What is the gene ID?
\item
  From the gene page, scroll down to the \textbf{Genome Data Viewer}
  button, and click it. From the download button, hover over ``Download
  FASTA'' and click the ``FASTA (visible range)'' option. This will
  download the ADH gene region for \emph{D. grimshawi}.
\end{enumerate}

\begin{center}
\includegraphics[width=0.5\linewidth,height=\textheight,keepaspectratio]{images/Lab3_GDV.png}
\end{center}

\begin{enumerate}
\def\labelenumi{\arabic{enumi}.}
\setcounter{enumi}{3}
\item
  Repeat the previous steps to download the gene sequence for ADH for
  \emph{D. melanogaster}. For this species, you will notice there are
  multiple results from this gene. Select the one with the gene symbol
  ``Adh''.
\item
  Go to
  \textbf{\href{https://blast.ncbi.nlm.nih.gov/Blast.cgi?BLAST_SPEC=blast2seq&LINK_LOC=align2seq&PAGE_TYPE=BlastSearch}{blast2seq}}
  and:

  \begin{itemize}
  \tightlist
  \item
    Paste \emph{D. melanogaster} ADH into the \textbf{Query}.
  \item
    Paste \emph{D. grimshawi} ADH into the \textbf{Subject}.
  \item
    Under \textbf{Program Selection}, choose \textbf{blastn} for
    somewhat similar sequences.
  \item
    Click the
    
\includegraphics[width=0.15\linewidth,height=\textheight,keepaspectratio]{images/Lab3_Blast.png}
    button.
  \end{itemize}
\item
  From the BLAST Results Page, record the following information:

  \begin{longtable}[]{@{}ll@{}}
  \toprule\noalign{}
  Blast result features & Recorded Results \\
  \midrule\noalign{}
  \endhead
  \bottomrule\noalign{}
  \endlastfoot
  Query length & \\
  Subject length & \\
  Percent identity between the two species & \\
  Percent query cover & \\
  E value & \\
  Total score & \\
  \end{longtable}
\item
  Under the \textbf{Alignments} tab, find the alignment with the highest
  percent identity and record:

  \begin{longtable}[]{@{}ll@{}}
  \toprule\noalign{}
  Alignment information & Recorded Results \\
  \midrule\noalign{}
  \endhead
  \bottomrule\noalign{}
  \endlastfoot
  Number of matches (Identities) & \\
  Number of Gaps & \\
  Orientation (Strand) & \\
  \end{longtable}
\item
  Now select the \textbf{Dot Plot} tab and describe:

  \begin{itemize}
  \tightlist
  \item
    How do your calculations above compare to the visual of hits in the
    dot plot?
  \end{itemize}
\item
  In a new tab, return to FlyBase. At the top, under \textbf{Tools},
  hover over ``Search/Browse Portals'', and select \textbf{Sequenced
  Species}. Explore the page and answer the following questions:

  \begin{itemize}
  \tightlist
  \item
    Based on the tree, when did \emph{D. melanogaster} and \emph{D.
    grimshawi} last share a common ancestor?
  \item
    Based on your inferred evolutionary divergence time between \emph{D.
    melanogaster} and \emph{D. grimshawi}, does the observed percent
    identity from your previous BLAST search make sense?
  \end{itemize}
\item
  Predict whether you expect a higher or lower percent difference
  between \emph{D. pseudoobscura} and \emph{D. melanogaster}.

  \begin{itemize}
  \tightlist
  \item
    If time permits, test this by repeating the two-way BLAST with ADH
    from \emph{D. pseudoobscura} as the subject and \emph{D.
    melanogaster} as the query.
  \end{itemize}
\end{enumerate}

\section{3. Compare synteny among Muller
Elements}\label{compare-synteny-among-muller-elements}

On FlyBase, under the \textbf{Tools} tab, hover over \textbf{Genomics
Tools} and select \textbf{Synteny Table}. \emph{Drosophila} chromosomes
are labeled as \textbf{Muller's Elements}, and comparative genomics
across \emph{Drosophila} species has been studied since the 1930s using
polytene chromosome squashes.

Answer the following:

\begin{enumerate}
\def\labelenumi{\arabic{enumi}.}
\item
  What are the major changes at the chromosome level across these
  \emph{Drosophila} species? Specifically, do any species have pieces of
  multiple Muller elements spanning a single chromosome?
\item
  Compare the below (Figure 3 from\textsuperscript{209}) to the synteny
  table and mark where Muller elements have shuffled to form a combined
  element.
\end{enumerate}

\begin{center}
\pandocbounded{\includegraphics[keepaspectratio]{images/Lab3_synteny.jpeg}}
\end{center}

\begin{enumerate}
\def\labelenumi{\arabic{enumi}.}
\setcounter{enumi}{2}
\tightlist
\item
  What other kinds of changes can you find?
\end{enumerate}

\section{4. Whole genome alignments on
UCSC}\label{whole-genome-alignments-on-ucsc}

\begin{enumerate}
\def\labelenumi{\arabic{enumi}.}
\item
  Go to the
  \textbf{\href{https://genome.ucsc.edu/cgi-bin/hgGateway}{UCSC Genome
  Browser}} and select the \emph{D. melanogaster} genome. \emph{At the
  top left, in the ``Popular Species'' region, click the button that
  says ``Fruitfly''.}
\item
  After selecting the ``dm6'' assembly, in the ``Position/Search Term''
  field, search for the \textbf{rdl} gene and select the transcript
  variant A , then zoom out 3×.
\item
  Under \textbf{Comparative Genomics}, click \textbf{Conservation} and:

  \begin{itemize}
  \item
    Turn on \textbf{MultiZ Alignment} for these three \emph{Drosophila}
    species:

    \begin{itemize}
    \tightlist
    \item
      \emph{D. simulans}
    \item
      \emph{D. pseudoobscura}
    \item
      \emph{D. grimshawi}
    \end{itemize}
  \item
    Under \textbf{Others}, select all available options.
  \item
    Under \textbf{Subtracks}, select all and display them as
    \textbf{dense}.
  \item
    Set the MultiZ track to \textbf{full}.
  \item
    Click \textbf{Submit} to return to the genome viewer.
  \end{itemize}
\item
  Do the different conservation tracks (\textbf{phyloP},
  \textbf{phastCons}, \textbf{Cons Elements}, and \textbf{MultiZ Align})
  agree with each other?
\item
  Underneath the Cons Elements track, examine the individual alignments
  for \emph{D. melanogaster} and the other Drosphila species previously
  selected.

  \begin{itemize}
  \tightlist
  \item
    Do these tracks confirm what you saw in Part 2 with the pairwise
    alignment between \emph{D. grimshawi} and \emph{D. pseudoobscura}
    based on the species relationships on FlyBase?
  \item
    Would you guess \emph{D. simulans} is more closely related to
    \emph{D. melanogaster} than \emph{D. pseudoobscura} or more
    distantly related? Return to the species tree in FlyBase to confirm.
  \end{itemize}
\item
  Repeat the setup for the \textbf{DSX} gene (you can type DSX in the
  top search panel and click \textbf{go}; your configuration should
  persist) and answer:

  \begin{itemize}
  \tightlist
  \item
    Do the different conservation tracks still agree for this gene?
  \item
    Why might different genes show different degrees of conservation?
  \end{itemize}
\end{enumerate}

\section{5. Orthologs, paralogues, and gene trees on
Ensembl}\label{orthologs-paralogues-and-gene-trees-on-ensembl}

\begin{enumerate}
\def\labelenumi{\arabic{enumi}.}
\item
  Go to the \textbf{Ensembl} genome browser and select the
  \textbf{Human} genome from the species drop-down menu.
\item
  In the search bar, type \textbf{PRDM9} and click \textbf{Go}.
\end{enumerate}

\begin{center}
\pandocbounded{
\includegraphics[keepaspectratio]{images/Lab3_Ensembl.jpeg}}
\end{center}

\begin{enumerate}
\def\labelenumi{\arabic{enumi}.}
\setcounter{enumi}{2}
\item
  From the gene page, select \textbf{Orthologues} under the
  \textbf{Comparative Genomics} section in the left-hand panel.
\item
  Answer:

  \begin{itemize}
  \tightlist
  \item
    How many orthologues are predicted for this gene in
    \textbf{Primates}?
  \item
    How many are predicted for \textbf{Sauropsida}? Which species?
  \item
    Note the \textbf{Target \%id} and \textbf{Query \%id} for these
    orthologues.
  \end{itemize}
\item
  Compare the counts of 1:1 orthologues versus 1:many orthologues.

  \begin{itemize}
  \tightlist
  \item
    What do you think the ``many'' indicate biologically?
  \end{itemize}
\item
  Set \textbf{Species set} to \textbf{All}, then in the left-hand panel
  choose \textbf{Gene tree}.

  \begin{itemize}
  \tightlist
  \item
    Select \textbf{Tetrapods} on the gene tree, then click \textbf{View
    sub-tree} using the \textbf{Wasabi} viewer.
  \item
    What do you notice about these gene sequences?
  \item
    Is this a protein sequence alignment or a DNA sequence alignment?
  \end{itemize}
\end{enumerate}

\begin{center}
\pandocbounded{
\includegraphics[keepaspectratio]{images/Lab3_export.jpeg}}
\end{center}

\begin{enumerate}
\def\labelenumi{\arabic{enumi}.}
\setcounter{enumi}{6}
\tightlist
\item
  Along the top, click \textbf{Export} and:

  \begin{itemize}
  \tightlist
  \item
    Make sure \textbf{ALL} species are included.
  \item
    Download the tree as a \textbf{Newick} file.
  \end{itemize}
\end{enumerate}

\begin{center}
\pandocbounded{\includegraphics[keepaspectratio]{images/Lab3_Newick.jpeg}}
\end{center}

\begin{enumerate}
\def\labelenumi{\arabic{enumi}.}
\setcounter{enumi}{7}
\item
  (Optional) Open the Newick file in \textbf{FigTree}:

  \begin{itemize}
  \tightlist
  \item
    Explore settings to configure your tree (e.g., change tip names via
    the annotate feature, collapse groups).
  \item
    When satisfied, export the tree as a PDF.
  \end{itemize}
\end{enumerate}

\section{Reflection Questions}\label{reflection-questions-11}

\begin{enumerate}
\def\labelenumi{\arabic{enumi}.}
\item
  What challenges did you encounter while navigating the different
  genome databases, and how did you overcome them?
\item
  How did the different genome browsers compare in terms of ease of use
  and the type of information available?
\item
  What new insights did you gain about genome structure and evolution by
  using these tools?
\item
  When exploring the gene annotation tracks, how does the evidence from
  each track support your existing knowledge of how the gene functions?
\item
  How does the percent identity in your BLAST search reflect
  evolutionary relationships between species, and how did this compare
  to the conservation tracks in UCSC?
\item
  How do synteny comparisons help infer evolutionary history?
\item
  Were you surprised by the number of orthologs and paralogs across
  different vertebrate groups for PRDM9? Do you think the accuracy of
  these counts differs across clades?
\end{enumerate}

\section{References}\label{references-12}

\chapter{Lab04: Genome Sequence Indexing and
Alignment}\label{lab04-genome-sequence-indexing-and-alignment}

For this lab, you will work with the data and scripts in the
``Sample\_NGS\_Workflow'' shared folder to index a reference genome,
split the fastq files into multiple lanes, and do an initial genome
alignment. You will end by collecting a few QC metrics to assess the
quality of the input fastq files and the alignment file. This lab
activity is designed to give you real world experience with the tools we
have been discussing in class so far. These are some of the first steps
in the genome analysis workflow. We will work through later steps
throughout the semester.

\begin{enumerate}
\def\labelenumi{\arabic{enumi}.}
\item
  Login to your ASC account and make a folder for Lab 4 called
  \texttt{Lab4}.
\item
  Copy the script files listed below from the shared folder
  \texttt{Sample\_NGS\_Workflow/scripts} into your \texttt{Lab4} folder.
\end{enumerate}

\begin{itemize}
\tightlist
\item
  1\_index\_sacCer3.sh\\
\item
  2\_sample\_sra.sh\\
\item
  3\_sample\_trimming.sh\\
\item
  4\_BWA\_example.sh
\end{itemize}

\begin{enumerate}
\def\labelenumi{\arabic{enumi}.}
\setcounter{enumi}{2}
\tightlist
\item
  Look over each script file and try to understand which tools they are
  using and what the scripts are doing. Draw out a diagram of the
  workflow from script to script.
\end{enumerate}

\section{Step 1 : Index your genome}\label{step-1-index-your-genome}

\begin{enumerate}
\def\labelenumi{\arabic{enumi}.}
\tightlist
\item
  For the 1st script \texttt{1\_index\_sacCer3.sh}, edit line 8 to
  change the variable \texttt{path} to point to your Lab4 directory and
  submit it to the queue using the following parameters:
\end{enumerate}

\begin{itemize}
\tightlist
\item
  Number of CPUs = 1
\item
  RAM/Memory = 2GB
\item
  Queue = class
\item
  Time = 01:00:00 (1 hour)
\end{itemize}

\begin{enumerate}
\def\labelenumi{\arabic{enumi}.}
\setcounter{enumi}{1}
\item
  Inspect the job standard error file. It will follow the naming
  convention of your job name followed by a period and the letter ``o''
  and the job number (e.g.~\texttt{1indexsacCer3shSCR.o8618}). Any
  errors with your job will always be in this file, so get in the habit
  of checking here FIRST after your job finishes! This includes both the
  standard out and standard error of your job.
\item
  This script will index the reference genome fasta file. First, it will
  use \href{https://bio-bwa.sourceforge.net/bwa.shtml}{bwa index}, then
  samtools. Each index is for a different part of the pipeline. Look at
  the files that have been generated.
\item
  Read over the
  \href{https://www.htslib.org/doc/samtools-faidx.html}{samtools faidx}
  documentation to understand this file: ``sacCer3.masked.fa.fai''. Then
  explore the file using \texttt{less}.
\item
  You will notice that it lists the chromosomes and their lengths in the
  1st and 2nd columns. How can you use command line tools to calculate
  the total genome size from these files? \emph{Remember, this
  information is useful in calculating genome coverage!}
\end{enumerate}

\section{Step 2 : Download reads from
NCBI}\label{step-2-download-reads-from-ncbi}

\begin{enumerate}
\def\labelenumi{\arabic{enumi}.}
\item
  Inspect the next script you copied over ``2\_sample\_sra.sh''. As you
  did with the first script, edit line 10 to change the path variable so
  that it runs in your Lab4 directory. Before submitting it to the
  queue, we will need to configure your account for
  \href{https://hpcdocs.asc.edu/content/sra-sequence-read-archive}{SRAtoolkit}:

\begin{verbatim}
module purge
module load sra
perl /apps/x86-64/apps/sra_2.3.5/sratoolkit.2.3.5-centos_linux64/bin/configuration-assistant.perl
Y
Y
\end{verbatim}

  \begin{quote}
  This should end with the message ``Configuration is correct''. This
  creates a subdirectory named .ncbi The user-settings.mkfg file in that
  directory is slightly different from the one made by the graphical
  configuration tool in that it does not define which directory to
  download files into.
  \end{quote}
\item
  Once you have your account configured, then you are ready to submit
  this script to the queue using the following parameters:
\end{enumerate}

\begin{itemize}
\tightlist
\item
  Number of CPUs = 1
\item
  RAM/Memory = 1GB
\item
  Queue = class
\item
  Time = 01:00:00 (1 hour)
\end{itemize}

\begin{enumerate}
\def\labelenumi{\arabic{enumi}.}
\setcounter{enumi}{2}
\item
  Inspect the job output and make any tweaks as needed to the script to
  get it to work if it did not work. A good troubleshooting tip is to
  find the FIRST error in your code, at the top of the standard out.
  Once you fix that, re-run it and work your way down until you have no
  errors. A lot of times, a fatal early error can cause a lot of
  downstream errors, so it is not productive to start at the bottom -
  always work from the top-down!
\item
  What new files have appeared? Use zcat and less to examine these
  files. How are they different and what information was used to split
  them? \emph{Hint: refer back to the Fastq section of Chapter 3 for
  more detail on this file format!}
\item
  How many reads are in the FASTQ file? \emph{Hint: refer back to FASTQ
  section of Chapter 3 for details of how many lines are there per entry
  in this type of file. Then check the job output file to see the
  results of the line count in the script}
\end{enumerate}

\section{Step 3 : Trimming}\label{step-3-trimming}

\begin{enumerate}
\def\labelenumi{\arabic{enumi}.}
\item
  Inspect the next script you copied over ``3\_sample\_trimming.sh''.
  Look over the
  \href{http://www.usadellab.org/cms/uploads/supplementary/Trimmomatic/TrimmomaticManual_V0.32.pdf}{trimmomatic
  tool} documentation.
\item
  As you did with the first script, make any necessary edits for it to
  run in your Lab4 directory. You will need to edit the username on line
  14 as well. Then, submit it to the queue using the following
  parameters:
\end{enumerate}

\begin{itemize}
\tightlist
\item
  Number of CPUs = 2
\item
  RAM/Memory = 1GB
\item
  Queue = class
\item
  Time = 01:00:00 (1 hour)
\end{itemize}

\begin{enumerate}
\def\labelenumi{\arabic{enumi}.}
\setcounter{enumi}{2}
\tightlist
\item
  Inspect the job output. Do you see any errors? If so, work to fix them
  from the top-down as you did above!
\end{enumerate}

\section{Step 4 : Alignment}\label{step-4-alignment}

\begin{enumerate}
\def\labelenumi{\arabic{enumi}.}
\item
  Inspect the next script you copied over ``4\_BWA\_example.sh''. This
  script uses the tool
  \href{https://bio-bwa.sourceforge.net/bwa.shtml}{bwa mem} to align the
  raw reads to the reference genome.
\item
  As you did with the first few scripts, make any necessary edits for it
  to run in your Lab4 directory and submit it to the queue using the
  following parameters:
\end{enumerate}

\begin{itemize}
\tightlist
\item
  Number of CPUs = 4
\item
  RAM/Memory = 2GB
\item
  Queue = class
\item
  Time = 01:00:00 (1 hour)
\end{itemize}

\begin{enumerate}
\def\labelenumi{\arabic{enumi}.}
\setcounter{enumi}{2}
\tightlist
\item
  Inspect the job output. This script does a LOT, so look it over
  carefully to make sure you understand what the output means. Feel free
  to refer back to the initial script - notice where the echo statements
  are to help find your place in the output. Look up any of the commands
  or tools to make sure you understand the output.
\end{enumerate}

\section{Step 5 : Quality assessment}\label{step-5-quality-assessment}

\begin{enumerate}
\def\labelenumi{\arabic{enumi}.}
\tightlist
\item
  At the bottom of the scripts ``3\_sample\_trimming.sh'' and
  ``4\_BWA\_example.sh'', you will see that they each collect some QC
  metrics - specifically, the first one runs
  \href{https://www.bioinformatics.babraham.ac.uk/projects/fastqc/}{fastQC}
  on both the input fastq file and the trimmed one. Download the
  resulting zip files locally so you can inspect the results. For a
  refresher on uploading and downloading files from the ASC, see this
  \href{https://auburn.instructure.com/courses/1709268/pages/getting-on-the-alabama-super-computer-asc-2}{Canvas
  Page}.
\end{enumerate}

\begin{itemize}
\tightlist
\item
  Was trimming necessary?
\item
  Does the data look better after trimming?
\end{itemize}

\begin{enumerate}
\def\labelenumi{\arabic{enumi}.}
\setcounter{enumi}{1}
\item
  Additionally, the latter script collects coverage using
  \texttt{samtools\ depth} into an output file
  \texttt{SRR1693691.coverage\_summary.txt}. Use \texttt{less} to
  examine the file these statistics have printed to.
\item
  Now, look at the
  \href{https://www.htslib.org/doc/samtools-depth.html}{samtools depth}
  documentation for this program ``depth'' and parse the \texttt{awk}
  command that it uses to generate the summary you just inspected.
  Notice the use of the variable \texttt{\$genome\_size}. This uses to
  the file you made earlier in script \#1, ``sacCer3.masked.fa.fai'', to
  programmatically calculate the genome size and store it in a variable
  for this calculation.
\item
  Another way of gathering QC metrics is to use the tool
  \href{https://www.htslib.org/doc/samtools-flagstat.html}{samtools
  flagstat}. See Section \textbf{Verification of Sam/Bam files using
  Samtools} in Chapter 3 for a sample output of this program. Write a
  script to collect this additional information from your merged sorted
  bam file. \emph{Note: When adapting this main script for your
  research, you can add this code to the main script so it is all done
  in one step!}
\item
  From the \texttt{flagstat} output, locate the section that reports the
  percentage of reads that mapped to the reference genome. Include this
  in your summary.
\end{enumerate}

Congratulations! You have now successfully indexed a reference genome
and mapped reads to it! You have also collected various QC metrics of
this process. Based on how you have seen these types of methods written
in papers we have read, \emph{write a short methods summary of what you
did, including the quality metrics from the raw reads and the resulting
coverage and \% mapped calculations you made. (no more than 250 words!)}

\chapter{Lab05: BLAST on the Command
Line}\label{lab05-blast-on-the-command-line}

For this lab, you will learn to use BLAST (Basic Local Alignment Search
Tool) on the command line to identify sequence homology and annotate
unknown sequences. While web-based BLAST is convenient for single
queries, command-line BLAST is essential for analyzing large datasets,
automating repetitive tasks, and integrating BLAST into bioinformatics
pipelines. This lab will walk you through creating BLAST databases,
running searches, and parsing results---skills that are fundamental to
genome annotation, comparative genomics, and evolutionary studies.

\section{Background: Why Command Line
BLAST?}\label{background-why-command-line-blast}

BLAST is the most widely used bioinformatics tool, with the original
publication cited nearly 80,000 times\textsuperscript{63}. While the web
interface at NCBI is user-friendly, command-line
BLAST\textsuperscript{64} offers several advantages:

\begin{itemize}
\tightlist
\item
  \textbf{Batch processing}: Query thousands of sequences automatically
\item
  \textbf{Custom databases}: Search against your own sequences, not just
  NCBI's databases
\item
  \textbf{Reproducibility}: Scripts can be shared and rerun exactly
\item
  \textbf{Integration}: Combine BLAST with other tools in analysis
  pipelines
\item
  \textbf{Control}: Fine-tune parameters and output formats for
  downstream analysis
\end{itemize}

In this lab, you'll work with a \textbf{realistic scenario}: you've
assembled transcripts from a newly sequenced plant genome and need to
annotate them using the well-characterized \emph{Arabidopsis thaliana}
genome. Your goal is to determine which transcripts correspond to
mitochondrial, chloroplast, or nuclear genes.

\section{Setup}\label{setup}

\begin{enumerate}
\def\labelenumi{\arabic{enumi}.}
\item
  Login to your ASC account and make a folder for Lab 5 called
  \texttt{Lab5}.
\item
  Copy all files from the shared folder
  \texttt{\textasciitilde{}/bioinf\_shared/Lab5\_BLAST\_CL} into your
  \texttt{Lab5} directory. \textbf{Important}: Copy these files, do not
  move them!
\item
  List the files in your directory. You should see:

  \begin{itemize}
  \tightlist
  \item
    \texttt{test.fasta} (your query sequences - the unknown transcripts)
  \item
    \texttt{ATcp.fasta} (\emph{Arabidopsis} chloroplast genome)
  \item
    \texttt{ATmt.fasta} (\emph{Arabidopsis} mitochondrial genome)
  \item
    \texttt{ATchrV.fasta} (\emph{Arabidopsis} chromosome V - a sample
    nuclear chromosome)
  \item
    \texttt{Example\_script.sh} (template for your commands)
  \end{itemize}
\item
  Open \texttt{Example\_script.sh} in your preferred text editor. This
  script will be your working file throughout the lab. You'll add
  commands to it as you progress, building a complete, reusable BLAST
  workflow by the end.
\item
  Replace your username at the top:

\begin{Shaded}
\begin{Highlighting}[]
\VariableTok{DIR}\OperatorTok{=}\StringTok{"/home/aubclsfXXXX/Lab5"} \CommentTok{\# Replace with your username}
\end{Highlighting}
\end{Shaded}
\end{enumerate}

\section{Step 1: Examine Your Query
Sequences}\label{step-1-examine-your-query-sequences}

Before running BLAST, you should understand your input data.

\begin{enumerate}
\def\labelenumi{\arabic{enumi}.}
\item
  Use \texttt{less} to view the \texttt{test.fasta} file. Recall that
  FASTA format consists of a header line starting with
  \texttt{\textgreater{}} followed by the sequence on subsequent lines.
\item
  Count how many sequences are in \texttt{test.fasta}.

  \textbf{Hint}: The \texttt{\textgreater{}} symbol marks each new
  sequence header. Use \texttt{grep} with the \texttt{-c} flag to count
  matching lines. Check \texttt{man\ grep} if needed.
\item
  Record this number---you'll need it later to determine how many
  sequences had no BLAST hits.
\end{enumerate}

\section{Step 2: Create BLAST
Databases}\label{step-2-create-blast-databases}

When you use BLAST on the NCBI website, you're searching against
pre-built databases. To use BLAST on the command line with custom
sequences, you must first build your own databases.

\begin{enumerate}
\def\labelenumi{\arabic{enumi}.}
\item
  Load the BLAST module to access the BLAST+ tools:

\begin{Shaded}
\begin{Highlighting}[]
\ExtensionTok{module}\NormalTok{ load blast+/2.15.0}
\end{Highlighting}
\end{Shaded}
\item
  Examine the database creation tool:

\begin{Shaded}
\begin{Highlighting}[]
\ExtensionTok{makeblastdb} \AttributeTok{{-}help}
\end{Highlighting}
\end{Shaded}

  Look through the documentation to identify the required parameters.
  Key options include:

  \begin{itemize}
  \tightlist
  \item
    \texttt{-in}: Input FASTA file
  \item
    \texttt{-dbtype}: Type of sequences (\texttt{nucl} for nucleotides,
    \texttt{prot} for proteins)
  \item
    \texttt{-out}: Output database name (optional; defaults to input
    filename)
  \end{itemize}

  \begin{quote}
  Documentation on this module can be found on the
  \href{https://hpcdocs.asc.edu/content/blast}{ASC Documentation Site}.
  Additional documentation on this tool can be found
  \href{https://www.ncbi.nlm.nih.gov/books/NBK52640/}{at NCBI}.
  \end{quote}
\item
  Determine the sequence type in your FASTA files. Open one of the
  \texttt{AT*.fasta} files with \texttt{less}. Are these nucleotide or
  protein sequences?
\item
  Create a BLAST database for each of the three \emph{Arabidopsis}
  sequence files (\texttt{ATcp.fasta}, \texttt{ATmt.fasta}, and
  \texttt{ATchrV.fasta}). Do \textbf{not} create a database from
  \texttt{test.fasta}---this will be your query.

  For example, for the chloroplast:

\begin{Shaded}
\begin{Highlighting}[]
\ExtensionTok{makeblastdb} \AttributeTok{{-}in}\NormalTok{ ATcp.fasta }\AttributeTok{{-}dbtype}\NormalTok{ nucl}
\end{Highlighting}
\end{Shaded}

  \textbf{Tip}: To make your script more flexible, use variables as
  shown in the example script. For instance:

\begin{Shaded}
\begin{Highlighting}[]
\VariableTok{cp}\OperatorTok{=}\NormalTok{ATcp.fasta}
\ExtensionTok{makeblastdb} \AttributeTok{{-}in} \VariableTok{$cp} \AttributeTok{{-}dbtype}\NormalTok{ nucl}
\end{Highlighting}
\end{Shaded}
\item
  Add these three \texttt{makeblastdb} commands to your
  \texttt{Example\_script.sh} file.
\item
  Submit your script to the ASC queue with these parameters:

  \begin{itemize}
  \tightlist
  \item
    Number of CPUs = 1
  \item
    RAM/Memory = 50MB
  \item
    Queue = class
  \item
    Time = 01:00:00 (1 hour)
  \end{itemize}
\item
  After the job completes, check the standard output file (ending in
  \texttt{.o} followed by the job number) for errors. Then list your
  directory contents.
\item
  You should now see multiple new files for each database (with
  extensions like \texttt{.nhr}, \texttt{.nin}, \texttt{.nsq}, etc.).
  These are the indexed database files that BLAST will use for rapid
  searching. Count how many new files were created per input FASTA file.
\item
  Check your job efficiency using \texttt{jobinfo\ -j\ JOBNUMBER}. Note
  the memory and CPU usage---this information helps you optimize future
  jobs. What was your memory efficiency? CPU efficiency?
\end{enumerate}

\section{Step 3: Run Your First BLAST
Search}\label{step-3-run-your-first-blast-search}

Now you'll compare your unknown transcripts \texttt{test.fasta} against
each of the three databases you created. First, examine
Figure~\ref{fig-blast} to determine which program to use within BLAST.

\begin{tcolorbox}[enhanced jigsaw, rightrule=.15mm, leftrule=.75mm, colbacktitle=quarto-callout-note-color!10!white, colframe=quarto-callout-note-color-frame, toprule=.15mm, opacitybacktitle=0.6, title=\textcolor{quarto-callout-note-color}{\faInfo}\hspace{0.5em}{Programs within BLAST}, left=2mm, colback=white, bottomtitle=1mm, toptitle=1mm, bottomrule=.15mm, coltitle=black, titlerule=0mm, breakable, opacityback=0, arc=.35mm]

\begin{figure}[H]

\centering{

\includegraphics[width=0.65\linewidth,height=\textheight,keepaspectratio]{images/Figure4_2.png}

}

\caption{\label{fig-blast}Overview of BLAST five main Algorithms}

\end{figure}%

There are five major algorithms within BLAST. In making your database,
you already know your input sequences are nucelotides. Example the input
file \texttt{test.fasta} to determine which of the five programs you
will use when running your BLAST search.

\end{tcolorbox}

\begin{enumerate}
\def\labelenumi{\arabic{enumi}.}
\item
  Since both your query and database sequences are nucleotides, you'll
  use the \texttt{blastn} program. Review the available options:

\begin{Shaded}
\begin{Highlighting}[]
\ExtensionTok{blastn} \AttributeTok{{-}help}
\end{Highlighting}
\end{Shaded}

  The essential parameters are:

  \begin{itemize}
  \tightlist
  \item
    \texttt{-query}: Your input FASTA file
  \item
    \texttt{-db}: The database to search against
  \item
    \texttt{-out}: Output file name
  \end{itemize}
\item
  Run BLAST to compare \texttt{test.fasta} against the chloroplast
  database. For example:

\begin{Shaded}
\begin{Highlighting}[]
\ExtensionTok{blastn} \AttributeTok{{-}query}\NormalTok{ test.fasta }\AttributeTok{{-}db}\NormalTok{ ATcp.fasta }\AttributeTok{{-}out}\NormalTok{ test\_vs\_cp.blastn\_out}
\end{Highlighting}
\end{Shaded}

  Use a descriptive output filename that indicates the query and
  database (e.g., \texttt{test\_vs\_cp.blastn\_out} for test query
  versus chloroplast database).
\item
  Repeat for the other two databases (mitochondrial and chromosome V).
  Add all three BLAST commands to your script.
\item
  Submit your updated script to the queue:

  \begin{itemize}
  \tightlist
  \item
    Number of CPUs = 1
  \item
    RAM/Memory = 100MB
  \item
    Queue = class
  \item
    Time = 01:00:00
  \end{itemize}

  \begin{quote}
  \textbf{Tip}: After successfully creating the blast databases, comment
  out the \texttt{makeblastdb} commands (add \texttt{\#} at the start of
  each line) so they don't rerun unnecessarily as we continue to add to
  this script.
  \end{quote}
\item
  Once the job completes, examine one of the output files using
  \texttt{less}. This is BLAST's default ``pairwise'' format, similar to
  what you see on the NCBI website. While human-readable, it's difficult
  to parse computationally.
\end{enumerate}

\section{Step 4: Optimize Output Format for
Parsing}\label{step-4-optimize-output-format-for-parsing}

The default BLAST output is verbose and hard to parse with command-line
tools. BLAST offers alternative formats that are more suitable for
automated analysis.

\begin{enumerate}
\def\labelenumi{\arabic{enumi}.}
\item
  Review the \texttt{-outfmt} option in the \texttt{blastn\ -help}
  output. You'll see numbered format options (0-18). Formats 6 and 7 are
  tabular and particularly useful for parsing.
\item
  Rerun one of your BLAST searches using \texttt{-outfmt\ 6}:

\begin{Shaded}
\begin{Highlighting}[]
\ExtensionTok{blastn} \AttributeTok{{-}query}\NormalTok{ test.fasta }\AttributeTok{{-}db}\NormalTok{ ATcp.fasta }\AttributeTok{{-}outfmt}\NormalTok{ 6 }\AttributeTok{{-}out}\NormalTok{ test\_vs\_cp.tab}
\end{Highlighting}
\end{Shaded}
\item
  Examine the output with \texttt{less}. Format 6 produces tab-delimited
  columns:

  \begin{itemize}
  \tightlist
  \item
    Column 1: Query sequence ID
  \item
    Column 2: Subject (database) sequence ID
  \item
    Column 3: Percent identity
  \item
    Column 4: Alignment length
  \item
    Column 5: Number of mismatches
  \item
    Columns 6-12: Additional alignment statistics including E-value
  \end{itemize}
\item
  Compare format 6 to format 7 (which includes column headers). Which do
  you prefer? As a class, decide on one format to use consistently.
\item
  Update your script to use your chosen output format for all three
  BLAST searches. Update the queue parameters based on your previous
  job's runtime and memory usage (check with \texttt{jobinfo}):

  \begin{itemize}
  \tightlist
  \item
    Number of CPUs = 1
  \item
    RAM/Memory = ? (what did you actually use last time?)
  \item
    Queue = class
  \item
    Time = ? (how long did it actually take?)
  \end{itemize}
\item
  Comment out the previous BLAST commands and add the new versions.
  Submit to the queue.
\end{enumerate}

\section{Step 5: Understanding Database Size Effects on
E-values}\label{step-5-understanding-database-size-effects-on-e-values}

The E-value (expect value) represents the number of hits you'd expect to
see by chance in a database of this size. Larger databases generally
produce larger E-values for the same match.

\begin{enumerate}
\def\labelenumi{\arabic{enumi}.}
\item
  To avoid duplicating large files, I have pre-downloaded Chr1 of the
  \emph{Arabidopsis} genome. You'll create a symbolic link to the
  database folder (make sure to replace your username as you did during
  setup):

\begin{Shaded}
\begin{Highlighting}[]
\FunctionTok{ln} \AttributeTok{{-}s}\NormalTok{ /home/aubclsfXXXX/bioinf\_shared/Arabidopsis\_thaliana/CHR\_I}
\end{Highlighting}
\end{Shaded}

  A symbolic link acts like a shortcut---the files appear in your
  directory but remain stored in the original location.
\item
  BLAST allows you to search multiple databases simultaneously by
  listing them in quotes:

\begin{Shaded}
\begin{Highlighting}[]
\ExtensionTok{blastn} \AttributeTok{{-}db} \StringTok{"ATcp.fasta ATchrV.fasta CHR\_I/NC\_003070.fasta ATmt.fasta"} \DataTypeTok{\textbackslash{}}
  \AttributeTok{{-}query}\NormalTok{ test.fasta }\AttributeTok{{-}outfmt}\NormalTok{ 6 }\AttributeTok{{-}out}\NormalTok{ combined\_search.txt}
\end{Highlighting}
\end{Shaded}

  This searches all four databases (chloroplast, chrV, chrI, and
  mitochondrial) in one run.

  \begin{quote}
  Note: If you named your BLAST databases something different than the
  default names, then you will need to edit this to match your database
  naming conventions here.
  \end{quote}
\item
  Add this command to your script, comment out previous lines, and
  submit to the queue with updated parameters.
\item
  After completion, compare the E-values for a specific hit:

  \begin{itemize}
  \tightlist
  \item
    Find a transcript that matched the chloroplast in your earlier
    single-database search
  \item
    Look up the same transcript in your combined search output
  \item
    How did the E-value change? Why? \textbf{Consider}: The combined
    database is much larger than the chloroplast alone.
  \end{itemize}
\end{enumerate}

\section{Step 6: Parse BLAST Output to Count Hits Per
Database}\label{step-6-parse-blast-output-to-count-hits-per-database}

Now you'll use command-line text processing tools to extract meaningful
summaries from your BLAST results. The goal is to count how many of your
query transcripts match each database (mitochondrial, chloroplast, or
nuclear chromosomes). You will be making a long one-liner with multiple
piped steps that you will build incrementally below.

\begin{enumerate}
\def\labelenumi{\arabic{enumi}.}
\item
  Your BLAST output columns 1 and 2 contain the query name and database
  hit. Extract just these columns using \texttt{awk}:

\begin{Shaded}
\begin{Highlighting}[]
\ExtensionTok{blastn} \AttributeTok{{-}db} \StringTok{"ATcp.fasta ATchrV.fasta CHR\_I/NC\_003070.fasta ATmt.fasta"} \DataTypeTok{\textbackslash{}}
  \AttributeTok{{-}query}\NormalTok{ test.fasta }\AttributeTok{{-}outfmt}\NormalTok{ 6 }\KeywordTok{|} \FunctionTok{awk} \StringTok{\textquotesingle{}\{print $1,$2\}\textquotesingle{}}
\end{Highlighting}
\end{Shaded}

  The \texttt{\textbar{}} (pipe) sends BLAST output directly to
  \texttt{awk} without creating an intermediate file. In the next
  several steps, add each command via pipe to this command line
  argument. In the end, you will redirect the output to a file.

  \begin{quote}
  Keep in mind that if you run the command above by itself, the output
  will be LARGE! So if you want to test each step of your piped
  one-liner, simplify the output by piping it to \texttt{head}
  periodically. But \textbf{beware}! Do not accidentally leave a
  \texttt{\textbar{}\ head\ \textbar{}} in the middle of your one-liner
  which will effectively remove a vast majority of your output hits!
  \end{quote}
\item
  Each query might have multiple hits to the same database sequence.
  Remove duplicates with another pipe to \texttt{sort} and then
  \texttt{uniq} to yield unique query-database pairs (one line per
  transcript-database match).
\end{enumerate}

\begin{quote}
Hint: Do not forget to use the man page when needed to remind yourself
how to use these tools.
\end{quote}

\begin{enumerate}
\def\labelenumi{\arabic{enumi}.}
\setcounter{enumi}{2}
\item
  Now focus only on which databases were hit by adding a pipe to awk to
  isolate only column 2. This will remove the query names.
\item
  Next, count how many times each database appears by piping the output
  to \texttt{uniq\ -c} (which requires sorted input).
\item
  Sort the counts numerically and save to a file called
  \texttt{RawCounts.txt}.
\item
  Examine \texttt{RawCounts.txt}. You should see counts for each
  database. Did you get more hits to chromosome I (the full chromosome)
  or to chromosome V (partial)? Why might that be?
\item
  Both chrV and chrI should be ``NT'', but since we used a symbolic link
  for chrI, it shows up differently in our output. Use \texttt{sed} to
  convert chrI hits to \textbf{\emph{NT}} so we can combine hits for
  these two chromosomes.
\item
  Add this complete one-liner to your script. This is a powerful example
  of piping multiple tools together---a common pattern in
  bioinformatics.
\end{enumerate}

\section{Step 7: Identify Sequences with No BLAST
Hits}\label{step-7-identify-sequences-with-no-blast-hits}

Some of your query transcripts may not match any database, suggesting
they could be novel genes, misassembled sequences, or contamination.

\begin{enumerate}
\def\labelenumi{\arabic{enumi}.}
\item
  Sequences with no BLAST hits won't appear in your output file at all.
  Calculate the number:

  \begin{itemize}
  \tightlist
  \item
    Total sequences in \texttt{test.fasta} (you counted this in Step 1)
  \item
    Minus: Number of sequences with at least one hit
  \end{itemize}
\item
  To count sequences with hits, extract unique query IDs from your BLAST
  output:

\begin{Shaded}
\begin{Highlighting}[]
\FunctionTok{awk} \StringTok{\textquotesingle{}\{print $1\}\textquotesingle{}}\NormalTok{ combined\_search.txt }\KeywordTok{|} \FunctionTok{sort} \KeywordTok{|} \FunctionTok{uniq} \KeywordTok{|} \FunctionTok{wc} \AttributeTok{{-}l}
\end{Highlighting}
\end{Shaded}
\item
  Calculate:
  \texttt{(total\ sequences)\ -\ (sequences\ with\ hits)\ =\ sequences\ with\ no\ hits}
\item
  Alternatively, if you want to know which specific sequences had no
  hits, you could:

  \begin{itemize}
  \tightlist
  \item
    Extract all sequence IDs from \texttt{test.fasta}:
    \texttt{grep\ "\textgreater{}"\ test.fasta\ \textgreater{}\ all\_queries.txt}
  \item
    Extract hit IDs from BLAST output:
    \texttt{awk\ \textquotesingle{}\{print\ \$1\}\textquotesingle{}\ combined\_search.txt\ \textbar{}\ sort\ \textbar{}\ uniq\ \textgreater{}\ hit\_queries.txt}
  \item
    Find the difference:
    \texttt{comm\ -23\ all\_queries.txt\ hit\_queries.txt}
  \end{itemize}
\item
  Add the count of sequences with no hits to the bottom of your
  \texttt{RawCounts.txt} file using \texttt{echo}:

\begin{Shaded}
\begin{Highlighting}[]
\BuiltInTok{echo} \StringTok{"No hits: [your count]"} \OperatorTok{\textgreater{}\textgreater{}}\NormalTok{ RawCounts.txt}
\end{Highlighting}
\end{Shaded}
\end{enumerate}

\section{Step 8: Finalize Your
Script}\label{step-8-finalize-your-script}

\begin{enumerate}
\def\labelenumi{\arabic{enumi}.}
\item
  Review your \texttt{Example\_script.sh}. It should now contain:

  \begin{itemize}
  \tightlist
  \item
    Commands to create three BLAST databases
  \item
    Commands to run BLAST searches (with appropriate output formatting)
  \item
    A one-liner to parse results and count hits per database
  \item
    Calculation of sequences with no hits
  \end{itemize}
\item
  Add comments throughout your script to explain what each section does.
\item
  Test that your script runs successfully from start to finish. You can
  do this by:

  \begin{itemize}
  \tightlist
  \item
    Removing all output files
  \item
    Resubmitting the script (with all commands uncommented)
  \item
    Verifying all expected outputs are generated
  \end{itemize}
\end{enumerate}

This script is now a reusable tool. With minor modifications, you could
use it to annotate sequences from any organism against any reference
genome.

\section{Reflection Questions}\label{reflection-questions-12}

Now that you've completed the lab, reflect on the following:

\begin{enumerate}
\def\labelenumi{\arabic{enumi}.}
\item
  \textbf{Command-line advantages}: What are some specific scenarios
  where command-line BLAST would be preferable to web-based BLAST?
  Consider factors like:

  \begin{itemize}
  \tightlist
  \item
    Number of sequences
  \item
    Database selection
  \item
    Reproducibility
  \item
    Integration with other tools
  \end{itemize}
\item
  \textbf{Nucleotide vs.~protein BLAST}: For identifying orthologs
  across distantly related species, is it better to use nucleotide
  sequences (blastn) or protein sequences (blastp/blastx)? Why?
  \emph{Hint: Consider how DNA vs.~protein sequences evolve over time.}
\item
  \textbf{Novel sequences}: Based on your \texttt{RawCounts.txt}
  results:

  \begin{itemize}
  \tightlist
  \item
    How many transcripts matched the chloroplast genome?
  \item
    How many matched the mitochondrial genome?
  \item
    How many matched nuclear chromosomes?
  \item
    How many had no hits in any database?
  \end{itemize}

  For the sequences with no hits, what are possible biological
  explanations? What would you do next to characterize these sequences?
\item
  \textbf{E-value interpretation}: You observed how database size
  affects E-values. If you're comparing BLAST results from different
  studies that used different reference databases, why is it important
  to consider database size? How might this affect your interpretation
  of ``significant'' hits?
\item
  \textbf{Applications}: Beyond annotation, describe two other
  biological questions or research applications where command-line BLAST
  would be useful. Be specific about what you'd search and why.
\end{enumerate}

\section{Deliverables}\label{deliverables-1}

Upload the following to Canvas:

\begin{enumerate}
\def\labelenumi{\arabic{enumi}.}
\tightlist
\item
  Your completed \texttt{Example\_script.sh} file (with comments
  explaining each section)
\item
  Your \texttt{RawCounts.txt} file
\item
  Answers to the five reflection questions above (in a Word document or
  PDF)
\end{enumerate}

\section{Acknowledgement}\label{acknowledgement}

This lab assignment has been adapted from an activity initially composed
by Dr.~Nathan Hall from Auburn University. It has been further modified
through multiple cohorts of students and more recently with the
assistance of AI.

\section{Additional Resources}\label{additional-resources}

\textbf{Extending Your BLAST Skills:}

\begin{itemize}
\item
  \textbf{Primer design}: BLAST can verify primer specificity. For
  large-scale primer design (\textgreater5 loci), automating BLAST
  checks saves time. Python's
  \href{https://primer3.org/manual.html}{primer3} wrapper can be
  integrated with command-line BLAST.
\item
  \textbf{Ortholog discovery}:
  \href{http://seqanswers.com/forums/showthread.php?t=27735}{SeqAnswers
  Thread} discusses automating BLAST for ortholog identification across
  multiple genomes---useful for comparative genomics projects.
\item
  \textbf{Alternative alignment tools}:

  \begin{itemize}
  \tightlist
  \item
    \href{http://www.bx.psu.edu/~rsharris/lastz/}{LASTZ}: Better for
    whole-genome alignments
  \item
    \href{http://bio-bwa.sourceforge.net/}{BWA}: Optimized for mapping
    short reads to reference genomes
  \item
    \href{http://darlinglab.org/mauve/mauve.html}{Mauve}: Multiple
    genome alignment and visualization
  \end{itemize}
\item
  \textbf{Visualization in R}: The
  \href{http://genoplotr.r-forge.r-project.org/}{genoPlotR} package
  creates publication-quality figures from BLAST and Mauve outputs,
  useful for comparing genome organization.
\item
  \textbf{BLAST+ Documentation}: Comprehensive guide to all BLAST+ tools
  and parameters:
  \href{https://www.ncbi.nlm.nih.gov/books/NBK279690/}{NCBI BLAST+
  Command Line Applications User Manual}
\end{itemize}

\section{References}\label{references-13}

\chapter{Lab06: HMMs in Predicting Gene
Models}\label{lab06-hmms-in-predicting-gene-models}

Hidden Markov models (HMMs) are a set of mathematical tools that can be
used to draw inferences about genomic and evolutionary processes. In
general, an HMM calculates the probability of each scenario that could
have resulted in an observed data set, then uses statistical methods to
determine which scenarios are most likely to have occurred. For example,
HMMs can be applied to DNA sequence data from different species to infer
the likely pattern of relationships among those species.

\section{Learning Goals}\label{learning-goals-1}

\begin{itemize}
\tightlist
\item
  Students will understand the basic structure of an HMM, the types of
  data used in ab initio gene prediction, and its consequent
  limitations.
\item
  Students will understand the process of gene exon-intron boundary
  prediction, having altered the values for the key parameters and noted
  the impact on the prediction.
\item
  Students will consider how the exon-intron boundary model might be
  designed to be more accurate, thinking about the biological features
  of the process and how they might be represented in a computer
  program.
\end{itemize}

\section{Learning Objectives}\label{learning-objectives-2}

\begin{itemize}
\tightlist
\item
  Students will be able to process unannotated genomic data using ab
  initio gene finders as well as other inputs.
\item
  Students will be able to defend the proposed gene annotation.
\item
  Students will reflect on the other uses for HMMs.
\end{itemize}

\section{Step-by-Step Instructions:}\label{step-by-step-instructions}

\subsection{Part 1: Reviewing the
basics}\label{part-1-reviewing-the-basics}

The first part of this lab activity reviews some of the concepts
introduced in \hyperref[hidden-markov-models-in-biology]{Chapter 9} on
how HMMs work to predict a specific gene model based on a DNA sequence.
Specifically, you will use the simple Hidden Markov Model introduced in
Figure~\ref{fig-toy-hmm} to calculate the various transition and
emmission probabilities of given State Paths. This approach is used to
identify the most likely path, thus identifying the location of the
transition from exon to 5'SS to intron, or gene model.

\begin{enumerate}
\def\labelenumi{\arabic{enumi}.}
\tightlist
\item
  On a separate sheet of paper, draw a state diagram for an HMM that
  will scan an un-annotated genome for genes with a start codon, a stop
  codon, exons, introns, a 5' splice site and a 3' splice site. Use
  Figure~\ref{fig-toy-hmm} as a guide.
\end{enumerate}

\begin{itemize}
\tightlist
\item
  What are the ``states'' of the model that you have drawn?
\end{itemize}

\begin{enumerate}
\def\labelenumi{\arabic{enumi}.}
\setcounter{enumi}{1}
\tightlist
\item
  Using the simple state HMM you have drawn out, calculate the
  probability of EACH of the state paths listed in the table below given
  the observed DNA sequence shown in the top row.
\end{enumerate}

\begin{quote}
To calculate the probability of each possible state path, multiply all
of the transition and emission probabilities for each base in the
sequence above given the state in each part of the path. Take the log of
the product of all probabilities across all states in a path to obtain
the log~P (see Note~\ref{nte-hmm}).
\end{quote}

\begin{longtable}[]{@{}llllllll@{}}
\toprule\noalign{}
& & A & G & T & G & A & \\
\midrule\noalign{}
\endhead
\bottomrule\noalign{}
\endlastfoot
Path 1 & \emph{Start} & E & E & E & 5 & I & \emph{End} \\
Path 2 & \emph{Start} & E & E & 5 & I & I & \emph{End} \\
Path 3 & \emph{Start} & E & 5 & I & I & I & \emph{End} \\
\end{longtable}

\begin{itemize}
\item
  \textbf{Which of the state paths below is the most likely one to
  annotate the DNA sequence at the top correctly?}
\item
  Note that each state path has the 5' splice site at a different
  position in the sequence. Which 5' splice site is likely to be the
  correct one?
\item
  Given what you know about splice sites in eukaryotic genes, does
  anything in this model seem problematic?
\end{itemize}

\begin{enumerate}
\def\labelenumi{\arabic{enumi}.}
\setcounter{enumi}{2}
\tightlist
\item
  HMMs assume that each base is an independent observation. For
  instance, they assume that the probability of the ith base being in an
  intron does not affect whether the (i+1)th base is in an intron.
  \textbf{Give an example of a genomic feature where this assumption
  does not hold. Explain your reasoning.}
\end{enumerate}

\subsection{Part 2: Exploring an HMM}\label{part-2-exploring-an-hmm}

For the next part of the lab, we will use an excel spreadsheet provided
by Anton Weisstein to explore properties of Hidden Markov Models. The
Excel workbook ``Hidden Markov Model'' illustrates the mathematical
workings of an HMM, using the example of locating the 5' splice site
within a DNA sequence (see
\hyperref[hidden-markov-models-in-biology]{Chapter 9} for more detail).
By manipulating model parameters such as the average length and base
composition of introns vs.~exons, the user can study how such changes
influence the model's decisions about where the splice site is likeliest
to occur. For ease of demonstration, the workbook begins by analyzing a
very short sequence (6 bp) before re-creating Eddy's full, 26-bp
model\textsuperscript{4}.

\begin{enumerate}
\def\labelenumi{\arabic{enumi}.}
\item
  Download and open the excel file
  ``\href{datasets/Hidden\%20Markov\%20Model.intron.xls}{HiddenMarkovModel.intron.xls}''.
  Before changing any of the sliders on the spreadsheet, save a copy of
  the file under a different name.
\item
  The first sheet in the workbook, appropriately named ``Simple Model'',
  demonstrates the calculations involved in building an HMM to determine
  the exact location of an exon/intron boundary. The user enters a short
  DNA sequence and sets model parameters. The workbook then uses these
  values to compute the likelihood of each potential 5' splice site
  location.

  \begin{itemize}
  \tightlist
  \item
    Cells \textbf{B2--B4} contain the probabilities that the DNA
    sequence begins in different regions of a gene. In HMM parlance,
    these regions are described as states of the model.
  \item
    Cells \textbf{F4--I6} give the probability per nucleotide of moving
    from one state to another. For example, if G4 = 0.60, then a
    nucleotide in the gene's exon has a 60\% chance of being followed by
    a nucleotide in the splice site. This model focuses only on 5'
    splice sites, so states will always occur in the order Exon --
    Splice -- Intron -- End.
  \item
    Cells \textbf{F12--I14} give the state-specific probability of
    producing, or emitting, a particular nucleotide. In the initial
    model, all four nucleotides are equally likely to occur in an exon,
    but introns are GC-poor. For the sake of simplicity, the splice site
    is considered to be only one nucleotide long: in the initial model,
    that nucleotide is always a guanine.\\
  \item
    Finally, Cells \textbf{B17--G17} contain the nucleotide sequence to
    be analyzed.
  \end{itemize}
\item
  Cells \textbf{B21--L51} show the model's calculations. For each
  potential state path, the model computes:

  \begin{enumerate}
  \def\labelenumii{\roman{enumii}.}
  \item
    the probability that the DNA sequence starts in the appropriate
    state,
  \item
    the probability of precisely following the state path, and
  \item
    the probability of emitting the observed DNA sequence if the state
    path is followed.
  \end{enumerate}
\end{enumerate}

For example, state path 1 (Cells \textbf{B21--H21}) posits that the
first nucleotide, a cytosine, occurs in an intron. The initial model
gives a 0\% probability of starting in an intron (Cell \textbf{B4}), and
a 10\% frequency of cytosines within an intron (Cell \textbf{G14}), so
these values are reported. The second nucleotide, a guanine, is also
proposed to occur in an intron. The initial model therefore reports the
90\% probability of remaining within an intron (Cell \textbf{H6}) and
the 10\% frequency of guanines within an intron (Cell \textbf{H14}).
This process continues through all six nucleotides, then reports the
10\% probability of transitioning from the last nucleotide's state ---
an intron --- to the end of the sequence (Cell \textbf{I6}).

\begin{itemize}
\item
  HMMs assume that each nucleotide is \textbf{independent} of the
  others, much as the result of a coin flip does not depend on previous
  or future flips. This assumption of independence is biologically
  unrealistic, but it greatly simplifies the math: each nucleotide's
  individual probabilities can be simply \textbf{multiplied} together to
  calculate the total probability of the overall sequence. For state
  path 1, this combined probability is reported in Cell \textbf{J22}.
  Unsurprisingly, given that one of the numbers being multiplied was
  zero, the resulting product is also zero, meaning that this state path
  is not consistent with the observed sequence. In fact, under the
  initial model, state paths 3 and 5 are the only ones possible, because
  only they correctly place the splice site at a G.
\item
  Computationally, it is faster to add numbers together than to multiply
  them. This can make a huge difference when running a large HMM that
  may contain many millions of state paths. For this reason, most HMMs
  calculate the combined probability of a state path not by multiplying
  the individual nucleotide probabilities, but by \textbf{adding the
  logarithms} of those probabilities, using the identity. The resulting
  log-probabilities are reported in Column K. Note that ln(0) is
  undefined, so state paths with probability zero yield an error message
  in this column.
\item
  If the model correctly includes all possible state paths, adding up
  all of the probabilities in \textbf{Column J} should give the total
  probability p of emitting the observed sequence. But by definition,
  that sequence has indeed been observed, so the individual state path
  probabilities should be calibrated to \textbf{sum to one}. (In much
  the same way, if I have prior knowledge that a poker player is holding
  either a flush or a full house, each of those two hands is much
  likelier than if I lacked such knowledge.) Each state path's
  probability is therefore divided by the total probability p to compute
  the path's likelihood given observed DNA sequence. The results are
  displayed in \textbf{Column L} (which does sum to 1) and graphed in
  the \textbf{likelihood profile} at the bottom of the worksheet.
\end{itemize}

\textbf{\emph{Now that we have reviewed how this worksheet works, let's
perturb it to better understand WHY it works!}}

\begin{enumerate}
\def\labelenumi{\arabic{enumi}.}
\setcounter{enumi}{3}
\item
  Use the first slider to change the \textbf{transition probability}
  from exon to splice site (top row). For each value (0.2, 0.4, 0.6),
  use the second slider to vary the transition probability from intron
  to intron (bottom row) across three different values (0.7, 0.9, 1.0).
  Record in the table below how these perturbations affect the
  likelihood profile of the 5' splice site position (describe which
  position(s) show highest likelihood and approximate probability values
  from the bar graph at the bottom of the spreadsheet). \textbf{Why do
  these changes have these effects?}

  \begin{longtable}[]{@{}
    >{\centering\arraybackslash}p{(\linewidth - 8\tabcolsep) * \real{0.0986}}
    >{\centering\arraybackslash}p{(\linewidth - 8\tabcolsep) * \real{0.0986}}
    >{\centering\arraybackslash}p{(\linewidth - 8\tabcolsep) * \real{0.2817}}
    >{\centering\arraybackslash}p{(\linewidth - 8\tabcolsep) * \real{0.1690}}
    >{\raggedright\arraybackslash}p{(\linewidth - 8\tabcolsep) * \real{0.3521}}@{}}
  \toprule\noalign{}
  \begin{minipage}[b]{\linewidth}\centering
  Exon→Splice transition
  \end{minipage} & \begin{minipage}[b]{\linewidth}\centering
  Intron→Intron transition
  \end{minipage} & \begin{minipage}[b]{\linewidth}\centering
  Predicted splice position(s)
  \end{minipage} & \begin{minipage}[b]{\linewidth}\centering
  Likelihood (\%)
  \end{minipage} & \begin{minipage}[b]{\linewidth}\raggedright
  Biological interpretation
  \end{minipage} \\
  \midrule\noalign{}
  \endhead
  \bottomrule\noalign{}
  \endlastfoot
  0.2 & 0.7 & & & \\
  0.2 & 0.9 & & & \\
  0.2 & 1.0 & & & \\
  0.4 & 0.7 & & & \\
  0.4 & 0.9 & & & \\
  0.4 & 1.0 & & & \\
  0.6 & 0.7 & & & \\
  0.6 & 0.9 & & & \\
  0.6 & 1.0 & & & \\
  \end{longtable}

  When you are done, reset to the original 0.6 exon→splice and 0.9
  intron→intron transition probabilities.
\item
  Now systematically vary the \textbf{emission probabilities} for the
  splice site state. Starting with the default model where G = 1.0 at
  the splice site, change the emission probabilities to make the splice
  site less specific. \textbf{Record in the table below how reducing the
  G content at the splice site affects the model's predictions. What
  does this tell you about the importance of strong splice site
  consensus sequences?}

  \begin{longtable}[]{@{}
    >{\centering\arraybackslash}p{(\linewidth - 6\tabcolsep) * \real{0.4474}}
    >{\centering\arraybackslash}p{(\linewidth - 6\tabcolsep) * \real{0.2632}}
    >{\centering\arraybackslash}p{(\linewidth - 6\tabcolsep) * \real{0.1579}}
    >{\raggedright\arraybackslash}p{(\linewidth - 6\tabcolsep) * \real{0.1316}}@{}}
  \toprule\noalign{}
  \begin{minipage}[b]{\linewidth}\centering
  Splice site emission probabilities
  \end{minipage} & \begin{minipage}[b]{\linewidth}\centering
  Predicted splice position(s)
  \end{minipage} & \begin{minipage}[b]{\linewidth}\centering
  Max likelihood (approx.)
  \end{minipage} & \begin{minipage}[b]{\linewidth}\raggedright
  Comments
  \end{minipage} \\
  \midrule\noalign{}
  \endhead
  \bottomrule\noalign{}
  \endlastfoot
  A=0.00, C=0.00, G=1.00, T=0.00 (default) & & & Strong consensus \\
  A=0.05, C=0.00, G=0.95, T=0.00 & & & \\
  A=0.10, C=0.00, G=0.90, T=0.00 & & & \\
  A=0.10, C=0.10, G=0.70, T=0.10 & & & \\
  A=0.25, C=0.25, G=0.25, T=0.25 & & & No consensus \\
  \end{longtable}

  When you are done, reset to the original splice site emission
  probabilities.
\item
  If your model predicts multiple positions in the sequence have the
  same likelihood to be a splice site, how could you use RNA-Seq data to
  identify the best splice site candidate?
\item
  Consider the notion that transition probabilities in an evidence-based
  gene prediction algorithm can be calculated based on the length of a
  genetic feature (exons, introns, etc). That is, transition
  probabilities are a function of average gene length. \textbf{If an HMM
  were trained using the genome of an organism with many short genes and
  few long genes, would you expect this HMM to predict more long genes
  or more short genes on an un-annotated genome? Why?}
\end{enumerate}

\subsection{Part 3: Model
Perturbations}\label{part-3-model-perturbations}

In this part of the lab, you will explore how modifications to the HMM
affect its predictions. Choose \textbf{at least one} of the following
explorations (items 1-4), document the changes you make, and record the
resulting effects on the model's splice site predictions. For each
exploration you choose, you should test multiple parameter values to
observe trends, not just a single change.

\begin{enumerate}
\def\labelenumi{\arabic{enumi}.}
\item
  \textbf{Change the observed DNA sequence.} Alterations in GC richness
  and in the position of any G nucleotides may influence the likelihood
  of the 5' splice site occurring at a specific position. Note that the
  sequence must consist only of standard DNA nucleotides (A, C, G, and
  T).

  \emph{Suggestion: Try sequences with different numbers of G's and/or
  different GC content (e.g., GC-rich vs.~AT-rich).}
\item
  \textbf{Change the base emission probabilities for exon and/or intron
  states.} Alterations in the base composition of exons vs.~introns will
  affect the model's inferences about the best splice site position. In
  particular, changing the emission probabilities so that exons and
  introns have similar base compositions will severely degrade the
  model's ability to predict the splice site location with any degree of
  confidence.

  \begin{quote}
  Note that each position within a given gene region must be occupied by
  one of the four nucleotides, so each row of the emission probability
  matrix must sum to one. If this condition is violated, the offending
  row will be highlighted in red.
  \end{quote}

  \emph{Suggestion: Try making introns more GC-rich or making exons and
  introns compositionally similar.}
\item
  \textbf{Change the state transition probabilities using the built-in
  scrollbars.} Again, each row of the transition matrix must sum to one.

  \begin{quote}
  Note that the transition probability (p) for the exon state to the
  splice site state is the reciprocal of the average exon length. (This
  is because the mean of a geometric distribution is 1/p.) For example,
  if the transition probability for the exon-to-splice site state is
  10\%, then the average exon will be 10 bp long. Similarly, the
  intron-to-end transition probability is the reciprocal of the average
  intron length. Adjusting these parameters will therefore shift the
  most likely splice site position to earlier or later in the sequence.
  The model assumes that the splice site consists of a single
  nucleotide, so the probabilities in the middle row cannot be changed.
  \end{quote}

  \emph{Suggestion: Experiment with different average exon and intron
  lengths by adjusting transition probabilities. What happens with very
  short vs.~very long predicted gene features?}
\item
  \textbf{Change the vector of initial state probabilities.} The
  original model assumes that the DNA sequence starts in an exon (100\%
  probability in cell B2). However, the start of the DNA sequence could
  be part of an intron or even at the splice site itself. A better model
  should account for this possibility and assign non-zero initial
  probabilities to all three states.

  \begin{quote}
  Note that, by definition, the sequence must start in one of these
  three states and these initial probabilities must sum to one.
  \end{quote}

  \emph{Suggestion: Try splitting the initial probability equally among
  all three states, or giving higher probability to intron starts.}
\end{enumerate}

\textbf{For the exploration you choose to conduct document:}

\begin{itemize}
\tightlist
\item
  The specific parameter(s) you changed and the values you tested
\item
  The effect on the predicted splice site position(s)
\item
  The effect on the likelihood values
\item
  Your interpretation of why these changes had the observed effects
\end{itemize}

\subsection{Part 4: Exploration of the Full
Model}\label{part-4-exploration-of-the-full-model}

The second sheet of the ``Hidden Markov Model'' workbook contains an
exact replica of the HMM described by Eddy (2004)\textsuperscript{4}.
This sheet can be reached using the navigation tabs at the bottom of the
Excel workspace. Eddy's HMM is designed to handle DNA sequences up to 26
bp in length; moreover, under the initial settings, the splice site can
occur at either an A or a G, although the latter is far more likely.

\textbf{Key differences from the Simple Model:}

\begin{itemize}
\tightlist
\item
  Handles sequences up to 26 nucleotides instead of just 6
\item
  The splice site can occur at either an A or a G (though G is far more
  likely)
\item
  You can paste sequences directly into Cell E16 rather than entering
  individual nucleotides
\item
  The likelihood profile shows more nuanced predictions across a longer
  sequence
\item
  Additional calculations are available in separate worksheets for those
  interested in the mathematical details
\end{itemize}

This sheet's controls work similarly to the Simple Model. The interface
is streamlined to focus on the model's controls and output graph. If you
want to examine the detailed calculations, scroll down to Rows 44--183
or navigate to the ``Additional Calculations'' sheet.

\textbf{Exploratory Activity:} Choose \textbf{ONE} exploration from the
options below. Test multiple parameter values or sequences to observe
trends and patterns. Document what you changed, what happened, and your
interpretation of the results.

\subsubsection{Option 1: Base Composition
Effects}\label{option-1-base-composition-effects}

Investigate how the compositional difference between exons and introns
affects splice site prediction. Start with the default settings where
exons have equal base frequencies (25\% each) and introns are AT-rich
(A=40\%, T=40\%, G=10\%, C=10\%).

Systematically change the emission probabilities to make introns less
AT-rich (more similar to exons). Test at least 3-4 different scenarios
ranging from strong contrast to no contrast. For each scenario, record
which position is predicted as the splice site and the maximum
likelihood value.

\textbf{Key questions to explore:} How much compositional difference is
needed for confident predictions? What happens when exons and introns
have identical base composition?

\subsubsection{Option 2: Sequence
Variation}\label{option-2-sequence-variation}

Test how different DNA sequences are interpreted by the model while
keeping all parameters constant. Generate or find at least 3-4 different
26 bp sequences with varying characteristics (GC-rich, AT-rich, multiple
G's in different positions, random sequences, or real genomic
sequences).

For each sequence, note its GC content, where G nucleotides are located,
and where the model predicts the splice site.

\textbf{Key questions to explore:} Does the model always pick the first
G? What happens with closely spaced G's? How important is the position
of G nucleotides versus overall sequence composition?

\subsubsection{Option 3: Species-Specific Gene
Architecture}\label{option-3-species-specific-gene-architecture}

Research typical exon and intron lengths for a specific organism
(budding yeast, fruit fly, human, arabidopsis, or another of your
choice). Use online databases like NCBI Gene, Ensembl, or model organism
databases to find this information for several genes.

Calculate average lengths and convert them to transition probabilities
using the formula: self-transition probability ≈ (average length - 1) /
(average length). Enter these organism-specific values into the model
and observe how predictions change.

\textbf{Key questions to explore:} How do your organism's gene
architecture parameters differ from the defaults? What would happen if
you used a yeast-trained HMM to predict human genes, or vice versa?

\subsubsection{Option 4: Statistical Power
Analysis}\label{option-4-statistical-power-analysis}

Explore the relationship between base composition difference and
prediction confidence. Create a series of tests where you systematically
vary the AT content difference between exons and introns (e.g., 30\%,
20\%, 10\%, 5\%, 0\% difference).

For each level of difference, record the maximum likelihood value from
the profile. You might create a simple graph (by hand or in Excel)
showing how prediction confidence changes as compositional difference
decreases.

\textbf{Key questions to explore:} At what compositional difference does
prediction confidence drop below 50\%? Is there a threshold where the
model essentially ``gives up'' and spreads probability across many
positions?

\begin{center}\rule{0.5\linewidth}{0.5pt}\end{center}

\textbf{For your chosen exploration, document:}

\begin{itemize}
\tightlist
\item
  Which option you selected and why
\item
  Specific parameters or sequences you tested (at least 3-4 different
  conditions)
\item
  What happened in each test (predicted positions and approximate
  likelihoods)
\item
  Patterns or trends you observed
\item
  Your interpretation: Why did the model behave this way?
\end{itemize}

\section{Reflection Questions}\label{reflection-questions-13}

Now that you've completed the lab, reflect on the following questions.
Your answers should demonstrate understanding of HMM principles and
their application to gene prediction.

\begin{enumerate}
\def\labelenumi{\arabic{enumi}.}
\item
  \textbf{Parameter sensitivity}: Based on your systematic explorations
  in Part 2 (transition probabilities and emission probabilities), which
  type of parameter seems to have the strongest effect on splice site
  prediction accuracy? Explain your reasoning using specific examples
  from your completed tables. How does this relate to the way real gene
  prediction software is trained?
\item
  \textbf{Model complexity trade-offs}: The toy model in this lab uses
  just three states (Exon, Splice, Intron) and models the splice site as
  a single nucleotide with strong G preference. Real gene finders like
  AUGUSTUS use dozens of states and model splice sites with
  position-specific weight matrices spanning multiple nucleotides. What
  are the advantages and disadvantages of increasing model complexity?
  Consider factors like:

  \begin{itemize}
  \tightlist
  \item
    Prediction accuracy
  \item
    Amount of training data required
  \item
    Computational cost
  \item
    Risk of overfitting
  \item
    Interpretability
  \end{itemize}
\item
  \textbf{Integration with other data}: In question 6 of Part 2, you
  considered how RNA-Seq data could help resolve ambiguous splice site
  predictions. Describe in detail how you would integrate RNA-Seq
  evidence with HMM-based predictions. What would you look for in the
  RNA-Seq data? How would you handle cases where they disagree?
\item
  \textbf{Simple versus Full Model}: Compare your experience working
  with the Simple Model (6 bp, Part 2) versus the Full Model (26 bp,
  Part 4). What additional insights did the longer sequence provide?
  What made the Full Model more challenging or interesting to work with?
  How does the increased sequence length affect the clarity of
  predictions?
\item
  \textbf{Beyond gene prediction}: HMMs are used throughout
  bioinformatics for diverse applications including protein domain
  identification (profile HMMs), sequence alignment (pair HMMs), and
  ancestry inference (as discussed in
  \hyperref[hidden-markov-models-in-biology]{Chapter 9}). Choose one
  application of HMMs from the chapter that is NOT gene prediction and
  explain:

  \begin{itemize}
  \tightlist
  \item
    What are the ``states'' in that application?
  \item
    What are the ``observations'' (emissions)?
  \item
    Why is an HMM appropriate for that problem?
  \item
    How is it similar to or different from the gene prediction HMM you
    worked with in this lab?
  \end{itemize}
\end{enumerate}

\section{Deliverables}\label{deliverables-2}

Upload the following to Canvas:

\begin{enumerate}
\def\labelenumi{\arabic{enumi}.}
\tightlist
\item
  \textbf{Part 1} (can be photo/scan of handwritten work or typed):

  \begin{itemize}
  \tightlist
  \item
    Your hand-drawn state diagram for a complete gene (with start codon,
    stop codon, exons, introns, 5' SS, and 3' SS)
  \item
    Calculations for the three state paths (showing your work for at
    least one path in detail)
  \item
    Answers to the follow-up questions about which path is most likely
    and what seems problematic
  \end{itemize}
\item
  \textbf{Part 2 Completed Tables}:

  \begin{itemize}
  \tightlist
  \item
    Completed table from question 4 (transition probability
    perturbations) with all cells filled in
  \item
    Completed table from question 5 (emission probability perturbations)
    with all cells filled in
  \item
    Brief description of the key patterns you observed across these
    systematic perturbations (2-3 sentences)
  \end{itemize}
\item
  \textbf{Part 2 Short Answer Questions}:

  \begin{itemize}
  \tightlist
  \item
    Answers to questions 3, 6, and 7 from Part 2
  \end{itemize}
\item
  \textbf{Part 3 Exploration Summary} (2-3 sentences):

  \begin{itemize}
  \tightlist
  \item
    For your chosen exploration of the simple model:

    \begin{itemize}
    \tightlist
    \item
      Description of parameter changes you made (be specific about
      values)
    \item
      Table or organized list of results showing how predictions changed
    \item
      Interpretation of findings (why did the model behave this way?)
    \end{itemize}
  \item
    Include screenshots or sketches if they help illustrate your results
  \end{itemize}
\item
  \textbf{Part 4 Exploration Summary} (1 paragraph,
  \textasciitilde150-250 words):

  \begin{itemize}
  \tightlist
  \item
    State which exploration option you chose
  \item
    Briefly describe what you tested
  \item
    Summarize your key findings
  \item
    Include one insight about what this taught you regarding HMM-based
    gene prediction
  \end{itemize}
\item
  \textbf{Reflection Questions}: Answers to all five reflection
  questions above (in a Word document or PDF).
\end{enumerate}

\begin{center}\rule{0.5\linewidth}{0.5pt}\end{center}

Congratulations! The goal of this lab is to give you real-world
experience with an HMM and to help you understand what properties lead
to variation across software tools in gene model predictions.

\section{Acknowledgement}\label{acknowledgement-1}

This lab exercise was adapted from one written by Dr.~Anton Weisstein
from Truman State University and Zane Goodwin from Washington University
in St.~Louis in 2013\textsuperscript{148}.

\section{References}\label{references-14}

\chapter{Lab07: Variant Calling via GATK Best
Practices}\label{lab07-variant-calling-via-gatk-best-practices}

\section{Overview and Learning
Objectives}\label{overview-and-learning-objectives}

The primary goal of most NGS workflows is to generate high-quality
genetic variants for downstream biological analysis. However, as you
learned in \textbf{Chapter 10}, raw NGS data is inherently noisy,
containing sequencing errors, PCR duplicates, and alignment artifacts
that can generate thousands of false-positive variant calls. Without
careful post-processing and filtering, these technical artifacts can
lead to incorrect biological inferences---for example, mistakenly
identifying a sequencing error as a true adaptive mutation in an
experimental evolution study.

This lab walks you through the \textbf{GATK Best Practices workflow}
(see \textbf{Figure 10.1} in Chapter 10) using the \emph{Saccharomyces
cerevisiae} experimental evolution dataset introduced in
\textbf{Appendix C} and first analyzed in \textbf{Lab 4}. You will
complete three major steps:

\begin{enumerate}
\def\labelenumi{\arabic{enumi}.}
\tightlist
\item
  \textbf{Post-processing and variant calling}: Mark PCR duplicates,
  call variants using GATK HaplotypeCaller, and assess alignment
  quality.
\item
  \textbf{Hard filtering and quality assessment}: Apply GATK's
  recommended filter thresholds, visualize variant quality
  distributions, and evaluate whether standard cutoffs are appropriate
  for your yeast dataset.
\item
  \textbf{Iterative refinement and comparison}: Adjust filter parameters
  based on quality plots, generate a second filtered callset, and
  compare all three VCFs (raw, GATK-filtered, custom-filtered) using an
  UpSet plot to understand how filtering decisions affect the final
  variant set.
\end{enumerate}

By the end of this lab, you will understand why \textbf{filtering is
iterative and context-dependent}, and you will have hands-on experience
making the kinds of decisions that real genomics researchers face when
balancing sensitivity (retaining true variants) against specificity
(excluding false positives).

\section{Step 1: Post-processing and variant
calling}\label{step-1-post-processing-and-variant-calling}

\subsection{Goal}\label{goal}

Learn how to complete post-processing (duplicate marking), run quality
control checks, and call variants using GATK HaplotypeCaller. This step
corresponds to Figure~\ref{fig-gatk-workflow2} (boxes for ``Mark
Duplicates'' and ``GATK HaplotypeCaller'') and the concepts in
\textbf{Chapter 10, Sections 10.3--10.4}.

\begin{enumerate}
\def\labelenumi{\arabic{enumi}.}
\item
  First, return to the directory in your ASC home directory titled
  \texttt{Lab4} that you made previously. Let's continue working here
  since all the files we need are in here.
\item
  Copy the contents of the folder
  \texttt{\textasciitilde{}/bioinf\_shared/Lab7\_VariantCalling/} to
  your working directory. Examine carefully the script
  ``5\_GATK\_Variant\_Calling.sh'' to walk through what it is doing and
  the tools it uses (refer to Figure~\ref{fig-gatk-workflow2} as you
  read).
\item
  Using your favorite text editor, open the script
  ``5\_GATK\_Variant\_Calling.sh''. Edit this script in the following
  ways:
\end{enumerate}

\begin{itemize}
\tightlist
\item
  Change the path to your current working directory
\item
  Add a QC step to the script like you did in Lab 4 to run
  \texttt{samtools\ flagstat} on the ``*.markdup.bam'' output. This is
  the last bam formatted file in your pipeline.
\end{itemize}

\begin{quote}
Make sure to include the loading of the samtools module in your script
(see script \#4 as an example), as well as redirect the output to a
file.
\end{quote}

\begin{enumerate}
\def\labelenumi{\arabic{enumi}.}
\setcounter{enumi}{3}
\tightlist
\item
  Submit this script to the ASC class queue using the following
  parameters:
\end{enumerate}

\begin{itemize}
\tightlist
\item
  Number of CPUs = 1
\item
  RAM/Memory = 2GB
\item
  Queue = class
\item
  Time = default (press enter to keep default time)
\end{itemize}

\begin{enumerate}
\def\labelenumi{\arabic{enumi}.}
\setcounter{enumi}{4}
\tightlist
\item
  Compare the flagstat output to the one
  \href{https://github.com/AUBioInformatics26/BamQC/blob/main/flagstat.md}{here}.
\end{enumerate}

\begin{quote}
Note: if this link does not work for you, make sure you are logged into
GitHub. This link goes to a private repository in the GitHub
organization for our class this semester. If you have not yet joined,
please refer back to an earlier announcement from the LA Chipper Johnson
on this topic.
\end{quote}

\section{Step 2: Variant Filtering}\label{step-2-variant-filtering}

\subsection{Goal}\label{goal-1}

Learn how to filter variants using GATK's \textbf{VariantFiltration}
tool, visualize the distributions of key quality metrics, and evaluate
whether the standard GATK-recommended thresholds are appropriate for
your yeast dataset. This step corresponds to \textbf{Section 10.5 (Hard
Filtering)} and builds your intuition for \textbf{why filtering is
iterative and data-dependent}.

\subsection{Why Filter Variants?}\label{why-filter-variants}

The raw VCF from HaplotypeCaller (\texttt{SRR1693691.SNPs.vcf}) contains
thousands of putative SNPs, but many are false positives caused by
sequencing errors, mapping artifacts, or systematic biases (e.g., strand
bias, where all variant-supporting reads come from one DNA strand).
\textbf{Section 10.5} describes the key quality metrics used for
filtering.

GATK provides \textbf{recommended thresholds} for these metrics (based
on human whole-genome sequencing), but these may not be optimal for
yeast or other organisms. Your task in this step is to apply the
recommended filters, examine the distributions of quality scores, and
decide whether adjustments are needed.

\begin{enumerate}
\def\labelenumi{\arabic{enumi}.}
\item
  Examine carefully the script ``6\_GATK\_Variant\_Filtering.sh'' to
  walk through what it is doing and the tools it uses. This script will
  perform hard filtering based on the filtering criteria recommended by
  GATK. It then runs a custom Perl script to produce density plots of
  the various parameters that are output as PDFs.
\item
  Using your favorite text editor, open the script
  ``6\_GATK\_Variant\_Filtering.sh''. Edit this script in the following
  ways:
\end{enumerate}

\begin{itemize}
\tightlist
\item
  Change the path to your current working directory
\item
  Copy the vcftools command to get depth statistics on the filtered VCF
  file and run it also on the unfiltered file ``SRR1693691.SNPs.vcf'' so
  you can compare later.
\item
  Add another QC step to the script. Here use the tool
  \href{https://hpcdocs.asc.edu/search/node/multiqc}{multiqc} to point
  to your directory to generate a comprehensive summary of ALL the QC
  files you have made along the way:

  \begin{itemize}
  \tightlist
  \item
    \texttt{module\ load\ multiqc}
  \item
    \texttt{multiqc\ \$\{path\}/*}
  \end{itemize}
\end{itemize}

\begin{quote}
MultiQC will scan your directory for FastQC reports, flagstat outputs,
GATK logs, and VCFtools summaries, and compile them into a single
interactive HTML report (\texttt{multiqc\_report.html}). This gives you
a bird's-eye view of the entire Lab 4 → Lab 7 workflow.
\end{quote}

\begin{enumerate}
\def\labelenumi{\arabic{enumi}.}
\setcounter{enumi}{2}
\tightlist
\item
  Submit this script to the ASC class queue using the following
  parameters:
\end{enumerate}

\begin{itemize}
\tightlist
\item
  Number of CPUs = 1
\item
  RAM/Memory = 2GB
\item
  Queue = class
\item
  Time = default (press enter to keep default time)
\end{itemize}

\begin{enumerate}
\def\labelenumi{\arabic{enumi}.}
\setcounter{enumi}{3}
\item
  After inspecting the standard out job file to make sure all ran OK,
  download the Html multiqc report and the resulting PDFs to your local
  computer. Examine the multiqc report. Fill in the table below based on
  what you see with both the result as well as which tool was used to
  get this information:

  Fill in the table below based on the MultiQC html report downloaded
  from the ASC:

  \begin{longtable}[]{@{}lll@{}}
  \toprule\noalign{}
  \textbf{QC Metric} & \textbf{Result} & \textbf{Tool Used} \\
  \midrule\noalign{}
  \endhead
  \bottomrule\noalign{}
  \endlastfoot
  \% reads mapped to reference & & \\
  \% duplicate reads & & \\
  \end{longtable}

  \begin{tcolorbox}[enhanced jigsaw, rightrule=.15mm, leftrule=.75mm, colbacktitle=quarto-callout-note-color!10!white, colframe=quarto-callout-note-color-frame, toprule=.15mm, opacitybacktitle=0.6, title=\textcolor{quarto-callout-note-color}{\faInfo}\hspace{0.5em}{Finding metrics in MultiQC}, left=2mm, colback=white, bottomtitle=1mm, toptitle=1mm, bottomrule=.15mm, coltitle=black, titlerule=0mm, breakable, opacityback=0, arc=.35mm]

  \begin{itemize}
  \tightlist
  \item
    \textbf{\% mapped}: Look under the ``SAMtools flagstat'' section.
    The report will show a bar chart for each BAM file with the
    percentage of mapped reads.
  \item
    \textbf{\% duplicates}: Also in the SAMtools flagstat section. This
    is the fraction of reads marked as duplicates by GATK
    MarkDuplicates.
  \end{itemize}

  Typical values for this yeast dataset: \textbf{\textasciitilde99\%
  mapped}, \textbf{\textasciitilde10--15\% duplicates}.

  \end{tcolorbox}
\item
  Next, compare the number of raw SNPs to the number of filtered SNPs to
  see how stringent you are currently:

  Fill in the table below:

  \begin{longtable}[]{@{}lll@{}}
  \toprule\noalign{}
  \textbf{File} & \textbf{\# SNPs} & \textbf{Mean Depth} \\
  \midrule\noalign{}
  \endhead
  \bottomrule\noalign{}
  \endlastfoot
  SRR1693691.SNPs.vcf (raw) & & \\
  SRR1693691.SNPs.filtered.vcf (GATK filters) & & \\
  \end{longtable}

  \textbf{Interpretation}: If filtering removes 30--50\% of raw
  variants, this is typical---many low-quality sites are being excluded.
  If filtering removes \textless10\%, the thresholds may be too
  permissive. If filtering removes \textgreater70\%, the thresholds may
  be too strict, and you risk losing true variants.
\item
  Transfer all PDF files (named like \texttt{QD\_density.pdf},
  \texttt{FS\_density.pdf}, etc.) to your computer. Open each plot and
  inspect the distribution of quality scores. For each metric, ask
  yourself:

  \begin{itemize}
  \tightlist
  \item
    \textbf{Where does the GATK threshold fall?} (The script annotates
    the threshold as a vertical red line on each plot.)
  \item
    \textbf{Do most of my variants pass or fail?} If 90\% of variants
    are above the threshold, the filter is not very stringent. If 50\%
    fail, the filter is aggressive.
  \item
    \textbf{Does the distribution have outliers?} A few extreme values
    might be real rare variants or artifacts---context matters.
  \end{itemize}
\end{enumerate}

Based on your understanding of these density plots, decide whether to
adjust the filtering or not. See Section 10.4.2 for a description of
each of these criteria and the GATK recommended cutoffs.

\section{Step 3: Refilter and
Compare}\label{step-3-refilter-and-compare}

\subsection{Goal}\label{goal-2}

Apply your adjusted filter criteria, generate a second filtered callset,
and compare all three VCFs (raw, GATK-filtered, custom-filtered) using
an \textbf{UpSet plot} to visualize the overlaps and differences between
filtering strategies. This step reinforces the key lesson of
\textbf{Section 10.5}: filtering is iterative, context-dependent, and
involves trade-offs.

\subsection{Why Compare Multiple Filtered
Callsets?}\label{why-compare-multiple-filtered-callsets}

In real research, you rarely get filtering right on the first try. By
generating multiple filtered VCFs and comparing them, you can see:

\begin{itemize}
\tightlist
\item
  \textbf{Which variants are robust} (present in all callsets,
  regardless of filter stringency)
\item
  \textbf{Which variants are borderline} (present in permissive filters
  but not strict ones)
\item
  \textbf{How much overlap exists} between GATK's recommended filters
  and your custom filters
\end{itemize}

UpSet plots (introduced in \textbf{Chapter 10} as a superior alternative
to Venn diagrams) show these relationships clearly, especially when
comparing more than two sets.

\subsection{Instructions}\label{instructions}

\begin{enumerate}
\def\labelenumi{\arabic{enumi}.}
\item
  \textbf{Decide on your adjusted filter criteria}:

  Based on the density plots and tables from Step 2, choose new
  thresholds for \textbf{one or two metrics}. For example:

  \begin{itemize}
  \tightlist
  \item
    If the GATK QD threshold (QD \textgreater{} 2.0) looks too
    permissive (most variants have QD \textgreater{} 10), raise it to QD
    \textgreater{} 5.0.
  \item
    If the FS threshold (FS \textless{} 60.0) excludes very few
    variants, tighten it to FS \textless{} 40.0.
  \end{itemize}

  Write down your new thresholds---you will need them in the next step.
\item
  \textbf{Edit the filtering script to create a second filtered VCF}:

  Open \texttt{6\_GATK\_Variant\_Filtering.sh} again and find the
  \texttt{VariantFiltration} command. \textbf{Copy the entire command}
  and paste it below the original, then modify:

  \begin{itemize}
  \tightlist
  \item
    Change the output filename to \texttt{SRR1693691.SNPs.filtered2.vcf}
    (so it doesn't overwrite the first filtered VCF).
  \item
    Adjust the filter expressions (\texttt{-\/-filter-expression} and
    \texttt{-\/-filter-name}) to reflect your new thresholds.
  \end{itemize}
\item
  Run the edited script on the ASC. See the suggested queue parameters
  from the first script and make adjustments based on that job run time
  and usage.
\item
  Once you are done, you should now have 3 VCF files:
\end{enumerate}

\begin{itemize}
\tightlist
\item
  Raw unfiltered VCF: SRR1693691.SNPs.vcf.gz
\item
  GATK filtered VCF: SRR1693691.SNPs.filtered.vcf
\item
  Your adjusted filtering VCF: SRR1693691.SNPs.filtered2.vcf
\end{itemize}

Now that you have 3 VCF files, you can compare the sets of variants
between the three files using \texttt{vcf-compare}, a tool within
VCFtools. This will generate a text file that can be used to make VENN
diagrams. But of course, we \textbf{\emph{HATE}} VENN diagrams, so
instead we will make an \href{https://caleydo.org/tools/upset/}{UpSet}
plot using the R package
\href{http://gehlenborglab.org/research/projects/upsetr/}{UpSetR}.

\begin{enumerate}
\def\labelenumi{\arabic{enumi}.}
\setcounter{enumi}{4}
\item
  Using \texttt{less}, examine the script \texttt{UpSetR.sh}. This
  script will first use VCFtools to remove all of the sites that were
  filtered out (GATK only marks them, but does not remove them). It then
  properly compresses each VCF and indexed them for compatability with
  \texttt{vcf-compare}. These two steps use the tool \texttt{tabix}
  which is one of the modules loaded at the beginning of the script.
  Submit this script to the ASC queue.
\item
  You will notice that the end of this script uses the tools grep and
  cut to simplify the output from vcf-compare. This is printed into the
  file ``intersection.4upsetR.venn''. Examine that file.
\item
  \textbf{Generate an UpSet plot in R} (on your local computer):

  \begin{enumerate}
  \def\labelenumii{\alph{enumii}.}
  \item
    Download the R script
    \href{datasets/UpSetR_plot.R}{\texttt{UpSetR.R}}.
  \item
    Open RStudio on your computer and load the script.
  \item
    Install the \texttt{UpSetR} package if needed:
  \end{enumerate}

\begin{Shaded}
\begin{Highlighting}[]
\FunctionTok{install.packages}\NormalTok{(}\StringTok{"UpSetR"}\NormalTok{)}
\end{Highlighting}
\end{Shaded}

  \begin{enumerate}
  \def\labelenumii{\alph{enumii}.}
  \setcounter{enumii}{3}
  \item
    Copy the counts from \texttt{intersection.4upsetR.venn} into the
    appropriate lines of the R script (the script has placeholder
    values---replace them with your actual counts).
  \item
    Run the script to generate an UpSet plot. The plot shows:
  \end{enumerate}

  \begin{itemize}
  \tightlist
  \item
    \textbf{Vertical bars}: Size of each intersection (e.g., how many
    variants are in raw only, how many are in all three sets)
  \item
    \textbf{Connected dots below}: Which sets are included in each
    intersection
  \end{itemize}

  \begin{enumerate}
  \def\labelenumii{\alph{enumii}.}
  \setcounter{enumii}{5}
  \tightlist
  \item
    Export the plot as a PDF or PNG for your submission.
  \end{enumerate}
\item
  Interpret the Upset Plot and answer the following reflection
  questions:
\end{enumerate}

\begin{itemize}
\item
  Depending on how you adjusted the filtering, which two files have the
  MOST overlap?
\item
  Which two files have the least overlap?
\item
  Does this make sense based on the differences in the stringency of
  these two sets of filtering parameters you used?
\end{itemize}

\begin{tcolorbox}[enhanced jigsaw, rightrule=.15mm, leftrule=.75mm, colbacktitle=quarto-callout-tip-color!10!white, colframe=quarto-callout-tip-color-frame, toprule=.15mm, opacitybacktitle=0.6, title=\textcolor{quarto-callout-tip-color}{\faLightbulb}\hspace{0.5em}{What if my script fails?}, left=2mm, colback=white, bottomtitle=1mm, toptitle=1mm, bottomrule=.15mm, coltitle=black, titlerule=0mm, breakable, opacityback=0, arc=.35mm]

Common issues:

\begin{itemize}
\tightlist
\item
  \textbf{Path errors}: Make sure the \texttt{path} variable points to
  your actual working directory and that all input files (sorted BAM,
  reference genome FASTA and indices) exist.
\item
  \textbf{Module loading errors}: The script loads GATK and SAMtools
  modules---if you get ``module not found,'' check that the module names
  are correct for the ASC.
\item
  \textbf{Memory errors}: If the job is killed for exceeding memory,
  resubmit with 4GB instead of 2GB.
\end{itemize}

Check the \texttt{.e} error file for detailed messages.

\end{tcolorbox}

\begin{center}\rule{0.5\linewidth}{0.5pt}\end{center}

\section{Deliverables and Submission}\label{deliverables-and-submission}

Upload the following files to Canvas for grading:

\begin{enumerate}
\def\labelenumi{\arabic{enumi}.}
\tightlist
\item
  \textbf{PDF density plots} from Step 2 (QD, FS, MQ, SOR, MQRankSum,
  ReadPosRankSum)
\item
  \textbf{MultiQC HTML report} and the \texttt{multiqc\_data/} folder
\item
  \textbf{Filled-in tables} from Step 2:

  \begin{itemize}
  \tightlist
  \item
    QC Metrics table (\% mapped, \% duplicates)
  \item
    SNP count and depth table (raw vs.~filtered)
  \end{itemize}
\item
  \textbf{UpSet plot} (PDF or PNG) from Step 3
\item
  \textbf{Reflection answers} for Step 3 (see questions above)
\end{enumerate}

\begin{center}\rule{0.5\linewidth}{0.5pt}\end{center}

\section{Congratulations!}\label{congratulations}

You have successfully completed the \textbf{GATK Best Practices
workflow} for variant calling and filtering. This lab provides a
real-world experience of the challenges of variant filtering: being too
permissive can lead to false positives (mistaking sequencing errors for
true variants), while being too stringent can exclude real biological
variation. The best practice is to \textbf{experiment with filter
thresholds, examine quality distributions, and err on the side of
caution}---especially for downstream applications where false positives
are costly (e.g., functional validation experiments).

\subsection{Alternative Approaches to
Filtering}\label{alternative-approaches-to-filtering}

As mentioned in \textbf{Section 10.5} and in class, hard filtering is
just one approach to variant quality control. Other methods include:

\begin{itemize}
\item
  \textbf{Variant Quality Score Recalibration (VQSR)}: A
  machine-learning approach that uses a set of known high-quality
  variants to train a model that distinguishes true variants from
  artifacts. VQSR is the gold standard for human genomics but requires a
  large database of validated variants, making it impractical for
  non-model organisms like yeast. It is also computationally intensive,
  so we do not use it in this course.
\item
  \textbf{Ensemble calling}: Running multiple variant callers (e.g.,
  GATK, FreeBayes, bcftools) on the same data and taking the
  intersection of callsets. Variants called by multiple independent
  algorithms are more likely to be real. You will explore this approach
  in \textbf{Lab 8} when you read benchmarking papers that compare
  different callers.
\item
  \textbf{Empirical validation}: For high-stakes projects (e.g.,
  clinical diagnostics, identifying causal mutations), researchers
  validate a subset of variants using orthogonal methods such as Sanger
  sequencing, high-coverage targeted resequencing, or alternative
  sequencing platforms (e.g., PacBio long reads). If you identify
  100,000 SNPs and want to validate 1\%, that requires validating 1,000
  sites---costly and time-intensive, but sometimes necessary.
\end{itemize}

For most research applications, \textbf{hard filtering combined with
visual inspection in IGV} (which you will do in Lab 9) provides a good
balance between rigor and practicality.

\subsection{Looking Ahead to Lab 9}\label{looking-ahead-to-lab-9}

In the next lab, you will load the filtered VCFs you generated here into
\textbf{IGV (Integrative Genomics Viewer)} along with the aligned BAM
files and the yeast reference genome. You will navigate to specific
SNPs, zoom in to see individual read alignments, and evaluate the
evidence for each variant.

\chapter{Lab08: Beyond GATK -- Benchmarking Variant Calling
Pipelines}\label{lab08-beyond-gatk-benchmarking-variant-calling-pipelines}

\section{Overview}\label{overview-2}

Lab 8 is an \textbf{in-class discussion and synthesis activity} focused
on recent benchmarking studies that compare variant calling pipelines
across different sequencing platforms, organisms, and data types. Over
the past decade, the genomics community has moved beyond asking ``which
tool is best?'' to understanding \emph{when} and \emph{why} different
callers excel, and how upstream choices (alignment, preprocessing,
population ancestry) shape final variant quality. In this lab, you will
critically evaluate five recent papers (2021--2026) that benchmark
variant callers on diverse datasets, extract practical recommendations,
and synthesize guidelines for choosing pipelines in your own
research\textsuperscript{5--9}.

Each student has been assigned one paper based on the first three
letters of their last name (see assignments below). You are expected to
read your assigned paper carefully before class, guided by the questions
provided, and come prepared to contribute to small-group and whole-class
discussions. During class, you will work with peers who read the same
paper to fill out a structured summary table, present key findings to
the class, and collectively build a decision framework for variant
calling that addresses common questions facing genomics researchers.

\begin{tcolorbox}[enhanced jigsaw, rightrule=.15mm, leftrule=.75mm, colbacktitle=quarto-callout-important-color!10!white, colframe=quarto-callout-important-color-frame, toprule=.15mm, opacitybacktitle=0.6, title=\textcolor{quarto-callout-important-color}{\faExclamation}\hspace{0.5em}{Attendance and Participation Required}, left=2mm, colback=white, bottomtitle=1mm, toptitle=1mm, bottomrule=.15mm, coltitle=black, titlerule=0mm, breakable, opacityback=0, arc=.35mm]

The entire Lab 8 assignment takes place during the scheduled lab period.
Your grade is based on active participation in group discussions and
submission of the completed class summary tables (photographed and
uploaded to Canvas before leaving). If you cannot attend, you must
notify the instructor in advance to arrange an alternative assignment.

\end{tcolorbox}

\section{Paper Assignment by Last
Name}\label{paper-assignment-by-last-name}

Each student is responsible for reading ONE of the five papers for the
Lab 8 discussion. Find your assigned paper based on the \textbf{first
three letters of your last name} as listed in Canvas:

\begin{longtable}[]{@{}
  >{\raggedright\arraybackslash}p{(\linewidth - 6\tabcolsep) * \real{0.3137}}
  >{\raggedright\arraybackslash}p{(\linewidth - 6\tabcolsep) * \real{0.3137}}
  >{\raggedright\arraybackslash}p{(\linewidth - 6\tabcolsep) * \real{0.3137}}
  >{\raggedright\arraybackslash}p{(\linewidth - 6\tabcolsep) * \real{0.0588}}@{}}
\toprule\noalign{}
\begin{minipage}[b]{\linewidth}\raggedright
Last Name Range
\end{minipage} & \begin{minipage}[b]{\linewidth}\raggedright
Assigned Paper
\end{minipage} & \begin{minipage}[b]{\linewidth}\raggedright
Reference Citation
\end{minipage} & \begin{minipage}[b]{\linewidth}\raggedright
Number of Pages
\end{minipage} \\
\midrule\noalign{}
\endhead
\bottomrule\noalign{}
\endlastfoot
\textbf{A -- Da} & Tan et al.~2025, \emph{Briefings in Bioinformatics} &
\textsuperscript{5} & 11 \\
\textbf{Db -- L} & Abdelwahab et al.~2023, \emph{BMC Bioinformatics} &
\textsuperscript{6} & 13 \\
\textbf{M -- Pa} & Betschart et al.~2022, \emph{Scientific Reports} &
\textsuperscript{7} & 11 \\
\textbf{Pb -- Pz} & Pinto et al.~2026, \emph{PLoS ONE} &
\textsuperscript{8} & 22 \\
\textbf{Q -- Z} & Guille et al.~2024, \emph{Briefings in Bioinformatics}
& \textsuperscript{9} & 11 \\
\end{longtable}

\section{Pre-Class Preparation: Guiding
Questions}\label{pre-class-preparation-guiding-questions}

As you read your assigned paper, take notes on the following questions.
You will use these notes during the in-class small-group synthesis.

\subsection{Study Design and Scope}\label{study-design-and-scope}

\begin{enumerate}
\def\labelenumi{\arabic{enumi}.}
\tightlist
\item
  What is the main goal of the benchmarking study? (e.g., comparing
  germline vs.~somatic callers, evaluating AI-based tools, assessing
  performance on non-model organisms)
\item
  Which datasets were used? (e.g., Genome in a Bottle reference samples,
  monozygotic twins, cancer exomes, specific organisms)
\item
  Which sequencing platforms and read lengths were analyzed? (Illumina
  short reads, PacBio HiFi, Oxford Nanopore, or a mix)
\end{enumerate}

\subsection{Tools and Pipelines}\label{tools-and-pipelines}

\begin{enumerate}
\def\labelenumi{\arabic{enumi}.}
\setcounter{enumi}{3}
\tightlist
\item
  Which \textbf{aligners} were tested (if applicable)? (e.g., BWA-MEM,
  Bowtie2, minimap2)
\item
  Which \textbf{variant callers} were compared? (e.g., GATK
  HaplotypeCaller, bcftools, FreeBayes, Platypus, DeepVariant, others)
\item
  Did the study evaluate \textbf{ensemble or voting-based methods}? If
  so, how were results combined?
\end{enumerate}

\subsection{Performance Metrics and
Results}\label{performance-metrics-and-results}

\begin{enumerate}
\def\labelenumi{\arabic{enumi}.}
\setcounter{enumi}{6}
\tightlist
\item
  Which metrics were used to assess performance? (e.g., precision,
  recall, F1-score, concordance with truth sets, computational time,
  memory usage)
\item
  Which tools performed best for \textbf{SNPs}? Which for
  \textbf{indels}?
\item
  Were there differences in performance by \textbf{genomic context}
  (e.g., exonic vs.~intergenic, high-GC regions, repetitive regions)?
\item
  Did computational cost (runtime, memory) play a role in the
  recommendations?
\end{enumerate}

\subsection{Key Findings and
Recommendations}\label{key-findings-and-recommendations}

\begin{enumerate}
\def\labelenumi{\arabic{enumi}.}
\setcounter{enumi}{10}
\tightlist
\item
  What are the paper's main recommendations? (e.g., ``use GATK for
  high-coverage human genomes,'' ``ensemble methods reduce false
  positives in somatic calling'')
\item
  Did the authors identify biases or limitations? (e.g., training on
  European-ancestry genomes, poor performance in non-model systems,
  platform-specific errors)
\item
  Are the recommendations generalizable to other organisms, coverages,
  or research questions? Why or why not?
\end{enumerate}

\begin{tcolorbox}[enhanced jigsaw, rightrule=.15mm, leftrule=.75mm, colbacktitle=quarto-callout-tip-color!10!white, colframe=quarto-callout-tip-color-frame, toprule=.15mm, opacitybacktitle=0.6, title=\textcolor{quarto-callout-tip-color}{\faLightbulb}\hspace{0.5em}{Tips for Success}, left=2mm, colback=white, bottomtitle=1mm, toptitle=1mm, bottomrule=.15mm, coltitle=black, titlerule=0mm, breakable, opacityback=0, arc=.35mm]

\begin{itemize}
\tightlist
\item
  \textbf{Read your paper thoroughly} before class. Annotate it and
  bring notes---you cannot participate meaningfully without preparation.
\item
  \textbf{Be ready to explain, not just summarize.} Your classmates will
  rely on you to interpret your paper's findings.
\item
  \textbf{Ask questions during presentations.} Benchmarking studies
  often use different metrics or datasets---clarify what makes results
  comparable.
\item
  \textbf{Think critically about generalizability.} A tool that works
  well for human exomes may fail for non-model organisms or low-coverage
  data.
\end{itemize}

\end{tcolorbox}

\section{In-Class Activity Structure (75
minutes)}\label{in-class-activity-structure-75-minutes}

\subsection{Part 1: Small-Group Synthesis (20
minutes)}\label{part-1-small-group-synthesis-20-minutes}

Meet with others who read the same paper. As a group, discuss your
answers to the guiding questions and work together to fill out the
\textbf{Paper Summary Table} for your study (template below). One member
of your group should cast their screen so the table can be filled and
photographed later. Focus on extracting key details that will help the
class compare across papers:

\textbf{Paper Summary Table Template:}

\begin{longtable}[]{@{}
  >{\raggedright\arraybackslash}p{(\linewidth - 4\tabcolsep) * \real{0.3500}}
  >{\raggedright\arraybackslash}p{(\linewidth - 4\tabcolsep) * \real{0.3250}}
  >{\raggedright\arraybackslash}p{(\linewidth - 4\tabcolsep) * \real{0.3250}}@{}}
\toprule\noalign{}
\begin{minipage}[b]{\linewidth}\raggedright
\textbf{Category}
\end{minipage} & \begin{minipage}[b]{\linewidth}\raggedright
\textbf{Details}
\end{minipage} & \begin{minipage}[b]{\linewidth}\raggedright
\textbf{Your Paper}
\end{minipage} \\
\midrule\noalign{}
\endhead
\bottomrule\noalign{}
\endlastfoot
\textbf{Study Citation} & Author, Year, Journal & \\
\textbf{Study Goal} & Brief statement of main objective & \\
\textbf{Datasets Used} & Organism(s), reference samples, sample size
& \\
\textbf{Read Length(s)} & e.g., 150 bp paired-end, 15 kb HiFi & \\
\textbf{Sequencing Platform(s)} & Illumina, PacBio, ONT, etc. & \\
\textbf{Tools/Pipelines Tested} & Aligner(s) + Caller(s) & \\
\textbf{Variant Callers Compared} & List key callers (GATK, bcftools,
FreeBayes, DeepVariant, etc.) & \\
\textbf{Ensemble Strategy (if any)} & Voting, intersection, machine
learning meta-classifier & \\
\textbf{Performance Metrics} & Precision, recall, F1, concordance,
runtime, etc. & \\
\textbf{Best Tool(s) for SNPs} & Tool name and brief reason & \\
\textbf{Best Tool(s) for Indels} & Tool name and brief reason & \\
\textbf{Key Recommendation} & One-sentence practical takeaway & \\
\textbf{Limitations/Caveats} & What contexts or questions are NOT
covered? & \\
\end{longtable}

Your group will present this table to the class in Part 2.

\subsection{Part 2: Group Presentations (15
minutes)}\label{part-2-group-presentations-15-minutes}

Each paper group gives a \textbf{3-minute presentation} summarizing
their study. Focus on:

\begin{itemize}
\tightlist
\item
  Study goal and dataset (30 seconds)
\item
  Tools compared and metrics used (1 minute)
\item
  Key findings: which tools performed best for SNPs and indels, and
  under what conditions (1 minute)
\item
  Main recommendation and limitations (30 seconds)
\end{itemize}

While other groups present, \textbf{take notes} on their summary
tables---you will need this information for Part 3.

\subsection{Part 3: Synthesis Across Papers (20
minutes)}\label{part-3-synthesis-across-papers-20-minutes}

Now that all five papers have been presented, work as a class to fill
out the \textbf{Cross-Study Synthesis Table} on the board. The
instructor will lead a discussion to populate this table by extracting
common themes and differences across the five benchmarking studies.

\textbf{Cross-Study Synthesis Table Template:}

\begin{longtable}[]{@{}
  >{\raggedright\arraybackslash}p{(\linewidth - 10\tabcolsep) * \real{0.0848}}
  >{\raggedright\arraybackslash}p{(\linewidth - 10\tabcolsep) * \real{0.1636}}
  >{\raggedright\arraybackslash}p{(\linewidth - 10\tabcolsep) * \real{0.2061}}
  >{\raggedright\arraybackslash}p{(\linewidth - 10\tabcolsep) * \real{0.1939}}
  >{\raggedright\arraybackslash}p{(\linewidth - 10\tabcolsep) * \real{0.1697}}
  >{\raggedright\arraybackslash}p{(\linewidth - 10\tabcolsep) * \real{0.1818}}@{}}
\toprule\noalign{}
\begin{minipage}[b]{\linewidth}\raggedright
\textbf{Question}
\end{minipage} & \begin{minipage}[b]{\linewidth}\raggedright
\textbf{Insights from Tan2025}
\end{minipage} & \begin{minipage}[b]{\linewidth}\raggedright
\textbf{Insights from Abdelwahab2023}
\end{minipage} & \begin{minipage}[b]{\linewidth}\raggedright
\textbf{Insights from Betschart2022}
\end{minipage} & \begin{minipage}[b]{\linewidth}\raggedright
\textbf{Insights from Pinto2026}
\end{minipage} & \begin{minipage}[b]{\linewidth}\raggedright
\textbf{Insights from Guille2024}
\end{minipage} \\
\midrule\noalign{}
\endhead
\bottomrule\noalign{}
\endlastfoot
\textbf{When should I use GATK HaplotypeCaller?} & & & & & \\
\textbf{When should I use bcftools/samtools?} & & & & & \\
\textbf{When is an ensemble approach worth the effort?} & & & & & \\
\textbf{How do I choose filter thresholds?} & & & & & \\
\textbf{What are pitfalls of applying human-optimized pipelines to
non-model systems?} & & & & & \\
\textbf{Do AI-based callers (e.g., DeepVariant) outperform traditional
tools?} & & & & & \\
\textbf{Does alignment choice matter as much as caller choice?} & & & &
& \\
\end{longtable}

This table will be filled in collaboratively during discussion. One or
more students should photograph the final version for submission.

\subsection{Part 4: Case Study Challenge (15
minutes)}\label{part-4-case-study-challenge-15-minutes}

Break into mixed groups (each group should have at least one
representative from each paper). Your challenge:

\textbf{Design a variant calling benchmark for a non-model insect
species} with the following constraints:

\begin{itemize}
\tightlist
\item
  No high-quality reference panel (like Genome in a Bottle) exists.
\item
  You have 30× coverage Illumina short reads from 10 individuals.
\item
  Your goal is to identify population-level SNPs for a GWAS study.
\end{itemize}

Discuss and sketch answers to:

\begin{enumerate}
\def\labelenumi{\arabic{enumi}.}
\tightlist
\item
  Which variant caller(s) would you use, and why?
\item
  How would you validate your calls without a truth set?
\item
  What filtering strategy would you adopt?
\item
  Would you use an ensemble approach? Why or why not?
\end{enumerate}

Be prepared to share one key insight from your group during wrap-up.

\subsection{Part 5: Wrap-Up and Reflection (5
minutes)}\label{part-5-wrap-up-and-reflection-5-minutes}

Class discussion:

\begin{itemize}
\tightlist
\item
  Which study provided the most actionable recommendations?
\item
  What gaps remain in current benchmarking literature? (e.g.,
  underrepresented organisms, low-coverage data, structural variants)
\item
  If you were designing a variant calling pipeline for your own
  research, what is the one takeaway from today that will influence your
  decision?
\end{itemize}

\section{Deliverable and Grading}\label{deliverable-and-grading}

Your Lab 8 grade is based on \textbf{attendance and active
participation}. To receive credit, you must:

\begin{enumerate}
\def\labelenumi{\arabic{enumi}.}
\tightlist
\item
  \textbf{Attend the full lab session} and actively participate in all
  four parts of the activity.
\item
  \textbf{Photograph and upload} the completed \textbf{Paper Summary
  Table} (from your small group in Part 1) and the \textbf{Cross-Study
  Synthesis Table} (from the whole-class discussion in Part 3) to Canvas
  \textbf{before leaving lab}.
\item
  \textbf{Case study challenge}: briefly write up your tables design for
  the case study challenge and upload it to Canvas \textbf{before
  leaving lab}.
\end{enumerate}

\begin{tcolorbox}[enhanced jigsaw, rightrule=.15mm, leftrule=.75mm, colbacktitle=quarto-callout-warning-color!10!white, colframe=quarto-callout-warning-color-frame, toprule=.15mm, opacitybacktitle=0.6, title=\textcolor{quarto-callout-warning-color}{\faExclamationTriangle}\hspace{0.5em}{Submission Instructions}, left=2mm, colback=white, bottomtitle=1mm, toptitle=1mm, bottomrule=.15mm, coltitle=black, titlerule=0mm, breakable, opacityback=0, arc=.35mm]

\begin{itemize}
\tightlist
\item
  Take clear, legible photos of both tables (or a single photo if they
  fit on one board/page).
\item
  Upload to Canvas under ``Lab 8 Submission'' before the end of the lab
  period.
\item
  If working digitally, you may submit a shared Google Doc or screenshot
  instead.
\item
  \textbf{All group members} must submit their own photo/screenshot,
  even if the content is identical---this confirms individual
  attendance.
\item
  AI can be used to help with the initial paper summary and parsing of
  information, but the tables and case study design should be based on
  what you and your group generated during class without the use of
  generative AI.
\end{itemize}

\end{tcolorbox}

\section{References}\label{references-15}

\chapter{Lab09: IGV Genome Visualization Lab
Activity}\label{lab09-igv-genome-visualization-lab-activity}

\section{Overview}\label{overview-3}

Throughout this course, you have been building a complete NGS
variant-calling workflow using a real experimental evolution dataset. In
Lab 4, you downloaded whole-genome sequencing reads from the
\emph{Saccharomyces cerevisiae} population
SRR1693691\textsuperscript{24}, trimmed them, and aligned them to the
sacCer3 reference genome using BWA to produce a sorted BAM file. In Lab
7, you ran the GATK best practices pipeline on those alignments: you
marked PCR duplicates, called variants with HaplotypeCaller, examined
quality-score distributions using VCFtools density plots, and applied
hard filters to separate high-confidence SNPs from technical noise.
Together, those two labs took you from raw FASTQ files all the way to a
filtered VCF.

This lab is the final step of that journey. Here you will load the same
files---the pre-filtering BAM, the duplicate-marked BAM, and the
unfiltered VCF---into the \textbf{Integrative Genomics Viewer (IGV)} and
look at the raw read evidence underneath each variant call. As discussed
in \emph{Chapter 10}, automated filtering is powerful but imperfect:
visual inspection of alignments at candidate variant sites remains an
essential validation step, letting you recognize patterns such as strand
bias, duplicate stacking, and low mapping quality that are difficult to
catch with thresholds alone\textsuperscript{159}. You will also connect
the sequence changes you observe to the biological dataset described in
\textbf{Appendix C}---a pooled yeast population that experienced
hundreds of generations of selection---and reason about whether specific
SNPs could represent adaptive allele-frequency shifts.

By the end of this lab you will be able to:

\begin{itemize}
\tightlist
\item
  Load and configure a multi-track IGV session combining a reference
  genome, BAM alignment files, and a VCF.
\item
  Distinguish high-confidence SNPs from likely artifacts based on visual
  read evidence (strand balance, mapping quality, duplicate patterns,
  base quality).
\item
  Apply Human Genome Variation Society (HGVS) standardized variant
  nomenclature to describe DNA and protein sequence changes.
\item
  Interpret whether a called SNP falls in a coding region and assess its
  potential amino acid and functional consequences.
\item
  Reflect on how your filtering choices in Lab 7 relate to what you see
  in the read pileups.
\end{itemize}

\begin{center}\rule{0.5\linewidth}{0.5pt}\end{center}

\section{Step 1 --- Download and Install
IGV}\label{step-1-download-and-install-igv}

The software IGV is available from the Broad Institute (which also hosts
GATK). IGV stands for \textbf{Integrative Genomics Viewer} and allows
you to upload a variety of genome file formats to view on your desktop
machine. It runs on Java, so you may need to update your version of Java
to run it.

\begin{enumerate}
\def\labelenumi{\arabic{enumi}.}
\item
  Download IGV from
  \href{https://software.broadinstitute.org/software/igv/downloadLinks}{here}
  (make sure to scroll to the appropriate distribution based on your
  OS). Linux users will download the binary distribution, which can be
  launched directly from the terminal using Java.
\item
  Follow the instructions to install the software on your main OS.
\item
  Make sure after you install it that it loads properly before getting
  started on the exercise below.
\end{enumerate}

\section{Step 2 --- Load Data into IGV}\label{step-2-load-data-into-igv}

The goal of this step is to load the datasets we have worked with in
Labs 4 and 7 into IGV so we can visualize the variants and the read
evidence that supports them. Recall from
\textbf{\href{appendixC.qmd}{Appendix C}} that these files represent a
single \emph{S. cerevisiae} population sequenced from a pooled,
evolve-and-resequence experiment; because the sample is pooled rather
than a single diploid individual, you will see a range of allele
frequencies at variant sites rather than clean 50/100\% het/hom
patterns.

\begin{enumerate}
\def\labelenumi{\arabic{enumi}.}
\item
  Download files from
  \href{https://auburn.instructure.com/files/275453972/download?download_frd=1}{IGV\_files.zip}
  from Canvas.
\item
  Open IGV. Refer to the slides from this
  \href{https://auburn.hosted.panopto.com/Panopto/Pages/Viewer.aspx?id=51da88a9-e569-4453-b414-b29d0125a0ac}{tutorial}
  if needed throughout.
\item
  First, we need to load the genome file. At the top left, select the
  dropdown menu with the human genome and select the option ``Click for
  more''.
\item
  From the pop-up menu, scroll to the genome for \emph{S. cerevisiae
  (sacCer3)} and click OK to load it into the current session. This is
  the same sacCer3 reference (NCBI assembly R64) you indexed with BWA in
  Lab 4 and used for all downstream analyses.
\item
  Next, at the top left, select ``File'' and ``Load from File'' to
  import your raw sequence alignment (pre-filtering). The file should be
  named \texttt{SRR1693691.sorted.bam} and have a corresponding index
  file named \texttt{SRR1693691.sorted.bam.bai}. Make sure not to select
  the index, but only verify it exists in the same directory. IGV will
  search for it and give an error if it is not there.
\item
  Repeat step 5 on the duplicate-marked BAM file (post-filtering). The
  file should be named \texttt{SRR1693691.markdup.bam} and have a
  corresponding index file named \texttt{SRR1693691.markdup.bai}. As
  discussed in
  \hyperref[variant-calling-and-filtering-with-gatk]{Chapter 10},
  \emph{marking} duplicates rather than removing them lets you see the
  flagged reads in IGV and verify that the duplicate-calling algorithm
  behaved reasonably.
\item
  Finally, you will load the VCF file you generated (pre-filtering). At
  the top left, select ``File'' and ``Load from File'' to import the
  file \texttt{SRR1693691.SNPs.vcf.gz}. This file does not have an index
  but should still load.
\item
  Configure your view in the following ways:

  \begin{itemize}
  \item
    In the sorted BAM panel, right-click and scroll to ``Color
    alignments by'' and in the pull-out menu, select ``read strand''.
    This color-codes reads by forward and reverse, which is essential
    for detecting strand bias---a key artifact signal discussed in
    \hyperref[variant-calling-and-filtering-with-gatk]{Chapter 10}.
  \item
    In the markdup BAM panel, right-click and scroll to ``Duplicates''
    and select ``texture duplicates''. This makes PCR and optical
    duplicates visually distinct rather than hiding them, so you can see
    how many reads at each site are flagged.
  \item
    In the Sequence panel, click to show the standard translation panel
    of the amino acids. This will allow you to determine immediately
    whether a coding SNP changes the amino acid sequence.
  \end{itemize}
\item
  When you are done configuring your session, save it so you can return
  to this configuration in the future. Under ``File'' select ``Save
  Session'' and name your session Lab 9.
\end{enumerate}

\section{Step 3 --- Examine Variants in
Depth}\label{step-3-examine-variants-in-depth}

The goal of this step is to navigate to specific variants and evaluate
the raw read evidence that supports---or does not support---each as a
real, biologically relevant SNP. This mirrors the visual validation
workflow described in
\hyperref[variant-calling-and-filtering-with-gatk]{Chapter 10}, where
IGV is used to recognize artifact patterns (duplicate stacking,
strand-biased reads, low mapping quality, base-end clustering) versus
signatures of a high-confidence true variant (balanced strands, high MQ,
consistent base quality throughout reads).

\begin{enumerate}
\def\labelenumi{\arabic{enumi}.}
\item
  Now that you have the genome and two BAM tracks loaded into IGV, use
  the tutorial slides as a guide and navigate to the following regions.
  Zoom in enough to see individual reads and any SNP variants validated
  in your VCF:

  \begin{itemize}
  \tightlist
  \item
    Region 1: chrII:598117--598156
  \item
    Region 2: chrV:430725--431072
  \item
    Region 3: chrXV:468250--468449
  \end{itemize}
\item
  For each region, answer: how many apparent mismatches (potential
  variants) are visible in the \textbf{raw sorted BAM}? How many remain
  in the \textbf{markdup BAM}? And how many of those are called in the
  \textbf{VCF}? Think about why these three numbers differ.
\item
  For each variant present in the VCF within each window, assess the
  read evidence:

  \begin{itemize}
  \tightlist
  \item
    What is the depth of coverage at that site?
  \item
    Are reads supporting the alternate allele present on both forward
    and reverse strands?
  \item
    Are the supporting bases near the middle of reads (high confidence)
    or clustered at read ends (potential artifact)?
  \end{itemize}

  To better evaluate, sort the reads by base at each variant position:
\item
  Are there any apparent mismatches (or potential sequencing errors) in
  the BAM that are \textbf{not} ``called'' as variants in the VCF? Which
  features of those sites---coverage, strand distribution, read
  position, base quality---might explain why HaplotypeCaller did not
  ``call'' them as variants?
\item
  For the variants that \textbf{are} called in the VCF, based on the
  visual evidence (strand balance, mapping quality, duplicate patterns),
  do you expect that each variant ``Passed'' the hard filters you
  applied at the end of Lab 7? Recall the key thresholds from
  \hyperref[variant-calling-and-filtering-with-gatk]{Chapter 10}: QD
  \textgreater{} 2.0, MQ ≥ 40, FS \textless{} 60, SOR \textless{} 3.0,
  and ReadPosRankSum \textgreater{} −8.0.
\item
  Pick one variant and test your prediction by using \texttt{less} on
  the command line to search for that chromosome and position in your
  filtered VCF file.
\item
  Does the filtered VCF support or reject your prediction? Based on this
  comparison, would you adjust any of your filtering thresholds? In what
  direction, and why?
\item
  For each called SNP in the VCF, return to IGV: is the SNP located
  within a protein-coding exon? Use the gene annotation track and amino
  acid translation track to check.
\item
  If a SNP falls in a coding region, does the nucleotide change alter
  the amino acid? What effect---if any---do you think this amino acid
  change might have on the function of the gene? Remember that variants
  rising to high frequency in this experimental evolution dataset may
  reflect adaptive changes selected over hundreds of generations.
\item
  Using the HGVS naming conventions introduced in
  \hyperref[sequence-alignment]{Chapter 5}, write out the sequence
  changes for each variant you identified:

  \begin{itemize}
  \tightlist
  \item
    A prefix of \textbf{\texttt{g.}} is used for a genomic sequence
    change (e.g., g.69A\textgreater T for an A-to-T substitution at
    genomic position 69).
  \item
    A prefix of \textbf{\texttt{p.}} is used for a protein sequence
    change using the one-letter amino acid code (e.g., p.K23Y for a
    lysine-to-tyrosine substitution at position 23 in the protein).
  \item
    If the SNP is synonymous (silent), note that as well---synonymous
    changes are still described with \texttt{g.} notation.
  \end{itemize}
\item
  Feel free to explore additional regions on your own:

  \begin{itemize}
  \tightlist
  \item
    chrII:486633--487123
  \item
    chrV:553052--553399
  \item
    chrXIV:39075--39311
  \end{itemize}
\end{enumerate}

\begin{center}\rule{0.5\linewidth}{0.5pt}\end{center}

\section{Deliverables}\label{deliverables-3}

Upload the following to Canvas for a final grade:

\begin{itemize}
\tightlist
\item
  \textbf{Screenshots} of each of the three required regions showing
  your configured view (strand coloring, duplicate texturing, and
  translation track visible).
\item
  \textbf{Written answers} to the questions in Step 3 (questions 2--9),
  clearly numbered.
\item
  \textbf{HGVS variant notation} for each variant you identified in the
  three required regions (question 10), indicating both genomic
  (\texttt{g.}) and, where applicable, protein (\texttt{p.}) level
  changes.
\end{itemize}

\begin{center}\rule{0.5\linewidth}{0.5pt}\end{center}

IGV is a powerful way to visually validate specific variants and develop
intuition for what real biological signal looks like versus noise. As
you complete this lab, you are closing the loop on the entire
variant-calling workflow you have built across this course: from raw
FASTQ reads in Lab 4, through alignment and GATK processing in Lab 7, to
biological interpretation here. The ability to move fluidly between a
filtered VCF, the raw read alignments behind it, and the gene annotation
context is a core skill in modern genomics.

\chapter{Lab10: Population Genomics using the 1000 Genomes
Dataset}\label{lab10-population-genomics-using-the-1000-genomes-dataset}

\section{Lab 10: Population Genomics using the 1000 Genomes
Dataset}\label{lab-10-population-genomics-using-the-1000-genomes-dataset}

Today you will be working with human population genomics data from the
\href{https://www.internationalgenome.org/}{1000 genomes project}. The
goal of this exercise is to provide real world experience using
bioinformatics to examine genetic differences within and between
populations.

First, you will be assigned one of five regions (\textasciitilde2-10 Mb
each) that have been targeted due to known signatures of importance in
human evolution. Then, in Part 1, you will explore between population
divergence specific to your region. The examination of these plots and
statistics will be important in selecting populations to examine further
for within population diversity and selection as part of Part 2.
Finally, in Part 3, you will explore your region on the UCSC genome
browser to examine potential genes that may contribute to the population
genetic signal you see in your plots.

\section{Part 1: Examine region and identify target
populations}\label{part-1-examine-region-and-identify-target-populations}

\begin{enumerate}
\def\labelenumi{\arabic{enumi}.}
\item
  The 1000 genomes project sampled individuals from 26 human population,
  but today we will only be working with the four listed below. Each
  population has a unique three letter abbreviation:

  \begin{itemize}
  \tightlist
  \item
    CEU: Utah Residents (CEPH) w/ Northern \& Western European Ancestry
  \item
    YRI: Yoruba in Ibadan, Nigeria
  \item
    CHB: Han Chinese in Bejing, China
  \item
    PEL: Peruvians from Lima, Peru
  \end{itemize}
\item
  For today's exercise, I have created an
  \href{https://lstevison.shinyapps.io/Lab10-PopGen/}{R shiny app}. Go
  to this link and in the ``Between Populations'' tab, select your
  assigned region from the drop down menu and click the button ``Plot
  Fst''.
\item
  Examine the initial plot for your assigned region. What are the X and
  Y axes? These plots are based on a 100kb sliding window. How do you
  think different window sizes might influence the plot you are seeing?
\item
  Based on what you learned in class or find online, what is
  \href{https://en.wikipedia.org/wiki/Fixation_index}{Fst}? What does
  the estimate in the plot tell you about this region? What are the
  typical upper and lower bounds for Fst and why?
\item
  Next, use the second drop down menu ``Select Population Comparison''
  to examine FST statistics for each comparison for the whole region.
  These values represent the weighted Weir FST estimate from a
  bioinformatics software
  \href{https://vcftools.github.io/man_latest.html}{VCFtools}. Write
  down the max and average values for each pairwise comparison in the
  table below:
\end{enumerate}

\begin{longtable}[]{@{}
  >{\raggedright\arraybackslash}p{(\linewidth - 12\tabcolsep) * \real{0.1429}}
  >{\raggedright\arraybackslash}p{(\linewidth - 12\tabcolsep) * \real{0.1429}}
  >{\raggedright\arraybackslash}p{(\linewidth - 12\tabcolsep) * \real{0.1429}}
  >{\raggedright\arraybackslash}p{(\linewidth - 12\tabcolsep) * \real{0.1429}}
  >{\raggedright\arraybackslash}p{(\linewidth - 12\tabcolsep) * \real{0.1429}}
  >{\raggedright\arraybackslash}p{(\linewidth - 12\tabcolsep) * \real{0.1429}}
  >{\raggedright\arraybackslash}p{(\linewidth - 12\tabcolsep) * \real{0.1429}}@{}}
\toprule\noalign{}
\begin{minipage}[b]{\linewidth}\raggedright
Metric
\end{minipage} & \begin{minipage}[b]{\linewidth}\raggedright
CEU vs.~YRI
\end{minipage} & \begin{minipage}[b]{\linewidth}\raggedright
YRI vs.~CHB
\end{minipage} & \begin{minipage}[b]{\linewidth}\raggedright
CEU vs.~CHB
\end{minipage} & \begin{minipage}[b]{\linewidth}\raggedright
YRI vs.~PEL
\end{minipage} & \begin{minipage}[b]{\linewidth}\raggedright
CEU vs.~PEL
\end{minipage} & \begin{minipage}[b]{\linewidth}\raggedright
CHB vs.~PEL
\end{minipage} \\
\midrule\noalign{}
\endhead
\bottomrule\noalign{}
\endlastfoot
Average Fst & & & & & & \\
Max Fst & & & & & & \\
\end{longtable}

\begin{enumerate}
\def\labelenumi{\arabic{enumi}.}
\setcounter{enumi}{4}
\tightlist
\item
  In examining the 6 different plots, which two populations are most
  different across your region (consider both peak height and width)?
  Based on the Fst and the known upper/lower bounds, how important is
  this difference likely to be? \textgreater If multiple populations
  have similar values, select the two comparisons that you want to
  investigate further for the next part of this assignment based on the
  overall peak height of the highest value in your plot.
\end{enumerate}

\begin{itemize}
\tightlist
\item
  Max Fst difference: \_\_\_\_\_\_ v. \_\_\_\_\_\_
\end{itemize}

\begin{enumerate}
\def\labelenumi{\arabic{enumi}.}
\setcounter{enumi}{5}
\tightlist
\item
  Where (roughly) in this region is the peak of differentiation?
\end{enumerate}

\begin{itemize}
\tightlist
\item
  Region of max Fst: \_\_\_\_\_\_\_\_\_\_
\end{itemize}

\begin{enumerate}
\def\labelenumi{\arabic{enumi}.}
\setcounter{enumi}{6}
\tightlist
\item
  Right click on the Fst plot for the selected pair of populations to
  include in your submitted assignment that you will turn in.
\end{enumerate}

\section{Part 2: Investigate selected populations
further}\label{part-2-investigate-selected-populations-further}

\begin{enumerate}
\def\labelenumi{\arabic{enumi}.}
\item
  For Part 2, select the tab labeled ``Within Populations''.
\item
  Use the drop down menus to select your region and the two populations
  that you selected at the end of Part 1 to move forward with.
\item
  Click the button ``Plot Pi and Tajima's D''.
\item
  Scroll to the table at the bottom of the page. \textbf{How many
  individuals are in each of these two populations?}
\end{enumerate}

\begin{itemize}
\tightlist
\item
  Pop 1 Size: \_\_\_\_\_\_\_\_\_\_
\item
  Pop 2 Size: \_\_\_\_\_\_\_\_\_\_
\end{itemize}

\begin{enumerate}
\def\labelenumi{\arabic{enumi}.}
\setcounter{enumi}{4}
\item
  This set of four plots were generated again using the output of
  VCFtools, which was used here to calculate within population
  diversity, known as
  \href{https://en.wikipedia.org/wiki/Nucleotide_diversity}{pi} plotted
  in blue. What is pi?
\item
  In addition to Pi, you should also see two sets of plots with
  \href{https://youtu.be/wiyay4YMq2A?si=LlrqVCvglyxwKDhq}{Tajima's D} in
  red across the region in both populations.
  \href{https://vcftools.github.io/man_latest.html\#OUTPUT\%20OPTIONS}{Vcftools}
  was also used to generate the data used for these plots. What does
  Tajima's D tell you about the populations? What does a positive value
  mean? Negative value?
\item
  Like Fst, these statistics are also estimated based on 100kb windows?
  Do you think the results would be different in the window size were
  bigger? Smaller?
\item
  Based on the number of individuals in your two selected populations,
  how reliable do you feel these estimates are in capturing the
  population genetic signatures? How does sample size influence
  estimates of pi and Tajima's D? Why?
\item
  Examine the four plots considering also the plots you examined in Part
  1.
\end{enumerate}

\begin{enumerate}
\def\labelenumi{\alph{enumi}.}
\item
  What similarities do you see between these plots and the one you made
  earlier?
\item
  Does the sub-region you recorded at the end of Part 1 stand out in
  either of these plots?
\item
  Make a hypothesis as to which population may be driving this result?
  Hint: which population has the lowest value of Tajima's D?
\item
  Based on \href{https://youtu.be/wiyay4YMq2A?si=LlrqVCvglyxwKDhq}{this
  video}, interpret the results of the Tajima's D statistic in this
  plot.
\end{enumerate}

\begin{enumerate}
\def\labelenumi{\arabic{enumi}.}
\setcounter{enumi}{9}
\tightlist
\item
  Right click on the image of Pi and Tajima's D for this population to
  include in your submitted assignment that you will turn in.
\end{enumerate}

\section{Part 3: Examine region in Genome
Browser}\label{part-3-examine-region-in-genome-browser}

\begin{enumerate}
\def\labelenumi{\arabic{enumi}.}
\item
  For Part 3, click the third tab in the app titled ``Explore''. Again,
  select your region from the drop down menu and click on the button
  ``Load Region''.
\item
  Record the chromosome and genomic coordinates of your region below and
  then click the link to \href{https://genome.ucsc.edu/}{View in UCSC}
  which will take you to the UCSC Genome Browser. These data were mapped
  to the human genome reference hg19.
\item
  Zoom into the region from Part 1 of Max Fst (\textasciitilde500kb
  window is fine). Look at the genes present. Make a list of the genes
  present in the region:
\item
  Take a screen shot to include with your submitted assignment that you
  will turn in.
\item
  Is there anything special about any of the genes in your sub-region
  that can explain why they differ? Click on the genes and read about
  them, but also do some google investigative work to find out!
\end{enumerate}

\begin{itemize}
\tightlist
\item
  Based on what you know about human evolution, does any particular gene
  support your hypothesis in Part 2 as to which of the two populations
  might be driving this pattern? What are some known differences about
  these two human populations?
\end{itemize}

\begin{enumerate}
\def\labelenumi{\arabic{enumi}.}
\setcounter{enumi}{5}
\item
  Based on what you know about human evolution, what additional
  regions/genes in the human genome would you like to analyze like this
  in the future?
\item
  When you are done, compare your results to others in the class and the
  other regions. When you are done, write up a short (\textless200
  words) summary of what you did in this lab and what you learned from
  your region. Include the plots from Part 1 with your selected pair of
  populations, and the 2 plots from Part 2 for the population you
  hypothesize to be driving this signature. Finally, include the screen
  shot of your zoomed in region with the selected gene from the UCSC.
\end{enumerate}

\begin{center}\rule{0.5\linewidth}{0.5pt}\end{center}

This exercise should provide a sense of how sliding window approaches
are used to identify regions of interest in population genetics studies.
Typically, these scans are done genome-wide along with other statistics
to identify regions of interest. Then, gene annotations are used to find
overlapping functional regions that can explain the genetic patterns.
This can be done programmatically by downloading annotations from UCSC
and using \texttt{bedtools} to intersect with the regions of interest.
Then,
\href{https://genome.ucsc.edu/goldenPath/help/customTrack.html}{custom
tracks} can be used to upload FST, pi, and Tajima's D results to the
genome browser. We will discuss sliding window analyses in more depth in
two weeks.

\cleardoublepage
\phantomsection
\addcontentsline{toc}{part}{Appendices}
\appendix

\part{Appendices}

\chapter{Appendix: How to Read Scientific
Literature}\label{appendix-how-to-read-scientific-literature}

\section{Reading Primary Scientific Literature}\label{sec-papers}

This appendix is a refresher on how to read and interpret original
research articles. It is especially relevant for Lab01 (finding and
describing a genome analysis paper) and for later labs and projects
where you will evaluate genomics studies in the context of the three
perspectives from Chapter 1 and the ethical themes introduced
there\textsuperscript{1,2,10}.

Sharing the results of scientific experiments is one of the most
important aspects of science. Early scientists shared their observations
and results in open letters to their colleagues that were hand copied
and passed from person to person. Pretty soon, scientists developed
better ways of sharing their results with each other, and our current
knowledge of science and the world around us is dependent upon current
scientists sharing their results.

Scientists share their discoveries with each other in several ways.
Informally, scientific colleagues may call or email each other about
their latest research. But once the experiments have been done and the
analysis carried out, a scientist needs to share this new information
with a wider audience in a more formal manner. Other scientists can then
find this information to help them design their own experiments, to put
their own research in context, and to add to our knowledge of the world.

One of the most common ways for scientists to formally share their
research is to publish a ``journal article.'' The journal article is a
part of the formal written record of the scientific process, and you've
probably come across them before. The typical format includes a
discussion of the authors' original research, offers thoughtful analysis
of the results, and cites relevant papers from other authors that relate
to the research. When the authors are discussing research and
experiments that they carried out, and giving you an outline of the
experiment, this is called an original research article. We will discuss
other types of writing in part 3 of this assignment.

\section{Part 1: Breaking down the
paper}\label{part-1-breaking-down-the-paper}

In this activity, you will read and interpret an original research
article. You will be asked to identify the hypothesis and the main
ideas. From that, you will work to make a ``concept map'' of the
experiment. Finally, you will think of ideas for follow up experiments
based on these results.

The article you will read was published in 2012 in the Journal of
Heredity. It can be accessed via the web for\textsuperscript{16}. If you
haven't already done so, read over the paper. This quick
guide\textsuperscript{210} can be reference for how to approach
scientific literature.

\subsection{Step 1: Hypothesis}\label{step-1-hypothesis}

A scientific hypothesis (plural hypotheses) is the initial building
block in the scientific method. It is a proposed explanation for an
observation. In other words, it is a prediction or explanation that can
be tested by an experiment. Observations and experiments may support a
scientific hypothesis, but can never entirely prove one.

Typically, a hypothesis is an if-then proposition, written in the form,
``If X, then Y.'' Spend a few minutes reading through the article and
try to write out what you think the hypothesis is based on the article.

\subsection{Step 2: Concept Map}\label{step-2-concept-map}

A concept map is a visual tool that helps show the relationship between
ideas. Concept maps begin with a main idea and then branch out to show
how the idea can be broken down into sub-topics. Concept maps can help
you organize information from your reading. The concept map arranges
related ideas in a hierarchy. You start off broad and the sub-topics
will get more and more specific. Watch the following
\href{https://www.youtube.com/watch?v=sZJj6DwCqSU&t=3s}{brief video} on
concept mapping.

After reading the paper, try to identify the main ideas. On scratch
paper, make a list of these ideas and start to construct a concept map
of how these are related. First, you'll want to start with a central
idea from your list. Next, what is that central idea related to from
your list? Your concept map can be as small or big as you need it to be
in order to cover your topic. If your concept map ends up being too big,
don't worry! The purpose of the map was to help you organize your ideas
to better understand the article.

\subsection{Step 3: What's Missing}\label{step-3-whats-missing}

Sometimes, we can have difficulty interpreting a scientific article.
First, look up anything you don't understand. This may require
additional reading, but the more you read, the less you will look up and
over time you will get more and more comfortable with this process.
Second, accept that this may not be your fault! Often scientists may
leave out details they assume readers will know making it harder for
people to understand. This is not to say this work is not properly
vetted. A typical scholarly, scientific journal article is undergoes a
rigorous editing process known as peer review (more on this later).
Remember, your ability to critique scientific work is directly related
to your ability to understand the findings. It is important that you
take the time needed to look up what you don't know and try to identify
additional pieces of information you need to properly assess the work.

\textbf{\emph{In re-reading the article, what were some pieces you felt
were missing that would have made it easier to read had they been
included? Your list need not be exhaustive, but try to think of 1 or 2
things that would help you read the article better.}}

\subsection{Step 4: What's Next}\label{step-4-whats-next}

After reading a scientific article, you may have lots of follow up
questions. Every scientific experiment is just one step in a larger
puzzle. Perhaps one of the most exciting (or maybe exhausting?) parts of
scientific research is that there is always more to uncover. Our work is
never done! In this spirit, one way to reflect on reading can be to
think about what you would do next if you were the researcher? Design a
follow-up experiment to this study based on what you learned.

On your scratch paper, write out how you would conduct the experiment.
Be sure to apply the scientific method. Present a clearly defined
hypothesis and how you would go about testing that hypothesis.

Some questions to think about in designing your experiment:

\begin{itemize}
\tightlist
\item
  What resources would you need to conduct the study?
\item
  How long would it take to get results?
\item
  Who might be interested in your findings?
\end{itemize}

Your experiment does not have to be 100\% related to this study, but it
should be inspired to some degree by reading and discussing this work.

\section{Part 2: Further Exploration}\label{part-2-further-exploration}

Now that you've read the paper over and made a concept map, take some
time to explore the findings further to reinforce your understanding.

For this part of the activity, you will explore OCA on OMIM, learn about
other animal models for this disease, explore other databases, download
and annotate the tyrosinase gene, and finally, find out what's happening
with Migaloo the whale.

\subsection{Step 1: Oculocutaneous Albinism
(OCA)}\label{step-1-oculocutaneous-albinism-oca}

In the article, it lists the MIM number for OCA. Look this number up in
the paper and go to the \href{https://www.omim.org/}{OMIM database} and
type it in. As you may recall from an earlier in-class activity, OMIM is
a comprehensive, authoritative compendium of human genes and genetic
phenotypes that is freely available and updated daily. The overviews in
OMIM contain information on all known mendelian disorders and over
15,000 genes with references to additional primary literature. OMIM
focuses on the relationship between phenotype and genotype. Once you
find the record with OCA corresponding to the number in the paper, click
it.

\emph{What is the cytogenetic location for the gene, tyrosinase?}

\subsection{Step 2: Animal Models}\label{step-2-animal-models}

While OMIM is focused on Human Diseases, it has a section for animal
models. On the left-hand side of the page, navigate to the Animal Model
section. Notice that Migaloo is listed here in reference to this study!

\emph{In addition to Migaloo, what other animal models for OCA1 are
listed?}

\subsection{Step 3: Albinism database}\label{step-3-albinism-database}

In the third paragraph, the article also references the albinism
database. Remove the date from the URL and go to the website. You'll
notice that you get rerouted. Explore this website.

\emph{Is this database likely to have different information since the
publication of this article?}

\subsection{Step 4: Gene Annotation}\label{step-4-gene-annotation}

GenBank is a nucleotide database curated as part of the NCBI. The
journal article provides a GenBank Accession number for the tyrosinase
gene. Navigate to GenBank and pull up the accession listed in the
article. Walk through the following steps to annotate the PCR primer
sequence from the article to see where it is located in the gene
sequence. 1. Copy and paste the mRNA sequence into a word document on
your computer. 2. Using the search function, find the forward PCR primer
listed in the article in this pasted mRNA sequence. You may need to only
search 5-6 base pairs at a time as the sequence has space gaps in this
format making it impossible to find the full-length sequence without
some tedious edits. Also note that the primer sequence has an R in it.
You may want to change that to a G. This is what is called DNA code
ambiguity (learn more
\href{https://droog.gs.washington.edu/parc/images/iupac.html}{here}!).
3. Once you find the primer in the larger sequence, mark it in some way
with word edits (e.g.~bold, italic, highlight, change text color, etc.).
4. Return to the GenBank website and look at the FEATURES section of the
record for this sequence. Here, you will see the location of each
attribute of the gene. For example, the full-length gene can be found
from position 1 to 2096. 5. Identify the feature that overlaps the
primer sequence you annotated in Step 2.

\emph{Based on where the PCR primer is located in the sequence, which
feature of the mRNA does this primer include?}

\subsection{Step 5: Migaloo, where are
you?}\label{step-5-migaloo-where-are-you}

Migaloo the white whale is rather famous now. He even has his own
twitter handle
\href{https://twitter.com/migaloo1}{\textbackslash@Migaloo1} with over
10,000 followers! There are also several other websites dedicated to
following Migaloo and other white whales. Do a bit of internet research
of Migaloo sightings.

\emph{Based on your research, when was Migaloo last sighted?}

\section{Part 3: Publishing \& types of scientific
writing}\label{part-3-publishing-types-of-scientific-writing}

In Part 1 of this activity, we focused on original or primary scientific
literature. A slightly different type of journal article (called a
``REVIEW ARTICLE'') will not report on original research, but will
outline the current state of research in a particular field, citing the
appropriate literature and connecting the various pieces of research
together. Like primary research articles, review articles are generally
peer reviewed.

Review articles and original research articles can often look the same
at first glance, and most search engines or databases won't tell you
what type of an article it is. To tell them apart, you need to identify
whether the authors are discussing their own research and experiments or
someone else's. Often, the ``Materials and methods'' section (aka
``Experimental procedures'' or something similar) will be your best
clue. This section is occasionally stored online, separate from the
article as a part of the ``supplementary materials''. A review article
may also read more like a textbook chapter than an experiment.

Both type of journal articles can be found online and/or in print.

When you are searching for information on a research topic, you may also
run across some other types of information. Shorter news articles (1-2
pages) may appear in some scientific and popular publications reporting
on recent developments in a particular field, or reporting on a
particular piece of research. These news articles are not peer reviewed,
and are normally written by science journalists, not researchers. The
news articles may be easier to read, but since they are normally one or
two steps removed from the original research, a news article may not be
the best source for your paper or project. However, news articles can
lead you to a piece of original research, and can help you easily stay
informed about recent research developments.

If you conduct your searches online, via Google, Yahoo or another
popular search engine, you may find journal articles, but you may also
come across other scientific information that can take many forms.
Wikis, blogs and personal websites can often contain a lot of scientific
information, but these resources are generally very far removed from the
original research where the ideas were first developed. Each of these
sources needs to be evaluated very carefully to determine if the
information is credible, and these sources won't be suitable for a
research paper. There is a lot of great scientific information on the
web, but there is also a lot of bad science, pseudo-science, and
non-science-pretending-to-be-science available and distinguishing them
can be tricky. We will talk more on finding credible sources in the next
part of this activity.

For this part of the activity, you will be provided a popular news
article on the same topic as the original article you read in Part 1.
You will be asked to compare and contrast the two types of writing.

\subsection{Step 1: Compare and
Contrast}\label{step-1-compare-and-contrast}

Watch this
\href{https://www.youtube.com/watch?v=TRfu48motcQ&t=108s}{brief video}
on the difference between journals versus magazines.

\emph{What are some of the key differences between magazines and
journals?}

\subsection{Step 2: See the
Difference?}\label{step-2-see-the-difference}

Now, read this \href{https://www.pacificwhale.org/blog/migaloo/}{blog
article} on Migaloo the whale. Afterward, think about how this work
compares to the journal article you read in Part 1.

Using the differences listed in response to the last question, how do
some of these compare for this specific set of journal versus magazine
article? You may also want to contrast each piece with the style of
writing in textbooks.

\subsection{Step 3: Peer Review}\label{step-3-peer-review}

Peer review is the process that allows scientists to trust the
reliability of published journal articles. Here's how it works:

\begin{enumerate}
\def\labelenumi{\arabic{enumi}.}
\tightlist
\item
  A scientist submits an article to a journal saying ``please publish
  this article.''
\item
  The journal finds 2 or 3 people who know a lot about the research
  topic, called REVIEWERS or REFEREES, and asks them to look at the
  article.
\item
  The reviewers look at the article carefully. They check to see if the
  experiment is designed and conducted well, they look at the analysis
  of the data, they see whether the conclusions are justified by the
  data, and they make sure the article can be understood by other
  scientists. They also make a judgment about how ``important'' the
  article is. Some journals only accept really innovative and important
  research, other journals accept research that advances the field just
  a little bit.
\item
  The reviewers say ``yes, we should publish this article'', ``no, we
  shouldn't publish this article'' or ``if the author makes some
  changes, maybe we should publish this article''
\item
  If the article is published, we can say that it has been PEER
  REVIEWED.
\end{enumerate}

Scientists rely on their colleagues, the reviewers, to make sure that
good science is given a wide audience and that not-so-good science stays
out of the science journals. Because blogs, wikis and personal websites
don't automatically have this expert filter, you have to do a lot more
digging to determine if the information is reliable.

The only way to tell if a journal article has been peer reviewed is to
look for information about the journal itself, normally on the
publishers website. Most databases won't indicate if an article is peer
reviewed or not.

For the article you read in Part 1, go to the
\href{https://academic.oup.com/jhered/pages/Policies}{journal's website}
and read over their peer review process. It can be slightly different
for different journals.

\emph{Continue reading about this journal on their website. When was the
Journal of Heredity established?}

\subsection{Step 4: Timeline to
Publication}\label{step-4-timeline-to-publication}

The information on peer review is usually included with the original
article. Explore the article history to see how long the article spent
in peer review.

\emph{Based on when the article was received versus when it was
accepted, how many months did it spend before being accepted?}

\section{Concluding Remarks}\label{concluding-remarks}

This activity was designed to familiarize you with reading and
understanding scientific literature. In your research in our lab at
Auburn, you will need to read scientific articles for your project and
for lab meeting paper discussions. Hopefully this introduction makes
that task seem less daunting.

One thing not covered here was how to find scientific literature on your
own. This process of vetting articles is also important. To help get you
started, I recommend attending the
\href{https://libguides.auburn.edu/savvy_researcher/workshops}{Research
Savvy Bootcamps} offered every semester.

From here, you can find the resources to help continue your journey in
reading, understanding, and finding scientific literature.

Additionally, you can request a meeting with a
\href{https://libguides.auburn.edu/subjectlibrarians}{subject librarian}
if you want to learn more about the resources available through the
library and how best to use them.

As you move into Lab01 and later assignments, you can reuse the same
skills:

\begin{itemize}
\tightlist
\item
  Identify hypotheses and concept maps for genomics papers you select.
\item
  Ask ``what's missing?'' when authors gloss over methods or QC details.
\item
  Propose follow‑up experiments or analyses that build on published
  genomics work.
\end{itemize}

This is exactly the kind of reading and reasoning you will need for the
semester‑long PRDM9 project (Appendix: Semester Research Project) and
for your own future research\textsuperscript{1,2,24}.

\chapter{Appendix: Semester Research
Project}\label{appendix-semester-research-project}

\section{Group Research Project Overview}\label{sec-prdm9}

The majority of this course will be based on a semester-long genome
analysis. In the beginning of the semester, you will choose a group
based on the projects listed below. Once the projects are selected, you
will form groups based on the projects in class. Students may choose to
work individually if desired. Otherwise, groups sizes should be based on
the scale of the project selected (see descriptions below). Groups will
be required to meet OUTSIDE of class during the semester to discuss
their project goals and complete their research plan (see below).

\subsection{Learning Objectives}\label{learning-objectives-3}

\begin{itemize}
\tightlist
\item
  Experience with real data and research questions
\item
  Reinforcement of the scientific method
\item
  Learn common file types used for raw sequence data, alignments to
  reference genomes, and variant files
\item
  Learn how to assess the quality of NGS data
\item
  Learn a variety of bioinformatics tools for conducting genome analysis
\item
  Develop proficiency in scientific communication skills and
  reproducibility of research
\item
  Perform basic statistical analysis of NGS data using R
\item
  Use git for collaboration on a research project
\item
  Become proficient in shell scripting
\item
  Work together as a team to conduct a research project (see below)
\end{itemize}

\section{The Project: Evolution of the recombination regulation protein
PRDM9 across non-avian
reptiles}\label{the-project-evolution-of-the-recombination-regulation-protein-prdm9-across-non-avian-reptiles}

Since 2010, the protein PRDM9 has been of interest to scientists for its
role in modifying the meiotic recombination rate
landscape\textsuperscript{211}. Specifically, in taxa without a
functional copy of this protein, recombination initiates near gene
promoter regions. However, in taxa with a functional copy of this
protein, recombination initiates near binding sites. Further, this
protein has been shown to be involved in both male and female fertility
across mammals, and speciation within mice\textsuperscript{212,213}.

PRDM9 has four major domains, including a zinc-finger domain that allows
it to bind to DNA and act as a transcription factor to initiate
recombination. This dramatically alters the recombination landscape
across the genome. Further, because this protein is rapidly evolving,
the binding motif also changes rapidly, leaving recombination rates to
evolve more rapidly in taxa with functional copies than without. This
protein has been well studied in mammals, and particularly in
primates\textsuperscript{214}. While we know it has been lost in birds
and crocodiles, its role in other non-avian reptiles (e.g.~testudines
and squamates) is largely unknown. Specifically, recent work in snakes
has revealed non-canonical patterns of recombination regulation that
warrant a deeper look in this clade\textsuperscript{215}. Further, there
is major variation across this clade in sex-determination and
reproductive modes that could impact its role in fertility in unknown
ways.

A previous study surveyed a large number of vertebrate taxa, including
reptiles, revealing loss of function across several major clades within
reptiles, including complete loss of function in birds and
crocodiles\textsuperscript{216}, and a recent loss in
Anoles\textsuperscript{217}. However, they were limited in their
analysis due to a lack of available reptile genomes at the
time\textsuperscript{218}. There are now over 165 published non-avian
reptile genomes with this number growing
rapidly\textsuperscript{219--221} making this a ripe time to conduct a
detailed portrait on the evolution of this protein in non-avian
reptiles. Our goal this semester in Bioinformatics Class will be to
follow up on this interesting question using publicly available genomic
data in two major clades of non-avian reptiles -- testudines
(\href{https://www.ncbi.nlm.nih.gov/datasets/genome/?taxon=8459&reference_only=true}{N=59})
and squamates
(\href{https://www.ncbi.nlm.nih.gov/datasets/genome/?taxon=8509&reference_only=true}{N=307}).

This project was started last spring in this course, with much insight
gained to make the work this semester run smoother. But it is worth
noting that with research, there is always an uncertainty that can lead
to frustration. It can also seem like we ``do not know what we are
doing''. This feeling should be embraced - you are not doing busy work
on some known dataset, you are delving into the unknown and that is
exciting, but also remember that uncharted territory is often rugged. So
in the same way that you would dress appropriately for a trek through
the Amazon, approach this project with curiosity and enthusiasm.

\section{Project 1: All students done
invidually}\label{project-1-all-students-done-invidually}

Early in the semester, each student will be assigned a genome from a
non-avian reptile where there is no annotation of the reference genome.
Your job will be to manually annotate only the PRDM9 gene. You will use
web-based tools to identify the rough location of the ortholog of PRDM9
using BLAST. You will then use an exon-by-exon mapping approach to
develop a complete gene with start and end positions as well as the
location of each exon/intron. The goal will be to submit a gene report
mid-way through the semester of your proposed gene model for PRDM9 in
your assigned species. This will include a GFF of the gene model, a
screenshot of the genome view with your gene model and BLAST results
visible, as well as a two-way sequence alignment between your predicted
protein sequence and the PRDM9 protein from a close relative. This
protein sequence will be added to a growing database of non-avian PRDM9
sequences that will be used for this project.

\section{Project 2: Semester-long group
project}\label{project-2-semester-long-group-project}

Early in the semester, you will form groups (students may also choose to
work alone) to each investigate a key question towards the larger
research goal. At the end of the semester, we will compile what we find
collectively and hopefully be able to draft a manuscript of our findings
for scientific publication. Below I have outlined a set of open
questions that I think would lead to fruitful independent as well as
collective results. If students want to pick a question that is NOT on
this list, that would be OK, but they should meet up with me during
office hours to discuss their plan early in the semester.

Each group will put together a github repository of their research
project and give a presentation during finals week. You will also be
required to provide feedback for both the GitHub and the presentation in
the form of peer review. The peer review will be INDIVIDUAL feedback to
the groups on their work.

The data analysis may employ a variety of bioinformatic tools you have
learned during the semester or elsewhere. Some projects will rely on the
data your group collects, so be mindful of this collaborative nature
among groups and plan accordingly. Each project will address a
scientific question or hypothesis. I have set aside a few class periods
to work on these in class so that we can ALL brainstorm on the data
analysis for EVERY project together.

I encourage you to start each step early to make sure you have enough
progress for it to be completed. For example, a preliminary analysis and
your GitHub repository are due just after Spring Break. For this part of
the assignment, you will get credit for completion ONLY, BUT it will be
used for peer review. This means the depth of feedback from your peers
is dependent upon the scale of your progress thus far.

\section{Summary of Student Led Group Project
Options:}\label{summary-of-student-led-group-project-options}

\textbf{Project 1: Is PRDM9 expressed in the germline of the selected
taxa?}: Because of its role in meiosis, expression in the germline is a
huge component of validating the functionality of PRDM9. This project
will likely start with surveying available RNA-seq datasets (See
Supplementary File 2 from\textsuperscript{218}) and then, based on
availability of data, it may be used to help narrow down the selected
taxa for our downstream investigation. Students who have already taken
functional genomics are particularly well-suited to this project. A
major goal of this would be to validate systems where there is no
evidence of PRDM9 based on our larger genome survey by comparing
expression of PRDM9 to other genes involved in recombination rate as
positive controls (see methods from\textsuperscript{217}).

\textbf{Project 2: Is there evidence for positive selection in the zinc
finger residues of PRDM9 across non-avian reptiles?} There are four
amino acid residues within each zinc finger that make contact when
binding to DNA. Therefore, these residues are essential for
investigating the rate of evolution of this protein. The goal here would
be to generate a multiple sequence alignment of all the zinc fingers and
then use PAML to estimate dN/dS (omega) to investigate the rate of
evolution with a few different bioinformatics tools (see Figure 2
of\textsuperscript{214}). Students should consider alternative
approaches as well (such as newer methods/tools,
e.g.~\href{https://academic.oup.com/evolut/advance-article-abstract/doi/10.1093/evolut/qpaf240/8378101}{MoleRate}).

\textbf{Project 3: What are the evolutionary relationships between the
SET domain of PRDM9 across non-avian reptiles?} This domain has specific
residues with known catalytic activity that is necessary for opening the
chromatin for the zinc finger array to bind to DNA. This project would
investigate the SET domain in depth and explore the evolutionary
relationships of this domain similar to the project above.

\textbf{Project 4: What are the evolutionary relationships between the
SSXRD domain of PRDM9 across non-avian reptiles?} Related to the two
above questions, this project would investigate the SSXRD domain. This
is the shortest functional domain in PRDM9, so it is unclear how
necessary it is for function of this protein. Variation in
presence/absence of this domain on its own would be a novel finding.

\textbf{Project 5: What are the evolutionary relationships between the
KRAB domain of PRDM9 across non-avian reptiles?} Related to the three
above questions, this project would investigate the KRAB domain which is
proposed to recruit recombination machinery after DNA binding occurs.
\href{https://pubmed.ncbi.nlm.nih.gov/37716846}{Suggested additional
reading}.

\emph{Depending on interest and the group size, projects 2-5 could
easily be combined. Otherwise, each would be suitable for an independent
project.}

\textbf{Project 6: How do the predicted binding motifs of PRDM9 and the
number of predicted binding sites across the genome differ across
non-avian reptiles?} This project would use the zinc finger residue to
predict the DNA binding sequence and examine how degenerate it is.
Further, the predicted binding sequence matrix would be searched across
the reference genome for each selected taxa to compare the number of
predicted hits and how they cluster in the genome.

\textbf{Project 7: What are the potential positional effects of PRDM9 in
the genome that impact its evolution?} One thing our work has recently
revealed in investigating PRDM9 along so many genome is that it seems to
be located towards chromosome ends very often. Further, the synteny of
this gene is not well-maintained, which could be a contributing factor
to its rapid evolution. The goal of this project would be to use genomes
with chromosome-level assemblies to investigate this pattern more
robustly. It may also include diving into neighboring genes a bit to see
how this gene moves within the genomes across the tree of life. You may
consider the high-quality genomic assembly standard introduced in the
first day of class reading
(e.g.~\href{https://www.earthbiogenome.org/report-on-assembly-standards}{6.C.Q40}).

\textbf{Project 8: For the newly identified losses in testudines and
squamates, what is the timing of these events?} Dogs are the most recent
reported loss in the literature dated at around 14 million years. But
our work since Spring 2025 has identified what appear to be much more
recent losses that have yet to be formally dated. Thus, this project
will involve more of a literature-based approach that will revolve
around narrowing down the time frame for each of the putative losses
previously identified in these two clades. This may also require
additional manual annotation of additional close relatives as a
continuation of the first project earlier in the semester.

\textbf{Project 9: Are there any phenotypic traits that correspond to
presence/absence of PRDM9?} For this project, students will zoom out to
the entire vertebrate tree and examine various databases to characterize
various phenotypic and life history traits across species with and
without PRDM9. This will then be used in a statistical framework to
identify potential explanatory variables to this gene. Potential tools
to analyze this should be explored, including
\href{https://academic.oup.com/bioinformatics/article/35/22/4815/5514536}{R
Package RER Converge}. See
\href{https://www.science.org/doi/10.1126/sciadv.ads8488}{Figure 2A in
this additional paper} as an example of this type of analysis in fishes.

\section{References}\label{references-16}

\chapter{Appendix: Yeast Experimental Evolution
Dataset}\label{appendix-yeast-experimental-evolution-dataset}

\section{Overview of the yeast dataset}\label{sec-yeast}

In the labs, you will work with whole-genome sequencing data from a
single \emph{Saccharomyces cerevisiae} population that was part of an
experimental evolution study on standing genetic variation and
adaptation in outcrossing yeast. The dataset you will analyze is one run
from this experiment deposited in the NCBI Sequence Read Archive (SRA)
under accession \textbf{SRR1693691}. This sample represents pooled
genomic DNA from many yeast cells, capturing allele frequency changes
that occurred over hundreds of generations under defined selection
conditions in the original study.\textsuperscript{24}

The reference genome used throughout the course is the sacCer3 build of
the \emph{S. cerevisiae} S288C genome (NCBI assembly R64), which is a
compact, well-annotated eukaryotic genome of roughly 12 Mb spread across
16 chromosomes plus mitochondrial DNA. The small genome size makes it
feasible to run a full variant-calling workflow (alignment,
post-processing, GATK variant calling, and filtering) on the
high-performance computing (HPC) system within a single semester, while
still exposing you to tools and file formats that are standard in modern
human and biomedical genomics.

\section{Where the data come from}\label{where-the-data-come-from}

The sequencing data for SRR1693691 were generated as part of an
``evolve-and-resequence'' experiment in which yeast populations were
founded from a genetically diverse base population and then propagated
for many generations with periodic outcrossing.@burke\_standing\_2014 In
this design, recombination reshuffles standing genetic variants across
the genome, allowing adaptation to proceed from existing variation
rather than relying primarily on new mutations, and making evolutionary
responses more repeatable across replicate populations.

The original publication (``Standing Genetic Variation Drives Repeatable
Experimental Evolution in Outcrossing Populations of \emph{Saccharomyces
cerevisiae}'') used whole-genome resequencing of evolved populations to
identify genomic regions where allele frequencies changed consistently
across replicates under selection. In this course, you will not attempt
to fully reproduce the paper's population-genetic analyses; instead, you
will focus on learning how to process raw sequencing reads into
high-quality variant calls and then interpret those variants in a
biological and biomedical context using a subset of the data.

\section{Core files and formats you will
see}\label{core-files-and-formats-you-will-see}

You will repeatedly encounter a small set of core file types as you move
from raw data to filtered variants, and part of the course goal is to
help you learn to ``read'' these files as biological objects as well as
technical ones.

\begin{itemize}
\item
  \textbf{Raw sequencing reads (FASTQ)}:\\
  The SRA run SRR1693691 can be converted into one or more compressed
  FASTQ files (for example, \texttt{SRR1693691.fastq.gz} or lane-split
  files) that contain the raw read sequences and their per-base quality
  scores. In Lab 4 you will download the reads from NCBI using provided
  scripts, inspect the FASTQ structure, and run \textbf{FastQC} to
  assess base quality, GC content, adapter contamination, and other
  quality metrics before and after trimming.
\item
  \textbf{Reference genome (FASTA and index files)}:\\
  The sacCer3 reference genome is provided as a FASTA file (for example,
  \texttt{sacCer3.masked.fa}), which you will index with \textbf{bwa}
  and \textbf{samtools} to support alignment and downstream operations.
  Index files such as \texttt{sacCer3.masked.fa.bwt} (for BWA) and
  \texttt{sacCer3.masked.fa.fai} (for samtools) allow fast lookups of
  sequence coordinates, and you will use the \texttt{.fai} file to
  compute the total genome size and estimate coverage in Lab 4.
\item
  \textbf{Aligned reads (BAM and BAI)}:\\
  After trimming, you will align the reads to sacCer3 using BWA and
  convert the output into a coordinate-sorted BAM file, for example
  \texttt{SRR1693691.sorted.bam}, plus an index file
  \texttt{SRR1693691.sorted.bam.bai} that enables rapid random access to
  specific genomic regions. Later, as part of the GATK post-processing
  workflow, you will mark PCR and optical duplicates to produce a
  duplicate-marked BAM (for example, \texttt{SRR1693691.markdup.bam})
  that better reflects the true independent read depth at each site.
\item
  \textbf{Variant calls (VCF)}:\\
  GATK will be used to call single nucleotide variants, producing a VCF
  file such as \texttt{SRR1693691.SNPs.vcf} (or a compressed
  \texttt{SRR1693691.SNPs.vcf.gz}) that lists putative variant sites
  along with quality scores and annotations. You will then apply hard
  filters with GATK (using recommended thresholds and your own
  adjustments) to generate one or more filtered VCFs (for example,
  \texttt{SRR1693691.SNPs.filtered.vcf} and
  \texttt{SRR1693691.SNPs.filtered2.vcf}) and compare how filtering
  choices affect the number and quality of retained variants.
\item
  \textbf{Quality summary files (flagstat, depth, MultiQC, plots)}:\\
  Along the way you will generate summary text files from
  \texttt{samtools\ flagstat}, depth-of-coverage summaries from
  \textbf{vcftools}, and an integrated \textbf{MultiQC} HTML report that
  aggregates QC metrics across the workflow. Custom scripts will also
  create density plots of key variant-quality metrics (for example,
  quality-by-depth or mapping quality) to help you decide on appropriate
  filtering thresholds.
\end{itemize}

\section{How the dataset flows through the
labs}\label{how-the-dataset-flows-through-the-labs}

The genome dataset SRR1693691 is the thread that ties together multiple
labs into a coherent end-to-end workflow, from raw data to
interpretation. Each lab revisits the same sample and reference genome
but focuses on a different set of tools and questions, so that by the
end of the course you will have seen every major step of a standard
variant-calling pipeline in the context of a real experimental evolution
dataset.

\begin{itemize}
\tightlist
\item
  In \textbf{\hyperref[lab04-genome-sequence-indexing-and-alignment]{Lab
  4} (Genome Sequence Indexing and Alignment)}, you will:

  \begin{itemize}
  \tightlist
  \item
    Index the sacCer3 reference genome with BWA and samtools and examine
    the \texttt{.fai} index to calculate total genome size.\\
  \item
    Download the SRR1693691 reads from NCBI using a provided shell
    script, inspect the FASTQ files, and run FastQC before and after
    trimming.\\
  \item
    Trim the reads and align them to sacCer3, producing a sorted BAM
    file plus coverage estimates and initial alignment QC metrics.
  \end{itemize}
\item
  In
  \textbf{\hyperref[lab07-variant-calling-via-gatk-best-practices]{Lab
  7} (Variant Calling via GATK Best Practices)}, you will:

  \begin{itemize}
  \tightlist
  \item
    Start from the aligned BAM and run GATK's post-processing and
    variant-calling steps using a provided script, including duplicate
    marking and variant calling for SNPs.\\
  \item
    Add additional QC steps (such as \texttt{samtools\ flagstat},
    vcftools depth summaries, and MultiQC) and inspect density plots of
    variant-quality metrics to decide how stringent your filters should
    be.
  \item
    Generate multiple filtered VCFs with different cutoff choices and
    compare the resulting sets of SNPs using VCFtools and an UpSet plot,
    learning how filtering decisions trade off between retaining true
    variants and excluding noise.
  \end{itemize}
\item
  In \textbf{\hyperref[lab09-igv-genome-visualization-lab-activity]{Lab
  9} (Genome Visualization with IGV)}, you will:

  \begin{itemize}
  \tightlist
  \item
    Load the sacCer3 genome, the pre-filtering BAM
    (\texttt{SRR1693691.sorted.bam}), the duplicate-marked BAM
    (\texttt{SRR1693691.markdup.bam}), and the VCF into IGV, configuring
    track coloring and display options to highlight read strand,
    duplicates, and amino acid translations.\\
  \item
    Navigate to specific genomic windows that contain called SNPs and
    visually evaluate the evidence for each variant: read depth, base
    quality patterns, strand balance, and consistency across reads.
  \item
    Use coordinate-based searches in both IGV and the filtered VCF to
    decide whether certain variants are likely real biological changes,
    consider whether they fall in protein-coding regions, and reason
    about possible amino acid changes and functional impacts.
  \end{itemize}
\end{itemize}

Throughout these labs, most of the heavy lifting is handled by template
shell scripts that you will lightly edit and submit to the HPC queue, so
you can concentrate on understanding the logic of each step and
interpreting the outputs rather than on writing a full pipeline from
scratch. By repeatedly touching the same dataset at different
stages---raw reads, alignments, unfiltered and filtered variants, and
visual inspection in IGV---you will gain practical experience with how
real genomic data look and behave and how bioinformatics decisions shape
the final set of variants that biologists use to draw conclusions about
evolution and disease.

\phantomsection\label{refs}
\begin{CSLReferences}{0}{0}
\bibitem[\citeproctext]{ref-hotaling2021}
\CSLLeftMargin{1. }%
\CSLRightInline{Hotaling, S., Kelley, J. L. \& Frandsen, P. B.
\href{https://doi.org/10.1073/pnas.2109019118}{Toward a genome sequence
for every animal: Where are we now?} \emph{Proceedings of the National
Academy of Sciences} \textbf{118}, e2109019118 (2021).}

\bibitem[\citeproctext]{ref-marks2021}
\CSLLeftMargin{2. }%
\CSLRightInline{Marks, R. A., Hotaling, S., Frandsen, P. B. \& VanBuren,
R. \href{https://doi.org/10.1038/s41477-021-01031-8}{Representation and
participation across 20 years of plant genome sequencing}. \emph{Nature
Plants} \textbf{7}, 1571--1578 (2021).}

\bibitem[\citeproctext]{ref-bogan2025}
\CSLLeftMargin{3. }%
\CSLRightInline{Bogan, S. N. \emph{et al.}
\href{https://doi.org/10.1093/molbev/msaf219}{Temperature and pressure
shaped the evolution of antifreeze proteins in polar and deep sea
zoarcoid fishes}. \emph{Molecular Biology and Evolution} \textbf{42},
msaf219 (2025).}

\bibitem[\citeproctext]{ref-eddy_what_2004}
\CSLLeftMargin{4. }%
\CSLRightInline{Eddy, S. R.
\href{https://doi.org/10.1038/nbt1004-1315}{What is a hidden {Markov}
model?} \emph{Nature Biotechnology} \textbf{22}, 1315--1316 (2004).}

\bibitem[\citeproctext]{ref-tan_comprehensive_2025}
\CSLLeftMargin{5. }%
\CSLRightInline{Tan, C. \emph{et al.}
\href{https://doi.org/10.1093/bib/bbaf652}{Comprehensive assessment of
whole-genome sequencing-based variant calling pipelines in monozygotic
twins}. \emph{Briefings in Bioinformatics} \textbf{26}, bbaf652 (2025).}

\bibitem[\citeproctext]{ref-abdelwahab_performance_2023}
\CSLLeftMargin{6. }%
\CSLRightInline{Abdelwahab, O., Belzile, F. \& Torkamaneh, D.
\href{https://doi.org/10.1186/s12859-023-05596-3}{Performance analysis
of conventional and {AI}-based variant callers using short and long
reads}. \emph{BMC bioinformatics} \textbf{24}, 472 (2023).}

\bibitem[\citeproctext]{ref-betschart_comparison_2022}
\CSLLeftMargin{7. }%
\CSLRightInline{Betschart, R. O. \emph{et al.}
\href{https://doi.org/10.1038/s41598-022-26181-3}{Comparison of calling
pipelines for whole genome sequencing: An empirical study demonstrating
the importance of mapping and alignment}. \emph{Scientific Reports}
\textbf{12}, 21502 (2022).}

\bibitem[\citeproctext]{ref-pinto_variant_2026}
\CSLLeftMargin{8. }%
\CSLRightInline{Pinto, V., Sousa, L. \& Silva, C.
\href{https://doi.org/10.1371/journal.pone.0339891}{Variant calling in
genomics: {A} comparative performance analysis and decision guide}.
\emph{PloS One} \textbf{21}, e0339891 (2026).}

\bibitem[\citeproctext]{ref-guille_benchmarking_2024}
\CSLLeftMargin{9. }%
\CSLRightInline{Guille, A. \emph{et al.}
\href{https://doi.org/10.1093/bib/bbae697}{A benchmarking study of
individual somatic variant callers and voting-based ensembles for
whole-exome sequencing}. \emph{Briefings in Bioinformatics} \textbf{26},
bbae697 (2024).}

\bibitem[\citeproctext]{ref-pevsner2015}
\CSLLeftMargin{10. }%
\CSLRightInline{Pevsner, J. \emph{Bioinformatics and Functional
Genomics}. (Wiley, 2015).}

\bibitem[\citeproctext]{ref-collins2003}
\CSLLeftMargin{11. }%
\CSLRightInline{Collins, F. S., Green, E. D., Guttmacher, A. E. \&
Guyer, M. S. \href{https://doi.org/10.1038/nature01626}{A vision for the
future of genomics research}. \emph{Nature} \textbf{422}, 835--847
(2003).}

\bibitem[\citeproctext]{ref-akalin2020}
\CSLLeftMargin{12. }%
\CSLRightInline{Akalin, A.
\emph{\href{https://compgenomr.github.io/book/}{Computational Genomics
with r}}. (CRC Press, 2020).}

\bibitem[\citeproctext]{ref-green2011}
\CSLLeftMargin{13. }%
\CSLRightInline{Green, E. D., Guyer, M. S. \& NHGRI.
\href{https://doi.org/10.1038/nature09764}{Charting a course for genomic
medicine from base pairs to bedside}. \emph{Nature} \textbf{470},
204--213 (2011).}

\bibitem[\citeproctext]{ref-yanai_computational_2017}
\CSLLeftMargin{14. }%
\CSLRightInline{Yanai, I. \& Chmielnicki, E.
\href{https://doi.org/10.1186/s13059-017-1357-1}{Computational
biologists: Moving to the driver's seat}. \emph{Genome Biology}
\textbf{18}, 223 (2017).}

\bibitem[\citeproctext]{ref-yanai2020}
\CSLLeftMargin{15. }%
\CSLRightInline{Yanai, I. \& Lercher, M. J.
\href{https://doi.org/10.1186/s13059-020-02133-w}{A hypothesis is a
liability}. \emph{Genome Biology} \textbf{21}, 231 (2020).}

\bibitem[\citeproctext]{ref-Migaloo1}
\CSLLeftMargin{16. }%
\CSLRightInline{Polanowski, A. M., Robinson-Laverick, S. M., Paton, D.
\& Jarman, S. N. \href{https://doi.org/10.1093/jhered/esr108}{Variation
in the {Tyrosinase} {Gene} {Associated} with a {White} {Humpback}
{Whale} ({Megaptera} novaeangliae)}. \emph{Journal of Heredity}
\textbf{103}, 130--133 (2012).}

\bibitem[\citeproctext]{ref-belmont1979}
\CSLLeftMargin{17. }%
\CSLRightInline{National Commission for the Protection of Human Subjects
of Biomedical and Behavioral Research.
\href{https://www.hhs.gov/ohrp/regulations-and-policy/belmont-report/read-the-belmont-report/index.html}{The
belmont report: Ethical principles and guidelines for the protection of
human subjects of research}. \emph{Federal Register} \textbf{44},
23191--23197 (1979).}

\bibitem[\citeproctext]{ref-reverby2009tuskegee}
\CSLLeftMargin{18. }%
\CSLRightInline{Reverby, S. M. Examining tuskegee: The infamous syphilis
study and its legacy. \emph{University of North Carolina Press} (2009).}

\bibitem[\citeproctext]{ref-marks2016hela}
\CSLLeftMargin{19. }%
\CSLRightInline{Marks, J. H.
\href{https://doi.org/10.1146/annurev-genom-083115-022536}{Lessons from
{HeLa} cells: The ethics and policy of biospecimens}. \emph{Annual
Review of Genomics and Human Genetics} \textbf{17}, 395--417 (2016).}

\bibitem[\citeproctext]{ref-skloot2010henrietta}
\CSLLeftMargin{20. }%
\CSLRightInline{Skloot, R. The immortal life of henrietta lacks.
\emph{Crown Publishing Group} (2010).}

\bibitem[\citeproctext]{ref-lacks2022history}
\CSLLeftMargin{21. }%
\CSLRightInline{Washington, H. A.
\href{https://doi.org/10.1089/heq.2021.0151}{Henrietta lacks and
america's dark history of research involving black women}. \emph{Health
Equity} \textbf{6}, e1--e3 (2022).}

\bibitem[\citeproctext]{ref-havasupai2012genomicjustice}
\CSLLeftMargin{22. }%
\CSLRightInline{Garrison, N. A.
\href{https://doi.org/10.1177/0162243912470009}{Genomic justice for
native americans: Impact of the {Havasupai} case on genetic research}.
\emph{Science, Technology, \& Human Values} \textbf{38}, 201--223
(2013).}

\bibitem[\citeproctext]{ref-havasupai2013irb}
\CSLLeftMargin{23. }%
\CSLRightInline{Harmon, A. \emph{et al.}
\href{https://doi.org/10.1080/21507716.2013.770104}{{IRB} and researcher
reflections on the {Havasupai} lawsuit}. \emph{AJOB Primary Research}
\textbf{4}, 55--63 (2013).}

\bibitem[\citeproctext]{ref-burke_standing_2014}
\CSLLeftMargin{24. }%
\CSLRightInline{Burke, M. K., Liti, G. \& Long, A. D.
\href{https://doi.org/10.1093/molbev/msu256}{Standing {Genetic}
{Variation} {Drives} {Repeatable} {Experimental} {Evolution} in
{Outcrossing} {Populations} of {Saccharomyces} cerevisiae}.
\emph{Molecular Biology and Evolution} \textbf{31}, 3228--3239 (2014).}

\bibitem[\citeproctext]{ref-vandijk2018}
\CSLLeftMargin{25. }%
\CSLRightInline{Dijk, E. L. van, Jaszczyszyn, Y., Naquin, D. \& Thermes,
C. \href{https://doi.org/10.1016/j.tig.2018.05.008}{The third revolution
in sequencing technology}. \emph{Trends in Genetics} \textbf{34},
666--681 (2018).}

\bibitem[\citeproctext]{ref-vandijk2023}
\CSLLeftMargin{26. }%
\CSLRightInline{Dijk, E. L. van \emph{et al.}
\href{https://doi.org/10.1016/j.tig.2023.04.006}{Genomics in the
long-read sequencing era}. \emph{Trends in Genetics} \textbf{39},
649--671 (2023).}

\bibitem[\citeproctext]{ref-scarano2024}
\CSLLeftMargin{27. }%
\CSLRightInline{Scarano, G., Nigro, G., Marengo, L., \emph{et al.}
\href{https://doi.org/10.3390/biom14040568}{The third-generation
sequencing challenge: Novel insights for the omic sciences}.
\emph{Biomolecules} \textbf{14}, 568 (2024).}

\bibitem[\citeproctext]{ref-midha_long-read_2019}
\CSLLeftMargin{28. }%
\CSLRightInline{Midha, M. K., Wu, M. \& Chiu, K.-P.
\href{https://doi.org/10.1007/s00439-019-02064-y}{Long-read sequencing
in deciphering human genetics to a greater depth}. \emph{Human Genetics}
\textbf{138}, 1201--1215 (2019).}

\bibitem[\citeproctext]{ref-simao2015}
\CSLLeftMargin{29. }%
\CSLRightInline{Simão, F. A., Waterhouse, R. M., Ioannidis, P.,
Kriventseva, E. V. \& Zdobnov, E. M.
\href{https://doi.org/10.1093/bioinformatics/btv351}{BUSCO: Assessing
genome assembly and annotation completeness with single-copy orthologs}.
\emph{Bioinformatics} \textbf{31}, 3210--3212 (2015).}

\bibitem[\citeproctext]{ref-lipman_rapid_1985}
\CSLLeftMargin{30. }%
\CSLRightInline{Lipman, D. J. \& Pearson, W. R.
\href{https://doi.org/10.1126/science.2983426}{Rapid and sensitive
protein similarity searches}. \emph{Science} \textbf{227}, 1435--1441
(1985).}

\bibitem[\citeproctext]{ref-pearson_improved_1988}
\CSLLeftMargin{31. }%
\CSLRightInline{Pearson, W. R. \& Lipman, D. J.
\href{https://doi.org/10.1073/pnas.85.8.2444}{Improved tools for
biological sequence comparison}. \emph{Proceedings of the National
Academy of Sciences} \textbf{85}, 2444--2448 (1988).}

\bibitem[\citeproctext]{ref-cock2010}
\CSLLeftMargin{32. }%
\CSLRightInline{Cock, P. J. A., Fields, C. J., Goto, N., Heuer, M. L. \&
Rice, P. M. \href{https://doi.org/10.1093/nar/gkp1137}{The {Sanger}
{FASTQ} file format for sequences with quality scores, and the
{Solexa}/{Illumina} {FASTQ} variants}. \emph{Nucleic Acids Research}
\textbf{38}, 1767--1771 (2010).}

\bibitem[\citeproctext]{ref-fastqc}
\CSLLeftMargin{33. }%
\CSLRightInline{Institute, B.
\href{https://www.bioinformatics.babraham.ac.uk/projects/fastqc/}{Babraham
{Bioinformatics} - {FastQC} {A} {Quality} {Control} tool for {High}
{Throughput} {Sequence} {Data}}.}

\bibitem[\citeproctext]{ref-danecek_twelve_2021}
\CSLLeftMargin{34. }%
\CSLRightInline{Danecek, P. \emph{et al.}
\href{https://doi.org/10.1093/gigascience/giab008}{Twelve years of
{SAMtools} and {BCFtools}}. \emph{GigaScience} \textbf{10}, giab008
(2021).}

\bibitem[\citeproctext]{ref-danecek_variant_2011}
\CSLLeftMargin{35. }%
\CSLRightInline{Danecek, P. \emph{et al.}
\href{https://doi.org/10.1093/bioinformatics/btr330}{The variant call
format and {VCFtools}}. \emph{Bioinformatics (Oxford, England)}
\textbf{27}, 2156--2158 (2011).}

\bibitem[\citeproctext]{ref-learning_vcf}
\CSLLeftMargin{36. }%
\CSLRightInline{\href{https://davetang.github.io/learning_vcf_file/}{Learning
the {VCF} format}.}

\bibitem[\citeproctext]{ref-mckenna_genome_2010}
\CSLLeftMargin{37. }%
\CSLRightInline{McKenna, A. \emph{et al.}
\href{https://doi.org/10.1101/gr.107524.110}{The {Genome} {Analysis}
{Toolkit}: A {MapReduce} framework for analyzing next-generation {DNA}
sequencing data}. \emph{Genome Research} \textbf{20}, 1297--1303
(2010).}

\bibitem[\citeproctext]{ref-depristo_framework_2011}
\CSLLeftMargin{38. }%
\CSLRightInline{DePristo, M. A. \emph{et al.}
\href{https://doi.org/10.1038/ng.806}{A framework for variation
discovery and genotyping using next-generation {DNA} sequencing data}.
\emph{Nature Genetics} \textbf{43}, 491--498 (2011).}

\bibitem[\citeproctext]{ref-kent_human_2002}
\CSLLeftMargin{39. }%
\CSLRightInline{Kent, W. J. \emph{et al.}
\href{https://doi.org/10.1101/gr.229102}{The human genome browser at
{UCSC}}. \emph{Genome Research} \textbf{12}, 996--1006 (2002).}

\bibitem[\citeproctext]{ref-quinlan_bedtools_2010}
\CSLLeftMargin{40. }%
\CSLRightInline{Quinlan, A. R. \& Hall, I. M.
\href{https://doi.org/10.1093/bioinformatics/btq033}{{BEDTools}: A
flexible suite of utilities for comparing genomic features}.
\emph{Bioinformatics} \textbf{26}, 841--842 (2010).}

\bibitem[\citeproctext]{ref-kent_bigwig_2010}
\CSLLeftMargin{41. }%
\CSLRightInline{Kent, W. J., Zweig, A. S., Barber, G., Hinrichs, A. S.
\& Karolchik, D.
\href{https://doi.org/10.1093/bioinformatics/btq351}{{BigWig} and
{BigBed}: Enabling browsing of large distributed datasets}.
\emph{Bioinformatics} \textbf{26}, 2204--2207 (2010).}

\bibitem[\citeproctext]{ref-herrero_ensembl_2016}
\CSLLeftMargin{42. }%
\CSLRightInline{Herrero, J. \emph{et al.}
\href{https://doi.org/10.1093/database/bav096}{Ensembl comparative
genomics resources}. \emph{Database: The Journal of Biological Databases
and Curation} \textbf{2016}, bav096 (2016).}

\bibitem[\citeproctext]{ref-aken_ensembl_2017}
\CSLLeftMargin{43. }%
\CSLRightInline{Aken, B. L. \emph{et al.}
\href{https://doi.org/10.1093/nar/gkw1104}{Ensembl 2017}. \emph{Nucleic
Acids Research} \textbf{45}, D635--D642 (2017).}

\bibitem[\citeproctext]{ref-larsson_aliview_2014}
\CSLLeftMargin{44. }%
\CSLRightInline{Larsson, A.
\href{https://doi.org/10.1093/bioinformatics/btu531}{{AliView}: A fast
and lightweight alignment viewer and editor for large datasets}.
\emph{Bioinformatics} \textbf{30}, 3276--3278 (2014).}

\bibitem[\citeproctext]{ref-Goecks2010GalaxyGenomeBiol}
\CSLLeftMargin{45. }%
\CSLRightInline{Goecks, J., Nekrutenko, A., Taylor, J. \& Team, T. G.
\href{https://doi.org/10.1186/gb-2010-11-8-r86}{Galaxy: A comprehensive
approach for supporting accessible, reproducible, and transparent
computational research in the life sciences}. \emph{Genome Biology}
\textbf{11}, R86 (2010).}

\bibitem[\citeproctext]{ref-Blankenberg2010FastqGalaxy}
\CSLLeftMargin{46. }%
\CSLRightInline{Blankenberg, D. \emph{et al.}
\href{https://doi.org/10.1093/bioinformatics/btq281}{Manipulation of
{FASTQ} data with {Galaxy}}. \emph{Bioinformatics} \textbf{26},
1783--1785 (2010).}

\bibitem[\citeproctext]{ref-Blankenberg2010GalaxyExperimentalists}
\CSLLeftMargin{47. }%
\CSLRightInline{Blankenberg, D. \emph{et al.}
\href{https://doi.org/10.1002/0471142727.mb1910s89}{Galaxy: A web-based
genome analysis tool for experimentalists}. \emph{Current Protocols in
Molecular Biology} \textbf{89}, 19.10.1--19.10.21 (2010).}

\bibitem[\citeproctext]{ref-jones_bermuda_2018}
\CSLLeftMargin{48. }%
\CSLRightInline{Jones, K. M., Ankeny, R. A. \& Cook-Deegan, R.
\href{https://doi.org/10.1007/s10739-018-9538-7}{The {Bermuda}
{Triangle}: {The} {Pragmatics}, {Policies}, and {Principles} for {Data}
{Sharing} in the {History} of the {Human} {Genome} {Project}}.
\emph{Journal of the history of biology} \textbf{51}, 693--805 (2018).}

\bibitem[\citeproctext]{ref-nicol_data_2024}
\CSLLeftMargin{49. }%
\CSLRightInline{Nicol, D., Nielsen, J. \& Archer, M.
\href{https://doi.org/10.1038/s41598-024-72653-z}{Data access
arrangements in genomic research consortia}. \emph{Scientific Reports}
\textbf{14}, 21685 (2024).}

\bibitem[\citeproctext]{ref-cook-deegan_moving_2017}
\CSLLeftMargin{50. }%
\CSLRightInline{Cook-Deegan, R. \& McGuire, A. L.
\href{https://doi.org/10.1101/gr.216911.116}{Moving beyond {Bermuda}:
Sharing data to build a medical information commons}. \emph{Genome
Research} \textbf{27}, 897--901 (2017).}

\bibitem[\citeproctext]{ref-cook-deegan_sharing_2017}
\CSLLeftMargin{51. }%
\CSLRightInline{Cook-Deegan, R., Ankeny, R. A. \& Jones, K. M.
\href{https://doi.org/10.1146/annurev-genom-083115-022515}{Sharing
{Data} to {Build} a {Medical} {Information} {Commons}: {From} {Bermuda}
to the {Global} {Alliance}}. \emph{Annual Review of Genomics and Human
Genetics} \textbf{18}, 389--415 (2017).}

\bibitem[\citeproctext]{ref-NIH_GDSP}
\CSLLeftMargin{52. }%
\CSLRightInline{\href{https://grants.nih.gov/grants/guide/notice-files/not-od-14-124.html}{{NOT}-{OD}-14-124:
{NIH} {Genomic} {Data} {Sharing} {Policy}}.}

\bibitem[\citeproctext]{ref-rogers_comparative_2019}
\CSLLeftMargin{53. }%
\CSLRightInline{Rogers, J. \emph{et al.}
\href{https://doi.org/10.1126/sciadv.aau6947}{The comparative genomics
and complex population history of {Papio} baboons}. \emph{Science
Advances} \textbf{5}, eaau6947 (2019).}

\bibitem[\citeproctext]{ref-wall_genomewide_2016}
\CSLLeftMargin{54. }%
\CSLRightInline{Wall, J. D. \emph{et al.}
\href{https://doi.org/10.1111/mec.13684}{Genomewide ancestry and
divergence patterns from low-coverage sequencing data reveal a complex
history of admixture in wild baboons}. \emph{Molecular Ecology}
\textbf{25}, 3469--3483 (2016).}

\bibitem[\citeproctext]{ref-batra_accurate_2020}
\CSLLeftMargin{55. }%
\CSLRightInline{Batra, S. S. \emph{et al.}
\href{https://doi.org/10.1093/gigascience/giaa134}{Accurate assembly of
the olive baboon ({Papio} anubis) genome using long-read and {Hi}-{C}
data}. \emph{GigaScience} \textbf{9}, giaa134 (2020).}

\bibitem[\citeproctext]{ref-park1998}
\CSLLeftMargin{56. }%
\CSLRightInline{Park, J., Teichmann, S. A., Hubbard, T. \& Chothia, C.
\href{https://doi.org/10.1006/jmbi.1998.1712}{Intermediate sequences
increase the detection of homology between sequences}. \emph{Journal of
Molecular Biology} \textbf{278}, 847--867 (1998).}

\bibitem[\citeproctext]{ref-needleman1970}
\CSLLeftMargin{57. }%
\CSLRightInline{Needleman, S. B. \& Wunsch, C. D.
\href{https://doi.org/10.1016/0022-2836(70)90057-4}{A general method
applicable to the search for similarities in the amino acid sequence of
two proteins}. \emph{Journal of Molecular Biology} \textbf{48}, 443--453
(1970).}

\bibitem[\citeproctext]{ref-smith1981}
\CSLLeftMargin{58. }%
\CSLRightInline{Smith, T. F. \& Waterman, M. S.
\href{https://doi.org/10.1016/0022-2836(81)90087-5}{Identification of
common molecular subsequences}. \emph{Journal of Molecular Biology}
\textbf{147}, 195--197 (1981).}

\bibitem[\citeproctext]{ref-dayhoff1978}
\CSLLeftMargin{59. }%
\CSLRightInline{Dayhoff, M. O., Schwartz, R. M. \& Orcutt, B. C. A model
of evolutionary change in proteins. \emph{Atlas of Protein Sequence and
Structure} \textbf{5}, 345--352 (1978).}

\bibitem[\citeproctext]{ref-henikoff1992}
\CSLLeftMargin{60. }%
\CSLRightInline{Henikoff, S. \& Henikoff, J. G.
\href{https://doi.org/10.1073/pnas.89.22.10915}{Amino acid substitution
matrices from protein blocks}. \emph{Proceedings of the National Academy
of Sciences of the USA} \textbf{89}, 10915--10919 (1992).}

\bibitem[\citeproctext]{ref-deborah_thurtle-schmidt_global_2018}
\CSLLeftMargin{61. }%
\CSLRightInline{Deborah Thurtle-Schmidt.
\href{https://www.youtube.com/watch?v=LhpGz5--isw}{Global sequence
alignment}. (2018).}

\bibitem[\citeproctext]{ref-tel-aviv_tutorials_how_2019}
\CSLLeftMargin{62. }%
\CSLRightInline{Tel-Aviv Tutorials.
\href{https://www.youtube.com/watch?v=F1YId_z4pc8}{{HOW} {TO} {Sequence}
{Alignment} ({Local}) - {SMITH} {WATERMAN}}. (2019).}

\bibitem[\citeproctext]{ref-altschul1990}
\CSLLeftMargin{63. }%
\CSLRightInline{Altschul, S. F., Gish, W., Miller, W., Myers, E. W. \&
Lipman, D. J. \href{https://doi.org/10.1016/S0022-2836(05)80360-2}{Basic
local alignment search tool}. \emph{Journal of Molecular Biology}
\textbf{215}, 403--410 (1990).}

\bibitem[\citeproctext]{ref-camacho2009}
\CSLLeftMargin{64. }%
\CSLRightInline{Camacho, C. \emph{et al.}
\href{https://doi.org/10.1186/1471-2105-10-421}{{BLAST+}: Architecture
and applications}. \emph{BMC Bioinformatics} \textbf{10}, 421 (2009).}

\bibitem[\citeproctext]{ref-chothia1986}
\CSLLeftMargin{65. }%
\CSLRightInline{Chothia, C. \& Lesk, A. M. The relation between the
divergence of sequence and structure in proteins. \emph{EMBO Journal}
\textbf{5}, 823--826 (1986).}

\bibitem[\citeproctext]{ref-finn2014}
\CSLLeftMargin{66. }%
\CSLRightInline{Finn, R. D., Bateman, A., Clements, J., \emph{et al.}
\href{https://doi.org/10.1093/nar/gkt1223}{Pfam: The protein families
database}. \emph{Nucleic Acids Research} \textbf{42}, D222--D230
(2014).}

\bibitem[\citeproctext]{ref-letunic2012}
\CSLLeftMargin{67. }%
\CSLRightInline{Letunic, I., Doerks, T. \& Bork, P.
\href{https://doi.org/10.1093/nar/gkr931}{SMART 7: Recent updates to the
protein domain annotation resource}. \emph{Nucleic Acids Research}
\textbf{40}, D302--D305 (2012).}

\bibitem[\citeproctext]{ref-hunter2012}
\CSLLeftMargin{68. }%
\CSLRightInline{Hunter, S., Jones, P., Mitchell, A., \emph{et al.}
\href{https://doi.org/10.1093/nar/gkr948}{InterPro in 2011: New
developments in the family and domain prediction database}.
\emph{Nucleic Acids Research} \textbf{40}, D306--D312 (2012).}

\bibitem[\citeproctext]{ref-marchler2015}
\CSLLeftMargin{69. }%
\CSLRightInline{Marchler-Bauer, A., Derbyshire, M. K., Gonzales, N. R.,
\emph{et al.} \href{https://doi.org/10.1093/nar/gku1221}{CDD: {NCBI}'s
conserved domain database}. \emph{Nucleic Acids Research} \textbf{43},
D222--D226 (2015).}

\bibitem[\citeproctext]{ref-siepel2005}
\CSLLeftMargin{70. }%
\CSLRightInline{Siepel, A., Bejerano, G., Pedersen, J. S., \emph{et al.}
\href{https://doi.org/10.1101/gr.3715005}{Evolutionarily conserved
elements in vertebrate, insect, worm, and yeast genomes}. \emph{Genome
Research} \textbf{15}, 1034--1050 (2005).}

\bibitem[\citeproctext]{ref-pollard2010}
\CSLLeftMargin{71. }%
\CSLRightInline{Pollard, K. S., Hubisz, M. J., Rosenbloom, K. R. \&
Siepel, A. \href{https://doi.org/10.1101/gr.097857.109}{Detection of
nonneutral substitution rates on mammalian phylogenies}. \emph{Genome
Research} \textbf{20}, 110--121 (2010).}

\bibitem[\citeproctext]{ref-davydov2010}
\CSLLeftMargin{72. }%
\CSLRightInline{Davydov, E. V., Goode, D. L., Sirota, M., \emph{et al.}
\href{https://doi.org/10.1371/journal.pcbi.1001025}{Identifying a high
fraction of the human genome to be under selective constraint using
{GERP++}}. \emph{PLoS Computational Biology} \textbf{6}, e1001025
(2010).}

\bibitem[\citeproctext]{ref-blanchette2004}
\CSLLeftMargin{73. }%
\CSLRightInline{Blanchette, M., Kent, W. J., Riemer, C., \emph{et al.}
\href{https://doi.org/10.1101/gr.1933104}{Aligning multiple genomic
sequences with the threaded blockset aligner}. \emph{Genome Research}
\textbf{14}, 708--715 (2004).}

\bibitem[\citeproctext]{ref-paten2008}
\CSLLeftMargin{74. }%
\CSLRightInline{Paten, B., Herrero, J., Beal, M. K., Fitzgerald, S. \&
Birney, E. \href{https://doi.org/10.1101/gr.076554.108}{Enredo and
pecan: Genome-wide mammalian consistency-based multiple alignment with
paralogs}. \emph{Genome Research} \textbf{18}, 1814--1828 (2008).}

\bibitem[\citeproctext]{ref-feng1987}
\CSLLeftMargin{75. }%
\CSLRightInline{Feng, D.-F. \& Doolittle, R. F.
\href{https://doi.org/10.1007/BF02603120}{Progressive sequence alignment
as a prerequisite to correct phylogenetic trees}. \emph{Journal of
Molecular Evolution} \textbf{25}, 351--360 (1987).}

\bibitem[\citeproctext]{ref-feng1990}
\CSLLeftMargin{76. }%
\CSLRightInline{Feng, D.-F. \& Doolittle, R. F.
\href{https://doi.org/10.1016/0076-6879(90)83024-9}{Progressive sequence
alignment and phylogenetic tree construction with the {CLUSTAL}
program}. \emph{Methods in Enzymology} \textbf{183}, 375--387 (1990).}

\bibitem[\citeproctext]{ref-thompson1994}
\CSLLeftMargin{77. }%
\CSLRightInline{Thompson, J. D., Higgins, D. G. \& Gibson, T. J.
\href{https://doi.org/10.1093/nar/22.22.4673}{{CLUSTAL W}: Improving the
sensitivity of progressive multiple sequence alignment through sequence
weighting, position-specific gap penalties and weight matrix choice}.
\emph{Nucleic Acids Research} \textbf{22}, 4673--4680 (1994).}

\bibitem[\citeproctext]{ref-larkin2007}
\CSLLeftMargin{78. }%
\CSLRightInline{Larkin, M. A., Blackshields, G., \emph{et al.}
\href{https://doi.org/10.1093/bioinformatics/btm404}{{Clustal W} and
{Clustal X} version 2.0}. \emph{Bioinformatics} \textbf{23}, 2947--2948
(2007).}

\bibitem[\citeproctext]{ref-edgar2004a}
\CSLLeftMargin{79. }%
\CSLRightInline{Edgar, R. C.
\href{https://doi.org/10.1093/nar/gkh340}{{MUSCLE}: Multiple sequence
alignment with high accuracy and high throughput}. \emph{Nucleic Acids
Research} \textbf{32}, 1792--1797 (2004).}

\bibitem[\citeproctext]{ref-katoh2005}
\CSLLeftMargin{80. }%
\CSLRightInline{Katoh, K., Kuma, K.-I., Toh, H. \& Miyata, T.
\href{https://doi.org/10.1093/nar/gki198}{MAFFT version 5: Improvement
in accuracy of multiple sequence alignment}. \emph{Nucleic Acids
Research} \textbf{33}, 511--518 (2005).}

\bibitem[\citeproctext]{ref-katoh2009}
\CSLLeftMargin{81. }%
\CSLRightInline{Katoh, K. \& Toh, H.
\href{https://doi.org/10.1093/bib/bbn013}{Recent developments in the
{MAFFT} multiple sequence alignment program}. \emph{Briefings in
Bioinformatics} \textbf{9}, 286--298 (2008).}

\bibitem[\citeproctext]{ref-edgar2004b}
\CSLLeftMargin{82. }%
\CSLRightInline{Edgar, R. C.
\href{https://doi.org/10.1186/1471-2105-5-113}{{MUSCLE}: A multiple
sequence alignment method with reduced time and space complexity}.
\emph{BMC Bioinformatics} \textbf{5}, 113 (2004).}

\bibitem[\citeproctext]{ref-notredame2000}
\CSLLeftMargin{83. }%
\CSLRightInline{Notredame, C., Higgins, D. G. \& Heringa, J.
\href{https://doi.org/10.1006/jmbi.2000.4042}{{T-Coffee}: A novel method
for fast and accurate multiple sequence alignment}. \emph{Journal of
Molecular Biology} \textbf{302}, 205--217 (2000).}

\bibitem[\citeproctext]{ref-do2005}
\CSLLeftMargin{84. }%
\CSLRightInline{Do, C. B., Mahabhashyam, M. S. P., Brudno, M. \&
Batzoglou, S. \href{https://doi.org/10.1101/gr.2821705}{ProbCons:
Probabilistic consistency-based multiple sequence alignment}.
\emph{Genome Research} \textbf{15}, 330--340 (2005).}

\bibitem[\citeproctext]{ref-armougom2006}
\CSLLeftMargin{85. }%
\CSLRightInline{Armougom, F. \emph{et al.}
\href{https://doi.org/10.1093/nar/gkl092}{{Expresso}: Automatic
incorporation of structural information in multiple sequence alignments
using {3D}-{Coffee}}. \emph{Nucleic Acids Research} \textbf{34},
W604--W608 (2006).}

\bibitem[\citeproctext]{ref-thompson1999}
\CSLLeftMargin{86. }%
\CSLRightInline{Thompson, J. D., Plewniak, F. \& Poch, O.
\href{https://doi.org/10.1093/nar/27.13.2682}{A comprehensive comparison
of multiple sequence alignment programs}. \emph{Nucleic Acids Research}
\textbf{27}, 2682--2690 (1999).}

\bibitem[\citeproctext]{ref-lassmann2005}
\CSLLeftMargin{87. }%
\CSLRightInline{Lassmann, T. \& Sonnhammer, E. L. L.
\href{https://doi.org/10.1016/j.febslet.2005.11.036}{Quality assessment
of multiple alignment programs}. \emph{FEBS Letters} \textbf{579},
681--685 (2005).}

\bibitem[\citeproctext]{ref-blankenberg2011}
\CSLLeftMargin{88. }%
\CSLRightInline{Blankenberg, D., Taylor, J., Nekrutenko, A. \& Galaxy
Team, the. \href{https://doi.org/10.1093/bioinformatics/btr398}{Making
whole genome multiple alignments usable for biologists}.
\emph{Bioinformatics} \textbf{27}, 2426--2428 (2011).}

\bibitem[\citeproctext]{ref-flicek2014}
\CSLLeftMargin{89. }%
\CSLRightInline{Flicek, P., Amode, M. R., Barrell, D., \emph{et al.}
\href{https://doi.org/10.1093/nar/gkt1196}{Ensembl 2014}. \emph{Nucleic
Acids Research} \textbf{42}, D749--D755 (2014).}

\bibitem[\citeproctext]{ref-Reinert2015AlignmentNGS}
\CSLLeftMargin{90. }%
\CSLRightInline{Reinert, K., Langmead, B., Weese, D. \& Evers, D. J.
\href{https://doi.org/10.1146/annurev-genom-090413-025358}{Alignment of
next-generation sequencing reads}. \emph{Annual Review of Genomics and
Human Genetics} \textbf{16}, 133--151 (2015).}

\bibitem[\citeproctext]{ref-Li2009BWA}
\CSLLeftMargin{91. }%
\CSLRightInline{Li, H. \& Durbin, R.
\href{https://doi.org/10.1093/bioinformatics/btp324}{Fast and accurate
short read alignment with burrows--wheeler transform}.
\emph{Bioinformatics} \textbf{25}, 1754--1760 (2009).}

\bibitem[\citeproctext]{ref-Li2024GenomeAssemblyT2T}
\CSLLeftMargin{92. }%
\CSLRightInline{Li, H. \& Durbin, R.
\href{https://doi.org/10.1038/s41576-024-00718-w}{Genome assembly in the
telomere-to-telomere era}. \emph{Nature Reviews Genetics} \textbf{25},
658--670 (2024).}

\bibitem[\citeproctext]{ref-Wolff2023DataLiteracyGenomeResearch}
\CSLLeftMargin{93. }%
\CSLRightInline{Wolff, K., Friedhoff, R., Schwarzer, F. \& Pucker, B.
\href{https://doi.org/10.1515/jib-2023-0033}{Data literacy in genome
research}. \emph{Journal of Integrative Bioinformatics} \textbf{20},
20230033 (2023).}

\bibitem[\citeproctext]{ref-international2004}
\CSLLeftMargin{94. }%
\CSLRightInline{Consortium, I. H. G. S.
\href{https://doi.org/10.1038/nature03001}{Finishing the euchromatic
sequence of the human genome}. \emph{Nature} \textbf{431}, 931--945
(2004).}

\bibitem[\citeproctext]{ref-Pucker2022PlantAssemblyReview}
\CSLLeftMargin{95. }%
\CSLRightInline{Pucker, B., Irisarri, I., Vries, J. de \& Xu, B.
\href{https://doi.org/10.1017/qpb.2022.5}{Plant genome sequence assembly
in the era of long reads: Progress, challenges and future directions}.
\emph{Quantitative Plant Biology} \textbf{3}, e5 (2022).}

\bibitem[\citeproctext]{ref-Wolff2020ReadTrimmingEvaluation}
\CSLLeftMargin{96. }%
\CSLRightInline{Fabbro, C. D., Scalabrin, S., Morgante, M. \& Giorgi, F.
M. \href{https://doi.org/10.1371/journal.pone.0085024}{An {Extensive}
{Evaluation} of {Read} {Trimming} {Effects} on {Illumina} {NGS} {Data}
{Analysis}}. \emph{PLOS ONE} \textbf{8}, e85024 (2013).}

\bibitem[\citeproctext]{ref-Li2010ReviewNGSAlignment}
\CSLLeftMargin{97. }%
\CSLRightInline{Li, H. \& Homer, N.
\href{https://doi.org/10.1093/bib/bbq015}{A survey of sequence alignment
algorithms for next-generation sequencing}. \emph{Briefings in
Bioinformatics} \textbf{11}, 473--483 (2010).}

\bibitem[\citeproctext]{ref-ben_langmead_burrows-wheeler_2014}
\CSLLeftMargin{98. }%
\CSLRightInline{Ben Langmead.
\href{https://www.youtube.com/watch?v=4n7NPk5lwbI}{Burrows-{Wheeler}
{Transform}}. (2014).}

\bibitem[\citeproctext]{ref-Li2024WhyBwaAlnBlog}
\CSLLeftMargin{99. }%
\CSLRightInline{Li, H. Why is bwa-aln used for ancient DNA reads?
\url{https://lh3.github.io/2024/09/28/why-is-bwa-aln-still-used}
(2024).}

\bibitem[\citeproctext]{ref-Li2009SAMtools}
\CSLLeftMargin{100. }%
\CSLRightInline{Li, H. \emph{et al.}
\href{https://doi.org/10.1093/bioinformatics/btp352}{The sequence
alignment/map format and SAMtools}. \emph{Bioinformatics} \textbf{25},
2078--2079 (2009).}

\bibitem[\citeproctext]{ref-Sahlin2023SurveyLongReadMapping}
\CSLLeftMargin{101. }%
\CSLRightInline{Sahlin, K., Baudeau, T., Cazaux, B. \& Marchet, C.
\href{https://doi.org/10.1186/s13059-023-02972-3}{A survey of mapping
algorithms in the long-reads era}. \emph{Genome Biology} \textbf{24},
(2023).}

\bibitem[\citeproctext]{ref-LoTempio2023BenchmarkLongReadAligners}
\CSLLeftMargin{102. }%
\CSLRightInline{LoTempio, J., Delot, E. \& Vilain, E.
\href{https://doi.org/10.7717/peerj.16515}{Benchmarking long-read genome
sequence alignment tools for human genomics applications}. \emph{PeerJ}
\textbf{11}, e16515 (2023).}

\bibitem[\citeproctext]{ref-Nurk2022T2T}
\CSLLeftMargin{103. }%
\CSLRightInline{Nurk, S. \emph{et al.}
\href{https://doi.org/10.1126/science.abj6987}{The complete sequence of
a human genome}. \emph{Science} \textbf{376}, 44--53 (2022).}

\bibitem[\citeproctext]{ref-noble_quick_2009}
\CSLLeftMargin{104. }%
\CSLRightInline{Noble, W. S.
\href{https://doi.org/10.1371/journal.pcbi.1000424}{A {Quick} {Guide} to
{Organizing} {Computational} {Biology} {Projects}}. \emph{PLOS
Computational Biology} \textbf{5}, e1000424 (2009).}

\bibitem[\citeproctext]{ref-buffalo_bioinformatics}
\CSLLeftMargin{105. }%
\CSLRightInline{Buffalo, V.
\emph{\href{https://www.oreilly.com/library/view/bioinformatics-data-skills/9781449367480/}{Bioinformatics
{Data} {Skills}}}.}

\bibitem[\citeproctext]{ref-mirchandani_fast_2024}
\CSLLeftMargin{106. }%
\CSLRightInline{Mirchandani, C. D. \emph{et al.}
\href{https://doi.org/10.1093/molbev/msad270}{A {Fast}, {Reproducible},
{High}-throughput {Variant} {Calling} {Workflow} for {Population}
{Genomics}}. \emph{Molecular Biology and Evolution} \textbf{41}, msad270
(2024).}

\bibitem[\citeproctext]{ref-rother_toolbox_2012}
\CSLLeftMargin{107. }%
\CSLRightInline{Rother, K. \emph{et al.}
\href{https://doi.org/10.1093/bib/bbr035}{A toolbox for developing
bioinformatics software}. \emph{Briefings in Bioinformatics}
\textbf{13}, 244--257 (2012).}

\bibitem[\citeproctext]{ref-dudley_quick_2009}
\CSLLeftMargin{108. }%
\CSLRightInline{Dudley, J. T. \& Butte, A. J.
\href{https://doi.org/10.1371/journal.pcbi.1000589}{A {Quick} {Guide}
for {Developing} {Effective} {Bioinformatics} {Programming} {Skills}}.
\emph{PLOS Computational Biology} \textbf{5}, e1000589 (2009).}

\bibitem[\citeproctext]{ref-sandve_ten_2013}
\CSLLeftMargin{109. }%
\CSLRightInline{Sandve, G. K., Nekrutenko, A., Taylor, J. \& Hovig, E.
\href{https://doi.org/10.1371/journal.pcbi.1003285}{Ten {Simple} {Rules}
for {Reproducible} {Computational} {Research}}. \emph{PLOS Computational
Biology} \textbf{9}, e1003285 (2013).}

\bibitem[\citeproctext]{ref-ferenc_improving_2024}
\CSLLeftMargin{110. }%
\CSLRightInline{Ferenc, K. \emph{et al.}
\href{https://doi.org/10.1093/bioinformatics/btae632}{Improving
bioinformatics software quality through teamwork}. \emph{Bioinformatics}
\textbf{40}, btae632 (2024).}

\bibitem[\citeproctext]{ref-kallet_how_2004}
\CSLLeftMargin{111. }%
\CSLRightInline{Kallet, R. H. How to write the methods section of a
research paper. \emph{Respiratory Care} \textbf{49}, 1229--1232 (2004).}

\bibitem[\citeproctext]{ref-pop_use_2015}
\CSLLeftMargin{112. }%
\CSLRightInline{Pop, M. \& Salzberg, S. L.
\href{https://doi.org/10.1186/s12859-015-0668-z}{Use and mis-use of
supplementary material in science publications}. \emph{BMC
Bioinformatics} \textbf{16}, 237 (2015).}

\bibitem[\citeproctext]{ref-perkel_democratic_2016}
\CSLLeftMargin{113. }%
\CSLRightInline{Perkel, J.
\href{https://doi.org/10.1038/538127a}{Democratic databases: Science on
{GitHub}}. \emph{Nature} \textbf{538}, 127--128 (2016).}

\bibitem[\citeproctext]{ref-nielsen_genotype_2011}
\CSLLeftMargin{114. }%
\CSLRightInline{Nielsen, R., Paul, J. S., Albrechtsen, A. \& Song, Y. S.
\href{https://doi.org/10.1038/nrg2986}{Genotype and {SNP} calling from
next-generation sequencing data}. \emph{Nat Rev Genet} \textbf{12},
443--51 (2011).}

\bibitem[\citeproctext]{ref-auton_global_2015}
\CSLLeftMargin{115. }%
\CSLRightInline{Auton, A. \emph{et al.}
\href{https://doi.org/10.1038/nature15393}{A global reference for human
genetic variation}. \emph{Nature} \textbf{526}, 68--74 (2015).}

\bibitem[\citeproctext]{ref-sims_sequencing_2014}
\CSLLeftMargin{116. }%
\CSLRightInline{Sims, D., Sudbery, I., Ilott, N. E., Heger, A. \&
Ponting, C. P. \href{https://doi.org/10.1038/nrg3642}{Sequencing depth
and coverage: Key considerations in genomic analyses}. \emph{Nat Rev
Genet} \textbf{15}, 121--32 (2014).}

\bibitem[\citeproctext]{ref-orlando_ancient_2021}
\CSLLeftMargin{117. }%
\CSLRightInline{Orlando, L. \emph{et al.}
\href{https://doi.org/10.1038/s43586-020-00011-0}{Ancient {DNA}
analysis}. \emph{Nature Reviews Methods Primers} \textbf{1}, 14 (2021).}

\bibitem[\citeproctext]{ref-wang_high-coverage_2023}
\CSLLeftMargin{118. }%
\CSLRightInline{Wang, K. \emph{et al.}
\href{https://doi.org/10.1016/j.xgen.2023.100377}{High-coverage genome
of the {Tyrolean} {Iceman} reveals unusually high {Anatolian} farmer
ancestry}. \emph{Cell Genomics} \textbf{3}, 100377 (2023).}

\bibitem[\citeproctext]{ref-mahmoud_structural_2019}
\CSLLeftMargin{119. }%
\CSLRightInline{Mahmoud, M. \emph{et al.}
\href{https://doi.org/10.1186/s13059-019-1828-7}{Structural variant
calling: The long and the short of it}. \emph{Genome Biology}
\textbf{20}, 246 (2019).}

\bibitem[\citeproctext]{ref-whitford_evaluation_2019}
\CSLLeftMargin{120. }%
\CSLRightInline{Whitford, W., Lehnert, K., Snell, R. G. \& Jacobsen, J.
C. \href{https://doi.org/10.1016/j.jbi.2019.103174}{Evaluation of the
performance of copy number variant prediction tools for the detection of
deletions from whole genome sequencing data}. \emph{Journal of
Biomedical Informatics} \textbf{94}, 103174 (2019).}

\bibitem[\citeproctext]{ref-yilmaz2024}
\CSLLeftMargin{121. }%
\CSLRightInline{Yilmaz, F. \emph{et al.}
\href{https://doi.org/10.1126/science.adn0609}{Reconstruction of the
human amylase locus reveals ancient duplications seeding modern-day
variation}. \emph{Science} \textbf{386}, eadn0609 (2024).}

\bibitem[\citeproctext]{ref-yuan_comprehensive_2024}
\CSLLeftMargin{122. }%
\CSLRightInline{Yuan, N. \& Jia, P.
\href{https://doi.org/10.1093/bib/bbae441}{Comprehensive assessment of
long-read sequencing platforms and calling algorithms for detection of
copy number variation}. \emph{Briefings in Bioinformatics} \textbf{25},
bbae441 (2024).}

\bibitem[\citeproctext]{ref-uyar2020}
\CSLLeftMargin{123. }%
\CSLRightInline{Uyar, B.
\emph{\href{https://compgenomr.github.io/book/rnaseqanalysis.html}{Chapter
8 {RNA}-Seq {Analysis} {\textbar} {Computational} {Genomics} with
{R}}}.}

\bibitem[\citeproctext]{ref-dawadi_bench_2025}
\CSLLeftMargin{124. }%
\CSLRightInline{Dawadi, P. \emph{et al.}
\href{https://doi.org/10.3389/fgene.2025.1697922}{From bench to bytes: A
practical guide to {RNA} sequencing data analysis}. \emph{Frontiers in
Genetics} \textbf{16}, 1697922 (2025).}

\bibitem[\citeproctext]{ref-berge_rna_2019}
\CSLLeftMargin{125. }%
\CSLRightInline{Berge, K. V. den \emph{et al.}
\href{https://doi.org/10.1146/annurev-biodatasci-072018-021255}{{RNA}
{Sequencing} {Data}: {Hitchhiker}'s {Guide} to {Expression} {Analysis}}.
\emph{Annual Review of Biomedical Data Science} \textbf{2}, 139--173
(2019).}

\bibitem[\citeproctext]{ref-austin-tse_best_2022}
\CSLLeftMargin{126. }%
\CSLRightInline{Austin-Tse, C. A. \emph{et al.}
\href{https://doi.org/10.1038/s41525-022-00295-z}{Best practices for the
interpretation and reporting of clinical whole genome sequencing}.
\emph{npj Genomic Medicine} \textbf{7}, 27 (2022).}

\bibitem[\citeproctext]{ref-bamshad_exome_2011}
\CSLLeftMargin{127. }%
\CSLRightInline{Bamshad, M. J. \emph{et al.}
\href{https://doi.org/10.1038/nrg3031}{Exome sequencing as a tool for
{Mendelian} disease gene discovery}. \emph{Nature Reviews Genetics}
\textbf{12}, 745--755 (2011).}

\bibitem[\citeproctext]{ref-franke2020}
\CSLLeftMargin{128. }%
\CSLRightInline{Franke, V.
\emph{\href{https://compgenomr.github.io/book/chipseq.html}{Chapter 9
{ChIP}-Seq Analysis {\textbar} {Computational} {Genomics} with {R}}}.}

\bibitem[\citeproctext]{ref-bailey_practical_2013}
\CSLLeftMargin{129. }%
\CSLRightInline{Bailey, T. \emph{et al.}
\href{https://doi.org/10.1371/journal.pcbi.1003326}{Practical guidelines
for the comprehensive analysis of {ChIP}-seq data}. \emph{PLoS
computational biology} \textbf{9}, e1003326 (2013).}

\bibitem[\citeproctext]{ref-singer_practical_2019}
\CSLLeftMargin{130. }%
\CSLRightInline{Singer, B. D.
\href{https://doi.org/10.1165/rcmb.2019-0150TR}{A {Practical} {Guide} to
the {Measurement} and {Analysis} of {DNA} {Methylation}}. \emph{American
Journal of Respiratory Cell and Molecular Biology} \textbf{61}, 417--428
(2019).}

\bibitem[\citeproctext]{ref-gong_analysis_2022}
\CSLLeftMargin{131. }%
\CSLRightInline{Gong, T. \emph{et al.}
\href{https://doi.org/10.1002/smtd.202101251}{Analysis and {Performance}
{Assessment} of the {Whole} {Genome} {Bisulfite} {Sequencing} {Data}
{Workflow}: {Currently} {Available} {Tools} and a {Practical} {Guide} to
{Advance} {DNA} {Methylation} {Studies}}. \emph{Small Methods}
\textbf{6}, e2101251 (2022).}

\bibitem[\citeproctext]{ref-yan_reads_2020}
\CSLLeftMargin{132. }%
\CSLRightInline{Yan, F., Powell, D. R., Curtis, D. J. \& Wong, N. C.
\href{https://doi.org/10.1186/s13059-020-1929-3}{From reads to insight:
A hitchhiker's guide to {ATAC}-seq data analysis}. \emph{Genome Biology}
\textbf{21}, 22 (2020).}

\bibitem[\citeproctext]{ref-grandi_chromatin_2022}
\CSLLeftMargin{133. }%
\CSLRightInline{Grandi, F. C., Modi, H., Kampman, L. \& Corces, M. R.
\href{https://doi.org/10.1038/s41596-022-00692-9}{Chromatin
accessibility profiling by {ATAC}-seq}. \emph{Nature Protocols}
\textbf{17}, 1518--1552 (2022).}

\bibitem[\citeproctext]{ref-shi_fundamental_2022}
\CSLLeftMargin{134. }%
\CSLRightInline{Shi, P. \emph{et al.}
\href{https://doi.org/10.1007/s42994-022-00082-5}{Fundamental and
practical approaches for single-cell {ATAC}-seq analysis}.
\emph{aBIOTECH} \textbf{3}, 212--223 (2022).}

\bibitem[\citeproctext]{ref-huang2022}
\CSLLeftMargin{135. }%
\CSLRightInline{Huang, Y., Shukla, H. \& Lee, Y. C. G.
\href{https://doi.org/10.7554/eLife.81567}{Species-specific chromatin
landscape determines how transposable elements shape genome evolution}.
\emph{eLife} \textbf{11}, e81567 (2022).}

\bibitem[\citeproctext]{ref-amarasinghe_opportunities_2020}
\CSLLeftMargin{136. }%
\CSLRightInline{Amarasinghe, S. L. \emph{et al.}
\href{https://doi.org/10.1186/s13059-020-1935-5}{Opportunities and
challenges in long-read sequencing data analysis}. \emph{Genome Biology}
\textbf{21}, 30 (2020).}

\bibitem[\citeproctext]{ref-liu_practical_2021}
\CSLLeftMargin{137. }%
\CSLRightInline{Liu, Y.-X. \emph{et al.}
\href{https://doi.org/10.1007/s13238-020-00724-8}{A practical guide to
amplicon and metagenomic analysis of microbiome data}. \emph{Protein \&
Cell} \textbf{12}, 315--330 (2021).}

\bibitem[\citeproctext]{ref-ranjan_analysis_2016}
\CSLLeftMargin{138. }%
\CSLRightInline{Ranjan, R., Rani, A., Metwally, A., McGee, H. S. \&
Perkins, D. L.
\href{https://doi.org/10.1016/j.bbrc.2015.12.083}{Analysis of the
microbiome: {Advantages} of whole genome shotgun versus {16S} amplicon
sequencing}. \emph{Biochemical and Biophysical Research Communications}
\textbf{469}, 967--977 (2016).}

\bibitem[\citeproctext]{ref-joseph_introduction_2021}
\CSLLeftMargin{139. }%
\CSLRightInline{Joseph, T. A. \& Pe'er, I.
\href{https://doi.org/10.1007/978-1-0716-1103-6_6}{An {Introduction} to
{Whole}-{Metagenome} {Shotgun} {Sequencing} {Studies}}. \emph{Methods in
Molecular Biology} \textbf{2243}, 107--122 (2021).}

\bibitem[\citeproctext]{ref-proctor_integrative_2019}
\CSLLeftMargin{140. }%
\CSLRightInline{Proctor, L. M. \emph{et al.}
\href{https://doi.org/10.1038/s41586-019-1238-8}{The {Integrative}
{Human} {Microbiome} {Project}}. \emph{Nature} \textbf{569}, 641--648
(2019).}

\bibitem[\citeproctext]{ref-turnbaugh_human_2007}
\CSLLeftMargin{141. }%
\CSLRightInline{Turnbaugh, P. J. \emph{et al.}
\href{https://doi.org/10.1038/nature06244}{The {Human} {Microbiome}
{Project}}. \emph{Nature} \textbf{449}, 804--810 (2007).}

\bibitem[\citeproctext]{ref-gligorijevic_methods_2015}
\CSLLeftMargin{142. }%
\CSLRightInline{Gligorijević, V. \& Pržulj, N.
\href{https://doi.org/10.1098/rsif.2015.0571}{Methods for biological
data integration: Perspectives and challenges}. \emph{Journal of The
Royal Society Interface} \textbf{12}, 20150571 (2015).}

\bibitem[\citeproctext]{ref-athieniti_guide_2023}
\CSLLeftMargin{143. }%
\CSLRightInline{Athieniti, E. \& Spyrou, G. M.
\href{https://doi.org/10.1016/j.csbj.2022.11.050}{A guide to multi-omics
data collection and integration for translational medicine}.
\emph{Computational and Structural Biotechnology Journal} \textbf{21},
134--149 (2023).}

\bibitem[\citeproctext]{ref-ronen2020}
\CSLLeftMargin{144. }%
\CSLRightInline{Ronen, J.
\emph{\href{https://compgenomr.github.io/book/multiomics.html}{Chapter
11 Multi-Omics Analysis {\textbar} {Computational} {Genomics} with
{R}}}.}

\bibitem[\citeproctext]{ref-yoon_hidden_2009}
\CSLLeftMargin{145. }%
\CSLRightInline{Yoon, B.-J.
\href{https://doi.org/10.2174/138920209789177575}{Hidden {Markov}
{Models} and their {Applications} in {Biological} {Sequence}
{Analysis}}. \emph{Current Genomics} \textbf{10}, 402--415 (2009).}

\bibitem[\citeproctext]{ref-stanke_augustus_2005}
\CSLLeftMargin{146. }%
\CSLRightInline{Stanke, M. \& Morgenstern, B.
\href{https://doi.org/10.1093/nar/gki458}{{AUGUSTUS}: A web server for
gene prediction in eukaryotes that allows user-defined constraints}.
\emph{Nucleic Acids Research} \textbf{33}, W465--467 (2005).}

\bibitem[\citeproctext]{ref-brent_how_2007}
\CSLLeftMargin{147. }%
\CSLRightInline{Brent, M. R.
\href{https://doi.org/10.1038/nbt0807-883}{How does eukaryotic gene
prediction work?} \emph{Nature Biotechnology} \textbf{25}, 883--885
(2007).}

\bibitem[\citeproctext]{ref-weisstein_hands-introduction_2021}
\CSLLeftMargin{148. }%
\CSLRightInline{Weisstein, A. E. \emph{et al.} A {Hands}-on
{Introduction} to {Hidden} {Markov} {Models}.
\url{https://doi.org/10.24918/cs.2016.8} (2021)
doi:\href{https://doi.org/10.24918/cs.2016.8}{10.24918/cs.2016.8}.}

\bibitem[\citeproctext]{ref-sun_opportunities_2025}
\CSLLeftMargin{149. }%
\CSLRightInline{Sun, Q. \emph{et al.}
\href{https://doi.org/10.1016/j.ajhg.2025.03.004}{Opportunities and
challenges of local ancestry in genetic association analyses}. \emph{The
American Journal of Human Genetics} \textbf{112}, 727--740 (2025).}

\bibitem[\citeproctext]{ref-iasi_neanderthal_2024}
\CSLLeftMargin{150. }%
\CSLRightInline{Iasi, L. N. M. \emph{et al.}
\href{https://doi.org/10.1126/science.adq3010}{Neanderthal ancestry
through time: {Insights} from genomes of ancient and present-day
humans}. \emph{Science} \textbf{386}, eadq3010 (2024).}

\bibitem[\citeproctext]{ref-sumer_earliest_2025}
\CSLLeftMargin{151. }%
\CSLRightInline{Sümer, A. P. \emph{et al.}
\href{https://doi.org/10.1038/s41586-024-08420-x}{Earliest modern human
genomes constrain timing of {Neanderthal} admixture}. \emph{Nature}
\textbf{638}, 711--717 (2025).}

\bibitem[\citeproctext]{ref-bqsr_2024}
\CSLLeftMargin{152. }%
\CSLRightInline{\href{https://gatk.broadinstitute.org/hc/en-us/articles/360035890531-Base-Quality-Score-Recalibration-BQSR}{Base
{Quality} {Score} {Recalibration} ({BQSR})}. \emph{GATK} (2024).}

\bibitem[\citeproctext]{ref-haplotypecaller_2025}
\CSLLeftMargin{153. }%
\CSLRightInline{\href{https://gatk.broadinstitute.org/hc/en-us/articles/360037225632-HaplotypeCaller}{{HaplotypeCaller}}.
\emph{GATK} (2025).}

\bibitem[\citeproctext]{ref-logic_2024}
\CSLLeftMargin{154. }%
\CSLRightInline{\href{https://gatk.broadinstitute.org/hc/en-us/articles/360035890431-The-logic-of-joint-calling-for-germline-short-variants}{The
logic of joint calling for germline short variants}. \emph{GATK}
(2024).}

\bibitem[\citeproctext]{ref-gvcf_2024}
\CSLLeftMargin{155. }%
\CSLRightInline{\href{https://gatk.broadinstitute.org/hc/en-us/articles/360035531812-GVCF-Genomic-Variant-Call-Format}{{GVCF}
- {Genomic} {Variant} {Call} {Format}}. \emph{GATK} (2024).}

\bibitem[\citeproctext]{ref-how_2025}
\CSLLeftMargin{156. }%
\CSLRightInline{\href{https://gatk.broadinstitute.org/hc/en-us/articles/360035531112--How-to-Filter-variants-either-with-VQSR-or-by-hard-filtering}{({How}
to) {Filter} variants either with {VQSR} or by hard-filtering}.
\emph{GATK} (2025).}

\bibitem[\citeproctext]{ref-variant_2024}
\CSLLeftMargin{157. }%
\CSLRightInline{\href{https://gatk.broadinstitute.org/hc/en-us/articles/360035531612-Variant-Quality-Score-Recalibration-VQSR}{Variant
{Quality} {Score} {Recalibration} ({VQSR})}. \emph{GATK} (2024).}

\bibitem[\citeproctext]{ref-hard-filtering_2024}
\CSLLeftMargin{158. }%
\CSLRightInline{\href{https://gatk.broadinstitute.org/hc/en-us/articles/360035890471-Hard-filtering-germline-short-variants}{Hard-filtering
germline short variants}. \emph{GATK} (2024).}

\bibitem[\citeproctext]{ref-Robinson2011IGV}
\CSLLeftMargin{159. }%
\CSLRightInline{Robinson, J. T. \emph{et al.}
\href{https://doi.org/10.1038/nbt.1754}{Integrative genomics viewer}.
\emph{Nature Biotechnology} \textbf{29}, 24--26 (2011).}

\bibitem[\citeproctext]{ref-meynert_variant_2014}
\CSLLeftMargin{160. }%
\CSLRightInline{Meynert, A. M., Ansari, M., FitzPatrick, D. R. \&
Taylor, M. S. \href{https://doi.org/10.1186/1471-2105-15-247}{Variant
detection sensitivity and biases in whole genome and exome sequencing}.
\emph{BMC bioinformatics} \textbf{15}, 247 (2014).}

\bibitem[\citeproctext]{ref-cheng_assessing_2014}
\CSLLeftMargin{161. }%
\CSLRightInline{Cheng, A. Y., Teo, Y.-Y. \& Ong, R. T.-H.
\href{https://doi.org/10.1093/bioinformatics/btu067}{Assessing single
nucleotide variant detection and genotype calling on whole-genome
sequenced individuals}. \emph{Bioinformatics} \textbf{30}, 1707--1713
(2014).}

\bibitem[\citeproctext]{ref-wall_estimating_2014}
\CSLLeftMargin{162. }%
\CSLRightInline{Wall, J. D. \emph{et al.}
\href{https://doi.org/10.1101/gr.168393.113}{Estimating genotype error
rates from high-coverage next-generation sequence data}. \emph{Genome
Research} \textbf{24}, 1734--1739 (2014).}

\bibitem[\citeproctext]{ref-breton_comparison_2021}
\CSLLeftMargin{163. }%
\CSLRightInline{Breton, G., Johansson, A. C. V., Sjödin, P., Schlebusch,
C. M. \& Jakobsson, M.
\href{https://doi.org/10.1186/s12859-021-04407-x}{Comparison of
sequencing data processing pipelines and application to underrepresented
{African} human populations}. \emph{BMC bioinformatics} \textbf{22}, 488
(2021).}

\bibitem[\citeproctext]{ref-barbitoff_systematic_2022}
\CSLLeftMargin{164. }%
\CSLRightInline{Barbitoff, Y. A., Abasov, R., Tvorogova, V. E., Glotov,
A. S. \& Predeus, A. V.
\href{https://doi.org/10.1186/s12864-022-08365-3}{Systematic benchmark
of state-of-the-art variant calling pipelines identifies major factors
affecting accuracy of coding sequence variant discovery}. \emph{BMC
genomics} \textbf{23}, 155 (2022).}

\bibitem[\citeproctext]{ref-lefouili_evaluation_2022}
\CSLLeftMargin{165. }%
\CSLRightInline{Lefouili, M. \& Nam, K.
\href{https://doi.org/10.1038/s41598-022-15563-2}{The evaluation of
{Bcftools} mpileup and {GATK} {HaplotypeCaller} for variant calling in
non-human species}. \emph{Scientific Reports} \textbf{12}, 11331
(2022).}

\bibitem[\citeproctext]{ref-lam_performance_2012}
\CSLLeftMargin{166. }%
\CSLRightInline{Lam, H. Y. K. \emph{et al.}
\href{https://doi.org/10.1038/nbt.2065}{Performance comparison of
whole-genome sequencing platforms}. \emph{Nature Biotechnology}
\textbf{30}, 78--82 (2012).}

\bibitem[\citeproctext]{ref-poplin_universal_2018}
\CSLLeftMargin{167. }%
\CSLRightInline{Poplin, R. \emph{et al.}
\href{https://doi.org/10.1038/nbt.4235}{A universal {SNP} and
small-indel variant caller using deep neural networks}. \emph{Nature
Biotechnology} \textbf{36}, 983--987 (2018).}

\bibitem[\citeproctext]{ref-noauthor_history_2025}
\CSLLeftMargin{168. }%
\CSLRightInline{\href{https://www.nih.gov/about-nih/organization/history}{History
{\textbar} {National} {Institutes} of {Health} ({NIH})}. (2025).}

\bibitem[\citeproctext]{ref-noauthor_national_nodate}
\CSLLeftMargin{169. }%
\CSLRightInline{Mazuzan, G. T.
\href{https://www.nsf.gov/about/history/narrative}{The {National}
{Science} {Foundation}: {A} {Brief} {History}}. (1994).}

\bibitem[\citeproctext]{ref-noauthor_nsf_history_2024}
\CSLLeftMargin{170. }%
\CSLRightInline{\href{https://www.nsf.gov/about/history}{History - about
{NSF}}. (2024).}

\bibitem[\citeproctext]{ref-daniel_bolnick_asn_2025}
\CSLLeftMargin{171. }%
\CSLRightInline{Daniel Bolnick.
\href{https://www.youtube.com/watch?v=soPdRtN67hQ}{{ASN} {Presidential}
{Address} - {EVOLUTION} 2025}. (2025).}

\bibitem[\citeproctext]{ref-king_science_2025}
\CSLLeftMargin{172. }%
\CSLRightInline{King, R. Science {Spending} {Is} an {Investment}.
\emph{The Physiologist Magazine}
\url{https://www.physiology.org/publications/news/the-physiologist-magazine/policy-iq/science-spending-is-an-investment}
(2025).}

\bibitem[\citeproctext]{ref-smith_overview_2022}
\CSLLeftMargin{173. }%
\CSLRightInline{Smith, D., Chowdhury, A. S. \& Oxford, J. T.
\href{https://doi.org/10.1002/cpz1.545}{Overview on {Grant} {Writing}
for {Graduate} {Student} {Research}}. \emph{Current Protocols}
\textbf{2}, e545 (2022).}

\bibitem[\citeproctext]{ref-bhattacharya_how_2024}
\CSLLeftMargin{174. }%
\CSLRightInline{Bhattacharya, S. \& Saha, V.
\href{https://doi.org/10.1016/j.ijmmb.2023.100482}{How to write a
research grant proposal}. \emph{Indian Journal of Medical Microbiology}
\textbf{49}, 100482 (2024).}

\bibitem[\citeproctext]{ref-gilpin_nih_nodate}
\CSLLeftMargin{175. }%
\CSLRightInline{Gilpin, N. W.
\href{https://doi.org/10.7554/eLife.106710}{The {NIH} is a sound
investment for the {US} taxpayer}. \emph{eLife} \textbf{14}, e106710
(2025).}

\bibitem[\citeproctext]{ref-hegde_politics_2008}
\CSLLeftMargin{176. }%
\CSLRightInline{Hegde, D. \& Mowery, D. C.
\href{https://doi.org/10.1126/science.1158562}{Politics and {Funding} in
the {U}.{S}. {Public} {Biomedical} {R}\&{D} {System}}. \emph{Science}
\textbf{322}, 1797--1798 (2008).}

\bibitem[\citeproctext]{ref-read_analysis_2025}
\CSLLeftMargin{177. }%
\CSLRightInline{Read, 5. M.
\href{https://stories.tamu.edu/news/2025/06/20/analysis-government-funded-rd-leads-to-long-term-economic-growth}{Analysis:
{Government}-{Funded} {R}\&{D} {Leads} {To} {Long}-{Term} {Economic}
{Growth}}. \emph{Texas A\&M Stories} (2025).}

\bibitem[\citeproctext]{ref-preuss_return_2016}
\CSLLeftMargin{178. }%
\CSLRightInline{Preuss, M.
\href{https://files.eric.ed.gov/fulltext/EJ1134098.pdf}{Return {On}
{Investment} and {Grants}: {A} {Review} of {Present} {Understandings}
and {Recommendations} {For} {Change}}. \emph{Research Management Review}
\textbf{21}, (2016).}

\bibitem[\citeproctext]{ref-holloway_proposal_2003}
\CSLLeftMargin{179. }%
\CSLRightInline{Holloway, B. R. \emph{Proposal {Writing} {Across} the
{Disciplines}}. (Pearson, Upper Saddle River, NJ, 2003).}

\bibitem[\citeproctext]{ref-noauthor_guide_nodate}
\CSLLeftMargin{180. }%
\CSLRightInline{\href{https://www.learnerassociates.net/proposal/}{Guide
for {Writing} a {Funding} {Proposal}}.}

\bibitem[\citeproctext]{ref-noauthor_grant_nodate}
\CSLLeftMargin{181. }%
\CSLRightInline{\href{https://writingcenter.unc.edu/tips-and-tools/grant-proposals-or-give-me-the-money/}{Grant
{Proposals} (or {Give} me the money!)}. \emph{The Writing Center}.}

\bibitem[\citeproctext]{ref-noauthor_write_nodate}
\CSLLeftMargin{182. }%
\CSLRightInline{\href{https://catalyst.harvard.edu/writing-communication-center/get-funded/write-your-grant-proposal/}{Write
{Your} {Grant} {Proposal}}.}

\bibitem[\citeproctext]{ref-nih_funding_nih_2023}
\CSLLeftMargin{183. }%
\CSLRightInline{NIH Funding.
\href{https://www.youtube.com/watch?v=fAFth5aiBpU}{{NIH} {Peer}
{Review}: {Mock} {Study} {Section}}. (2023).}

\bibitem[\citeproctext]{ref-rori_matthew_2025}
\CSLLeftMargin{184. }%
\CSLRightInline{(RoRI), R. on R. I. \emph{et al.} The {Matthew} effect
and early-career setbacks in research funding - a replication study
({RoRI} {Working} {Paper} {No}. 16). (2025)
doi:\href{https://doi.org/10.6084/m9.figshare.29302004.v2}{10.6084/m9.figshare.29302004.v2}.}

\bibitem[\citeproctext]{ref-bol_matthew_2018}
\CSLLeftMargin{185. }%
\CSLRightInline{Bol, T., Vaan, M. de \& Rijt, A. van de.
\href{https://doi.org/10.1073/pnas.1719557115}{The {Matthew} effect in
science funding}. \emph{Proceedings of the National Academy of Sciences}
\textbf{115}, 4887--4890 (2018).}

\bibitem[\citeproctext]{ref-kumar_lessons_2026}
\CSLLeftMargin{186. }%
\CSLRightInline{Kumar, J., Worosz, M. R., Shiu, S.-H., Templer, P. H. \&
McNeal, K. S.
\href{https://doi.org/10.1371/journal.pone.0343307}{Lessons learned from
interdisciplinary {US} national science foundation research
traineeship-supported graduate programs}. \emph{PLOS ONE} \textbf{21},
e0343307 (2026).}

\bibitem[\citeproctext]{ref-noauthor_survey_nodate}
\CSLLeftMargin{187. }%
\CSLRightInline{\href{https://ncses.nsf.gov/surveys/graduate-students-postdoctorates-s-e/2024}{Survey
of {Graduate} {Students} and {Postdoctorates} in {Science} and
{Engineering} ({GSS}) 2024 {\textbar} {NCSES} {\textbar} {NSF}}.}

\bibitem[\citeproctext]{ref-ellegren_genome_2014}
\CSLLeftMargin{188. }%
\CSLRightInline{Ellegren, H.
\href{https://doi.org/10.1016/j.tree.2013.09.008}{Genome sequencing and
population genomics in non-model organisms}. \emph{Trends in Ecology \&
Evolution} \textbf{29}, 51--63 (2014).}

\bibitem[\citeproctext]{ref-coop_popgen_notes}
\CSLLeftMargin{189. }%
\CSLRightInline{Coop, G.
\href{https://github.com/cooplab/popgen-notes/releases}{Population and
quantitative genetics, 3rd ed.} (2020).}

\bibitem[\citeproctext]{ref-charlesworth_hubby_2016}
\CSLLeftMargin{190. }%
\CSLRightInline{Charlesworth, B.
\href{https://doi.org/10.1038/hdy.2015.102}{Is the {Öld Riddlë} of {H}.
J. Muller and j. L. Hubby solved?} \emph{Heredity} \textbf{116},
139--143 (2016).}

\bibitem[\citeproctext]{ref-rutherford_cautionary_2021}
\CSLLeftMargin{191. }%
\CSLRightInline{Rutherford, A.
\href{https://doi.org/10.1126/science.abm4415}{A cautionary history of
eugenics}. \emph{Science} \textbf{373}, 1419--1419 (2021).}

\bibitem[\citeproctext]{ref-NIH_eugenics}
\CSLLeftMargin{192. }%
\CSLRightInline{\href{https://www.genome.gov/about-genomics/fact-sheets/Eugenics-and-Scientific-Racism}{Eugenics
and {Scientific} {Racism}}.}

\bibitem[\citeproctext]{ref-stern_californias_2017}
\CSLLeftMargin{193. }%
\CSLRightInline{Stern, A. M. \emph{et al.}
\href{https://doi.org/10.2105/AJPH.2016.303489}{California's
{Sterilization} {Survivors}: {An} {Estimate} and {Call} for {Redress}}.
\emph{American Journal of Public Health} \textbf{107}, 50--54 (2017).}

\bibitem[\citeproctext]{ref-theobald_1970_2019}
\CSLLeftMargin{194. }%
\CSLRightInline{Theobald, B.
\href{https://time.com/5737080/native-american-sterilization-history/}{A
1970 {Law} {Led} to the {Mass} {Sterilization} of {Native} {American}
{Women}. {That} {History} {Still} {Matters}}. \emph{TIME} (2019).}

\bibitem[\citeproctext]{ref-johnson_female_2014}
\CSLLeftMargin{195. }%
\CSLRightInline{Johnson, C. G.
\href{https://revealnews.org/article/female-prison-inmates-sterilized-illegally-california-audit-confirms/}{Female
prison inmates sterilized illegally, {California} audit confirms}.
\emph{Reveal} (2014).}

\bibitem[\citeproctext]{ref-hawkins_tenn_2017}
\CSLLeftMargin{196. }%
\CSLRightInline{Hawkins, D. Tenn. Judge reprimanded for offering reduced
jail time in exchange for sterilization. \emph{The Washington Post}
\url{https://www.washingtonpost.com/news/morning-mix/wp/2017/11/21/tenn-judge-reprimanded-for-offering-reduced-jail-time-in-exchange-for-sterilization/}
(2017).}

\bibitem[\citeproctext]{ref-allen_eugenics_1986}
\CSLLeftMargin{197. }%
\CSLRightInline{Allen, G. E.
\href{https://www.jstor.org/stable/301835}{The {Eugenics} {Record}
{Office} at {Cold} {Spring} {Harbor}, 1910-1940: {An} {Essay} in
{Institutional} {History}}. \emph{Osiris} \textbf{2}, 225--264 (1986).}

\bibitem[\citeproctext]{ref-bodmer_outstanding_2021}
\CSLLeftMargin{198. }%
\CSLRightInline{Bodmer, W. \emph{et al.}
\href{https://doi.org/10.1038/s41437-020-00394-6}{The outstanding
scientist, {R}.{A}. {Fisher}: His views on eugenics and race}.
\emph{Heredity} \textbf{126}, 565--576 (2021).}

\bibitem[\citeproctext]{ref-muller_our_1950}
\CSLLeftMargin{199. }%
\CSLRightInline{Muller, H. J.
\href{https://pmc.ncbi.nlm.nih.gov/articles/PMC1716299/}{Our load of
mutations}. \emph{American Journal of Human Genetics} \textbf{2},
111--176 (1950).}

\bibitem[\citeproctext]{ref-mazumdar_reform_2002}
\CSLLeftMargin{200. }%
\CSLRightInline{Mazumdar, P. M. H.
\href{https://doi.org/10.1016/S0168-9525(01)02551-3}{{``{Reform}''}
eugenics and the decline of {Mendelism}}. \emph{Trends in Genetics}
\textbf{18}, 48--52 (2002).}

\bibitem[\citeproctext]{ref-lewontin_apportionment_1972}
\CSLLeftMargin{201. }%
\CSLRightInline{Lewontin, R. C. The {Apportionment} of {Human}
{Diversity}. in \emph{Evolutionary {Biology}: {Volume} 6} (eds.
Dobzhansky, T., Hecht, M. K. \& Steere, W. C.) 381--398 (Springer US,
New York, NY, 1972).
doi:\href{https://doi.org/10.1007/978-1-4684-9063-3_14}{10.1007/978-1-4684-9063-3\_14}.}

\bibitem[\citeproctext]{ref-UNESCO_race}
\CSLLeftMargin{202. }%
\CSLRightInline{\href{https://unesdoc.unesco.org/ark:/48223/pf0000128291}{The
{Race} question - {UNESCO} {Digital} {Library}}.}

\bibitem[\citeproctext]{ref-donovan_ending_2021}
\CSLLeftMargin{203. }%
\CSLRightInline{Donovan, B. M.
\href{https://doi.org/10.1016/j.xhgg.2021.100058}{Ending genetic
essentialism through genetics education}. \emph{Human Genetics and
Genomics Advances} \textbf{3}, 100058 (2021).}

\bibitem[\citeproctext]{ref-paul_eugenics_1997}
\CSLLeftMargin{204. }%
\CSLRightInline{Paul, D.
\href{https://doi.org/10.1017/S0898030600005844}{From {Eugenics} to
{Medical} {Genetics}}. \emph{Journal of Policy History} \textbf{9},
96--116 (1997).}

\bibitem[\citeproctext]{ref-leffler_revisiting_2012}
\CSLLeftMargin{205. }%
\CSLRightInline{Leffler, E. M. \emph{et al.}
\href{https://doi.org/10.1371/journal.pbio.1001388}{Revisiting an old
riddle: What determines genetic diversity levels within species?}
\emph{PLoS Biology} \textbf{10}, e1001388 (2012).}

\bibitem[\citeproctext]{ref-schaeffer_muller_2018}
\CSLLeftMargin{206. }%
\CSLRightInline{Schaeffer, S. W.
\href{https://doi.org/10.1534/genetics.118.301084}{Muller
{``{Elements}''} in {Drosophila}: {How} the {Search} for the {Genetic}
{Basis} for {Speciation} {Led} to the {Birth} of {Comparative}
{Genomics}}. \emph{Genetics} \textbf{210}, 3--13 (2018).}

\bibitem[\citeproctext]{ref-haasl_payseur_2016}
\CSLLeftMargin{207. }%
\CSLRightInline{Haasl, R. J. \& Payseur, B. A.
\href{https://doi.org/10.1111/mec.13339}{Fifteen years of genomewide
scans for selection: Trends, lessons and unaddressed genetic sources of
complication}. \emph{Molecular Ecology} \textbf{25}, 5--23 (2016).}

\bibitem[\citeproctext]{ref-sethuraman_missing_2026}
\CSLLeftMargin{208. }%
\CSLRightInline{Sethuraman, A.
\href{https://doi.org/10.1093/g3journal/jkaf269}{The missing data
problem in population genomics and statistical methods to address them}.
\emph{G3: Genes, Genomes, Genetics} \textbf{16}, jkaf269 (2026).}

\bibitem[\citeproctext]{ref-bhutkar_chromosomal_2008}
\CSLLeftMargin{209. }%
\CSLRightInline{Bhutkar, A. \emph{et al.}
\href{https://doi.org/10.1534/genetics.107.086108}{Chromosomal
{Rearrangement} {Inferred} {From} {Comparisons} of 12 {Drosophila}
{Genomes}}. \emph{Genetics} \textbf{179}, 1657--1680 (2008).}

\bibitem[\citeproctext]{ref-carey_ten_2020}
\CSLLeftMargin{210. }%
\CSLRightInline{Carey, M. A., Steiner, K. L. \& Petri, W. A.
\href{https://doi.org/10.1371/journal.pcbi.1008032}{Ten simple rules for
reading a scientific paper}. \emph{PLoS Computational Biology}
\textbf{16}, e1008032 (2020).}

\bibitem[\citeproctext]{ref-baudat_prdm9_2010}
\CSLLeftMargin{211. }%
\CSLRightInline{Baudat, F. \emph{et al.}
\href{https://doi.org/10.1126/science.1183439}{{PRDM9} {Is} a {Major}
{Determinant} of {Meiotic} {Recombination} {Hotspots} in {Humans} and
{Mice}}. \emph{Science} \textbf{327}, 836--840 (2010).}

\bibitem[\citeproctext]{ref-paigen_prdm9_2018}
\CSLLeftMargin{212. }%
\CSLRightInline{Paigen, K. \& Petkov, P. M.
\href{https://doi.org/10.1016/j.tig.2017.12.017}{{PRDM9} and {Its}
{Role} in {Genetic} {Recombination}}. \emph{Trends in Genetics}
\textbf{34}, 291--300 (2018).}

\bibitem[\citeproctext]{ref-grey_prdm9_2018}
\CSLLeftMargin{213. }%
\CSLRightInline{Grey, C., Baudat, F. \& Massy, B. de.
\href{https://doi.org/10.1371/journal.pgen.1007479}{{PRDM9}, a driver of
the genetic map}. \emph{PLOS Genetics} \textbf{14}, e1007479 (2018).}

\bibitem[\citeproctext]{ref-schwartz_primate_2014}
\CSLLeftMargin{214. }%
\CSLRightInline{Schwartz, J. J., Roach, D. J., Thomas, J. H. \&
Shendure, J. \href{https://doi.org/10.1038/ncomms5370}{Primate evolution
of the recombination regulator {PRDM9}}. \emph{Nature Communications}
\textbf{5}, 4370 (2014).}

\bibitem[\citeproctext]{ref-hoge_patterns_2024}
\CSLLeftMargin{215. }%
\CSLRightInline{Hoge, C. \emph{et al.}
\href{https://doi.org/10.1126/science.adj7026}{Patterns of recombination
in snakes reveal a tug-of-war between {PRDM9} and promoter-like
features}. \emph{Science} \textbf{383}, eadj7026 (2024).}

\bibitem[\citeproctext]{ref-baker_repeated_2017}
\CSLLeftMargin{216. }%
\CSLRightInline{Baker, Z. \emph{et al.}
\href{https://doi.org/10.7554/eLife.24133}{Repeated losses of
{PRDM9}-directed recombination despite the conservation of {PRDM9}
across vertebrates}. \emph{eLife} \textbf{6}, e24133 (2017).}

\bibitem[\citeproctext]{ref-cavassim_prdm9_2022}
\CSLLeftMargin{217. }%
\CSLRightInline{Cavassim, M. I. A. \emph{et al.}
\href{https://doi.org/10.1073/pnas.2114401119}{{PRDM9} losses in
vertebrates are coupled to those of paralogs {ZCWPW1} and {ZCWPW2}}.
\emph{Proceedings of the National Academy of Sciences of the United
States of America} \textbf{119}, e2114401119 (2022).}

\bibitem[\citeproctext]{ref-kosch_amphibian_2024}
\CSLLeftMargin{218. }%
\CSLRightInline{Kosch, T. A. \emph{et al.}
\href{https://doi.org/10.1186/s12864-024-10899-7}{The {Amphibian}
{Genomics} {Consortium}: Advancing genomic and genetic resources for
amphibian research and conservation}. \emph{BMC Genomics} \textbf{25},
1025 (2024).}

\bibitem[\citeproctext]{ref-gable_state_2023}
\CSLLeftMargin{219. }%
\CSLRightInline{Gable, S. M. \emph{et al.}
\href{https://doi.org/10.3390/genes14071387}{The {State} of {Squamate}
{Genomics}: {Past}, {Present}, and {Future} of {Genome} {Research} in
the {Most} {Speciose} {Terrestrial} {Vertebrate} {Order}}. \emph{Genes}
\textbf{14}, (2023).}

\bibitem[\citeproctext]{ref-gable_genomic_2022}
\CSLLeftMargin{220. }%
\CSLRightInline{Gable, S. M., Byars, M. I., Literman, R. \& Tollis, M.
\href{https://doi.org/10.1093/sysbio/syac019}{A {Genomic} {Perspective}
on the {Evolutionary} {Diversification} of {Turtles}}. \emph{Systematic
Biology} \textbf{71}, 1331--1347 (2022).}

\bibitem[\citeproctext]{ref-card_genome_2023}
\CSLLeftMargin{221. }%
\CSLRightInline{Card, D. C. \emph{et al.}
\href{https://doi.org/10.3390/ani13030471}{Genome {Evolution} and the
{Future} of {Phylogenomics} of {Non}-{Avian} {Reptiles}}. \emph{Animals}
\textbf{13}, (2023).}

\end{CSLReferences}




\end{document}
